\documentclass[11pt,a4paper,openright,twoside]{book}

% ============ 宏包 ============
\usepackage[utf8]{inputenc}
\usepackage[T1]{fontenc}
\usepackage{lmodern}
\usepackage{xcolor}
\usepackage{ctex}
\usepackage{amsmath,amssymb,amsthm,amsfonts,mathtools}
\usepackage{mathrsfs}
\usepackage{graphicx}
\usepackage{fancyhdr}
\usepackage{titlesec}
\usepackage{enumerate}
\usepackage{imakeidx}
\usepackage{thmtools}
\usepackage{tikz-cd}
\usepackage{etoolbox}
\usepackage{hyperref}
\usepackage{cleveref}
\usepackage{chngcntr}

% 设置英文名称
\ctexset{
  contentsname      = Contents,
  listfigurename    = List of Figures,
  listtablename     = List of Tables,
  bibname           = References,
  indexname         = Index,
  appendixname      = Appendix,
}

% 页面设置
\usepackage{geometry}
\geometry{
    top=2cm,
    bottom=2cm,
    left=1.5cm,
    right=1.5cm,
    marginparwidth=2.5cm,
    marginparsep=0.5cm,
    headheight=15pt,
}

% 行间距
\linespread{1.35}

% ============ 定理环境 ============
% 原书结构: 11 个 chapter (每章开头有一段引言), 章下设全局编号的 section (§1-§33).
% 编号约定: Theorem 7.1 指 §7 的第 1 个定理 (见原书 Conventions (8)),
% 故定理计数器挂在 section 上; section 编号全局连续 (1-33, 不随 chapter 重置).
\counterwithout{section}{chapter}
\renewcommand{\thesection}{\arabic{section}}
% 原书中 Theorem/Proposition/Corollary 共用一个计数器, Lemma 单独计数(每节从 1 开始),
% Definition 与 Example 不编号或每节独立编号, Remark 不编号.
\theoremstyle{plain}
\newtheorem{theorem}{Theorem}[section]
\newtheorem*{theorem*}{Theorem}
\newtheorem{proposition}[theorem]{Proposition}
\newtheorem{corollary}[theorem]{Corollary}
\newtheorem*{corollary*}{Corollary}
\newtheorem{corollarynum}{Corollary}[section]
\renewcommand{\thecorollarynum}{\arabic{corollarynum}} % 如 Theorem 4.3 后的 Corollary 1-4
\newtheorem{lemma}{Lemma}[section]
\renewcommand{\thelemma}{\arabic{lemma}} % 原书 Lemma 按节独立编号: Lemma 1, Lemma 2
\newtheorem*{lemma*}{Lemma}
\newtheorem{exercise}{Exercise}[section]
\theoremstyle{definition}
\newtheorem*{definition}{Definition}
\newtheorem*{definition*}{Definition}
\newtheorem{example}{Example}[section]
\renewcommand{\theexample}{\arabic{example}} % 原书 Example 按节独立编号
\newtheorem*{example*}{Example}
\theoremstyle{remark}
\newtheorem*{remark}{Remark}
\newtheorem*{remarks}{Remarks}

% 证明环境
\renewcommand{\proofname}{Proof}
\renewcommand{\qedsymbol}{$\blacksquare$}

% ============ 超链接 ============
\definecolor{linkblue}{RGB}{0,0,180}
\hypersetup{
    colorlinks=true,
    linkcolor=linkblue,
    hyperindex=true,
    pdfauthor={Hideyuki Matsumura},
    pdftitle={Commutative Ring Theory},
    pdfsubject={Commutative Algebra}
}

% ============ 页眉页脚 ============
\pagestyle{fancy}
\fancyhf{}
\fancyhead[LE,RO]{\thepage}
\fancyhead[RE]{\leftmark}
\fancyhead[LO]{\emph{Commutative Ring Theory --- H. Matsumura}}
\renewcommand{\headrulewidth}{0.4pt}
\renewcommand{\footrulewidth}{0pt}

% ============ 自定义命令 ============
\newcommand{\Z}{\mathbb{Z}}
\newcommand{\Q}{\mathbb{Q}}
\newcommand{\R}{\mathbb{R}}
\newcommand{\C}{\mathbb{C}}
\newcommand{\N}{\mathbb{N}}
\newcommand{\ann}{\operatorname{ann}}
\newcommand{\Hom}{\operatorname{Hom}}
\newcommand{\End}{\operatorname{End}}
\newcommand{\Aut}{\operatorname{Aut}}
\newcommand{\rad}{\operatorname{rad}}
\newcommand{\nil}{\operatorname{nil}}
\newcommand{\Spec}{\operatorname{Spec}}
\newcommand{\Zar}{\operatorname{Zar}}
\newcommand{\Ass}{\operatorname{Ass}}
\newcommand{\Att}{\operatorname{Att}}
\newcommand{\Supp}{\operatorname{Supp}}
\newcommand{\Ker}{\operatorname{Ker}}
\newcommand{\Coker}{\operatorname{Coker}}
\newcommand{\Imf}{\operatorname{Im}}
\newcommand{\lcm}{\operatorname{l.c.m.}}
\newcommand{\hcf}{\operatorname{h.c.f.}}
\newcommand{\trdeg}{\operatorname{tr.deg}}
% 注意: 不可重定义 \ht (TeX 原语, \LaTeX 徽标依赖它), 故用 \htop 表示高度算子 ht
\newcommand{\htop}{\operatorname{ht}}
\newcommand{\coht}{\operatorname{coht}}
\newcommand{\depth}{\operatorname{depth}}
\newcommand{\grade}{\operatorname{grade}}
\newcommand{\injdim}{\operatorname{inj\,dim}}
\newcommand{\projdim}{\operatorname{proj\,dim}}
\newcommand{\codepth}{\operatorname{codepth}}
\newcommand{\rank}{\operatorname{rank}}
\newcommand{\Frac}{\operatorname{Frac}}
\newcommand{\Max}{\operatorname{Max}}
\newcommand{\Tor}{\operatorname{Tor}}
\newcommand{\Ext}{\operatorname{Ext}}
\newcommand{\gr}{\operatorname{gr}}
\newcommand{\ord}{\operatorname{ord}}
\newcommand{\gldim}{\operatorname{gl\,dim}}
\newcommand{\embdim}{\operatorname{emb\,dim}}
\newcommand{\Pit}{\operatorname{Pit}}
\newcommand{\cl}{\operatorname{cl}}
\newcommand{\divisor}{\operatorname{div}}
\newcommand{\CM}{\operatorname{CM}}
\newcommand{\Gor}{\operatorname{Gor}}
\newcommand{\Reg}{\operatorname{Reg}}
\newcommand{\Der}{\operatorname{Der}}
\newcommand{\tr}{\operatorname{tr}}
\newcommand{\chark}{\operatorname{char}}
\newcommand{\Nor}{\operatorname{Nor}}
\providecommand{\blank}{{-}}

% 索引
\makeindex

% ============ 文档开始 ============
\begin{document}
\frontmatter

% 封面
\begin{titlepage}
    \begin{center}
        \vspace*{2cm}
        {\Huge\bfseries Commutative Ring Theory\par}
        \vspace{1.5cm}
        {\Large\bfseries HIDEYUKI MATSUMURA\par}
        \vspace{0.5cm}
        {\large Department of Mathematics, Faculty of Sciences\par}
        {\large Nagoya University, Nagoya, Japan\par}
        \vspace{0.7cm}
        {\large Translated by M.\ Reid\par}
        \vfill
        {\large CAMBRIDGE UNIVERSITY PRESS\par}
        \vspace{0.5cm}
        {\large Cambridge \textbullet{} New York \textbullet{} New Rochelle\par}
        {\large Melbourne \textbullet{} Sydney\par}
    \end{center}
\end{titlepage}

\thispagestyle{empty}
\begin{center}
    \small
    \noindent\textbf{关于本重排版 / About this edition}\\[0.5em]
    本文件为《Commutative Ring Theory》一书的中文非官方 \LaTeX{} 重排版, 内容忠实于\\
    Cambridge University Press 1986 年英文版 (Miles Reid 翻译), 未作任何内容删改.\\
    重排版仅用于学习交流, 原书版权归 Hideyuki Matsumura 与 Cambridge University Press 所有.\\
    \vspace{0.5cm}
    Original edition: \textit{Commutative ring theory}, translated by M.\ Reid,\\
    Cambridge Studies in Advanced Mathematics 8, Cambridge University Press, 1986.
\end{center}

% 重排前言
\newpage
\pagestyle{plain}
\chapter*{重排前言}
\addcontentsline{toc}{chapter}{重排前言}
松村英之 (Hideyuki Matsumura) 的《交换环论》(\textit{Commutative Ring Theory}, Cambridge University Press, 1986, 由 Miles Reid 翻译) 是交换代数领域公认的经典专著. 全书以精炼而严谨的笔法, 系统地覆盖了从理想与模的基本理论、局部化、平坦性、完备化, 到维数理论、正则环、导子与光滑性等现代交换代数的核心内容, 是代数几何与代数数论研究者不可或缺的参考书.

然而, 由于原书排版年代较早, 市面流通的扫描版电子文件普遍存在字体模糊、公式不清、版面拥挤的问题, 长期阅读容易造成视觉疲劳. 为了让读者能够更舒适地研读这部经典, 我们发起了本次非官方重排项目.

本次重排工作的要点如下:
\begin{enumerate}
  \item \textbf{完全忠于原书内容}: 节 (§) 的划分、定理编号 (如 Theorem~\ref{thm:7.1} 指 §7 的第 1 个定理)、证明过程、习题安排与解答均严格遵从 1986 年英文版, 未作任何删改.
  \item \textbf{采用 \LaTeX{} 专业排版}: 数学公式清晰锐利, 正文行距与页边距经过仔细调整, 适合屏幕阅读和打印.
  \item \textbf{补充超链接}: 在 PDF 版本中, 目录、交叉引用、定理编号、参考文献与索引页码均为可点击超链接.
  \item \textbf{分章推进}: 全书共 11 章、33 节及 3 个附录, 每章开头附有原书的章节引言, 本重排版按章节逐次完成.
\end{enumerate}

\noindent\textbf{关于版权}\\
本项目严格遵循"只重排形式、不改变内容"的原则:
\begin{itemize}
  \item 原书版权归作者及原出版社所有, 本重排版并非官方版本, 也不声称拥有任何内容版权.
  \item 重排版的 \LaTeX{} 源码开源共享, 其目的仅在于提供一种现代化的排版实现, 方便数学爱好者自行打印和学习.
  \item 本项目不涉及任何商业行为. 若您需要正式出版物, 请支持正版原书.
\end{itemize}

\vspace{1cm}
\begin{flushright}
  重排者: Vespera Zephyr\\[0.3em]
  \today
\end{flushright}

% 原书前言、引言、记号
% ============ 原书前言 Preface ============
\chapter*{Preface}
\addcontentsline{toc}{chapter}{Preface}

In publishing this English edition I have tried to make a rather extensive
revision. Most of the mistakes and insufficiencies in the original edition
have, I hope, been corrected, and some theorems have been improved.
Some topics have been added in the form of Appendices to individual
sections. Only Appendices A, B and C are from the original. The final
section, \S33, of the original edition was entitled `Kunz' Theorems'
and did not substantially differ from a section in the second edition of
my previous book \textit{Commutative Algebra} (Benjamin, 2nd edn 1980) so I
have replaced it by the present \S33. The bibliography at the end of
the book has been considerably enlarged, although it is obviously
impossible to do justice to all of the ever-increasing literature.

Dr.\ Miles Reid has done excellent work of translation. He also pointed
out some errors and proposed some improvements. Through his efforts this
new edition has become, I believe, more readable than the original. To him,
and to the staff of Cambridge University Press and Kyoritsu Shuppan Co.,
Tokyo, who cooperated to make the publication of this English edition
possible, I express here my heartfelt gratitude.

\begin{flushright}
Hideyuki Matsumura\\
Nagoya
\end{flushright}

% ============ 原书引言 Introduction ============
\chapter*{Introduction}
\addcontentsline{toc}{chapter}{Introduction}

In addition to being a beautiful and deep theory in its own right,
commutative ring theory is important as a foundation for algebraic
geometry and complex analytic geometry. Let us start with a historical
survey of its development.

The most basic commutative rings are the ring $\Z$ of rational integers, and
the polynomial rings over a field. $\Z$ is a principal ideal ring, and so is too
simple to be ring-theoretically very interesting, but it was in the course of
studying its extensions, the rings of integers of algebraic number fields, that
Dedekind first introduced the notion of an ideal in the 1870s. For it was
realised that only when prime ideals are used in place of prime numbers do
we obtain the natural generalisation of the number theory of $\Z$.

Meanwhile, in the second half of the 19th century, polynomial rings
gradually came to be studied both from the point of view of algebraic
geometry and of invariant theory. In his famous papers of the 1890s on
invariants, Hilbert proved that ideals in polynomial rings are finitely
generated, as well as other fundamental theorems. After the turn of the
present century had seen the deep researches of Lasker and Macaulay on
primary decomposition of polynomial ideals came the advent of the age of
abstract algebra. A forerunner of the abstract treatment of commutative
ring theory was the Japanese Shozij\=o Sono (On congruences, I--IV, Mem.\
Coll.\ Sci.\ Kyoto, 2 (1917), 3 (1918--19)); in particular he succeeded in giving
an axiomatic characterisation of Dedekind rings. Shortly after this Emmy
Noether discovered that primary decomposition of ideals is a consequence
of the ascending chain condition (1921), and gave a different system of
axioms for Dedekind rings (1927), in work which was to have a decisive
influence on the direction of subsequent development of commutative ring
theory. The central position occupied by Noetherian rings in commutative
ring theory became evident from her work.

However, the credit for raising abstract commutative ring theory to a
substantial branch of science belongs in the first instance to Krull (1899--
1970). In the 1920s and 30s he established the dimension theory of
Noetherian rings, introduced the methods of localisation and completion,
and the notion of a regular local ring, and went beyond the framework of
Noetherian rings to create the theory of general valuation rings and Krull
rings. The contribution of Akizuki in the 1930s was also considerable; in
particular, a counter-example, which he obtained after a year's hard
struggle, of an integral domain whose integral closure is not finite as a
module was to become the model for many subsequent counter-examples.

In the 1940s Krull's theory was applied to algebraic geometry by
Chevalley and Zariski, with remarkable success. Zariski applied general
valuation theory to the resolution of singularities and the theory of
birational transformations, and used the notion of regular local ring to give
an algebraic formulation of the theory of simple (non-singular) point of a
variety. Chevalley initiated the theory of multiplicities of local rings, and
applied it to the computation of intersection multiplicities of varieties.
Meanwhile, Zariski's student I.\,S.\ Cohen proved the structure theorem for
complete local rings~\cite{Coh1}, underlining the importance of completion.

The 1950s opened with the profound work of Zariski on the problem of
whether the completion of a normal local ring remains normal (Sur la
normalit\'e analytique des vari\'et\'es normales, Ann.\ Inst.\ Fourier 2 (1950)),
taking Noetherian ring theory from general theory deeper into precise
structure theorems. Multiplicity theory was given new foundations by
Samuel and Nagata, and became one of the powerful tools in the theory of
local rings. Nagata, who was the most outstanding research worker of the
1950s, also created the theory of Hensel rings, constructed examples of non-
catenary Noetherian rings and counter-examples to Hilbert's 14th prob-
lem, and initiated the theory of Nagata rings (which he called pseudo-
geometric rings). Y.\ Mori carried out a deep study of the integral closure
of Noetherian integral domains.

However, in contrast to Nagata and Mori's work following the Krull
tradition, there was at the same time a new and completely different
movement, the introduction of homological algebra into commutative ring
theory by Auslander and Buchsbaum in the USA, Northcott and Rees in
Britain, and Serre in France, among others. In this direction, the theory of
regular sequences and depth appeared, giving a new treatment of Cohen--
Macaulay rings, and through the homological characterisation of regular
local rings there was dramatic progress in the theory of regular local rings.

The early 1960s saw the publication of Bourbaki's \textit{Alg\`ebre commutative},
which emphasised flatness, and treated primary decomposition from a new
angle. However, without doubt, the most characteristic aspect of this
decade was the activity of Grothendieck. His scheme theory created a
fusion of commutative ring theory and algebraic geometry, and opened up
ways of applying geometric methods in ring theory. His local cohomology
is an example of this kind of approach, and has become one of the
indispensable methods of modern commutative ring theory. He also
initiated the theory of Gorenstein rings. In addition, his systematic
development, in Chapter IV of EGA, of the study of formal fibres, and the
theory of excellent rings arising out of it, can be seen as a continuation and a
final conclusion of the work of Zariski and Nagata in the 1950s.

In the 1960s commutative ring theory was to receive another two
important gifts from algebraic geometry. Hironaka's great work on the
resolution of singularities~\cite{Hir1} contained an extremely original piece of
work within the ideal theory of local rings, the ring-theoretical significance
of which is gradually being understood. The theorem on resolution of
singularities has itself recently been used by Rotthaus in the study of
excellent rings. Secondly, in 1969 M.\ Artin proved his famous approxim-
ation theorem; roughly speaking, this states that if a system of simultaneous
algebraic equations over a Hensel local ring $A$ has a solution in the
completion $\hat{A}$, then there exist arbitrarily close solutions in $A$ itself. This
theorem has a wide variety of applications both in algebraic geometry and
in ring theory. A new homology theory of commutative rings constructed
by M.\ Andr\'e and Quillen is a further important achievement of the 1960s.

The 1970s was a period of vigorous research in homological directions by
many workers. Buchsbaum, Eisenbud, Northcott and others made detailed
studies of properties of complexes, while techniques discovered by Peskine
and Szpiro~\cite{PS1} and Hochster~\cite{H} made ingenious use of the Frobenius
map and the Artin approximation theorem. Cohen--Macaulay rings,
Gorenstein rings, and most recently Buchsbaum rings have been studied in
very concrete ways by Hochster, Stanley, Kei-ichi Watanabe and S.\ Goto
among others. On the other hand, classical ideal theory has shown no sign
of dying off, with Ratliff and Rotthaus obtaining extremely deep results.

To give the three top theorems of commutative ring theory in order of
importance, I have not much doubt that Krull's dimension theorem
(Theorem~\ref{thm:13.5}) has pride of place. Next perhaps is I.\,S.\ Cohen's structure
theorem for complete local rings (Theorems 28.3, 29.3 and 29.4). The fact
that a complete local ring can be expressed as a quotient of a well-
understood ring, the formal power series ring over a field or a discrete
valuation ring, is something to feel extremely grateful for. As a third, I
would give Serre's characterisation of a regular local ring (Theorem 19.2);
this grasps the essence of regular local rings, and is also an important
meeting-point of ideal theory and homological algebra.

This book is written as a genuine textbook in commutative algebra, and
is as self-contained as possible. It was also the intention to give some
thought to the applications to algebraic geometry. However, both for
reasons of space and limited ability on the part of the author, we are not
able to touch on local cohomology, or on the many subsequent results of
the cohomological work of the 1970s. There are readable accounts of these
subjects in~\cite{G6} and~\cite{H}, and it would be useful to read these after this
book.

This book was originally to have been written by my distinguished friend
Professor Masao Narita, but since his tragic early death through illness, I
have taken over from him. Professor Narita was an exact contemporary of
mine, and had been a close friend ever since we met at the age of 24. Well-
respected and popular with all, he was a man of warm character, and it was
a sad loss when he was prematurely called to a better place while still in his
forties. Believing that, had he written the book, he would have included
topics which were characteristic of him, UFDs, Picard groups, and so on,
I have used part of his lectures in \S20 as a memorial to him. I could
wish for nothing better than to present this book to Professor Narita and to
hear his criticism.

\begin{flushright}
Hideyuki Matsumura\\
Nagoya
\end{flushright}

% ============ 记号与约定 Conventions and terminology ============
\chapter*{Conventions and terminology}
\addcontentsline{toc}{chapter}{Conventions and terminology}

\begin{enumerate}
  \item[(1)] Some basic definitions are given in Appendixes A--C. The index contains
  references to all definitions, including those of the appendixes.
  \item[(2)] In this book, by a \emph{ring} we always understand a commutative ring with
  unit; ring homomorphisms $A \to B$ are assumed to take the unit element
  of $A$ into the unit element of $B$. When we say that $A$ is a subring of $B$
  it is understood that the unit elements of $A$ and $B$ coincide.
  \item[(3)] If $f\colon A \to B$ is a ring homomorphism and $J$ is an ideal of $B$, then
  $f^{-1}(J)$ is an ideal of $A$, and we denote this by $A \cap J$; if $A$ is a subring of
  $B$ and $f$ is the inclusion map then this is the same as the usual set-theoretic
  notion of intersection. In general this is not true, but confusion does not
  arise. Moreover, if $I$ is an ideal of $A$, we will write $IB$ for the ideal $f(I)B$ of $B$.
  \item[(4)] If $A$ is a ring and $a_1, \ldots, a_n$ elements of $A$, the ideal of $A$ generated by
  these is written in any of the following ways: $a_1A + a_2A + \cdots + a_nA$, $\sum a_iA$,
  $(a_1, \ldots, a_n)$ or $(a_1, \ldots, a_n)$.
  \item[(5)] The sign $\subset$ is used for inclusion of a subset, including the possibility of
  equality; in~\cite{M} the sign $\subset$ was used for this purpose. However, when we
  say that `$M_1 \subset M_2 \subset \cdots$ is an ascending chain', $M_1 \subsetneq M_2 \subsetneq \cdots$ is intended.
  \item[(6)] When we say that $R$ is a ring of characteristic $p$, or write $\operatorname{char} R = p$, we
  always mean that $p > 0$ is a prime number.
  \item[(7)] In the exercises we generally omit the instruction `prove that'. Solutions
  or hints are provided at the end of the book for most of the exercises. Many of
  the exercises are intended to supplement the material of the main text, so it
  is advisable at least to glance through them.
  \item[(8)] The numbering Theorem~\ref{thm:7.1} refers to Theorem 1 of \S\ref{sec:7}; within
  one paragraph we usually just refer to Theorem 1, omitting the section
  number.
\end{enumerate}


% 目录
\tableofcontents

% ============ 主内容 ============
% 原书结构: 11 个 chapter, 每章下设全局编号的 section (§1-§33), 章首有引言.
% 每章一个 chapterNN.tex: 章标题 + 原书引言 + 各节 \input.
\mainmatter
\pagestyle{fancy}

% ============ Chapter 1. Commutative rings and modules (原书 pp.1-20) ============
\chapter{Commutative rings and modules}\label{ch:1}

This chapter discusses the very basic definitions and results.
\S\ref{sec:1} centres around the question of the existence of prime ideals. In \S\ref{sec:2}
we treat Nakayama's lemma, modules over local rings and modules of
finite presentation; we give a complete proof, following Kaplansky, of the
fact that a projective module over a local ring is free (Theorem~\ref{thm:2.5}),
although, since we will not make any subsequent use of this in the infinitely
generated case, the reader may pass over it. In \S\ref{sec:3} we give a detailed
treatment of finiteness conditions in the form of Emmy Noether's chain
condition, discussing among other things Akizuki's theorem, I.\,S.\ Cohen's
theorem and Formanek's proof of the Eakin--Nagata theorem.

% ============ §1 Ideals (原书 pp.1-6) ============
\section{Ideals}\label{sec:1}\index{Ideals, section 1}

If $A$ is a ring and $I$ an ideal of $A$, it is often important to consider
the residue class ring $A/I$. Set $\bar{A} = A/I$, and write $f\colon A \to \bar{A}$ for the
natural map; then ideals $\bar{J}$ of $\bar{A}$ and ideals $J = f^{-1}(\bar{J})$ of $A$ containing $I$ are
in one-to-one correspondence, with $\bar{J} = J/I$ and $\bar{A}/\bar{J} \cong A/J$. Hence, when
we just want to think about ideals of $A$ containing $I$, it is convenient to
shift attention to $A/I$. (If $I'$ is any ideal of $A$ then $f(I')$ is an ideal of $\bar{A}$,
with $f^{-1}(f(I')) = I + I'$, and $\overline{f(I')} = (I + I')/I$.)

$A$ is itself an ideal of $A$, often written $(1)$ since it is generated by the
identity element $1$. An ideal distinct from $(1)$ is called a \emph{proper ideal}.\index{proper ideal}
An element $a \in A$ which has an inverse in $A$ (that is, for which there exists
$a' \in A$ with $aa' = 1$) is called a \emph{unit}\index{unit} (or \emph{invertible element}) of $A$; this holds
if and only if the principal ideal $(a)$ is equal to $(1)$. If $a$ is a unit and $x$ is
nilpotent then $a + x$ is again a unit: indeed, if $x^n = 0$ then setting $y = -a^{-1}x$,
we have $y^n = 0$; now
\[
  (1 - y)(1 + y + \cdots + y^{n-1}) = 1 - y^n = 1,
\]
so that $a + x = a(1 - y)$ has an inverse.

In a ring $A$ we are allowed to have $1 = 0$, but if this happens then it
follows that $a = 1 \cdot a = 0 \cdot a = 0$ for every $a \in A$, so that $A$ has only one
element $0$; in this case we write $A = 0$. In definitions and theorems about
rings, it may sometimes happen that the condition $A \neq 0$ is omitted even
when it is actually necessary. A ring $A$ is an \emph{integral domain}\index{integral domain} (or simply a
\emph{domain}) if $A \neq 0$, and if $A$ has no zero-divisors other than $0$. If $A$ is an
integral domain and every non-zero element of $A$ is a unit then $A$ is a
\emph{field}.\index{field} A field is characterised by the fact that it is a ring having exactly
two ideals $(0)$ and $(1)$.

An ideal which is maximal among all proper ideals is called a \emph{maximal ideal};\index{maximal ideal}
an ideal $m$ of $A$ is maximal if and only if $A/m$ is a field. Given a
proper ideal $I$, let $\mathcal{M}$ be the set of ideals containing $I$ and not containing
$1$, ordered by inclusion; then Zorn's lemma can be applied to $\mathcal{M}$. Indeed,
$I \in \mathcal{M}$ so that $\mathcal{M}$ is non-empty, and if $\mathcal{L} \subset \mathcal{M}$ is a totally ordered subset
then the union of all the ideals belonging to $\mathcal{L}$ is an ideal of $A$ and obviously
belongs to $\mathcal{M}$, so is the least upper bound of $\mathcal{L}$ in $\mathcal{M}$. Thus by Zorn's
lemma $\mathcal{M}$ has got a maximal element. This proves the following theorem.

\begin{theorem}\label{thm:1.1}
If $I$ is a proper ideal then there exists at least one maximal ideal
containing $I$.
\end{theorem}

An ideal $P$ of $A$ for which $A/P$ is an integral domain is called a \emph{prime ideal}.\index{prime ideal}
In other words, $P$ is prime if it satisfies
\[
  \text{(i) } P \neq A \quad \text{and} \quad \text{(ii) } x, y \notin P \Rightarrow xy \notin P \text{ for } x, y \in A.
\]
A field is an integral domain, so that a maximal ideal is prime.

If $I$ and $J$ are ideals and $P$ a prime ideal, then
\[
  IJ \not\subset P \iff I \not\subset P \text{ or } J \not\subset P.
\]
Indeed, taking $x \in I$ and $y \in J$ with $x, y \notin P$, we have $xy \in IJ$ but $xy \notin P$.

A subset $S$ of $A$ is \emph{multiplicative}\index{multiplicative set} if it satisfies
\[
  \text{(i) } x, y \in S \Rightarrow xy \in S, \quad \text{and} \quad \text{(ii) } 1 \in S;
\]
(here condition (ii) is not crucial: given a subset $S$ satisfying (i), there will
usually not be any essential change on replacing $S$ by $S \cup \{1\}$). If $I$ is
an ideal disjoint from $S$, then exactly as in the proof of Theorem 1 we see
that the set of ideals containing $I$ and disjoint from $S$ has a maximal
element. If $P$ is an ideal which is maximal among ideals disjoint from $S$
then $P$ is prime. For if $x \notin P$, $y \notin P$, then since $P + xA$ and $P + yA$ both
meet $S$, the product $(P + xA)(P + yA)$ also meets $S$. However,
\[
  (P + xA)(P + yA) \subset P + xyA,
\]
so that we must have $xy \notin P$. We have thus obtained the following theorem.

\begin{theorem}\label{thm:1.2}
Let $S$ be a multiplicative set and $I$ an ideal disjoint from $S$;
then there exists a prime ideal containing $I$ and disjoint from $S$.
\end{theorem}

If $I$ is an ideal of $A$ then the set of elements of $A$, some power of which
belongs to $I$, is an ideal of $A$ (for $x^n \in I$ and $y^m \in I \Rightarrow (x + y)^{n+m-1} \in I$
and $(ax)^n \in I$). This set is called the \emph{radical}\index{radical} of $I$, and is sometimes written $\sqrt{I}$:
\[
  \sqrt{I} = \{a \in A \mid a^n \in I \text{ for some } n > 0\}.
\]
If $P$ is a prime ideal containing $I$ then $x^n \in I \subset P$ implies that $x \in P$, and
hence $\sqrt{I} \subset P$; conversely, if $x \notin \sqrt{I}$ then $S = \{1, x, x^2, \ldots\}$ is a multi-
plicative set disjoint from $I$, and by the previous theorem there exists a
prime ideal containing $I$ and not containing $x$. Thus, the radical of $I$ is
the intersection of all prime ideals containing $I$:
\[
  \sqrt{I} = \bigcap_{P \supset I} P.
\]
In particular if we take $I = (0)$ then $\sqrt{(0)}$ is the set of all nilpotent elements
of $A$, and is called the \emph{nilradical}\index{nilradical} of $A$; we will write $\nil(A)$ for this. Then
$\nil(A)$ is the intersection of all the prime ideals of $A$. When $\nil(A) = 0$ we say
that $A$ is \emph{reduced}.\index{reduced ring} For any ring $A$ we write $A_{red}$ for $A/\nil(A)$; $A_{red}$ is of
course reduced.

The intersection of all maximal ideals of a ring $A$ ($\neq 0$) is called the
\emph{Jacobson radical}\index{Jacobson radical}, or simply the radical of $A$, and written $\rad(A)$. If
$x \in \rad(A)$ then for any $a \in A$, $1 + ax$ is an element of $A$ not contained in
any maximal ideal, and is therefore a unit of $A$ by Theorem 1. Conversely
if $x \in A$ has the property that $1 + Ax$ consists entirely of units of $A$ then
$x \in \rad(A)$ (prove this!).

A ring having just one maximal ideal is called a \emph{local ring}\index{local ring}, and a
(non-zero) ring having only finitely many maximal ideals a \emph{semilocal ring}.\index{semilocal ring}
We often express the fact that $A$ is a local ring with maximal ideal $m$ by
saying that $(A, m)$ is a local ring; if this happens then the field $k = A/m$ is
called the \emph{residue field}\index{residue field} of $A$. We will say that $(A, m, k)$ is a local ring to
mean that $A$ is a local ring, $m = \rad(A)$ and $k = A/m$. If $(A, m)$ is a local
ring then the elements of $A$ not contained in $m$ are units; conversely a
(non-zero) ring $A$ whose non-units form an ideal is a local ring.

In general the product $II'$ of two ideals $I$, $I'$ is contained in $I \cap I'$, but
does not necessarily coincide with it. However, if $I + I' = (1)$ (in which case
we say that $I$ and $I'$ are \emph{coprime})\index{coprime ideals}, then $II' = I \cap I'$; indeed, then
\[
  I \cap I' = (I \cap I')(I + I') \subset II' \subset I \cap I'.
\]
Moreover, if $I$ and $I'$, as well as $I$ and $I''$, are coprime, then $I$ and $I'I''$ are coprime:
\[
  (1) = (I + I')(I + I'') \subset I + I'I'' \subset (1).
\]
By induction we obtain the following theorem.

\begin{theorem}\label{thm:1.3}
If $I_1, I_2, \ldots, I_n$ are ideals which are coprime in pairs then
\[
  I_1I_2 \cdots I_n = I_1 \cap I_2 \cap \cdots \cap I_n.
\]
In particular if $A$ is a semilocal ring and $m_1, \ldots, m_n$ are all of its maximal
ideals then
\[
  m_1 \cdots m_n = m_1 \cap \cdots \cap m_n = \rad(A).
\]
\end{theorem}

Furthermore, if $I + I' = (1)$ then
\[
  A/II' \cong A/I \times A/I'.
\]
To see this it is enough to prove that the natural injective map from $A/II' = A/(I \cap I')$ to
$A/I \times A/I'$ is surjective; taking $e \in I$, $e' \in I'$ such that $e + e' = 1$, we
have $ae' + a'e \equiv a \pmod I$, $ae' + a'e \equiv a' \pmod{I'}$ for any $a, a' \in A$,
giving the surjectivity. By induction we get the following theorem.

\begin{theorem}\label{thm:1.4}
If $I_1, \ldots, I_n$ are ideals which are coprime in pairs then
\[
  A/I_1 \cdots I_n \cong A/I_1 \times \cdots \times A/I_n.
\]
\end{theorem}

\begin{example}\label{ex:1.1}
Let $A$ be a ring, and consider the ring $A[[X]]$ of formal power
series over $A$. A power series $f = a_0 + a_1X + a_2X^2 + \cdots$ with $a_i \in A$ is a
unit of $A[[X]]$ if and only if $a_0$ is a unit of $A$. Indeed, if there exists an
inverse $f^{-1} = b_0 + b_1X + \cdots$ then $a_0b_0 = 1$; and conversely if $a_0$ is a unit of $A$, then
\[
  ff^{-1} = a_0b_0 + (a_0b_1 + a_1b_0)X + (a_0b_2 + a_1b_1 + a_2b_0)X^2 + \cdots = 1
\]
can be solved for $b_0, b_1, \ldots$: we just find $b_0, b_1, \ldots$ successively from
$a_0b_0 = 1$, $a_0b_1 + a_1b_0 = 0, \ldots$

Since the formal power series ring in several variables $A[[X_1, \ldots, X_n]]$ can be thought of as
$(A[[X_1, \ldots, X_{n-1}]])[[X_n]]$, here also $f = a_0 + \sum a_iX_i + \sum a_{ij}X_iX_j + \cdots$ is a unit
if and only if the constant term $a_0$ is a unit of $A$; from this we see that
if $g \in (X_1, \ldots, X_n)$ then $1 + gh$ is a unit for any power series $h$, so that
$g \in \rad(A[[X_1, \ldots, X_n]])$, and hence
\[
  (X_1, \ldots, X_n) \subset \rad(A[[X_1, \ldots, X_n]]).
\]
If $k$ is a field then $k[[X_1, \ldots, X_n]]$ is a local ring with maximal ideal
$(X_1, \ldots, X_n)$. If $A$ is any ring and we set $B = A[[X_1, \ldots, X_n]]$, then since
any maximal ideal of $B$ contains $(X_1, \ldots, X_n)$, it corresponds to a maximal
ideal of $B/(X_1, \ldots, X_n) \cong A$, and so is of the form $mB + (X_1, \ldots, X_n)$, where
$m$ is a maximal ideal of $A$. If we write $m$ for this then $m \cap A = m$.

By contrast the case of polynomial rings is quite complicated; here it
is just not true that a maximal ideal of $A[X]$ must contain $X$. For example,
$X - 1$ is a non-unit of $A[X]$, and so there exists a maximal ideal $m$
containing it, and $X \notin m$. Also, if $m$ is a maximal ideal of $A[X]$, it does
not necessarily follow that $m \cap A$ is a maximal ideal of $A$.

If $A$ is an integral domain then so are both $A[X]$ and $A[[X]]$: if
$f = a_rX^r + a_{r+1}X^{r+1} + \cdots$ and $g = b_sX^s + b_{s+1}X^{s+1} + \cdots$ with $a_r \neq 0$,
$b_s \neq 0$ then $fg = a_rb_sX^{r+s} + \cdots \neq 0$. If $I$ is an ideal of $A$ we write $I[X]$ or
$I[[X]]$ for the set of polynomials or power series with coefficients in $I$;
these are ideals of $A[X]$ or $A[[X]]$, the kernels of the homomorphisms
\[
  A[X] \to (A/I)[X] \quad \text{or} \quad A[[X]] \to (A/I)[[X]]
\]
obtained by reducing coefficients modulo $I$. Hence
\[
  \sqrt{I[X]} = (\sqrt{I})[X], \quad \text{and} \quad A[X]/I[X] \cong (A/I)[X];
\]
in particular if $P$ is a prime ideal then $P[X]$ and $P[[X]]$ are prime ideals
of $A[X]$ and $A[[X]]$, respectively.

If $I$ is finitely generated, that is $I = a_1A + \cdots + a_lA$, then $I[[X]] =
a_1A[[X]] + \cdots + a_lA[[X]] = I \cdot A[[X]]$; however, if $I$ is not finitely generated
then $I[[X]]$ is bigger than $I \cdot A[[X]]$. In the polynomial ring this distinction
does not arise, and we always have $I[X] = I \cdot A[X]$.
\end{example}

\begin{example}\label{ex:1.2}
For a ring $A$ and $a, b \in A$, we have $aA \subset bA$ if and only if $a$
is divisible by $b$, that is $a = bc$ for some $c \in A$. We assume that $A$ is an
integral domain in what follows. An element $a \in A$ is said to be \emph{irreducible}\index{irreducible element}
if $a$ is not a unit of $A$ and satisfies the condition
\[
  a = bc \Rightarrow b \text{ or } c \text{ is a unit of } A.
\]
This is equivalent to saying that $aA$ is maximal among proper principal ideals. If $aA$ is a
prime ideal then $a$ is said to be \emph{prime}.\index{prime element} As one sees easily,
a prime element is irreducible, but the converse does not always hold.

Suppose that an element $a$ has two expressions as products of prime elements:
\[
  a = p_1p_2 \cdots p_n = p_1' \cdots p_m', \quad \text{with } p_i \text{ and } p_i' \text{ prime}.
\]
Then $n = m$, and after a suitable reordering of the $p_i'$ we have $p_iA = p_i'A$;
for $p_1 \cdots p_n$ is divisible by $p_1'$, and so one of the factors, say $p_i$, is
divisible by $p_1'$. Now since both $p_1'$ and $p_i$ are irreducible, $p_1'A = p_iA$,
hence $p_i = up_1'$, with $u$ a unit, and $p_2 \cdots p_n = up_2' \cdots p_m'$. We can replace
$p_i$ by $up_i$, and induction on $n$ completes the proof. In this sense,
factorisation into prime elements (whenever possible) is unique.

An integral domain in which any element which is neither $0$ nor a unit
can be expressed as a product of prime elements is called a \emph{unique factorisation
domain}\index{unique factorisation domain} (abbreviated to UFD), or a \emph{factorial ring}. It is well
known that a principal ideal domain, that is an integral domain in which
every ideal is principal, is a UFD (see Ex.~\ref{ex:1.4-ex}). If $A$ is a principal ideal
domain then the prime ideals are of the form $(0)$ or $pA$ with $p$ a prime
element, and the latter are maximal ideals.

If $k$ is a field then $k[X_1, \ldots, X_n]$ is a UFD, as is well-known (see Ex. 20.2).
If $f(X_1, \ldots, X_n)$ is an irreducible polynomial then $(f)$ is a prime ideal,
but is not maximal if $n > 1$ (see \S\ref{sec:5}).

$\Z[\sqrt{-5}]$ is not a UFD; indeed if $z = n + m\sqrt{-5}$ with $n, m \in \Z$ then
$z\bar{z} = n^2 + 5m^2$, and since $2 = n^2 + 5m^2$ has no integer solutions it
follows that $2$ is an irreducible element of $\Z[\sqrt{-5}]$, but we see from
$2 \cdot 3 = (1 + \sqrt{-5})(1 - \sqrt{-5})$ that $2$ is not a prime element. We write
\[
  A = \Z[\sqrt{-5}] = \Z[X]/(X^2 + 5);
\]
then setting $k = \Z/2\Z$ we have
\[
  A/2A = \Z[X]/(2, X^2 + 5) = k[X]/(X^2 - 1) = k[X]/(X - 1)^2.
\]
Then $P = (2, 1 - \sqrt{-5})$ is a maximal ideal of $A$ containing $2$.
\end{example}

% ============ 习题 ============
\section*{Exercises to \S\ref{sec:1}}
\emph{Prove the following propositions.}

\begin{exercise}\label{ex:1.1-ex}
Let $A$ be a ring, and $I \subset \nil(A)$ an ideal made up of nilpotent elements;
if $a \in A$ maps to a unit of $A/I$ then $a$ is a unit of $A$.
\end{exercise}

\begin{exercise}\label{ex:1.2-ex}
Let $A_1, \ldots, A_n$ be rings; then the prime ideals of $A_1 \times \cdots \times A_n$ are of the form
\[
  A_1 \times \cdots \times A_{i-1} \times P_i \times A_{i+1} \times \cdots \times A_n,
\]
where $P_i$ is a prime ideal of $A_i$.
\end{exercise}

\begin{exercise}\label{ex:1.3-ex}
Let $A$ and $B$ be rings, and $f\colon A \to B$ a surjective homomorphism.
\begin{enumerate}[(a)]
  \item Prove that $f(\rad A) \subset \rad B$, and construct an example where the
        inclusion is strict.
  \item Prove that if $A$ is a semilocal ring then $f(\rad A) = \rad B$.
\end{enumerate}
\end{exercise}

\begin{exercise}\label{ex:1.4-ex}
Let $A$ be an integral domain. Then $A$ is a UFD if and only if every
irreducible element is prime and the principal ideals of $A$ satisfy the
ascending chain condition. (Equivalently, every non-empty family of
principal ideals has a maximal element.)
\end{exercise}

\begin{exercise}\label{ex:1.5-ex}
Let $\{P_\lambda\}_{\lambda \in \Lambda}$ be a non-empty family of prime ideals, and suppose that the $P_\lambda$
are totally ordered by inclusion; then $\bigcap P_\lambda$ is a prime ideal. Also, if $I$ is
any proper ideal, the set of prime ideals containing $I$ has a minimal element.
\end{exercise}

\begin{exercise}\label{ex:1.6-ex}
Let $A$ be a ring, $I, P_1, \ldots, P_n$ ideals of $A$, and suppose that $P_1, \ldots, P_n$ are
prime, and that $I$ is not contained in any of the $P_i$; then there exists an
element $x \in I$ not contained in any $P_i$.
\end{exercise}

% ============ §2 Modules (原书 pp.6-14) ============
\section{Modules}\label{sec:2}\index{Modules, section 2}

Let $A$ be a ring and $M$ an $A$-module. Given submodules $N$, $N'$
of $M$, the set $\{a \in A \mid aN' \subset N\}$ is an ideal of $A$, which we write $N:N'$
or $(N:N')_A$. Similarly, if $I \subset A$ is an ideal then $\{x \in M \mid Ix \subset N\}$ is a
submodule of $M$, which we write $N:I$ or $(N:I)_A$. For $a \in A$ we define
$N:a$ similarly. The ideal $0:M$ is called the \emph{annihilator}\index{annihilator} of $M$, and written
$\ann(M)$. We can consider $M$ as a module over $A/\ann(M)$. If $\ann(M) = 0$
we say that $M$ is a \emph{faithful} $A$-module.\index{faithful module} For $x \in M$ we write
$\ann(x) = \{a \in A \mid ax = 0\}$.

If $M$ and $M'$ are $A$-modules, the set of $A$-linear maps from $M$ to $M'$
is written $\Hom_A(M, M')$. This becomes an $A$-module if we define the sum
$f + g$ and the scalar product $af$ by
\[
  (f + g)(x) = f(x) + g(x), \qquad (af)(x) = af(x);
\]
(the fact that $af$ is $A$-linear depends on $A$ being commutative).

To say that $M$ is an $A$-module is to say that $M$ is an Abelian group
under addition, and that a scalar product $ax$ is defined for $a \in A$ and
$x \in M$ such that the following hold:
\[
  (*) \quad a(x + y) = ax + ay, \quad (ab)x = a(bx), \quad (a + b)x = ax + bx, \quad 1x = x;
\]
for fixed $a \in A$ the map $x \mapsto ax$ is an endomorphism of $M$ as an additive
group. Let $E$ be the set of endomorphisms of the additive group $M$;
defining the sum and product of $\lambda, \mu \in E$ by
\[
  (\lambda + \mu)(x) = \lambda(x) + \mu(x), \qquad (\lambda\mu)(x) = \lambda(\mu(x))
\]
makes $E$ into a ring (in general non-commutative), and giving $M$ an
$A$-module structure is the same thing as giving a homomorphism
$A \to E$. Indeed, if we write $a_L$ for the element of $E$ defined by $x \mapsto ax$ then
$(*)$ becomes
\[
  (ab)_L = a_Lb_L, \qquad (a + b)_L = a_L + b_L, \qquad (1)_L = 1_L.
\]
We can express the fact that $\varphi\colon M \to M$ is $A$-linear by saying that
$\varphi \in E$ and that $\varphi$ commutes with $a_L$ for $a \in A$, that is $a_L\varphi = \varphi a_L$. Since
$A$ is commutative, $a_L$ is itself an $A$-linear map of $M$ for $a \in A$. We normally
write simply $a\colon M \to M$ for the map $a_L$.

If $M$ is a $B$-module and $f\colon A \to B$ a ring homomorphism, then
we can make $M$ into an $A$-module by defining $a \cdot x = f(a) \cdot x$ for $a \in A$ and
$x \in M$. This is the $A$-module structure defined by the composite of
$f\colon A \to B$ with $B \to E$, where $E$ is the endomorphism ring of the additive
group of $M$, and $B \to E$ is the map defining the $B$-module structure of $M$.

If $M$ is finitely generated as an $A$-module we say simply that $M$ is a
\emph{finite $A$-module}\index{finite module}, or is \emph{finite over $A$}. A standard technique applicable to
finite $A$-modules is the `determinant trick', one form of which is as follows
(taken from Atiyah and Macdonald~\cite{AM}).

\begin{theorem}\label{thm:2.1}
Suppose that $M$ is an $A$-module generated by $n$ elements,
and that $\varphi \in \Hom_A(M, M)$; let $I$ be an ideal of $A$ such that $\varphi(M) \subset IM$.
Then there is a relation of the form
\[
  (**) \quad \varphi^n + a_1\varphi^{n-1} + \cdots + a_{n-1}\varphi + a_n = 0,
\]
with $a_i \in I^i$ for $1 \leq i \leq n$ (where both sides are considered as endomorphisms
of $M$).
\end{theorem}

\begin{proof}
Let $M = A\omega_1 + \cdots + A\omega_n$; by the assumption $\varphi(M) \subset IM$ there exist
$a_{ij} \in I$ such that $\varphi(\omega_i) = \sum_{j=1}^{n} a_{ij}\omega_j$. This can be rewritten
\[
  \sum_{j=1}^{n} (\varphi\delta_{ij} - a_{ij})\omega_j = 0 \quad (\text{for } 1 \leq i \leq n),
\]
where $\delta_{ij}$ is the Kronecker symbol. The coefficients of this system of linear
equations can be viewed as a square matrix $(\varphi\delta_{ij} - a_{ij})$ of elements of $A'[\varphi]$,
the commutative subring of the endomorphism ring $E$ of $M$ generated by
the image $A'$ of $A$ together with $\varphi$; let $b_{ij}$ denote its $(i,j)$th cofactor, and
$d$ its determinant. By multiplying the above equation through by $b_{ik}$ and
summing over $i$, we get $d\omega_k = 0$ for $1 \leq k \leq n$. Hence $d \cdot M = 0$, so that $d = 0$
as an element of $E$. Expanding the determinant $d$ gives a relation of the
form $(**)$.
\end{proof}

\begin{remark}
As one sees from the proof, the left-hand side of $(**)$ is the
characteristic polynomial of $(a_{ij})$,
\[
  f(X) = \det(X\delta_{ij} - a_{ij})
\]
with $\varphi$ substituted for $X$. If $M$ is the free $A$-module with basis
$\omega_1, \ldots, \omega_n$, and $I = A$, the above result is nothing other than the classical
Cayley--Hamilton theorem: let $f(X)$ be the characteristic polynomial of
the square matrix $\varphi = (a_{ij})$; then $f(\varphi) = 0$.
\end{remark}

\begin{theorem}[NAK]\label{thm:2.2}
Let $M$ be a finite $A$-module and $I$ an ideal of $A$. If
$M = IM$ then there exists $a \in A$ such that $aM = 0$ and $a \equiv 1 \pmod I$. If in
addition $I \subset \rad(A)$ then $M = 0$.
\end{theorem}

\begin{proof}
Setting $\varphi = 1$ in the previous theorem gives the relation
$a = 1 + a_1 + \cdots + a_n = 0$ as endomorphisms of $M$, that is $aM = 0$, and
$a \equiv 1 \pmod I$. If $I \subset \rad(A)$ then $a$ is a unit of $A$, so that on multiplying
both sides of $aM = 0$ by $a^{-1}$ we get $M = 0$.
\end{proof}

\begin{remark}
This theorem is usually referred to as Nakayama's lemma, but
the late Professor Nakayama maintained that it should be referred to as
a theorem of Krull and Azumaya; it is in fact difficult to determine which
of these three first had the result in the case of commutative rings, so we
refer to it as NAK in this book. Of course, this result can easily be proved
without using determinants, by induction on the number of generators
of $M$.
\end{remark}

\begin{corollary*}
Let $A$ be a ring and $I$ an ideal contained in $\rad(A)$. Suppose
that $M$ is an $A$-module and $N \subset M$ a submodule such that $M/N$ is finite
over $A$. Then $M = N + IM$ implies $M = N$.
\end{corollary*}

\begin{proof}
Setting $\bar{M} = M/N$ we have $\bar{M} = I\bar{M}$, so that, by the
theorem, $\bar{M} = 0$.
\end{proof}

If $W$ is a set of generators of an $A$-module $M$ which is minimal, in the
sense that any proper subset of $W$ does not generate $M$, then $W$ is said
to be a \emph{minimal basis}\index{minimal basis} of $M$. Two minimal bases do not necessarily have
the same number of elements; for example, when $M = A$, if $x$ and $y$ are
non-units of $A$ such that $x + y = 1$ then both $\{1\}$ and $\{x, y\}$ are
minimal bases of $A$. However, if $A$ is a local ring then the situation is clear:

\begin{theorem}\label{thm:2.3}
Let $(A, m, k)$ be a local ring and $M$ a finite $A$-module; set
$\bar{M} = M/mM$. Now $\bar{M}$ is a finite-dimensional vector space over $k$, and we
write $n$ for its dimension. Then:
\begin{enumerate}[(i)]
  \item If we take a basis $\{\bar{\omega}_1, \ldots, \bar{\omega}_n\}$ for $\bar{M}$ over $k$, and choose an
        inverse image $\omega_i \in M$ of each $\bar{\omega}_i$ then $\{\omega_1, \ldots, \omega_n\}$ is a minimal basis
        of $M$;
  \item conversely every minimal basis of $M$ is obtained in this way, and
        so has $n$ elements.
  \item If $\{\omega_1, \ldots, \omega_n\}$ and $\{\omega_1', \ldots, \omega_n'\}$ are both minimal bases of $M$, and
        $\omega_i = \sum_j a_{ij}\omega_j'$ with $a_{ij} \in A$ then $\det(a_{ij})$ is a unit of $A$, so that $(a_{ij})$ is an
        invertible matrix.
\end{enumerate}
\end{theorem}

\begin{proof}
(i) $M = \sum A\omega_i + mM$, and $M$ is finitely generated (hence also
$M/\sum A\omega_i$), so that by the above corollary $M = \sum A\omega_i$. If $\{\omega_1, \ldots, \omega_n\}$ is
not minimal, so that, for example, $\{\omega_2, \ldots, \omega_n\}$ already generates $M$ then
$\{\bar{\omega}_2, \ldots, \bar{\omega}_n\}$ generates $\bar{M}$, which is a contradiction. Hence $\{\omega_1, \ldots, \omega_n\}$
is a minimal basis.

(ii) If $\{\omega_1, \ldots, \omega_n\}$ is a minimal basis of $M$ and we set $\bar{\omega}_i$ for the image
of $\omega_i$ in $\bar{M}$, then $\bar{\omega}_1, \ldots, \bar{\omega}_n$ generate $\bar{M}$, and are linearly independent
over $k$; indeed, otherwise some proper subset of $\{\bar{\omega}_1, \ldots, \bar{\omega}_n\}$ would be
a basis of $\bar{M}$, and then by (i) a proper subset of $\{\omega_1, \ldots, \omega_n\}$ would
generate $M$, which is a contradiction.

(iii) Write $\bar{a}_{ij}$ for the image in $k$ of $a_{ij}$, so that $\bar{\omega}_i = \sum \bar{a}_{ij}\bar{\omega}_j$ holds
in $\bar{M}$. Since $(\bar{a}_{ij})$ is the matrix transforming one basis of the vector space
$\bar{M}$ into another, its determinant is non-zero. Since $\det(a_{ij}) \bmod m =
\det(\bar{a}_{ij}) \neq 0$ it follows that $\det(a_{ij})$ is a unit of $A$. By Cramer's formula
the inverse matrix of $(a_{ij})$ exists as a matrix with entries in $A$.
\end{proof}

We give another interesting application of NAK, the proof of which is
due to Vasconcelos~\cite{Vas2}.

\begin{theorem}\label{thm:2.4}
Let $A$ be a ring and $M$ a finite $A$-module. If $f\colon M \to M$
is an $A$-linear map and $f$ is surjective then $f$ is also injective, and is thus
an automorphism of $M$.
\end{theorem}

\begin{proof}
Since $f$ commutes with scalar multiplication by elements of $A$, we
can view $M$ as an $A[X]$-module by setting $X \cdot m = f(m)$ for $m \in M$. Then by
assumption $XM = M$, so that by NAK there exists $Y \in A[X]$ such that
$(1 + XY)M = 0$. Now for $u \in \Ker(f)$ we have $0 = (1 + XY)(u) = u + Yf(u) = u$, so
that $f$ is injective.
\end{proof}

\begin{theorem}\label{thm:2.5}
Let $(A, m)$ be a local ring; then a projective module over $A$
is free (for the definition of projective module, see Appendix B, p.~277).
\end{theorem}

\begin{proof}
\renewcommand{\qedsymbol}{}% 原书此证明延续至 Lemma 1/2 及文末, 中间无 QED
This is easy when $M$ is finite: choose a minimal basis $\omega_1, \ldots, \omega_n$ of $M$
and define a surjective map $\varphi\colon F \to M$ from the free module
$F = Ae_1 \oplus \cdots \oplus Ae_n$ to $M$ by $\varphi(\sum a_ie_i) = \sum a_i\omega_i$. If we set $K = \Ker(\varphi)$ then, from
the minimal basis property,
\[
  \sum a_i\omega_i = 0 \Rightarrow a_i \in m \text{ for all } i.
\]
Thus $K \subset mF$. Because $M$ is projective, there exists $\psi\colon M \to F$ such that
$F = \psi(M) \oplus K$, and it follows that $K = mK$. On the other hand, $K$ is a
quotient of $F$, therefore finite over $A$, so that $K = 0$ by NAK and $F \cong M$.

The result was proved by Kaplansky~\cite{Kap2} without the assumption that
$M$ is finite. He proves first of all the following lemma, which holds for
any ring (possibly non-commutative).
\end{proof}

\begin{lemma}\label{lem:2.1}
Let $R$ be any ring, and $F$ an $R$-module which is a direct sum of
countably generated submodules; if $M$ is an arbitrary direct summand of $F$
then $M$ is also a direct sum of countably generated submodules.
\end{lemma}

\begin{proof}[Proof of Lemma 1]
\renewcommand{\qedsymbol}{}% 同上: 原书无中间 QED
Suppose that $F = M \oplus N$, and that $F = \bigoplus_{\lambda \in \Lambda} E_\lambda$, where
each $E_\lambda$ is countably generated. By transfinite induction, we construct a
well-ordered family $\{F_\alpha\}$ of submodules of $F$ with the following properties:
\begin{enumerate}[(i)]
  \item if $\alpha < \beta$ then $F_\alpha \subset F_\beta$,
  \item $F = \bigcup_\alpha F_\alpha$,
  \item if $\alpha$ is a limiting ordinal then $F_\alpha = \bigcup_{\beta < \alpha} F_\beta$,
  \item $F_{\alpha+1}/F_\alpha$ is countably generated,
  \item $F_\alpha = M_\alpha \oplus N_\alpha$, where $M_\alpha = M \cap F_\alpha$, $N_\alpha = N \cap F_\alpha$,
  \item each $F_\alpha$ is a direct sum of $E_\lambda$ taken over a suitable subset of $\Lambda$.
\end{enumerate}
We now construct such a family $\{F_\alpha\}$. Firstly, set $F_0 = (0)$. For an
ordinal $\alpha$, assume that $F_\beta$ has been defined for all ordinals $\beta < \alpha$. If $\alpha$ is
a limiting ordinal, set $F_\alpha = \bigcup_{\beta<\alpha} F_\beta$. If $\alpha$ is of the form $\alpha = \beta + 1$, let $Q_1$
be any one of the $E_\lambda$ not contained in $F_\beta$ (if $F_\beta = F$ then the construction
stops at $F_\beta$). Take a set $x_{11}, x_{12}, \ldots$ of generators of $Q_1$, and decompose $x_{11}$
into its $M$- and $N$-components; now let $Q_2$ be the direct sum of the
finitely many $E_\lambda$ which are necessary to write each of these two components
in the decomposition $F = \bigoplus E_\lambda$, and let $x_{21}, x_{22}, \ldots$ be generators of $Q_2$.
Next decompose $x_{12}$ into its $M$- and $N$-components, let $Q_3$ be the direct
sum of the finitely many $E_\lambda$ needed to write these components, and let
$x_{31}, x_{32}, \ldots$ be generators of $Q_3$. Then carry out the same procedure with
$x_{21}$, getting $x_{41}, x_{42}, \ldots$, then do the same for $x_{13}$. Carrying out the same
procedure for each of the $x_{ij}$ in the order $x_{11}, x_{12}, x_{21}, x_{13}, x_{22}, x_{31}, \ldots$ we get
countably many elements $x_{ij}$. We let $F_\alpha$ be the submodule of $F$ generated
by $F_\beta$ and the $x_{ij}$, and this satisfies all our requirements. This gives the
family $\{F_\alpha\}$.

Now $M = \bigcup_\alpha M_\alpha$, with each $M_\alpha$ a direct summand of $F_\alpha$, and
$M_{\alpha+1} \supset M_\alpha$, so that $M_\alpha$ is also a direct summand of $M_{\alpha+1}$. Moreover,
\[
  F_{\alpha+1}/F_\alpha = (M_{\alpha+1}/M_\alpha) \oplus (N_{\alpha+1}/N_\alpha),
\]
and hence $M_{\alpha+1}/M_\alpha$ is countably generated. Thus we can write
\[
  M_{\alpha+1} = M_\alpha \oplus M'_{\alpha+1}, \quad \text{with } M'_{\alpha+1} \text{ countably generated}.
\]
When $\alpha$ is a limit ordinal, since $M_\alpha = \bigcup_{\beta<\alpha} M_\beta$, we set $M'_\alpha = 0$. Then finally
we can write
\[
  M = \bigoplus M'_\alpha, \quad \text{with } M'_\alpha \text{ (at most) countably generated.}
\]
\end{proof}

Of course a free module satisfies the assumption of Lemma~\hyperref[lem:2.1]{1}, so that,
in particular, we see that any projective module is a direct sum of countably
generated projective modules. Thus in the proof of Theorem~\ref{thm:2.5} we can
assume that $M$ is countably generated.

\begin{lemma}\label{lem:2.2}
Let $M$ be a projective module over a local ring $A$, and $x \in M$.
Then there exists a direct summand of $M$ containing $x$ which is a free
module.
\end{lemma}

\begin{proof}[Proof of Lemma 2]
\renewcommand{\qedsymbol}{}% 同上: 原书无中间 QED
We write $M$ as a direct summand of a free module
$F = M \oplus N$. Choose a basis $B = \{u_i\}_{i \in I}$ of $F$ such that the given element
$x$ has the minimum possible number of non-zero coordinates when
expressed in this basis. Then if $x = u_1a_1 + \cdots + u_na_n$ with $0 \neq a_i \in A$, we have
\[
  a_i \notin \sum_{j \neq i} Aa_j \quad \text{for } i = 1, 2, \ldots, n;
\]
indeed, if, say, $a_1 = \sum_{i=2}^n b_ia_i$ then $x = \sum_{i=2}^n (u_i + u_1b_i)a_i$, which contradicts
the choice of $B$. Now set $u_i = y_i + z_i$ with $y_i \in M$ and $z_i \in N$; then
\[
  x = \sum a_iu_i = \sum a_iy_i.
\]
If we write $y_i = \sum_{j=1}^n c_{ij}u_j + t_i$, with $t_i$ linear combinations of elements of
$B$ other than $u_1, \ldots, u_n$, we get relations $a_i = \sum_{j=1}^n a_jc_{ji}$, and, hence, in view
of what we have seen above, we must have
\[
  1 - c_{ii} \in m \quad \text{and} \quad c_{ij} \in m \text{ for } i \neq j.
\]
It follows that the matrix $(c_{ij})$ has an inverse (this can be seen from the
fact that the determinant is $\equiv 1 \pmod m$, or by elimination). Thus replacing
$u_1, \ldots, u_n$ by $y_1, \ldots, y_n$ in $B$, we still have a basis of $F$. Hence, $F_1 = \sum y_iA$ is
a direct summand of $F$, and hence also of $M$, and satisfies all the
requirements of Lemma~\hyperref[lem:2.2]{2}.
\end{proof}

To prove the theorem, let $M$ be a countably generated projective module,
$M = A\omega_1 + A\omega_2 + \cdots$. By Lemma~\hyperref[lem:2.2]{2}, there exists a free module $F_1$ such that
$\omega_1 \in F_1$, and $M = F_1 \oplus M_1$, where $M_1$ is a projective module. Let $\omega_2'$
be the $M_1$-component of $\omega_2$ in the decomposition $M = F_1 \oplus M_1$, and take
a free module $F_2$ such that $\omega_2' \in F_2$ and $M_1 = F_2 \oplus M_2$, where $M_2$ is a
projective module. Let $\omega_3'$ be the $M_2$-component of $\omega_3$ in $M = F_1 \oplus F_2 \oplus M_2$;
proceeding in the same way, we get
\[
  M = F_1 \oplus F_2 \oplus \cdots,
\]
so that $M$ is a free module. $\blacksquare$

We say that an $A$-module $M \neq 0$ is a \emph{simple module}\index{simple module} if it has no
submodules other than $0$ and $M$ itself. For any $0 \neq \omega \in M$, we then have
$M = A\omega$. Now $A\omega \cong A/\ann(\omega)$, but in order for this to be simple, $\ann(\omega)$
must be a maximal ideal of $A$. Hence, any simple $A$-module is isomorphic
to $A/m$ with $m$ a maximal ideal, and conversely an $A$-module of this form
is simple. If $M$ is an $A$-module, a chain
\[
  M = M_0 \supset M_1 \supset \cdots \supset M_r = 0
\]
of submodules of $M$ is called a \emph{composition series}\index{composition series} of $M$ if every $M_i/M_{i+1}$ is
simple; $r$ is called the \emph{length} of the composition series. If a composition series
of $M$ exists, its length is an invariant of $M$ independent of the choice of
composition series. More precisely, if $M$ has a composition series of length
$r$, and if $M_1 \supset N_1 \supset \cdots \supset N_s$ is a strictly descending chain of sub-
modules, then we have $s < r$. This invariance corresponds to part of the
basic Jordan--H\"older theorem in group theory, but it can easily be proved
on its own by induction, and the reader might like to do this as an exercise.
The length of a composition series of $M$ is called the \emph{length of $M$}\index{length of a module}, and written
$l(M)$; if $M$ does not have a composition series we set $l(M) = \infty$. A necessary
and sufficient condition for the existence of a composition series of $M$ is that
the submodules of $M$ should satisfy both the ascending and descending
chain conditions (for which see \S\ref{sec:3}). In general, if $N \subset M$ is a submodule,
we have
\[
  l(M) = l(N) + l(M/N).
\]
If $0 \to M_1 \to M_2 \to \cdots \to M_r \to 0$ is an exact sequence of $A$-modules and
each $M_i$ has finite length then
\[
  \sum_{i=1}^{r} (-1)^i l(M_i) = 0.
\]

If $m$ is a maximal ideal of $A$ and is finitely generated over $A$ then
$l(A/m^\nu) < \infty$. In fact,
\[
  l(A/m^\nu) = l(A/m) + l(m/m^2) + \cdots + l(m^{\nu-1}/m^\nu);
\]
now each $m^i/m^{i+1}$ is a finite-dimensional vector space over the field
$k = A/m$, and since its $A$-submodules are the same thing as its vector
subspaces, $l(m^i/m^{i+1})$ is equal to the dimension of $m^i/m^{i+1}$ as $k$-vector
space. (This shows that $A/m^\nu$ is an Artinian ring, see \S\ref{sec:3}.)

Considering $l(A/m^\nu)$ for all $\nu$, we get a function of $\nu$ which is intimately
related to the ring structure of $A$, and which also plays a role in the resolution
of singularities in algebraic and complex analytic geometry; this is studied in
Chapter~\ref{ch:5}.

We say that an $A$-module $M$ is \emph{of finite presentation}\index{module of finite presentation} if there exists an
exact sequence of the form
\[
  A^p \to A^q \to M \to 0.
\]
This means that $M$ can be generated by $q$ elements $\omega_1, \ldots, \omega_q$ in such a way
that the module $R = \{(a_1, \ldots, a_q) \in A^q \mid \sum a_i\omega_i = 0\}$ of linear relations
holding between the $\omega_i$ can be generated by $p$ elements.

\begin{theorem}\label{thm:2.6}
Let $A$ be a ring, and suppose that $M$ is an $A$-module of finite
presentation. If
\[
  0 \to K \to N \to M \to 0
\]
is an exact sequence and $N$ is finitely generated then so is $K$.
\end{theorem}

\begin{proof}
By assumption there exists an exact sequence of the form
$L_2 \xrightarrow{\alpha} L_1 \xrightarrow{\varphi} M \to 0$, where $L_1$ and $L_2$ are free modules of finite rank.
From this we get the following commutative diagram (see Appendix B):
\[
  \begin{tikzcd}[column sep=large, row sep=large]
    & L_2 \arrow[r, "\alpha"] \arrow[d, "\beta"] & L_1 \arrow[r, "\varphi"] \arrow[d, "u"] & M \arrow[r] \arrow[d, equal] & 0 \\
    0 \arrow[r] & K \arrow[r, "\gamma"] & N \arrow[r, "f"] & M \arrow[r] & 0
  \end{tikzcd}
\]
If we write $N = At_1 + \cdots + At_q$, then there exist $u_i \in L_1$ such that
$\varphi(u_i) = f(t_i)$. Set $t_i' = t_i - u(u_i)$; then $f(t_i') = 0$, so that we can write
$t_i' = \gamma(\eta_i)$ with $\eta_i \in K$. Let us now prove that
\[
  K = \beta(L_2) + A\eta_1 + \cdots + A\eta_q.
\]
For any $\eta \in K$, set $\psi(\eta) = \sum a_i t_i$; then
\[
  \psi\left(\eta - \sum a_i t_i'\right) = \sum a_i (t_i - t_i') = u\left(\sum a_i u_i\right),
\]
and since $0 = f\left(\sum a_i t_i\right) = \varphi\left(\sum a_i u_i\right)$, we can write $\sum a_i u_i = \alpha(w)$
with $w \in L_2$. Now
\[
  \eta - \sum a_i t_i' = u\left(\sum a_i u_i\right) = \beta(w),
\]
so that $\eta = \beta(w) + \sum a_i\eta_i$, and this proves our assertion.
\end{proof}

% ============ 习题 ============
\section*{Exercises to \S\ref{sec:2}}
\emph{Prove the following propositions.}

\begin{exercise}\label{ex:2.1-ex}
Let $A$ be a ring and $I$ a finitely generated ideal satisfying $I = I^2$; then $I$ is
generated by an idempotent $e$ (an element $e$ satisfying $e^2 = e$).
\end{exercise}

\begin{exercise}\label{ex:2.2-ex}
Let $A$ be a ring, $I$ an ideal of $A$ and $M$ a finite $A$-module; then
\[
  \sqrt{\ann(M/IM)} = \sqrt{\ann(M) + I}.
\]
\end{exercise}

\begin{exercise}\label{ex:2.3-ex}
Let $M$ and $N$ be submodules of an $A$-module $L$. If $M + N$ and $M \cap N$ are
finitely generated then so are $M$ and $N$.
\end{exercise}

\begin{exercise}\label{ex:2.4-ex}
Let $A$ be a (commutative) ring, $A \neq 0$. An $A$-module is said to be \emph{free of rank $n$}
if it is isomorphic to $A^n$.
\begin{enumerate}[(a)]
  \item If $A^n \cong A^m$ then $n = m$; prove this by reducing to the case of a field.
        (Note that there are counter-examples to this for non-commutative rings.)
  \item Let $C = (c_{ij})$ be an $n \times m$ matrix over $A$, and suppose that $C$ has a non-
        zero $r \times r$ minor, but that all the $(r + 1) \times (r + 1)$ minors are $0$. Show then
        that if $r < m$, the $m$ column vectors of $C$ are linearly dependent. (Hint: you
        can assume that $m = r + 1$.) Deduce from this an alternative proof of (a).
  \item If $A$ is a local ring, any minimal basis of the free module $A^n$ is a basis
        (that is, a linearly independent set of generators).
\end{enumerate}
\end{exercise}

\begin{exercise}\label{ex:2.5-ex}
Let $A$ be a ring, and $0 \to L \to M \to N \to 0$ an exact sequence of $A$-
modules.
\begin{enumerate}[(a)]
  \item If $L$ and $N$ are both of finite presentation then so is $M$.
  \item If $L$ is finitely generated and $M$ is of finite presentation then $N$ is of
        finite presentation.
\end{enumerate}
\end{exercise}

% ============ §3 Chain conditions (原书 pp.14-19) ============
\section{Chain conditions}\label{sec:3}\index{Chain conditions, section 3}

The following two conditions on a partially ordered set $\Gamma$ are
equivalent:
\begin{itemize}
  \item[(*)] any non-empty subset of $\Gamma$ has a maximal element;
  \item[(**)] any ascending chain $\gamma_1 < \gamma_2 < \cdots$ of elements of $\Gamma$ must stop after a
        finite number of steps.
\end{itemize}
The implication $(*) \Rightarrow (**)$ is obvious. We prove $(**) \Rightarrow (*)$. Let $\Gamma'$ be a non-
empty subset of $\Gamma$. If $\Gamma'$ does not have a maximal element, then by the axiom
of choice, for each $\gamma \in \Gamma'$ we can choose a bigger element of $\Gamma'$, say $\varphi(\gamma)$. Now if
we choose any $\gamma_1 \in \Gamma'$ and set $\gamma_2 = \varphi(\gamma_1)$, $\gamma_3 = \varphi(\gamma_2), \ldots$ then we get an infinite
ascending chain $\gamma_1 < \gamma_2 < \cdots$, contradicting $(**)$.

When these conditions are satisfied we say that $\Gamma$ has the \emph{ascending chain
condition}\index{ascending chain condition} (a.c.c.), or the \emph{maximal condition}. Reversing the order we can define
the \emph{descending chain condition}\index{descending chain condition} (d.c.c.), or \emph{minimal condition} in the same way.

If the set of ideals of a ring $A$ has the a.c.c., we say that $A$ is a \emph{Noetherian
ring}\index{Noetherian ring}, and if it has the d.c.c., that $A$ is an \emph{Artinian ring}.\index{Artinian ring} If $A$ is Noetherian (or
Artinian) and $B$ is a quotient of $A$ then $B$ has the same property; this is
obvious, since the set of ideals of $B$ is order-isomorphic to a subset of
that of $A$.\footnote{The a.c.c. and d.c.c. were first used in a paper of Emmy Noether (1882--1935), \textit{Idealtheorie in Ringbereichen}, Math.\ Ann., 83 (1921). Emil Artin (1898--1962) was, together with Emmy Noether, one of the founders of modern abstract algebra. As well as studying non-commutative rings whose one-sided ideals satisfy the d.c.c., he also discovered the Artin--Rees lemma, which will turn up in \S\ref{sec:8}.}

In the same way, we say that a module is \emph{Noetherian} or \emph{Artinian} if its set of
submodules satisfies the a.c.c. or the d.c.c. If $M$ has either of these properties,
then so do both its quotient modules and its submodules. (A subring of a
Noetherian or Artinian ring does not necessarily have the same property:
why not?)

A ring $A$ is Noetherian if and only if every ideal of $A$ is finitely generated.
(Proof, `only if': given an ideal $I$, consider a maximal element of the set of
finitely generated ideals contained in $I$; this must coincide with $I$. `If': given
an ascending chain $I_1 \subset I_2 \subset \cdots$ of ideals, $\bigcup I_n$ is also an ideal, so that by
assumption it can be generated by finitely many elements $a_1, \ldots, a_l$. There is
some $I_n$ which contains all the $a_i$, and the chain must stop there.)

In exactly the same way, an $A$-module $M$ is Noetherian if and only if
every submodule of $M$ is finitely generated. In particular $M$ itself must be
finitely generated, and if $A$ is Noetherian then this is also sufficient. Thus we
have the well-known fact that finite modules over a Noetherian ring are
Noetherian; we now give a proof of this in a more general form.

\begin{theorem}\label{thm:3.1}
Let $A$ be a ring and $M$ an $A$-module.
\begin{enumerate}[(i)]
  \item Let $M' \subset M$ be a submodule and $\varphi\colon M \to M/M'$ the natural map. If
        $N_1$ and $N_2$ are submodules of $M$ such that $N_1 \subset N_2$, $N_1 \cap M' = N_2 \cap M'$
        and $\varphi(N_1) = \varphi(N_2)$ then $N_1 = N_2$.
  \item Let $0 \to M' \to M \to M'' \to 0$ be an exact sequence of $A$-modules;
        if $M'$ and $M''$ are both Noetherian (or both Artinian), then so is $M$.
  \item Let $M$ be a finite $A$-module; then if $A$ is Noetherian (or Artinian),
        so is $M$.
\end{enumerate}
\end{theorem}

\begin{proof}
(i) is easy, and we leave it to the reader.

(ii) is obtained by applying (i) to an ascending (respectively descending)
chain of submodules of $M$.

(iii) If $M$ is generated by $n$ elements then it is a quotient of the free module
$A^n$, so that it is enough to show that $A^n$ is Noetherian (respectively Artinian).
However, this is clear from (ii) by induction on $n$.
\end{proof}

For a module $M$, it is equivalent to say that $M$ has both the a.c.c. and
the d.c.c., or that $M$ has finite length. Indeed, if $l(M) < \infty$ then $l(M_1) < l(M_2)$
for any two distinct submodules $M_1 \subset M_2 \subset M$, so that the two chain
conditions are clear. Conversely, if $M$ has the d.c.c. then we let $M_1$ be a
minimal non-zero submodule of $M$, let $M_2$ be a minimal element among all
submodules of $M$ strictly containing $M_1$, and proceed in the same way to
obtain an ascending chain $0 = M_0 \subset M_1 \subset M_2 \subset \cdots$; if $M$ also has the a.c.c.
then this chain must stop by arriving at $M$, so that $M$ has a composition
series.

Every submodule of the $\Z$-module $\Z$ is of the form $n\Z$, so that $\Z$ is
Noetherian, but not Artinian. Let $p$ be a prime, and write $W$ for the $\Z$-
module of rational numbers whose denominator is a power of $p$; then
the $\Z$-module $W/\Z$ is not Noetherian, but it is Artinian, since every proper
submodule of $W/\Z$ is either $0$ or is generated by $p^{-n}$ for $n = 1, 2, \ldots$ . This
shows that the a.c.c. and d.c.c. for modules are independent conditions,
but this is not the case for rings, as shown by the following result.

\begin{theorem}[Y.\ Akizuki]\label{thm:3.2}
An Artinian ring is Noetherian.
\end{theorem}

\begin{proof}
Let $A$ be an Artinian ring. It is sufficient to prove that $A$ has finite
length as an $A$-module. First of all, $A$ has only finitely many maximal ideals.
Indeed, if $p_1, p_2, \ldots$ is an infinite set of distinct maximal ideals then it is easy
to see that $p_1 \supset p_1p_2 \supset p_1p_2p_3 \supset \cdots$ is an infinite descending chain of ideals,
which contradicts the assumption. Thus, we let $p_1, p_2, \ldots, p_r$ be all the
maximal ideals of $A$ and set $I = p_1p_2 \cdots p_r = \rad(A)$. The descending chain
$I \supset I^2 \supset \cdots$ stops after finitely many steps, so that there is an $s$ such that
$I^s = I^{s+1}$. If we set $(0:I^s) = J$ then
\[
  (J:I) = ((0:I^s):I) = (0:I^{s+1}) = J;
\]
let's prove that $J = A$. By contradiction, suppose that $J \neq A$; then there exists
an ideal $J'$ which is minimal among all ideals strictly bigger than $J$. For any
$x \in J' - J$ we have $J' = Ax + J$. Now $I = \rad(A)$ and $J \neq J'$, so that by NAK
$J' \neq Ix + J$, and hence by minimality of $J'$ we have $Ix + J = J$, and this gives
$Ix \subset J$. Thus $x \in (J:I) = J$, which is a contradiction. Therefore $J = A$, so that
$I^s = 0$. Now consider the chain of ideals
\[
  A \supset p_1 \supset p_1p_2 \supset \cdots \supset p_1 \cdots p_r \supset (p_1 \cdots p_r)p_1 \supset \cdots \supset (p_1 \cdots p_r)^s = 0.
\]
Let $M$ and $Mp_i$ be any two consecutive terms in this chain; then $M/Mp_i$ is a
vector space over the field $A/p_i$, and since it is Artinian, it must be finite-
dimensional. Hence, $l(M/Mp_i) < \infty$, and therefore the sum $l(A)$ of these
terms is also finite.
\end{proof}

\begin{remark}
This theorem is sometimes referred to as Hopkins' theorem, but it was proved in the
above form by Akizuki [2] in 1935. It was rediscovered four years later by Hopkins [1], and
he proved it for non-commutative rings (a left-Artinian ring with unit is also left-Noetherian).
\end{remark}

\begin{theorem}\label{thm:3.3}
If $A$ is Noetherian then so are $A[X]$ and $A[[X]]$.
\end{theorem}

\begin{proof}
The statement for $A[X]$ is the well-known Hilbert basis theorem
(see, for example Lang, \textit{Algebra}, or [AM], p.~81), and we omit the proof.
We now briefly run through the proof for $A[[X]]$. Set $B = A[[X]]$, and
let $I$ be an ideal of $B$; we will prove that $I$ is finitely generated. Write
$I(r)$ for the ideal of $A$ formed by the leading coefficients $a_r$ of
$f = a_rX^r + a_{r+1}X^{r+1} + \cdots$ as $f$ runs through $I \cap X^rB$; then we have
\[
  I(0) \subset I(1) \subset I(2) \subset \cdots .
\]
Since $A$ is Noetherian, there is an $s$ such that $I(s) = I(s+1) = \cdots$; moreover,
each $I(i)$ is finitely generated. For each $i$ with $0 \le i \le s$ we take finitely many
elements $a_{i\nu} \in A$ generating $I(i)$, and choose $g_{i\nu} \in I \cap X^iB$ having $a_{i\nu}$ as the
coefficient of $X^i$. These $g_{i\nu}$ now generate $I$. Indeed, for $f \in I$ we can take
a linear combination $g_0$ of the $g_{0\nu}$ with coefficients in $A$ such that
$f - g_0 \in I \cap XB$, then take a linear combination $g_1$ of the $g_{1\nu}$ with
coefficients in $A$ such that $f - g_0 - g_1 \in I \cap X^2B$, and proceeding in
the same way we get
\[
  f - g_0 - g_1 - \cdots - g_s \in I \cap X^{s+1}B.
\]
Now $I(s+1) = I(s)$, so we can take a linear combination $g_{s+1}$ of the $Xg_{s\nu}$
with coefficients in $A$ such that
\[
  f - g_0 - g_1 - \cdots - g_{s+1} \in I \cap X^{s+2}B.
\]
We now proceed in the same way to get $g_{s+2}, \ldots$. For $i \le s$, each $g_i$
is a linear combination of the $g_{i\nu}$ with coefficients in $A$, and, for
$i > s$, a combination of the elements $X^{i-s}g_{s\nu}$. For each $i > s$ we write
$g_i = \sum_\nu a_{i\nu}X^{i-s}g_{s\nu}$, and then for each $\nu$ we set
$h_\nu = \sum_{i \ge s} a_{i\nu}X^{i-s}$; $h_\nu$ is an element of $B$, and
\[
  f = g_0 + \cdots + g_{s-1} + \sum_\nu h_\nu g_{s\nu}.
\]
\end{proof}

A ring $A[b_1, \ldots, b_n]$ which is finitely generated as a ring over a
Noetherian ring $A$ is a quotient of a polynomial ring $A[X_1, \ldots, X_n]$, and
so by the Hilbert basis theorem is again Noetherian. We now give some
other criteria for a ring to be Noetherian.

\begin{theorem}[I.\,S.\ Cohen]\label{thm:3.4}
If all the prime ideals of a ring $A$ are finitely
generated then $A$ is Noetherian.
\end{theorem}

\begin{proof}
Write $\Gamma$ for the set of ideals of $A$ which are not finitely generated.
If $\Gamma \neq \varnothing$ then by Zorn's lemma $\Gamma$ contains a maximal element $I$. Then $I$ is
not a prime ideal, so that there are elements $x, y \in A$ with $x \notin I$, $y \notin I$ but
$xy \in I$. Now $I + Ay$ is bigger than $I$, and hence is finitely generated, so
that we can choose $u_1, \ldots, u_n \in I$ such that
\[
  I + Ay = (u_1, \ldots, u_n, y).
\]
Moreover, $I:y = \{a \in A \mid ay \in I\}$ contains $x$, and is thus bigger than
$I$, so it has a finite system of generators $(v_1, \ldots, v_m)$. Finally, it is
easy to check that $I = (u_1, \ldots, u_n, yu_1, \ldots, yu_n, v_1y, \ldots, v_my)$; hence, $I \notin \Gamma$, which is a
contradiction. Therefore $\Gamma = \varnothing$.
\end{proof}

\begin{theorem}\label{thm:3.5}
Let $A$ be a ring and $M$ an $A$-module. Then if $M$ is a
Noetherian module, $A/\ann(M)$ is a Noetherian ring.
\end{theorem}

\begin{proof}
If we set $\bar{A} = A/\ann(M)$ and view $M$ as an $\bar{A}$-module, then the
submodules of $M$ as an $A$-module or $\bar{A}$-module coincide, so that $M$ is
also Noetherian as an $\bar{A}$-module. We can thus replace $A$ by $\bar{A}$, and then
$\ann(M) = (0)$. Now letting $M = A\omega_1 + \cdots + A\omega_n$, we can embed $A$ in $M^n$
by means of the map $a \mapsto (a\omega_1, \ldots, a\omega_n)$. By Theorem 1, $M^n$ is a Noetherian
module, so that its submodule $A$ is also Noetherian. (This theorem can
be expressed by saying that a ring having a faithful Noetherian module
is Noetherian.)
\end{proof}

\begin{theorem}[E.\ Formanek~\cite{Fo1}]\label{thm:3.6}
Let $A$ be a ring, and $B$ an $A$-module
which is finitely generated and faithful over $A$. Assume that the set of
submodules of $B$ of the form $IB$ with $I$ an ideal of $A$ satisfies the a.c.c.;
then $A$ is Noetherian.
\end{theorem}

\begin{proof}
It will be enough to show that $B$ is a Noetherian $A$-module. By
contradiction, suppose that it is not; then the set
\[
  \{IB \mid I \text{ is an ideal of } A \text{ and } B/IB \text{ is non-Noetherian as } A\text{-module}\}
\]
contains $(0)$ and so is non-empty, so that by assumption it contains a
maximal element. Let $IB$ be one such maximal element; then replacing $B$
by $B/IB$ and $A$ by $A/\ann(B/IB)$ we see that we can assume that $B$ is a
non-Noetherian $A$-module, but for any non-zero ideal $I$ of $A$ the quotient
$B/IB$ is Noetherian.

Next we set
\[
  \mathcal{I} = \{N \mid N \text{ is a submodule of } B \text{ and } B/N \text{ is a faithful } A\text{-module}\}.
\]
If $B = Ab_1 + \cdots + Ab_n$ then for a submodule $N$ of $B$,
\[
  N \in \mathcal{I} \iff \forall a \in A - \{0\}, \quad \{ab_1, \ldots, ab_n\} \nsubseteq N.
\]
From this, one sees at once that Zorn's lemma applies to $\mathcal{I}$; hence there
exists a maximal element $N_0$ of $\mathcal{I}$. If $B/N_0$ is Noetherian then $A$ is a
Noetherian ring, and thus $B$ is Noetherian, which contradicts our
hypothesis. It follows that on replacing $B$ by $B/N_0$ we arrive at a module
$B$ with the following properties:
\begin{enumerate}[(1)]
  \item $B$ is non-Noetherian as an $A$-module;
  \item for any ideal $I \neq (0)$ of $A$, $B/IB$ is Noetherian;
  \item for any submodule $N \neq (0)$ of $B$, $B/N$ is not faithful as an $A$-module.
\end{enumerate}
Now let $N$ be any non-zero submodule of $B$. By (3) there is an element
$a \in A$ with $a \neq 0$ such that $a(B/N) = 0$, that is such that $aB \subset N$. By (2)
$B/aB$ is a Noetherian module, so that $N/aB$ is finitely generated; but since
$B$ is finitely generated so is $aB$, and hence $N$ itself is finitely generated.
Thus, $B$ is a Noetherian module, which contradicts (1).
\end{proof}

As a corollary of this theorem we get the following result.

\begin{theorem}\label{thm:3.7}
\begin{enumerate}[(i)]
  \item (Eakin--Nagata theorem). Let $B$ be a Noetherian ring, and $A$ a subring
        of $B$ such that $B$ is finite over $A$; then $A$ is also a Noetherian ring.
  \item Let $B$ be a non-commutative ring whose right ideals have the a.c.c.,
        and let $A$ be a commutative subring of $B$. If $B$ is finitely generated as a
        left $A$-module then $A$ is a Noetherian ring.
  \item Let $B$ be a non-commutative ring whose two-sided ideals have the
        a.c.c., and let $A$ be a subring contained in the centre of $B$; if $B$ is finitely
        generated as an $A$-module then $A$ is a Noetherian ring.
\end{enumerate}
\end{theorem}

\begin{proof}
$B$ has a unit, so is faithful as an $A$-module. Hence it is enough to
apply the previous theorem.
\end{proof}

\begin{remark}
Part (i) of Theorem 7 was proved in Eakin's thesis [1] in 1968,
and the same result was obtained independently by Nagata [9] a little
later. Subsequently many alternative proofs and extensions to the non-
commutative case were published; the most transparent of these seems to
be Formanek's result~\cite{Fo1}, which we have given above in the form of
Theorem 6. However, this also goes back to the idea of the proofs of
Eakin and Nagata.
\end{remark}

% ============ 习题 ============
\section*{Exercises to \S\ref{sec:3}}
\emph{Prove the following propositions.}

\begin{exercise}\label{ex:3.1-ex}
Let $I_1, \ldots, I_n$ be ideals of a ring $A$ such that $I_1 \cap \cdots \cap I_n = (0)$; if each $A/I_i$ is a
Noetherian ring then so is $A$.
\end{exercise}

\begin{exercise}\label{ex:3.2-ex}
Let $A$ and $B$ be Noetherian rings, and $f\colon A \to C$ and $g\colon B \to C$ ring
homomorphisms. If both $f$ and $g$ are surjective then the fibre product
$A \times_C B$ (that is, the subring of the direct product $A \times B$ given by
$\{(a, b) \in A \times B \mid f(a) = g(b)\}$) is a Noetherian ring.
\end{exercise}

\begin{exercise}\label{ex:3.3-ex}
Let $A$ be a local ring such that the maximal ideal $m$ is principal and
$\bigcap_{n=1}^{\infty} m^n = (0)$. Then $A$ is Noetherian, and every non-zero ideal of $A$ is a
power of $m$.
\end{exercise}

\begin{exercise}\label{ex:3.4-ex}
Let $A$ be an integral domain with field of fractions $K$. A \emph{fractional ideal} $I$ of
$A$ is an $A$-submodule $I$ of $K$ such that $I \neq 0$ and $aI \subset A$ for some $0 \neq a \in K$.
The product of two fractional ideals is defined in the same way as the
product of two ideals. If $I$ is a fractional ideal of $A$ we set $I^{-1} = \{\alpha \in K \mid \alpha I
\subset A\}$; this is also a fractional ideal, and $II^{-1} \subset A$. In the particular case
that $II^{-1} = A$ we say that $I$ is \emph{invertible}. An invertible fractional ideal of $A$
is finitely generated as an $A$-module.
\end{exercise}

\begin{exercise}\label{ex:3.5-ex}
If $A$ is a UFD, the only ideals of $A$ which are invertible as fractional ideals
are the principal ideals.
\end{exercise}

\begin{exercise}\label{ex:3.6-ex}
Let $A$ be a Noetherian ring, and $\varphi\colon A \to A$ a homomorphism of rings.
Then if $\varphi$ is surjective it is also injective, and hence an automorphism of $A$.
\end{exercise}

\begin{exercise}\label{ex:3.7-ex}
If $A$ is a Noetherian ring then any finite $A$-module is of finite presentation,
but if $A$ is non-Noetherian then $A$ must have finite $A$-modules which are
not of finite presentation.
\end{exercise}


% ============ Chapter 2. Prime ideals (原书 pp.20-45) ============
\chapter{Prime ideals}\label{ch:2}

The notion of prime ideal is central to commutative ring theory. The set
$\Spec A$ of prime ideals of a ring $A$ is a topological space, and the
`localisation' of rings and modules with respect to this topology is an
important technique for studying them. These notions are discussed in \S\ref{sec:4}.
Starting with the topology of $\Spec A$, we can define the dimension of $A$ and
the height of a prime ideal as notions with natural geometrical content. In
\S\ref{sec:5} we treat elementary dimension theory using only field theory, developing
especially the dimension theory of ideals in polynomial rings, including the
Hilbert Nullstellensatz. We also discuss, as an example of an application of the
notion of dimension, the theory of Forster and Swan on estimates for the
number of generators of a module. (Dimension theory will be the subject of
a detailed study in Chapter~\ref{ch:5} using methods of ring theory). In \S\ref{sec:6} we discuss
the classical theory of primary decomposition as modernised by Bourbaki.

% ============ §4 Localisation and Spec of a ring (原书 pp.20-29) ============
\section{Localisation and Spec of a ring}\label{sec:4}\index{Localisation, section 4}\index{Spec of a ring, section 4}

Let $A$ be a ring and $S \subset A$ a multiplicative set; that is (as in
\S\ref{sec:1}) suppose that
\[
  \text{(i) } x, y \in S \Rightarrow xy \in S, \quad \text{and} \quad \text{(ii) } 1 \in S.
\]

\begin{definition}
Suppose that $f\colon A \to B$ is a ring homomorphism satisfying
the two conditions
\begin{enumerate}[(1)]
  \item $f(x)$ is a unit of $B$ for all $x \in S$;
  \item if $g\colon A \to C$ is a homomorphism of rings taking every element of $S$ to
        a unit of $C$ then there exists a unique homomorphism
        $h\colon B \to C$ such that $g = hf$;
\end{enumerate}
then $B$ is uniquely determined up to isomorphism, and is called the
\emph{localisation}\index{localisation} or the \emph{ring of fractions} of $A$ with respect to $S$. We write
$B = S^{-1}A$ or $A_S$, and call $f\colon A \to A_S$ the canonical map.
\end{definition}

We prove the existence of $B$ as follows: define the relation $\sim$ on the
set $A \times S$ by
\[
  (a, s) \sim (b, s') \iff \exists t \in S \text{ such that } t(s'a - sb) = 0;
\]
it is easy to check that this is an equivalence relation (if we just had
$s'a = sb$ in the definition, the transitive law would fail when $S$ has
zero-divisors). Write $a/s$ for the equivalence class of $(a, s)$ under $\sim$, and
let $B$ be the set of these; sums and products are defined in $B$ by the usual
rules for calculating with fractions:
\[
  a/s + b/s' = (as' + bs)/ss', \qquad (a/s) \cdot (b/s') = ab/ss'.
\]
This makes $B$ into a ring, and defining $f\colon A \to B$ by $f(a) = a/1$ we see that
$f$ is a homomorphism of rings satisfying (1) and (2) above. Indeed, if $s \in S$
then $f(s) = s/1$ has the inverse $1/s$; and if $g\colon A \to C$ is as in (2) then we just
have to set $h(a/s) = g(a)g(s)^{-1}$ (the reader should check that $a/s = b/s'$
implies $g(a)g(s)^{-1} = g(b)g(s')^{-1}$). From this construction we see that the
kernel of the canonical map $f\colon A \to A_S$ is given by
\[
  \Ker f = \{a \in A \mid sa = 0 \text{ for some } s \in S\}.
\]
Hence $f$ is injective if and only if $S$ does not contain any zero-divisors
of $A$. In particular, the set of all non-zero-divisors of $A$ is a multiplicative
set; the ring of fractions with respect to $S$ is called the \emph{total ring of fractions}
of $A$. If $A$ is an integral domain then its total ring of fractions is the same
thing as its field of fractions.

In general, let $f\colon A \to B$ be any ring homomorphism, $I$ an ideal of $A$
and $J$ an ideal of $B$. According to the conventions at the beginning of the
book, we write $IB$ for the ideal $f(I)B$ of $B$. This is called the \emph{extension} of
$I$ to $B$, or the \emph{extended ideal}\index{extended ideal}, and is sometimes also written $I^e$. Moreover,
we write $J \cap A$ for the ideal $f^{-1}(J)$ of $A$. This is called the \emph{contracted ideal}\index{contracted ideal}
of $J$, and is sometimes also written $J^c$. In this notation, the inclusions
\[
  I^{ec} \supset I \quad \text{and} \quad J^{ce} \subset J
\]
follow immediately from the definitions; from the first inclusion we get
$I \subset I^e$, but substituting $J = I^e$ in the second gives $I^{ece} \subset I^e$, and hence
\[
  (*) \quad I^{ece} = I^e, \quad \text{and similarly} \quad J^{cec} = J^c.
\]
This shows that there is a canonical bijection between the sets $\{IB \mid I$
is an ideal of $A\}$ and $\{J \cap A \mid J$ is an ideal of $B\}$.

If $P$ is a prime ideal of $B$ then $B/P$ is an integral domain, and since
$A/P^c$ can be viewed as a subring of $B/P$ it is also an integral domain, so
that $P^c$ is a prime ideal of $A$. (The extended ideal of a prime ideal does
not have to be prime.)

An ideal $J$ of $B$ is said to be \emph{primary}\index{primary ideal} if it satisfies the two conditions:
(1) $1 \notin J$, and (2) for $x, y \in B$, if $xy \in J$ and $x \notin J$ then $y^n \in J$ for some
$n > 0$; in other words, all zero-divisors of $B/J$ are nilpotent. The property
that all zero-divisors are nilpotent passes to subrings, so that just as for
prime ideals we see that the contraction of a primary ideal remains primary.
If $J$ is primary then $\sqrt{J}$ is a prime ideal (see Ex.~\ref{ex:4.1-ex}).

The importance of rings of fractions for ring theory stems mainly from
the following theorem.

\begin{theorem}\label{thm:4.1}
\begin{enumerate}[(i)]
  \item All the ideals of $A_S$ are of the form $IA_S$, with $I$ an ideal of $A$.
  \item Every prime ideal of $A_S$ is of the form $pA_S$ with $p$ a prime ideal of
        $A$ disjoint from $S$, and conversely, $pA_S$ is prime in $A_S$ for every such $p$;
        exactly the same holds for primary ideals.
\end{enumerate}
\end{theorem}

\begin{proof}
(i) If $J$ is an ideal of $A_S$, set $I = J \cap A$. If $x = a/s \in J$ then
$x \cdot f(s) = f(a) \in J$, so that $a \in I$, and then $x = (1/s) \cdot f(a) \in IA_S$. The converse
inclusion $IA_S \subset J$ is obvious, so that $J = IA_S$.

(ii) If $P$ is a prime ideal of $A_S$ and we set $p = P \cap A$, then $p$ is a prime
ideal of $A$, and from the above proof $P = pA_S$. Moreover, since $P$ does
not contain units of $A_S$, we have $p \cap S = \emptyset$. Conversely, if $p$ is a prime
ideal of $A$ disjoint from $S$ then
\[
  \frac{a}{s} \cdot \frac{b}{t} \in pA_S \quad \text{with } s, t \in S \Rightarrow rab \in p \text{ for some } r \in S,
\]
and since $r \notin p$ we must have $a \in p$ or $b \in p$, so that $a/s$ or $b/t \in pA_S$. One
also sees easily that $1 \notin pA_S$, so that $pA_S$ is a prime ideal of $A_S$.

For primary ideals the argument is exactly the same: if $p$ is a primary
ideal of $A$ disjoint from $S$ and if $rab \in p$ with $r \in S$, then since no power
of $r$ is in $p$ we have $ab \in p$. From this we get either $a/s \in pA_S$ or $(b/t)^n \in pA_S$
for some $n$.
\end{proof}

\begin{corollary*}
If $A$ is Noetherian (or Artinian) then so is $A_S$.
\end{corollary*}

\begin{proof}
This follows from (i) of the theorem.
\end{proof}

We now give examples of rings of fractions $A_S$ for various multiplicative
sets $S$.

\begin{example}\label{ex:4.1}
Let $a \in A$ be an element which is not nilpotent, and set
$S = \{1, a, a^2, \ldots\}$. In this case we sometimes write $A_a$ for $A_S$. (The reason
for not allowing $a$ to be nilpotent is so that $0 \notin S$. In general if $0 \in S$ then
from the construction of $A_S$ it is clear that $A_S = 0$, which is not very
interesting.) The prime ideals of $A_a$ correspond bijectively with the prime
ideals of $A$ not containing $a$.
\end{example}

\begin{example}\label{ex:4.2}
Let $p$ be a prime ideal of $A$, and set $S = A - p$. In this case
we usually write $A_p$ for $A_S$. (Writing $A_p$ and $A_{A-p}$ to denote the same thing
is totally illogical notation, and the Bourbaki school avoids $A_S$, writing
$S^{-1}A$ instead; however, the notation $A_S$ does not lead to any confusion.)
The localisation $A_p$ is a local ring with maximal ideal $pA_p$. Indeed, as we saw
in Theorem 1, $pA_p$ is a prime ideal of $A_p$, and furthermore, if $J \subset A_p$ is any
proper ideal then $I = J \cap A$ is an ideal of $A$ disjoint from $A - p$, and so $I \subset p$,
giving $J = IA_p \subset pA_p$. The prime ideals of $A_p$ correspond bijectively with
the prime ideals of $A$ contained in $p$.
\end{example}

\begin{example}\label{ex:4.3}
Let $I$ be a proper ideal of $A$ and set $S = 1 + I = \{1 + x \mid
x \in I\}$. Then $S$ is a multiplicative set, and the prime ideals of $A_S$
correspond bijectively with the prime ideals $p$ of $A$ such that $I + p \neq A$.
\end{example}

\begin{example}\label{ex:4.4}
Let $S$ be a multiplicative set, and set $\bar{S} = \{a \in A \mid ab \in S \text{ for
some } b \in A\}$. Then $\bar{S}$ is also a multiplicative set, called the \emph{saturation}\index{saturation (of a multiplicative set)}
of $S$. Since quite generally a divisor of a unit is again a unit, we see from
the definition of the ring of fractions that $A_S = A_{\bar{S}}$, and $\bar{S}$ is maximal among
multiplicative sets $T$ such that $A_S = A_T$. Indeed, one sees easily that $\bar{S} =
\{a \in A \mid a/1 \text{ is a unit in } A_S\}$. The multiplicative set $S = A - p$ of
Example 2 is already saturated.
\end{example}

\begin{theorem}\label{thm:4.2}
Localisation commutes with passing to quotients by ideals.
More precisely, let $A$ be a ring, $S \subset A$ a multiplicative set, $I$ an ideal of $A$
and $\bar{S}$ the image of $S$ in $A/I$; then
\[
  A_S/IA_S \cong (A/I)_{\bar{S}}.
\]
\end{theorem}

\begin{proof}
Both sides have the universal property for ring homomorphisms
$g\colon A \to C$ such that
\begin{enumerate}[(1)]
  \item every element of $S$ maps to a unit of $C$,
  \item every element of $I$ maps to $0$;
\end{enumerate}
the isomorphism follows by the uniqueness of the solution to a universal
mapping problem. In concrete terms the isomorphism is given by
\[
  a/s \bmod IA_S \mapsto \bar{a}/\bar{s}, \quad \text{where } \bar{a} = a + I, \; \bar{s} = s + I.
\]
\end{proof}

In particular, if $p$ is a prime ideal of $A$ then
\[
  A_p/pA_p \cong (A/p)_{A-p}.
\]
The left-hand side is the \emph{residue field} of the local ring $A_p$, whereas the
right-hand side is the field of fractions of the integral domain $A/p$. This
field is written $\kappa(p)$ and called the \emph{residue field of $p$}.\index{residue field of a prime ideal}

\begin{theorem}\label{thm:4.3}
Let $A$ be a ring, $S \subset A$ a multiplicative set, and
$f\colon A \to A_S$ the canonical map. If $B$ is a ring, with ring homomorphisms
$g\colon A \to B$ and $h\colon B \to A_S$ satisfying
\begin{enumerate}[(1)]
  \item $f = hg$,
  \item for every $b \in B$ there exists $s \in S$ such that $g(s) \cdot b \in g(A)$,
\end{enumerate}
then $A_S$ can also be regarded as a ring of fractions of $B$. More precisely,
\[
  A_S = B_{g(S)} = B_T, \quad \text{where } T = \{t \in B \mid h(t) \text{ is a unit of } A_S\}.
\]
\end{theorem}

\begin{proof}
We can factorise $h$ as $B \to B_T \to A_S$; write $\alpha\colon B_T \to A_S$ for
the second of these maps. Now $g(S) \subset T$, so that the composite $A \to B
\to B_T$ factorises as $A \to A_S \to B_T$; write $\beta\colon A_S \to B_T$ for the second of
these maps. Then
\[
  \alpha\beta = 1 \quad \text{and} \quad \beta\alpha = 1,
\]
so that $\beta\alpha = 1$, the identity map of $A_S$. Moreover by assumption, for $b \in B$
there exist $a \in A$ and $s \in S$ such that $bg(s) = g(a)$. Hence, $\beta(a/s) = g(a)/g(s) =
b/1$. In particular for $t \in T$, if we take $a \in A$ such that $t/1 = \beta(a)$ then
$a = \alpha\beta(a) = \alpha(t/1) = h(t)$, so that $a$ is a unit of $A_S$. Hence, $b/t = \beta(a/s)\beta(a^{-1})$,
and $\beta$ is surjective. Thus $\alpha$ and $\beta$ are mutually inverse, giving an
isomorphism $A_S \cong B_T$. The fact that $A_S \cong B_{g(S)}$ can be proved similarly.
(Alternatively, this follows since $T$ is the saturation of the multiplicative
set $g(S)$. The reader should check this for himself.)
\end{proof}

\begin{corollarynum}\label{cor:4.1}
If $p$ is a prime ideal of $A$, $S = A - p$ and $B$ satisfies the
conditions of the theorem, then setting $P = pA_S \cap B$ we have $A_S = B_P$.
\end{corollarynum}

\begin{proof}
Under these circumstances the $T$ in the theorem is exactly $B - P$.
\end{proof}

\begin{corollarynum}\label{cor:4.2}
Let $S \subset A$ be a multiplicative set not containing any zero-
divisors; then $A$ can be viewed as a subring of $A_S$, and for any intermediate
ring $A \subset B \subset A_S$, the ring $A_S$ is a ring of fractions of $B$.
\end{corollarynum}

\begin{corollarynum}\label{cor:4.3}
If $S$ and $T$ are two multiplicative subsets of $A$ with $S \subset T$,
then writing $\bar{T}$ for the image of $T$ in $A_S$, we have $(A_S)_{\bar{T}} = A_T$.
\end{corollarynum}

\begin{corollarynum}\label{cor:4.4}
If $S \subset A$ is a multiplicative set and $P$ is a prime ideal of $A$
disjoint from $S$ then $(A_S)_{PA_S} = A_P$. In particular if $P \subset Q$ are prime ideals of
$A$, then
\[
  (A_P)_{QA_P} = A_Q.
\]
\end{corollarynum}

\begin{definition}
The set of all prime ideals of a ring $A$ is called the \emph{spectrum}\index{spectrum of a ring}
of $A$, and written $\Spec A$; the set of maximal ideals of $A$ is called the
\emph{maximal spectrum} of $A$, and written m-Spec $A$.
\end{definition}

By Theorem~\ref{thm:1.1} we have
\[
  A \neq 0 \iff \text{m-Spec } A \neq \emptyset \iff \Spec A \neq \emptyset.
\]
If $I$ is any ideal of $A$, we set
\[
  V(I) = \{p \in \Spec A \mid p \supset I\}.
\]
Then
\[
  V(I) \cup V(I') = V(I \cap I') = V(II'),
\]
and for any family $\{I_\lambda\}_{\lambda \in \Lambda}$ of ideals we have
\[
  \bigcap_{\lambda \in \Lambda} V(I_\lambda) = V\left(\sum_{\lambda \in \Lambda} I_\lambda\right).
\]
From this it follows that $\mathcal{F} = \{V(I) \mid I \text{ is an ideal of } A\}$ is closed under
finite unions and arbitrary intersections, so that there is a topology on
$\Spec A$ for which $\mathcal{F}$ is the set of closed sets. This is called the \emph{Zariski
topology}\index{Zariski topology}. From now on we will usually consider the spectrum of a ring
together with its Zariski topology. m-Spec $A$ will be considered with the
subspace topology, which we will also call the Zariski topology.

For $a \in A$ we set $D(a) = \{p \in \Spec A \mid a \notin p\}$; this is the complement
of $V(aA)$, and so is an open set. Conversely, any open set of $\Spec A$ can
be written as the union of open sets of the form $D(a)$. Indeed, if
$U = \Spec A - V(I)$ then $U = \bigcup_{a \in I} D(a)$. Hence, the open sets of the form
$D(a)$ form a basis for the topology of $\Spec A$.

Let $f\colon A \to B$ be a ring homomorphism. For $P \in \Spec B$, the ideal
$P \cap A = f^{-1}(P)$ is a point of $\Spec A$. The map $\Spec B \to \Spec A$ defined by
taking $P$ into $P \cap A$ is written ${}^a f$. As one sees easily, $({}^a f)^{-1}(V(I)) = V(IB)$,
so that ${}^a f$ is continuous. If $g\colon B \to C$ is another ring homomorphism then
obviously ${}^a(gf) = {}^a f \cdot {}^a g$. Hence, the correspondence $A \mapsto \Spec A$ and
$f \mapsto {}^a f$ defines a contravariant functor from the category of rings to the
category of topological spaces. If ${}^a f(P) = p$, that is if $P \cap A = p$, we say
that $P$ \emph{lies over} $p$.\index{lies over}

\begin{remark}
For $P$ a maximal ideal of $B$ it does not necessarily follow that
$P \cap A$ is a maximal ideal of $A$; for an example we need only consider the
natural inclusion $A \to B$ of an integral domain $A$ in its field of fractions
$B$. Thus the correspondence $A \mapsto$ m-Spec $A$ is not functorial. This is one
reason for thinking of $\Spec A$ as more important than m-Spec $A$. On the
other hand, one could say that $\Spec A$ contains too many points; for
example, the set $\{p\}$ consisting of a single point is closed in $\Spec A$ if and
only if $p$ is a maximal ideal (in general the closure of $\{p\}$ coincides with
$V(p)$), so that $\Spec A$ almost never satisfies the separation axiom $T_1$.
\end{remark}

Let $M$ be an $A$-module and $S \subset A$ a multiplicative set; we define the
localisation $M_S$ of $M$ in the same way as $A_S$. That is,
\[
  M_S = \left\{\frac{m}{s} \;\middle|\; m \in M, \; s \in S\right\},
\]
and
\[
  \frac{m}{s} = \frac{m'}{s'} \iff t(s'm - sm') = 0 \text{ for some } t \in S.
\]
If we define addition in $M_S$ and scalar multiplication by elements
of $A_S$ by
\[
  m/s + m'/s' = (s'm + sm')/ss' \quad \text{and} \quad (a/s) \cdot (m/s') = am/ss',
\]
then $M_S$ becomes an $A_S$-module, and a canonical $A$-linear map $M \to M_S$
is given by $m \mapsto m/1$; the kernel is $\{m \in M \mid sm = 0 \text{ for some } s \in S\}$. If $S =
A - p$ is the complement of a prime ideal $p$ of $A$ we write $M_p$ for $M_S$. The set
$\{p \in \Spec A \mid M_p \neq 0\}$ is called the \emph{support} of $M$,\index{support of a module} and written $\Supp(M)$. If $M$ is
finitely generated, and we let $M = A\omega_1 + \cdots + A\omega_n$, then
\[
  p \in \Supp(M) \iff M_p \neq 0 \iff \exists i \text{ such that } \omega_i \neq 0 \text{ in } M_p
\]
\[
  \iff \exists i \text{ such that } \ann(\omega_i) \subset p \iff \ann(M) = \bigcap_{i=1}^{n} \ann(\omega_i) \subset p,
\]
so that $\Supp(M)$ coincides with the closed subset $V(\ann(M))$ of $\Spec A$.

\begin{theorem}\label{thm:4.4}
$M_S \cong M \otimes_A A_S$.
\end{theorem}

\begin{proof}
The map $M \times A_S \to M_S$ defined by $(m, a/s) \mapsto am/s$ is $A$-bilinear,
so that there exists a linear map $\alpha\colon M \otimes A_S \to M_S$ such that $\alpha(m \otimes a/s) =
am/s$. Conversely we can define $\beta\colon M_S \to M \otimes A_S$ by $\beta(m/s) = m \otimes (1/s)$;
indeed, if $m/s = m'/s'$ then $ts'm = tsm'$ for some $t \in S$, and so
\[
  m \otimes (1/s) = m \otimes (ts'/tss') = ts'm \otimes (1/tss') = tsm' \otimes (1/tss')
                 = m' \otimes (1/s').
\]
Now it is easy to check that $\alpha$ and $\beta$ are mutually inverse $A_S$-linear
maps. Hence, $M_S$ and $M \otimes_A A_S$ are isomorphic as $A_S$-modules.
\end{proof}

\begin{theorem}\label{thm:4.5}
$M \mapsto M_S$ is an exact (covariant) functor from the category
of $A$-modules to the category of $A_S$-modules. That is, for a morphism of
$A$-modules $f\colon M \to N$ there is a morphism $f_S\colon M_S \to N_S$ of $A_S$-modules
such that
\[
  (1_M)_S = 1_{M_S}, \qquad (gf)_S = g_S f_S,
\]
and such that an exact sequence $0 \to M' \to M \to M'' \to 0$ goes into an
exact sequence $0 \to M'_S \to M_S \to M''_S \to 0$.
\end{theorem}

\begin{proof}
To prove the exactness of $0 \to M'_S \to M_S$ on the last line, view $M'$
as a submodule of $M$; then for $x \in M'$ and $s \in S$,
\[
  x/s = 0 \text{ in } M_S \iff tx = 0 \text{ for some } t \in S
\]
\[
  \iff x/s = 0 \text{ in } M'_S,
\]
as required. The remaining assertions follow from the properties of the
tensor product (see Appendix A) and from the previous theorem. (Of
course they can easily be proved directly.)
\end{proof}

It follows from this that localisation commutes with $\otimes$ and with $\Tor$,
and we will treat all this together in the section on flatness in \S\ref{sec:7}.

Let $A$ be a ring, $M$ an $A$-module and $p \in \Spec A$. There are at least
two interpretations of what it should mean that some property of $A$ or $M$
holds `locally at $p$'. Namely, this could mean that $A_p$ (or $M_p$) has the
property, or it could mean that $A_q$ (or $M_q$) has the property for all $q$ in
some neighbourhood $U$ of $p$ in $\Spec A$. The first of these is more commonly
used, but there are cases when the two interpretations coincide. In any
case, we now prove a number of theorems which assert that a local property
implies a global one.

\begin{theorem}\label{thm:4.6}
Let $A$ be a ring, $M$ an $A$-module and $x \in M$. If $x = 0$ in $M_p$
for every maximal ideal $p$ of $A$, then $x = 0$.
\end{theorem}

\begin{proof}
\[
  x = 0 \text{ in } M_p \iff sx = 0 \text{ for some } s \in A - p \iff \ann(x) \not\subset p.
\]
However, if $1 \notin \ann(x)$ then by Theorem~\ref{thm:1.1}, there must exist a maximal
ideal containing $\ann(x)$. Therefore $1 \in \ann(x)$, that is $x = 0$.
\end{proof}

\begin{theorem}\label{thm:4.7}
Let $A$ be an integral domain with field of fractions $K$; set
$\mathcal{X} = \text{m-Spec } A$. We consider any ring of fractions of $A$ as a subring of $K$.
Then in this sense we have
\[
  A = \bigcap_{m \in \mathcal{X}} A_m.
\]
\end{theorem}

\begin{proof}
For $x \in K$ the set $I = \{a \in A \mid ax \in A\}$ is an ideal of $A$. Now
$x \in A_m$ is equivalent to $I \not\subset m$, so that if $x \in A_m$ for every maximal ideal
$m$ then $1 \in I$, that is $x \in A$.
\end{proof}

\begin{remark}
The above $I$ is the ideal consisting of all possible denominators
of $x$ when written as a fraction of elements of $A$, together with $0$, and this
can be called the \emph{ideal of denominators} of $x$; similarly $Ix$ can be called
the \emph{ideal of numerators} of $x$.
\end{remark}

\begin{theorem}\label{thm:4.8}
Let $A$ be a ring and $M$ a finite $A$-module. If $M \otimes_A \kappa(m) = 0$
for every maximal ideal $m$ then $M = 0$.
\end{theorem}

\begin{proof}
$\kappa(m) = A_m/mA_m$, so that $M \otimes_A \kappa(m) = M_m/mM_m$, and by NAK
(Theorem~\ref{thm:2.2}), $M \otimes_A \kappa(m) = 0 \iff M_m = 0$. Thus the assertion follows from
Theorem~\ref{thm:4.6}.
\end{proof}

The theorem just proved is easy, but we can weaken the assumption
that $M$ is finite over $A$; we have the following result.

\begin{theorem}\label{thm:4.9}
Let $f\colon A \to B$ be a homomorphism of rings, and $M$ a finite
$B$-module; if $M \otimes_A \kappa(p) = 0$ for every $p \in \Spec A$, then $M = 0$.
\end{theorem}

\begin{proof}
If $M \neq 0$ then by Theorem 6 there is a maximal ideal $P$ of $B$ such
that $M_P \neq 0$, so that by NAK, $M_P/PM_P \neq 0$. If we now set $p = P \cap A$ then,
since $pM_P \subset PM_P$, we have $M_P/pM_P \neq 0$. Set $T = B - P$ and $S = A - p$;
then the localisation $M_S = M_p$ of $M$ as an $A$-module and the localisation
$M_T$ of $M$ as a $B$-module coincide (both of them are $\{m/s \mid m \in M, \; s \in S\}$).
We have $f(S) \subset T$, so that
\[
  M_p = M_T \neq 0 \quad \text{and hence} \quad M_p/PM_p = (M_p/PM_p)_T = (M \otimes_A \kappa(p))_T \neq 0;
\]
it follows that $M \otimes_A \kappa(p) \neq 0$.
\end{proof}

\begin{remark}
In Theorem~\ref{thm:4.9} we cannot restrict $p$ to be a maximal ideal of
$A$. As one sees from the proof we have just given, $M = 0$ provided that
$M \otimes \kappa(p) = 0$ for every ideal $p$ which is the restriction of a maximal ideal
of $B$. However, if for example $(A, m)$ is a local integral domain with field
of fractions $B$, and $M = B$, then $M \otimes \kappa(m) = B/mB = 0$, but $M \neq 0$.
\end{remark}

\begin{theorem}\label{thm:4.10}
Let $A$ be a ring and $M$ a finite $A$-module.
\begin{enumerate}[(i)]
  \item For any non-negative integer $r$ set
        \[
          U_r = \{p \in \Spec A \mid M_p \text{ can be generated over } A_p \text{ by } r \text{ elements}\};
        \]
        then $U_r$ is an open subset of $\Spec A$.
  \item If $M$ is a module of finite presentation then the set
        \[
          U = \{p \in \Spec A \mid M_p \text{ is a free } A_p\text{-module}\}
        \]
        is open in $\Spec A$.
\end{enumerate}
\end{theorem}

\begin{proof}
(i) Suppose that $M_p = A_p\omega_1 + \cdots + A_p\omega_r$. Each $\omega_i$ is of the form
$\omega_i = m_i/s_i$ with $s_i \in A - p$ and $m_i \in M$, but since $s_i$ is a unit of $A_p$, we can
replace $\omega_i$ by $m_i$, and so assume that $\omega_i$ is (the image in $M_p$ of) an element
of $M$. Define a linear map $\varphi\colon A^r \to M$ from the direct sum of $r$ copies of
$A$ to $M$ by $(a_1, \ldots, a_r) \mapsto \sum a_i\omega_i$, and write $C$ for the cokernel of $\varphi$. Localising
the exact sequence $A^r \to M \to C \to 0$ at a prime ideal $q$, we get an
exact sequence
\[
  A_q^r \to M_q \to C_q \to 0,
\]
and when $q = p$ we get $C_p = 0$. $C$ is a quotient of $M$, so is finitely generated,
so that the support $\Supp(C)$ is a closed set, and hence there is an open
neighbourhood $V$ of $p$ such that $C_q = 0$ for $q \in V$. This means that $V \subset U_r$.
(In short, if $\omega_1, \ldots, \omega_r \in M$ generate $M_p$ at $p$ then they generate $M_q$ for all $q$
in a neighbourhood of $p$.)

(ii) Suppose that $M_p$ is a free $A_p$-module, and let $\omega_1, \ldots, \omega_r$ be a basis.
As in (i) there is no loss of generality in assuming that $\omega_i \in M$. Moreover,
if we choose a suitable $D(a)$ as a neighbourhood of $p$ in $\Spec A$, $\omega_1, \ldots, \omega_r$
generate $M_q$ for every $q \in D(a)$. Thus, replacing $A$ by $A_a$ and $M$ by $M_a$,
we can assume that the elements $\omega_1, \ldots, \omega_r$ satisfy $M_q = \sum A_q\omega_i$ for every
prime ideal $q$ of $A$. Then by Theorem 6,
\[
  M/\sum A\omega_i = 0, \quad \text{that is } M = A\omega_1 + \cdots + A\omega_r.
\]
(We think of replacing $A$ by $A_a$ as shrinking $\Spec A$ down to the
neighbourhood $D(a)$ of $p$.) Now, defining $\varphi\colon A^r \to M$ as above, and letting
$K$ be its kernel, we get the exact sequence
\[
  0 \to K \to A^r \to M \to 0;
\]
moreover, $K_p = 0$. By Theorem~\ref{thm:2.6}, $K$ is finitely generated, so that applying
(i) with $r = 0$, we have that $K_q = 0$ for every $q$ in a neighbourhood $V$ of
$p$; this gives $(A_q)^r \cong M_q$, so that $V \subset U$.
\end{proof}

% ============ 习题 ============
\section*{Exercises to \S\ref{sec:4}}
\emph{Prove the following propositions.}

\begin{exercise}\label{ex:4.1-ex}
The radical of a primary ideal is prime; also, if $I$ is a proper ideal containing
a power $m^\nu$ of a maximal ideal $m$ then $I$ is primary and $\sqrt{I} = m$.
\end{exercise}

\begin{exercise}\label{ex:4.2-ex}
If $P$ is a prime ideal of a ring $A$ then the \emph{symbolic $n$th power} of $P$ is the ideal
$P^{(n)}$ given by
\[
  P^{(n)} = P^nA_P \cap A.
\]
This is a primary ideal with radical $P$.
\end{exercise}

\begin{exercise}\label{ex:4.3-ex}
If $S$ is a multiplicative set of a ring $A$ then $\Spec(A_S)$ is homeomorphic to
the subspace $\{p \mid p \cap S = \emptyset\} \subset \Spec A$; this is in general neither open nor
closed in $\Spec A$.
\end{exercise}

\begin{exercise}\label{ex:4.4-ex}
If $I$ is an ideal of $A$ then $\Spec(A/I)$ is homeomorphic to the closed subset
$V(I)$ of $\Spec A$.
\end{exercise}

\begin{exercise}\label{ex:4.5-ex}
The spectrum of a ring $\Spec A$ is quasi-compact, that is, given an open
covering $\{U_\lambda\}_{\lambda \in \Lambda}$ of $X = \Spec A$ (with $X = \bigcup_\lambda U_\lambda$), a finite number of the
$U_\lambda$ already cover $X$.
\end{exercise}

\begin{exercise}\label{ex:4.6-ex}
If $\Spec A$ is disconnected then $A$ contains an idempotent $e$ (an element $e$
satisfying $e^2 = e$) distinct from $0$ and $1$.
\end{exercise}

\begin{exercise}\label{ex:4.7-ex}
If $A$ and $B$ are rings then $\Spec(A \times B)$ can be identified with the disjoint
union $\Spec A \sqcup \Spec B$, with both of these open and closed in
$\Spec(A \times B)$.
\end{exercise}

\begin{exercise}\label{ex:4.8-ex}
If $M$ is an $A$-module, $N$ and $N'$ submodules of $M$, and $S \subset A$ a
multiplicative set, then $N_S \cap N'_S = (N \cap N')_S$, where both sides are consi-
dered as subsets of $M_S$.
\end{exercise}

\begin{exercise}\label{ex:4.9-ex}
A topological space is said to be \emph{Noetherian} if the closed sets satisfy the
descending chain condition. If $A$ is a Noetherian ring then $\Spec A$ is a
Noetherian topological space. (Note that the converse is not true in
general.)
\end{exercise}

\begin{exercise}\label{ex:4.10-ex}
We say that a non-empty closed set $V$ in a topological space is \emph{reducible} if
it can be expressed as a union $V = V_1 \cup V_2$ of two strictly smaller closed
sets $V_1$ and $V_2$, and \emph{irreducible} if it does not have any such expression. If
$p \in \Spec A$ then $V(p)$ is an irreducible closed set, and conversely every
irreducible closed set of $\Spec A$ can be written as $V(p)$ for some $p \in \Spec A$.
\end{exercise}

\begin{exercise}\label{ex:4.11-ex}
Any closed subset of a Noetherian topological space can be written as a
union of finitely many irreducible closed sets.
\end{exercise}

\begin{exercise}\label{ex:4.12-ex}
Use the results of the previous two exercises to prove the following: for $I$ a
proper ideal of a Noetherian ring, the set of prime ideals containing $I$ has
only finitely many minimal elements.
\end{exercise}

% ============ §5 The Hilbert Nullstellensatz and first steps in dimension theory (原书 pp.30-37) ============
\section{The Hilbert Nullstellensatz and first steps in dimension theory}\label{sec:5}\index{Hilbert Nullstellensatz, section 5}\index{dimension theory, section 5}

Let $X$ be a topological space; we consider strictly decreasing (or
strictly increasing) chains $Z_0 \supsetneq Z_1 \supsetneq \cdots \supsetneq Z_r$ of length $r$ of irreducible closed
subsets of $X$. The supremum of the lengths, taken over all such chains, is
called the \emph{combinatorial dimension}\index{combinatorial dimension} of $X$ and denoted $\dim X$. If $X$ is a
Noetherian space then there are no infinite strictly decreasing chains, but it
can nevertheless happen that $\dim X = \infty$.

Let $Y$ be a subspace of $X$. If $S \subset Y$ is an irreducible closed subset of $Y$ then
its closure in $X$ is an irreducible closed subset $S' \subset X$ such that $S' \cap Y = S$.
Indeed, if $S' = V \cup W$ with $V$ and $W$ closed in $X$ then $S = (V \cap Y) \cup (W \cap Y)$,
so that say $S = V \cap Y$, but then $V = S'$. It follows easily from this that
$\dim Y \le \dim X$.

Let $A$ be a ring. The supremum of the lengths $r$, taken over all strictly
decreasing chains $p_0 \supsetneq p_1 \supsetneq \cdots \supsetneq p_r$ of prime ideals of $A$, is called the \emph{Krull
dimension}\index{Krull dimension}, or simply the \emph{dimension} of $A$, and denoted $\dim A$. As one sees
easily from Ex.~\ref{ex:4.10-ex}, the Krull dimension of $A$ is the same thing as the
combinatorial dimension of $\Spec A$. For a prime ideal $p$ of $A$, the supremum
of the lengths, taken over all strictly decreasing chains of prime ideals
$p = p_0 \supsetneq p_1 \supsetneq \cdots \supsetneq p_r$ starting from $p$, is called the \emph{height} of $p$,\index{height of a prime ideal} and denoted
$\htop p$; (if $A$ is Noetherian it will be proved in Theorem~\ref{thm:13.5} that $\htop p < \infty$).
Moreover, the supremum of the lengths, taken over all strictly increasing
chains of prime ideals $p = p_0 \subsetneq p_1 \subsetneq \cdots \subsetneq p_r$ starting from $p$, is called the
\emph{coheight} of $p$, and written $\coht p$. It follows from the definitions that
\[
  \htop p = \dim A_p, \quad \coht p = \dim A/p \quad \text{and} \quad \htop p + \coht p \le \dim A.
\]

\begin{remark}
In more old-fashioned terminology $\htop p$ was usually called the
\emph{rank} of $p$, and $\coht p$ the \emph{dimension} of $p$; in addition, Nagata [N1] calls
$\dim A$ the \emph{altitude} of $A$.
\end{remark}

\begin{example}\label{ex:5.1}
The prime ideals in the ring $\Z$ of rational integers are the ideals
$p\Z$ generated by the primes $p = 2, 3, 5, \ldots$, together with $(0)$. Hence, every $p\Z$
is a maximal ideal, and $\dim \Z = 1$. More generally, any principal ideal
domain which is not a field is one-dimensional.
\end{example}

\begin{example}\label{ex:5.2}
An Artinian ring is zero-dimensional; indeed, we have seen in
the proof of Theorem~\ref{thm:3.2} that there are only a finite number of maximal
ideals $p_1, \ldots, p_n$, and that the product of all of these is nilpotent. If then $p$ is a
prime ideal, $p \supset (0) = (p_1 \cdots p_n)^\nu$ so that $p \supset p_i$ for some $i$; hence, $p = p_i$, so
that every prime ideal is maximal.
\end{example}

\begin{example}\label{ex:5.3}
A zero-dimensional integral domain is just a field.
\end{example}

\begin{example}\label{ex:5.4}
The polynomial ring $k[X_1, \ldots, X_n]$ over a field $k$ is an integral
domain, and since
\[
  k[X_1, \ldots, X_n]/(X_1, \ldots, X_i) \cong k[X_{i+1}, \ldots, X_n],
\]
$(X_1, \ldots, X_i)$ is a prime ideal of $k[X_1, \ldots, X_n]$. Thus
\[
  (0) \subset (X_1) \subset (X_1, X_2) \subset \cdots \subset (X_1, \ldots, X_n)
\]
is a chain of prime ideals of length $n$, and $\dim k[X_1, \ldots, X_n] \ge n$. In fact we
will shortly be proving that equality holds.
\end{example}

For an ideal $I$ of a ring $A$ we define the \emph{height of $I$} to be the infimum of the
heights of prime ideals containing $I$:
\[
  \htop I = \inf\{\htop p \mid I \subset p \in \Spec A\}.
\]
Here also we have the inequality
\[
  \htop I + \dim A/I \le \dim A.
\]
If $M$ is an $A$-module we define the \emph{dimension of $M$} by
\[
  \dim M = \dim(A/\ann(M)).
\]
If $M$ is finitely generated then $\dim M$ is the combinatorial dimension of the
closed subspace $\Supp(M) = V(\ann(M))$ of $\Spec A$.

A strictly increasing (or decreasing) chain $p_0, p_1, \ldots$ of prime ideals is said
to be \emph{saturated}\index{saturated chain} if there do not exist prime ideals strictly contained between
any two consecutive terms. We say that $A$ is a \emph{catenary ring}\index{catenary ring} if the following
condition is satisfied: for any prime ideals $p$ and $p'$ of $A$ with $p \subset p'$, there
exists a saturated chain of prime ideals starting from $p$ and ending at $p'$, and
all such chains have the same (finite) length.

If a local domain $(A, m)$ is catenary then for any prime ideal $p$ we have
$\htop p + \coht p = \dim A$. Conversely, if $A$ is a Noetherian local domain and
this equality holds for all $p$ then $A$ is catenary (Ratliff [3], 1972); the proof
of this is difficult, and we postpone it to Theorem 31.4. Practically all the
important Noetherian rings arising in applications are known to be
catenary; the first example of a non-catenary Noetherian ring was
discovered in 1956 by Nagata [15].

We now spend some time discussing the elementary theory of dimensions
of rings which are finitely generated over a field $k$.

\begin{theorem}\label{thm:5.1}
Let $k$ be a field, $L$ an algebraic extension of $k$ and
$a_1, \ldots, a_n \in L$; then
\begin{enumerate}[(i)]
  \item $k[a_1, \ldots, a_n] = k(a_1, \ldots, a_n)$.
  \item Write $\varphi\colon k[X_1, \ldots, X_n] \to k[a_1, \ldots, a_n]$ for the homomorphism
        over $k$ which maps $X_i$ to $a_i$; then $\Ker \varphi$ is the maximal ideal generated
        by $n$ elements of the form $f_1(X_1), f_2(X_1, X_2), \ldots, f_n(X_1, \ldots, X_n)$, where
        each $f_i$ can be taken to be monic in $X_i$.
\end{enumerate}
\end{theorem}

\begin{proof}
Let $f_1(X)$ be the minimal polynomial of $a_1$ over $k$; then $(f_1(X_1))$ is a
maximal ideal of $k[X_1]$, so that $k[a_1] \cong k[X_1]/(f_1(X_1))$ is a field, and hence
$k[a_1] = k(a_1)$. Now let $\varphi_2(X)$ be the minimal polynomial of $a_2$ over $k(a_1)$;
then since $k(a_1) = k[a_1]$, the coefficients of $\varphi_2$ can be expressed as
polynomials in $a_1$, and there is a polynomial $f_2 \in k[X_1, X_2]$, monic in $X_2$,
such that $\varphi_2(X_2) = f_2(a_1, X_2)$. Thus
\[
  k[a_1, a_2] = k[a_1][a_2] = k(a_1)[a_2] \cong k(a_1)[X_2]/(f_2(a_1, X_2)).
\]
Proceeding in the same way, for $1 \le i \le n$ there is an
$f_i(X_1, \ldots, X_i) \in k[X_1, \ldots, X_i]$, monic in $X_i$, such that
\[
  k[a_1, \ldots, a_i] = k(a_1, \ldots, a_i)
\]
\[
  \cong k(a_1, \ldots, a_{i-1})[X_i]/(f_i(a_1, \ldots, a_{i-1}, X_i)).
\]

Now if $P(X) \in k[X_1, \ldots, X_n]$ is in the kernel of $\varphi$, we have $\varphi(P) =
P(a_1, \ldots, a_n) = 0$, so that $P(a_1, \ldots, a_{n-1}, X_n)$ is divisible by $f_n(a_1, \ldots,
a_{n-1}, X_n)$; dividing $P(X_1, \ldots, X_n)$ as a polynomial in $X_n$ by the
monic polynomial $f_n(X_1, \ldots, X_n)$ and letting $R_n(X_1, \ldots, X_n)$ be the remain-
der, we can write $P = Q_n f_n + R_n$, with $R_n(a_1, \ldots, a_{n-1}, X_n) = 0$. Similarly,
dividing $R_n(X_1, \ldots, X_n)$ as a polynomial in $X_{n-1}$ by $f_{n-1}(X_1, \ldots, X_{n-1})$ and
letting $R_{n-1}(X_1, \ldots, X_n)$ be the remainder, we get
\[
  R_n = Q_{n-1}f_{n-1} + R_{n-1},
\]
with
\[
  R_{n-1}(a_1, \ldots, a_{n-2}, X_{n-1}, X_n) = 0.
\]
Proceeding in the same way we get $P = \sum Q_i f_i + R$, with $R(X_1, \ldots, X_n) = 0$;
that is $R = 0$ and $P = \sum Q_i f_i$, so that $\Ker \varphi = (f_1, f_2, \ldots, f_n)$.
\end{proof}

The following theorem can be regarded as a converse of Theorem 1, (i).

\begin{theorem}\label{thm:5.2}
Let $k$ be a field and $A = k[\alpha_1, \ldots, \alpha_n]$ an integral domain, and
write $r = \trdeg_k A$ for the transcendence degree of $A$ (that is, of its field of
fractions) over $k$. Then if $r > 0$, $A$ is not a field.
\end{theorem}

\begin{proof}
Suppose that $\alpha_1, \ldots, \alpha_r$ is a transcendence basis for $A$ over $k$, and set
$K = k(\alpha_1, \ldots, \alpha_r)$. Then since $\alpha_{r+1}, \ldots, \alpha_n$ are algebraic over $K$, there exist
polynomials $f_i(X_{r+1}, \ldots, X_i) \in K[X_{r+1}, \ldots, X_i]$, monic of degree $d_i$ in $X_i$, such that
\[
  k[\alpha_{r+1}, \ldots, \alpha_n] \cong K[X_{r+1}, \ldots, X_n]/(f_{r+1}, \ldots, f_n)
\]
and
\[
  d_i = [K(\alpha_{r+1}, \ldots, \alpha_i) : K(\alpha_{r+1}, \ldots, \alpha_{i-1})] \quad \text{for } i > r.
\]
The coefficients of $f_i$ are in $K$, so that for suitable $0 \neq g \in k[\alpha_1, \ldots, \alpha_r]$ we
have $gf_i \in k[\alpha_1, \ldots, \alpha_r][X_{r+1}, \ldots, X_n]$. In other words, if we set $B =
k[\alpha_1, \ldots, \alpha_r, g^{-1}]$ then the $f_i$ are polynomials with coefficients in $B$. We
are now going to show that $A[g^{-1}] = B[\alpha_{r+1}, \ldots, \alpha_n]$ is a free module over
$B$ with basis $\{\prod_{i>r} \alpha_i^{e_i} \mid 0 \le e_i < d_i\}$. Every element of $B[\alpha_{r+1}, \ldots, \alpha_n]$ can
be written as $P(\alpha_{r+1}, \ldots, \alpha_n)$ for some $P \in B[X_{r+1}, \ldots, X_n]$; dividing $P$ as a
polynomial in $X_n$ by $f_n$, and replacing $P$ by the remainder, we can assume
that $P$ has degree at most $d_n - 1$ in $X_n$; then dividing $P$ as a polynomial
in $X_{n-1}$ by $f_{n-1}$, and replacing by the remainder, we can assume that $P$
has degree at most $d_{n-1} - 1$ in $X_{n-1}$. Proceeding in the same way, we
can assume that $P$ has degree at most $d_i - 1$ in $X_i$ for each $i$; in addition,
the elements $\{\alpha_i^e \mid 0 \le e < d_i\}$ are linearly independent over $K(\alpha_{r+1}, \ldots, \alpha_{i-1})$. Hence $A[g^{-1}]$ is a free $B$-module. However, $B$ is not a field; for
$k[a_1, \ldots, a_r]$ is a polynomial ring in $r$ variables over $k$, and hence it contains
infinitely many irreducible polynomials (the proof of this is exactly the
same as Euclid's proof that there exist infinitely many primes). Hence,
there is an irreducible polynomial $h \in k[a_1, \ldots, a_r]$ which does not divide
$g$, and then obviously $h^{-1} \notin k[a_1, \ldots, a_r, g^{-1}]$. Therefore $B$ contains an
ideal $I$ with $I \neq 0, B$, and since $A[g^{-1}]$ is a free module over $B$, $IA[g^{-1}]$
is a proper ideal of $A[g^{-1}]$. Thus $A[g^{-1}]$ is not a field. But if $A$ were a
field then we would have $A[g^{-1}] = A$, and hence $A$ is not a field.
\end{proof}

\begin{theorem}\label{thm:5.3}
Let $k$ be a field, and let $m$ be any maximal ideal of the
polynomial ring $k[X_1, \ldots, X_n]$; then the residue class field $k[X_1, \ldots, X_n]/m$
is algebraic over $k$. Hence $m$ can be generated by $n$ elements, and in
particular if $k$ is algebraically closed then $m$ is of the form $m =
(X_1 - a_1, \ldots, X_n - a_n)$ for $a_i \in k$.
\end{theorem}

\begin{proof}
Set $k[X_1, \ldots, X_n]/m = K$, and write $\alpha_i$ for the image of $X_i$ in $K$; then
$K = k[\alpha_1, \ldots, \alpha_n]$. By the previous theorem, since $K$ is a field it is algebraic
over $k$, and then by Theorem 1, (ii), $m$ is generated by $n$ elements. If $k$ is
algebraically closed then $k = K$, so that each $X_i$ is congruent modulo $m$ to
some $a_i \in k$; then $(X_1 - a_1, \ldots, X_n - a_n) \subset m$. On the other hand
$(X_1 - a_1, \ldots, X_n - a_n)$ is obviously a maximal ideal, so that equality must
hold.
\end{proof}

Let $k$ be a field and $\bar{k}$ its algebraic closure. Suppose that $\Phi \subset
k[X_1, \ldots, X_n]$ is a subset. An $n$-tuple $\alpha = (a_1, \ldots, a_n)$ of elements $a_i \in \bar{k}$ is an
\emph{algebraic zero}\index{algebraic zero} of $\Phi$ if it satisfies $f(\alpha) = 0$ for every $f(X) \in \Phi$.

\begin{theorem}[The Hilbert Nullstellensatz]\label{thm:5.4}
\begin{enumerate}[(i)]
  \item If $\Phi$ is a subset of $k[X_1, \ldots, X_n]$ which does not have any algebraic
        zeros then the ideal generated by $\Phi$ contains $1$.
  \item Given a subset $\Phi$ of $k[X_1, \ldots, X_n]$ and an element $f \in k[X_1, \ldots, X_n]$,
        suppose that $f$ vanishes at every algebraic zero of $\Phi$. Then some power of $f$
        belongs to the ideal generated by $\Phi$, that is there exist $\nu > 0$,
        $g_i \in k[X_1, \ldots, X_n]$ and $h_i \in \Phi$ such that $f^\nu = \sum g_i h_i$.
\end{enumerate}
\end{theorem}

\begin{proof}
(i) Let $I$ be the ideal generated by $\Phi$; if $1 \notin I$ then there exists a
maximal ideal $m$ containing $I$. By the previous theorem, $k[X_1, \ldots, X_n]/m$
is algebraic over $k$, so that it has a $k$-linear isomorphic embedding $\theta$ into
$\bar{k}$. If we set $\theta(X_i \bmod m) = a_i$ then for all $g(X) \in m$ we have $0 = \theta(g(X)) =
g(a)$, therefore $\alpha = (a_1, \ldots, a_n)$ is an algebraic zero of $m$, and
hence also of $\Phi$. This contradicts the hypothesis; hence, $1 \in I$.

(ii) Inside $k[X_1, \ldots, X_n, Y]$ we consider the set $\Phi \cup \{1 - Yf(X)\}$; then this
set has no algebraic zeros, so that by (i) it generates the ideal $(1)$. Therefore
there exists a relation of the form
\[
  1 = \sum p_i(X, Y)h_i(X) + Q(X, Y)(1 - Yf(X)),
\]
with $h_i(X) \in \Phi$. This is an identity in $X_1, \ldots, X_n$ and $Y$, so that it still holds if
we substitute $Y = f(X)^{-1}$. Hence we have
\[
  1 = \sum p_i(X, f^{-1})h_i(X),
\]
so that multiplying through by a suitable power of $f$ and clearing
denominators gives $f^\nu = \sum g_i(X)h_i(X)$, with $g_i \in k[X_1, \ldots, X_n]$ and
$h_i \in \Phi$.
\end{proof}

\begin{remark}
The above proof of (ii) is a classical idea due to Rabinowitch
[1]. In a modern form it can be given as follows: let $I \subset k[X_1, \ldots, X_n] = A$
be the ideal generated by $\Phi$; then in the localisation $A_f$ with respect to $f$ (see
\S\ref{sec:4}, Example 1), we have $IA_f = A_f$, so that a power of $f$ is in $I$.
\end{remark}

\begin{theorem}\label{thm:5.5}
Let $k$ be a field, $A$ a ring which is finitely generated over $k$, and
$I$ a proper ideal of $A$; then the radical of $I$ is the intersection of all maximal
ideals containing $I$, that is $\sqrt{I} = \bigcap_{I \subset m} m$.
\end{theorem}

\begin{proof}
Let $A = k[a_1, \ldots, a_n]$, so that $A$ is a quotient of $k[X_1, \ldots, X_n]$.
Considering the inverse image of $I$ in $k[X]$ reduces to the case $A = k[X]$,
and the assertion follows from Theorem 4, (ii).
\end{proof}

Compared to the result $\sqrt{I} = \bigcap_{I \subset P} P$ proved in \S\ref{sec:1}, the conclusion of
Theorem 5 is much stronger. It is equivalent to the condition on a ring
that every prime ideal $P$ should be expressible as an intersection of maximal
ideals. Rings for which this holds are called \emph{Hilbert rings}\index{Hilbert ring} or \emph{Jacobson rings}\index{Jacobson ring},
and they have been studied independently by O. Goldmann [1] and
W. Krull [7]. See also Kaplansky [K] and Bourbaki [B-5].

\begin{theorem}\label{thm:5.6}
Let $k$ be a field and $A$ an integral domain which is finitely
generated over $k$; then
\[
  \dim A = \trdeg_k A.
\]
\end{theorem}

\begin{proof}
Let $A = k[X_1, \ldots, X_n]/P$, and set $r = \trdeg_k A$. To prove that
$r \ge \dim A$ it is enough to show that if $P$ and $Q$ are prime ideals of
$k[X] = k[X_1, \ldots, X_n]$ with $Q \supset P$ and $Q \neq P$ then
\[
  \trdeg_k k[X]/Q < \trdeg_k k[X]/P.
\]
The $k$-algebra homomorphism $k[X]/P \to k[X]/Q$ is onto, so that $\trdeg_k
k[X]/Q \le \trdeg_k k[X]/P$ is obvious. Suppose that equality holds. Let
$k[X]/P = k[\alpha_1, \ldots, \alpha_n]$ and $k[X]/Q = k[\beta_1, \ldots, \beta_n]$; we can assume that
$\beta_1, \ldots, \beta_r$ is a transcendence basis for $k(\beta)$ over $k$. Then $\alpha_1, \ldots, \alpha_r$ are also
algebraically independent over $k$, so that they form a transcendence basis
for $k(\alpha)$ over $k$. Now set $S = k[X_1, \ldots, X_n] - \{0\}$; $S$ is a multiplicative set in
$k[X]$ with $P \cap S = \emptyset$ and $Q \cap S = \emptyset$. Setting $R = k[X_1, \ldots, X_n]$ and
$K = k(X_1, \ldots, X_n)$, we have $R_S = K[X_{r+1}, \ldots, X_n]$, and
\[
  R_S/PR_S \cong k(\alpha_1, \ldots, \alpha_r)[\alpha_{r+1}, \ldots, \alpha_n],
\]
so that $R_S/PR_S$ is algebraic over $K = k(X_1, \ldots, X_n) \cong k(\alpha_1, \ldots, \alpha_r)$, and
therefore by Theorem 1, $PR_S$ is a maximal ideal of $R_S$; but this contradicts
the assumptions $P \subset Q$ with $P \neq Q$ and $Q \cap S = \emptyset$.

Now let us prove that $r \le \dim A$ by induction on $r$. If $r = 0$ then, by
Theorem 1, $A$ is a field, so $\dim A = 0$ and the assertion holds. Now let $r > 0$,
and suppose that $A = k[\alpha_1, \ldots, \alpha_n]$ with $\alpha_1$ transcendental over $k$; setting
$S = k[X_1] - \{0\}$ and $R = k[X_1, \ldots, X_n]$ we get
\[
  R_S = k(X_1)[X_2, \ldots, X_n] \quad \text{and} \quad R_S/PR_S \cong k(X_1)[\alpha_2, \ldots, \alpha_n].
\]
Hence $R_S/PR_S$ has transcendence degree $r - 1$ over $k(X_1)$, so that by
induction $\dim R_S/PR_S \ge r - 1$. Thus there exists a strictly increasing chain
$PR_S = Q_0 \subsetneq Q_1 \subsetneq \cdots \subsetneq Q_{r-1}$ of prime ideals of $R_S$. If we set $P_i = Q_i \cap R$ then
$P_i$ is a prime ideal of $R$ disjoint from $S$; in particular, the residue class of
$X_1$ in $R/P_{r-1}$ is not algebraic over $k$, and so $\trdeg_k R/P_{r-1} \ge 1$. Then
$P_{r-1}$ is not a maximal ideal of $R$ by Theorem 3, and therefore $R$ has a
maximal ideal $P_r$ strictly bigger than $P_{r-1}$. Hence $\dim A = \coht P \ge r$.
\end{proof}

\begin{corollary*}
If $k$ is a field then $\dim k[X_1, \ldots, X_n] = n$.
\end{corollary*}

We now turn to a different topic, the theorem of Forster and Swan on the
number of generators of a module. Let $A$ be a ring and $M$ a finite $A$-module;
for $p \in \Spec A$, write $\kappa(p)$ for the residue field of $A_p$, and let $\mu(p, M)$ denote the
dimension over $\kappa(p)$ of the vector space $M \otimes \kappa(p) = M_p/pM_p$ (in the usual
sense of linear algebra). This is the cardinality of a minimal basis of the
$A_p$-module $M_p$. Hence, if $p \supset p'$ then $\mu(p, M) \ge \mu(p', M)$.

In 1964 the young function-theorist O. Forster surprised the experts in
algebra by proving the following theorem [1].

\begin{theorem}\label{thm:5.7}
Let $A$ be a Noetherian ring and $M$ a finite $A$-module. Set
\[
  b(M) = \sup\{\mu(p, M) + \coht p \mid p \in \Supp M\};
\]
then $M$ can be generated by at most $b(M)$ elements.
\end{theorem}

This theorem is a very important link between the number of local and
global generators. However, there was room for improvement in the bound
for the number of generators, and in no time R. Swan obtained a better
bound. We will prove Swan's bound. For this we need the concept of \emph{j-Spec
A}\index{j-Spec of a ring} introduced by Swan. This is a space having the same irreducible closed
subsets as m-Spec $A$, but has the advantage of having a `generic point', not
present in m-Spec $A$, for every irreducible closed subset.

A prime ideal which can be expressed as an intersection of maximal
ideals is called a \emph{$j$-prime ideal}\index{j-prime ideal}, and we write j-Spec $A$ for the set of all
$j$-prime ideals. We consider j-Spec $A$ also with its topology as a subspace
of $\Spec A$. Set $M = \text{m-Spec } A$ and $J = \text{j-Spec } A$. If $F$ is a closed subset of
$J$ then there is an ideal $I$ of $A$ such that $F = V(I) \cap J$. One sees easily that
a prime ideal $P$ belongs to $F$ if and only if $P$ can be expressed as an
intersection of elements of $F \cap M = V(I) \cap M$. Hence $F$ is determined by
$F \cap M$, so that there is a natural one-to-one correspondence between
closed subsets of $J$ and of $M$. It follows that if $M$ is Noetherian so is $J$,
and they both have the same combinatorial dimension. Now let $B$ be an
irreducible closed subset of $J$, and let $P$ be the intersection of all the
elements of $B$. If $B = V(I) \cap J$ then $I \subset P$ and we can also write
$B = V(P) \cap J$. If $P$ is not a prime ideal then there exist $f, g \in A$ such that
$f \notin P$, $g \notin P$ and $fg \in P$; but then
\[
  B = (V(P + fA) \cap J) \cup (V(P + gA) \cap J),
\]
and by definition of $P$ there is a $Q \in B$ not containing $f$ and a $Q' \in B$ not
containing $g$, which implies that $B$ is reducible, a contradiction. Therefore
$P$ is a prime ideal. Hence $P \in B$ and $B = V(P) \cap J$. This $P$ is called the
\emph{generic point}\index{generic point} of $B$. Conversely if $P$ is any element of $J$ then $V(P) \cap J$ is
an irreducible closed subset of $J$, and is the closure in $J$ of $\{P\}$. We will
write \emph{$j$-dim $p$} for the combinatorial dimension of $V(p) \cap J$.

For a finite $A$-module $M$ and $p \in J$ we set
\[
  b(p, M) =
  \begin{cases}
    0 & \text{if } M_p = 0 \\
    \text{$j$-dim } p + \mu(p, M) & \text{if } M_p \neq 0.
  \end{cases}
\]

\begin{theorem}[Swan~\cite{Sw1}]\label{thm:5.8}
Let $A$ be a ring, and suppose that m-Spec $A$ is
a Noetherian space. Let $M$ be a finite $A$-module. If
\[
  \sup\{b(p, M) \mid p \in \text{j-Spec } A\} = r < \infty
\]
then $M$ is generated by at most $r$ elements.
\end{theorem}

\begin{proof}
\textbf{Step 1.} For $p \in \Spec A$ and $x \in M$, we will say that $x$ is \emph{basic} at $p$ if $x$
has non-zero image in $M \otimes \kappa(p)$. It is easy to see that this condition is
equivalent to $\mu(p, M/Ax) = \mu(p, M) - 1$.

\begin{lemma*}
Let $M$ be a finite $A$-module and $p_1, \ldots, p_n \in \Supp(M)$. Then there
exists $x \in M$ which is basic at each of $p_1, \ldots, p_n$.
\end{lemma*}

\begin{proof}[Proof of the lemma]
By reordering $p_1, \ldots, p_n$ we assume that $p_i$ is maximal among
$\{p_i, p_{i+1}, \ldots, p_n\}$ for each $i$. By induction on $n$ suppose that $x' \in M$ is basic at
$p_1, \ldots, p_{n-1}$. If $x'$ is basic at $p_n$, then we can take $x = x'$. Suppose then that $x'$
is not basic at $p_n$. By assumption $M_{p_n} \neq 0$ so that we can choose some $y \in M$
which is basic at $p_n$. We have $p_1 \cdots p_{n-1} \not\subset p_n$, so that if we take an element
$a \in p_1 \cdots p_{n-1}$ not belonging to $p_n$, and set $x = x' + ay$, this $x$ satisfies our
requirements. This proves the lemma.
\end{proof}

\textbf{Step 2.} Setting $\sup\{b(p, M) \mid p \in \text{j-Spec } A\} = r$, we now show that there are
just a finite number of primes $p$ such that $b(p, M) = r$. Indeed, for
$n = 1, 2, \ldots$, the subset $X_n = \{p \in \text{j-Spec } A \mid \mu(p, M) \ge n\}$ is closed in j-Spec $A$ by
Theorem~\ref{thm:4.10}; it has a finite number of irreducible components (by
Ex.~\ref{ex:4.11-ex}), and we let $p_{ni}$ (for $1 \le i \le v_n$) be their generic points. If $M$ is
generated by $s$ elements then $X_n = \emptyset$ for $n > s$, so that the set $\{p_{ni}\}_{n,i}$ is
finite. Let us prove that if $b(p, M) = r$ then $p \in \{p_{ni}\}_{n,i}$. Suppose
$\mu(p, M) = n$; then $p \in X_n$, so that by construction $p \supset p_{ni}$ for some $i$. But
if $p \neq p_{ni}$ then j-dim $p$ < j-dim $p_{ni}$, and since $\mu(p, M) = n = \mu(p_{ni}, M)$ we have
$b(p, M) < b(p_{ni}, M)$, which is a contradiction. Hence $p = p_{ni}$.

\textbf{Step 3.} Let us choose an element $x \in M$ which is basic for each of the
finitely many primes $p$ with $b(p, M) = r$, and set $\bar{M} = M/Ax$; then clearly
$b(p, \bar{M}) \le r - 1$ for every $p \in \text{j-Spec } A$. Hence by induction $\bar{M}$ is generated by
$r - 1$ elements, and therefore $M$ by $r$ elements.
\end{proof}

Swan's paper contains a proof of the following generalisation to non-
commutative rings: let $A$ be a commutative ring, $\Lambda$ a possibly non-
commutative $A$-algebra and $M$ a finite left $\Lambda$-module. Suppose that
m-Spec $A$ is Noetherian, and that for every maximal ideal $p$ of $A$ the
$\Lambda_p$-module $M_p$ is generated by at most $r$ elements; then $M$ is generated as a
$\Lambda$-module by at most $r + d$ elements, where $d$ is the combinatorial
dimension of m-Spec $A$.

The Forster--Swan theorem is a statement that local properties imply
global ones; remarkable results in this direction have been obtained by
Mohan Kumar [2] (see also Cowsik--Nori [1] and Eisenbud--Evans [1],
[2]). The number of generators of ideals in local rings is the subject of
a nice book by J. Sally [Sa].

% ============ 习题 ============
\section*{Exercises to \S\ref{sec:5}}
\emph{Prove the following propositions.}

\begin{exercise}\label{ex:5.1-ex}
Let $k$ be a field, $R = k[X_1, \ldots, X_n]$ and let $P \in \Spec R$; then $\htop P +
\coht P = n$.
\end{exercise}

\begin{exercise}\label{ex:5.2-ex}
A zero-dimensional Noetherian ring is Artinian (the converse to Example
2 above).
\end{exercise}

% ============ §6 Associated primes and primary decomposition (原书 pp.37-42) ============
\section{Associated primes and primary decomposition}\label{sec:6}\index{associated primes, section 6}\index{primary decomposition, section 6}

Most readers will presumably have come across primary decom-
position of ideals in Noetherian rings. This was the first big theorem
obtained by Emmy Noether in her abstract treatment of commutative
rings. Nowadays, as exemplified by Bourbaki [B4], the notion of associated
prime is considered more important than primary decomposition itself.

Let $A$ be a ring and $M$ an $A$-module. A prime ideal $P$ of $A$ is called an
\emph{associated prime ideal}\index{associated prime ideal} of $M$ if $P$ is the annihilator $\ann(x)$ of some $x \in M$. The
set of associated primes of $M$ is written $\Ass(M)$ or $\Ass_A(M)$. For $I$ an ideal
of $A$, the associated primes of the $A$-module $A/I$ are referred to as the
\emph{prime divisors}\index{prime divisor} of $I$. We say that $a \in A$ is a \emph{zero-divisor} for $M$ if there is a non-zero
$x \in M$ such that $ax = 0$, and otherwise that $a$ is \emph{$M$-regular}.\index{regular element}

\begin{theorem}\label{thm:6.1}
Let $A$ be a Noetherian ring and $M$ a non-zero $A$-module.
\begin{enumerate}[(i)]
  \item Every maximal element of the family of ideals $\mathcal{F} = \{\ann(x) \mid 0 \neq x \in M\}$
        is an associated prime of $M$, and in particular $\Ass(M) \neq \emptyset$.
  \item The set of zero-divisors for $M$ is the union of all the associated primes
        of $M$.
\end{enumerate}
\end{theorem}

\begin{proof}
(i) We have to prove that if $\ann(x)$ is a maximal element of $\mathcal{F}$ then it is
prime: if $a, b \in A$ are such that $abx = 0$ but $bx \neq 0$ then by maximality
$\ann(bx) = \ann(x)$; hence, $ax = 0$.

(ii) If $ax = 0$ for some $x \neq 0$ then $a \in \ann(x) \in \mathcal{F}$, and by (i) there is an
associated prime of $M$ containing $\ann(x)$.
\end{proof}

\begin{theorem}\label{thm:6.2}
Let $S \subset A$ be a multiplicative set, and $N$ an $A_S$-module.
Viewing $\Spec(A_S)$ as a subset of $\Spec A$, we have $\Ass_A(N) = \Ass_{A_S}(N)$. If $A$
is Noetherian then for an $A$-module $M$ we have $\Ass(M_S) =
\Ass(M) \cap \Spec(A_S)$.
\end{theorem}

\begin{proof}
For $x \in N$ we have $\ann_A(x) = \ann_{A_S}(x) \cap A$, so that if $P \in \Ass_{A_S}(N)$
then $P \cap A \in \Ass_A(N)$. Conversely if $p \in \Ass_A(N)$ and we choose $x \in N$
such that $p = \ann_A(x)$ then $x \neq 0$, and hence, $p \cap S = \emptyset$ and $pA_S$ is a prime
ideal of $A_S$ with $pA_S = \ann_{A_S}(x)$. For the second part, if $p \in \Ass(M) \cap
\Spec(A_S)$ then $p \cap S = \emptyset$, and $p = \ann_A(x)$ for some $x \in M$; if $(a/s)x = 0$ in
$M_S$ then there is a $t \in S$ such that $tax = 0$ in $M$, and $t \notin p$, $ta \in p$ gives
$a \in p$, so that $\ann_{A_S}(x) = pA_S$ and $pA_S \in \Ass(M_S)$. Conversely, if $P \in
\Ass(M_S)$ then without loss of generality we have $P = \ann_{A_S}(x)$ with $x \in M$.
Setting $p = P \cap A$ we have $P = pA_S$. Now $p$ is finitely generated since $A$ is
Noetherian, and it follows that there exists some $t \in S$ such that $p =
\ann_A(tx)$. Therefore $p \in \Ass_A(M)$.
\end{proof}

\begin{corollary*}
For a Noetherian ring $A$, an $A$-module $M$ and a prime ideal $P$ of
$A$ we have
\[
  P \in \Ass_A(M) \iff PA_P \in \Ass_{A_P}(M_P).
\]
\end{corollary*}

\begin{theorem}\label{thm:6.3}
Let $A$ be a ring and $0 \to M' \to M \to M'' \to 0$ an exact
sequence of $A$-modules; then
\[
  \Ass(M) \subset \Ass(M') \cup \Ass(M'').
\]
\end{theorem}

\begin{proof}
If $P \in \Ass(M)$ then $M$ contains a submodule $N$ isomorphic to $A/P$.
Since $P$ is prime, for any non-zero element $x$ of $N$ we have $\ann(x) = P$.
Therefore if $N \cap M' \neq 0$ we have $P \in \Ass(M')$. If $N \cap M' = 0$ then the image
of $N$ in $M''$ is also isomorphic to $A/P$, so that $P \in \Ass(M'')$.
\end{proof}

\begin{theorem}\label{thm:6.4}
Let $A$ be a Noetherian ring and $M \neq 0$ a finite $A$-module.
Then there exists a chain $0 = M_0 \subset M_1 \subset \cdots \subset M_r = M$ of submodules of $M$
such that for each $i$ we have $M_i/M_{i-1} \cong A/P_i$ with $P_i \in \Spec A$.
\end{theorem}

\begin{proof}
Choose any $P_1 \in \Ass(M)$; then there exists a submodule $M_1$ of $M$
with $M_1 \cong A/P_1$. If $M_1 \neq M$ and we choose any $P_2 \in \Ass(M/M_1)$ then
there exists $M_2 \subset M$ such that $M_2/M_1 \cong A/P_2$. Continuing in the same way and
using the ascending chain condition, we eventually arrive at $M_r = M$.
\end{proof}

\begin{theorem}\label{thm:6.5}
Let $A$ be a Noetherian ring and $M$ a finite $A$-module.
\begin{enumerate}[(i)]
  \item $\Ass(M)$ is a finite set.
  \item $\Ass(M) \subset \Supp(M)$.
  \item The set of minimal elements of $\Ass(M)$ and of $\Supp(M)$ coincide.
\end{enumerate}
\end{theorem}

\begin{proof}
(i) follows from the previous two theorems; we need only note
that $\Ass(A/P) = \{P\}$. For (ii), if $0 \to A/P \to M$ is exact then so is $0 \to
A_P/PA_P \to M_P$, and therefore $M_P \neq 0$. For (iii) it is enough to show that
if $P$ is a minimal element of $\Supp(M)$ then $P \in \Ass(M)$. We have $M_P \neq 0$
so that by Theorem 2 and (ii),
\[
  \emptyset \neq \Ass(M_P) = \Ass(M) \cap \Spec(A_P) \subset \Supp(M) \cap \Spec(A_P)
             = \{P\}.
\]
Therefore we must have $P \in \Ass(M)$.
\end{proof}

Let $A$ be a Noetherian ring and $M$ a finite $A$-module. Let $P_1, \ldots, P_r$ be the
minimal elements of $\Supp(M)$; then $\Supp(M) = V(P_1) \cup \cdots \cup V(P_r)$, and the
$V(P_i)$ are the irreducible components of the closed set $\Supp(M)$ (see
Ex.~\ref{ex:4.11-ex}). The prime ideals $P_1, \ldots, P_r$ are called the \emph{isolated associated primes}\index{isolated associated prime}
of $M$, and the remaining associated primes of $M$ are called \emph{embedded}
\emph{primes}.\index{embedded prime} If $I$ is an ideal of $A$ then $\Supp_A(A/I)$ is the set of prime ideals
containing $I$, and the \emph{minimal prime divisors} of $I$ (that is the minimal
associated primes of the $A$-module $A/I$) are precisely the minimal prime
ideals containing $I$. We have seen in Ex.~\ref{ex:4.12-ex} that there are only a finite
number of such primes, and Theorem 5 now gives a new proof of this.
(For examples of embedded primes see Ex.~\ref{ex:6.6-ex} and Ex.~\ref{ex:8.9-ex}.)

\begin{definition}
Let $A$ be a ring, $M$ an $A$-module and $N \subset M$ a submodule. We
say that $N$ is a \emph{primary submodule}\index{primary submodule} of $M$ if the following condition holds for
all $a \in A$ and $x \in M$:
\[
  x \notin N \text{ and } ax \in N \Rightarrow a^\nu M \subset N \text{ for some } \nu.
\]
This definition in fact only depends on the quotient module $M/N$. It can be
restated as
\[
  \text{if } a \in A \text{ is a zero-divisor for } M/N \text{ then } a \in \sqrt{\ann(M/N)}.
\]
\end{definition}

A primary ideal is just a primary submodule of the $A$-module $A$. One
might wonder about trying to set up a notion of prime submodule
generalising prime ideal, but this does not turn out to be useful.

\begin{theorem}\label{thm:6.6}
Let $A$ be a Noetherian ring and $M$ a finite $A$-module. Then a
submodule $N \subset M$ is primary if and only if $\Ass(M/N)$ consists of one
element only. In this case, if $\Ass(M/N) = \{P\}$ and $\ann(M/N) = I$ then $I$ is
primary and $\sqrt{I} = P$.
\end{theorem}

\begin{proof}
If $\Ass(M/N) = \{P\}$ then by the previous theorem $\Supp(M/N) =
V(P)$, so that $P = \sqrt{\ann(M/N)}$. Now if $a \in A$ is a zero-divisor for $M/N$ it
follows from Theorem 1 that $a \in P$, so that $a \in \sqrt{\ann(M/N)}$; hence, $N$ is a
primary submodule of $M$. Conversely, if $N$ is a primary submodule and
$P \in \Ass(M/N)$ then every $a \in P$ is a zero-divisor for $M/N$, so that by
assumption $a \in \sqrt{I}$, where $I = \ann(M/N)$. Hence $P \subset \sqrt{I}$, but from the
definition of associated prime we obviously have $I \subset P$, and hence $\sqrt{I} \subset P$,
so that $P = \sqrt{I}$. Thus $\Ass(M/N)$ has just one element $\sqrt{I}$. We prove that in
this case $I$ is a primary ideal: let $a, b \in A$ with $b \notin I$; if $ab \in I$ then $ab(M/N) = 0$,
but $b(M/N) \neq 0$, so that $a$ is a zero-divisor for $M/N$, and therefore
$a \in P = \sqrt{I}$.
\end{proof}

\begin{definition}
If $\Ass(M/N) = \{P\}$ we say that $N \subset M$ is a \emph{$P$-primary} sub-
module, or a \emph{primary submodule belonging to $P$}.
\end{definition}

\begin{theorem}\label{thm:6.7}
If $N$ and $N'$ are $P$-primary submodules of $M$ then so is $N \cap N'$.
\end{theorem}

\begin{proof}
We can embed $M/(N \cap N')$ as a submodule of $(M/N) \oplus (M/N')$, so
that
\[
  \Ass(M/(N \cap N')) \subset \Ass(M/N) \cup \Ass(M/N') = \{P\}.
\]
\end{proof}

If $N \subset M$ is a submodule, we say that $N$ is \emph{reducible}\index{reducible submodule} if it can be written as
an intersection $N = N_1 \cap N_2$ of two submodules $N_1, N_2$ with $N_i \neq N$, and
otherwise that $N$ is \emph{irreducible}; note that this has nothing to do with the
notion of irreducible modules in representation theory (= no submodules
other than $0$ and $M$), which is a condition on $M$ only.

If $M$ is a Noetherian module then any submodule $N$ of $M$ can be written
as a finite intersection of irreducible submodules. Proof: let $\mathcal{F}$ be the set
of submodules $N \subset M$ having no such expression. If $\mathcal{F} \neq \emptyset$ then it has a
maximal element $N_0$. Then $N_0$ is reducible, so that $N_0 = N_1 \cap N_2$, and
$N_i \notin \mathcal{F}$. Now each of the $N_i$ is an intersection of a finite number of
irreducible submodules, and hence so is $N_0$. This is a contradiction.

\begin{remark}
The representation as an intersection of irreducible submodules is
in general not unique. For example, if $A$ is a field and $M$ an $n$-dimensional
vector space over $A$ then the irreducible submodules of $M$ are just its
$(n - 1)$-dimensional subspaces. An $(n - 2)$-dimensional subspace can be
written in lots of ways as an intersection of $(n - 1)$-dimensional subspaces.
\end{remark}

In general we say that an expression of a set $N$ as an intersection
$N = N_1 \cap \cdots \cap N_r$ is \emph{irredundant}\index{irredundant decomposition} if we cannot omit any $N_i$, that is if
$N \neq N_1 \cap \cdots \cap N_{i-1} \cap N_{i+1} \cap \cdots \cap N_r$. If $M$ is an $A$-module, we call an
expression $N = N_1 \cap \cdots \cap N_r$ of a submodule $N$ as an intersection of a finite
number of submodules $N_i \subset M$ a \emph{decomposition} of $N$; if each of the $N_i$ is
irreducible we speak of an \emph{irreducible decomposition}, if primary of a \emph{primary decomposition}.
Let $N = N_1 \cap \cdots \cap N_r$ be an irredundant primary decompo-
sition with $\Ass(M/N_i) = \{P_i\}$; if $P_i = P_j$ then $N_i \cap N_j$ is again primary, so
that grouping together all of the $N_i$ belonging to the same prime ideal we get
a primary decomposition such that $P_i \neq P_j$ for $i \neq j$. A decomposition with
this property will be called a \emph{shortest primary decomposition}, and the $N_i$
appearing in it the \emph{primary components} of $N$; if $N_i$ belongs to a prime $P$ we
sometimes say that $N_i$ is the \emph{$P$-primary component} of $N$.

\begin{theorem}\label{thm:6.8}
Let $A$ be a Noetherian ring and $M$ a finite $A$-module.
\begin{enumerate}[(i)]
  \item An irreducible submodule of $M$ is a primary submodule.
  \item If
        \[
          N = N_1 \cap \cdots \cap N_l \quad \text{with } \Ass(M/N_i) = \{P_i\}
        \]
        is an irredundant primary decomposition of a proper submodule $N \subset M$
        then $\Ass(M/N) = \{P_1, \ldots, P_l\}$.
  \item Every proper submodule $N$ of $M$ has a primary decomposition. If $N$ is
        a proper submodule of $M$ and $P$ is a minimal associated prime of $M/N$ then
        the $P$-primary component of $N$ is $\varphi_P^{-1}(N_P)$, where $\varphi_P\colon M \to M_P$ is the
        canonical map, and therefore it is uniquely determined by $M$, $N$ and $P$.
\end{enumerate}
\end{theorem}

\begin{proof}
(i) It is enough to prove that a submodule $N \subset M$ which is not
primary is reducible: replacing $M$ by $M/N$ we can assume that $N = 0$. By
Theorem 6, $\Ass(M)$ has at least two elements $P_1$ and $P_2$. Then $M$ contains
submodules $K_i$ isomorphic to $A/P_i$ for $i = 1, 2$. Now since $\ann(x) = P_i$ for
any non-zero $x \in K_i$, we must have $K_1 \cap K_2 = 0$, and hence $0$ is reducible.

(ii) We can again assume that $N = 0$. If $0 = N_1 \cap \cdots \cap N_l$ then $M$ is
isomorphic to a submodule of $M/N_1 \oplus \cdots \oplus M/N_l$, so that
\[
  \Ass(M) \subset \bigcup_i \Ass(M/N_i) = \{P_1, \ldots, P_l\}.
\]
On the other hand $N_2 \cap \cdots \cap N_l \neq 0$, and taking $0 \neq x \in N_2 \cap \cdots \cap N_l$, we
have $\ann(x) = 0:x \subset N_1:x$. But $N_1:M$ is a primary ideal belonging to $P_1$, so
that $P_1^\nu M \subset N_1$ for some $\nu > 0$. Therefore $P_1^\nu x = 0$; hence there exists $i \ge 0$
such that $P_1^i x \neq 0$ but $P_1^{i+1} x = 0$, and choosing $0 \neq y \in P_1^i x$ we have
$P_1y = 0$. However, since $y \in N_2 \cap \cdots \cap N_l$ it follows that $y \notin N_1$, and by
the definition of primary submodule $\ann(y) \subset P_1$, so that $P_1 = \ann(y)$ and
$P_1 \in \Ass(M)$. The same works for the other $P_i$, and this proves that
$\{P_1, \ldots, P_l\} \subset \Ass(M)$.

(iii) We have already seen that a proper submodule has an irreducible
decomposition, so that by (i) it has a primary decomposition. Suppose that
$N = N_1 \cap \cdots \cap N_l$ is a shortest primary decomposition, and that $N_1$ is the
$P$-primary component with $P = P_1$. By Ex.~\ref{ex:4.8-ex} we know that $N_P =
(N_1)_P \cap \cdots \cap (N_l)_P$, and for $i > 1$ a power of $P_i$ is contained in $\ann(M/N_i)$;
then since $P_i \neq P_1$ we have $(M/N_i)_P = 0$, and therefore $(N_i)_P = M_P$. Thus
$N_P = (N_1)_P$ and hence $\varphi_P^{-1}(N_P) = \varphi_P^{-1}((N_1)_P)$; it is easy to check that
the right-hand side is $N_1$.
\end{proof}

\begin{remark}
The uniqueness of the $P$-primary component $N_1$ proved in (iii) for
minimal primes $P$ does not hold in general; see Ex.~\ref{ex:6.6-ex}.
\end{remark}

% ============ 习题 ============
\section*{Exercises to \S\ref{sec:6}}

\begin{exercise}\label{ex:6.1-ex}
Find $\Ass(M)$ for the $\Z$-module $M = \Z \oplus (\Z/3\Z)$.
\end{exercise}

\begin{exercise}\label{ex:6.2-ex}
If $M$ is a finite module over a Noetherian ring $A$, and $M_1, M_2$ are
submodules of $M$ with $M = M_1 + M_2$ then can we say that $\Ass(M) =
\Ass(M_1) \cup \Ass(M_2)$?
\end{exercise}

\begin{exercise}\label{ex:6.3-ex}
Let $A$ be a Noetherian ring and let $x \in A$ be an element which is neither a
unit nor a zero-divisor; prove that the ideals $xA$ and $x^nA$ for $n = 1, 2, \ldots$
have the same prime divisors:
\[
  \Ass_A(A/xA) = \Ass_A(A/x^nA).
\]
\end{exercise}

\begin{exercise}\label{ex:6.4-ex}
Let $I$ and $J$ be ideals of a Noetherian ring $A$. Prove that if $JA_P \subset IA_P$ for
every $P \in \Ass_A(A/I)$ then $J \subset I$.
\end{exercise}

\begin{exercise}\label{ex:6.5-ex}
Prove that the total ring of fractions of a reduced Noetherian ring $A$ is a
direct product of fields.
\end{exercise}

\begin{exercise}\label{ex:6.6-ex}
(Taken from [Nor 1], p. 30.) Let $k$ be a field. Show that in $k[X, Y]$ we have
\[
  (X^2, XY) = (X) \cap (X^2, Y) = (X) \cap (X^2, XY, Y^2).
\]
\end{exercise}

\begin{exercise}\label{ex:6.7-ex}
Let $f\colon A \to B$ be a homomorphism of Noetherian rings, and $M$ a finite $B$-
module. Write ${}^a f\colon \Spec B \to \Spec A$ as in \S\ref{sec:4}. Prove that ${}^a f(\Ass_B(M)) =
\Ass_A(M)$. (Consequently, $\Ass_A(M)$ is a finite set for such $M$.)
\end{exercise}

% ============ Appendix to §6. Secondary representations of a module (原书 pp.42-44) ============
\section*{Appendix to \S\ref{sec:6}. Secondary representations of a module}
\phantomsection
\addcontentsline{toc}{section}{Appendix to \S\ref{sec:6}. Secondary representations of a module}

I.G. Macdonald~\cite{MacD1} has developed the theory of \emph{attached prime ideals}\index{attached prime ideal}
and \emph{secondary representations}\index{secondary representation} of a module, which is in a certain sense
dual to the theory of associated prime ideals and primary decompositions.
This theory was successfully applied to the theory of local cohomology
by him and R.Y. Sharp (Macdonald \& Sharp~\cite{MS1}, Sharp~\cite{Shar7}).

Let $A$ be a commutative ring. An $A$-module $M$ is said to be \emph{secondary}\index{secondary module}
if $M \neq 0$ and, for each $a \in A$, the endomorphism $\varphi_a\colon M \to M$ defined
by $\varphi_a(m) = am$ (for $m \in M$) is either surjective or nilpotent. If $M$ is secondary,
then $P = \sqrt{\ann M}$ is a prime ideal, and $M$ is said to be \emph{$P$-secondary}. Any
non-zero quotient of a $P$-secondary module is $P$-secondary.

\begin{example}\label{ex:6.5}
If $A$ is an integral domain, its quotient field $K$ is a $(0)$-secondary
$A$-module.
\end{example}

\begin{example}\label{ex:6.6}
Let $W = \Z[p^{-1}]$, where $p$ is a prime number, and consider the
Artinian $\Z$-module $W/\Z$ (see \S\ref{sec:3}). This is also a $(0)$-secondary $\Z$-module.
\end{example}

\begin{example}\label{ex:6.7}
If $A$ is a local ring with maximal ideal $P$ and if every element of
$p$ is nilpotent, then $A$ itself is a $P$-secondary $A$-module.
\end{example}

\begin{example}\label{ex:6.8}
If $P$ is a maximal ideal of $A$, then $A/P^n$ is a $P$-secondary $A$-
module for every $n > 0$.
\end{example}

A \emph{secondary representation} of an $A$-module $M$ is an expression of $M$ as a
finite sum of secondary submodules:
\[
  (*) \quad M = N_1 + \cdots + N_r.
\]
The representation is \emph{minimal}\index{minimal secondary representation} if (1) the prime ideals $P_i := \sqrt{\ann N_i}$ are all
distinct, and (2) none of the $N_i$ is redundant. It is easy to see that the sum of
two $P$-secondary submodules is again $P$-secondary, hence if $M$ has a
secondary representation then it has a minimal one.

A prime ideal $P$ is called an \emph{attached prime ideal} of $M$ if $M$ has a $P$-
secondary quotient. The set of the attached prime ideals of $M$ is denoted by
$\Att(M)$.

\begin{theorem}\label{thm:6.9}
If $(*)$ is a minimal secondary representation of $M$ and $P_i =
\sqrt{\ann N_i}$, then $\Att(M) = \{P_1, \ldots, P_n\}$.
\end{theorem}

\begin{proof}
Since $M/(N_1 + \cdots + N_{i-1} + N_{i+1} + \cdots + N_n)$ is a non-zero quotient
of $N_i$, it is a $P_i$-secondary module. Thus $\{P_1, \ldots, P_n\} \subset \Att(M)$. Conver-
sely, let $P \in \Att(M)$ and let $W$ be a $P$-secondary quotient of $M$. Then $W =
W_1 + \cdots + W_n$, where $W_i$ is the image of $N_i$ in $W$. From this we obtain a
minimal secondary representation $W = W'_1 + \cdots + W'_m$, and then
$\Att(W) \supset \{P'_1, \ldots, P'_m\}$. On the other hand $\Att(W) = \{P\}$ since $W$ is
$P$-secondary. Therefore $P = P_i$ for some $i$.
\end{proof}

\begin{theorem}\label{thm:6.10}
If $0 \to M' \to M \to M'' \to 0$ is an exact sequence of
$A$-modules, then $\Att(M'') \subset \Att(M) \subset \Att(M') \cup \Att(M'')$.
\end{theorem}

\begin{proof}
The first inclusion is trivial from the definition. For the second, let
$P \in \Att(M)$ and let $N$ be a submodule such that $M/N$ is $P$-secondary. If
$M' + N = M$ then $M/N$ is a non-trivial quotient of $M'$, hence $P \in \Att(M')$.
If $M' + N \neq M$ then $M/(M' + N)$ is a non-trivial quotient of $M''$ as well
as of $M/N$, hence $M''$ has a $P$-secondary quotient and $P \in \Att(M'')$.
\end{proof}

An $A$-module $M$ is said to be \emph{sum-irreducible}\index{sum-irreducible module} if it is neither zero nor the
sum of two proper submodules.

\begin{lemma*}\label{lem:6.1}
If $M$ is Artinian and sum-irreducible, then it is secondary.
\end{lemma*}

\begin{proof}
Suppose $M$ is not secondary. Then there is $a \in A$ such that $M \neq aM$
and $a^nM \neq 0$ for all $n > 0$. Since $M$ is Artinian, we have $a^nM = a^{n+1}M$ for
some $n$. Set $K = \{x \in M \mid a^nx = 0\}$. Then it is immediate that $M = K + aM$,
and so $M$ is not sum-irreducible.
\end{proof}

\begin{theorem}\label{thm:6.11}
If $M$ is Artinian, then it has a secondary representation.
\end{theorem}

\begin{proof}
Similar to the proof of Theorem~\ref{thm:6.8}, (iii).
\end{proof}

The class of modules which have secondary representations is larger
than that of Artinian modules. Sharp [8] proved that an injective module
over a Noetherian ring has a secondary representation.

% ============ 习题 ============
\section*{Exercises to Appendix to \S\ref{sec:6}}

\begin{exercise}\label{ex:6.8-ex}
An $A$-module $M$ is \emph{coprimary}\index{coprimary module} if $\Ass(M)$ has just one element. Show that a
finite module $M \neq 0$ over a Noetherian ring $A$ is coprimary if and only if
the following condition is satisfied: for every $a \in A$, the endomorphism
$a\colon M \to M$ is either injective or nilpotent. In this case $\Ass M = \{P\}$,
where $P = \sqrt{\ann M}$.
\end{exercise}

\begin{exercise}\label{ex:6.9-ex}
Show that if $M$ is an $A$-module of finite length then $M$ is coprimary if and
only if it is secondary. Show also that such a module $M$ is a direct sum of
secondary modules belonging to maximal ideals, and $\Ass(M) = \Att(M)$.
\end{exercise}


% ============ Chapter 3. Properties of extension rings (原书 pp.45-71) ============
\chapter{Properties of extension rings}\label{ch:3}

Flatness was formulated by Serre in the 1950s and quickly grew into one of
the basic tools of both algebraic geometry and commutative algebra. This is
an algebraic notion which is hard to grasp geometrically. Flatness is defined
quite generally for modules, but is particularly important for extensions of
rings. The model case is that of completion. Complete local rings have a
number of wonderful properties, and passing to the completion of a local
ring is an effective technique in many cases; this is analogous to studying an
algebraic variety as an analytic space. The theory of integral extensions of
rings had been studied by Krull, and he discovered the so-called going-up
and going-down theorems. We show that the going-down theorem also
holds for flat extensions, and gather together flatness, completion and
integral extensions in this chapter. We will use more sophisticated argu-
ments to study flatness over Noetherian rings in Chapter 8, and completion
in Chapter 10. % TODO: 第 8、10 章重排完成后改为 \ref{ch:8} / \ref{ch:10}

% ============ §7 Flatness (原书 pp.45-53) ============
\section{Flatness}\label{sec:7}\index{flatness, section 7}

Let $A$ be a ring and $M$ an $A$-module. Writing $\mathcal{Y}$ to stand for a
sequence $\cdots \to N' \to N \to N'' \to \cdots$ of $A$-modules and linear
maps, we let $\mathcal{Y} \otimes_A M$, or simply $\mathcal{Y} \otimes M$ stand for the induced sequence $\cdots
\to N' \otimes_A M \to N \otimes_A M \to N'' \otimes_A M \to \cdots$.

\begin{definition}
$M$ is \emph{flat}\index{flat module} over $A$ if for every exact sequence $\mathcal{Y}$ the sequence
$\mathcal{Y} \otimes_A M$ is again exact. We sometimes shorten this to $A$-flat.

$M$ is \emph{faithfully flat}\index{faithfully flat module} if for every sequence $\mathcal{Y}$,
\[
  \mathcal{Y} \text{ is exact} \iff \mathcal{Y} \otimes_A M \text{ is exact}.
\]
\end{definition}

Any exact sequence $\mathcal{Y}$ can be broken up into short exact sequences of
the form $0 \to N_1 \to N_2 \to N_3 \to 0$, so that in the definition of flatness
we need only consider short exact sequences $\mathcal{Y}$. Moreover, in view of the
right-exactness of tensor product (see Appendix A, Formula 8) we can
restrict attention to exact sequences $\mathcal{Y}$ of the form $0 \to N_1 \to N_2$ and
need only check the exactness of $\mathcal{Y} \otimes M$: $0 \to N_1 \otimes M \to N_2 \otimes M$.

If $f\colon A \to B$ is a homomorphism of rings and $B$ is flat as an $A$-module,
we say that $f$ is a \emph{flat homomorphism}, or that $B$ is a \emph{flat $A$-algebra}. For
example, the localisation $A_S$ of $A$ is a flat $A$-algebra (Theorems 4.4 and 4.5).

\emph{Transitivity.} Let $B$ be an $A$-algebra and $M$ a $B$-module. Then the following
hold;
\begin{enumerate}[(1)]
  \item $B$ is flat over $A$ and $M$ is flat over $B \Rightarrow M$ is flat over $A$;
  \item $B$ is faithfully flat over $A$ and $M$ is faithfully flat over $B \Rightarrow M$ is
        faithfully flat over $A$;
  \item $M$ is faithfully flat over $B$ and flat over $A \Rightarrow B$ is flat over $A$;
  \item $M$ is faithfully flat over both $A$ and $B \iff B$ is faithfully flat over $A$.
\end{enumerate}
Each of these follows easily from the fact that $(\mathcal{Y} \otimes_A B) \otimes_B M = \mathcal{Y} \otimes_A M$ for
any sequence of $A$-modules $\mathcal{Y}$.

\emph{Change of coefficient ring.} Let $B$ be an $A$-algebra and $M$ an $A$-module.
Then the following hold:
\begin{enumerate}[(1)]
  \item $M$ is flat over $A \Rightarrow M \otimes_A B$ is flat over $B$;
  \item $M$ is faithfully flat over $A \Rightarrow M \otimes_A B$ is faithfully flat over $B$.
\end{enumerate}
These follow from that fact that $\mathcal{Y} \otimes_A (B \otimes_A M) = \mathcal{Y} \otimes_A M$ for any
sequence of $B$-modules $\mathcal{Y}$.

\begin{theorem}\label{thm:7.1}
Let $A \to B$ be a homomorphism of rings and $M$ a $B$-
module. A necessary and sufficient condition for $M$ to be flat over $A$ is that
for every prime ideal $P$ of $B$, the localisation $M_P$ is flat over $A_p$, where
$p = P \cap A$ (or the same condition for every maximal ideal $P$ of $B$).
\end{theorem}

\begin{proof}
First of all we make the following observation: if $S \subset A$ is a
multiplicative set and $M, N$ are $A_S$-modules, then $M \otimes_{A_S} N = M \otimes_A N$.
This follows from the fact that in $N \otimes_A M$ we have
\[
  \frac{a}{s}x \otimes y = x \otimes \frac{a}{s}y \quad \text{for } x \in M, y \in N, a \in A \text{ and } s \in S.
\]
(In general, if $B$ is an $A$-algebra and $M$ and $N$ are $B$-modules, it can be seen from the construction of the
tensor product that $M \otimes_B N$ is the quotient of $M \otimes_A N$ by the submodule generated
by $\{bx \otimes y - x \otimes by \mid x \in M, y \in N \text{ and } b \in B\}$.)

Assume now that $M$ is $A$-flat. The map $A \to B$ induces $A_p \to B_P$,
and $M_P$ is a $B_P$-module, therefore an $A_p$-module. Let $\mathcal{Y}$ be an exact
sequence of $A_p$-modules; then, by the above observation,
\[
  \mathcal{Y} \otimes_{A_p} M_P = \mathcal{Y} \otimes_A M_P = (\mathcal{Y} \otimes_A M) \otimes_B B_P,
\]
and the right-hand side is an exact sequence, so that $M_P$ is $A_p$-flat.

Next, suppose that $M_P$ is $A_p$-flat for every maximal ideal $P$ of $B$.
Let $0 \to N' \to N$ be an exact sequence of $A$-modules, and write $K$
for the kernel of the $B$-linear map $N' \otimes_A M \to N \otimes_A M$, so that
$0 \to K \to N' \otimes_A M \to N \otimes_A M$ is an exact sequence of $B$-modules. For
any $P \in \text{m-Spec } B$ the localisation
\[
  0 \to K_P \to N' \otimes_A M_P \to N \otimes_A M_P
\]
is exact, and since $N' \otimes_A M_P = N' \otimes_A (A_p \otimes_{A_p} M_P) = N'_p \otimes_{A_p} M_P$, and
similarly $N \otimes_A M_P = N_p \otimes_{A_p} M_P$, we have $K_P = 0$ by hypothesis. There-
fore by Theorem~\ref{thm:4.6} we have $K = 0$, and this is what we have to prove.
\end{proof}

\begin{theorem}\label{thm:7.2}
Let $A$ be a ring and $M$ an $A$-module. Then the following
conditions are equivalent:
\begin{enumerate}[(1)]
  \item $M$ is faithfully flat over $A$;
  \item $M$ is $A$-flat, and $N \otimes_A M \neq 0$ for any non-zero $A$-module $N$;
  \item $M$ is $A$-flat, and $mM \neq M$ for every maximal ideal $m$ of $A$.
\end{enumerate}
\end{theorem}

\begin{proof}
$(1) \Rightarrow (2)$. Let $\mathcal{Y}$ be the sequence $0 \to N \to 0$. If $N \otimes M = 0$ then
$\mathcal{Y} \otimes M$ is exact, so $\mathcal{Y}$ is exact, and therefore $N = 0$.

$(2) \Rightarrow (3)$. This is clear from $M/mM = (A/m) \otimes_A M$.

$(3) \Rightarrow (2)$. If $N \neq 0$ and $0 \neq x \in N$ then $Ax \cong A/\ann(x)$, so that taking a
maximal ideal $m$ containing $\ann(x)$, we have $M \neq mM \supset \ann(x) \cdot M$;
hence, $Ax \otimes M \neq 0$. By the flatness assumption, $Ax \otimes M \to N \otimes M$ is
injective, so that $N \otimes M \neq 0$.

$(2) \Rightarrow (1)$. Consider a sequence of $A$-modules
\[
  \mathcal{Y}\colon \quad N' \xrightarrow{f} N \xrightarrow{g} N''.
\]
If
\[
  \mathcal{Y} \otimes M\colon \quad N' \otimes M \xrightarrow{f_M} N \otimes M \xrightarrow{g_M} N'' \otimes M
\]
is exact then $g_M f_M = (gf)_M = 0$, so that by flatness, $\Im(gf) \otimes M =
\Im(g_M f_M) = 0$. By assumption we then have $\Im(gf) = 0$, that is $gf = 0$;
hence $\Ker g \supset \Im f$. If we set $H = \Ker g/\Im f$ then by flatness,
\[
  H \otimes M = \Ker(g_M)/\Im(f_M) = 0,
\]
so that the assumption gives $H = 0$. Therefore $\mathcal{Y}$ is exact.
\end{proof}

A ring homomorphism $f\colon A \to B$ induces a map ${}^a f\colon \Spec B \to \Spec A$,
under which a point $p \in \Spec A$ has an inverse image ${}^a f^{-1}(p) =
\{P \in \Spec B \mid P \cap A = p\}$ which is homeomorphic to $\Spec(B \otimes_A \kappa(p))$.
Indeed, setting $C = B \otimes_A \kappa(p)$ and $S = A - p$, and defining $g\colon B \to C$ by
$g(b) = b \otimes 1$, then since $\kappa(p) = (A/p) \otimes_A A_p$, we have
\[
  C = B \otimes_A \kappa(p) \cong B \otimes_A (A/p) \otimes_A A_p = (B/pB) \otimes_A A_p = (B/pB)_{\bar{S}},
\]
where $\bar{S}$ is the image of $S$ in $A/p$. Thus ${}^a g\colon \Spec C \to \Spec B$ has the image
\[
  \{P \in \Spec B \mid P \supset pB \text{ and } P \cap f(S) = \emptyset\}
   = \{P \in \Spec B \mid P \cap A = p\},
\]
which is ${}^a f^{-1}(p)$, and ${}^a g$ induces a homeomorphism of $\Spec C$ with
${}^a f^{-1}(p)$. For this reason we call $\Spec C = \Spec(B \otimes \kappa(p))$ the \emph{fibre} over $p$.
The inverse map ${}^a f^{-1}(p) \to \Spec C$ takes $P \in {}^a f^{-1}(p)$ into $PC = PB_P/pB_P$.

For $P^* \in \Spec C$ we set $P = P^* \cap B$; then by Theorems 4.2 and 4.3, we have
\[
  P^* = PC \quad \text{and} \quad C_{P^*} = (B_P/pB_P)_{P^*} = B_P/pB_P = B_P \otimes_{A_P} \kappa(p).
\]

\begin{theorem}\label{thm:7.3}
Let $f\colon A \to B$ be a ring homomorphism and $M$ a $B$-
module. Then
\begin{enumerate}[(i)]
  \item $M$ is faithfully flat over $A \iff {}^a f(\Supp(M)) = \Spec A$.
  \item If $M$ is a finite $B$-module then
        \[
          M \text{ is $A$-flat and } {}^a f(\Supp(M)) \supset \text{m-Spec } A \iff M \text{ is faithfully flat over } A.
        \]
\end{enumerate}
\end{theorem}

\begin{proof}
(i) For $p \in \Spec A$, by faithful flatness we have $M \otimes_A \kappa(p) \neq 0$. Hence,
if we set $C = B \otimes_A \kappa(p)$ and $M' = M \otimes_A \kappa(p) = M \otimes_B C$, the $C$-module
$M' \neq 0$, so that there is a $P^* \in \Spec C$ such that $M'_{P^*} \neq 0$. Now set
$P = P^* \cap B$; then
\[
  M_P = M \otimes_B C_{P^*} = M \otimes_B (B_P \otimes_{B_P} C_{P^*}) = M_P \otimes_{B_P} C_{P^*},
\]
so that $M_P \neq 0$, that is $P \in \Supp(M)$. But $P^* \in \Spec(B \otimes \kappa(p))$, so that as
we have seen $P \cap A = p$. Therefore $p \in {}^a f(\Supp(M))$.

(ii) It is enough to show that $M/mM \neq 0$ for any maximal ideal $m$ of $A$. By
assumption there is a prime ideal $P$ of $B$ such that $P \cap A = m$ and
$M_P \neq 0$. By NAK, since $M_P$ is finite over $B_P$, we have $M_P/PM_P \neq 0$, and
a fortiori $M_P/mM_P = (M/mM)_P \neq 0$, so that $M/mM \neq 0$.
\end{proof}

Let $(A, m)$ and $(B, n)$ be local rings, and $f\colon A \to B$ a ring homomor-
phism; $f$ is said to be a \emph{local homomorphism}\index{local homomorphism} if $f(m) \subset n$. If this happens then
by Theorem 2, or by Theorem 3, (ii), we see that it is equivalent to say that $f$
is flat or faithfully flat.

Let $S$ be a multiplicative set of $A$. Then it is easy to see that $\Spec(A_S)
\to \Spec A$ is surjective only if $S$ consists of units, that is $A = A_S$. Thus
from the above theorem, if $A \neq A_S$ then $A_S$ is flat but not faithfully flat over
$A$.

\begin{theorem}\label{thm:7.4}
\begin{enumerate}[(i)]
  \item Let $A$ be a ring, $M$ a flat $A$-module, and $N_1, N_2$ two submodules of an
        $A$-module $N$. Then as submodules of $N \otimes M$ we have
        \[
          (N_1 \cap N_2) \otimes M = (N_1 \otimes M) \cap (N_2 \otimes M).
        \]
  \item Let $A \to B$ be a flat ring homomorphism, and let $I_1$ and $I_2$ be
        ideals of $A$. Then
        \[
          (I_1 \cap I_2)B = I_1B \cap I_2B.
        \]
  \item If in addition $I_2$ is finitely generated then
        \[
          (I_1:I_2)B = I_1B:I_2B.
        \]
\end{enumerate}
\end{theorem}

\begin{proof}
(i) Define $\varphi\colon N \to N/N_1 \oplus N/N_2$ by $\varphi(x) = (x + N_1, x + N_2)$; then
$0 \to N_1 \cap N_2 \to N \to N/N_1 \oplus N/N_2$ is exact, and hence so is
\[
  0 \to (N_1 \cap N_2) \otimes M \to N \otimes M \to
      (N \otimes M)/(N_1 \otimes M) \oplus (N \otimes M)/(N_2 \otimes M).
\]
This is the assertion in (i).

(ii) This is a particular case of (i) with $N = A$, $M = B$, in view of the fact
that for an ideal $I$ of $A$ the subset $I \otimes_A B$ of $A \otimes_A B = B$ coincides with $IB$.

(iii) If $I_2 = Aa_1 + \cdots + Aa_n$ then since $(I_1:I_2) = \bigcap_i (I_1:a_i)$, we can use
(ii) to reduce to the case that $I_2$ is principal. For $a \in A$ we have the exact
sequence
\[
  0 \to (I_1:a) \to A \xrightarrow{a} A/I_1,
\]
and tensoring this with $B$ gives the assertion.
\end{proof}

\begin{example*}
Let $k$ be a field, and consider the subring $A = k[x^2, x^3]$ of the
polynomial ring $B = k[x]$ in an indeterminate $x$. Then $x^2A \cap x^3A$ is the set
of polynomials made up of terms of degree $\ge 5$ in $x$, so that $(x^2A \cap x^3A)B
= x^5B$, but on the other hand $x^2B \cap x^3B = x^3B$. Therefore by the above
theorem, $B$ is not flat over $A$.
\end{example*}

\begin{theorem}\label{thm:7.5}
Let $f\colon A \to B$ be a faithfully flat ring homomorphism.
\begin{enumerate}[(i)]
  \item For any $A$-module $M$, the map $M \to M \otimes_A B$ defined by $m \mapsto m \otimes 1$
        is injective; in particular $f\colon A \to B$ is itself injective.
  \item If $I$ is an ideal of $A$ then $IB \cap A = I$.
\end{enumerate}
\end{theorem}

\begin{proof}
(i) Let $0 \neq m \in M$. Then $(Am) \otimes B$ is a $B$-submodule of $M \otimes B$ which
can be identified with $(m \otimes 1)B$. But by Theorem 2, $(Am) \otimes B \neq 0$, so that
$m \otimes 1 \neq 0$.

(ii) follows by applying (i) to $M = A/I$, using $(A/I) \otimes B = B/IB$.
\end{proof}

\begin{theorem}\label{thm:7.6}
Let $A$ be a ring and $M$ a flat $A$-module. If $a_{ij} \in A$ and $x_j \in M$
(for $1 \le i \le l$ and $1 \le j \le n$) satisfy
\[
  \sum_j a_{ij}x_j = 0 \quad \text{for all } i,
\]
then there exists an integer $s$ and $b_{jk} \in A$, $y_k \in M$ (for $1 \le j \le n$ and $1 \le k \le s$)
such that
\[
  \sum_j a_{ij}b_{jk} = 0 \text{ for all } i, k, \quad \text{and} \quad x_j = \sum_k b_{jk}y_k \text{ for all } j.
\]
Thus the solutions in a flat module $M$ of a system of simultaneous linear
equations with coefficients in $A$ can be expressed as a linear combination of
solutions in $A$. Conversely, if the above conclusion holds for the case of a
single equation (that is for $l = 1$), then $M$ is flat.
\end{theorem}

\begin{proof}
Set $\varphi\colon A^n \to A^l$ for the linear map defined by the matrix $(a_{ij})$, and
let $\varphi_M\colon M^n \to M^l$ be the same thing for $M$; then $\varphi_M = \varphi \otimes 1$, where $1$ is
the identity map of $M$. Setting $K = \Ker \varphi$ and tensoring the exact sequence
\[
  0 \to K \to A^n \xrightarrow{\varphi} A^l
\]
with $M$, we get the exact sequence
\[
  0 \to K \otimes M \to M^n \xrightarrow{\varphi_M} M^l.
\]
By assumption $\varphi_M(x_1, \ldots, x_n) = 0$, so that we can write
\[
  (x_1, \ldots, x_n) = (i \otimes 1)\left(\sum_k \beta_k \otimes y_k\right) \quad \text{with } \beta_k \in K \text{ and } y_k \in M.
\]
If we write out $\beta_k$ as an element of $A^n$ in the form $\beta_k = (b_{1k}, \ldots, b_{nk})$ with
$b_{jk} \in A$ then the conclusion follows. The converse will be proved after the
next theorem.
\end{proof}

\begin{theorem}\label{thm:7.7}
Let $A$ be a ring and $M$ an $A$-module. Then $M$ is flat over $A$ if
and only if for every finitely generated ideal $I$ of $A$ the canonical map
$I \otimes_A M \to A \otimes_A M$ is injective, and therefore $I \otimes M \cong IM$.
\end{theorem}

\begin{proof}
The `only if' is obvious, and we prove the `if'. Firstly, every ideal
of $A$ is the direct limit of the finitely generated ideals contained in it, so
that by Theorems A1 and A2 of Appendix A, $I \otimes M \to M$ is injective for
every ideal $I$. Moreover, if $N$ is an $A$-module and $N' \subset N$ a submodule,
then since $N$ is the direct limit of modules of the form $N' + F$, with $F$
finitely generated, to prove that $N' \otimes M \to N \otimes M$ is injective we can
assume that $N = N' + A\omega_1 + \cdots + A\omega_n$. Then setting $N_i = N' + A\omega_1
+ \cdots + A\omega_i$ (for $1 \le i \le n$), we need only show that each step in the
chain
\[
  N' \otimes M \to N_1 \otimes M \to N_2 \otimes M \to \cdots \to N \otimes M
\]
is injective, and finally that if $N = N' + A\omega$ then $N' \otimes M \to N \otimes M$
is injective. Now we set $I = \{a \in A \mid a\omega \in N'\}$, and get the exact sequence
\[
  0 \to N' \to N \to A/I \to 0.
\]
This induces a long exact sequence (see Appendix B, p. 279)
\[
  \cdots \to \Tor^A(M, A/I) \to N' \otimes M \to N \otimes M \to (A/I) \otimes M \to 0;
\]
hence it is enough to prove that
\[
  (*) \quad \Tor^A_1(M, A/I) = 0.
\]
For this consider the short exact sequence
\[
  0 \to I \to A \to A/I \to 0
\]
and the induced long exact sequence
\[
  \Tor^A_1(M, A) = 0 \to \Tor^A_1(M, A/I) \to I \otimes M \to M \to \cdots;
\]
since $I \otimes M \to M$ is injective, $(*)$ must hold.
\end{proof}

From this theorem we can prove the converse of Theorem 6. Indeed, if
$I = Au_1 + \cdots + Au_n$ is a finitely generated ideal of $A$ then an element $\xi$ of
$I \otimes M$ can be written as $\xi = \sum_i a_i \otimes m_i$ with $m_i \in M$. Suppose that $\xi$ is $0$ in $M$,
that is that $\sum a_im_i = 0$. Now if the conclusion of Theorem 6 holds for $M$,
there exist $b_{ij} \in A$ and $y_j \in M$ such that
\[
  \sum_i u_ib_{ij} = 0 \text{ for all } j, \quad \text{and} \quad m_i = \sum_j b_{ij}y_j \text{ for all } i.
\]
Then $\xi = \sum a_i \otimes m_i = \sum_{i,j} u_ib_{ij} \otimes y_j = 0$, so that $I \otimes M \to M$ is injec-
tive, and therefore $M$ is flat.

\begin{theorem}\label{thm:7.8}
Let $A$ be a ring and $M$ an $A$-module. The following
conditions are equivalent:
\begin{enumerate}[(1)]
  \item $M$ is flat;
  \item for every $A$-module $N$ we have $\Tor^A_i(M, N) = 0$ for $i > 0$;
  \item $\Tor^A_1(M, A/I) = 0$ for every finitely generated ideal $I$.
\end{enumerate}
\end{theorem}

\begin{proof}
$(1) \Rightarrow (2)$ If we let $\cdots \to L_i \to L_{i-1} \to \cdots \to L_0 \to N \to 0$
be a projective resolution of $N$ then
\[
  \cdots \to L_i \otimes M \to L_{i-1} \otimes M \to \cdots \to L_0 \otimes M
\]
is exact, so that $\Tor^A_i(M, N) = 0$ for all $i > 0$.

$(2) \Rightarrow (3)$ is obvious.

$(3) \Rightarrow (1)$ The short exact sequence $0 \to I \to A \to A/I \to 0$ induces a
long exact sequence
\[
  \Tor^A_1(M, A/I) = 0 \to I \otimes M \to M \to M \otimes A/I \to 0,
\]
and hence $I \otimes M \to M$ is injective; therefore by the previous theorem $M$
is flat.
\end{proof}

\begin{theorem}\label{thm:7.9}
Let $0 \to M' \to M \to M'' \to 0$ be an exact sequence of $A$-
modules; then if $M'$ and $M''$ are both flat, so is $M$.
\end{theorem}

\begin{proof}
For any $A$-module $N$ the sequence $\Tor_1(M', N) \to \Tor_1(M, N)
\to \Tor_1(M'', N)$ is exact, and since the first and third groups are zero,
also $\Tor_1(M, N) = 0$. Therefore by the previous theorem $M$ is flat.
\end{proof}

A free module is obviously faithfully flat (if $F$ is free and $\mathcal{Y}$ is a sequence of
$A$-modules then $\mathcal{Y} \otimes F$ is just a sum of copies of $\mathcal{Y}$ in number equal to the
cardinality of a basis of $F$). Conversely, over a local ring the following
theorem holds, so that for finite modules flat, faithfully flat and free are
equivalent conditions.

\begin{theorem}\label{thm:7.10}
Let $(A, m)$ be a local ring and $M$ a flat $A$-module. If
$x_1, \ldots, x_n \in M$ are such that their images $\bar{x}_1, \ldots, \bar{x}_n$ in $\bar{M} = M/mM$ are
linearly independent over the field $A/m$ then $x_1, \ldots, x_n$ are linearly
independent over $A$. Hence if $M$ is finite, or if $m$ is nilpotent, then any
minimal basis of $M$ (see \S\ref{sec:2}) is a basis of $M$, and $M$ is a free module.
\end{theorem}

\begin{proof}
By induction on $n$. If $n = 1$, and $a \in A$ is such that $ax_1 = 0$ then
by Theorem 6 there are $b_1, \ldots, b_s \in A$ such that $ab_i = 0$ and $x_1 \in \sum b_iM$; by
assumption $x_1 \notin mM$, so that among the $b_i$ there must be one not contained
in $m$. This $b_i$ is then a unit, so that we must have $a = 0$.

For $n > 1$, let $\sum a_ix_i = 0$; then there are $b_{ij} \in A$ and $y_j \in M$ (for $1 \le j \le s$)
such that $\sum a_ib_{ij} = 0$ and $x_i = \sum_j b_{ij}y_j$. Now $x_n \notin mM$, so that among the $b_{nj}$ at
least one is a unit. Hence $a_n$ is a linear combination of $a_1, \ldots, a_{n-1}$, that
is $a_n = \sum_{i=1}^{n-1} a_ic_i$ for some $c_i \in A$. Therefore we have
\[
  a_1(x_1 + c_1x_n) + \cdots + a_{n-1}(x_{n-1} + c_{n-1}x_n) = 0;
\]
however, the $(n - 1)$ elements $\bar{x}_1 + \bar{c}_1\bar{x}_n, \ldots, \bar{x}_{n-1} + \bar{c}_{n-1}\bar{x}_n$ of $\bar{M}$ are
linearly independent over $A/m$, so that by induction, $a_1 = \cdots = a_{n-1} = 0$. Hence also
$a_n = 0$.
\end{proof}

\begin{theorem}\label{thm:7.11}
Let $A$ be a ring, $M$ and $N$ two $A$-modules, and $B$ a flat $A$-
algebra. If $M$ is of finite presentation then we have
\[
  \Hom_A(M, N) \otimes_A B = \Hom_B(M \otimes_A B, N \otimes_A B).
\]
\end{theorem}

\begin{proof}
Fixing $N$ and $B$, we define contravariant functors $F$ and $G$ of an
$A$-module $M$ by
\[
  F(M) = \Hom_A(M, N) \otimes_A B
\]
and
\[
  G(M) = \Hom_B(M \otimes_A B, N \otimes_A B);
\]
then we can define a morphism of functors $\lambda\colon F \to G$ by
\[
  \lambda(f \otimes b) = b \cdot (f \otimes 1_B) \quad \text{for } f \in \Hom_A(M, N) \text{ and } b \in B.
\]
Both $F$ and $G$ are left-exact functors.

Now if $M$ is of finite presentation there is an exact sequence of the form
$A^p \to A^q \to M \to 0$, and from this we get a commutative diagram
\[
  \begin{tikzcd}[column sep=large, row sep=large]
    0 \arrow[r] & F(M) \arrow[r] \arrow[d, "\lambda"] & F(A^q) \arrow[r] \arrow[d, "\lambda"] & F(A^p) \arrow[d, "\lambda"] \\
    0 \arrow[r] & G(M) \arrow[r] & G(A^q) \arrow[r] & G(A^p)
  \end{tikzcd}
\]
having two exact rows. Now $F(A^p) = N^p \otimes B$ and $G(A^p) = (N \otimes B)^p$, so that
the right-hand $\lambda$ is an isomorphism, and similarly the middle $\lambda$ is an
isomorphism. Thus, as one sees easily, the left-hand $\lambda$ is also an
isomorphism.
\end{proof}

\begin{corollary*}
Let $A$, $M$ and $N$ be as in the theorem, and let $p$ be a prime ideal of
$A$. Then
\[
  \Hom_A(M, N) \otimes_A A_p = \Hom_{A_p}(M_p, N_p).
\]
\end{corollary*}

\begin{theorem}\label{thm:7.12}
Let $A$ be a ring and $M$ an $A$-module of finite presentation.
Then $M$ is a projective $A$-module if and only if $M_m$ is a free $A_m$-module for
every maximal ideal $m$ of $A$.
\end{theorem}

\begin{proof}[Proof of `only if']
If $M$ is projective it is a direct summand of a free module
and this property is preserved by localisation, so that $M_m$ is projective over
$A_m$, and is therefore free by Theorem~\ref{thm:2.5}.
\end{proof}

\begin{proof}[Proof of `if']
Let $N_1 \to N_2 \to 0$ be an exact sequence of $A$-modules.
Write $C$ for the cokernel of
\[
  \Hom_A(M, N_1) \to \Hom_A(M, N_2);
\]
then for any maximal ideal $m$ of $A$ we have
\[
  C_m = \Coker\{\Hom_{A_m}(M_m, (N_1)_m) \to \Hom_{A_m}(M_m, (N_2)_m)\} = 0.
\]
Hence $C = 0$ by Theorem~\ref{thm:4.6}, and this is what we had to prove.
\end{proof}

\begin{corollary*}
If $A$ is a ring and $M$ is an $A$-module of finite presentation, then
$M$ is flat if and only if it is projective.
\end{corollary*}

\begin{proof}
This follows from Theorems 1, 12 and 10.
\end{proof}

% ============ 习题 ============
\section*{Exercises to \S\ref{sec:7}}
\emph{Prove the following propositions.}

\begin{exercise}\label{ex:7.1-ex}
If $B$ is a faithfully flat $A$-algebra then for an $A$-module $M$ we have
\[
  B \otimes_A M \text{ is $B$-flat} \iff M \text{ is $A$-flat},
\]
and similarly for faithfully flat.
\end{exercise}

\begin{exercise}\label{ex:7.2-ex}
Let $A$ and $B$ be integral domains with $A \subset B$, and suppose that $A$ and $B$
have the same field of fractions; if $B$ is faithfully flat over $A$ then $A = B$.
\end{exercise}

\begin{exercise}\label{ex:7.3-ex}
Let $B$ be a faithfully flat $A$-algebra; for an $A$-module $M$ we can view $M$ as a
submodule of $B \otimes_A M$ (by Theorem~\ref{thm:7.5}). Then if $\{m_\lambda\}$ is a subset of $M$
which generates $B \otimes M$ over $B$, it also generates $M$ over $A$.
\end{exercise}

\begin{exercise}\label{ex:7.4-ex}
Let $A$ be a Noetherian ring and $\{M_\lambda\}_{\lambda \in \Lambda}$ a family of flat $A$-modules; then
the direct product module $\prod_{\lambda \in \Lambda} M_\lambda$ is also flat. In particular the formal
power series $A[[X_1, \ldots, X_n]]$ is a flat $A$-algebra (Chase [1]).
\end{exercise}

\begin{exercise}\label{ex:7.5-ex}
Let $A$ be a ring and $N$ a flat $A$-module; if $a \in A$ is $A$-regular, it is also $N$-
regular.
\end{exercise}

\begin{exercise}\label{ex:7.6-ex}
Let $A$ be a ring, and $C_\bullet$ a complex of $A$-modules; for an $A$-module $N$ we
write $C_\bullet \otimes N$ for the complex $\cdots \to C_{i+1} \otimes N \to C_i \otimes N \to \cdots$. If
$N$ is flat over $A$ then $H_i(C_\bullet) \otimes N = H_i(C_\bullet \otimes N)$ for all $i$.
\end{exercise}

\begin{exercise}\label{ex:7.7-ex}
Let $A$ be a ring and $B$ a flat $A$-algebra; then if $M$ and $N$ are $A$-modules,
\[
  \Tor^A_i(M, N) \otimes_A B = \Tor^B_i(M \otimes B, N \otimes B) \quad \text{for all } i.
\]
If in addition $M$ is finitely generated and $A$ is Noetherian then
\[
  \Ext^i_A(M, N) \otimes_A B = \Ext^i_B(M \otimes_A B, N \otimes_A B) \quad \text{for all } i.
\]
\end{exercise}

\begin{exercise}\label{ex:7.8-ex}
Theorem~\ref{thm:7.4}, (i) does not hold for the intersection of infinitely many
submodules; explain why, and construct a counter-example.
\end{exercise}

\begin{exercise}\label{ex:7.9-ex}
If $B$ is a faithfully flat $A$-algebra and $B$ is Noetherian then $A$ is Noetherian.
\end{exercise}

% ============ Appendix to §7. Pure submodules (原书 pp.53-54) ============
\section*{Appendix to \S\ref{sec:7}. Pure submodules}
\phantomsection
\addcontentsline{toc}{section}{Appendix to \S\ref{sec:7}. Pure submodules}

Let $A$ be a ring and $M$ an $A$-module. A submodule $N$ of $M$ is said to be \emph{pure}\index{pure submodule}
if the sequence $0 \to N \otimes E \to M \otimes E$ is exact for every $A$-module $E$. Since
tensor product and exactness commute with inductive limits, we need only
consider $A$-modules $E$ of finite presentation.

\begin{example}\label{ex:7.1}
If $M/N$ is a flat $A$-module, then $N$ is a pure submodule of
$M$. This follows from the exact sequence $\Tor^A_1(M/N, E) \to N \otimes E
\to M \otimes E$.
\end{example}

\begin{example}\label{ex:7.2}
Any direct summand of $M$ is a pure submodule.
\end{example}

\begin{example}\label{ex:7.3}
If $A = \Z$, a submodule $N$ of $M$ is pure if and only if $N \cap mM
= mN$ for all $m > 0$. In fact the condition is equivalent to the exactness
of $0 \to N \otimes \Z/m\Z \to M \otimes \Z/m\Z$, and every finitely generated $\Z$-module
is a direct sum of cyclic modules.
\end{example}

\begin{theorem}\label{thm:7.13}
A submodule $N$ of $M$ is pure if and only if the following
condition holds: if $x_i = \sum_{j=1}^{s} a_{ij}m_j$ (for $1 \le i \le r$), with $m_j \in M$, $x_i \in N$ and
$a_{ij} \in A$, then there exist $y_j \in N$ (for $1 \le j \le s$) such that $x_i = \sum_{j=1}^{s} a_{ij}y_j$
(for $1 \le i \le r$).
\end{theorem}

\begin{proof}
Suppose $N$ is pure in $M$. Consider the free module $A^r$ with basis
$e_1, \ldots, e_r$ and let $D$ be the submodule of $A^r$ generated by $\sum_i a_{ij}e_i$, $1 \le j \le s$. Set
$E = A^r/D$, and let $\bar{e}_i$ denote the image of $e_i$ in $E$. Then in $M \otimes E$ we have
\[
  \sum_i x_i \otimes \bar{e}_i = \sum_{i,j} a_{ij}m_j \otimes \bar{e}_i = \sum_j m_j \otimes \left(\sum_i a_{ij}\bar{e}_i\right) = 0,
\]
hence $\sum_i x_i \otimes \bar{e}_i = 0$ in $N \otimes E$ by purity. But this means that, in $N \otimes A^r$, the
element $\sum_i x_i \otimes e_i$ is of the form $\sum_j y_j \otimes \sum_i a_{ij}e_i$ for some $y_j \in N$.

Conversely, suppose the condition is satisfied. Let $E$ be an $A$-module of
finite presentation. Then we can write $E = A^r/D$ with $D$ generated by a finite
number of elements of $A^r$, say $\sum_i a_{ij}e_i$, $1 \le j \le s$. Then reversing the
preceding argument we can see that $N \otimes E \to M \otimes E$ is injective.
\end{proof}

\begin{theorem}\label{thm:7.14}
If $N$ is a pure submodule and $M/N$ is of finite presentation,
then $N$ is a direct summand of $M$.
\end{theorem}

\begin{proof}
We will prove that $0 \to N \xrightarrow{i} M \xrightarrow{p} M/N \to 0$ splits, where $i$
and $p$ are the natural maps. For this we need only construct a linear map
$f\colon M/N \to M$ such that $pf$ is the identity map of $M/N$. Let $\{t_1, \ldots, t_l\}$
be a set of generators of $M/N$, so that $M/N \cong A^l/R$, where $R$ is the
submodule of relations among the $t_j$; let $\{(a_{i1}, \ldots, a_{il}) \mid 1 \le i \le s\}$ be a
set of generators of $R$. Choose a pre-image $\bar{t}_j$ of $t_j$ in $M$ for each $j$. Then
set $\eta_i = \sum_j a_{ij}\bar{t}_j \in N$ (for $1 \le i \le s$). By the preceding theorem there exist
$\zeta_j \in N$ such that $\eta_i = \sum_j a_{ij}\zeta_j$ (for $1 \le i \le s$). Then $\sum_j a_{ij}(\bar{t}_j - \zeta_j) = 0$ (for
$1 \le i \le s$), and setting $f(t_j) = \bar{t}_j - \zeta_j$, we obtain a linear map $f\colon M/N \to M$
which satisfies the requirement.
\end{proof}

% ============ §8 Completion and the Artin-Rees lemma (原书 pp.55-64) ============
\section{Completion and the Artin--Rees lemma}\label{sec:8}\index{completion, section 8}\index{Artin-Rees lemma, section 8}

Let $A$ be a ring and $M$ an $A$-module; for a directed set $\Lambda$, suppose
that $\mathcal{Y} = \{M_\lambda\}_{\lambda \in \Lambda}$ is a family of submodules of $M$ indexed by $\Lambda$ and such
that $\lambda < \mu \Rightarrow M_\lambda \supset M_\mu$. Then taking $\mathcal{Y}$ as a system of neighbourhoods
of $0$ makes $M$ into a topological group under addition. In this topology, for any
$x \in M$ a system of neighbourhoods of $x$ is given by $\{x + M_\lambda\}_{\lambda \in \Lambda}$. In $M$
addition and subtraction are continuous, as is scalar multiplication $x \mapsto ax$
for any $a \in A$. When $M = A$ each $M_\lambda$ is an ideal, so that multiplication is also
continuous:
\[
  (a + M_\lambda)(b + M_\lambda) \subset ab + M_\lambda.
\]
This type of topology is called a \emph{linear topology}\index{linear topology} on $M$; it is separated
(that is, Hausdorff) if and only if $\bigcap_\lambda M_\lambda = 0$. Each $M_\lambda \subset M$ is an open set,
each coset $x + M_\lambda$ is again open, and the complement $M - M_\lambda$ of $M_\lambda$ is a
union of cosets, so is also open. Hence $M_\lambda$ is an open and closed subset;
the quotient module $M/M_\lambda$ is then discrete in the quotient topology.

$M/\bigcap_\lambda M_\lambda$ is called the \emph{separated module} associated with $M$. Moreover,
since for $\lambda < \mu$ there is a natural linear map $\varphi_{\mu\lambda}\colon M/M_\mu \to M/M_\lambda$,
we can construct the inverse system $\{M/M_\lambda; \varphi_{\mu\lambda}\}$ of $A$-modules; its
inverse limit $\varprojlim M/M_\lambda$ is called the \emph{completion} of $M$, and is written $\hat{M}$.
We give each $M/M_\lambda$ the discrete topology, the direct product $\prod_\lambda M/M_\lambda$ the
product topology, and $\hat{M}$ the subspace topology in $\prod_\lambda M/M_\lambda$.
Let $\varphi\colon M \to \hat{M}$ be the natural $A$-linear map; then $\varphi$ is continuous, and $\varphi(M)$ is
dense in $\hat{M}$. Write $p_\lambda\colon \hat{M} \to M/M_\lambda$ for the projection, and set $\Ker p_\lambda =
\hat{M}_\lambda$; it is easy to see that the topology of $\hat{M}$ coincides with the linear
topology defined by $\hat{\mathcal{Y}} = \{\hat{M}_\lambda\}_{\lambda \in \Lambda}$. The map $p_\lambda$ is surjective (in fact
$p_\lambda(\varphi(M)) = M/M_\lambda$), so that $\hat{M}/\hat{M}_\lambda = M/M_\lambda$, and the completion of $\hat{M}$
coincides with $\hat{M}$ itself. If $\varphi\colon M \to \hat{M}$ is an isomorphism, we say
that $M$ is \emph{complete}\index{complete module}. (Caution: in Bourbaki terminology this is `complete and
separated'; we shorten this to `complete' throughout.)

If $\mathcal{Y}' = \{M'_\gamma\}_{\gamma \in \Gamma}$ is another family of submodules of $M$ indexed by a
directed set $\Gamma$, then $\mathcal{Y}$ and $\mathcal{Y}'$ give the same topology on $M$ if and only if for
each $M_\lambda$ there is a $\gamma \in \Gamma$ such that $M'_\gamma \subset M_\lambda$, and for every $M'_\gamma$, there is a $\lambda \in \Lambda$
such that $M_\lambda \subset M'_\gamma$. It is then easy to see that there is an isomorphism of
topological modules $\varprojlim M/M_\lambda \cong \varprojlim M/M'_\gamma$. Thus $\hat{M}$ depends only on
the topology of $M$, as does the question of whether $M$ is complete.

When $M = A$, $\{M/M_\lambda; \varphi_{\mu\lambda}\}$ becomes an inverse system of rings, $\hat{M} = \hat{A}$
is a ring, and $\varphi\colon A \to \hat{A}$ a ring homomorphism. $\hat{M}_\lambda \subset \hat{A}$ is not just
an $A$-submodule, but an ideal of $\hat{A}$; this is clear from the fact that
$p_\lambda\colon \hat{A} \to A/M_\lambda$ is a ring homomorphism.

If $N \subset M$ is a submodule, then the closure $\bar{N}$ of $N$ in $M$ is given by the
following formula:
\[
  \bar{N} = \bigcap_\lambda (N + M_\lambda).
\]
Indeed,
\[
  x \in \bar{N} \iff (x + M_\lambda) \cap N \neq \emptyset \text{ for all } \lambda
\]
\[
  \iff x \in N + M_\lambda \text{ for all } \lambda.
\]
If we write $\bar{M}_\lambda$ for the image of $M_\lambda$ in the quotient module $M/N$, the
quotient topology of $M/N$ is just the linear topology defined by $\{\bar{M}_\lambda\}_{\lambda \in \Lambda}$.
In fact, let $G \subset M$ be the inverse image of $G' \subset M/N$; then
\begin{align*}
  G' \text{ is open in the quotient topology of } M/N
  &\iff G \text{ is open in } M \\
  &\iff \forall x \in G \ \exists M_\lambda,\ x + M_\lambda \subset G \\
  &\iff \forall x' \in G' \ \exists \bar{M}'_\lambda,\ x' + \bar{M}'_\lambda \subset G'.
\end{align*}
Hence the condition for $M/N$ to be separated is that $\bigcap_\lambda \bar{M}_\lambda = 0$, that is
$\bigcap_\lambda (N + M_\lambda) = N$, or in other words, that $N$ is closed in $M$. Moreover, the
subspace topology of $N$ is clearly the same thing as the linear topology
defined by $\{N \cap M_\lambda\}_{\lambda \in \Lambda}$. Set $\bar{M} = M/N$; then
\[
  0 \to N/(N \cap M_\lambda) \to M/M_\lambda \to \bar{M}/\bar{M}_\lambda = M/(N + M_\lambda) \to 0
\]
is an exact sequence, so that taking the inverse limit, we see that
\[
  0 \to \hat{N} \to \hat{M} \to \widehat{(M/N)}
\]
is exact. If we view $N$ as a submodule of $\hat{M}$, the condition that $\xi = (\xi_\lambda)_{\lambda \in \Lambda} \in \hat{M}$
belongs to $\hat{N}$ is that each $\xi_\lambda$ can be represented by an element of $N$, or in
other words that $\xi \in \varphi(N) + \hat{M}_\lambda$ for each $\lambda$. Hence $\hat{N}$ is the same thing as the
closure of $\varphi(N)$ in $\hat{M}$. In general it is not clear whether $\hat{M} \to \widehat{(M/N)}$ is
surjective, but this holds in the case $\Lambda = \{1, 2, \ldots\}$. In fact then
\[
  \widehat{(M/N)} = \varprojlim M/(N + M_n);
\]
given an element $\xi' = (\xi'_1, \xi'_2, \ldots) \in \widehat{(M/N)}$ with $\xi'_n \in M/(N + M_n)$, let
$x_1 \in M$ be an inverse image of $\xi'_1$, and $y_2 \in M$ an inverse image of $\xi'_2$;
then $y_2 - x_1 \in N + M_1$, so that we can write
\[
  y_2 - x_1 = t + m_1 \quad \text{with } t \in N \text{ and } m_1 \in M_1.
\]
If we set $x_2 = y_2 - t$ then $x_2 \in M$ is also an inverse image of $\xi'_2$, and satisfies
$x_2 - x_1 \in M_1$. Similarly we can successively choose inverse images $x_n \in M$ of
the $\xi'_n$ in such a way that for $n = 1, 2, \ldots$, we have $x_{n+1} - x_n \in M_n$. If we set
$\zeta_n \in M/M_n$ for the image of $x_n$, then by construction $\zeta = (\zeta_1, \zeta_2, \ldots)$ is an
element of $\varprojlim M/M_n = \hat{M}$ which maps to $\xi'$ in $\widehat{(M/N)}$. This proves the
following theorem.

\begin{theorem}\label{thm:8.1}
Let $A$ be a ring, $M$ an $A$-module with a linear topology, and
$N \subset M$ a submodule. We give $N$ the subspace topology, and $M/N$ the
quotient topology. Then these are both linear topologies, and we have:
\begin{enumerate}[(i)]
  \item $0 \to \hat{N} \to \hat{M} \to \widehat{(M/N)}$ is an exact sequence, and $\hat{N}$ is the closure
        of $\varphi(N)$ in $\hat{M}$, where $\varphi\colon M \to \hat{M}$ is the natural map.
  \item If moreover the topology of $M$ is defined by a decreasing chain of
        submodules $M_1 \supset M_2 \supset \cdots$, then
        \[
          0 \to \hat{N} \to \hat{M} \to \widehat{(M/N)} \to 0
        \]
        is exact. In other words, $\widehat{(M/N)} = \hat{M}/\hat{N}$.
\end{enumerate}
\end{theorem}

Now suppose that $M$ and $N$ are two $A$-modules with linear topologies,
and let $f\colon M \to N$ be a continuous linear map. If the topologies of $M$ and
$N$ are given by $\{M_\lambda\}_{\lambda \in \Lambda}$ and $\{N_\gamma\}_{\gamma \in \Gamma}$, then for any $\gamma \in \Gamma$ there
exists $\lambda \in \Lambda$ such that $M_\lambda \subset f^{-1}(N_\gamma)$.
Define $\varphi_\gamma\colon \hat{M} \to N/N_\gamma$ as the composite
$\hat{M} \to M/M_\lambda \to N/N_\gamma$, where the first arrow is the natural
map, and the second is induced by $f$; one sees at once that $\varphi_\gamma$ does not
depend on the choice of $\lambda$ for which $M_\lambda \subset f^{-1}(N_\gamma)$.
Also, for $\gamma < \gamma'$ if we
let $\psi_{\gamma\gamma'}$ denote the natural map $N/N_\gamma \to N/N_{\gamma'}$, it is easy to see that
$\varphi_\gamma = \psi_{\gamma\gamma'} \varphi_{\gamma'}$; hence there is a continuous linear map $\hat{f}\colon \hat{M} \to \hat{N}$
defined by the $(\varphi_\gamma)_{\gamma \in \Gamma}$, and the following diagram is commutative (the
vertical arrows are the natural maps):
\[
  \begin{tikzcd}[column sep=large, row sep=large]
    M \arrow[r, "f"] \arrow[d] & N \arrow[d] \\
    \hat{M} \arrow[r, "\hat{f}"] & \hat{N}
  \end{tikzcd}
\]
Moreover, $\hat{f}$ is determined uniquely by this diagram and by continuity.
Similarly, if $A$ and $B$ are rings with linear topologies, and $f\colon A \to B$ is
a continuous ring homomorphism, then $f$ induces a continuous ring
homomorphism $\hat{f}\colon \hat{A} \to \hat{B}$.

Among the linear topologies, those defined by ideals are of particular
importance. Let $I$ be an ideal of $A$ and $M$ an $A$-module; the topology on $M$
defined by $\{I^nM\}_{n=1,2,\ldots}$ is called the \emph{$I$-adic topology}\index{I-adic topology}. If we also give $A$ the $I$-
adic topology, the completions $\hat{A}$ and $\hat{M}$ of $A$ and $M$ are called \emph{$I$-adic
completions}\index{I-adic completion}; it is easy to see that $\hat{M}$ is an $\hat{A}$-module: for $\alpha = (a_1, a_2, \ldots)
\in \hat{A}$ with $a_n \in A/I^n$ and $\xi = (x_1, x_2, \ldots) \in \hat{M}$ with $x_n \in M/I^nM$ (for all $n$), we can
just set
\[
  \alpha\xi = (a_1x_1, a_2x_2, \ldots) \in \hat{M}.
\]

As one can easily check, to say that $M$ is complete for the $I$-adic topology is
equivalent to saying that for every sequence $x_1, x_2, \ldots$ of elements of $M$
satisfying $x_i - x_{i+1} \in I^iM$ for all $i$, there exists a unique $x \in M$ such that
$x - x_i \in I^iM$ for all $i$. We can define a Cauchy sequence in $M$ in the usual
way ($\{x_i\}$ is Cauchy if and only if for every positive integer $r$ there is an
$n_0$ such that $x_{n+1} - x_n \in I^rM$ for $n > n_0$), and completeness can then be
expressed as saying that a Cauchy sequence has a unique limit.

\begin{theorem}\label{thm:8.2}
Let $A$ be a ring, $I$ an ideal, and $M$ an $A$-module.
\begin{enumerate}[(i)]
  \item If $A$ is $I$-adically complete then $I \subset \rad(A)$;
  \item If $M$ is $I$-adically complete and $a \in I$, then multiplication by $1 + a$ is an
        automorphism of $M$.
\end{enumerate}
\end{theorem}

\begin{proof}
(i) For $a \in I$, $1 - a + a^2 - a^3 + \cdots$ converges in $A$, and provides an
inverse of $1 + a$; hence $1 + a$ is a unit of $A$. This means (see \S\ref{sec:1}) that
$I \subset \rad(A)$.

(ii) $\hat{M}$ is also an $\hat{A}$-module, and $1 + a$ (or rather, its image in $\hat{A}$) is a unit in
$\hat{A}$, so that this is clear.
\end{proof}

The following two results show the usefulness of completeness.

\begin{theorem}[Hensel's lemma]\label{thm:8.3}
Let $(A, m, k)$ be a local ring, and suppose
that $A$ is $m$-adically complete. Let $F(X) \in A[X]$ be a monic polynomial, and
let $\bar{F} \in k[X]$ be the polynomial obtained by reducing the coefficients of $F$
modulo $m$. If there are monic polynomials $g, h \in k[X]$ with $(g, h) = 1$
and such that $\bar{F} = gh$, then there exist monic polynomials $G, H$ with coeffi-
cients in $A$ such that $F = GH$, $\bar{G} = g$ and $\bar{H} = h$.
\end{theorem}

\begin{proof}
If we take polynomials $G_1, H_1 \in A[X]$ such that $g = \bar{G}_1$, $h = \bar{H}_1$, then
$F \equiv G_1H_1 \bmod m[X]$. Suppose by induction that monic polynomials $G_n$,
$H_n$ have been constructed such that $F \equiv G_nH_n \bmod m^n[X]$, and $\bar{G}_n = g$,
$\bar{H}_n = h$; then we can write
\[
  F - G_nH_n = \sum_i \delta_iU_i(X), \quad \text{with } \delta_i \in m^n \text{ and } \deg U_i < \deg F.
\]
Since $(g, h) = 1$ we can find $u_i, w_i \in k[X]$ such that $\bar{U}_i = gu_i + hw_i$. Replacing
$u_i$ by its remainder modulo $h$, and making the corresponding correction to
$w_i$ we can assume $\deg u_i < \deg h$. Then
\[
  \deg hw_i = \deg(\bar{U}_i - gu_i) < \deg F, \quad \text{hence } \deg w_i < \deg g.
\]
Choosing $V_i, W_i \in A[X]$ such that $\bar{V}_i = v_i$, $\deg V_i = \deg v_i$, $\bar{W}_i = w_i$,
$\deg W_i = \deg w_i$, and setting $G_{n+1} = G_n + \sum \delta_iW_i$, $H_{n+1} = H_n + \sum \delta_i V_i$,
we get
\[
  F - G_{n+1}H_{n+1} \in m^{n+1}[X].
\]
We construct in this way sequences of polynomials $G_n, H_n$ for $n = 1, 2, \ldots$;
then $\lim G_n = G$ and $\lim H_n = H$ clearly exist and satisfy $F = GH$. Obvi-
ously, $\bar{G} = \bar{G}_1 = g$, $\bar{H} = \bar{H}_1 = h$.
\end{proof}

\begin{theorem}\label{thm:8.4}
Let $A$ be a ring, $I$ an ideal, and $M$ an $A$-module. Suppose
that $A$ is $I$-adically complete, and $M$ is separated for the $I$-adic topology. If
$M/IM$ is generated over $A/I$ by $\bar{\omega}_1, \ldots, \bar{\omega}_n$, and $\omega_i \in M$ is an arbitrary
inverse image of $\bar{\omega}_i$ in $M$, then $M$ is generated over $A$ by $\omega_1, \ldots, \omega_n$.
\end{theorem}

\begin{proof}
By assumption $M = \sum A\omega_i + IM$, so that $M = \sum A\omega_i + I(\sum A\omega_i +
IM) = \sum A\omega_i + I^2M$, and similarly, $M = \sum A\omega_i + I^\nu M$ for all $\nu \ge 0$. For
any $\xi \in M$, write $\xi = \sum a_{i,1}\omega_i + \xi_1$ with $\xi_1 \in IM$, then $\xi_1 = \sum a_{i,2}\omega_i + \xi_2$ with
$a_{i,2} \in I$ and $\xi_2 \in I^2M$, and choose successively $a_{i,\nu} \in I^\nu$ and $\xi_\nu \in I^\nu M$ to satisfy
\[
  \xi_\nu = \sum a_{i,\nu+1}\omega_i + \xi_{\nu+1} \quad \text{for } \nu = 1, 2, \ldots.
\]
Then $a_i + a_{i,2} + a_{i,3} + \cdots$ converges in $A$. If we set $b_i$ for this sum then
\[
  \xi = \sum b_i\omega_i. \qquad\blacksquare
\]
\end{proof}

This theorem is extremely handy for proving the finiteness of $M$. For a
Noetherian ring $A$, the $I$-adic topology has several more important
properties, which are based on the following theorem, proved independ-
ently by E. Artin and D. Rees.

\begin{theorem}[the Artin--Rees lemma]\label{thm:8.5}
Let $A$ be a Noetherian ring, $M$ a
finite $A$-module, $N \subset M$ a submodule, and $I$ an ideal of $A$. Then there exists
a positive integer $c$ such that for every $n > c$, we have
\[
  I^nM \cap N = I^{n-c}(I^cM \cap N).
\]
\end{theorem}

\begin{proof}
The inclusion $\supset$ is obvious, so that we only have to prove $\subset$.
Suppose that $I$ is generated by $r$ elements $a_1, \ldots, a_r$, and $M$ by $s$ elements
$\omega_1, \ldots, \omega_s$. An element of $I^nM$ can be written as $\sum_i f_i(a)\omega_i$, where
$f_i(X) = f_i(X_1, \ldots, X_r)$ is a homogeneous polynomial of degree $n$ with
coefficients in $A$. Now set $A[X_1, \ldots, X_r] = B$, and for each $n > 0$ set
\[
  U_n = \left\{(f_1, \ldots, f_s) \in B^s \;\middle|\;
    \begin{array}{l}
      \text{the } f_i \text{ are homogeneous of degree } n \\
      \text{and } \sum_i f_i(a)\omega_i \in N
    \end{array} \right\};
\]
let $C \subset B^s$ be the $B$-submodule generated by $\bigcup_{n, \gamma} U_n$. Since $B$ is Noetherian,
$C$ is a finite $B$-module, so that $C = \sum_{i=1}^{t} Bu_j$, where each $u_j$ is a linear
combination of elements of $\bigcup_n U_n$; therefore $C$ is generated by finitely
many elements of $\bigcup_n U_n$. Suppose
\[
  C = Bu_1 + \cdots + Bu_t, \quad \text{where } u_j = (u_{j1}, \ldots, u_{js}) \in U_{d_j} \text{ for } 1 \le j \le t.
\]
Set $c = \max(d_1, \ldots, d_t)$. Now if $\eta \in I^nM \cap N$, we can write $\eta = \sum_i f_i(a)\omega_i$ with
$(f_1, \ldots, f_s) \in U_n$, and hence
\[
  (f_1, \ldots, f_s) = \sum_j P_j(X)u_j, \quad \text{with } P_j \in B = A[X_1, \ldots, X_r].
\]
The left-hand side is a vector made up of homogeneous polynomials of
degree $n$ only, so that the terms of degree other than $n$ on the right-hand side
must cancel out to give $0$. Hence we can suppose that the $P_j(X)$ are
homogeneous of degree $n - d_j$. Then $\eta = \sum_i f_i(a)\omega_i = \sum_j P_j(a) \sum_i u_{ji}(a)\omega_i$,
and $\sum_i u_{ji}(a)\omega_i \in I^{d_j}M \cap N$, so that if $n > c$, $P_j(a) \in I^{n-d_j} \supset I^{n-c}$, giving
\[
  \eta \in I^{n-c}(I^cM \cap N) \quad \text{for any } n > c. \qquad\blacksquare
\]
\end{proof}

\begin{theorem}\label{thm:8.6}
In the notation of the above theorem, the $I$-adic topology of
$N$ coincides with the topology induced by the $I$-adic topology of $M$ on the
subspace $N \subset M$.
\end{theorem}

\begin{proof}
By the previous theorem, for $n > c$, we have $I^nN \subset I^nM \cap N
\subset I^{n-c}N$. The topology of $N$ as a subspace of $M$ is the linear topology
defined by $\{I^nM \cap N\}_{n=1,2,\ldots}$, and the above formula says that this defines the
same topology as $\{I^nN\}_{n=1,2,\ldots}$.
\end{proof}

\begin{theorem}\label{thm:8.7}
Let $A$ be a Noetherian ring, $I$ an ideal, and $M$ a finite $A$-
module. Writing $\hat{A}$, $\hat{M}$ for the $I$-adic completions of $M$ and $A$ we have
\[
  \hat{M} = M \otimes_A \hat{A}.
\]
Hence if $A$ is $I$-adically complete, so is $M$.
\end{theorem}

\begin{proof}
By Theorems 1 and 6, the $I$-adic completion of an exact sequence of
finite $A$-modules is again exact. Now given $M$, let $A^n \to A^m \to M \to 0$ be
an exact sequence; the commutative diagram
\[
  \begin{tikzcd}[column sep=large, row sep=large]
    \hat{A}^n \arrow[r] \arrow[d] & \hat{A}^m \arrow[r] \arrow[d] & \hat{M} \arrow[r] \arrow[d] & 0 \\
    A^n \otimes_A \hat{A} \arrow[r] & A^m \otimes_A \hat{A} \arrow[r] & M \otimes_A \hat{A} \arrow[r] & 0
  \end{tikzcd}
\]
has exact rows. Here the vertical arrows are the natural maps; since
completion commutes with direct sums, the two left-hand arrows are
obviously isomorphisms, and hence the right-hand arrow is an
isomorphism, as required.
\end{proof}

\begin{theorem}\label{thm:8.8}
Let $A$ be a Noetherian ring, $I$ an ideal, and $\hat{A}$ the $I$-adic
completion of $A$; then $\hat{A}$ is flat over $A$.
\end{theorem}

\begin{proof}
By Theorem~\ref{thm:7.7} it is enough to show that $a \otimes \hat{A} \to \hat{A}$ is injective
for every ideal $a \subset A$; but $a \otimes \hat{A} = \hat{a}$, and by Theorems 1 and 6, $\hat{a} \to \hat{A}$ is
injective.
\end{proof}

\begin{theorem}[Krull]\label{thm:8.9}
Let $A$ be a Noetherian ring, $I$ an ideal, and $M$
a finite $A$-module; set $\bigcap_{n \ge 0} I^nM = N$. Then there exists $a \in A$ such that
$a \equiv 1 \bmod I$ and $aN = 0$.
\end{theorem}

\begin{proof}
By NAK, it is enough to show that $N = IN$. By the Artin--Rees
lemma, $I^nM \cap N \subset IN$ for sufficiently large $n$; now by definition of $N$, the
left-hand side coincides with $N$.
\end{proof}

\begin{theorem}[the Krull intersection theorem]\label{thm:8.10}
\begin{enumerate}[(i)]
  \item Let $A$ be a Noetherian ring and $I$ an ideal of $A$ with $I \subset \rad A$; then for
        any finite $A$-module the $I$-adic topology is separated, and any submodule is
        a closed set.
  \item If $A$ is a Noetherian integral domain and $I \subset A$ a proper ideal, then
        \[
          \bigcap_{n \ge 1} I^n = (0).
        \]
\end{enumerate}
\end{theorem}

\begin{proof}
(i) In this case the $a$ of the previous theorem is a unit of $A$, so that
$N = 0$, and $M$ is separated. If $M' \subset M$ is a submodule then $M/M'$ is also
$I$-adically separated, which is the same as saying that $M'$ is closed in $M$.

(ii) Setting $M = A$ in the previous theorem, from $1 \notin I$ we get that $a \neq 0$,
and since $a$ is not a zero-divisor, $N = 0$.
\end{proof}

\begin{theorem}\label{thm:8.11}
Let $A$ be a Noetherian ring, $I$ and $J$ ideals of $A$, and $M$ a
finite $A$-module; write $\hat{\;}$ for the completion of an $A$-module in the $I$-adic
topology, and $\varphi\colon M \to \hat{M}$ for the natural map. Then
\[
  \widehat{(JM)} = J\hat{M} = \text{the closure of } \varphi(JM) \text{ in } \hat{M},
\]
and
\[
  \widehat{(M/JM)} = \hat{M}/J\hat{M}.
\]
\end{theorem}

\begin{proof}
By Theorems 1 and 6, $\widehat{(JM)}$ is the kernel of $\hat{M} \to \widehat{(M/JM)}$, and
this is equal to the closure of $\varphi(JM)$ in $\hat{M}$ by Theorem 1. Now suppose $J
= \sum_{i=1}^{r} a_iA$ and define $\varphi\colon M^r \to M$ by $(\xi_1, \ldots, \xi_r) \mapsto \sum a_i\xi_i$. Then the
sequence
\[
  M^r \xrightarrow{\varphi} M \xrightarrow{\mu} M/JM \to 0,
\]
where $\mu$ is the natural map, is exact. The $I$-adic completion,
\[
  \hat{M}^r \xrightarrow{\hat{\varphi}} \hat{M} \xrightarrow{\hat{\mu}} \widehat{(M/JM)} \to 0,
\]
is again exact. On the other hand $\hat{\varphi}$ is given by the same formula
$(\xi_1, \ldots, \xi_r) \mapsto \sum a_i\xi_i$ as $\varphi$, hence $\widehat{(JM)} = \Ker(\hat{\mu}) = \Im(\hat{\varphi}) = \sum a_i\hat{M} = J\hat{M}$.
\end{proof}

As is easily seen, the $(X_1, \ldots, X_n)$-adic completion of $A[X_1, \ldots, X_n]$
is the formal power series ring
\[
  A[[X_1, \ldots, X_n]].
\]
Using this we get the following theorem.

\begin{theorem}\label{thm:8.12}
Let $A$ be a Noetherian ring, and $I = (a_1, \ldots, a_n)$ an ideal of
$A$. Then the $I$-adic completion $\hat{A}$ of $A$ is isomorphic to
$A[[X_1, \ldots, X_n]]/(X_1 - a_1, \ldots, X_n - a_n)$. Hence $\hat{A}$ is a Noetherian ring.
\end{theorem}

\begin{proof}
Let $B = A[X_1, \ldots, X_n]$, and set $I' = \sum X_iB$, $J = \sum (X_i - a_i)B$;
then $B/J \cong A$, and the $I'$-adic topology on $A$ considered as the $B$-module
$B/J$ coincides with the $I$-adic topology of $A$. Now writing $\hat{\;}$ for the $I'$-adic
completion of $B$-modules, we have
\[
  \hat{A} = \widehat{B/J} = \hat{B}/J\hat{B} = A[[X_1, \ldots, X_n]]/(X_1 - a_1, \ldots, X_n - a_n). \qquad\blacksquare
\]
\end{proof}

\begin{theorem}\label{thm:8.13}
Let $A$ be a Noetherian ring, $I$ an ideal, $M$ a finite $A$-module,
and $\hat{M}$ the $I$-adic completion of $M$; then the topology of $\hat{M}$ is the $I$-adic
topology of $\hat{M}$ as an $A$-module, and is the $I\hat{A}$-adic topology of $\hat{M}$ as
an $\hat{A}$-module.
\end{theorem}

\begin{proof}
If we let $\hat{M}_n$ be the kernel of the map from $\hat{M} = \varprojlim (M/I^nM)$ to
$M/I^nM$, the topology of $\hat{M}$ is that defined by $\{\hat{M}_n\}$. Thus it is enough to
prove that $\hat{M}_n = I^n\hat{M}$. Since $M/I^nM$ is discrete in the $I$-adic topology, we
have $\widehat{(M/I^nM)} = M/I^nM$, and the kernel of $\hat{M} \to \widehat{(M/I^nM)}$ is $I^n\hat{M}$ by
Theorem 11. Therefore $\hat{M}_n = I^n\hat{M}$. Moreover, $I^n\hat{M}$ can also be written
$(I^n\hat{A})\hat{M}$, and $I^n\hat{A} = \widehat{I^nA}$, so that the topology of $\hat{M}$ is also the $I\hat{A}$-adic
topology.
\end{proof}

\begin{theorem}\label{thm:8.14}
Let $A$ be a Noetherian ring and $I$ an ideal. If we consider $A$
with the $I$-adic topology, the following conditions are equivalent:
\begin{enumerate}[(1)]
  \item $I \subset \rad(A)$;
  \item every ideal of $A$ is a closed set;
  \item the $I$-adic completion $\hat{A}$ of $A$ is faithfully flat over $A$.
\end{enumerate}
\end{theorem}

\begin{proof}
We have already seen $(1) \Rightarrow (2)$.

$(2) \Rightarrow (3)$ Since $\hat{A}$ is flat over $A$, we need only prove that $m\hat{A} \neq \hat{A}$ for every
maximal ideal $m$ of $A$. By assumption, $(0)$ is closed in $A$, so that we can
assume that $A \subset \hat{A}$, and by Theorem 11, $m\hat{A}$ is the closure of $m$ in $\hat{A}$.
However, $m$ is closed in $A$, so that $m\hat{A} \cap A = m$, and so $m\hat{A} \neq \hat{A}$.

$(3) \Rightarrow (1)$ By Theorem~\ref{thm:7.5}, $m\hat{A} \cap A = m$ for every maximal ideal $m$ of $A$.
Now $m\hat{A} \subset \hat{A}$ is a closed set by Theorems 2, (i) and 10, (i), and since the
natural map $A \to \hat{A}$ is continuous, $m = m\hat{A} \cap A$ is closed in $A$. If $I \nsubseteq m$,
then $I^n + m = A$ for every $n > 0$, so that $m$ is not closed. Thus $I \subset m$.
\end{proof}

If the conditions of the above theorem are satisfied, the topological ring $A$
is said to be a \emph{Zariski ring}\index{Zariski ring}, and $I$ an \emph{ideal of definition} of $A$. An
ideal of definition is not uniquely determined; any ideal defining the same
topology will do. The most important example of a Zariski ring is a Noetherian
local ring $(A, m)$ with the $m$-adic topology. When discussing the completion of a
local ring, we will mean the $m$-adic completion unless otherwise specified.

\begin{theorem}\label{thm:8.15}
Let $A$ be a semilocal ring with maximal ideals $m_1, \ldots, m_n$,
and set $I = \rad(A) = m_1m_2 \cdots m_n$. Then the $I$-adic completion $\hat{A}$ of $A$
decomposes as a direct product
\[
  \hat{A} \cong \hat{A}_1 \times \cdots \times \hat{A}_n,
\]
where $A_i = A_{m_i}$, and $\hat{A}_i$ is the completion of the local ring $A_i$.
\end{theorem}

\begin{proof}
Since for $i \neq j$ and any $n > 0$ we have $m_i^n + m_j^n = A$, Theorem~\ref{thm:1.4}
gives
\[
  A/I^n = A/m_1^n \times \cdots \times A/m_n^n \quad \text{for } n > 0.
\]
Hence taking the limit we get
\[
  \hat{A} = \varprojlim A/I^n = \left(\varprojlim A/m_1^n\right) \times \cdots \times \left(\varprojlim A/m_n^n\right).
\]
If we set $A_i$ for the localisation of $A$ at $m_i$, then, since $A/m_i^n$ is already local,
\[
  A/m_i^n = (A/m_i^n)_{m_i} = A_i/m_i^nA_i,
\]
and so $\varprojlim A/m_i^n$ can be identified with $\hat{A}_i$.
\end{proof}

We now summarise the main points proved in this section for a local
Noetherian ring. Let $(A, m)$ be a local Noetherian ring; then we have:
\begin{enumerate}[(1)]
  \item $\bigcap_{n \ge 0} m^n = (0)$.
  \item For $M$ a finite $A$-module and $N \subset M$ a submodule,
        \[
          \bigcap_{n \ge 0} (N + m^nM) = N.
        \]
  \item The completion $\hat{A}$ of $A$ is faithfully flat over $A$; hence $A \subset \hat{A}$, and
        $I\hat{A} \cap A = I$ for any ideal $I$ of $A$.
  \item $\hat{A}$ is again a Noetherian local ring, with maximal ideal $m\hat{A}$, and it has
        the same residue class field as $A$; moreover, $\hat{A}/m^n\hat{A} = A/m^n$ for all $n > 0$.
  \item If $A$ is a complete local ring, then for any ideal $I \neq A$, $A/I$ is again a
        complete local ring.
\end{enumerate}

\begin{remark}
Even if $A$ is complete, the localisation $A_p$ of $A$ at a prime $p$ may
not be.
\end{remark}

\begin{remark}
An Artinian local ring $(A, m)$ is complete; in fact, it is clear from
the proof of Theorem~\ref{thm:3.2} that there exists a $\nu$ such that $m^\nu = 0$, so that
$A = \varprojlim A/m^n = A$.
\end{remark}

% ============ 习题 ============
\section*{Exercises to \S\ref{sec:8}}
\emph{Prove the following propositions.}

\begin{exercise}\label{ex:8.1-ex}
If $A$ is a Noetherian ring, $I$ and $J$ are ideals of $A$, and $A$ is complete both for
the $I$-adic and $J$-adic topologies, then $A$ is also complete for the $(I + J)$-
adic topology.
\end{exercise}

\begin{exercise}\label{ex:8.2-ex}
Let $A$ be a Noetherian ring, and $I, J$ ideals of $A$; if $A$ is $I$-adically
complete, it is also $J$-adically complete.
\end{exercise}

\begin{exercise}\label{ex:8.3-ex}
Let $A$ be a Zariski ring and $\hat{A}$ its completion. If $a \subset A$ is an ideal such that
$a\hat{A}$ is principal, then $a$ is principal.
\end{exercise}

\begin{exercise}\label{ex:8.4-ex}
According to Theorem~\ref{thm:8.12}, if $y \in \bigcap_{n} I^n$ then
\[
  y \in \sum_{i=1}^{n} (X_i - a_i)A[[X_1, \ldots, X_n]].
\]
Verify this directly in the special case $I = eA$, where $e^2 = e$.
\end{exercise}

\begin{exercise}\label{ex:8.5-ex}
Let $A$ be a Noetherian ring and $I$ a proper ideal of $A$; consider the
multiplicative set $S = 1 + I$ as in \S\ref{sec:4}, Example 3. Then $A_S$ is a Zariski ring
with ideal of definition $IA_S$, and its completion coincides with the $I$-adic
completion of $A$.
\end{exercise}

\begin{exercise}\label{ex:8.6-ex}
If $A$ is $I$-adically complete then $B = A[X]$ is $(IB + XB)$-adically complete.
\end{exercise}

\begin{exercise}\label{ex:8.7-ex}
Let $(A, m)$ be a complete Noetherian local ring, and $a_1 \supset a_2 \supset \cdots$ a chain of
ideals of $A$ for which $\bigcap_\nu a_\nu = (0)$; then for each $n$ there exists $\nu(n)$ for which
$a_{\nu(n)} \subset m^n$. In other words, the linear topology defined by $\{a_\nu\}_{\nu=1,2,\ldots}$ is
stronger than or equal to the $m$-adic topology (Chevalley's theorem).
\end{exercise}

\begin{exercise}\label{ex:8.8-ex}
Let $A$ be a Noetherian ring, $a_1, \ldots, a_r$ ideals of $A$; if $M$ is a finite $A$-module
and $N \subset M$ a submodule, then there exists $c > 0$ such that
\[
  a_1^{n_1} \cdots a_r^{n_r}M \cap N \subset a_1^{n_1-c} \cdots a_r^{n_r-c}(a_1^c \cdots a_r^c M \cap N).
\]
\end{exercise}

\begin{exercise}\label{ex:8.9-ex}
Let $A$ be a Noetherian ring and $P \in \Ass(A)$. Then there is an integer $c > 0$
such that $P \in \Ass(A/I)$ for every ideal $I \subset P^c$ (hint: localise at $P$).
\end{exercise}

\begin{exercise}\label{ex:8.10-ex}
Show by example that the conclusion of Ex.~\ref{ex:8.7-ex} does not necessarily hold
if $A$ is not complete.
\end{exercise}

% ============ §9 Integral extensions (原书 pp.64-70) ============
\section{Integral extensions}\label{sec:9}\index{integral extensions, section 9}

If $A$ is a subring of a ring $B$ we say that $B$ is an \emph{extension ring} of $A$.
In this case, an element $b \in B$ is said to be \emph{integral}\index{integral element} over $A$ if $b$ is a root of
a monic polynomial with coefficients in $A$, that is if there is a relation of the
form $b^n + a_1b^{n-1} + \cdots + a_n = 0$ with $a_i \in A$. If every element of $B$ is integral
over $A$ we say that $B$ is \emph{integral over $A$}, or that $B$ is an \emph{integral extension} of $A$.

\begin{theorem}\label{thm:9.1}
Let $A$ be a ring and $B$ an extension of $A$.
\begin{enumerate}[(i)]
  \item An element $b \in B$ is integral over $A$ if and only if there exists a ring $C$
        with $A \subset C \subset B$ and $b \in C$ such that $C$ is finitely generated as an $A$-module.
  \item Let $\bar{A} \subset B$ be the set of elements of $B$ integral over $A$; then $\bar{A}$ is
        a subring of $B$.
\end{enumerate}
\end{theorem}

\begin{proof}
(i) If $b$ is a root of $f(X) = X^n + a_1X^{n-1} + \cdots + a_n$, for any $P(X) \in
A[X]$ let $r(X)$ be the remainder of $P$ on dividing by $f$; then $P(b) = r(b)$
and $\deg r < n$. Hence
\[
  A[b] = A + Ab + \cdots + Ab^{n-1},
\]
so that we can take $C$ to be $A[b]$. Conversely if an extension ring $C$ of $A$ is a
finite $A$-module then every element of $C$ is integral over $A$: for if $C = A\omega_1
+ \cdots + A\omega_n$ and $b \in C$ then
\[
  b\omega_i = \sum_j a_{ij}\omega_j \quad \text{with } a_{ij} \in A,
\]
so that by Theorem~\ref{thm:2.1} we get a relation $b^n + a_1b^{n-1} + \cdots + a_n = 0$. (The left-
hand side is obtained by expanding out $\det(b\delta_{ij} - a_{ij})$.)

(ii) If $b, b' \in \bar{A}$ then we see easily that $A[b, b']$ is finitely generated as an $A$-
module, so that its elements $bb'$ and $b + b'$ are integral over $A$.
\end{proof}

The $\bar{A}$ appearing in (ii) above is called the \emph{integral closure}\index{integral closure} of $A$ in $B$; if
$A = \bar{A}$ we say that $A$ is \emph{integrally closed} in $B$. In particular, if $A$ is an integral
domain, and is integrally closed in its field of fractions, we say that $A$ is an
\emph{integrally closed domain}\index{integrally closed domain}. If for every prime ideal $p$ of $A$ the localisation $A_p$ is
an integrally closed domain we say that $A$ is a \emph{normal ring}\index{normal ring}.

\begin{remark}
`Normal ring' is often used to mean `integrally closed domain'; in
this book we follow the usage of Serre and Grothendieck. If $A$ is a
Noetherian ring which is normal in our sense, and $p_1, \ldots, p_r$ are all
the minimal prime ideals of $A$ then it can be seen (see Ex.~\ref{ex:9.11-ex}) that
$A = A/p_1 \times \cdots \times A/p_r$, and then each $A/p_i$ is an integrally closed domain
(see Theorem~\ref{thm:4.7}). Conversely, the direct product of a finite number of
integrally closed domains is normal (see Example 3 below).
\end{remark}

Let $A \subset C \subset B$ be a chain of ring extensions; if an element $b \in B$ is integral
over $C$ and $C$ is integral over $A$ then $b$ is integral over $A$. Indeed, if
$b^n + c_1b^{n-1} + \cdots + c_n = 0$ with $c_i \in C$ then
\[
  A[c_1, \ldots, c_n, b] = \sum_{\nu=0}^{n-1} A[c_1, \ldots, c_n]b^\nu,
\]
and since $A[c_1, \ldots, c_n]$ is a finite $A$-module, so is $A[c_1, \ldots, c_n, b]$. In
particular, if we take $C$ to be the integral closure $\bar{A}$ of $A$ in $B$ we see that
$\bar{A}$ is integrally closed in $B$.

\begin{example}\label{ex:9.1}
A UFD is an integrally closed domain -- the proof is easy.
\end{example}

\begin{example}\label{ex:9.2}
Let $k$ be a field and $t$ an indeterminate over $k$; set $A =
k[t^2, t^3] \subset B = k[t]$. Then $A$ and $B$ both have the same field of fractions
$K = k(t)$. Since $B$ is a UFD, it is integrally closed; but $t$ is integral over
$A$, so that $B$ is the integral closure of $A$ in $K$.

Note that in this example $A \cong k[X, Y]/(Y^2 - X^3)$. Thus $A$ is the
coordinate ring of the plane curve $Y^2 = X^3$, which has a singularity at the
origin. The fact that $A$ is not integrally closed is related to the existence of
this singularity.
\end{example}

\begin{example}\label{ex:9.3}
If $B$ is an extension ring of $A$, $S \subset A$ is a multiplicative set, and $\bar{A}$
is the integral closure of $A$ in $B$, then the integral closure of $A_S$ in $B_S$ is $\bar{A}_S$.
The proof is again easy. It follows from this that if $A$ is an integrally closed
domain, so is $A_S$.
\end{example}

\begin{theorem}\label{thm:9.2}
Let $A$ be an integrally closed domain, $K$ the field of fractions
of $A$, and $L$ an algebraic extension of $K$. Then an element $\alpha \in L$ is integral
over $A$ if and only if its minimal polynomial over $K$ has all its coefficients
in $A$.
\end{theorem}

\begin{proof}
Let $f(X) = X^n + a_1X^{n-1} + \cdots + a_n$ be the minimal polynomial of $\alpha$
over $K$. We have $f(\alpha) = 0$, so that if all the $a_i$ are in $A$ then $\alpha$ is integral over
$A$. Conversely, if $\alpha$ is integral over $A$, then letting $\bar{L}$ be an algebraic closure of
$L$ we have a splitting $f(X) = (X - \alpha_1) \cdots (X - \alpha_n)$ of $f(X)$ in $\bar{L}[X]$ into linear
factors; each of the $\alpha_i$ is conjugate to $\alpha$ over $K$, so that there is an
isomorphism $K[\alpha] \cong K[\alpha_i]$ fixing the elements of $K$ and taking $\alpha$ into $\alpha_i$,
and therefore the $\alpha_i$ are also integral over $A$. Then $a_1, \ldots, a_n \in A[\alpha_1, \ldots, \alpha_n]$,
and hence they are integral over $A$; but $a_i \in K$ and $A$ is integrally closed,
so that finally $a_i \in A$.
\end{proof}

\begin{example}\label{ex:9.4}
Let $A$ be a UFD in which $2$ is a unit. Let $f \in A$ be square-free,
(that is, not divisible by the square of any prime of $A$). Then $A[\sqrt{f}]$ is an
integrally closed domain.
\end{example}

\begin{proof}
Let $\alpha$ be a square root of $f$. Let $K$ be the field of fractions of $A$; then
$A$ is integrally closed in $K$ by Example 1, so that if $\alpha \in K$ we have $\alpha \in A$ and
$A[\alpha] = A$, and the assertion is trivial. If $\alpha \notin K$ then the field of fractions of
$A[\alpha]$ is $K(\alpha) = K + K\alpha$, and every element $\xi \in K(\alpha)$ can be written in a
unique way as $\xi = x + y\alpha$ with $x, y \in K$. The minimal polynomial of $\xi$ over $K$
is $X^2 - 2xX + (x^2 - y^2f)$, so that using the previous theorem, if $\xi$ is integral
over $A$ we get $2x \in A$ and $x^2 - y^2f \in A$. By assumption, $2x \in A$ implies $x \in A$.
Hence $y^2f \in A$. From this, if some prime $p$ of $A$ divides the denominator of $y$
we get $p^2 \mid f$, which contradicts the square-free hypothesis. Thus $y \in A$, and
$\xi \in A + A\alpha = A[\alpha]$, so that $A[\alpha]$ is integrally closed in $K(\alpha)$.
\end{proof}

\begin{lemma}\label{lem:9.1}
Let $B$ be an integral domain and $A \subset B$ a subring such that $B$ is
integral over $A$. Then
\[
  A \text{ is a field} \iff B \text{ is a field}.
\]
\end{lemma}

\begin{proof}
($\Rightarrow$) If $0 \neq b \in B$ then there is a relation of the form $b^n + a_1b^{n-1} + \cdots
+ a_n = 0$ with $a_i \in A$, and since $B$ is an integral domain we can assume $a_n \neq 0$.
Then
\[
  b^{-1} = -a_n^{-1}(b^{n-1} + a_1b^{n-2} + \cdots + a_{n-1}) \in B.
\]
($\Leftarrow$) If $0 \neq a \in A$ then $a^{-1} \in B$, so that there is a relation $a^{-n} + c_1a^{-n+1}
+ \cdots + c_n = 0$ with $c_i \in A$. Then
\[
  a^{-1} = -(c_1 + c_2a + \cdots + c_na^{n-1}) \in A. \qquad\blacksquare
\]
\end{proof}

\begin{lemma}\label{lem:9.2}
Let $A$ be a ring, and $B$ an extension ring which is integral over $A$.
If $P$ is a maximal ideal of $B$ then $P \cap A$ is a maximal ideal of $A$. Conversely if
$p$ is a maximal ideal of $A$ then there exists a prime ideal $P$ of $B$ lying over $p$,
and any such $P$ is a maximal ideal of $B$.
\end{lemma}

\begin{proof}
For $P \in \Spec B$ let $P \cap A = p$; then the extension $A/p \subset B/P$ is
integral. Thus by Lemma~\hyperref[lem:9.1]{1} above, $P$ is maximal if and only if $p$ is maximal.
Next, to prove that there exists $P$ lying over a given maximal ideal $p$ of $A$, it
is enough to prove that $pB \neq B$. For then any maximal ideal $P$ of $B$
containing $pB$ will satisfy $P \cap A \supset p$ and $1 \notin P \cap A$, so that $P \cap A = p$. By
contradiction, assume that $pB = B$; then there is an expression $1 = \sum_{i=1}^{n} \pi_ib_i$
with $b_i \in B$ and $\pi_i \in p$. If we set $C = A[b_1, \ldots, b_n]$ then $C$ is finite over $A$ and
$pC = C$. Letting $C = A\omega_1 + \cdots + A\omega_n$ we get $\omega_i = \sum_j \pi_{ij}\omega_j$ for some $\pi_{ij} \in p$,
so that $\Delta = \det(\delta_{ij} - \pi_{ij})$ satisfies $\Delta\omega_j = 0$ for each $j$, and hence $\Delta C = 0$. But
$1 \in C$, so that $\Delta = 0$, and on the other hand $\Delta \equiv 1 \bmod p$; therefore $1 \in p$,
which is a contradiction.
\end{proof}

\begin{theorem}\label{thm:9.3}
Let $A$ be a ring, $B$ an extension ring which is integral over
$A$ and $p$ a prime ideal of $A$.
\begin{enumerate}[(i)]
  \item There exists a prime ideal of $B$ lying over $p$.
  \item There are no inclusions between prime ideals of $B$ lying over $p$.
  \item Let $A$ be an integrally closed domain, $K$ its field of fractions, and $L$ a
        normal field extension of $K$ in the sense of Galois theory (that is $K \subset L$ is
        algebraic, and for any $\alpha \in L$, all the conjugates of $\alpha$ over $K$ are in $L$); if $B$ is the
        integral closure of $A$ in $L$ then all the prime ideals of $B$ lying over $p$
        are conjugate over $K$.
\end{enumerate}
\end{theorem}

\begin{proof}
Localising the exact sequence $0 \to A \to B$ at $p$ gives an exact
sequence $0 \to A_p \to B_p = B \otimes_A A_p$ in which $B_p$ is an extension ring
integral over $A_p$. From the commutative diagram
\[
  \begin{tikzcd}[column sep=large, row sep=large]
    A \arrow[r] \arrow[d] & B \arrow[d] \\
    A_p \arrow[r] & B_p
  \end{tikzcd}
\]
we see that the prime ideals of $B$ lying over $p$ correspond bijectively with the
maximal ideals of $B_p$ lying over the maximal ideal $pA_p$ of $A_p$. Hence, to prove
(i) and (ii) it is enough to consider the case that $p$ is a maximal ideal, which
has already been done in Lemma~\hyperref[lem:9.2]{2}.

Now for (iii). Let $P_1$ and $P_2$ be prime ideals of $B$ lying over $p$. First of all
we consider the case $[L:K] < \infty$; let $G = \{\sigma_1, \ldots, \sigma_l\}$ be the group of $K$-
automorphisms of $L$. If $P_2 \neq \sigma_j^{-1}(P_1)$ for any $j$ then by (ii) we have
$P_2 \not\subset \sigma_j^{-1}(P_1)$, so that there is an element $x \in P_2$ not contained in any
$\sigma_j^{-1}(P_1)$ for $1 \le j \le l$ (see Ex.~\ref{ex:1.6-ex}). Set $y = (\prod_j \sigma_j(x))^q$, where $q = 1$ if $\operatorname{char}
K = 0$, and $q = p^\nu$ for a sufficiently large integer $\nu$ if $\operatorname{char} K = p > 0$. Then
$y \in K$, and is integral over $A$, so that $y \in A$. However, the identity map of $L$ is
contained among the $\sigma_j$, so that $y \in P_2$, and hence $y \in P_2 \cap A = p \subset P_1$. This
contradicts $\sigma_j(x) \notin P_1$ for all $j$. Therefore $P_2 = \sigma_j^{-1}(P_1)$ for some $j$.

If $[L:K] = \infty$ we need Galois theory for infinite extensions. Let $K' \subset L$ be
the fixed subfield of $G = \Aut(L/K)$; then $L$ is Galois over $K'$ and $K \subset K'$ is a
purely inseparable extension. If $K' \neq K$ we must have $\operatorname{char} K = p > 0$, and
setting $A'$ for the integral closure of $A$ in $K'$ we see easily that
\[
  p' = \{x \in A' \mid x^q \in p \text{ for some } q = p^r\}
\]
is the unique prime ideal of $A'$ lying over $p$. Thus replacing $K$ by $K'$ we can
assume that $L$ is a Galois extension of $K$. For any finite Galois extension
$K \subset L'$ contained in $L$ we now set
\[
  F(L') = \{\sigma \in G \mid \sigma(P_1) = P_2 \text{ in } L'\},
\]
then by the case of finite extensions we have just proved, $F(L') \neq \emptyset$.
Moreover, $F(L') \subset G$ is closed in the Krull topology. (Recall that the Krull
topology of $G$ is the topology induced by the inclusion of $G$ into the direct
product of finite groups $\prod_{L'} \Aut(L'/K)$; with respect to this topology, $G$ is
compact. For details see textbooks on field theory.) If $L_i$ for $1 \le i \le n$ are
finite Galois extensions of $K$ contained in $L$ then their composite $L''$ is
also a finite Galois extension of $K$, and $\bigcap_i F(L_i) \supset F(L'') \neq \emptyset$, so that the
family $\{F(L') \mid L' \subset L \text{ is a finite Galois extension of } K\}$ of closed subsets
of $G$ has the finite intersection property; since $G$ is compact,
$\bigcap F(L') \neq \emptyset$. Taking $\sigma \in \bigcap F(L')$ we obviously have $\sigma(P_1) = P_2$.
\end{proof}

For a ring $A$ and an $A$-algebra $B$, the following statement is called the
\emph{going-up theorem}\index{going-up theorem}: given two prime ideals $p \subset p'$ of $A$ and a prime ideal $P$ of $B$
lying over $p$, there exists $P' \in \Spec B$ such that $P \subset P'$ and $P' \cap A = p'$.
Similarly, the \emph{going-down theorem}\index{going-down theorem} is the following statement: given $p \subset p'$
and $P' \in \Spec B$ lying over $p'$, there exists $P \in \Spec B$ such that $P \subset P'$ and
$P \cap A = p$.

\begin{theorem}\label{thm:9.4}
\begin{enumerate}[(i)]
  \item If $B \supset A$ is an extension ring which is integral over $A$ then
        the going-up theorem holds.
  \item If in addition $B$ is an integral domain and $A$ is integrally closed, the
        going-down theorem also holds.
\end{enumerate}
\end{theorem}

\begin{proof}
(i) Suppose $p \subset p'$ and $P$ are given as above. Since $P \cap A = p$, we can
think of $B/P$ as an extension ring of $A/p$, and it is integral over $A/p$ because
the condition that an element of $B$ is integral over $A$ is preserved by the
homomorphism $B \to B/P$. By (i) of the previous theorem there is a prime
ideal of $B/P$ lying over $p'/p$, and writing $P'$ for its inverse image in $B$ we have
$P' \in \Spec B$ and $P' \cap A = p'$.

(ii) Let $K$ be the field of fractions of $A$, and let $L$ be a normal extension
field of $K$ containing $B$; set $C$ for the integral closure of $A$ in $L$. Suppose
given prime ideals $p \subset p'$ of $A$ and $P' \in \Spec B$ such that $P' \cap A = p'$, and
choose $Q' \in \Spec C$ such that $Q' \cap B = P'$. Choose also a prime ideal $Q$ of $C$
over $p$, so that using the going-up theorem we can find a prime ideal $Q_1$ of $C$
containing $Q$ and lying over $p'$. Both $Q_1$ and $Q'$ lie over $p'$, so that by (iii) of
the previous theorem there is an automorphism $\sigma \in \Aut(L/K)$ such that
$\sigma(Q_1) = Q'$. Setting $\sigma(Q) = Q_2$ we have $Q_2 \subset Q'$, and $Q_2 \cap A = Q \cap A = p$,
so that setting $P = Q_2 \cap B$ we get $P \cap A = p$, $P \subset Q' \cap B = P'$. (For
a different proof of (ii) which does not use Theorem 3, (iii), see [AM], (5.16),
or [Kunz].)
\end{proof}

We now treat another important case in which the going-down theorem
holds.

\begin{theorem}\label{thm:9.5}
Let $A$ be a ring and $B$ a flat $A$-algebra; then the going-down
theorem holds between $A$ and $B$.
\end{theorem}

\begin{proof}
Let $p \subset p'$ be prime ideals of $A$, and let $P'$ be a prime ideal of $B$ lying
over $p'$; then $B_{P'}$ is faithfully flat over $A_{p'}$, so that by Theorem~\ref{thm:7.3}
$\Spec(B_{P'}) \to \Spec(A_{p'})$ is surjective. Hence there is a prime ideal $Q_1$
of $B_{P'}$ lying over $pA_{p'}$, and setting $Q_1 \cap B = P$ we obviously have $P \subset P'$
and $P \cap A = p$.
\end{proof}

\begin{theorem}\label{thm:9.6}
Let $A \subset B$ be integral domains such that $A$ is integrally closed
and $B$ is integral over $A$; then the canonical map $f\colon \Spec B \to \Spec A$ is
open. More precisely, for $t \in B$, let $X^n + a_1X^{n-1} + \cdots + a_n$ be a monic poly-
nomial with coefficients in $A$ having $t$ as a root and of minimal degree; then
\[
  f(D(t)) = \bigcup_{i=1}^{n} D(a_i),
\]
where the notation $D(\;)$ is as in \S\ref{sec:4}.
\end{theorem}

\begin{proof}[Proof (H. Seydi [4])]
By Theorem 2, $F(X) = X^n + a_1X^{n-1} + \cdots + a_n$ is
irreducible over the field of fractions of $A$; if we set $C = A[t]$ then
$C \cong A[X]/(F(X))$, so $C$ is a free $A$-module with basis $1, t, t^2, \ldots, t^{n-1}$, and
is hence faithfully flat over $A$. Suppose that $P \in D(t)$, so that $P \in \Spec B$
with $t \notin P$, and set $p = P \cap A$; then $p \in \bigcup_i D(a_i)$, since otherwise $a_i \in p$ for
all $i$, and so $t^n \in P$, hence $t \in P$, which is a contradiction.

Conversely, given $p \in \bigcup_i D(a_i)$, suppose that $t \in \sqrt{pC}$; then for sufficiently
large $m$ we have $t^m = \sum_{i=1}^{n} b_it^{n-i}$ with $b_i \in p$. We can take $m > n$. Then $X^m
- \sum b_iX^{n-i}$ is divisible by $F(X)$ in $A[X]$, which implies that $\bar{X}^m$ is divisible
by $\bar{F}(X) = X^n + \sum \bar{a}_iX^{n-i}$ in $(A/p)[X]$; since at least one of the $\bar{a}_i$ is non-
zero, this is a contradiction. Thus $t \notin \sqrt{pC}$, so that there exists $Q \in \Spec C$
with $t \notin Q$ and $pC \subset Q$. Setting $Q \cap A = q$ we have $p \subset q$, so that by
the previous theorem there exists $P_1 \in \Spec C$ satisfying $P_1 \cap A = p$ and $P_1 \subset Q$.
Since $B$ is integral over $C$ there exists $P \in \Spec B$ lying over $P_1$. We have
$P \in D(t)$, since otherwise $t \in P \cap C = P_1 \subset Q$, which contradicts $t \notin Q$. This
proves that
\[
  f(D(t)) \supset \bigcup_{i=1}^{n} D(a_i).
\]
Any open set of $\Spec B$ is a union of open sets of the form $D(t)$, and hence
$f\colon \Spec B \to \Spec A$ is open.
\end{proof}

% ============ 习题 ============
\section*{Exercises to \S\ref{sec:9}}
\emph{Prove the following propositions.}

\begin{exercise}\label{ex:9.1-ex}
Let $A$ be a ring, $A \subset B$ an integral extension, and $p$ a prime ideal of $A$.
Suppose that $B$ has just one prime ideal $P$ lying over $p$; then $B_P = B_p$.
\end{exercise}

\begin{exercise}\label{ex:9.2-ex}
Let $A$ be a ring and $A \subset B$ an integral extension ring. Then $\dim A =
\dim B$.
\end{exercise}

\begin{exercise}\label{ex:9.3-ex}
Let $A$ be a ring, $A \subset B$ a finitely generated integral extension of $A$, and $p$ a
prime ideal of $A$. Then $B$ has only a finite number of prime ideals lying over
$p$.
\end{exercise}

\begin{exercise}\label{ex:9.4-ex}
Let $A$ be an integral domain and $K$ its field of fractions. We say that $x \in K$ is
\emph{almost integral} over $A$ if there exists $0 \neq a \in A$ such that $ax^n \in A$ for all $n > 0$.
If $x$ is integral over $A$ it is almost integral, and if $A$ is Noetherian the
converse holds.
\end{exercise}

\begin{exercise}\label{ex:9.5-ex}
Let $A \subset K$ be as in the previous question. Say that $A$ is \emph{completely}
\emph{integrally closed} if every $x \in K$ which is almost integral over $A$ belongs to $A$.
If $A$ is completely integrally closed, so is $A[[X]]$.
\end{exercise}

\begin{exercise}\label{ex:9.6-ex}
Let $A$ be an integrally closed domain, $K$ its field of fractions, and let
$f(X) \in A[X]$ be a monic polynomial. Then if $f(X)$ is reducible in $K[X]$ it is
also reducible in $A[X]$.
\end{exercise}

\begin{exercise}\label{ex:9.7-ex}
Let $m \in \Z$ be square-free, and write $A$ for the integral closure of $\Z$ in
$\Q(\sqrt{m})$. Then $A = \Z[(1 + \sqrt{m})/2]$ if $m \equiv 1 \bmod 4$, and $A =
\Z[\sqrt{m}]$ otherwise.
\end{exercise}

\begin{exercise}\label{ex:9.8-ex}
Let $A$ be a ring and $A \subset B$ an integral extension. If $P$ is a prime ideal of $B$
with $p = P \cap A$ then $\htop P \le \htop p$.
\end{exercise}

\begin{exercise}\label{ex:9.9-ex}
Let $A$ be a ring and $B$ an $A$-algebra, and suppose that the going-down
theorem holds between $A$ and $B$. If $P$ is a prime ideal of $B$ with $p = P \cap A$
then $\htop P \ge \htop p$.
\end{exercise}

\begin{exercise}\label{ex:9.10-ex}
Let $K$ be a field and $L$ an extension field of $K$. If $P$ is a prime ideal of
$L[X_1, \ldots, X_n]$ and $p = P \cap K[X_1, \ldots, X_n]$ then $\htop P \le \htop p$, and equality
holds if $L$ is an algebraic extension of $K$. Moreover, if two polynomials
$f(X), g(X) \in K[X_1, \ldots, X_n]$ have no common factors in $K[X_1, \ldots, X_n]$,
they have none in $L[X_1, \ldots, X_n]$.
\end{exercise}

\begin{exercise}\label{ex:9.11-ex}
Let $A$ be a Noetherian ring and $p_1, \ldots, p_r$ all the minimal prime ideals of $A$.
Suppose that $A_P$ is an integral domain for all $P \in \Spec A$. Then (i) $\Ass A
= \{p_1, \ldots, p_r\}$; (ii) $p_1 \cap \cdots \cap p_r = \nil(A) = 0$; (iii) $p_i + \bigcap_{j \neq i} p_j = A$ for all $i$. It
follows that $A \cong A/p_1 \times \cdots \times A/p_r$.
\end{exercise}


% ============ Chapter 4. Valuation rings (原书 pp.71-92) ============
\chapter{Valuation rings}\label{ch:4}

From Hensel's theory of $p$-adic numbers onwards, valuation theory has
been an important tool of number theory and the theory of function fields in
one variable; the main object of study was however the multiplicative
valuations which generalise the usual notion of absolute value of a number.
In contrast, Krull defined and studied valuation rings from a more ring-
theoretic point of view (\cite{Kr3}, 1931). His theory was immediately applied to
algebraic geometry by Zariski. In \S\ref{sec:10} we treat the elementary parts of their
theory. Discrete valuation rings (DVRs) and Dedekind rings, the classical
objects of study, are treated in the following \S\ref{sec:11}, which also includes the
Krull--Akizuki theorem, so that this section contains the theory of one-
dimensional Noetherian rings. In \S\ref{sec:12} we treat Krull rings, which should be
thought of as a natural extension of Dedekind rings; we go as far as a recent
theorem of J.\ Nishimura. This book is mainly concerned with Noetherian
rings, and general valuation rings and Krull rings are the most important
rings outside this category. The present chapter is intended as
complementary to the theory of Noetherian rings, and we have left out
quite a lot on valuation theory. The reader should consult [B6,7],~\cite{ZS} or
other textbooks for more information.

% ============ §10 General valuations (原书 pp.71-78) ============
\section{General valuations}\label{sec:10}\index{valuations, section 10}

An integral domain $R$ is a \emph{valuation ring}\index{valuation ring} if every element $x$ of its
field of fractions $K$ satisfies
\[
  x \notin R \Rightarrow x^{-1} \in R;
\]
if we write $R^{-1}$ for the set of inverses of non-zero elements of $R$ then this
condition can be expressed as $R \cup R^{-1} = K$. We also say that $R$ is a
\emph{valuation ring of $K$}. The case $R = K$ is the \emph{trivial} valuation ring.

If $R$ is a valuation ring then for any two ideals $I, J$ of $R$ either $I \subset J$
or $J \subset I$ must hold; indeed, if $x \in I$ and $x \notin J$ then for any $0 \neq y \in J$ we
have $x/y \notin R$, so that $y/x \in R$ and $y = x(y/x) \in I$, therefore $J \subset I$. Thus the
ideals of $R$ form a totally ordered set. In particular, since $R$ has only one
maximal ideal, $R$ is a local ring. We write $m$ for the maximal ideal of $R$.
Then as one sees easily,
\[
  K - R = \{x \in K^* \mid x^{-1} \in m\},
\]
where we write $K^*$ for the multiplicative group $K - \{0\}$. Thus $R$ is determined by $R$
and $m$.

If $R$ is a valuation ring of a field $K$ then any ring $R'$ with $R \subset R' \subset K$ is
obviously also a valuation ring, and in fact we have the following stronger
statement.

\begin{theorem}\label{thm:10.1}
Let $R \subset R' \subset K$ be as above, let $m$ be the maximal ideal of $R$
and $p$ that of $R'$, and suppose that $R \neq R'$. Then
\begin{enumerate}[(i)]
  \item $p \subset m \subset R \subset R'$ and $p \neq m$.
  \item $p$ is a prime ideal of $R$ and $R' = R_p$.
  \item $R/p$ is a valuation ring of the field $R'/p$.
  \item Given any valuation ring $\bar{S}$ of the field $R'/p$, let $S$ be its inverse image
        in $R'$. Then $S$ is a valuation ring having the same field of fractions $K$ as $R'$.
\end{enumerate}
\end{theorem}

\begin{proof}
(i) If $x \in p$ then $x^{-1} \notin R'$, so $x^{-1} \notin R$ and hence $x \in m$; $x$ is not a unit of $R$,
so that $x \in m$. Also, since $R \neq R'$ we have $p \neq m$.

(ii) We know that $p \subset R$, so that $p = p \cap R$, and this is a prime ideal of $R$.
Since $R - p \subset R' - p = \{\text{units of } R'\}$ we have $R_p \subset R'$, and moreover by
construction, the maximal ideal of $R_p$ is contained in the maximal ideal $p$ of
$R'$. Thus by (i), $R_p = R'$.

(iii) Write $\varphi\colon R' \to R'/p$ for the natural map; then for $x \in R' - p$, if $x \in R$
we have $\varphi(x) \in R/p$, and if $x \notin R$ we have $\varphi(x)^{-1} = \varphi(x^{-1}) \in R/p$, and
therefore $R/p$ is a valuation ring of $R'/p$.

(iv) Note that $p \subset S$ and $S/p = \bar{S}$, so that if $x \in R'$ and $x \notin S$ then $x$ is
a unit of $R'$, and $\varphi(x) \notin \bar{S}$. Thus $\varphi(x^{-1}) = \varphi(x)^{-1} \in \bar{S}$, and hence $x^{-1} \in S$.
If on the other hand $x \in K - R'$ then $x^{-1} \in p \subset S$, so that we have proved
that $S \cup S^{-1} = K$.
\end{proof}

The valuation ring $S$ in (iv) is called the \emph{composite} of $R'$ and $\bar{S}$. According
to (iii), every valuation ring of $K$ contained in $R'$ is obtained as the
composite of $R'$ and a valuation ring of $R'/p$.

Quite generally, we write $m_R$ for the maximal ideal of a local ring $R$. If $R$
and $S$ are local rings with $R \supset S$ and $m_R \cap S = m_S$ we say that $R$ \emph{dominates} $S$,\index{domination}
and write $R \ge S$. If $R \ge S$ and $R \neq S$, we write $R > S$.

\begin{theorem}\label{thm:10.2}
Let $K$ be a field, $A \subset K$ a subring, and $p$ a prime ideal of $A$.
Then there exists a valuation ring $R$ of $K$ satisfying
\[
  R \supset A \quad \text{and} \quad m_R \cap A = p.
\]
\end{theorem}

\begin{proof}
Replacing $A$ by $A_p$ we can assume that $A$ is a local ring with $p = m_A$.
Now write $\mathcal{F}$ for the set of all subrings $B$ of $K$ containing $A$ and such that
$1 \notin pB$. Now $A \in \mathcal{F}$, and if $\mathcal{F}_0 \subset \mathcal{F}$ is a subset totally ordered by inclusion
then the union of all the elements of $\mathcal{F}_0$ is again an element of $\mathcal{F}$, so that, by
Zorn's lemma, $\mathcal{F}$ has an element $R$ which is maximal for inclusion. Since
$pR \neq R$ there is a maximal ideal $m$ of $R$ containing $pR$. Then $R \subset R_m \in \mathcal{F}$, so
that $R = R_m$, and $R$ is local. Also $p \subset m$ and $p$ is a maximal ideal of $A$, so that
$m \cap A = p$. Thus it only remains to prove that $R$ is a valuation ring of $K$. If
$x \in K$ and $x \notin R$ then since $R[x] \notin \mathcal{F}$ we have $1 \in pR[x]$, and there
exists a relation of the form
\[
  1 = a_0 + a_1x + \cdots + a_nx^n \quad \text{with } a_i \in pR.
\]
Since $1 - a_0$ is a unit of $R$ we can modify this to get a relation
\[
  (*) \quad 1 = b_1x + \cdots + b_nx^n \quad \text{with } b_i \in m.
\]
Among all such relations, choose one for which $n$ is as small as possible.
Similarly, if $x^{-1} \notin R$ we can find a relation
\[
  (**) \quad 1 = c_1x^{-1} + \cdots + c_mx^{-m} \quad \text{with } c_i \in m,
\]
and choose one for which $m$ is as small as possible. If $n \ge m$ we multiply
$(**)$ by $b_nx^n$ and subtract from $(*)$, and obtain a relation of the form $(*)$
but with a strictly smaller degree $n$, which is a contradiction; if $n < m$ then
we get the same contradiction on interchanging the roles of $x$ and $x^{-1}$.
Thus if $x \notin R$ we must have $x^{-1} \in R$.
\end{proof}

\begin{theorem}\label{thm:10.3}
A valuation ring is integrally closed.
\end{theorem}

\begin{proof}
Let $R$ be a valuation ring of a field $K$, and let $x \in K$ be integral over $R$,
so that $x^n + a_1x^{n-1} + \cdots + a_n = 0$ with $a_i \in R$. If $x \notin R$ then $x^{-1} \in m_R$, but then
\[
  1 + a_1x^{-1} + \cdots + a_nx^{-n} = 0,
\]
and we get $1 \in m_R$, which is a contradiction. Hence $x \in R$.
\end{proof}

\begin{theorem}\label{thm:10.4}
Let $K$ be a field, $A \subset K$ a subring, and let $B$ be the integral
closure of $A$ in $K$. Then $B$ is the intersection of all the valuation rings of $K$
containing $A$.
\end{theorem}

\begin{proof}
Write $B'$ for the intersection of all valuation rings of $K$ containing
$A$, so that by the previous theorem we have $B' \supset B$. To prove the opposite
inclusion it is enough to show that for any element $x \in K$ which is not inte-
gral over $A$ there is a valuation ring of $K$ containing $A$ but not $x$. Set
$x^{-1} = y$. The ideal $yA[y]$ of $A[y]$ does not contain $1$: for if $1 = a_1y +
a_2y^2 + \cdots + a_ny^n$ with $a_i \in A$ then $x$ would be integral over $A$, contradicting
the assumption. Therefore there is a maximal ideal $p$ of $A[y]$ containing
$yA[y]$, and by Theorem 2 there exists a valuation ring $R$ of $K$ such that
$R \supset A[y]$ and $m_R \cap A[y] = p$. Now $y = x^{-1} \in m_R$, so that $x \notin R$.
\end{proof}

Let $K$ be a field and $A \subset K$ a subring. If a valuation ring $R$ of $K$ contains $A$
we say that $R$ has a \emph{centre} in $A$,\index{centre} and the prime ideal $m_R \cap A$ of $A$ is called the
\emph{centre of $R$ in $A$}. The set of valuation rings of $K$ having a centre in $A$ is called
the \emph{Zariski space}\index{Zariski space} or the \emph{Zariski Riemann surface} of $K$ over $A$, and written
$\Zar(K, A)$. We will treat $\Zar(K, A)$ as a topological space, introducing a
topology as follows.

For $x_1, \ldots, x_n \in K$, set $U(x_1, \ldots, x_n) = \Zar(K, A[x_1, \ldots, x_n])$. Then since
\[
  U(x_1, \ldots, x_n) \cap U(y_1, \ldots, y_m) = U(x_1, \ldots, x_n, y_1, \ldots, y_m),
\]
the collection $\mathcal{B} = \{U(x_1, \ldots, x_n) \mid n \ge 0 \text{ and } x_i \in K\}$ is the basis for
the open sets of a topology on $\Zar(K, A)$. That is, we take as open sets the
subsets of $\Zar(K, A)$ which can be written as a union of elements of $\mathcal{B}$.
As in the case of $\Spec$, this topology is called the \emph{Zariski topology}.

\begin{theorem}\label{thm:10.5}
$\Zar(K, A)$ is quasi-compact.
\end{theorem}

\begin{proof}
It is enough to prove that if $\mathcal{A}$ is a family of closed sets of $\Zar(K, A)$
having the finite intersection property (that is, the intersection of any finite
number of elements of $\mathcal{A}$ is non-empty) then the intersection of all the
elements of $\mathcal{A}$ is non-empty. By Zorn's lemma there exists a maximal family
of closed sets $\mathcal{A}'$ having the finite intersection property and containing $\mathcal{A}$.
Since it is then enough to show that the intersection of all the elements of
$\mathcal{A}'$ is non-empty, we can take $\mathcal{A} = \mathcal{A}'$. Then it is easy to see that $\mathcal{A}$ has the
following properties:
\begin{enumerate}[($\alpha$)]
  \item $F_1, \ldots, F_r \in \mathcal{A} \Rightarrow F_1 \cap \cdots \cap F_r \in \mathcal{A}$,
  \item $Z_1, \ldots, Z_r$ are closed sets and $Z_1 \cup \cdots \cup Z_r \in \mathcal{A} \Rightarrow Z_i \in \mathcal{A}$ for some $i$;
  \item if a closed set $F$ contains an element of $\mathcal{A}$ then $F \in \mathcal{A}$.
\end{enumerate}
For a subset $F \subset \Zar(K, A)$ we write $F^c$ to denote the complement of $F$.
If $F \in \mathcal{A}$ and $F^c = \bigcup_n U_n$ then $F = \bigcap_n U_n^c$, and moreover if $U(x_1, \ldots, x_n)^c =
\bigcup_{i=1}^{n} U(x_i)^c \in \mathcal{A}$ then by ($\alpha$) above one of the $U(x_i)^c$ must belong to $\mathcal{A}$.
Hence the intersection of all the elements of $\mathcal{A}$ is the same thing as the
intersection of the sets of the form $U(x)^c$ belonging to $\mathcal{A}$. Set
\[
  \Gamma = \{y \mid U(y^{-1})^c \in \mathcal{A}\}.
\]
Now since the condition for $R \in \Zar(K, A)$ to belong to $U(y^{-1})^c$ is that
$y \in m_R$, the intersection of all the elements of $\mathcal{A}$ is equal to
\[
  \{R \in \Zar(K, A) \mid m_R \supset \Gamma\}.
\]
Write $I$ for the ideal of $A[\Gamma]$ generated by $\Gamma$. If $1 \notin I$ then by Theorem 2, the
above set is non-empty, as required to prove. But if $1 \in I$ then there is a
finite subset $\{y_1, \ldots, y_n\} \subset \Gamma$ such that $1 \in \sum y_iA[y_1, \ldots, y_n]$; but then
$U(y_1^{-1})^c \cap \cdots \cap U(y_n^{-1})^c = \emptyset$, which contradicts the finite intersection
property of $\mathcal{A}$.
\end{proof}

When $K$ is an algebraic function field over an algebraically closed field $k$
of characteristic $0$ (that is, $K$ is a finitely generated extension of $k$), Zariski
gave a classification of valuation rings of $K$ containing $k$, and using this and
the compactness result above, he succeeded in giving an algebraic proof in
characteristic $0$ of the resolution of singularities of algebraic varieties of
dimension $2$ and $3$. However, Hironaka's general proof of resolution of
singularities in characteristic $0$ in all dimensions was obtained by other
methods, without the use of valuation theory.

As we saw at the beginning of this section, the ideals of $R$ form a totally
ordered set under inclusion. This holds not just for ideals, but for all $R$-
modules contained in $K$. In particular, if we set
\[
  G = \{xR \mid x \in K \text{ and } x \neq 0\},
\]
then $G$ is a totally ordered set under inclusion; we will, however, give $G$ the
opposite order to that given by inclusion. That is, we define $\le$ by
\[
  xR \le yR \iff xR \supset yR.
\]
Moreover, $G$ is an Abelian group with product $(xR) \cdot (yR) = xyR$. In general,
an Abelian group $H$ written additively, together with a total order relation
$\le$ is called an \emph{ordered group}\index{ordered group} if the axiom
\[
  x \ge y, \; z \ge t \Rightarrow x + z \ge y + t
\]
holds. This axiom implies
\[
  \text{(1) } x > 0, \; y > 0 \Rightarrow x + y > 0, \qquad \text{and} \qquad \text{(2) } x > y \Rightarrow -y > -x.
\]
We make an ordered set $H \cup \{\infty\}$ by adding to $H$ an element $\infty$ bigger
than all the elements of $H$, and fix the conventions $\infty + \alpha = \infty$ for $\alpha \in H$
and $\alpha + \infty = \infty$. A map $v\colon K \to H \cup \{\infty\}$ from a field $K$ to $H \cup \{\infty\}$
is called an \emph{additive valuation}\index{additive valuation} or just a \emph{valuation} of $K$ if it satisfies the
conditions
\begin{enumerate}[(1)]
  \item $v(xy) = v(x) + v(y)$;
  \item $v(x + y) \ge \min\{v(x), v(y)\}$;
  \item $v(x) = \infty \iff x = 0$.
\end{enumerate}
If we write $K^*$ for the multiplicative group of $K$ then $v$ defines a
homomorphism $K^* \to H$; the image is a subgroup of $H$, called the \emph{value group} of $v$. We also set
\[
  R_v = \{x \in K \mid v(x) \ge 0\} \quad \text{and} \quad m_v = \{x \in K \mid v(x) > 0\},
\]
obtaining a valuation ring $R_v$ of $K$ with $m_v$ as its maximal ideal, and call $R_v$ the
\emph{valuation ring of $v$}, and $m_v$ the \emph{valuation ideal} of $v$. Conversely, if $R$ is a
valuation ring of $K$, then the group $G = \{xR \mid x \in K^*\}$ described above is an
ordered group, and we obtain an additive valuation of $K$ with value group
$G$ by defining $v\colon K \to G \cup \{\infty\}$ by $v(0) = \infty$ and $v(x) = xR$ for $x \in K^*$ (there
is no real significance in whether or not we rewrite the multiplication in
$G$ additively); the valuation ring of $v$ is $R$. The additive valuation corres-
ponding to a valuation ring $R$ is not quite unique, but if $v$ and $v'$ are
additive valuations of $K$ with value groups $H$ and $H'$ and both having
valuation ring $R$ then there exists an order-isomorphism
$\varphi\colon H \to H'$ such that $v' = \varphi v$ (prove this!). Thus we can think of valuation rings and
additive valuations as being two aspects of the same thing.

We now give some examples of ordered groups:
\begin{enumerate}[(1)]
  \item the additive group of real numbers $\R$ (this is isomorphic to the
        multiplicative group of positive reals), or any subgroup of this;
  \item the group $\Z$ of rational integers;
  \item the direct product $\Z^n$ of $n$ copies of $\Z$, with lexicographical order, that is
        \[
          (a_1, \ldots, a_n) < (b_1, \ldots, b_n) \iff \text{the first non-zero element of}
        \]
        \[
          b_1 - a_1, \ldots, b_n - a_n \text{ is positive}
        \]
\end{enumerate}

An ordered group $G$ is said to be \emph{Archimedean}\index{Archimedean ordered group} if it is order-isomorphic to
a suitable subgroup of $\R$. The name is explained by the following theorem:
the condition in it is known as the Archimedean axiom. (Note that our
usage is completely unrelated to that in number theory, where non-
Archimedean fields are p-adic fields, as opposed to subfields of $\R$ and $\C$ with
the usual `Archimedean' metrics.)

\begin{theorem}\label{thm:10.6}
Let $G$ be an ordered group; then $G$ is Archimedean if and
only if the following condition holds:
\[
  \text{if } a, b \in G \text{ with } a > 0, \text{ there exists a natural number } n \text{ such that } na > b.
\]
\end{theorem}

\begin{proof}
The condition is obviously necessary, and we prove sufficiency. If
$G = \{0\}$ then we can certainly embed $G$ in $\R$. Suppose that $G \neq \{0\}$.
Fix some $0 < \alpha \in G$. For any $\beta \in G$, there is a well-defined largest integer
$n$ such that $n\alpha \le \beta$ (if $\beta \ge 0$ this is clear by assumption; if $\beta < 0$, let $m$ be
the smallest integer such that $-\beta \le m\alpha$, and set $n = -m$). Let this be $n_0$.
Now set $\beta_1 = \beta - n_0\alpha$ and let $n_1$ be the largest integer $n$ such that $n\alpha \le 10\beta_1$;
we have $0 \le n_1 < 10$. Set $\beta_2 = 10\beta_1 - n_1\alpha$ and let $n_2$ be the largest integer
$n$ such that $n\alpha \le 10\beta_2$. Continuing in the same way, we find integers $n_0$,
$n_1, n_2, \ldots$, and set $\varphi(\beta) = \alpha_0$, where $\alpha_0$ is the real number given by the decimal
expression $n_0 + 0.n_1n_2n_3\cdots$. Then it can easily be checked that $\varphi\colon G \to \R$
is order-preserving, in the sense that $\beta < \beta'$ implies $\varphi(\beta) < \varphi(\beta')$.

We also see that $\varphi$ is injective. For this, we only need to observe that
if $\beta < \beta'$ then there exists a natural number $r$ such that $\alpha < 10^r(\beta' - \beta)$;
the details are left to the reader.

Finally we show that $\varphi$ is a group homomorphism. For $\beta \in G$, we write
$n/10^r$ with $n \in \Z$ to denote the number obtained by taking the first $r$
decimal places of $\varphi(\beta)$; the numerator $n$ is determined by the property that
$n\alpha \le 10^r\beta < (n + 1)\alpha$. For $\beta' \in G$, if $n'\alpha \le 10^r\beta' < (n' + 1)\alpha$ then we have
\[
  (n + n')\alpha \le 10^r(\beta + \beta') < (n + n' + 2)\alpha,
\]
and hence
\[
  \varphi(\beta + \beta') - (n + n') \cdot 10^{-r} < 2 \cdot 10^{-r},
\]
so that
\[
  |\varphi(\beta + \beta') - \varphi(\beta) - \varphi(\beta')| \le 4 \cdot 10^{-r},
\]
and since $r$ is arbitrary, $\varphi(\beta + \beta') = \varphi(\beta) + \varphi(\beta')$.
\end{proof}

A non-zero group $G$ order-isomorphic to a subgroup of $\R$ is said to have
\emph{rank 1}.\index{rank (of an ordered group)} The \emph{rational rank} of an ordered group $G$ of rank $1$ is the maximum
number of elements of $G$ (viewed as a subgroup of $\R$) which are linearly
independent over $\Z$. For example, the additive group $G = \Z + \Z\sqrt{2} \subset \R$ is
an ordered group of rank $1$ and rational rank $2$.

\begin{theorem}\label{thm:10.7}
Let $R$ be a valuation ring having value group $G$. Then $G$ has
rank $1 \iff R$ has Krull dimension $1$.
\end{theorem}

\begin{proof}
($\Rightarrow$) Since $G \neq 0$, $R$ is not a field. Suppose that $p$ is a prime ideal of $R$
distinct from $m_R$. Let $\xi \in m_R$ be such that $\xi \notin p$, and set $v(\xi) = x$, where $v$ is the
additive valuation corresponding to $R$. Suppose that $0 \neq \eta \in p$, and set $\psi =
v(\eta)$; then $\psi \in G$ and $x > 0$, so that $nx > \psi$ for some sufficiently large natural
number $n$. This means that $\xi^n/\eta \in R$, so that $\xi^n \in \eta R \subset p$; then since $p$ is prime
we have $\xi \in p$, which is a contradiction. Therefore $p = (0)$. The only prime
ideals of $R$ are $m_R$ and $(0)$, which means $\dim R = 1$.

($\Leftarrow$) If $0 \neq \eta \in m_R$ then $m_R$ is the unique prime ideal containing $\eta R$, and
hence $\sqrt{\eta R} = m_R$. Thus for any $\xi \in m_R$ there exists a natural number $n$ such
that $\xi^n \in \eta R$. From this one sees easily that $G$ satisfies the Archimedean
axiom.
\end{proof}

% ============ 习题 ============
\section*{Exercises to \S\ref{sec:10}}
\emph{Prove the following propositions.}

\begin{exercise}\label{ex:10.1-ex}
In a valuation ring any finitely generated ideal is principal.
\end{exercise}

\begin{exercise}\label{ex:10.2-ex}
If $R$ is a valuation ring then an $R$-module $M$ is flat if and only if it is torsion-
free (that is, $a \neq 0, x \neq 0 \Rightarrow ax \neq 0$ for $a \in R, x \in M$).
\end{exercise}

\begin{exercise}\label{ex:10.3-ex}
In Theorem~\ref{thm:10.4}, if $A$ is a local ring then $B$ is the intersection of the
valuation rings of $K$ dominating $A$.
\end{exercise}

\begin{exercise}\label{ex:10.4-ex}
If $R$ is a valuation ring of Krull dimension $\ge 2$ then the formal power series
ring $R[[X]]$ is not integrally closed ([B5], Ex. 27, p.~76 and Seidenberg~\cite{Sd1}).
\end{exercise}

\begin{exercise}\label{ex:10.5-ex}
If $R$ is a valuation ring of Krull dimension $1$ and $K$ its field of fractions then
there do not exist any rings intermediate between $R$ and $K$. In other words
$R$ is maximal among proper subrings of $K$. Conversely if a ring $R$, not a
field, is a maximal proper subring of a field $K$ then $R$ is a valuation ring of
Krull dimension $1$.
\end{exercise}

\begin{exercise}\label{ex:10.6-ex}
If $v$ is an additive valuation of a field $K$, and if $\alpha, \beta \in K$ are such that
$v(\alpha) \neq v(\beta)$ then $v(\alpha + \beta) = \min\{v(\alpha), v(\beta)\}$.
\end{exercise}

\begin{exercise}\label{ex:10.7-ex}
If $v$ is an additive valuation of a field $K$ and $\alpha_1, \ldots, \alpha_n \in K$ are such that
$\alpha_1 + \cdots + \alpha_n = 0$ then there exist two indices $i, j$ such that $i \neq j$ and
$v(\alpha_i) = v(\alpha_j)$.
\end{exercise}

\begin{exercise}\label{ex:10.8-ex}
Let $K \subset L$ be algebraic field extension of degree $[L:K] = n$, and let $S$ be
a valuation ring of $L$; set $R = S \cap K$. Write $k, k'$ for the residue fields of $S$
and $R$, and set $[k:k'] = f$. Now let $G$ be the value group of $S$, and let $G'$ be
the image of $K^*$ under the valuation map $L^* \to G$; set $[G:G'] = e$. Then
$ef \le n$. (The numbers $f$ and $e$ are called the \emph{degree} and the \emph{ramification}
\emph{index} of the valuation ring extension $S/R$.)
\end{exercise}

\begin{exercise}\label{ex:10.9-ex}
Let $L, K, S$ and $R$ be as in the previous question, and let $S_1 \neq S$ be a
valuation ring of $L$ such that $S_1 \cap K = R$. Then neither of $S$ or $S_1$ contains
the other.
\end{exercise}

\begin{exercise}\label{ex:10.10-ex}
Let $A$ be an integral domain with field of fractions $K$, and let $H$ be an
ordered group. If a map $v\colon A \to H \cup \{\infty\}$ satisfies conditions (1), (2) and (3)
of an additive valuation (on elements of $A$), then $v$ can be extended
uniquely to an additive valuation $K \to H \cup \{\infty\}$.
\end{exercise}

\begin{exercise}\label{ex:10.11-ex}
Let $k$ be a field, $X$ and $Y$ indeterminates, and suppose that $\alpha$ is a positive
irrational number. Then the map $v\colon k[X, Y] \to \R \cup \{\infty\}$ defined by
taking $\sum c_{nm}X^nY^m$ (with $c_{nm} \in k$) into $v(\sum c_{nm}X^nY^m) = \min\{n +
m\alpha \mid c_{nm} \neq 0\}$ determines a valuation of $k(X, Y)$ with value group
$\Z + \Z\alpha$.
\end{exercise}

% ============ §11 DVRs and Dedekind rings (原书 pp.78-86) ============
\section{DVRs and Dedekind rings}\label{sec:11}\index{discrete valuation rings, section 11}\index{Dedekind rings, section 11}

A valuation ring whose value group is isomorphic to $\Z$ is called
a \emph{discrete valuation ring} (DVR).\index{discrete valuation ring} Discrete refers to the fact that the value
group is a discrete subgroup of $\R$, and has nothing to do with the
$m$-adic topology of the local ring being discrete.

\begin{theorem}\label{thm:11.1}
Let $R$ be a valuation ring. Then the following conditions are
equivalent.
\begin{enumerate}[(1)]
  \item $R$ is a DVR;
  \item $R$ is a principal ideal domain;
  \item $R$ is Noetherian.
\end{enumerate}
\end{theorem}

\begin{proof}
Let $K$ be the field of fractions of $R$ and $m$ its maximal ideal.

$(1) \Rightarrow (2)$ Let $v_R$ the additive valuation of $R$ having value group $\Z$; this is
called the \emph{normalised additive valuation} corresponding to $R$. There exists
$t \in R$ such that $v_R(t) = 1$. For $0 \neq x \in m$, the valuation $v_R(x)$ is a positive
integer, say $v_R(x) = n$; then $v_R(x/t^n) = 0$, so that we can write $x = t^nu$ with $u$ a
unit of $R$. In particular $m = tR$. Let $I \neq (0)$ be any ideal of $R$; then
$\{v_R(a) \mid 0 \neq a \in I\}$ is a set of non-negative integers, and so has a smallest
element, say $n$. If $n = 0$ then $I$ contains a unit of $R$, so that $I = R$. If $n > 0$ then
there exists an $x \in I$ such that $v_R(x) = n$; then $I = xR = t^nR$. Therefore $R$ is a
principal ideal domain, and moreover every non-zero ideal of $R$ is a power
of $m = tR$.

$(2) \Rightarrow (3)$ is obvious.

$(3) \Rightarrow (2)$ In general, given any two ideals of a valuation ring, one contains
the other, so that any finitely generated ideal $a_1R + \cdots + a_nR$ is equal to one
of the $a_iR$, and therefore principal. Hence, if $R$ is Noetherian every ideal of $R$
is principal.

$(2) \Rightarrow (1)$ We can write $m = xR$ for some $x$. Now if we set $I = \bigcap_{\nu \ge 1} x^\nu R$
then this is also a principal ideal, so that we can write $I = yR$. If we set
$y = xz$, then from $y \in x^\nu R$ we get $z \in x^{\nu-1}R$, and since this holds for every
$\nu$, we have $z \in I$, hence we can write $z = yu$. Since $y = xz = xyu$, we have $y(1
- xu) = 0$, but then since $x \in m$ we must have $y = 0$, and therefore $I = (0)$.
Because of this, for every non-zero element $a \in R$, there is a well-defined
integer $\nu \ge 0$ such that $a \in x^\nu R$ but $a \notin x^{\nu+1}R$; we then set $v(a) = \nu$. It is not
difficult to see that if $a, b, c, d \in R - \{0\}$ satisfy $a/b = c/d$ then
\[
  v(a) - v(b) = v(c) - v(d);
\]
therefore setting $v(\xi) = v(a) - v(b)$ for $\xi = a/b \in K^*$ gives a map $v\colon K^* \to \Z$
which can easily be seen to be an additive valuation of $K$ whose valuation
ring is $R$. The value group of $v$ is clearly $\Z$, so that $R$ is a DVR.
\end{proof}

If $R$ is a DVR with maximal ideal $m$ then an element $t \in R$ such that
$m = tR$ is called a \emph{uniformising element} of $R$.\index{uniformising element}

\begin{remark}
A valuation ring $S$ whose maximal ideal $m_S$ is principal does not
have to be a DVR. To obtain a counter-example, let $K$ be a field, and $R$
a DVR of $K$; set $k = R/m_R$, and suppose that $\bar{R}$ is a DVR of $k$. Now let
$S$ be the composite of $R$ and $\bar{R}$. Let $f$ be a uniformising element of $R$, and
$g \in S$ be any element mapping to a uniformising element $\bar{g}$ of $\bar{R}$. Then
$\eta = fR \subset m_R \subset S \subset R$, and $m_S/m_R = \bar{g}\bar{R} = \sqrt{\bar{g}S/\bar{m}_R}$, and so
\[
  m_S = m_R + gS.
\]
On the other hand $g^{-1} \in R$, so that for any $h \in m_R$ we have $h/g \in m_R \subset S$, and
hence $m_R \subset gS$, so that
\[
  m_S = gS.
\]
However, $m_R = fR$ is not finitely generated as an ideal of $S$, being generated
by $f, fg^{-1}, fg^{-2}, \ldots$. The value group of $S$ is $\Z^2$, with the valuation
$v\colon K^* \to \Z^2$ given by
\[
  v(x) = (n, m), \quad \text{where } n = v_R(x) \text{ and } m = v_{\bar{R}}(\bar{\varphi}(xf^{-n})),
\]
where $\varphi\colon R \to R/m_R = k$ is the natural homomorphism.
\end{remark}

The previous theorem gives a characterisation of DVRs among valuation
rings; now we consider characterisations among all rings.

\begin{theorem}\label{thm:11.2}
Let $R$ be a ring; then the following conditions are
equivalent:
\begin{enumerate}[(1)]
  \item $R$ is a DVR;
  \item $R$ is a local principal ideal domain, and not a field;
  \item $R$ is a Noetherian local ring, $\dim R > 0$ and the maximal ideal $m_R$ is
        principal;
  \item $R$ is a one-dimensional normal Noetherian local ring.
\end{enumerate}
\end{theorem}

\begin{proof}
We saw $(1) \Rightarrow (2)$ in the previous theorem; $(2) \Rightarrow (3)$ is obvious.

$(3) \Rightarrow (1)$ Let $xR$ be the maximal ideal of $R$. If $x$ were nilpotent then we
would have $\dim R = 0$, and hence $x^\nu \neq 0$ for all $\nu$. By the Krull intersection
theorem (Theorem~\ref{thm:8.10}, (i)) we have $\bigcap_{\nu \ge 1} x^\nu R = (0)$, so that for $0 \neq y \in R$
there is a well-determined $\nu$ such that $y \in x^\nu R$ and $y \notin x^{\nu+1}R$. If $y = x^\nu u$, then
since $u \notin xR$ it must be a unit. Similarly, for $0 \neq z \in R$ we have $z = x^\mu v$, with $v$ a
unit. Therefore $yz = x^{\mu+\nu}uv \neq 0$, and so $R$ is an integral domain. Finally,
any element $\xi$ of the fraction field of $R$ can be written $\xi = x^\nu u$, with $u$ a unit of
$R$ and $\nu \in \Z$, and it is easy to see that setting $v(\xi) = \nu$ defines an additive
valuation of the field of fractions of $R$ whose valuation ring is $R$.

$(1) \Rightarrow (4)$ In a DVR the only ideals are $(0)$ and the powers of the maximal
ideal, so that the only prime ideals of $R$ are $(0)$ and $m_R$, and hence
$\dim R = 1$. By the previous theorem $R$ is Noetherian, and it is normal
because it is a valuation ring.

$(4) \Rightarrow (3)$ By assumption $R$ is an integral domain. Write $K$ for the field of
fractions and $m$ for the maximal ideal of $R$. Then $m \neq 0$, so that by Theorem~\ref{thm:8.10}, $m \neq m^2$; choose some $x \in m - m^2$. Since $\dim R = 1$ the only prime
ideals of $R$ are $(0)$ and $m$, so that $m$ must be a prime divisor of $xR$, and there
exists $y \in R$ such that $xR:y = m$. Set $a = yx^{-1}$; then $a \notin R$, but $am \subset R$. Now
we set $m^{-1} = \{\xi \in K \mid \xi m \subset R\}$, so that $R \subset m^{-1}$, and $R \neq m^{-1}$ since $a \in m^{-1}$.
Consider the ideal $m^{-1}m$ of $R$; since $R \subset m^{-1}$ we have $m \subset m^{-1}m$. If we
had $m = m^{-1}m$ then we would get $am \subset m$, and then $a$ would be integral
over $R$ by Theorem~\ref{thm:2.1}, so that $a \in R$, which is a contradiction. Hence we
must have $m^{-1}m = R$. Moreover, $xm^{-1} \subset R$ is an ideal, and if $xm^{-1} \subset m$
then we would have $xR = xm^{-1}m \subset m^2$, contradicting $x \notin m^2$. Therefore
$xm^{-1} = R$, and hence $xR = xm^{-1}m = m$, so that $m$ is principal.
\end{proof}

Quite generally, if $R$ is an integral domain and $K$ its fields of fractions, we
say that an $R$-submodule $I$ of $K$ is a \emph{fractional ideal}\index{fractional ideal} of $R$ if $I \neq 0$, and
there exists a non-zero element $\alpha \in R$ such that $\alpha I \subset R$ (see Ex.~\ref{ex:3.4-ex}). As
an $R$-module we have $I \cong \alpha I$, so that if $R$ is a Noetherian integral domain
then any fractional ideal is finitely generated. For $I$ a fractional ideal we
set $I^{-1} = \{\alpha \in K \mid \alpha I \subset R\}$; we say that $I$ is \emph{invertible}\index{invertible ideal} if $I^{-1}I = R$.

\begin{theorem}\label{thm:11.3}
Let $R$ be an integral domain and $I$ a fractional ideal of $R$.
Then the following conditions are equivalent:
\begin{enumerate}[(1)]
  \item $I$ is invertible;
  \item $I$ is a projective $R$-module;
  \item $I$ is finitely generated, and for every maximal ideal $P$ of $R$, the
        fractional ideal $I_P = IR_P$ of $R_P$ is principal.
\end{enumerate}
\end{theorem}

\begin{proof}
$(1) \Rightarrow (2)$ If $I^{-1}I = R$ then there exist $a_i \in I$ and $b_i \in I^{-1}$ such that
$\sum_i a_ib_i = 1$. Then $a_1, \ldots, a_n$ generate $I$, since for any $x \in I$ we have $\sum (xb_i)a_i =
x$, and $xb_i \in R$. Let $F = Re_1 + \cdots + Re_n$ be the free $R$-module with basis
$e_1, \ldots, e_n$; we define the $R$-linear map $\varphi\colon F \to I$ by $\varphi(e_i) = a_i$, so that $\varphi$ is
surjective. Then we define $\psi\colon I \to F$ by writing $\psi_i\colon I \to R$ for the map
$\psi_i(x) = b_ix$, and setting $\psi(x) = \sum \psi_i(x)e_i$. We then have $\varphi\psi(x) = x$, so that $\varphi$
splits, and $I$ is isomorphic to a direct summand of the free module $F$, and
therefore projective.

$(2) \Rightarrow (1)$ Every $R$-linear map from $I$ to $R$ is given by multiplication by
some element of $K$ (prove this!). If we let $\varphi\colon F \to I$ be a surjective map from
a free module $F = \bigoplus Re_i$, by assumption there exists a splitting $\psi\colon I \to F$
such that $\varphi\psi = 1$. Write $\psi(x) = \sum \lambda_i(x)e_i$ for $x \in I$; then by what we have said,
each $\lambda_i$ determines a $b_i \in K$ such that $\lambda_i(x) = b_ix$, and since for each $x$ there
are only finitely many $i$ such that $\lambda_i(x) \neq 0$, we have $b_i = 0$ for all but finitely
many $i$. Letting $b_1, \ldots, b_n$ be the non-zero ones, we have $\sum a_ib_ix = x$ for all
$x \in I$, where $a_i = \varphi(e_i)$. Thus $\sum_i a_ib_i = 1$. Moreover, since $\lambda_iI = \lambda_i \subset R$ we
have $b_i \in I^{-1}$, and therefore $I^{-1}I = R$.

$(1) \Rightarrow (3)$ As we have already seen, $I$ is finitely generated. Now if $\sum a_ib_i = 1$
and $P$ is any prime ideal then at least one of $a_ib_i$ must be a unit of $R_P$,
and $I_P = a_iR_P$. Hence $I_P$ is a principal fractional ideal.

$(3) \Rightarrow (1)$ If $I$ is finitely generated then $(I^{-1})_P = (I_P)^{-1}$. Indeed, the
inclusion $\subset$ holds for any ideal; for $\supset$, if $I = a_1R + \cdots + a_nR$ and $x \in (I_P)^{-1}$
then $xa_i \in R_P$, so there exist $c_i \in R - P$ such that $xa_ic_i \in R$, so that setting
$c = c_1 \cdots c_n$ we have $(cx)a_i \in R$ for all $i$, which gives $cx \in I^{-1}$ and $x \in (I^{-1})_P$.
From the fact that $I_P$ is principal, we get $I_P \cdot (I_P)^{-1} = R_P$. Now if $II^{-1} \neq R$
then we can take a maximal ideal $P$ such that $II^{-1} \subset P$, and then
$I_P \cdot (I_P)^{-1} = I_P(I^{-1})_P \subset PR_P$, which is a contradiction. Thus we must have
$II^{-1} = R$.
\end{proof}

\begin{theorem}\label{thm:11.4}
Let $R$ be a Noetherian integral domain, and $P$ a non-zero
prime ideal of $R$. If $P$ is invertible then $\htop P = 1$ and $R_P$ is a DVR.
\end{theorem}

\begin{proof}
If $P$ is invertible the maximal ideal $PR_P$ of $R_P$ is principal, and
condition (3) of Theorem 2 is satisfied; thus $R_P$ is a DVR, and so $\dim R_P = 1$.
\end{proof}

\begin{theorem}\label{thm:11.5}
Let $R$ be a normal Noetherian domain. Then we have
\begin{enumerate}[(i)]
  \item all the prime divisors of a non-zero principal ideal have height $1$;
  \item $R = \bigcap_{\htop p = 1} R_p$.
\end{enumerate}
\end{theorem}

\begin{proof}
(i) Suppose $0 \neq a \in R$ and that $P$ is one of the prime divisors of $aR$;
then there exists an element $b \in R$ such that $aR:b = P$. We set $PR_P = m$, and
then $aR_P:b = m$, so that $ba^{-1} \notin R_P$.
If $ba^{-1}m \subset m$ then by the determinant trick $ba^{-1}$ is integral over $R_P$, which
contradicts the fact that $R_P$ is integrally closed. Thus $ba^{-1}m = R_P$, so that $m^{-1}m = R_P$, and
then by the previous theorem we get $\htop m = \htop P = 1$.

(ii) It is sufficient to prove that for $a, b \in R$ with $a \neq 0$, $b \in aR_P$ for every
height $1$ prime $P \in \Spec R$ implies $b \in aR$. Let $P_1, \ldots, P_r$ be the prime
divisors of $aR$, and let $aR = q_1 \cap \cdots \cap q_r$ be a primary decomposition of
$aR$, where $q_i$ is a $P_i$-primary ideal for each $i$. Then since $\htop P_i = 1$, we have
$b \in aR_{P_i} \cap R = q_i$ for $i = 1, \ldots, n$, and therefore $b \in \bigcap q_i = aR$.
\end{proof}

\begin{corollary*}
Let $R$ be a Noetherian domain. The following two conditions
are necessary and sufficient for $R$ to be normal:
\begin{enumerate}[(a)]
  \item for $P$ a height $1$ prime ideal, $R_P$ is a DVR;
  \item all the prime divisors of a non-zero principal ideal of $R$ have height $1$.
\end{enumerate}
\end{corollary*}

\begin{proof}
We have already seen necessity. For sufficiency, note that the proof
of (ii) above shows that (b) implies $R = \bigcap_{\htop p = 1} R_p$. Then by (a) each $R_p$
is normal, so that $R$ is normal.
\end{proof}

\begin{definition}
An integral domain for which every non-zero ideal is invertible
is called a \emph{Dedekind ring}\index{Dedekind ring} (sometimes \emph{Dedekind domain}).
\end{definition}

\begin{theorem}\label{thm:11.6}
For an integral domain $R$ the following conditions are
equivalent:
\begin{enumerate}[(1)]
  \item $R$ is a Dedekind ring;
  \item $R$ is either a field or a one-dimensional Noetherian normal domain;
  \item every non-zero ideal of $R$ can be written as a product of a finite
        number of prime ideals.
\end{enumerate}
Moreover, the factorisation into primes in (3) is unique.
\end{theorem}

\begin{proof}
$(1) \Rightarrow (2)$ Every non-zero ideal is invertible, and therefore finitely
generated, so that $R$ is Noetherian. Let $P$ be a non-zero prime ideal of $R$;
then by Theorem 4, the local ring $R_P$ is a DVR and $\htop P = 1$, and therefore
either $R$ is a field or $\dim R = 1$. Also, by Theorem~\ref{thm:4.7} we know that $R$ is the
intersection of the $R_P$ as $P$ runs through all the maximal ideals of $R$, but
since each $R_P$ is a DVR it follows that $R$ is normal.

$(2) \Rightarrow (1)$ If $R$ is a field there is no problem. If $R$ is not a field then for every
maximal ideal $P$ of $R$ the local ring $R_P$ is a one-dimensional Noetherian
local ring and is normal, so that by Theorem 2 it is a principal ideal ring.
Thus by Theorem 3, $R$ is a Dedekind ring.

$(1) \Rightarrow (3)$ Let $I$ be a non-zero ideal. If $I = R$ then we can view it as the
product of zero ideals; if $I$ is itself maximal then it is the product of just one
prime ideal. We have already seen that $R$ is Noetherian, so that we can use
descending induction on $l(I)$, that is assume that $I \neq R$ and that every ideal
strictly bigger than $I$ is a product of prime ideals. If $I \neq R$ then there is a
maximal ideal $P$ containing $I$, and $I \subset IP^{-1} \subset R$. If $IP^{-1} = I$ then using
$P^{-1}P = R$ we would have $I = IP$, and by NAK this would lead to a
contradiction. Hence $IP^{-1} \neq I$, so that by induction we can write $IP^{-1} =
Q_1 \cdots Q_s$, with $Q_i \in \Spec R$. Multiplying both sides by $P$ gives $I = Q_1 \cdots Q_sP$.

The proof of $(3) \Rightarrow (1)$ is a little harder, and we break it up into four steps.

\textbf{Step 1.} In general, any non-zero principal ideal $aR$ of an integral domain $R$
is obviously invertible. Moreover, suppose that $I$ and $J$ are non-zero
fractional ideals and $B = IJ$; then obviously $I$ and $J$ invertible implies $B$
invertible, but the converse also holds. To see this, from $I^{-1}J^{-1}B \subset R$ we
get $I^{-1}J^{-1} \subset B^{-1}$, and also from $B^{-1}IJ \subset R$ we get $B^{-1}I \subset J^{-1}$ and
$B^{-1}J \subset I^{-1}$; now if $B$ is invertible then multiplying the last two inclusions
together we get $B^{-1} = B^{-1}B^{-1}IJ \subset I^{-1}J^{-1}$, and hence $B^{-1} = I^{-1}J^{-1}$.
Therefore
\[
  R = BB^{-1} = IJI^{-1}J^{-1} = (II^{-1})(JJ^{-1}),
\]
and we must have $II^{-1} = JJ^{-1} = R$.

\textbf{Step 2.} Let $P$ be a non-zero prime ideal. Let us prove that if $I$ is an ideal
strictly bigger than $P$ then $IP = P$. For this it is sufficient to show that if
$I = P + aR$ with $a \notin P$ then $P \subset IP$. Consider expressions of $P^2$ and $a^2R + P$
as product of prime ideals, $P^2 = P_1 \cdots P_r$ and $a^2R + P = Q_1 \cdots Q_s$. Then $P_i$
and $Q_j$ are prime ideals containing $P^2$, and so are prime ideals strictly bigger than $P$,
We now set $\bar{R} = R/P$, and write $-$ to denote the image in $\bar{R}$ of
elements or ideals of $R$. Then we have
\[
  (*) \quad \bar{P}_1 \cdots \bar{P}_r = \bar{a}^2\bar{R} = \bar{Q}_1 \cdots \bar{Q}_s,
\]
and applying Step 1 to the domain $\bar{R}$ we find that $\bar{P}_i$ and $\bar{Q}_j$ are all invertible,
and are prime ideals of $\bar{R}$. We can suppose that $\bar{P}_1$ is a minimal element of
the set $\{\bar{P}_1, \ldots, \bar{P}_r\}$. Moreover, at least one of $\bar{Q}_1, \ldots, \bar{Q}_s$ is contained
in $\bar{P}_1$, so that we can assume that $\bar{Q}_1 \subset \bar{P}_1$, and, on the other hand, since
$\bar{Q}_1$ is also a prime ideal and $\bar{P}_1 \cdots \bar{P}_r \subset \bar{Q}_1$ we must have $\bar{Q}_1 \supset \bar{P}_i$
for some $i$. Then $\bar{P}_i \subset \bar{Q}_1 \subset \bar{P}_1$, and by the minimality of $\bar{P}_1$ we have $\bar{P}_i =
\bar{P}_1 = \bar{Q}_1$. Multiplying through both sides of $(*)$ by $\bar{P}_1^{-1}$ gives
\[
  \bar{P}_2 \cdots \bar{P}_r = \bar{Q}_2 \cdots \bar{Q}_s;
\]
proceeding in the same way, we see that $r = s$, and that after reordering the
$\bar{Q}_j$ we can assume that $\bar{P}_i = \bar{Q}_i$ for $i = 1, \ldots, r$. From this we get $P_i = Q_i$, and
$a^2R + P = (P + aR)^2 = P^2 + aP + a^2R$. Thus any element $x \in P$ can be
written
\[
  x = y + az + a^2t \quad \text{with } y \in P^2, z \in P \text{ and } t \in R.
\]
Since $a \notin P$ we must have $t \in P$, and then as required we have
$P \subset P^2 + aP = (P + aR)P$.

\textbf{Step 3.} Let $b \in R$ be a non-zero element; then in the factorisation
$bR = P_1 \cdots P_r$, every $P_i$ is a maximal ideal of $R$. Indeed, if $I$ is any ideal
strictly greater than $P_i$ then $IP_i = P_i$, and by Step 1 $P_i$ is invertible, so that
$I = R$.

\textbf{Step 4.} Let $P$ be a prime ideal of $R$, and $0 \neq a \in P$. If $aR = P_1 \cdots P_r$ with
$P_i \in \Spec R$ then $P$ must contain one of the $P_i$, but from Step 3 we know that
$P_i$ is maximal, so that $P = P_i$. We deduce that $P$ is a maximal ideal and is
invertible. If every non-zero prime ideal is invertible then any non-zero ideal
of $R$ is invertible, since it can be written as a product of primes. This
completes the proof of $(3) \Rightarrow (1)$.

Finally, if (1), (2) and (3) hold, then as we have seen in Step 2 above, the
uniqueness of factorisation into primes is a consequence of the fact that
the prime ideals of $R$ are invertible.
\end{proof}

\begin{theorem}[the Krull--Akizuki theorem]\label{thm:11.7}
Let $A$ be a one-dimensional
Noetherian integral domain with field of fractions $K$, let $L$ be a finite
algebraic extension field of $K$, and $B$ a ring with $A \subset B \subset L$; then $B$ is a
Noetherian ring of dimension at most $1$, and if $J$ is a non-zero ideal of $B$ then
$B/J$ is an $A$-module of finite length.
\end{theorem}

\begin{proof}
We follow the method of proof of Akizuki [1] in the linear algebra
formulation of [BS]. First of all we prove the following lemma.

\begin{lemma}\label{lem:11.1}
Let $A$ and $K$ be as in the theorem, and let $M$ be a torsion-free
$A$-module (see Ex.~\ref{ex:10.2-ex}) of rank $r < \infty$. Then for $0 \neq a \in A$ we have
\[
  l(M/aM) \le r \cdot l(A/aA).
\]
\end{lemma}

\begin{remark}
The \emph{rank} of a module $M$ over an integral domain $A$ is the maximal
number of elements of $M$ linearly independent over $A$; this is equal to the
dimension of the $K$-vector space $M \otimes_A K$.
\end{remark}

\begin{proof}[Proof of the lemma]
First we assume that $M$ is finitely generated. Choose
elements $\xi_1, \ldots, \xi_r \in M$ linearly independent over $A$ and set $E = \sum A\xi_i$; then
for any $\eta \in M$ there exists $t \in A$ with $t \neq 0$ such that $t\eta \in E$. If we set $C = M/E$
then from the assumption on $M$ we see that $C$ is also finitely generated, so
that $tC = 0$ for suitable $0 \neq t \in A$. Applying Theorem~\ref{thm:6.4} to $C$, we can find
$C = C_0 \supset C_1 \supset \cdots \supset C_m = 0$, such that $C_{i-1}/C_i \cong A/p_i$, with $p_i \in \Spec A$.
Now $t \in p_i$, and since $A$ is one-dimensional each $p_i$ is maximal, so that
$l(C) = m < \infty$. If $0 \neq a \in A$ then the exact sequence
\[
  E/a^nE \to M/a^nM \to C/a^nC \to 0
\]
gives
\[
  (*) \quad l(M/a^nM) \le l(E/a^nE) + l(C) \quad \text{for all } n > 0.
\]
Now $E$ and $M$ are both torsion-free $A$-modules, and one sees easily that
$aM/a^{i+1}M \cong M/aM$, and similarly for $E$. Hence $(*)$ can be written
$n \cdot l(M/aM) \le n \cdot l(E/aE) + l(C)$ for all $n > 0$, which gives $l(M/aM) \le l(E/aE)$.
Since $E \cong A^r$ we have $l(E/aE) = r \cdot l(A/aA)$. This completes the proof in the
case that $M$ is finitely generated. If $M$ is not finitely generated, take a finitely
generated submodule $\bar{M}_0 = A\bar{\xi}_1 + \cdots + A\bar{\xi}_r$ of $\bar{M} = M/aM$. Then choos-
ing an inverse image $\omega_i$ in $M$ for each $\bar{\xi}_i$, and setting $M_0 = \sum A\omega_i$, we get
\[
  l(\bar{M}_0) = l(M_0/M_0 \cap aM) \le l(M_0/aM_0) \le r \cdot l(A/aA).
\]
The right-hand side is now independent of $\bar{M}_0$, so that $\bar{M}$ is in fact finitely
generated, and $l(\bar{M}) \le r \cdot l(A/aA)$.
\end{proof}

We return to the proof of the theorem. We can replace $L$ by the field of
fractions of $B$. Set $[L:K] = r$; then $B$ is a torsion-free $A$-module of rank $r$.
Hence by the lemma, for any $0 \neq a \in A$ we have $l_A(B/aB) < \infty$. Now if $J \neq 0$ is
an ideal of $B$ and $0 \neq b \in J$ then since $b$ is algebraic over $A$ it satisfies a
relation
\[
  a_0b^m + a_1b^{m-1} + \cdots + a_{m-1}b + a_m = 0 \quad \text{with } a_i \in A.
\]
$B$ is an integral domain, so that we can assume $a_0 \neq 0$. Then $0 \neq a_0 \in J \cap A$
and so
\[
  l_A(B/J) \le l_A(B/a_0B) < \infty.
\]
Moreover, one sees from $l_A(J/a_0B) \le l_A(B/a_0B) < \infty$ that $J/a_0B$
is a finite $B$-module; hence, $J$ itself is a finite $B$-module, and therefore $B$ is
Noetherian. If $P$ is a non-zero prime ideal of $B$ then $B/P$ is an Artinian
ring and an integral domain, and therefore a field. Thus $P$ is maximal and
$\dim B = 1$.
\end{proof}

\begin{corollary*}
Let $A$ be a one-dimensional Noetherian integral domain, $K$ its
field of fractions, and $L$ a finite algebraic extension field of $K$; write $B$ for the
integral closure of $A$ in $L$. Then $B$ is a Dedekind ring, and for any maximal
ideal $P$ of $A$ there are just a finite number of primes of $B$ lying over $P$.
\end{corollary*}

\begin{proof}
By the theorem $B$ is a one-dimensional Noetherian integral domain
and is normal by construction; hence it is a Dedekind ring. It is easy to see
that if we factorise $PB$ as a product $PB = Q_1^{e_1} \cdots Q_r^{e_r}$ of a finite number of
prime ideals, then $Q_1, \ldots, Q_r$ are all the prime ideals of $B$ lying over $P$.
\end{proof}

% ============ 习题 ============
\section*{Exercises to \S\ref{sec:11}}
\emph{Prove the following propositions.}

\begin{exercise}\label{ex:11.1-ex}
Let $A$ be a DVR, $K$ its field of fractions, and $\bar{K}$ an algebraic closure of $K$;
then any valuation ring of $\bar{K}$ dominating $A$ is a one-dimensional non-
discrete valuation ring.
\end{exercise}

\begin{exercise}\label{ex:11.2-ex}
Let $A$ be a DVR, $K$ its field of fractions, and $L$ a finite extension field of $K$;
then a valuation ring of $L$ dominating $A$ is a DVR.
\end{exercise}

\begin{exercise}\label{ex:11.3-ex}
Let $A$ be a DVR and $m$ its maximal ideal; then the $m$-adic completion $\hat{A}$ of
$A$ is again a DVR.
\end{exercise}

\begin{exercise}\label{ex:11.4-ex}
Let $v\colon K \to \R \cup \{\infty\}$ be an Archimedean additive valuation of a field $K$,
and let $c$ be a real number with $0 < c < 1$. For $\alpha, \beta \in K$, set $d(\alpha, \beta) = c^{v(\alpha-\beta)}$;
then $d$ satisfies the axioms for a metric on $K$ (that is $d(\alpha, \beta) \ge 0$, $d(\alpha, \beta)
= 0 \iff \alpha = \beta$, $d(\alpha, \beta) = d(\beta, \alpha)$ and $d(\alpha, \gamma) \le d(\alpha, \beta) + d(\beta, \gamma)$), and the to-
pology of $K$ defined by $d$ does not depend on the choice of $c$. Let $R$ be the
valuation ring of $v$ and $m$ its valuation ideal; if $R$ is a DVR then the
topology determined by $d$ restricts to the $m$-adic topology on $R$.
\end{exercise}

\begin{exercise}\label{ex:11.5-ex}
Any ideal in a Dedekind ring can be generated by at most two elements.
\end{exercise}

\begin{exercise}\label{ex:11.6-ex}
Let $A$ be the integral closure of $\Z$ in $\Q(\sqrt{10})$; then $A$ is a Dedekind ring but
not a principal ideal ring.
\end{exercise}

\begin{exercise}\label{ex:11.7-ex}
If a Dedekind ring $A$ is semilocal then it is a principal ideal ring.
\end{exercise}

\begin{exercise}\label{ex:11.8-ex}
A module over a Dedekind ring is flat if and only if it is torsion-free.
\end{exercise}

\begin{exercise}\label{ex:11.9-ex}
Let $A$ be an integral domain (not necessarily Noetherian). The following
two conditions are equivalent:
\begin{enumerate}[(1)]
  \item $A_P$ is a valuation ring for every maximal ideal $P$ of $A$;
  \item an $A$-module is flat if and only if it is torsion-free. (An integral
        domain satisfying these conditions is called a \emph{Pr\"ufer domain}.)
\end{enumerate}
\end{exercise}

\begin{exercise}\label{ex:11.10-ex}
A finite torsion-free module over a Dedekind ring is projective, and is
isomorphic to a direct sum of ideals.
\end{exercise}

% ============ §12 Krull rings (原书 pp.86-91) ============
\section{Krull rings}\label{sec:12}\index{Krull rings, section 12}

Let $A$ be an integral domain and $K$ its field of fractions. We write
$K^*$ for the multiplicative group of $K$. We say that $A$ is a \emph{Krull ring}\index{Krull ring} if there is
a family $\mathcal{F} = \{R_\lambda\}_{\lambda \in \Lambda}$ of DVRs of $K$ such that the following two conditions
hold, where we write $v_\lambda$ for the normalised additive valuation correspond-
ing to $R_\lambda$:
\begin{enumerate}[(1)]
  \item $A = \bigcap_\lambda R_\lambda$;
  \item for every $x \in K^*$ there are at most a finite number of $\lambda \in \Lambda$ such that
        $v_\lambda(x) \neq 0$.
\end{enumerate}
The family $\mathcal{F}$ of DVRs is said to be a \emph{defining family} of $A$. Since DVRs are
completely integrally closed (see Ex.~\ref{ex:9.5-ex}), so are Krull rings. If $A$ is a Krull
ring then for any subfield $K' \subset K$ the intersection $A \cap K'$ is again Krull.

\begin{theorem}\label{thm:12.1}
If $A$ is a Krull ring and $S \subset A$ a multiplicative set, then $A_S$ is
again Krull. If $\mathcal{F} = \{R_\lambda\}_{\lambda \in \Lambda}$ is a defining family of $A$ then the subfamily
$\{R_\lambda\}_{\lambda \in \Gamma}$, where $\Gamma = \{\lambda \in \Lambda \mid R_\lambda \supset A_S\}$, is a defining family of $A_S$.
\end{theorem}

\begin{proof}
Setting $m_\lambda$ for the maximal ideal of $R_\lambda$, we have
\[
  \bigcap_{\lambda \in \Gamma} m_\lambda = 0.
\]
Let $0 \neq x \in \bigcap_{\lambda \in \Gamma} R_\lambda$; there are at most finitely many $\lambda \in \Lambda$ such that $v_\lambda(x) < 0$;
let $\Delta = \{\lambda_1, \ldots, \lambda_n\}$ be the set of these. If $s \in S$ then $s \notin m_\lambda$, hence we can find
$t_\lambda \in m_\lambda \cap S$. Replacing $t_\lambda$ by a suitable power, we can assume that $v_\lambda(t_\lambda x) \ge 0$.
We then set $t = \prod_{\lambda \in \Delta} t_\lambda$, so that for every $\lambda \in \Lambda$ we have $v_\lambda(tx) \ge 0$, and
therefore $tx \in A$; but on the other hand $t \in S$ so that $x \in A_S$, and we have
proved that $A_S \supset \bigcap_{\lambda \in \Gamma} R_\lambda$. The opposite inclusion is obvious. The finiteness
condition (2) holds for $A$, so also for the subset $\Gamma$.
\end{proof}

Krull rings defined by a finite number of DVRs have a simple structure.

\begin{lemma}[Nagata]\label{lem:12.1}
Let $K$ be a field and $R_1, \ldots, R_n$ valuation rings of $K$;
set $A = \bigcap R_i$. Then for any given $a \in K$ there exists a natural number $s \ge 2$
such that
\[
  (1 + a + \cdots + a^{s-1})^{-1} \quad \text{and} \quad a \cdot (1 + a + \cdots + a^{s-1})^{-1}
\]
both belong to $A$.
\end{lemma}

\begin{proof}
We consider separately each $R_i$. Note first that $(1 - a)(1 + a + \cdots
+ a^{s-1}) = 1 - a^s$. If $a \notin R_i$ then $a^{-1} \in m_i$, and any $s \ge 2$ will do. If $a \in R_i$ then
provided that there does not exist $t$ such that $1 - a^t \in m_i$, any $s \ge 2$ will do.
If $1 - a \in m_i$ then any $s$ which is not a multiple of the characteristic of $R_i/m_i$
will do. If on the other hand $1 - a \notin m_i$ but $1 - a^t \in m_i$ for some $t > 2$, letting $t_0$
be the smallest value of $t$ for which this happens, we see that $1 - a^s \in m_i$
only for $s$ multiples of $t_0$, so that we only have to avoid these. Thus for each $i$
the bad values of $s$ (if any) are multiples of some number $d_i > 1$, so that
choosing $s$ not divisible by any of these $d_i$ we get the result.
\end{proof}

\begin{theorem}\label{thm:12.2}
Let $K$ be a field and $R_1, \ldots, R_n$ valuation rings of $K$ such
that $R_i \nsubseteq R_j$ for $i \neq j$; set $m_i = \rad(R_i)$. Then the intersection $A = \bigcap_{i=1}^{n} R_i$ is
a semilocal ring, having $p_i = m_i \cap A$ for $i = 1, \ldots, n$ as its only maximal
ideals; moreover $A_{p_i} = R_i$. If each $R_i$ is a DVR then $A$ is a principal ideal
ring.
\end{theorem}

\begin{proof}
The inclusion $A_{p_i} \subset R_i$ is obvious. For the opposite inclusion, let
$a \in R_i$; choosing $s \ge 2$ as in the lemma, and setting $u = (1 + a + \cdots + a^{s-1})^{-1}$
we get $u \in A$ and $au \in A$. Obviously $u$ is a unit of $R_i$, so that $u \in A - p_i$
and $a = (au)/u \in A_{p_i}$. This proves that $A_{p_i} = R_i$. It follows from this
that there are no inclusions among $p_1, \ldots, p_n$. If $I$ is an ideal of $A$ not
contained in any $p_i$ then (by Ex.~\ref{ex:1.6-ex}) there exists $x \in I$ not contained in
$\bigcup_{i=1}^{n} p_i$; then $x$ is a unit in each $R_i$, and hence in $A$, so that $I = A$. Thus
$p_1, \ldots, p_n$ are all the maximal ideals of $A$.

If each $R_i$ is a DVR then we have $m_i \neq m_i^{(2)}$, and hence $p_i \neq p_i^{(2)}$ (where $p^{(2)}$
denotes $p^2A_p \cap A$). Thus there exists $x_i \in p_i$ such that $x_i \notin p_i^{(2)}$, and $x_i \notin p_j$ for
$i \neq j$; then $p_i = x_iA$. If $I$ is any ideal of $A$ and $IR_i = x_i^{e_i}R_i$ for $i = 1, \ldots, n$ then it
is easy to see that $I = x_1^{e_1} \cdots x_n^{e_n}A$.
\end{proof}

If a Krull ring $A$ is defined by an infinite number of DVRs then the
defining family of DVRs is not necessarily unique, but the following
theorem tells us that among them there is a minimal family.

\begin{theorem}\label{thm:12.3}
Let $A$ be a Krull ring, $K$ its field of fractions, and $p$ a height $1$
prime ideal of $A$; then if $\mathcal{F} = \{R_\lambda\}_{\lambda \in \Lambda}$ is a family of DVRs of $K$ defining $A$, we
must have $A_p \in \mathcal{F}$. If we set $\mathcal{F}_0 = \{A_p \mid p \in \Spec A \text{ and } \htop p = 1\}$ then $\mathcal{F}_0$ is a
defining family of $A$. Thus $\mathcal{F}_0$ is the minimal defining family of DVRs of $A$.
\end{theorem}

\begin{proof}
By Theorem 1, $A_p$ is a Krull ring defined by the subfamily $\mathcal{F}_1 =
\{R_\lambda \mid A_p \subset R_\lambda\} \subset \mathcal{F}$; if $A_p \subset R_\lambda$ then the elements of $A - p$ are units of $R_\lambda$,
so that $p \supset m_\lambda \cap A$. If $m_\lambda \cap A = (0)$ then $R_\lambda \supset K$ which is a contradiction,
hence $m_\lambda \cap A \neq (0)$; since $\htop p = 1$, we must have $p = m_\lambda \cap A$. Thus if we fix
some $0 \neq x \in p$, then $v_\lambda(x) > 0$ for all $R_\lambda \in \mathcal{F}_1$, and hence $\mathcal{F}_1$ is a finite set,
Now by the previous theorem and Ex.~\ref{ex:10.5-ex}, the elements of $\mathcal{F}_1$ correspond
bijectively with maximal ideals of $A_p$, and $\mathcal{F}_1$ has just one element $A_p$. Thus
$A_p \in \mathcal{F}$, in other words $\mathcal{F}_0 \subset \mathcal{F}$.

To prove that $\mathcal{F}_0$ is a defining family of DVRs of $A$ it is enough to show that
$A \supset \bigcap_{\htop p = 1} A_p$. That is, it is enough to prove the implication
\[
  \text{for } a, b \in A \text{ with } a \neq 0, \quad b \in aA_p \text{ for all } A_p \in \mathcal{F}_0 \Rightarrow b \in aA.
\]
As one sees easily, this is equivalent to saying that $aA$ can be written as
the intersection of height $1$ primary ideals. The set of $R_\lambda \in \mathcal{F}$ such that
$aR_\lambda \neq R_\lambda$ is finite, so we write $R_1, \ldots, R_r$ for this. If we set
\[
  aR_i \cap A = q_i \quad \text{and} \quad \rad(R_i) \cap A = p_i
\]
then $q_i$ is a primary ideal belonging to $p_i$ for each $i$, and $aA = q_1 \cap \cdots \cap q_r$.
Eliminating redundant terms from this expression, we get an irredundant
expression, say $aA = q_1 \cap \cdots \cap q_r$. It is enough to show that then $\htop p_i = 1$
for $1 \le i \le r$. By contradiction, suppose that $\htop p_1 > 1$. By Theorem 1, $A_{p_1}$ is
a Krull ring with defining family $\mathcal{F}' = \{R_\lambda \in \mathcal{F} \mid A_{p_1} \subset R_\lambda\}$, but is not itself
a DVR, and hence by Theorem 2, $\mathcal{F}'$ is infinite. Thus there exists $R' \in \mathcal{F}'$
such that $aR' = R'$; we set $p' = \rad(R') \cap A$. We have $a \notin p'$, and $A_{p'} \subset R'$
implies that $p' \subset p_1$. Now by assumption $aA = q_1 \cap \cdots \cap q_r$ is irredundant, and $R_1$ is
a DVR, so that $(\rad(R_1))^\nu \subset aR_1$ for some $\nu > 0$, and hence $p_1^\nu \subset q_1$.
Therefore there exists an $i \ge 0$ such that
\[
  aA \neq p'^i \cap q_1 \cap \cdots \cap q_r \quad \text{and} \quad aA \supset p'^{i+1} \cap q_1 \cap \cdots \cap q_r.
\]
Hence there exist $b \in A$ such that $b \notin aA$ but $bp'^i \subset aA$. In particular $bp' \subset aA$,
but since $a$ is a unit of $R'$ we have
\[
  (b/a)p' \subset A \cap \rad(R') = p'.
\]
Taking $0 \neq c \in p'$ then for every $n > 0$ we have $(b/a)^nc \in p' \subset A$, and since $A$ is
completely integrally closed, $b/a \in A$. This is a contradiction, and it proves
that $\htop p_i = 1$ for $1 \le i \le r$.
\end{proof}

\begin{corollary*}
Let $A$ be a Krull ring and $\mathcal{P}$ the set of height $1$ prime ideals of $A$.
For $0 \neq a \in A$ set $v_p(a) = n_p$; then
\[
  aA = \bigcap_{p \in \mathcal{P}} p^{(n_p)},
\]
where $p^{(n)}$ denotes the symbolic $n$th power $p^nA_p \cap A$.
\end{corollary*}

\begin{proof}
According to the theorem we have $aA = \bigcap_p (aA_p \cap A)$,
but $aA_p = p^{n_p}A_p$, so that $aA_p \cap A = p^{(n_p)}$.
\end{proof}

\begin{theorem}\label{thm:12.4}
\begin{enumerate}[(i)]
  \item A Noetherian normal domain is a Krull ring.
  \item Let $A$ be an integral domain, $K$ its field of fractions, and $L$ an extension
        field of $K$. If $(A_i)_{i \in I}$ is a family of Krull rings contained in $L$ satisfying
        the two conditions (1) $A = \bigcap A_i$ and (2) given any $0 \neq a \in A$ we have
        $aA_i = A_i$ for all but finitely many $i$, then $A$ is a Krull ring.
  \item If $A$ is a Krull ring then so is $A[X]$ and $A[[X]]$.
\end{enumerate}
\end{theorem}

\begin{proof}
(i) This follows from Theorem~\ref{thm:11.5} and the fact that for any non-
zero $a \in A$ there are only finitely many height $1$ prime ideals containing $aA$
(because these are the prime divisors of $aA$).

(ii) is easy, and we leave it to the reader.

(iii) $K[X]$ is a principal ideal ring and therefore a Krull ring. Moreover,
if we let $\mathcal{P}$ be the set of height $1$ prime ideals of $A$ then for $p \in \mathcal{P}$
the ideal $p[X]$ is prime in $A[X]$, and by Theorem~\ref{thm:11.2}, (3), the local ring
$A[X]_{p[X]}$ is a DVR of $K(X)$. (If we write $v$ for the additive valuation of $K$
corresponding to the valuation ring $A_p$, we can extend $v$ to an additive
valuation of $K(X)$ by setting $v(F(X)) = \min\{v(a_i)\}$ for a polynomial
$F(X) = a_0 + a_1X + \cdots + a_rX^r$ (with $a_i \in K$),
and $v(F/G) = v(F) - v(G)$ for a rational function $F(X)/G(X)$; then the
valuation ring of $v$ in $K(X)$ is $A[X]_{p[X]}$.) Now we have
\[
  K[X] \cap A[X]_{p[X]} = A_p[X]
\]
(prove this!), and so
\[
  A[X] = \bigcap_{p \in \mathcal{P}} A_p[X],
\]
by (ii) this is a Krull ring.

Now for $A[[X]]$, let $\{R_\lambda\}_{\lambda \in \Lambda}$ be a family of DVRs of $K$ defining $A$; then
inside $K[[X]]$ we have $A[[X]] = \bigcap_\lambda R_\lambda[[X]]$; also by Ex.~\ref{ex:9.5-ex}, $R_\lambda[[X]]$ is an
integrally closed Noetherian ring, and is therefore a Krull ring by (i).
However, we cannot use (ii) as it stands, since $X$ is a non-unit of all the rings
$R_\lambda[[X]]$, so we set $R_\lambda[[X]][X^{-1}] = B_\lambda$ and note that $A[[X]] = K[[X]] \cap
(\bigcap_\lambda B_\lambda)$; now the hypothesis in (ii) is easily verified. Indeed,
\[
  F(X) = a_rX^r + a_{r+1}X^{r+1} + \cdots \in A[[X]] \quad \text{with } a_r \neq 0
\]
is a non-unit of $B_\lambda$ if and only if $a_r$ is a non-unit of $R_\lambda$, and there are only
finitely many such $\lambda$. Therefore $A[[X]]$ is a Krull ring.
\end{proof}

\begin{remark}
Note that the field of fractions of $A[[X]]$ is in general smaller
than the field of fractions of $K[[X]]$.
\end{remark}

\begin{remark}
The $B_\lambda$ occurring above are Euclidean rings ([B7], \S1, Ex. 9).
\end{remark}

\begin{theorem}\label{thm:12.5}
The notions of Dedekind ring and one-dimensional Krull
ring coincide.
\end{theorem}

\begin{proof}
A Dedekind ring is a normal Noetherian domain, and therefore a
Krull ring. Conversely, if $A$ is a one-dimensional Krull ring, let us prove that
$A$ is Noetherian. Let $I$ be a non-zero ideal of $A$, and let $0 \neq a \in I$. If we can
prove that $A/aA$ is Noetherian then $I/aA$ is finitely generated, and thus so is
$I$. By the corollary of Theorem 3 we can write $aA = q_1 \cap \cdots \cap q_r$, where
$q_i$ are symbolic powers of prime ideals $p_i$ and $p_i \neq p_j$ if $i \neq j$; now since
$\dim A = 1$ each $p_i$ is maximal and we have
\[
  A/aA = A/q_1 \times \cdots \times A/q_r
\]
by Theorem~\ref{thm:1.3} and Theorem~\ref{thm:1.4}. But $A/q_i$ is a local ring with maximal
ideal $p_i/q_i$, and hence $A/q_i \cong A_{p_i}/q_iA_{p_i}$; now since each $A_{p_i}$ is a DVR, $A/aA$ is
Noetherian (in fact even Artinian). Hence $A$ is one-dimensional Noetherian
integral domain, and is normal, and is therefore a Dedekind ring.
\end{proof}

\begin{theorem}\label{thm:12.6}
Let $A$ be a Krull ring, $K$ its field of fractions, and write $\mathcal{P}$
for the set of height $1$ prime ideals of $A$. Suppose given any $p_1, \ldots, p_r \in \mathcal{P}$
and $e_1, \ldots, e_r \in \Z$. Then there exists $x \in K$ satisfying
\[
  v_{p_i}(x) = e_i \quad \text{for } 1 \le i \le r
\]
and
\[
  v_p(x) \ge 0 \quad \text{for all } p \in \mathcal{P} - \{p_1, \ldots, p_r\}.
\]
Here $v_{p_i}$ and $v_p$ stand for the normalised additive valuations of $K$
corresponding to $p_i$ and $p$.
\end{theorem}

\begin{proof}
If $y_i \in A$ is chosen so that $y_i \in p_i$ but $y_i \notin p_i^{(2)}$ and $y_i \notin p_j$ for
$j \neq i$, then $v_{p_j}(y_i) = \delta_{ij}$ for $1 \le i, j \le r$. Similarly we choose $y_1, \ldots, y_r \in A$ such
that $v_{p_i}(y_j) = \delta_{ij}$. Then we set
\[
  y = \prod_{i=1}^{r} y_i^{e_i};
\]
let $p_1', \ldots, p_s'$ be all the primes $p \in \mathcal{P} - \{p_1, \ldots, p_r\}$ for which $v_p(y) \neq 0$.
Then choosing for each $j = 1, \ldots, s$ an element $t_j \in p_j'$ not belonging to
$p_1 \cup \cdots \cup p_r$, and taking $\nu$ to be sufficiently large, we see that
\[
  x = y(t_1 \cdots t_s)^\nu
\]
satisfies the requirements of the theorem.
\end{proof}

\begin{theorem}[Y. Mori and J. Nishimura]\label{thm:12.7}
Let $A$ and $\mathcal{P}$ be as in the previous
theorem. If $A/p$ is Noetherian for every $p \in \mathcal{P}$ then $A$ is Noetherian.
\end{theorem}

\begin{proof}[Proof (J. Nishimura)]
As in the proof of Theorem 5, it is enough to show
that $A/p^{(n)}$ is Noetherian for $p \in \mathcal{P}$ and any $n > 0$. If $n = 1$ this holds by
hypothesis. For $n > 1$ we proceed as follows. Applying the previous
theorem with $r = 1$ and $e = -1$ we can find an element $x$ in the field of
fractions of $A$ such that $v_p(x) = 1$ and $v_q(x) < 0$ for every other $q \in \mathcal{P}$. Set
$B = A[x]$. If $y \in p$ then $y/x \in A$ so that $y \in xB$, and conversely $B \subset A_p$, and
$xB \subset pA_p$, so that $p = xB \cap A$. Moreover, $B = A + xB$ and so
\[
  B/xB = A/p.
\]
Now $x^iB/x^{i+1}B \cong B/xB$ for each $i$, so that by induction on $i$ we see that
$B/x^iB$ is a Noetherian $B$-module for each $i$, and hence a Noetherian ring.
Now we have
\[
  A/(x^nB \cap A) \hookrightarrow B/x^nB,
\]
and $B/x^nB$ is finite over $A/(x^nB \cap A)$, being generated by the images of
$1, x, \ldots, x^{n-1}$, so that by the Eakin--Nagata theorem (Theorem~\ref{thm:3.7}),
$A/(x^nB \cap A)$ is a Noetherian ring; therefore its quotient $A/p^{(n)}$ is also
Noetherian.
\end{proof}

\begin{remark}
If $A$ is a Noetherian integral domain and $K$ its field of fractions,
then the integral closure of $A$ in $K$ is a possibly non-Noetherian Krull ring
([N1], (33.10)). This was proved by Y. Mori (1952) in the local case, and
in the general case by M. Nagata (1955). Theorem~\ref{thm:12.7} was proved by
Mori [1] in 1955 as a theorem on the integral closure of Noetherian
rings. His proof was correct (in spite of a number of easily rectifiable
inaccuracies), and was an extremely interesting piece of work, but due to
its difficulty, and the fact that it appeared in an inaccessible journal, the
result was practically forgotten. After Marot [1], (1973) applied it
successfully, Mori's work attracted attention once more, and J. Nishimura
[1], (1975) reformulated the result as above as a theorem on Krull rings
and gave an elegant proof.

More results on Krull rings can be found in [N1], [B7], [F], among
others.
\end{remark}

% ============ 习题 ============
\section*{Exercises to \S\ref{sec:12}}
\emph{Prove the following propositions.}

\begin{exercise}\label{ex:12.1-ex}
Let $K \subset L$ be a finite extension of fields, and $R$ a valuation ring of $K$. Then
there are a finite number of valuation rings of $L$ dominating $R$, and if $L$ is a
normal extension of $K$ then these are all conjugate to one another under
elements of the Galois group $\Aut_K(L)$.
\end{exercise}

\begin{exercise}\label{ex:12.2-ex}
Let $R$ be a valuation ring of a field $K$, and let $K \subset L$ be a (possibly infinite)
algebraic extension; write $\bar{R}$ for the integral closure of $R$ in $L$. Then the
localisation of $\bar{R}$ at a maximal ideal is a valuation ring dominating $R$, and
conversely every valuation ring of $L$ dominating $R$ is obtained in this way.
\end{exercise}

\begin{exercise}\label{ex:12.3-ex}
Let $A$ be a Krull ring, $K$ its field of fractions, and $K \subset L$ a finite extension
field; if $B$ is the integral closure of $A$ in $L$, then $B$ is also a Krull ring.
\end{exercise}

\begin{exercise}\label{ex:12.4-ex}
Let $A$ be an integral domain and $K$ its field of fractions. For $I$ a fractional
ideal, write $\tilde{I} = (I^{-1})^{-1}$. If $I = \tilde{I}$ we say that $I$ is \emph{divisorial}. If $A$ is a Krull
ring, then an ideal of $A$ is divisorial if and only if it can be expressed as the
intersection of a finite number of height $1$ primary ideals.
\end{exercise}


% ============ Chapter 5. Dimension theory (原书 pp.92-123) ============
\chapter{Dimension theory}\label{ch:5}

The dimension theory of Noetherian rings is probably the greatest of
Krull's many achievements; with his principal ideal theorem
(Theorem~\ref{thm:13.5}) the theory of Noetherian rings gained in mathematical
profundity. Then the theory of multiplicities was first treated rigorously
and in considerable generality by Chevalley, and was simplified by Samuel's
definition of multiplicity in terms of the Samuel function.

Here we follow the method of EGA, proving Theorem~\ref{thm:13.4} via the
Samuel function, and deducing the principal ideal theorem as a corollary.
The Samuel function is of importance as a measure of singularity in
Hironaka's resolution of singularities, but in this book we can only cover its
basic properties. In \S\ref{sec:15} we exploit the notion of systems of parameters to
discuss among other things the dimension of the fibres of a ring
homomorphism and the dimension formula for finitely generated extension
rings.

% ============ §13 Graded rings, the Hilbert function and the Samuel function (原书 pp.92-101) ============
\section{Graded rings, the Hilbert function and the Samuel function}\label{sec:13}\index{graded rings, section 13}\index{Hilbert function, section 13}\index{Samuel function, section 13}

Let $G$ be an Abelian semigroup with identity element $0$; (that is, $G$
is a set with an addition law $+$ satisfying associativity $(x + y) + z = x +
(y + z)$, commutativity $x + y = y + x$, and such that $0 + x = x$). A \emph{graded} (or
$G$-graded) ring is a ring $R$ together with a direct sum decomposition of $R$ as
an additive group $R = \bigoplus_{i \in G} R_i$ satisfying $R_iR_j \subset R_{i+j}$. Similarly, a \emph{graded}
$R$-module is an $R$-module $M$ together with a direct sum decomposition $M =
\bigoplus_{i \in G} M_i$ satisfying $R_iM_j \subset M_{i+j}$. An element $x \in M$ is \emph{homogeneous}\index{homogeneous element} if
$x \in M_i$ for some $i \in G$, and $i$ is then called the \emph{degree} of $x$. A general element
$x \in M$ can be written uniquely in the form $x = \sum_{i \in G} x_i$ with $x_i \in M_i$ and only
finitely many $x_i \neq 0$; $x_i$ is called the \emph{homogeneous term} of $x$ of degree $i$.

A submodule $N \subset M$ is called a \emph{homogeneous submodule}\index{homogeneous submodule} (or \emph{graded}
submodule) if it can be generated by homogeneous elements. This
condition is equivalent to either of the following two:
\begin{enumerate}[(1)]
  \item For $x \in M$, if $x \in N$ then each homogeneous term of $x$ is in $N$;
\end{enumerate}
For a homogeneous submodule $N \subset M$ we set $N_i = N \cap M_i$; then
$M/N = \bigoplus_{i \in G} M_i/N_i$ is again a graded $R$-module.

One sees from the definition that $R_0 \subset R$ is a subring, and that each
graded piece $M_i$ of a graded $R$-module $M$ is an $R_0$-module.

The notion of graded ring is most frequently used when $G$ is the
semigroup $\{0, 1, 2, \ldots\}$ of non-negative integers, which we denote by $\N$. In this
case, we set $R^+ = \bigoplus_{n > 0} R_n$; then $R^+$ is an ideal of $R$, with $R/R^+ \cong R_0$.

The polynomial ring $R = R_0[X_1, \ldots, X_r]$ over a ring $R_0$ is usually made
into an $\N$-graded ring by defining the degree of a monomial $X_1^{\alpha_1} \cdots X_r^{\alpha_r}$ as
the total degree $\alpha_1 + \cdots + \alpha_r$; however, $R$ has other useful gradings. For
example, $R$ has an $\N^r$-grading in which $X_1^{\alpha_1} \cdots X_r^{\alpha_r}$ has degree $(\alpha_1, \ldots, \alpha_r)$;
the value of systematically using this grading can be seen in Goto--
Watanabe [1]. Alternatively, giving each of the $X_i$ some suitable weight
$d_i$ and letting the monomial $X_1^{\alpha_1} \cdots X_r^{\alpha_r}$ have weight $\sum \alpha_i d_i$ defines an
$\N$-grading of $R$. For example, the ring $R_0[X, Y, Z]/(f)$, where $f = a_1X^{\alpha_1} +
a_2Y^{\alpha_2} + a_3Z^{\alpha_3}$ can be graded by giving the images of $X, Y, Z$ the weights
$\alpha_1, \alpha_2$ and $\alpha_3$, respectively.

A \emph{filtration}\index{filtration} of a ring $A$ is a descending chain $A = J_0 \supset J_1 \supset \cdots$ of ideals
such that $J_nJ_m \subset J_{n+m}$; the associated graded ring $\gr(A)$ is defined as
follows. First of all as a module we set $\gr_n(A) = J_n/J_{n+1}$ for $n \ge 0$, and
$\gr(A) = \bigoplus_{n \ge 0} \gr_n(A)$; then we define the product by
\[
  (x + J_{n+1})(y + J_{m+1}) = xy + J_{n+m+1} \quad \text{for } x \in J_n \text{ and } y \in J_m.
\]
It is easy to see that $\gr(A)$ becomes a graded ring. The filtration
$J_1 \supset J_2 \supset \cdots$ defines a linear topology on $A$ (see \S\ref{sec:8}), and the completion $\hat{A}$ of
$A$ in this topology has a filtration $\hat{J}_1 \supset \hat{J}_2 \supset \cdots$ such that $\hat{A}/\hat{J}_n \cong A/J_n$ for
all $n$, hence $\hat{J}_n/\hat{J}_{n+1} \cong J_n/J_{n+1}$, and
\[
  \gr(\hat{A}) = \gr(A).
\]

Let $A$ be a ring, $I$ an ideal, and let $B = \bigoplus_{n \ge 0} I^n/I^{n+1}$ be the graded ring
associated with the filtration $I \supset I^2 \supset \cdots$ of $A$ by powers of $I$; the various
notations $\gr_I(A)$, $\gr^I(A)$ and $G_I(I)$ are used to denote $B$ in the current
literature. An element of $B_n = I^n/I^{n+1}$ can be expressed as a linear
combination of products of $n$ elements of $B_1 = I/I^2$, so that $B$ is generated
over the subring $B_0 = A/I$ by elements of $B_1$. If $I = Ax_1 + \cdots + Ax_r$ and $\bar{x}_i$
denotes the image of $x_i$ in $B_1 = I/I^2$ then
\[
  B = \gr_I(A) = (A/I)[\bar{x}_1, \ldots, \bar{x}_r],
\]
and $B$ is a quotient of the polynomial ring $(A/I)[X_1, \ldots, X_r]$ as a graded
ring.

\begin{theorem}\label{thm:13.1}
An $\N$-graded ring $R = \bigoplus_{n \ge 0} R_n$ is Noetherian if and only if
$R_0$ is Noetherian and $R$ is finitely generated as a ring over $R_0$.
\end{theorem}

\begin{proof}
The `if' is obvious, and we prove the `only if': suppose that $R$
is Noetherian. Then since $R_0 \cong R/R^+$, $R_0$ is Noetherian. $R^+$ is a home-
geneous ideal, and is finitely generated, so that we can suppose that it is
generated by homogeneous elements $x_1, \ldots, x_r$. Then it is easy to see that
$R = R_0[x_1, \ldots, x_r]$; in fact it is enough to show that $R_n \subset R_0[x_1, \ldots, x_r]$
for every $n$. Now writing $d_i$ for the degree of $x_i$ we have
\[
  (*) \quad R_n = x_1R_{n-d_1} + x_2R_{n-d_2} + \cdots + x_rR_{n-d_r}.
\]
Indeed, for $y \in R_n$, write $y = \sum c_ix_i$ with $f_i \in R$; then setting $g_i$ for the
homogeneous term of degree $n - d_i$ of $f_i$ (with $g_i = 0$ if $n - d_i < 0$), we also
have $y = \sum x_ig_i$. From $(*)$ it follows by induction that $R_n \subset R_0[x_1, \ldots, x_r]$.
\end{proof}

Let $R = \bigoplus_{n \ge 0} R_n$ be a Noetherian graded ring; then if $M = \bigoplus_{n \ge 0} M_n$
is a finitely generated graded $R$-module, each $M_n$ is finitely generated
as $R_0$-module. In fact when $M = R$ this is clear from $(*)$ above. In the general
case $M$ can be generated by a finite number of homogeneous elements $\omega_i$: $M
= \sum R\omega_i$. Now letting $e_i$ be the degree of $\omega_i$, we have as above that
\[
  M_n = R_{n-e_1}\omega_1 + \cdots + R_{n-e_r}\omega_r \quad \text{(where } R_i = 0 \text{ for } i < 0),
\]
and hence $M_n$ is a finite $R_0$-module. In particular if $R_0$ is an Artinian ring,
then $l(M_n) < \infty$, where $l$ denotes the length of an $R_0$-module. In this case we
define the \emph{Hilbert series}\index{Hilbert series} $P(M, t)$ of $M$ by the formula:
\[
  P(M, t) = \sum_{n=0}^{\infty} l(M_n)t^n \in \Z[[t]].
\]
(In combinatorics it is a standard procedure to associate with a sequence of
numbers $a_0, a_1, a_2, \ldots$ the generating function $\sum a_it^i$.)

\begin{theorem}\label{thm:13.2}
Let $R = \bigoplus_{n \ge 0} R_n$ be a Noetherian graded ring with $R_0$
Artinian, and let $M$ be a finitely generated graded $R$-module. Suppose that
$R = R_0[x_1, \ldots, x_r]$ with $x_i$ of degree $d_i$, and that $P(M, t)$ is as above. Then
$P(M, t)$ is a rational function of $t$, and can be written
\[
  P(M, t) = f(t) \Big/ \prod_{i=1}^{r} (1 - t^{d_i}),
\]
where $f(t)$ is a polynomial with coefficients in $\Z$.
\end{theorem}

\begin{proof}
By induction on $r$, the number of generators of $R$. When $r = 0$ we
have $R = R_0$, so that for $n$ sufficiently large, $M_n = 0$, and the power series
$P(M, t)$ is a polynomial. When $r > 0$, multiplication by $x_r$ defines an $R_0$-
linear map $M_n \to M_{n+d_r}$; writing $K_n$ and $L_{n+d_r}$ for the kernel and cokernel,
we get an exact sequence
\[
  0 \to K_n \to M_n \to M_{n+d_r} \to L_{n+d_r} \to 0.
\]
Set $K = \bigoplus K_n$, and $L = \bigoplus L_n$. Then $K$ is a submodule of $M$, and $L = M/x_rM$,
so that $K$ and $L$ are finite $R$-modules; moreover $x_rK = x_rL = 0$, so that $K$
and $L$ can be viewed as $R/x_rR$-modules, and hence we can apply the
induction hypothesis to $P(K, t)$ and $P(L, t)$. Now from the above exact
sequence we get
\[
  l(K_n) - l(M_n) + l(M_{n+d_r}) - l(L_{n+d_r}) = 0.
\]
If we multiply this by $t^{n+d_r}$ and sum over $n$ this gives
\[
  t^{d_r}P(K, t) - t^{d_r}P(M, t) + P(M, t) - P(L, t) = g(t),
\]
where $g(t) \in \Z[t]$. The theorem follows at once from this.
\end{proof}

A lot of information on the values of $l(M_n)$ can be obtained from the
above theorem. Especially simple is the case $d_1 = \cdots = d_r = 1$, so that $R$ is
generated over $R_0$ by elements of degree $1$. In this case $P(M, t) = f(t)
(1 - t)^{-r}$; if $f(t)$ has $(1 - t)$ as a factor we can cancel to get $P$ in the form
\[
  P(M, t) = f(t)(1 - t)^{-d} \quad \text{with } f \in \Z[t], \; d \ge 0,
\]
and if $d > 0$ then $f(1) \neq 0$.
If this holds, we will write $d = d(M)$. Since $(1 - t)^{-1} = 1 + t + t^2 + \cdots$, we can
repeatedly differentiate both sides to get
\[
  (1 - t)^{-d} = \sum_{m \ge 0} \binom{m + d - 1}{d - 1} t^m.
\]
(This can of course be proved in other ways, for example by induction on $d$.)
If $f(t) = a_0 + a_1t + \cdots + a_st^s$ then
\[
  (*) \quad l(M_n) = a_0\binom{n+d-1}{d-1} + a_1\binom{n+d-2}{d-1} + \cdots + a_s\binom{n-s+d-1}{d-1},
\]
where we interpret $\binom{m}{d-1} = 0$ for $m < d - 1$. The right-hand side of $(*)$ can be
formally rearranged as a polynomial in $n$ with rational coefficients, say $\varphi(n)$;
then
\[
  \varphi(X) = \frac{f(1)}{(d-1)!}X^{d-1} + \text{(terms of lower degree)}.
\]
Since $\binom{m + d - 1}{d-1}$ coincides with the polynomial $m(m - 1) \cdots (m - d + 2)/
(d-1)!$ for $m \ge 0$, this implies the following result.

\begin{corollary*}
If $d_1 = \cdots = d_r = 1$ in Theorem 2, and $d = d(M)$ is defined
as above, then there is a polynomial $\varphi_M(X)$ of degree $d - 1$ with rational
coefficients such that for $n \ge s + 1 - d$ we have $l(M_n) = \varphi_M(n)$. Here $s$ is the
degree of the polynomial $(1 - t)^dP(M, t)$.
\end{corollary*}

The polynomial $\varphi_M$ appearing here is called the \emph{Hilbert polynomial}\index{Hilbert polynomial} of the
graded module $M$. The numerical function $l(M_n)$ itself is called the \emph{Hilbert}
\emph{function} of $M$; by the degree of a Hilbert function we mean the degree of the
corresponding Hilbert polynomial.

\begin{remark}
For general $d_1, \ldots, d_r$, it is no longer necessarily the case that $l(M_n)$
can be represented by one polynomial.
\end{remark}

\begin{example}\label{ex:13.1}
When $R = R_0[X_1, \ldots, X_r]$, the number of monomials of
degree $n$ is $\binom{n+r-1}{r-1}$, so that
\[
  l(R_n) = l(R_0)\binom{n+r-1}{r-1}
\]
holds for every $n \ge 0$, and the right-hand side is $\varphi_R(n)$. Thus $\varphi_R(X) =
(l(R_0)/r!)(X + r)(X + r - 1) \cdots (X + 1)$.
\end{example}

\begin{example}\label{ex:13.2}
Let $k$ be a field, and $F(X_1, \ldots, X_r)$ a homogeneous poly-
nomial of degree $s$; set $R = k[X_1, \ldots, X_r]/(F(X))$. Then for $n \ge s$,
\[
  l(R_n) = \binom{n+r-1}{r-1} - \binom{n-s+r-1}{r-1},
\]
and hence, setting $\binom{n+r-1}{r-1} = (1/r!)n^{r-1} + a_1n^{r-2} + \cdots$, we have
\[
  \varphi_R(X) = \frac{1}{r!}[X^{r-1} - (X - s)^{r-1}] + a_1[X^{r-2} - (X - s)^{r-2}] + \cdots
\]
\[
  = \frac{s}{(r-1)!}X^{r-2} + \text{(terms of lower degree)}.
\]
\end{example}

\begin{example}\label{ex:13.3}
Let $k$ be a field, and $R = k[X_1, \ldots, X_r]/P = k[\xi_1, \ldots, \xi_r]$,
where $P$ is a homogeneous prime ideal. Let $t$ be the transcendence degree
of $R$ over $k$, and suppose that $\xi_1, \ldots, \xi_t$ are algebraically independent
over $k$; then there are $\binom{n+t-1}{t-1}$ monomials of degree $n$ in the
$\xi_1, \ldots, \xi_t$, and these are linearly independent over $k$, so that
\[
  l(R_n) \ge \binom{n+t-1}{t-1},
\]
from which it follows that $d \ge t$. In fact we will prove
later (Theorem 8) that $d = t$.
\end{example}

A homogeneous ideal of the polynomial ring $k[X_1, \ldots, X_r]$ over a field
$k$ defines an algebraic variety in $r$-dimensional projective space $\mathbb{P}^r$, and
the Hilbert polynomial plays an important role in algebraic geometry.
For example, note that the numerator of the leading term of $\varphi_R$ in Example
2 is equal to the degree of $F$. This holds in more generality, but we must
leave details of this to textbooks on algebraic geometry [Ha].

The idea of using the construction of $\gr_m(A)$ to relate the study of a
general Noetherian local ring $(A, m)$ to the theory of ideals in a polynomial
ring over a field was one of the crucial ideas introduced by Krull in his
article `Dimension theory of local rings' [6], a work of monumental
significance for the theory of Noetherian rings. If $m$ is generated by $r$
elements then $\gr_m(A)$ is of the form $k[X_1, \ldots, X_r]/I$, where $k = A/m$ and $I$ is a
homogeneous ideal. However, the Hilbert function of this graded ring was
first used in the study of the multiplicity of $A$ by P. Samuel (1951).

\subsection*{Samuel functions}

In a little more generality, let $A$ be a Noetherian semilocal ring, and $m$
the Jacobson radical of $A$. If $I$ is an ideal of $A$ such that for some $\nu > 0$
we have $m^\nu \subset I \subset m$, we call $I$ an \emph{ideal of definition}\index{ideal of definition}; the $I$-adic and $m$-adic
topologies then coincide, so that `ideal of definition' means `ideal defining
the $m$-adic topology'. Let $M$ be a finite $A$-module. If we set
\[
  \gr_I(M) = \bigoplus_{n \ge 0} I^nM/I^{n+1}M
\]
then $\gr_I(M)$ is in a natural way a graded module over $\gr_I(A) = \bigoplus I^n/I^{n+1}$.
For brevity write $\gr_I(A) = A^*$ and $\gr_I(M) = M^*$. Then the ring $A^*_0 = A/I$ is
Artinian, and if $I = \sum_{i=1}^{r} x_iA$, and $\bar{x}_i$ is the image of $x_i$ in $I/I^2$, then
$A^* = A^*_0[\bar{x}_1, \ldots, \bar{x}_r]$. If also $M = \sum A\omega_i$ then $M^* = \sum A^*\bar{\omega}_i$ (where $\bar{\omega}_i$ is the
image of $\omega_i$ in $M_0^* = M/IM$), so that we can apply Theorem 2 and its
corollary to $M^*$. Noting that $l(M_n^*) = l(I^nM/I^{n+1}M)$ (where on the left-hand
side $l$ is the length as an $A^*_0$-module, on the right-hand side as an $A$-module),
we have
\[
  \sum_{n=0}^{\infty} l(M_n^*)t^n = \sum_{n=0}^{\infty} l(I^nM/I^{n+1}M)t^n.
\]
We now set $\chi_M^I(n) = l(M/I^{n+1}M)$. In particular we abbreviate $\chi_A^I(n)$ to
$\chi_A(n)$, and call it the \emph{Samuel function}\index{Samuel function} of the $A$-module $M$.

Repeatedly using the well-known formula
\[
  \binom{n}{r} = \binom{n-1}{r} + \binom{n-1}{r-1}
\]
we get
\[
  l(M/I^{n+1}M) = \sum_{k=0}^{n} l(I^kM/I^{k+1}M),
\]
so that from formula $(*)$ on p. 95 we get
\[
  \chi_M^I(n) = a_dn^d + a_{d-1}n^{d-1} + \cdots + a_0 \quad \text{with } a_i \in \Z.
\]
When $n \ge s$ this is a polynomial in $n$ of degree $d$. This degree $d$ is determined by $M$, and does not depend on $I$; to see this, if $I$ and $J$ are
both ideals of definition of $A$ then there exist natural numbers $a$ and $b$
such that $I^a \subset J$, $J^b \subset I$, so that
\[
  \chi_M^I(an + a - 1) \ge \chi_M^J(n) \quad \text{and} \quad \chi_M^J(bn + b - 1) \ge \chi_M^I(n).
\]
We thus write $d = d(M)$. It is natural to think of $d(M)$ as a measure of the size
of $M$.

\begin{theorem}\label{thm:13.3}
Let $A$ be a semilocal Noetherian ring, and $0 \to M' \to
M \to M'' \to 0$ an exact sequence of finite $A$-modules; then
\[
  d(M) = \max(d(M'), d(M'')).
\]
If $I$ is any ideal of definition of $A$, then $\chi_M^I$ and $\chi_{M''}^I$ have the same
leading coefficient.
\end{theorem}

\begin{proof}
We can assume $M'' = M/M'$. Then since $M''/I^nM'' = M/(M' + I^nM)$
we have
\[
  l(M/I^nM) = l(M/M' + I^nM) + l(M' + I^nM/I^nM)
\]
\[
  = l(M''/I^nM'') + l(M'/M' \cap I^nM).
\]
Thus setting $q(n) = l(M'/M' \cap I^{n+1}M)$, we have $\chi_M^I = \chi_{M''}^I + q$. Since
moreover both $\chi_{M''}^I$ and $q$ take on only positive values, $d(M)$ coincides with
whichever is the greater of $d(M'')$ and $\deg q$. However, by the Artin--Rees
lemma, there is a $c > 0$ such that
\[
  I^{n+1}M \cap M' \subset M' \cap I^{n+1}M \subset I^{n+1-c}M',
\]
and hence
\[
  \chi_{M''}^I(n) \ge q(n) \ge \chi_{M'}^I(n - c);
\]
therefore $q$ and $\chi_{M'}^I$ have the same leading coefficient.
\end{proof}

We now define a further measure $\delta(M)$ of the size of $M$: let $\delta(M)$ be the
smallest value of $n$ such that there exist $x_1, \ldots, x_n \in m$ for which $l(M/x_1M +
\cdots + x_nM) < \infty$. When $l(M) < \infty$, we interpret this as $\delta(M) = 0$. If $I$ is any
ideal of definition of $A$ then $l(M/IM) < \infty$, so that $\delta(M) \le$ number of
generators of $I$. Conversely, in the case that $A$ is a local ring and $M = A$,
then $l(A/I) < \infty$ implies that $I$ is an $m$-primary ideal. Therefore in this case
$\delta(A)$ is the minimum of the number of generators of $m$-primary ideals.

We have now arrived at the fundamental theorem of dimension theory.

\begin{theorem}\label{thm:13.4}
Let $A$ be a semilocal Noetherian ring and $M$ a finite
$A$-module; then we have
\[
  \dim M = d(M) = \delta(M).
\]
\end{theorem}

\begin{proof}
\textbf{Step 1.} Each of $d(M)$ and $\delta(M)$ are finite, but the finiteness of $\dim M$
has not yet been established. First of all, let us prove that $d(A) \ge \dim A$
for the case $M = A$, by induction on $d(A)$. Set $m = \rad(A)$. If $d(A) = 0$ then
$l(A/m^n)$ is constant for $n \gg 0$, so that for some $n$ we have $m^n = m^{n+1}$ and
by NAK, $m^n = 0$. Hence any prime ideal of $A$ is maximal and $\dim A = 0$.

Next suppose that $d(A) > 0$; if $\dim A = 0$ then we're done. If $\dim A > 0$,
consider a strictly increasing sequence $p_0 \subset p_1 \subset \cdots \subset p_e$ of prime ideals
of $A$, choose some element $x \in p_1 - p_0$, and set $B = A/(p_0 + xA)$; then by
the previous theorem applied to the exact sequence
\[
  0 \to A/p_0 \xrightarrow{x} A/p_0 \to B \to 0,
\]
we have $d(B) < d(A)$, and so by induction
\[
  \dim B \le d(B) \le d(A) - 1.
\]
(The values of $d(B)$ and $\dim B$ are independent of whether we consider $B$ as
an $A$-module or as a $B$-module, as is clear from the definitions.) In $B$, the
image of $p_1 \subset \cdots \subset p_e$ provides a chain of prime ideals of length $e - 1$, so that
\[
  e - 1 \le \dim B \le d(A) - 1;
\]
hence $e < d(A)$. Since this holds for any chain of prime ideals of $A$, this
proves $\dim A \le d(A)$. For general $M$, by Theorem~\ref{thm:6.4}, there are submodules
$M_i$ such that $0 = M_0 \subset M_1 \subset \cdots \subset M_r = M$ with
\[
  M_i/M_{i-1} \cong A/p_i, \quad \text{and} \quad p_i \in \Spec A.
\]
Since for an exact sequence $0 \to M' \to M \to M'' \to 0$ of finite $A$-
modules we have
\[
  \Supp(M) = \Supp(M') \cup \Supp(M'')
\]
and
\[
  \dim M = \max(\dim M', \dim M''),
\]
it is easy to see that
\[
  d(M) = \max\{d(A/p_i)\} \ge \max\{\dim(A/p_i)\} = \dim M.
\]

\textbf{Step 2.} We show that $\delta(M) \ge d(M)$. If $\delta(M) = 0$ then $l(M) < \infty$ so that $\chi_M(n)$
is bounded, hence $d(M) = 0$. Next suppose that $\delta(M) = s > 0$, choose
$x_1, \ldots, x_s \in m$ such that $l(M/x_1M + \cdots + x_sM) < \infty$, and set $M_i = M/x_1M +
\cdots + x_iM$; then clearly $\delta(M_i) = \delta(M) - i$. On the other hand,
\[
  l(M_i/m^\nu M_i) = l(M/x_iM + m^\nu M)
\]
\[
  = l(M/m^\nu M) - l(x_iM/x_iM \cap m^\nu M)
\]
\[
  = l(M/m^\nu M) - l(M/(m^\nu M:x_i))
\]
\[
  \ge l(M/m^\nu M) - l(M/m^{\nu-1}M),
\]
so that $d(M_i) \ge d(M) - 1$. Repeating this, we get $d(M_s) \ge d(M) - s$, but since
$\delta(M_s) = 0$ we have $d(M_s) = 0$, so that $s \ge d(M)$.

\textbf{Step 3.} We show that $\dim M \ge \delta(M)$, by induction on $\dim M$. If $\dim M = 0$
then $\Supp(M) \subset \text{m-Spec } A = V(m)$ so that for large enough $n$ we have
$m^n \subset \ann M$, and $l(M) < \infty$, therefore $\delta(M) = 0$. Next suppose that
$\dim M > 0$, and let $p_i$ for $1 \le i \le t$ be the minimal prime divisors of $\ann(M)$
with $\coht p_i = \dim M$; then the $p_i$ are not maximal ideals, so do not contain
$m$. Hence we can choose $x \in m$ not contained in any $p_i$. Setting
$M_1 = M/xM$ we get $\dim M_1 < \dim M$. Therefore by the inductive hypo-
thesis $\delta(M_1) \le \dim M_1$; but obviously $\delta(M) \le \delta(M_1) + 1$, so that $\delta(M) \le
\dim M_1 + 1 \le \dim M$.
\end{proof}

\begin{theorem}\label{thm:13.5}
Let $A$ be a Noetherian ring, and $I = (a_1, \ldots, a_r)$ an ideal
generated by $r$ elements; then if $p$ is a minimal prime divisor of $I$ we have
$\htop p \le r$. Hence the height of a proper ideal of $A$ is always finite.
\end{theorem}

\begin{proof}
The ideal $IA_p \subset A_p$ is a primary ideal belonging to the maximal
ideal, so that $\htop p = \dim A_p = \delta(A_p) \le r$.
\end{proof}

\begin{remark}
Krull proved this theorem by induction on $r$; the case $r = 1$ is
then the hardest part of the proof. Krull called the $r = 1$ case the
\emph{principal ideal theorem} (Hauptidealsatz), and the whole of Theorem 5 is sometimes
known by this name. Here Theorem 5 is merely a corollary of Theorem 4,
but one can also deduce the statement $\dim M = \delta(M)$ of Theorem 4 from
it. As far as proving Theorem 5 is concerned, Krull's proof, which does
not use the Samuel function, is easier. For this proof, see [N1] or [K].
More elementary proofs of the principal ideal theorem can be found in
Rees [3] and Caruth [1].

The definition of height is abstract, and even when one can find a lower
bound, one cannot expect an upper bound just from the definition, so
that this theorem is extremely important. The principal ideal theorem
corresponds to the familiar and obvious-looking proposition of geo-
metrical and physical intuition (which is strictly speaking not always true)
that `adding one equation can decrease the dimension of the space of
solutions by at most one'.
\end{remark}

\begin{theorem}\label{thm:13.6}
Let $P$ be a prime ideal of height $r$ in a Noetherian ring $A$.
Then
\begin{enumerate}[(i)]
  \item $P$ is a minimal prime divisor of some ideal $(a_1, \ldots, a_r)$ generated by
        $r$ elements;
  \item if $b_1, \ldots, b_s \in P$ we have $\htop P/(b_1, \ldots, b_s) \ge r - s$;
  \item if $a_1, \ldots, a_r$ are as in (i) we have
        \[
          \htop P/(a_1, \ldots, a_i) = r - i \quad \text{for } 1 \le i \le r.
        \]
\end{enumerate}
\end{theorem}

\begin{proof}
(i) $A_P$ is an $r$-dimensional local ring, so that by Theorem 4 we can
choose $r$ elements $a_1, \ldots, a_r \in PA_P$ such that $(a_1, \ldots, a_r)A_P$ is $PA_P$-primary.
Each $a_i$ is of the form an element of $P$ times a unit of $A_P$, so that without
loss of generality we can assume that $a_i \in P$. Then $P$ is a minimal prime
divisor of $(a_1, \ldots, a_r)A$.

(ii) Set $\bar{A} = A/(b_1, \ldots, b_s)$, $\bar{P} = P/(b_1, \ldots, b_s)$ and $\htop P = t$. Then by (i),
there exist $c_1, \ldots, c_t \in P$ such that $P$ is a minimal prime divisor of
$(c_1, \ldots, c_t)A$. Then $\bar{P}$ is a minimal prime divisor of $(\bar{b}_1, \ldots, \bar{b}_s, \bar{c}_1, \ldots, \bar{c}_t)$,
and hence $r \le s + t$ by Theorem 5.

(iii) In $\bar{A} = A/(a_1, \ldots, a_i)$, the ideal $\bar{P} = P/(a_1, \ldots, a_i)$ is a minimal
prime divisor of $(\bar{a}_{i+1}, \ldots, \bar{a}_r)$, hence
\[
  \htop P/(a_1, \ldots, a_i) \le r - i.
\]
The opposite inequality was proved in (ii).
\end{proof}

\begin{theorem}\label{thm:13.7}
Let $A = \bigoplus_{n \ge 0} A_n$ be a Noetherian graded ring.
\begin{enumerate}[(i)]
  \item If $I$ is a homogeneous ideal and $P$ is a prime divisor of $I$ then $P$ is also
        homogeneous.
  \item If $P$ is a homogeneous prime ideal of height $r$ then there exists a
        sequence
        \[
          P = P_r \supset P_{r-1} \supset \cdots \supset P_0
        \]
        of length $r$ consisting of homogeneous
        prime ideals.
\end{enumerate}
\end{theorem}

\begin{proof}
(i) $P$ can be expressed in the form $P = \ann(x)$ for a suitable element
$x$ of the graded $A$-module $A/I$. Let $a \in P$, and let $x = x_0 + x_1 + \cdots + x_n$
and $a = a_p + a_{p+1} + \cdots + a_q$ be decompositions into homogeneous terms.
Then since $ax = 0$,
\[
  a_px_0 = 0, \quad a_px_1 + a_{p+1}x_0 = 0, \quad a_px_2 + a_{p+1}x_1 + a_{p+2}x_0 = 0, \ldots,
\]
from which we get $a_p^2x_0 = 0$, $a_p^2x_1 = 0, \ldots$, and finally $a_p^{n+1}x = 0$. It follows
that $a_p^{n+1} \in P$, but since $P$ is prime, $a_p \in P$. Thus $a_{p+1} + \cdots + a_q \in P$, so
that in turn $a_{p+1} \in P$. Proceeding in the same way, we see that all the homo-
geneous terms of $a$ are in $P$, so that $P$ is a homogeneous ideal.

(ii) First of all note that we can assume that $A$ is an integral domain.
To see this, if we take a chain $P = P_r \supset \cdots \supset P_0$ of prime ideals of length
$r$ then $P_0$ is a minimal prime divisor of $(0)$, and so by (i) is a homogeneous
ideal; so we can replace $A$ by $A/P_0$. Now choose a homogeneous element
$0 \neq b_r \in P$; then by Theorem 6, $\htop (P/b_rA) = r - 1$, and so there is a minimal
prime divisor $Q$ of $b_rA$ such that $\htop (P/Q) = r - 1$; since $Q \neq (0)$, it is a
height $1$ homogeneous prime ideal. By the inductive hypothesis on $r$
applied to $P/Q$ there exists a chain $P = P_r \supset P_{r-1} \supset \cdots \supset P_1 = Q$ of homo-
geneous prime ideals of length $r - 1$, and adding on $P_0 = (0)$ we get a chain
of length $r$.
\end{proof}

Let us investigate more closely the relation between local rings and
graded rings.

\begin{theorem}\label{thm:13.8}
Let $k$ be a field, and $R = k[\xi_1, \ldots, \xi_r]$ a graded ring
generated by elements $\xi_1, \ldots, \xi_r$ of degree $1$; set $M = \sum \xi_iR$, $A = R_M$, and
$m = MA_M$.
\begin{enumerate}[(i)]
  \item Let $\chi$ be the Samuel function of the local ring $A$, and $\varphi$ the Hilbert
        function of the graded ring $R$; then $\varphi(n) = \chi(n) - \chi(n - 1)$;
  \item $\dim R = \htop M = \dim A = \deg \varphi + 1$;
  \item $\gr_m(A) \cong R$ as graded rings.
\end{enumerate}
\end{theorem}

\begin{proof}
$M$ is a maximal ideal of $R$ so that
\[
  m^n/m^{n+1} \cong M^n/M^{n+1} \cong R_n;
\]
hence $\chi(n) - \chi(n - 1) = l(m^n/m^{n+1}) = l(R_n) = \varphi(n)$, and so $\dim A = \deg \chi = 1 +
\deg \varphi$. Then since $A = R_M$, we have $\dim A = \htop M$. After this it is enough
to prove that $\dim R = \htop M$. First of all, assume that $R$ is an integral
domain, so that by Example 3 in the section on Hilbert functions and by
Theorem~\ref{thm:5.6}, we have
\[
  1 + \deg \varphi \ge \trdeg_k R = \dim R \ge \htop M;
\]
putting this together with $\htop M = \dim A = 1 + \deg \varphi$, we get $\dim R = \htop M$.
Next for general $R$, let $P_1, \ldots, P_s$ be the minimal prime ideals of $R$; then by
Theorem 7, these are all homogeneous ideals, and each $R/P_i$ is a graded
ring. Choosing $P_1$ such that $\dim R = \dim R/P_1$ and using the above result,
we get
\[
  \dim R = \dim R/P_1 = \htop M/P_1 \le \htop M \le \dim R,
\]
so that $\dim R = \htop M$ as required. We have $R_n \subset M^n \subset m^n$ with $m^n/m^{n+1} \cong
R_n$, and so taking an element $x$ of $R_n$ into its image in $m^n/m^{n+1}$ we obtain
a canonical one-to-one map $R \cong \gr_m A$, and it is clear from the definition
that this is a ring isomorphism.
\end{proof}

\begin{theorem}\label{thm:13.9}
Let $(A, m, k)$ be a Noetherian local ring, and set $G = \gr_m A$;
then $\dim A = \dim G$.
\end{theorem}

\begin{proof}
Letting $\varphi$ be the Hilbert polynomial of $G$, we have $\dim A =
1 + \deg \varphi$ (by Theorem 4), and by the previous theorem this is equal to
$\dim G$.
\end{proof}

In fact the following more general theorem holds: for $I$ a proper ideal in a
Noetherian local ring $A$, set $G = \gr_I(A)$; then $\dim A = \dim G$. This will be
proved a little later (Theorem~\ref{thm:15.7}).

% ============ 习题 ============
\section*{Exercises to \S\ref{sec:13}}
\emph{Prove the following propositions.}

\begin{exercise}\label{ex:13.1-ex}
Let $R = R_0 + R_1 + \cdots$ be a graded ring, and $u$ a unit of $R_0$. Then the map $T_u$,
defined by $T_u(x_0 + x_1 + \cdots + x_n) = x_0 + x_1u + \cdots + x_nu^n$ (where $x_i \in R_i$) is
an automorphism of $R$. If $R_0$ contains an infinite field $k$, then an ideal $I$ of $R$
is homogeneous if and only if $T_u(I) = I$ for every $u \in k$.
\end{exercise}

\begin{exercise}\label{ex:13.2-ex}
Let $R = R_0 + R_1 + \cdots$ be a graded ring, $I$ an ideal of $R$ and $t$ an
indeterminate over $R$. Set $R' = R[t, t^{-1}]$ and consider $R'$ as a graded ring
where $t$ has degree $0$ (that is, $R'_n = R_n[t, t^{-1}]$). Then an ideal $I$ of $R$ is
homogeneous if and only if $T_u(IR') = IR'$.
\end{exercise}

\begin{exercise}\label{ex:13.3-ex}
Let $A$ be a Noetherian ring having an embedded associated prime. If $a \in A$
is a non-zero divisor satisfying $\bigcap_{n \ge 1} a^nA = (0)$, then $A/(a)$ also has an
embedded associated prime.
\end{exercise}

\begin{exercise}\label{ex:13.4-ex}
Let $R = \bigoplus_{n \in \Z} R_n$ be a $\Z$-graded ring. For an ideal $I$ of $R$, let $I^*$ denote the
greatest homogeneous ideal of $R$ contained in $I$, that is the ideal of $R$
generated by all the homogeneous elements of $I$.
\begin{enumerate}[(i)]
  \item If $P$ is prime so is $P^*$.
  \item If $P$ is a homogeneous prime ideal and $Q$ is a $P$-primary ideal then $Q^*$
        is again $P$-primary.
\end{enumerate}
\end{exercise}

\begin{exercise}\label{ex:13.5-ex}
Let $R$ be a $\Z$-graded integral domain; write $S$ for the multiplicative set
consisting of all non-zero homogeneous elements of $R$. Then $R_S$ is a graded
ring, and its component of degree $0$ is a field $(R_S)_0 = K$; if $R \neq R_S$ then $R_S
\cong K[X, X^{-1}]$, where the degree of $X$ is the greatest common divisor of the
degrees of elements of $S$.
\end{exercise}

\begin{exercise}\label{ex:13.6-ex}
Let $R$ be an $\N$-graded ring and $P$ an inhomogeneous prime ideal of $R$; then
there are no prime ideals contained between $P^*$ and $P$. If $\htop P < \infty$ then
$\htop P = \htop P^* + 1$ (Matijevic--Roberts [1]).
\end{exercise}

% ============ Appendix to §13. Determinantal ideals (原书 pp.101-103) ============
\section*{Appendix to \S\ref{sec:13}. Determinantal ideals}
\phantomsection
\addcontentsline{toc}{section}{Appendix to \S\ref{sec:13}. Determinantal ideals}

Let $M = (a_{ij})$ be an $r \times s$ matrix $(r \le s)$ with elements $a_{ij}$ in a Noetherian ring
$A$, and let $I_t$ be the ideal of $A$ generated by the $t \times t$ minors (that is
$t$-determinants) of $M$. When $t = r$ and $A$ is a polynomial ring
$k[X_1, \ldots, X_s]$ over a field $k$, Macaulay proved that all the prime divisors of
$I_t$ have height $\le s - r + 1$ ([Mac], p. 54). In his Ph.D. thesis, Eagon
generalised this result as follows: for an arbitrary Noetherian ring $A$, every
minimal prime divisor of $I_t$ has height $\le (r - t + 1)(s - t + 1)$. The following
ingenious proof is taken from Eagon--Northcott [1]. We begin with some
preliminary observations.

The following operations on a matrix $M$ with elements in a ring $A$ are
called \emph{elementary row operations}: (1) permutation of the rows; (2) replacing
$C_i$ by $uC_i + vC_j$, where $C_i$ and $C_j$ $(i \neq j)$ are two distinct rows of $M$, $u$ is a unit
of $A$ and $v$ is an element of $A$; elementary column operations are defined
similarly. The ideal $I_t$ does not change under these operations. Now, if an
element of $M$ is a unit in $A$, we can transform $M$ by a finite number of
elementary row and column operations to the following form:
\[
  \begin{pmatrix}
    1 & 0 & \cdots & 0 \\
    0 & & & \\
    \vdots & & N & \\
    0 & & &
  \end{pmatrix},
\]
and $I_t$ is equal to the ideal of $A$ generated by the $(t - 1) \times (t - 1)$ minors of $N$.

\begin{lemma}\label{lem:13.1}
Let $(A, P)$ be a Noetherian local ring and set $B = A[X]$. Let $J$ be a
$P$-primary ideal of $A$ and $J'$ an ideal of $B$ such that $J' \subset PB$ and
$J' + XB = JB + XB$. Then $PB$ is a minimal prime divisor of $J'$.
\end{lemma}

\begin{proof}
$PB + XB$ is a maximal ideal of $B$, and is the radical of $JB + XB =
J' + XB$. Thus, in the ring $B/J'$ we have that $(PB + XB)/J'$ is a minimal
prime divisor of the principal ideal $(J' + XB)/J'$. Hence $\htop((PB + XB)/
J') = 1$. Since $PB \subset J'$ is a prime ideal in $B/J'$, we have $\htop(PB/J') = 0$.
\end{proof}

\begin{theorem}[Eagon]\label{thm:13.10}
Let $A$ be a Noetherian ring and $M$ be an $r \times s$
matrix $(r \le s)$ of elements of $A$. Let $I_t$ be the ideal of $A$ generated by the $t \times t$
minors of $M$. If $P$ is a minimal prime divisor of $I_t$, then we have
\[
  \htop P \le (r - t + 1)(s - t + 1).
\]
\end{theorem}

\begin{proof}
Induction on $r$. When $r = 1$ we have $t = 1$, and so $(r - t + 1)(s - t +
1) = s$. The ideal $I_t$ is generated by $s$ elements, so that the assertion is just
the principal ideal theorem (Theorem 5) in this case. Next assume that $r > 1$.
Localising at $P$ we may assume that $A$ is a local ring with maximal ideal $P$,
and that $I_t$ is $P$-primary.

If $t = 1$, then $I_t$ is generated by $rs$ elements and $(r - t + 1)(s - t + 1) = rs$,
so our assertion holds also for this case. Therefore we assume $t > 1$. If at
least one of the elements of $M$ is a unit of $A$, then by what we said above, $I_t$ is
generated by $(t - 1) \times (t - 1)$ minors of a $(r - 1) \times (s - 1)$ matrix, and again
we are done. Therefore we assume that all the elements of $M$ are in $P$. Now
comes the brilliant idea. Let $M'$ be the matrix with elements in $B = A[X]$
obtained from $M$ by replacing $a_{11}$ by $a_{11} + X$, and let $I'$ be the ideal of $B$
generated by the $t \times t$ minors of $M'$. Since $t > 1$ and $a_{ij} \in P$ for all $i$ and $j$ we
have $I' \subset PB$. We also have $I' + XB = I_tB + XB$ since both sides have the
same image in $B/XB = A$. Therefore $PB$ is a minimal prime divisor of $I'$ by
the lemma. Since the element $a_{11} + X$ of $M'$ is not in $PB$, we have
$\htop PB \le (r - t + 1)(s - t + 1)$ by our previous argument. Since $\htop PB = \htop P$,
as we can see by Theorems 4 and 5, we are done.
\end{proof}

% ============ §14 Systems of parameters and multiplicity (原书 pp.104-110) ============
\section{Systems of parameters and multiplicity}\label{sec:14}\index{systems of parameters, section 14}\index{multiplicity, section 14}

Let $(A, m)$ be an $r$-dimensional Noetherian local ring; by Theorem~\ref{thm:13.4}, there exists an $m$-primary ideal generated by $r$ elements, but none
generated by fewer. If $a_1, \ldots, a_r \in m$ generate an $m$-primary ideal,
$\{a_1, \ldots, a_r\}$ is said to be a \emph{system of parameters}\index{system of parameters} of $A$ (sometimes
abbreviated to s.o.p.). If $M$ is a finite $A$-module with $\dim M = s$, there exist
$y_1, \ldots, y_s \in m$ such that $l(M/(y_1, \ldots, y_s)M) < \infty$, and then $\{y_1, \ldots, y_s\}$ is
said to be a \emph{system of parameters of $M$}.

If we set $A/m = k$, the smallest number of elements needed to generate
$m$ itself is equal to $\rank_k m/m^2$; (here $\rank_k$ is the rank of a free module
over $k$, that is the dimension of $m/m^2$ as $k$-vector space). This number is
called the \emph{embedding dimension}\index{embedding dimension} of $A$, and is written emb dim $A$. In general
\[
  \dim A \le \text{emb dim } A,
\]
and equality holds when $m$ can be generated by $r$ elements; in this case
$A$ is said to be a \emph{regular local ring}\index{regular local ring}, and a system of parameters generating
$m$ is called a \emph{regular system of parameters}.

\begin{theorem}\label{thm:14.1}
Let $(A, m)$ be a Noetherian local ring, and $x_1, \ldots, x_r$ a system
of parameters. Then
\begin{enumerate}[(i)]
  \item $\dim A/(x_1, \ldots, x_i) = r - i$ for $1 \le i \le r$.
  \item although it is not true that $\htop (x_1, \ldots, x_i) = i$ for all $i$ for an arbitrary
        system of parameters, there exists a choice of $x_1, \ldots, x_r$ such that every
        subset $F \subset \{x_1, \ldots, x_r\}$ generates an ideal of $A$ of height equal to the
        number of elements of $F$.
\end{enumerate}
\end{theorem}

\begin{proof}
(i) is contained in Theorem~\ref{thm:13.6}. We now prove the second half of
(ii). If $r \le 1$ the assertion is obvious; suppose that $r > 1$. Let $p_{0j}$ (for $1 \le
j \le e_0$) be the prime ideals of $A$ of height $0$. Choosing $x_1 \in m$ not contained
in any $p_{0j}$, we have $\htop (x_1) = 1$. Next letting $p_{1j}$ (for $1 \le j \le e_1$) be the minimal
prime divisors of $(x_1)$, so that $\htop p_{1j} = 1$, and choosing $x_2 \in m$ not contained
in any $p_{0j}$ or any $p_{1j}$, we have $\htop(x_1) = 1$, $\htop(x_1, x_2) = 2$; if $r = 2$ we're
done. If $r > 2$ we choose $x_3 \in m$ not contained in any minimal prime
divisor of $(0)$, $(x_1)$, $(x_2)$, $(x_1, x_2)$, and proceed in the same way to obtain
the result.

We now give an example where $\htop (x_1, \ldots, x_i) < i$. Let $k$ be a field and
set $R = k[X, Y, Z]$; let $I = (X) \cap (Y, Z)$, and write $A = R/I$, and $x, y, z$ for
the images in $A$ of $X, Y, Z$. The minimal prime ideals of $A$ are $(x)$ and
$(y, z)$; now $A/(x) = R/(X) \cong k[Y, Z]$ is two-dimensional and $A/(y, z) \cong
R/(Y, Z) \cong k[X]$ is one-dimensional, so that $\dim A = 2$. $\{y, x + z\}$ is a
system of parameters of $A$; in fact $xy = xz = 0$, so that $x^2 = x(x + z) \in
(y, x + z)$ and $z^2 = z(x + z) \in (y, x + z)$. However, $y$ is contained in the
minimal prime ideal $(y, z)$ of $A$, and hence $\htop (y) = 0$.
\end{proof}

\begin{theorem}\label{thm:14.2}
Let $(R, m)$ be an $n$-dimensional regular local ring, and
$x_1, \ldots, x_i$ elements of $m$. Then the following conditions are equivalent:
\begin{enumerate}[(1)]
  \item $x_1, \ldots, x_i$ is a subset of a regular system of parameters of $R$;
  \item the images in $m/m^2$ of $x_1, \ldots, x_i$ are linearly independent over $R/m$;
  \item $R/(x_1, \ldots, x_i)$ is an $(n - i)$-dimensional regular local ring.
\end{enumerate}
\end{theorem}

\begin{proof}
$(1) \Rightarrow (2)$ If $x_1, \ldots, x_i, x_{i+1}, \ldots, x_n$ is a regular system of parameters then
their images generate $m/m^2$ over $k = R/m$, and since $\rank_k m/m^2 = n$ they
must be linearly independent over $k$.

$(1) \Rightarrow (3)$ We know that $\dim R/(x_1, \ldots, x_i) = n - i$, and the images of
$x_{i+1}, \ldots, x_n$ generate the maximal ideal of $R/(x_1, \ldots, x_i)$.

$(3) \Rightarrow (1)$ If the maximal ideal $m/(x_1, \ldots, x_i)$ of $R/(x_1, \ldots, x_i)$ is generated
by the images of $y_1, \ldots, y_{n-i} \in m$ then $m$ is generated by $x_1, \ldots, x_i$,
$y_1, \ldots, y_{n-i}$.
\end{proof}

\begin{remark}
The hypothesis that $R$ is regular is not needed for $(3) \Rightarrow (1)$.
\end{remark}

\begin{proof}[Proof of $(2) \Rightarrow (1)$]
Using $\rank_k m/m^2 = n$, if we choose $x_{i+1}, \ldots, x_n \in m$ such that
the images of $x_1, \ldots, x_n$ in $m/m^2$ form a basis then $x_1, \ldots, x_n$ generate $m$
and so forms a regular system of parameters.
\end{proof}

\begin{theorem}\label{thm:14.3}
A regular local ring is an integral domain.
\end{theorem}

\begin{proof}
Let $(R, m)$ be an $n$-dimensional regular local ring; we proceed by
induction on $n$. If $n = 0$ then $m$ is an ideal generated by $0$ elements, so
that $m = (0)$. This in turn means that $R$ is a field. Thus a zero-dimensional
regular local ring is just a field by another name.

When $n = 1$, the maximal ideal $m = xR$ is principal and $\htop m = 1$, so that
there exists a prime ideal $p \neq m$ with $m \supset p$. If $y \in p$ we can write $y = xa$ with
$a \in R$, and since $x \notin p$ we have $a \in p$; hence $p = xp$, and by NAK, $p = (0)$. This
proves that $R$ is an integral domain. (There is a slightly different proof in the
course of the proof of Theorem~\ref{thm:11.2}; as proved there, a one-dimensional
regular local ring is just a DVR by another name.)

When $n > 1$, let $p_1, \ldots, p_l$ be the minimal prime ideals of $R$; then since
$m \nsubseteq m^2$ and $m \nsubseteq p_i$ for all $i$, there exists an element $x \in m$ not contained
in any of $m^2, p_1, \ldots, p_l$ (see Ex.~\ref{ex:1.6-ex}). Then the image of $x$ in $m/m^2$ is non-zero,
so that by the previous theorem $R/xR$ is an $(n - 1)$-dimensional regular
local ring. By the induction hypothesis, $R/xR$ is an integral domain, in
other words, $xR$ is a prime ideal of $R$. If $p_1$ is one of the minimal prime ideals
contained in $xR$ then since $x \notin p_1$, the same argument as in the $n = 1$ case
shows that $p_1 = xp_1$, and hence $p_1 = (0)$.
\end{proof}

\begin{theorem}\label{thm:14.4}
Let $(A, m, k)$ be a $d$-dimensional regular local ring; then
\[
  \gr_m(A) \cong k[X_1, \ldots, X_d],
\]
and if $\chi(n)$ is the Samuel function of $A$ then
\[
  \chi(n) = \binom{n+d}{d} \quad \text{for all } n \ge 0.
\]
\end{theorem}

\begin{proof}
Since $m$ is generated by $d$ elements, $\gr_m(A)$ is of the form
$k[X_1, \ldots, X_d]/I$, where $I$ is a homogeneous ideal. Now if $I \neq (0)$ let $S \in I$ be a
non-zero homogeneous element of degree $r$; then for $n \ge r$ the homogeneous
piece of $k[X]/I$ of degree $n$ has length at most
\[
  \binom{n+r-1}{d-1} - \binom{n-1}{d-1},
\]
which is a polynomial of degree $d - 2$ in $n$. This
implies that the Samuel function of $A$ is of degree at most $d - 1$, and
contradicts $\dim A = d$. Hence $I = (0)$; the second assertion follows from
the first.
\end{proof}

Let $(A, m)$ be a Noetherian local ring. Elements $y_1, \ldots, y_l \in m$ are said to be
\emph{analytically independent}\index{analytically independent elements} if they have the following property; for every
homogeneous form $F(Y_1, \ldots, Y_l)$ with coefficients in $A$,
\[
  F(y_1, \ldots, y_l) = 0 \Rightarrow \text{the coefficients of } F \text{ are in } m.
\]
If $y_1, \ldots, y_l$ are analytically independent and $A$ contains a field $k$, then
$F(y) \neq 0$ for any non-zero homogeneous form $F(Y) \in k[Y_1, \ldots, Y_l]$.

\begin{theorem}\label{thm:14.5}
Let $(A, m)$ be a $d$-dimensional Noetherian local ring and
$x_1, \ldots, x_d$ a system of parameters of $A$; then $x_1, \ldots, x_d$ are analytically
independent.
\end{theorem}

\begin{proof}
Set $q = \sum x_iA$. Since $q$ is an ideal of definition of $A$, by
Theorem~\ref{thm:13.4}, $\chi_q(n) = l(A/q^n)$ is a polynomial of degree $d$ in $n$ for $n \gg 0$. Set
$\bar{A} = k$; we say that a homogeneous form $\bar{f}(X) \in k[X_1, \ldots, X_d]$ of degree $n$
is a \emph{null-form} of $q$ if $F(x_1, \ldots, x_d) \in q^{n+1}m$ for any homogeneous form
$F(X) \in A[X_1, \ldots, X_d]$ which reduces to $\bar{f}(X)$ modulo $m$. Write $\mathfrak{n}$ for the ideal
of $k[X_1, \ldots, X_d]$ generated by the null-forms of $q$. Then
\[
  k[X_1, \ldots, X_d]/\mathfrak{n} = \bigoplus_{n \ge 0} q^n/q^{n+1}m,
\]
and writing $\varphi$ for the Hilbert polynomial of $k[X]/\mathfrak{n}$, we have $\varphi(n)
= l(q^n/q^{n+1}m)$ for $n \gg 0$. The right-hand side is just the number of elements in a
minimal basis of $q^n$, so that $\varphi(n) \cdot l(A/q) \ge l(q^n/q^{n+1})$. Now
\[
  l(q^n/q^{n+1}) = \chi_q(n) - \chi_q(n - 1)
\]
is a polynomial in $n$ of degree $d - 1$, so that $\deg \varphi \ge d - 1$, but if $\mathfrak{n} \neq (0)$ this
is impossible. Thus $\mathfrak{n} = (0)$, and the statement in the theorem follows at
once.
\end{proof}

\subsection*{Multiplicity}

Let $(A, m)$ be a $d$-dimensional Noetherian local ring, $M$ a finite $A$-module,
and $q$ an ideal of definition of $A$ (that is, an $m$-primary ideal). As we saw
in \S\ref{sec:13}, the Samuel function $l(M/q^{n+1}M) = \chi_M^q(n)$ can be expressed for
$n \gg 0$ as a polynomial in $n$ with rational coefficients, and degree equal to
$\dim M$, and therefore at most $d$. In addition, this polynomial can only take
integer values for $n \gg 0$, so it is easy to see by induction on $d$ (using the fact
that $\chi(n + 1) - \chi(n)$ has the same property) that
\[
  \chi_M^q(n) = \frac{e}{d!}n^d + \text{(terms of lower order)},
\]
with $e \in \Z$. This integer will be written $e(q, M)$. By definition we have the
following property.

\smallskip
\noindent\textbf{Formula 14.1.}\quad $\displaystyle e(q, M) = \lim_{n \to \infty} \frac{d!}{n^d}l(M/q^nM)$, and in particular, if $d = 0$ then
$e(q, M) = l(M)$.

\smallskip
\noindent From this we see easily the following:
\noindent\textbf{Formula 14.2.}\quad $e(q, M) > 0$ if $\dim M = d$, and $e(q, M) = 0$ if $\dim M < d$;

\smallskip
\noindent\textbf{Formula 14.3.}\quad $e(q^r, M) = e(q, M)r^d$;

\smallskip
\noindent\textbf{Formula 14.4.}\quad If $q$ and $q'$ are both $m$-primary ideals and $q \supset q'$ then
$e(q, M) \le e(q', M)$.

\smallskip
\noindent We set $e(q, A) = e(q)$, and define this to be the \emph{multiplicity} of $q$. In addition,
we will refer to the multiplicity $e(m)$ of the maximal ideal as the \emph{multiplicity of the
local ring} $A$, and sometimes write $e(A)$ for it. For example, if $A$ is a
regular local ring then by Theorem 4, we can see that $e(A) = 1$.

\begin{theorem}\label{thm:14.6}
Let $0 \to M' \to M \to M'' \to 0$ be an exact sequence of
finite $A$-modules. Then
\[
  e(q, M) = e(q, M') + e(q, M'').
\]
\end{theorem}

\begin{proof}
We view $M'$ as a submodule of $M$. Then
\[
  l(M/q^nM) = l(M''/q^nM'') + l(M'/M' \cap q^nM),
\]
and obviously $q^nM' \subset M' \cap q^nM$. On the other hand by Artin--Rees, there
exists $c > 0$ such that
\[
  M' \cap q^nM \subset q^{n-c}M' \quad \text{for all } n > c.
\]
Hence
\[
  l(M'/q^{n-c}M') \le l(M'/M' \cap q^nM) \le l(M'/q^nM').
\]
From this and Formula 14.1 it follows easily that
\[
  e(q, M) - e(q, M'') = \lim_{n \to \infty} \frac{d!}{n^d}l(M'/M' \cap q^nM) = e(q, M'). \qquad\blacksquare
\]
\end{proof}

\begin{theorem}\label{thm:14.7}
Let $\{p_1, \ldots, p_t\}$ be all the minimal prime ideals of $A$ such
that $\dim A/p_i = d$; then
\[
  e(q, M) = \sum_{i=1}^{t} l(M_{p_i})e(\bar{q}, A/p_i),
\]
where $\bar{q}_i$ denotes the image of $q$ in $A/p_i$, and $l(M_{p_i})$ stands for the length of $M_{p_i}$ as
$A_{p_i}$-module.
\end{theorem}

\begin{proof}[Proof (taken from Nagata [N1])]
We write $\sigma = \sum l(M_{p_i})$ and proceed by
induction on $\sigma$. If $\sigma = 0$ then $\dim M < d$, so that the left-hand side is $0$, and
the right-hand side is obviously $0$; now suppose $\sigma > 0$. Now there is some
$p \in \{p_1, \ldots, p_t\}$ for which $M_p \neq 0$; then $p$ is a minimal element of
$\Supp(M)$. Hence $p \in \Ass(M)$, that is $M$ contains a submodule $N$ isomorphic
to $A/p$. Then
\[
  e(q, M) = e(q, N) + e(q, M/N).
\]
On the other hand, $N_{p_i} \cong A_p/pA_p$, and $N_{p_i} = 0$ for $p_i \neq p$, so that $l(N_{p_i}) = 1$,
and the value of $\sigma$ for $M/N$ has decreased by one, so that the theorem holds
for $M/N$. However, from the definition
\[
  e(q, N) = e(q, A/p) = e(\bar{q}, A/p), \quad \text{where } \bar{q} = (q + p)/p.
\]
Putting this together, we see that the theorem also holds for $M$.
\end{proof}

Theorem 7 allows us to reduce the study of $e(q, M)$ to the case that $A$ is an
integral domain and $M = A$. In particular, if $A$ is an integral domain then
$l(M_{p_i})$ is just the rank of $M$, so that we obtain the following theorem.

\begin{theorem}\label{thm:14.8}
Let $A$ be a Noetherian local integral domain, $q$ an ideal of
definition of $A$ and $M$ a finite $A$-module; then
\[
  e(q, M) = e(q) \cdot s, \quad \text{where } s = \rank M.
\]
\end{theorem}

\begin{theorem}\label{thm:14.9}
Let $(A, m)$ be a Noetherian local ring, $q$ an ideal of definition
of $A$, and $x_1, \ldots, x_d$ a system of parameters of $A$ contained in $q$. Suppose that
$x_i \in q^{\nu_i}$ for $1 \le i \le d$. Then for a finite $A$-module $M$ and $s = 1, \ldots, d$ we have
\[
  e(q/(x_1, \ldots, x_s), M/(x_1, \ldots, x_s)M) \ge \nu_1\nu_2 \cdots \nu_s e(q, M).
\]
In particular if $s = d$, we have
\[
  l(M/(x_1, \ldots, x_d)M) \ge \nu_1\nu_2 \cdots \nu_de(q, M).
\]
\end{theorem}

\begin{proof}
It is enough to prove the case $s = 1$. We set $A' = A/x_1A$, $q' = q/x_1A$,
$M' = M/x_1M$ and $\nu = \nu_1$. By Theorem 1, we have $\dim A' = d - 1$. On the
other hand,
\[
  l(M'/q'^nM') = l(M/x_1M + q^nM)
\]
\[
  = l(M/q^nM) - l(x_1M + q^nM/q^nM).
\]
In addition, in view of $(x_1M + q^nM)/q^nM \cong x_1M/x_1M \cap q^nM \cong M/
(q^nM:x_1)$ and $q^{n-\nu}M \subset q^nM:x_1$, we have
\[
  l(M'/q'^nM') \ge l(M/q^nM) - l(M/q^{n-\nu}M),
\]
and therefore
\[
  l(M'/q'^nM') \ge l(M/q^nM) - l(M/q^{n-\nu}M).
\]
When $n \gg 0$ the right-hand side is of the form
\[
  \frac{e(q, M)}{d!}[n^d - (n - \nu)^d] + \text{(polynomial of degree } d - 2 \text{ in } n)
\]
\[
  = \frac{e(q, M)\nu}{(d-1)!}n^{d-1} + \text{(polynomial of degree } d - 2 \text{ in } n),
\]
so that the assertion is clear.
\end{proof}

A case of the above theorem which is particularly simple, but important,
is the following.

\begin{theorem}\label{thm:14.10}
Let $(A, m)$ be a $d$-dimensional Noetherian local ring, let
$x_1, \ldots, x_d$ be a system of parameters of $A$, and set $q = (x_1, \ldots, x_d)$; then
\[
  l(A/q) \ge e(q) \ge e(m).
\]
and if in addition $x_i \in m^{\nu_i}$ for all $i$ then $l(A/q) \ge \nu_1 \cdots \nu_d e(m)$.
\end{theorem}

\begin{theorem}\label{thm:14.11}
Let $A, m, x_i$ and $q$ be as above. Let $M$ be a finite $A$-module,
and set $A' = A/x_1A$, $M' = M/x_1M$ and $q' = q/x_1A = \sum_{i \ge 2} x_iA'$. Then if $x_1$ is a
non-zero-divisor of $M$, we have the following equality
\[
  e(q', M') = e(q, M).
\]
\end{theorem}

\begin{proof}
Since $l(M'/q'^{n+1}M') = l(M/x_1M + q^{n+1}M)$ we have
\[
  l(M/q^{n+1}M) - l(M'/q'^{n+1}M') = l(x_1M + q^{n+1}M/q^{n+1}M)
\]
\[
  = l(x_1M/x_1M \cap q^{n+1}M) = l(M/(q^{n+1}M:x_1))
\]
\[
  = l(M/q^nM) - l((q^{n+1}M:x_1)/q^nM).
\]
On the other hand, setting $\mathfrak{a} = \sum_{i \ge 2} x_iA$ we have $q = x_1A + \mathfrak{a}$ and $q^{n+1} =
x_1q^n + \mathfrak{a}^{n+1}$, and therefore
\[
  q^{n+1}M:x_1 = q^nM + (\mathfrak{a}^{n+1}M:x_1).
\]
Moreover, by Artin--Rees, there is a $c > 0$ such that for $n > c$ we have
$\mathfrak{a}^{n+1}M \cap x_1M = \mathfrak{a}^{n-c}(\mathfrak{a}^{c+1}M \cap x_1M)$, and therefore $\mathfrak{a}^{n+1}M:x_1 \subset \mathfrak{a}^{n-c}M$.
Thus
\[
  (q^{n+1}M:x_1)/q^nM = (q^nM + (\mathfrak{a}^{n+1}M:x_1))/q^nM
\]
\[
  \subset (q^nM + \mathfrak{a}^{n-c}M)/q^nM
\]
\[
  \cong \mathfrak{a}^{n-c}M/\mathfrak{a}^{n-c}M \cap q^nM.
\]
Now $\mathfrak{a}^{n-c}M/\mathfrak{a}^{n-c}M \cap q^nM$ is a module over $A/q'$, and since $\mathfrak{a}$ is generated
by $d - 1$ elements, $\mathfrak{a}^{n-c}$ is generated by $\binom{n-c+d-2}{d-2}$ elements. Thus
for $n > c$ we have
\[
  l(M'/q'^{n+1}M') \ge l(M/q^{n+1}M) - m\binom{n-c+d-2}{d-2},
\]
where $m$ is the number of generators of $M$. The right-hand side is a
polynomial of degree $d - 2$ in $n$, so that
\[
  e(q', M') = (d - 1)! \lim_{n \to \infty} l(M'/q'^{n+1}M')/n^{d-1}
\]
\[
  = (d - 1)! \lim_{n \to \infty} [l(M/q^{n+1}M) - l(M/q^nM)]/n^{d-1}
\]
\[
  = e(q, M).
\]
\end{proof}

\begin{theorem}[Lech's lemma]\label{thm:14.12}
Let $A$ be a $d$-dimensional Noetherian
local ring, and $x_1, \ldots, x_d$ a system of parameters of $A$; set $q = (x_1, \ldots, x_d)$, and
suppose that $M$ is a finite $A$-module. Then
\[
  e(q, M) = \lim_{\min(\nu_i) \to \infty} \frac{l(M/(x_1^{\nu_1}, \ldots, x_d^{\nu_d})M)}{\nu_1 \cdots \nu_d}.
\]
\end{theorem}

\begin{proof}
If $d = 0$ then both sides are equal to $l(M)$. If $d = 1$ then the right-hand
side is exactly Formula 14.1 which defines $e(q, M)$. For $d > 1$ we use
induction on $d$.

Setting $N_\nu = \{m \in M \mid x_1^\nu m = 0\}$ we have $N_1 \subset N_2 \subset \cdots$ so that there is a
$c > 0$ such that $N_\nu = N_{\nu+1} = \cdots$. If we set $M' = x_1^cM$ then $x_1$ is a non-zero-
divisor for $M'$, and there is an exact sequence $0 \to N_\nu \to M \to M' \to 0$.
Since $N_\nu$ is a module over $A/x_1^cA$ we have $\dim N_\nu < d$, and therefore
$e(q, N_\nu) = e(q, M')$. On the other hand,
\[
  l(M/(x_1^{\nu_1}, \ldots, x_d^{\nu_d})M) - l(M'/(x_1^{\nu_1}, \ldots, x_d^{\nu_d})M')
\]
\[
  = l(N_\nu + (x_1^{\nu_1}, \ldots, x_d^{\nu_d})M/(x_1^{\nu_1}, \ldots, x_d^{\nu_d})M)
\]
\[
  = l(N_\nu/N_\nu \cap (x_1^{\nu_1}, \ldots, x_d^{\nu_d})M)
\]
\[
  \le l(N_\nu/(x_1^{\nu_1}, \ldots, x_d^{\nu_d})N_\nu).
\]
If $\nu_1 \ge c$ then $x_1^{\nu_1}N_\nu = 0$, and $N_\nu$ is a module over the $(d - 1)$-dimensional
local ring $A/x_1^cA$, so that by induction there is a constant $C$ such that as
$\min(\nu_i) \to \infty$ we have
\[
  l(N_\nu/(x_1^{\nu_1}, \ldots, x_d^{\nu_d})N_\nu) = l(N_\nu/(x_2^{\nu_2}, \ldots, x_d^{\nu_d})N_\nu) \le C\nu_2 \cdots \nu_d.
\]
Therefore,
\[
  \lim [l(M/(x_1^{\nu_1}, \ldots, x_d^{\nu_d})M) - l(M'/(x_1^{\nu_1}, \ldots, x_d^{\nu_d})M')]/\nu_1 \cdots \nu_d = 0.
\]
This means that we can replace $M$ by $M'$ in the theorem, and so we can
assume that $x_1$ is a non-zero-divisor in $M$. Then by the previous theorem we
have $e(q, M) = e(\bar{q}, \bar{M})$, with $\bar{q} = q/x_1A$ and $\bar{M} = M/x_1M$. If we furthermore set
\[
  E = (x_2^{\nu_2}, \ldots, x_d^{\nu_d})M \quad \text{and} \quad F = M/E
\]
then by Theorem 9, we have
\[
  e(\bar{q}, \bar{M}) \cdot \nu_2 \cdots \nu_d \le l(M/(x_1^{\nu_1}, \ldots, x_d^{\nu_d})M) = l(F/x_1^{\nu_1}F)
\]
\[
  = \sum_{i=1}^{\nu_1} l(x_1^{i-1}F/x_1^iF) \le \nu_1l(F/x_1F) = \nu_1l(M/(x_1, x_2^{\nu_2}, \ldots, x_d^{\nu_d})M).
\]
Then by induction on $d$ we have
\[
  \lim l(M/(x_1^{\nu_1}, \ldots, x_d^{\nu_d})M)/\nu_1 \cdots \nu_d = \lim l(\bar{M}/(\bar{x}_2^{\nu_2}, \ldots, \bar{x}_d^{\nu_d})\bar{M})/\nu_2 \cdots \nu_d
  = e(\bar{q}, \bar{M}). \qquad\blacksquare
\]
\end{proof}

Although we will not use it in this book, we state here without proof a
remarkable result of Serre which shows that multiplicity can be expressed
as the Euler characteristic of the homology groups of the Koszul complex
(discussed in \S\ref{sec:16}).

\begin{theorem*}
Let $A$ be a $d$-dimensional Noetherian local ring, and $x_1, \ldots, x_d$ a
system of parameters of $A$; set $q = (x_1, \ldots, x_d)$ and let $M$ be a finite $A$-
module. Then
\[
  e(q, M) = \sum_i (-1)^i l(H_i(x, M)).
\]
For a proof, see for example Auslander and Buchsbaum [2].
\end{theorem*}

As we have seen in several of the above theorems, the multiplicity of
ideals generated by systems of parameters enjoy various nice properties. We
are now going to see that in a certain sense the general case can be reduced
to this one. We follow the method of Northcott and Rees [1].

More generally, let $A$ be a ring and $\mathfrak{a}$ an ideal. We say that an ideal $\mathfrak{b}$ is a
\emph{reduction}\index{reduction of an ideal} of $\mathfrak{a}$ if it satisfies the following condition:
\[
  \mathfrak{b} \subset \mathfrak{a}, \quad \text{and for some } r \ge 0 \text{ we have } \mathfrak{a}^{r+1} = \mathfrak{b}\mathfrak{a}^r.
\]
If $\mathfrak{b}$ is a reduction of $\mathfrak{a}$ and $\mathfrak{a}^{r+1} = \mathfrak{b}\mathfrak{a}^r$ then for any $n \ge 0$ we have $\mathfrak{a}^{r+n} = \mathfrak{b}^n\mathfrak{a}^r$.

\begin{theorem}\label{thm:14.13}
Let $(A, m)$ be a Noetherian local ring, $q$ an $m$-primary ideal
and $\mathfrak{b}$ a reduction of $q$; then $\mathfrak{b}$ is also $m$-primary, and for any finite $A$-module
$M$ we have
\[
  e(q, M) = e(\mathfrak{b}, M).
\]
\end{theorem}

\begin{proof}
If $q^{r+1} = \mathfrak{b}q^r$ then $q^{r+1} \subset \mathfrak{b} \subset q$, hence $\mathfrak{b}$ is also $m$-primary.
Moreover,
\[
  l(M/\mathfrak{b}^{n+1}M) \ge l(M/q^{n+1}M) = l(M/\mathfrak{b}^nq^r) \ge l(M/\mathfrak{b}^{n+r}M),
\]
so that $e(q, M) = e(\mathfrak{b}, M)$ follows easily.
\end{proof}

\begin{theorem}\label{thm:14.14}
Let $(A, m)$ be a $d$-dimensional Noetherian local ring, and
suppose that $A/m$ is an infinite field; let $q = (u_1, \ldots, u_s)$ be an $m$-primary
ideal. Then if $y_i = \sum a_{ij}u_j$ for $1 \le i \le d$ are $d$ `sufficiently general' linear
combinations of $u_1, \ldots, u_s$, the ideal $\mathfrak{b} = (y_1, \ldots, y_d)$ is a reduction of $q$ and
$\{y_1, \ldots, y_d\}$ is a system of parameters of $A$.
\end{theorem}

\begin{proof}
If $d = 0$ then $q^r = (0)$ for some $r > 0$, hence $(0)$ is a reduction of $q$
so that the result holds. We suppose below that $d > 0$.

\textbf{Step 1.} Set $A/m = k$ and consider the polynomial ring $k[X_1, \ldots, X_s]$ (or
$k[X]$ for short). For a homogeneous form $\varphi(X) = \varphi(X_1, \ldots, X_s) \in A[X]$
of degree $n$, we write $\bar{\varphi}(X) \in k[X]$ for the polynomial obtained by
reducing the coefficients of $\varphi$ modulo $m$. As in the proof of Theorem 5 we say that
$\bar{\varphi}(X) \in k[X]$ is a null-form of $q$ if $\varphi(u_1, \ldots, u_s) \in q^{n+1}m$; this notion depends
not just on $q$, but also on $u_1, \ldots, u_s$. However, for fixed $\bar{\varphi}$ it does not
depend on the choice of $\varphi$. We write $Q$ for the ideal of $k[X]$ generated by
all the null-forms of $q$, and call $Q$ the \emph{ideal of null-forms} of $q$. One sees
easily that all the homogeneous elements of $Q$ are null-forms of $q$, and
that the graded ring $k[X]/Q$ has graded component of degree $n$ isomorphic
to $q^n/q^{n+1}m$, so that we have
\[
  k[X]/Q = \bigoplus_{n \ge 0} q^n/q^{n+1}m = \gr_q(A) \otimes_A k.
\]
Write $\varphi(n)$ for the Hilbert function of $k[X]/Q$; then
\[
  l(q^n/q^{n+1}) \le \varphi(n) = l(q^n/q^{n+1}m) \le l(A/q) \cdot l(q^n/q^{n+1})
\]
(see the proof of Theorem 5). We know that for $n \gg 0$, the function $l(q^n/q^{n+1})$
is a polynomial in $n$ of degree $d - 1$ (where $d = \dim A$). Thus from the
above inequality, $\varphi$ is also a polynomial of degree $d - 1$, so that by
Theorem~\ref{thm:13.8}, (ii), we have $\dim k[X]/Q = d$.

Now set $V = \bigoplus kX_i$, and let $P_1, \ldots, P_t$ be the minimal prime divisors of
$Q$. By the assumption that $d > 0$, we have $P_i \nsupseteq V$, so that $P_i \cap V$ is a
proper vector subspace of $V$. Since $k$ is an infinite field,
\[
  V \neq \bigcup_{i=1}^{t} (V \cap P_i).
\]
Hence we can take a linear form $l_1(X) \in V$ not belonging to any $P_i$. If
$d \ge 1$ then similarly we can take $l_2(X) \in V$ such that $l_2(X)$ is not contained
in any minimal prime divisor of $(Q, l_1(X))$, and, proceeding in the same
way, we get $l_1(X), \ldots, l_d(X) \in V$ such that $(Q, l_1, \ldots, l_d)$ is a primary ideal
belonging to $(X_1, \ldots, X_s)$.

\textbf{Step 2.} We let $\mathfrak{b}$ be the ideal of $A$ generated by linear combinations
\[
  L_i(u) = \sum_{j=1}^{s} a_{ij}u_j \quad \text{(for } 1 \le i \le d\text{)}
\]
of $u_1, \ldots, u_s$ with coefficients in $A$. Then if we
put $l_i(X) = \sum a_{ij}X_j$, a necessary and sufficient condition for $\mathfrak{b}$ to be
a reduction of $q$ is that the ideal $(Q, l_1, \ldots, l_d)$ of $k[X]$ is $(X_1, \ldots, X_s)$-
primary.

\emph{Proof of necessity.} Suppose that $\mathfrak{b}q^r = q^{r+1}$. Then if $M(X)$ is a
monomial of degree $r + 1$ in $X_1, \ldots, X_s$, we can write
\[
  M(u) - \sum_i L_i(u)F_i(u) \in q^{r+2}m
\]
where the $F_i(X)$ are homogeneous forms of degrees $r$ with coefficients in
$A$. Thus
\[
  \bar{M}(X) - \sum_i l_i(X)\bar{F}_i(X) \in Q.
\]
Hence
\[
  (X_1, \ldots, X_s)^{r+1} \subset (Q, l_1, \ldots, l_d).
\]

\emph{Proof of sufficiency.} We go through the same argument in reverse: if
$\bar{M}(X) - \sum l_i\bar{F}_i \in Q$ then
\[
  M(u) - \sum L_i(u)F_i(u) \in q^{r+2}m,
\]
so that $q^{r+1} \subset \mathfrak{b}q^r + q^{r+2}m$; thus by NAK, $q^{r+1} = \mathfrak{b}q^r$.

\textbf{Step 3.} Putting together Steps 1 and 2 we see that $q$ has a reduction
$\mathfrak{b} = (y_1, \ldots, y_d)$ generated by $d$ elements. Both $q$ and its reduction $\mathfrak{b}$ are
$m$-primary ideals, so that $y_1, \ldots, y_d$ is a system of parameters of $A$. We
are going to prove that there exists a finite number of polynomials $D_\alpha(Z_{ij})$
for $1 \le \alpha \le \mu$ in $sd$ indeterminates $Z_{ij}$ (for $1 \le i \le d$ and $1 \le j \le s$) such that
$d$ linear combinations $y_i = \sum a_{ij}u_j$ (for $1 \le i \le d$) generate a reduction ideal
of $q$ if and only if at least one of $D_\alpha(a_{ij}) \neq 0$. (The expression `$d$ sufficiently
general linear combinations' in the statement of the theorem is quite vague,
but in the present case it has a precise interpretation as above.)

Let $G_1(X), \ldots, G_t(X)$ be a set of generators of $Q$, with $G_j$ homogeneous
of degree $e_j$. For any $sd$ elements $a_{ij}$ of $k$ (for $1 \le i \le d$ and $1 \le j \le s$), set
$l_i(X) = \sum a_{ij}X_j$. We write $I_n$ for the homogeneous component of degree $n$
of a homogeneous ideal $I \subset k[X_1, \ldots, X_s]$, and in particular we write
$(X_1, \ldots, X_s)_n = V_n$, so that
\[
  (X_1, \ldots, X_s)_n \subset (Q, l_1, \ldots, l_d) \iff V_n = (Q, l_1, \ldots, l_d)_n.
\]
Set $c_n = \dim_k V_n$. We have
\[
  (Q, l_1, \ldots, l_d)_n = \left\{\sum l_iF_i + \sum G_jH_j \mid F_i \in V_{n-1} \text{ and } H_j \in V_{n-e_j}\right\}.
\]
Let $K_1, \ldots, K_w$ be the elements obtained as $\sum l_iF_i + \sum G_jH_j$ as the $F_i$ run
through a basis of $V_{n-1}$ and the $H_j$ run independently through a basis
of $V_{n-e_j}$; it is clear that they span $(Q, l_1, \ldots, l_d)_n$. Each of $K_1, \ldots, K_w$ is a
linear combination of the $c_n$ monomials of degree $n$ in the $X_i$, with linear
functions in the $a_{ij}$ as coefficients; we write out these coefficients in a
$c_n \times w$ matrix. If $\psi_{\nu n}(a_{ij})$ for $1 \le \nu \le \rho$ are the $c_n \times c_n$ minors of this matrix
then the necessary and sufficient condition for $(X_1, \ldots, X_s)_n \subset (Q, l_1, \ldots, l_d)$ to
hold is that at least one of the $\psi_{\nu n}(a_{ij})$ is non-zero. Therefore the ideal
$(Q, l_1, \ldots, l_d)$ will fail to be $(X_1, \ldots, X_s)$-primary if and only if the
quantities $a_{ij}$ satisfy $\psi_{\nu n}(a_{ij}) = 0$ for all $n$ and all $\nu$. However, the ring $k[Z_{ij}]$
is Noetherian, so that the ideal of $k[Z_{ij}]$ generated by all of the $\psi_{\nu n}(Z_{ij})$ is
generated by finitely many elements $D_\alpha(Z_{ij})$ for $1 \le \alpha \le \mu$. These $D_\alpha$ clearly
meet our requirements.
\end{proof}

\begin{remark}
The polynomials $D_\alpha(Z_{ij})$ obtained above are in fact the necessary
and sufficient conditions on the coefficients $a_{ij}$ for the system of homo-
geneous equations $l_1(X) = \cdots = l_d(X) = G_1(X) = \cdots = G_t(X) = 0$ to have a
non-trivial solution, and as such they are known as a \emph{system of resultants}.
Here we have avoided appealing to the classical theory of resultants by
following a method given in Shafarevich [Sh].
\end{remark}

If $k = A/m$ is a finite field then Theorem 14 cannot to be used as it
stands, but we can use the following trick. Let $x$ be an indeterminate over
$A$, and set $S = A[x] - m[x]$; then $S$ consists of polynomials having a unit
of $A$ among their coefficients, and so the composite of the canonical maps
$A \to A[x] \to A[x]_S$ is injective. (In fact $S$ does not contain any zero-
divisors of $A[x]$, so that $A \subset A[x] \subset A[x]_S$; for this see [AM], Chap. 1,
Ex. 2.) Following Nagata~\cite{N1} we write $A(x)$ for $A[x]_S$. This is a
Noetherian local ring containing $A$, with maximal ideal $mA(x)$, and the
residue class field $A(x)/mA(x)$ is the field of fractions of $A[x]/m[x] = k[x]$,
that is, the field $k(x)$ of rational functions over $k$; this is an infinite field.
If $q$ is an $m$-primary ideal of $A$ then $qA(x)$ is a primary ideal belonging
to $mA(x)$. Moreover, since $A(x)$ is flat over $A$, we see that quite generally
if $I \supset I'$ are ideals of $A$ such that $I/I' \cong k$, then
\[
  IA(x)/I'A(x) \cong (I/I') \otimes_A A(x) \cong k \otimes_A A(x) = A(x)/mA(x).
\]
This gives $l_A(A/q^n) = l_{A(x)}(A(x)/q^nA(x))$, so that
\[
  \dim A = \dim A(x) \quad \text{and} \quad e(q) = e(qA(x)).
\]
Thus there are many instances when we can discuss properties of
$e(q)$ in terms of $A(x)$, to which Theorem 14 applies.

% ============ 习题 ============
\section*{Exercises to \S\ref{sec:14}}
\emph{Prove the following propositions.}

\begin{exercise}\label{ex:14.1-ex}
Let $(A, m)$ be a Noetherian local ring and set $G = \gr_m(A)$.
\begin{enumerate}[(i)]
  \item If $G$ is an integral domain then so is $A$ (hence Theorem 3 also follows
        from Theorem 4).
  \item Let $k$ be a field, and $A = k[[X, Y]]/(Y^2 - X^3)$; then $A$ is an integral
        domain, but $G$ has nilpotents.
\end{enumerate}
\end{exercise}

\begin{exercise}\label{ex:14.2-ex}
Let $(A, m)$ and $G$ be as above. For $a \in A$, suppose that $a \in m^i$ but $a \notin m^{i+1}$,
and write $a^*$ for the image of $a$ in $m^i/m^{i+1}$, viewed as an element of $G$;
define $a^*$ to be the \emph{leading term} of $a$. Set $0^* = 0$. Then
\begin{enumerate}[(i)]
  \item if $a^*b^* \neq 0$ then $(ab)^* = a^*b^*$;
  \item if $a^*$ and $b^*$ have the same degree and $a^* + b^* \neq 0$ then $(a + b)^* =
        a^* + b^*$;
  \item let $I \subset m$ be an ideal of $A$. Write $I^*$ for the ideal of $G$ generated by all
        the leading terms of elements of $I$; then setting $B = A/I$ and $n = m/I$, we
        have $\gr_n(B) = G/I^*$.
\end{enumerate}
\end{exercise}

\begin{exercise}\label{ex:14.3-ex}
In the above notation, if $G$ is an integral domain and $I = aA$ then
$I^* = a^*G$. If $I = (a_1, \ldots, a_r)$ with $r > 1$ then it can happen that
$I^* \neq (a_1^*, \ldots, a_r^*)$. Construct an example.
\end{exercise}

\begin{exercise}\label{ex:14.4-ex}
Let $(A, m)$ be a regular local ring, and $K$ its field of fractions.
\begin{enumerate}[(i)]
  \item For $0 \neq a \in A$, set $v(a) = i$ if $a \in m^i$ but $a \notin m^{i+1}$; then $v$ extends to an
        additive valuation of $K$.
  \item Let $R$ be the valuation ring of $v$; then $R$ is a DVR of $K$ dominating $A$.
        Let $x_1, \ldots, x_d$ be a regular system of parameters of $A$, and set
        $B = A[x_2/x_1, \ldots, x_d/x_1]$ and $P = x_1B$; then $P$ is a prime ideal of $B$ and
        $R = B_P$.
\end{enumerate}
\end{exercise}

\begin{exercise}\label{ex:14.5-ex}
In the above notation, if $0 \neq f \in m$ then $v(f)$ is equal to the multiplicity of
$A/(f)$.
\end{exercise}

\begin{exercise}\label{ex:14.6-ex}
(Associativity formula for multiplicities.) Let $A$ be a $d$-dimensional Noether-
ian local ring, $x_1, \ldots, x_d$ a system of parameters of $A$, $q = (x_1, \ldots, x_d)$, and
for $s \le d$ let $\mathfrak{a} = (x_1, \ldots, x_s)$. Write $\Gamma$ for the set of all prime divisors of $\mathfrak{a}$
satisfying $\htop p = s$, $\coht p = d - s$. Let $M$ be a finite $A$-module. Use Lech's
lemma to prove the following formula:
\[
  e(q, M) = \sum_{p \in \Gamma} e(q + p/p, M_p),
\]
(in particular, it follows that $\Gamma \neq \emptyset$).
\end{exercise}

\begin{remark}
The name of the formula comes from its connection with the associativity
of intersection product in algebraic geometry. For details, see [S3], pp. 84--5.
\end{remark}

\begin{exercise}\label{ex:14.7-ex}
Let $(A, m)$ be an $n$-dimensional Noetherian local integral domain, with
$n > 1$. If $0 \neq f \in A$ then $A_f$ is a Jacobson ring (see p. 34).
\end{exercise}

% ============ §15 The dimension of extension rings (原书 pp.116-122) ============
\section{The dimension of extension rings}\label{sec:15}\index{dimension of extension rings, section 15}

\subsection*{1. Fibres}

Let $\varphi\colon A \to B$ be a ring homomorphism, and for $p \in \Spec A$, write $\kappa(p) =
A_p/pA_p$; then $\Spec(B \otimes \kappa(p))$ is called the \emph{fibre}\index{fibre} of $\varphi$ over $p$. As we
saw in \S\ref{sec:7}, it can be identified with the inverse image in $\Spec B$ of $p$ under the
map ${}^a\varphi\colon \Spec B \to \Spec A$ induced by $\varphi$. The ring $B \otimes \kappa(p)$ will be called the
\emph{fibre ring} over $p$. When $(A, m)$ is a local ring, $m$ is the unique closed point of
$\Spec A$, and so the spectrum of $B \otimes \kappa(m) = B/mB$ is called the \emph{closed fibre} of
$\varphi$. If $A$ is an integral domain and $K$ its field of fractions then the spectrum of
$B \otimes_A K = B \otimes_A \kappa(0)$ is called the \emph{generic fibre} of $\varphi$.

\begin{theorem}\label{thm:15.1}
Let $\varphi\colon A \to B$ be a homomorphism of Noetherian rings,
and $P$ a prime ideal of $B$; then setting $p = P \cap A$, we have
\begin{enumerate}[(i)]
  \item $\htop P \le \htop p + \dim B_P/pB_P$;
  \item if $\varphi$ is flat, or more generally if the going-down theorem holds between
        $A$ and $B$, then equality holds in (i).
\end{enumerate}
\end{theorem}

\begin{proof}
We can replace $A$ and $B$ by $A_p$ and $B_P$, and assume that $(A, m)$ and
$(B, n)$ are local rings, with $mB \subset n$. Rewriting (i) in the form
\[
  \dim B \le \dim A + \dim B/mB
\]
makes clear the geometrical content. To prove this, take a system of
parameters $x_1, \ldots, x_r$ of $A$, and choose $y_1, \ldots, y_s \in B$ such that their images
in $B/mB$ form a system of parameters of $B/mB$. Then for $\nu, \rho$ large enough
we have $n^\nu \subset mB + \sum y_iB$ and $m^\rho \subset \sum x_jA$, giving $n^{\nu\rho} \subset \sum y_iB + \sum x_jB$.
Hence $\dim B \le r + s$.

(ii) Let $\dim B/mB = s$, and let $n = P_0 \supset P_1 \supset \cdots \supset P_s$ be a strictly decreas-
ing chain of prime ideals of $B$ between $n$ and $mB$. Obviously we have
$P_i \cap A = m$ for $0 \le i \le s$. Now set $\dim A = r$ and let $m = p_0 \supset p_1 \supset \cdots \supset p_r$
be a strictly decreasing chain of prime ideals of $A$; by the going-down
theorem, we can construct a strictly decreasing chain of prime ideals of $B$
\[
  P_s \supset P_{s+1} \supset \cdots \supset P_{s+r} \quad \text{such that } P_{s+i} \cap A = p_i.
\]
Thus $\dim B \ge r + s$, and putting this together with (i) gives equality.
\end{proof}

\begin{theorem}\label{thm:15.2}
Let $\varphi\colon A \to B$ be a homomorphism of Noetherian rings,
and suppose that the going-up theorem holds between $A$ and $B$. Then if
$p$ and $q$ are prime ideals of $A$ such that $p \supset q$, we have
\[
  \dim B \otimes \kappa(p) \ge \dim B \otimes \kappa(q).
\]
\end{theorem}

\begin{proof}
Set $r = \dim B \otimes \kappa(q)$ and $s = \htop(p/q)$. We choose a strictly
increasing chain $Q_0 \subset Q_1 \subset \cdots \subset Q_r$ of prime ideals of $B$ lying over $q$ and
a strictly increasing chain $q = p_0 \subset p_1 \subset \cdots \subset p_s = p$ of prime ideals of $A$.
By the going-up theorem there exists a chain $Q_r \subset Q_{r+1} \subset \cdots \subset Q_{r+s}$ of
prime ideals of $B$ such that $Q_{r+i} \cap A = p_i$. We set $P = Q_{r+s}$; then
\[
  \htop(P/qB) \ge r + s \quad \text{and} \quad P \cap A = p.
\]
Thus applying the previous theorem to the homomorphism $A/q \to B/qB$
induced by $\varphi$ we get $r + s \le \htop(P/qB) \le s + \dim B_P/pB_P$, and therefore
\[
  r \le \dim B_P/pB_P \le \dim B \otimes \kappa(p). \qquad\blacksquare
\]
\end{proof}

\begin{theorem}\label{thm:15.3}
Let $\varphi\colon A \to B$ be a homomorphism of Noetherian rings, and
suppose that the going-down theorem holds between $A$ and $B$. If $p$ and $q$
are prime ideals of $A$ with $p \supset q$ then
\[
  \dim B \otimes \kappa(p) \le \dim B \otimes \kappa(q).
\]
\end{theorem}

\begin{proof}
We may assume that $\htop(p/q) = 1$, and it is enough to prove that,
given a chain $P_0 \subset P_1 \subset \cdots \subset P_r$ of prime ideals of $B$ lying over $p$ such that
$\htop(P_i/P_{i-1}) = 1$ we can construct a chain of prime ideals $Q_0 \subset Q_1 \subset \cdots \subset Q_r$
of $B$ lying over $q$ such that
\[
  Q_i \subset P_i \quad (0 \le i \le r) \quad \text{and} \quad \htop(Q_i/Q_{i-1}) = 1 \quad (0 < i \le r).
\]
We can find $Q_0$ by going down. If $r \ge 1$ then take $x \in p - q$ and let $T_1, \ldots, T_s$
be the minimal prime divisors of $Q_0 + xB$. Then $\htop(T_j/Q_0) = 1$, while
$\htop(P_r/Q_0) \ge 2$, hence we can choose
\[
  y \in P_r - \bigcup_{j=1}^{s} T_j.
\]
Let $Q_1$ be a minimal prime divisor of $Q_0 + yB$ contained in $P_r$. Then
$\htop(Q_1/Q_0) = 1$, and $y \notin T_i$ for all $i$, hence $x \notin Q_1$.
Therefore $Q_1 \cap A \neq p$, and since $\htop(p/q) = 1$ we must have $Q_1 \cap A = q$. By
the same method we can successively construct $Q_1, Q_2, \ldots, Q_r$.
\end{proof}

\subsection*{2. Polynomial and formal power series rings}

\begin{theorem}\label{thm:15.4}
Let $A$ be a Noetherian ring, and $X_1, \ldots, X_n$ indeterminates
over $A$. Then
\[
  \dim A[X_1, \ldots, X_n] = \dim A[[X_1, \ldots, X_n]] = \dim A + n.
\]
\end{theorem}

\begin{proof}
It is enough to consider the case $n = 1$. For any $p \in \Spec A$, the
fibre $A[X] \otimes_A \kappa(p) = \kappa(p)[X]$ is a principal ideal ring, and therefore one-
dimensional; also $A[X]$ is a free $A$-module, hence faithfully flat, so by
Theorem 1, (ii), $\dim A[X] = \dim A + 1$.

For $A[[X]]$ it is not true in general that $A[[X]] \otimes_A \kappa(p)$ and $\kappa(p)[[X]]$
coincide; however, if $m$ is a maximal ideal of $A$ we have
\[
  A[[X]] \otimes_A \kappa(m) = A[[X]] \otimes_A (A/m) = (A/m)[[X]],
\]
and this fibre ring is one-dimensional. Also, as we saw on p. 4, every
maximal ideal $\mathfrak{M}$ of $A[[X]]$ is of the form $\mathfrak{M} = (m, X)$, where $m = \mathfrak{M} \cap A$
is a maximal ideal of $A$. Thus for a maximal ideal $\mathfrak{M}$ of $A[[X]]$ we have
\[
  \htop \mathfrak{M} = \htop(m \cap A) + 1;
\]
conversely, if $m$ is a maximal ideal of $A$ then $\htop(m, X) = \htop m + 1$, and
putting these together gives $\dim A[[X]] = \dim A + 1$.
\end{proof}

\begin{remark}
It is not necessarily true that a maximal ideal of $A[X]$ lies
over a maximal ideal of $A$. For example, if $A$ is a DVR and $t$ a
uniformising element then $A[t^{-1}] = K$ is the field of fractions of $A$, so
that $A[X]/(tX - 1) \cong K$, and $(tX - 1)$ is a maximal ideal of $A[X]$; however,
$(tX - 1) \cap A = (0)$.
\end{remark}

\begin{remark}
It is quite common for fibre rings of $A \to A[[X_1, \ldots, X_n]]$
to have dimension strictly greater than $n$. For example, let $k$ be a field and
set $A = k[Y, Z]$. It is well-known that the field of fractions of $k[[X]]$ has
infinite transcendence degree over $k(X)$ (see [ZS], vol. II, p. 220). Let $u(X)$,
$v(X) \in k[[X]]$ be two elements algebraically independent over $k(X)$, and
define a $k$-homomorphism (continuous for the $X$-adic topology)
\[
  \varphi\colon A[[X]] = k[Y, Z][[X]] \to k[[X]]
\]
by $\varphi(X) = X$, $\varphi(Y) = u(X)$, $\varphi(Z) = v(X)$. If we set $\Ker \varphi = P$ then $P \cap A = (0)$,
and $A[[X]]/P \cong k[[X]]$ is one-dimensional. Now every maximal ideal of
$A[[X]]$ has height $3$, and, as we will see later, $A[[X]]$ is catenary, so that
$\htop P = 2$. Thus we see that the generic fibre of $A \to A[[X]]$ is two-
dimensional.
\end{remark}

\subsection*{3. The dimension inequality}

We say that a ring $A$ is \emph{universally catenary}\index{universally catenary ring} if $A$ is Noetherian and every
finitely generated $A$-algebra is catenary. Since any $A$-algebra generated
by $n$ elements is a quotient of $A[X_1, \ldots, X_n]$, and since a quotient of a
catenary ring is again catenary, a necessary and sufficient condition for a
Noetherian ring $A$ to be universally catenary is that $A[X_1, \ldots, X_n]$ is
catenary for every $n \ge 0$. (In fact it is known that it is sufficient for $A[X]$
to be catenary, compare Theorem 31.7.)

\begin{theorem}[I. S. Cohen~\cite{Coh3}]\label{thm:15.5}
Let $A$ be a Noetherian integral domain, and
$B$ an extension ring of $A$ which is an integral domain. Let $P \in \Spec B$ and
$p = P \cap A$; then we have
\[
  (*) \quad \htop P + \trdeg_{\kappa(p)} \kappa(P) \le \htop p + \trdeg_K B,
\]
where $\trdeg_K B$ is the transcendence degree of the field of quotients of $B$
over that of $A$.
\end{theorem}

\begin{proof}
We may assume that $B$ is finitely generated over $A$. For if the right
hand side is finite and $m$ and $t$ are non-negative integers such that $m < \htop P$
and $t < \trdeg_{\kappa(p)}\kappa(P)$, then there is a prime ideal chain
$P = P_0 \supset P_1 \supset \cdots \supset P_m$ in $B$. Take $a_i \in P_i - P_{i+1}$, $0 \le i < m$, and let
$c_1, \ldots, c_t \in B$ be such that their images modulo $P$ are algebraically
independent over $A/p$. Set $C = A[\{a_i\}, \{c_j\}]$. If the theorem holds for $C$,
then we have $m + t < \htop p + \trdeg_K C \le \htop p + \trdeg_K B$. Letting $m$ and $t$ vary
we see the validity of $(*)$.

We may furthermore assume, by induction, that $B$ is generated over $A$ by
a single element: $B = A[x]$. We can replace $A$ by $A_p$ and $B$ by $B_P = A_p[x]$,
and hence assume that $A$ is local and $p$ its maximal ideal. Set $k = A/p$ and
write $B = A[X]/Q$. If $Q = (0)$ then $B = A[X]$ and by Theorem 1 we have
$\htop P = \htop p + \htop(P/pB)$, and since $B/pB = k[X]$ we have either $P = pB$ or
$\htop(P/pB) = 1$. In both cases the equality holds in $(*)$.

If $Q \neq (0)$ then $\trdeg_K B = 0$. Since $A$ is a subring of $B$ we have
$Q \cap A = (0)$, so that writing $K$ for the field of fractions of $A$ we have
$\htop Q = \htop QK[X] = 1$. Let $P^*$ be the inverse image of $P$ in $A[X]$. Then
$P = P^*/Q$, $\kappa(P) = \kappa(P^*)$, and
\begin{align*}
  \htop P &\le \htop P^* - \htop Q = \htop P^* - 1 \\
          &= \htop p + 1 - \trdeg_{\kappa(p)}\kappa(P^*) - 1 = \htop p - \trdeg_{\kappa(p)}\kappa(P).
\end{align*}
\end{proof}

\begin{definition}
Suppose that $A$ and $B$ satisfy the conditions of the previous
theorem. We refer to $(*)$ as the \emph{dimension inequality}\index{dimension inequality}, and if the equality in $(*)$
holds for every $P \in \Spec B$, we say that the \emph{dimension formula}\index{dimension formula} holds between
$A$ and $B$. The above proof shows that dimension formula holds between $A$
and $A[x_1, \ldots, x_n]$.
\end{definition}

\begin{theorem}[Ratliff]\label{thm:15.6}
A Noetherian ring $A$ is universally catenary if and
only if the dimension formula holds between $A/p$ and $B$ for every prime
ideal $p$ of $A$ and every finitely generated extension ring $B$ of $A/p$ which
is an integral domain.
\end{theorem}

\begin{proof}[Proof of `only if']
If $A$ is universally catenary then so is $A/p$, so that we
need only consider the case that $A$ is an integral domain, and $B$ is a finitely
generated extension ring which is an integral domain. If
$B = A[X_1, \ldots, X_n]/Q$ and $P = P^*/Q$, then since $A[X_1, \ldots, X_n]$ is catenary
we have $\htop P = \htop P^* - \htop Q$, and an easy calculation proves our assertion.
\end{proof}

\begin{proof}[Proof of `if']
We suppose that $A$ is not universally catenary, so that there
exists a finitely generated $A$-algebra $B$ which is not catenary; without loss
of generality we can assume that $B$ is an integral domain. Write $p$ for the
kernel of the homomorphism $A \to B$. There exist prime ideals $P$ and $Q$
of $B$ such that
\[
  P \subset Q, \quad \htop(Q/P) = d \quad \text{but} \quad \htop Q > \htop P + d.
\]
We write $h = \htop P$, choose $a_1, \ldots, a_h \in P$ such that $\htop(a_1, \ldots, a_h) = h$, and
set $I = (a_1, \ldots, a_h)$, so that $P$ is a minimal prime divisor of $I$. Let
\[
  I = q_1 \cap \cdots \cap q_r
\]
be a shortest primary decomposition of $I$, with $P$ the minimal prime
divisor of $q_1$. Then for $b \in Q \cap q_2 \cap \cdots \cap q_r - P$ we have
\[
  I:b^\nu B = q_1 \quad \text{for } \nu = 1, 2, \ldots.
\]
We set $y_i = a_i/b$ for $1 \le i \le h$,
\[
  C = A[y_1, \ldots, y_h], \quad J = (y_1, \ldots, y_h)C \quad \text{and} \quad M = J + QC = J + Q.
\]
Every element of $C$ can be written in the form $u/b^k$ for suitable $k$, with
$u \in (I + bB)^k$, so that if $z \in J \cap B$ then $zb^\nu \in I$ holds for sufficiently large
$\nu$. Hence $z \in I:b^\nu = q_1$. The converse inclusion $q_1 \subset J \cap B$ is obvious, hence
$J \cap B = q_1$. Thus
\[
  M \cap B = (J + Q) \cap B = (J \cap B) + Q = Q,
\]
\[
  C/J \cong B/q_1, \quad \text{and} \quad C/M \cong B/Q.
\]
Therefore, $C_M/JC_M = B_Q/q_1B_Q$ is a $d$-dimensional local ring, and $J$ is
generated by $h$ elements, so that
\[
  \htop M = \dim C_M \le h + d \le \htop Q.
\]
Now $C$ and $B$ have the same field of fractions, and $\kappa(M) = \kappa(Q)$, so that
this inequality implies that the dimension formula does not hold between
$B$ and $C$. This is a contradiction, since we are assuming that the dimension
formula holds between $A/p$ and $B$ and between $A/p$ and $C$, and one sees
easily that it must then hold between $B$ and $C$.
\end{proof}

\subsection*{4. The Rees ring and gr${}_I$(A)}

Let $A$ be a ring, $I$ an ideal of $A$ and $t$ an indeterminate over $A$. Consider
$A[t]$ as a graded ring in the usual way. We obtain a graded ring $R_+ \subset A[t]$
by setting
\[
  R_+ = R_+(A, I) = \left\{\sum c_nt^n \mid c_n \in I^n\right\} = \bigoplus_{n \ge 0} I^nt^n \subset A[t].
\]
If $I = (a_1, \ldots, a_r)$ then $R_+$ can be written $R_+ = A[a_1t, \ldots, a_rt]$, so that $R_+$ is
Noetherian if $A$ is.

$R_+$ is related to the graded ring $\gr_I(A)$ associated with $A$ and $I$ by the fact
that
\[
  \gr_I(A) = \bigoplus_{n \ge 0} I^n/I^{n+1} \cong R_+/IR_+.
\]

Now let $u = t^{-1}$, and consider $A[t, u] = A[t, t^{-1}]$ as a $\Z$-graded ring in
the obvious way. The \emph{Rees ring}\index{Rees ring} $R(A, I)$ is the graded subring
\[
  R = R(A, I) = R_+[u] = \left\{\sum c_nt^n \in A[t, t^{-1}] \;\middle|\;
    \begin{array}{l}
      c_n \in I^n \text{ for } n < 0 \\
      c_n \in A \text{ for } n > 0
    \end{array} \right\}.
\]
Since
\[
  uR = \left\{\sum c_nt^n \;\middle|\; c_n \in I^{n+1} \text{ for } n \ge 0, \; c_n \in A \text{ for } n < 0\right\},
\]
we have $\gr_I(A) = R/uR$.

\smallskip
\noindent Set $S = \{1, u, u^2, \ldots\}$. Then $R_S = R[u^{-1}] = R[t] = A[t^{-1}, t]$, and
$R_S/(1 - u)R_S = A[t^{-1}, t]/(1 - t) = A$. But $R_S/(1 - u)R_S = (R/(1 - u)R)_{\bar{S}}$,
where $\bar{S}$ is the image of $S$ in $R/(1 - u)R$, and since $\bar{S} = 1$, we see that
$R_{\bar{S}}/(1 - u)R_{\bar{S}} = R/(1 - u)R$. Thus we have
\[
  R/(1 - u)R = A \quad \text{and} \quad R/uR = \gr_I(A),
\]
so that the graded ring $\gr_I(A)$ is a `deformation' of the original ring $A$,
with $R$ as `total space of the deformation', in the sense that $R$ contains a
parameter $u$ such that the values $u = 1$ and $0$ correspond to $A$ and $\gr_I(A)$,
respectively.

We also have
\[
  I^nR \cap A = I^n \quad \text{for all } n \ge 0
\]
and this property is often used to reduce problems about powers of $I$ to
the corresponding problems for powers of the principal ideal $uR$.

We conclude this section by applying the dimension inequality to the
study of the dimension of the Rees ring and $\gr_I(A)$.

Let $A$ be a Noetherian ring, $I = \sum_{i=1}^{r} a_iA$ a proper ideal of $A$, and $t$ an
indeterminate over $A$. We set
\[
  u = t^{-1}, \quad R = R(A, I) = A[u, a_1t, \ldots, a_rt] \quad \text{and} \quad G = \gr_I(A).
\]
We have $R \subset A[t, u]$ and $R/uR \cong G$. For any ideal $\mathfrak{a}$ of $A$, set
\[
  \mathfrak{a}' = \mathfrak{a}A[t, u] \cap R.
\]
Then $\mathfrak{a}' \cap A = \mathfrak{a}A[t, u] \cap A = \mathfrak{a}$, so that for $\mathfrak{a}_1 \neq \mathfrak{a}_2$ we have $\mathfrak{a}_1' \neq \mathfrak{a}_2'$.
Moreover, if $p$ is a prime ideal of $A$ then $p'$ is prime in $R$, and the same
thing goes for primary ideals. If $(0) = q_1 \cap \cdots \cap q_n$ is a primary decompo-
sition of $(0)$ in $A$ then $(0) = q_1' \cap \cdots \cap q_n'$ is a primary decomposition of $(0)$
in $R$. Hence if $p_{0i}$ (for $1 \le i \le m$) are all the minimal prime ideals of $A$ then
$\{p_{0i}'\}_{1 \le i \le m}$ is the set of all minimal prime ideals of $R$. Let $p$ be a prime
ideal of $A$ with $\htop p = h$, and let $p = p_0 \supset p_1 \supset \cdots \supset p_h$ be a strictly decreasing
chain of prime ideals of $A$; then $p' \supset p_1' \supset \cdots \supset p_h'$ is a strictly descending
chain of prime ideals of $R$, so that
\[
  \htop p \le \htop p'.
\]
Conversely, suppose that $P \in \Spec R$ and $P \cap A = p$. Let $p_{0i}$ be a minimal
prime of $R$ contained in $P$ and such that $\htop P = \htop (P/p_{0i}')$; then $R/p_{0i}' \cong
A/p_{0i}$, so that by the dimension inequality
\[
  \htop P = \htop (P/p_{0i}') \le \htop (p/p_{0i}) + 1 - \trdeg_{\kappa(p_{0i})}\kappa(P)
\]
\[
  \le \htop p + 1.
\]
Hence $\dim R \le \dim A + 1$. On the other hand $A[u, t] = R[u^{-1}]$ is a
localisation of $R$ so that $\dim R \ge \dim A[u, t] = \dim A + 1$, so that finally
\[
  \dim R = \dim A + 1.
\]
Moreover, for any $p \in \Spec A$ we set $\bar{a}_i = \text{image of } a_i \bmod p$, so that
\[
  R/p' = (A/p)[u, \bar{a}_1t, \ldots, \bar{a}_rt],
\]
and hence $\trdeg_{\kappa(p)} \kappa(p') = 1$; carrying out the above
calculation using the dimension inequality with $p'$ in place of $P$ we get
$\htop p' \le \htop p$, and so
\[
  \htop p = \htop p'.
\]
We now choose a maximal ideal $m$ of $A$ containing $I$; then since
$R/m' = (A/m)[u]$ we see that $\mathfrak{M} = (m', u)$ is a maximal ideal of $R$ and
$\mathfrak{M} \neq m'$, so that $\htop \mathfrak{M} > \htop m'$. However, by the dimension inequality, we
have $\htop \mathfrak{M} \le \htop m + 1 = \htop m' + 1$. Thus
\[
  \htop \mathfrak{M} = \htop m' + 1 = \htop m + 1.
\]
The element $u$ is a non-zero-divisor of $R$ so that considering a system of
parameters gives $\htop (\mathfrak{M}/uR) = \htop \mathfrak{M} - 1 = \htop m$. Thus providing that there
exists a maximal ideal such that $\htop m = \dim A$ containing $I$, (in particular
if $A$ is local), then we have
\[
  \dim G = \dim (R/uR) = \dim A.
\]
We summarise the above in the following theorem.

\begin{theorem}\label{thm:15.7}
Let $A$ be a Noetherian ring and $I$ a proper ideal; then
setting $R = R(A, I)$ and $G = \gr_I(A)$ we have
\[
  \dim R = \dim A + 1, \quad \dim G \le \dim A.
\]
If in addition $A$ is local, then
\[
  \dim G = \dim A.
\]
\end{theorem}

% ============ 习题 ============
\section*{Exercises to \S\ref{sec:15}}
\emph{Let $k$ be a field.}

\begin{exercise}\label{ex:15.1-ex}
Let $A = k[X, Y] \subset B = k[X, Y, X/Y]$, and $P = (Y, X/Y)B$, $p = (X, Y)A$;
then check that $P \cap A = p$, $\htop P = \htop p = 2$, and $\dim B_P/pB_P = 1$, and hence
that
\[
  \htop P < \htop p + \dim B_P/pB_P.
\]
Show also by a concrete example that the going-down theorem does not
hold between $A$ and $B$.
\end{exercise}

\begin{exercise}\label{ex:15.2-ex}
Does the going-up theorem hold between $A$ and $B$, where $A = k[x] \subset
B = k[x, Y]$?
\end{exercise}

\begin{exercise}\label{ex:15.3-ex}
In Theorem~\ref{thm:15.7}, construct an example where $\dim G < \dim A$.
\end{exercise}


% ============ Chapter 6. Regular sequences (原书 pp.123-153) ============
\chapter{Regular sequences}\label{ch:6}

In the 1950s homological algebra was introduced into commutative ring
theory, opening up new avenues of study. In this chapter we run through
some fundamental topics in this direction.

In \S\ref{sec:16} we define regular sequences, depth and the Koszul complex. The
notion of depth is not very geometric, and rather hard to grasp, but is
an extremely important invariant. It can be treated either in terms of
$\Ext$'s, or by means of the Koszul complex, and we give both versions. We
discuss the relation between regular and quasi-regular sequences in a
transparent treatment due to Rees. \S\ref{sec:17} contains the definition and principal
properties of Cohen--Macaulay (CM) rings. The theorem that quotients of
CM rings are always catenary is of great significance in dimension
theory. In \S\ref{sec:18} we treat a distinguished subclass of CM rings having even
nicer properties, the Gorenstein rings. In the famous paper of H.\ Bass~\cite{Ba1},
Gorenstein rings are discussed using Matlis' theory of injective modules.
Here we give an elementary treatment of Gorenstein rings following~\cite{G6}
before going through Matlis theory.

% ============ §16 Regular sequences and the Koszul complex (原书 pp.123-132) ============
\section{Regular sequences and the Koszul complex}\label{sec:16}\index{regular sequences, section 16}\index{Koszul complex, section 16}

Let $A$ be a ring and $M$ an $A$-module. An element $a \in A$ is said
to be $M$-regular\index{regular element} if $ax \neq 0$ for all $0 \neq x \in M$. A sequence $a_1, \ldots, a_n$ of elements
of $A$ is an \emph{$M$-sequence}\index{M-sequence} (or an \emph{$M$-regular sequence}) if the following two
conditions hold:
\begin{enumerate}[(1)]
  \item $a_1$ is $M$-regular, $a_2$ is $(M/a_1M)$-regular, $\ldots$, $a_n$ is $(M/\sum_{i=1}^{n-1} a_iM)$-
        regular;
  \item $M \neq \sum_{i=1}^{n} a_iM$.
\end{enumerate}
Note that, after permutation, the elements of an $M$-sequence may no longer
form an $M$-sequence.

\begin{theorem}\label{thm:16.1}
If $a_1, \ldots, a_n$ is an $M$-sequence then so is $a_1^{\nu_1}, \ldots, a_n^{\nu_n}$ for
positive integers $\nu_1, \ldots, \nu_n$.
\end{theorem}

\begin{proof}
It is sufficient to prove that if $a_1, \ldots, a_n$ is an $M$-sequence then so is
$a_1^{\nu_1}, a_2, \ldots, a_n$. Indeed, assuming this, we have in turn that $a_1^{\nu_1}, a_2, \ldots, a_n$
is an $M$-sequence, then setting $M_1 = M/a_1^{\nu_1}M$ we have that $a_2, a_3, \ldots, a_n$,
and hence also $a_2^{\nu_2}, a_3, \ldots, a_n$, is an $M_1$-sequence, and so on. Also, the
second condition $M \neq \sum a_i^{\nu_i}M$ is obvious.

Let us now prove by induction on $n$ that if $b_1, \ldots, b_n$ is an $M$-sequence,
and if $\sum_{i=1}^{n} b_it_i = 0$ with $t_i \in M$ then $t_i \in b_1M + \cdots + b_nM$ for all
$i$. First of all from the condition that $b_n$ is not a zero-divisor modulo
$b_1, \ldots, b_{n-1}$ we can write
\[
  t_n = \sum_{i=1}^{n-1} b_i\eta_i, \quad \text{with } \eta_i \in M.
\]
Therefore $\sum_{i=1}^{n-1} b_i(t_i + b_n\eta_i) = 0$, so that by induction we have
\[
  t_i + b_n\eta_i \in b_1M + \cdots + b_{n-1}M \quad \text{for } 1 \le i \le n - 1,
\]
giving $t_i \in b_1M + \cdots + b_nM$ for $1 \le i \le n - 1$. The condition for $t_n$ is already known.

Now assuming $\nu > 1$ we prove by induction on $\nu$ that $a_1^{\nu}, a_2, \ldots, a_n$ is
an $M$-sequence. Since $a_1$ is $M$-regular, so is $a_1^{\nu}$. For $i > 1$, suppose that
for some $\sigma \in M$ we have
\[
  a_i\sigma = a_1^{\nu}\xi_1 + a_2\xi_2 + \cdots + a_{i-1}\xi_{i-1} \quad \text{with } \xi_j \in M.
\]
Then since $a_1^{\nu-1}, a_2, \ldots, a_i$ is an $M$-sequence, we can write
\[
  \sigma = a_1^{\nu-1}\eta_1 + \cdots + a_{i-1}\eta_{i-1} \quad \text{with } \eta_j \in M.
\]
Hence we get
\[
  0 = a_1^{\nu-1}(a_i\eta_1 - a_i\xi_1) + a_2(\eta_2 - a_i\xi_2) + \cdots + a_{i-1}(\eta_{i-1} - a_i\xi_{i-1}).
\]
The above assertion gives $a_i\eta_1 - a_i\xi_1 \in a_1^{\nu-1}M + a_2M + \cdots + a_{i-1}M$,
and hence $a_i\eta_1 \in a_1M + a_2M + \cdots + a_{i-1}M$. Therefore $\eta_1 \in a_1M + \cdots +
a_{i-1}M$, and so as required we have $\sigma \in a_1^{\nu}M + a_2M + \cdots + a_{i-1}M$.
\end{proof}

Let $A$ be a ring, $X_1, \ldots, X_n$ indeterminates over $A$, and $M$ an $A$-module.
We can view elements of $M \otimes_A A[X_1, \ldots, X_n]$ as polynomials in the $X_i$
with coefficients in $M$,
\[
  F(X) = F(X_1, \ldots, X_n) = \sum t_{\alpha_1 \cdots \alpha_n}X_1^{\alpha_1} \cdots X_n^{\alpha_n}, \quad \text{with } t_{\alpha} \in M.
\]
For this reason we write $M[X_1, \ldots, X_n]$ for $M \otimes_A A[X_1, \ldots, X_n]$; we can
consider this either as an $A$-module or as an $A[X_1, \ldots, X_n]$-module. For
$a_1, \ldots, a_n \in A$ and $F \in M[X_1, \ldots, X_n]$, we can substitute the $a_i$ for $X_i$ to
get $F(a_1, \ldots, a_n) \in M$.

\begin{definition}
Let $a_1, \ldots, a_n \in A$, set $I = \sum a_iA$ and let $M$ be an $A$-module
with $IM \neq M$. We say that $a_1, \ldots, a_n$ is an $M$-quasi-regular sequence\index{quasi-regular sequence} if the
following condition holds for each $\nu$:
\[
  (*) \quad
  \begin{aligned}
    &F(X_1, \ldots, X_n) \in M[X_1, \ldots, X_n] \text{ is homogeneous of degree } \nu \text{ and }\\
    &F(a) \in I^{\nu+1}M \text{ implies that all the coefficients of } F \text{ are in } IM.
  \end{aligned}
\]
This notion is obviously independent of the order of $a_1, \ldots, a_n$.
\end{definition}

In the above definition it would not make any difference if we replaced
the condition that $F(a) \in I^{\nu+1}M$ by the condition $F(a) = 0$. Indeed, if $F$ is
homogeneous of degree $\nu$ and $F(a) \in I^{\nu+1}M$ then there exist a homogeneous
element $G(X) \in M[X_1, \ldots, X_n]$ of degree $\nu + 1$ such that $F(a) = G(a)$. Then
write $G(X) = \sum X_iG_i(X)$ with each $G_i$ homogeneous of degree $\nu$, and set
$F^*(X) = F(X) - \sum a_iG_i(X)$, so that $F^*$ is homogeneous of degree $\nu$ and
$F^*(a) = 0$. Moreover, if $F^*$ has coefficients in $IM$ then so does $F$.

We can define a map $\varphi\colon (M/IM)[X_1, \ldots, X_n] \to \gr_I M = \bigoplus_{\nu \ge 0}
I^\nu M/I^{\nu+1}M$ as follows: taking a homogeneous element $F(X) \in M[X]$ of
degree $\nu$ into the class of $F(a)$ in $I^\nu M/I^{\nu+1}M$ provides a homomorphism (of
additive groups) from $M[X]$ into $\gr_I M$ which preserves degrees. Since
$IM[X]$ is in the kernel, this induces a homomorphism
\[
  \varphi\colon M[X]/IM[X] = (M/IM)[X] \to \gr_I M,
\]
which is obviously surjective. Then $a_1, \ldots, a_n$ is a quasi-regular sequence
precisely when $\varphi$ is injective, and hence an isomorphism.

\begin{theorem}\label{thm:16.2}
Let $A$ be a ring, $M$ an $A$-module, and $a_1, \ldots, a_n \in A$; set
$I = (a_1, \ldots, a_n)A$. Then we have the following:
\begin{enumerate}[(i)]
  \item if $a_1, \ldots, a_n$ is an $M$-sequence then it is $M$-quasi-regular;
  \item if $a_1, \ldots, a_n$ is an $M$-quasi-regular sequence, and if $x \in A$ satisfies
        $IM:x = IM$ then $I^\nu M:x = I^\nu M$ for any $\nu > 0$.
\end{enumerate}
\end{theorem}

\begin{proof}[Proof (taken from Rees [5])]
First of all we prove (ii) by induction on $\nu$.
The case $\nu = 1$ is just the assumption; suppose that $\nu > 1$. For $\xi \in M$, if
$x\xi \in I^\nu M$ then also $x\xi \in I^{\nu-1}M$, so that by the inductive hypothesis
$\xi \in I^{\nu-1}M$, and hence we can write $\xi = F(a)$ with $F = F(X) \in M[X_1, \ldots,
X_n]$ homogeneous of degree $\nu - 1$. Now $x\xi = xF(a) \in I^\nu M$, so that by
definition of quasi-regular sequence each coefficient of $xF(X)$ belongs to
$IM$. Using $IM:x = IM$ once more we find that the coefficients of $F(X)$
also belong to $IM$, and therefore $\xi = F(a) \in I^\nu M$.

Now we prove (i) by induction on $n$. The case $n = 1$ can easily be
checked. Suppose that $n > 1$, and that the statement holds up to $n - 1$, so
that in particular $a_1, \ldots, a_{n-1}$ is $M$-quasi-regular. Now let $F(X) \in
M[X_1, \ldots, X_n]$ be homogeneous of degree $\nu$, such that $F(a) = 0$. We prove
by induction on $\nu$ that the coefficients of $F$ belong to $IM$. We separate
out $F(X)$ into terms containing $X_n$ and not containing $X_n$, writing
\[
  F(X) = G(X_1, \ldots, X_{n-1}) + X_nH(X_1, \ldots, X_n).
\]
Here $G$ is homogeneous of degree $\nu$ and $H$ of degree $\nu - 1$. Then, as we
proved in (ii),
\[
  H(a) \in (a_1, \ldots, a_{n-1})^\nu M:a_n = (a_1, \ldots, a_{n-1})^\nu M \subset I^\nu M,
\]
and hence, by induction on $\nu$, the coefficients of $H(X)$ belong to $IM$.
Moreover, by the above formula there is a homogeneous polynomial
$h(X_1, \ldots, X_{n-1})$ of degree $\nu$ with coefficients in $M$ such that $H(a) =
h(a_1, \ldots, a_{n-1})$, and so setting
\[
  g(X) = G(X_1, \ldots, X_{n-1}) + a_nh(X_1, \ldots, X_{n-1}),
\]
since $a_1, \ldots, a_{n-1}$ is $M$-quasi-regular, the coefficients of $g$ lie in
$(a_1, \ldots, a_{n-1})M$; therefore the coefficients of $G$ lie in
\[
  (a_1, \ldots, a_n)M.
\]
\end{proof}

This theorem holds for any $A$ and $M$, but as we will see in the next
theorem, under some conditions we can say that conversely, quasi-regular
implies regular. Then the notions of regular and quasi-regular sequences
for $M$ coincide, and so reordering an $M$-sequence gives again an $M$-
sequence.

\begin{theorem}\label{thm:16.3}
Let $A$ be a Noetherian ring, $M \neq 0$ an $A$-module, and
$a_1, \ldots, a_n \in A$; set $I = (a_1, \ldots, a_n)A$. Under the condition
\[
  (*) \quad \text{each of } M, M/a_1M, \ldots, M/(a_1, \ldots, a_{n-1})M \text{ is } I\text{-adically separated},
\]
if $a_1, \ldots, a_n$ is $M$-quasi-regular it is an $M$-sequence.
\end{theorem}

\begin{remark}
The hypothesis $(*)$ holds in either of the following cases:
\begin{enumerate}[($\beta$)]
  \item $M$ is finite and $I \subset \rad(A)$;
  \item $A$ is an $\N$-graded ring, $M$ an $\N$-graded module, and each $a_i$ is
        homogeneous of positive degree.
\end{enumerate}
However, for a non-Noetherian ring $A$ there are examples where the
theorem fails (Dieudonn\'e [1]) even if $A$ is local, $M = A$ and $I \subset \rad(A)$.
\end{remark}

\begin{proof}
We prove first that $a_1$ is $M$-regular. If $\xi \in M$ with $a_1\xi = 0$ then
by hypothesis $\xi \in IM$. Then setting $\xi = \sum a_i\eta_i$ we get $0 = \sum a_1a_i\eta_i$, so that
$\eta_i \in IM$. Proceeding in the same way we get $\xi \in \bigcap_{\nu \ge 1} I^\nu M = (0)$.

Now set $M_1 = M/a_1M$; if we prove that $a_2, \ldots, a_n$ is an $M_1$-quasi-regular
sequence then the theorem follows by induction on $n$. (If $M$ is $I$-adically
separated and $M \neq 0$ then $M \neq IM$.) So let $\bar{f}(X_2, \ldots, X_n)$ be a homogeneous
polynomial of degree $\nu$ with coefficients in $M_1$, such that $\bar{f}(a_2, \ldots, a_n) = 0$.
If $F(X_1, \ldots, X_n)$ is a homogeneous polynomial of degree $\nu$ with coefficients
in $M$ which reduces to $\bar{f}$ modulo $a_1M$, then $F(a_1, \ldots, a_n) \in a_1M$. Set
$F(a_1, \ldots, a_n) = a_1\sigma$; suppose that $\sigma \in I^iM$, so that we can write $\sigma = G_i(a)$
with $G_i(X) \in M[X_1, \ldots, X_n]$ homogeneous of degree $i$. Then
\[
  a_1G_i(a_2, \ldots, a_n) = a_1G_i(a_1, \ldots, a_n),
\]
and if $i < \nu - 1$ it follows that the coefficients of $G_i$ belong to $IM$, so that
$\sigma \in I^{i+1}M$; repeating this argument we see that $\sigma \in I^{\nu-1}M$. Setting $i = \nu - 1$
in the above formula, then since $X_1$ does not appear in $F$, we can
apply the definition of quasi-regular sequence to $F(X_1, \ldots, X_n) -
X_1G_{\nu-1}(X_1, \ldots, X_n)$ to deduce that the coefficients of $F$ belong to $IM$.
Hence, the coefficients of $\bar{f}$ belong to $IM_1$.
\end{proof}

\begin{corollary*}
Let $A$ be a Noetherian ring, $M$ an $A$-module and $a_1, \ldots, a_n$
an $M$-sequence. If conditions ($\alpha$) or ($\beta$) of the above remark hold then any
permutation of $a_1, \ldots, a_n$ is again an $M$-sequence.
\end{corollary*}

Here is an example where a permutation of an $M$-sequence fails to be
an $M$-sequence: let $k$ be a field, $A = k[X, Y, Z]$ and set $a_1 = X(Y - 1)$,
$a_2 = Y$, $a_3 = Z(Y - 1)$. Then $(a_1, a_2, a_3)A = (X, Y, Z)A \neq A$, and $a_1, a_2, a_3$
is an $A$-sequence, whereas $a_1, a_3, a_2$ is not.

\subsection*{The Koszul complex}

Given a ring $A$ and $x_1, \ldots, x_n \in A$, we define a complex $K_\bullet$ as follows:
set $K_0 = A$, and $K_p = 0$ if $p$ is not in the range $0 \le p \le n$. For $1 \le p \le n$,
let $K_p = \bigoplus A e_{i_1 \cdots i_p}$ be the free $A$-module of rank $\binom{n}{p}$ with basis
$\{e_{i_1, \ldots, i_p} \mid 1 \le i_1 < \cdots < i_p \le n\}$. The differential $d\colon K_p \to K_{p-1}$ is defined
by setting
\[
  d(e_{i_1 \cdots i_p}) = \sum_{r=1}^{p} (-1)^{r-1}x_{i_r}e_{i_1 \cdots \hat{i}_r \cdots i_p}
\]
(for $p = 1$, set $d(e_i) = x_i$). One checks easily that $dd = 0$. This complex is
called the \emph{Koszul complex}\index{Koszul complex}, and written $K_\bullet(x_1, \ldots, x_n)$ (alternatively,
$K_\bullet(x)$ or $K_{\bullet, x}$). For an $A$-module $M$ we set $K_\bullet(x, M) = K_\bullet(x) \otimes_A M$.
Moreover, for a complex $C_\bullet$ of $A$-modules we set $C_\bullet(x) = C_\bullet \otimes K_\bullet(x)$. In
particular, for $n = 1$ the complex $K_\bullet(x)$ is just
\[
  0 \to 0 \to A \xrightarrow{x} A \to 0,
\]
and it is easy to check that $K_\bullet(x_1, \ldots, x_n) = K_\bullet(x_1) \otimes \cdots \otimes K_\bullet(x_n)$. Since the
tensor product of complexes satisfies $L_\bullet \otimes M_\bullet \cong M_\bullet \otimes L_\bullet$, the Koszul com-
plex is invariant (up to isomorphism) under permutation of $x_1, \ldots, x_n$. The
Koszul complex $K_\bullet(x, M)$ has homology groups $H_p(K_\bullet(x, M))$, which we
abbreviate to $H_p(x, M)$. Quite generally we have
\[
  H_0(x, M) = M/\mathfrak{a}M,
\]
where $\mathfrak{a} = \sum x_iA$, and
\[
  H_n(x, M) \cong \{m \in M \mid x_1m = \cdots = x_nm = 0\}.
\]

\begin{theorem}\label{thm:16.4}
Let $C_\bullet$ be a complex of $A$-modules and $x \in A$. Then we obtain
an exact sequence of complexes
\[
  0 \to C_\bullet \to C_\bullet(x) \to C'_\bullet \to 0,
\]
where $C'_\bullet$ is the complex obtained by shifting the degrees in $C_\bullet$ up by $1$ (that is
$C'_p = C_{p-1}$, and the differential of $C'_\bullet$ is that of $C_\bullet$). The homology long
exact sequence obtained from this is
\[
  \cdots \to H_p(C_\bullet) \to H_p(C_\bullet(x)) \to H_{p-1}(C_\bullet) \xrightarrow{(-1)^{p-1}x}
  H_{p-1}(C_\bullet) \to \cdots;
\]
we have $x \cdot H_p(C_\bullet(x)) = 0$ for all $p$.
\end{theorem}

\begin{proof}
From the fact that $K_1(x) = Ax$ and $K_0(x) = A$ and the definition
of tensor product of complexes, we can identify $C_p(x)$ with $C_p \oplus C_{p-1}$, and
for $\xi \in C_p$, $\eta \in C_{p-1}$ we have
\[
  d(\xi, \eta) = (d\xi + (-1)^{p-1}x\eta, d\eta).
\]
The first assertion is clear from this. Moreover, $H_p(C'_\bullet) = H_{p-1}(C_\bullet)$ is also
clear, and if $\eta \in C_{p-1} = C'_p$ satisfies $d\eta = 0$ then in $C_\bullet(x)$ we have $d(0, \eta) =
((-1)^{p-1}x\eta, 0)$, so that the long exact sequence has the form indicated
in the theorem. Finally, if $d(\xi, \eta) = 0$ then $d\eta = 0$ and $d\xi = (-1)^px\eta$, so
that $x \cdot (\xi, \eta) = (x\xi, (-1)^px\eta) \in d(C_{p+1}(x))$, and therefore $x \cdot H_p(C_\bullet(x)) = 0$.
\end{proof}

Applying this theorem to $K_\bullet(x, M)$ and using the commutativity of
tensor product of complexes, we see that the ideal $(\mathfrak{x}) = (x_1, \ldots, x_n)$
generated by $\mathfrak{x}$ annihilates the homology groups $H_p(x, M)$:
\[
  (\mathfrak{x}) \cdot H_p(x, M) = 0 \quad \text{for all } p.
\]

\begin{theorem}\label{thm:16.5}
\begin{enumerate}[(i)]
  \item Let $A$ be a ring, $M$ an $A$-module, and $x_1, \ldots, x_n$ an $M$-sequence; then
        \[
          H_p(x, M) = 0 \text{ for } p > 0 \quad \text{and} \quad H_0(x, M) = M/\mathfrak{x}M.
        \]
  \item Suppose that one of the following two conditions ($\alpha$) or ($\beta$) holds:
        \begin{enumerate}[($\alpha$)]
          \item $(A, m)$ is a local ring, $x_1, \ldots, x_n \in m$ and $M$ is a finite $A$-module;
          \item $A$ is an $\N$-graded ring, $M$ is an $\N$-graded $A$-module, and $x_1, \ldots, x_n$ are
                homogeneous elements of degree $> 0$.
        \end{enumerate}
        Then the converse of (i) holds in the following strong form: if $H_p(x, M) = 0$
        and $M \neq 0$ then $x_1, \ldots, x_n$ is an $M$-sequence.
\end{enumerate}
\end{theorem}

\begin{proof}
We use induction on $n$.

(i) When $n = 1$ we have $H_1(x, M) = \{\xi \in M \mid x\xi = 0\} = 0$, so that
the assertion holds. When $n > 1$, for $p > 1$ the previous theorem provides
an exact sequence
\[
  0 = H_p(x_1, \ldots, x_{n-1}, M) \to H_p(x_1, \ldots, x_n, M)
  \to H_{p-1}(x_1, \ldots, x_{n-1}, M) = 0.
\]
so that $H_p(x_1, \ldots, x_n, M) = 0$. For $p = 1$, setting $M_i = M/(x_1, \ldots, x_i)M$ we
have an exact sequence
\[
  0 \to H_1(x_n, M_{n-1}) \to H_1(x_1, \ldots, x_n, M) \to H_1(x_1, \ldots, x_{n-1}, M) = M_{n-2}
\]
\[
  \xrightarrow{x_n} M_{n-1},
\]
and since $x_n$ is $M_{n-1}$-regular we have $H_1(x_n, M) = 0$.

(ii) Assuming either ($\alpha$) or ($\beta$), $M \neq 0$ implies that $M_i \neq 0$ for $1 \le i \le n$. By
hypothesis and by the previous theorem,
\[
  H_1(x_1, \ldots, x_{n-1}, M) \xrightarrow{x_n} H_1(x_1, \ldots, x_{n-1}, M) \to H_1(x, M) = 0;
\]
but quite generally $H_1(x_1, \ldots, x_{n-1}, M)$ is a finite $A$-module in case ($\alpha$), or an $\N$-graded
module in case ($\beta$), so that by NAK, $H_1(x_1, \ldots, x_{n-1}, M) = 0$. Thus by
induction $x_1, \ldots, x_{n-1}$ is an $M$-sequence. Now by the same exact sequence
and the first part of (i), we see that $x_n$ is $M_{n-1}$-regular, and therefore
$x_1, \ldots, x_n$ is an $M$-sequence.
\end{proof}

Let $A$ be a ring, $M$ an $A$-module and $I$ an ideal of $A$. If $a_1, \ldots, a_n$ are
elements of $I$, we say that they form a \emph{maximal $M$-sequence in $I$}\index{maximal M-sequence} if $a_1, \ldots, a_n$
is an $M$-sequence, and $a_1, \ldots, a_n, b$ is not an $M$-sequence for any $b \in I$. If
$a_1, \ldots, a_n$ is an $M$-sequence then $a_1M, (a_1, a_2)M, \ldots, (a_1, \ldots, a_n)M$ is strictly
increasing, so that the chain of ideals $(a_1) \subset (a_1, a_2) \subset \cdots$ is also strict-
ly increasing. If $A$ is Noetherian this cannot continue indefinitely,
so that any $M$-sequence can be extended until we arrive at a maximal
$M$-sequence.

In Theorems 6--8 below, the hypothesis that $M$ is a finite
$A$-module can be weakened to the statement that $M$ is a finite $B$-module for a
homomorphism $A \to B$ of Noetherian rings, as one sees on inspecting
the proof. The reason for this is that, if we set $\Ass_B(M) = \{P_1, \ldots, P_r\}$ and
$P_i \cap A = p_i$, then any ideal of $A$ consisting entirely of zero-divisors of $M$ is
contained in $\bigcup p_i$, and therefore contained in one of the $p_i$. Note that
according to [M], (9.A), we have $\Ass_A(M) = \{p_1, \ldots, p_r\}$.

\begin{theorem}\label{thm:16.6}
Let $A$ be a Noetherian ring, $M$ a finite $A$-module and $I$
an ideal of $A$; suppose that $IM \neq M$. For a given integer $n > 0$ the following
conditions are equivalent;
\begin{enumerate}[(1)]
  \item $\Ext^i_A(N, M) = 0$ for all $i < n$ and for any finite $A$-module $N$ with
        $\Supp(N) \subset V(I)$;
  \item $\Ext^i_A(A/I, M) = 0$ for all $i < n$;
  \item $\Ext^i_A(A/I, M) = 0$ for all $i < n$ and for some finite $A$-module $N$ with
        $\Supp(N) = V(I)$;
  \item there exists an $M$-sequence of length $n$ contained in $I$.
\end{enumerate}
\end{theorem}

\begin{proof}
$(1) \Rightarrow (2) \Rightarrow (3)$ are obvious. For $(2') \Rightarrow (3)$, if $I$ consists only of zero-
divisors of $M$ then there exists an associated prime $P$ of $M$ containing $I$
(this is where we need the finiteness of $M$). Hence there is an injective
map $A/P \to M$. Localising at $P$, we see that $\Hom_{A_P}(A_P/PA_P, M_P) \neq 0$, where
$(A/P)_P = A_P/PA_P$. Now $P \in V(I) = \Supp(N)$, so that $N_P \neq 0$, and hence by
NAK, $N_P/PN_P = N \otimes_A k \neq 0$. Thus $N \otimes k$ is a non-zero vector space
over $k$, and $\Hom_A(N \otimes k, k) \neq 0$. Putting together what we have said,
we can follow the composite $N_P \to N \otimes k \to k \to M_P$ to show that
$\Hom_{A_P}(N_P, M_P) \neq 0$. The left-hand side is equal to $(\Hom_A(N, M))_P$, so that
$\Hom_A(N, M) \neq 0$. But this contradicts $(2')$. Hence $I$ contains an
$M$-regular element $f$. By assumption, $M/IM \neq 0$, and if $n = 1$ then we are
done. If $n > 1$ we set $M_1 = M/fM$; then from the exact sequence
\[
  0 \to M \xrightarrow{f} M \to M_1 \to 0
\]
we get $\Ext^i_A(N, M_1) = 0$ for $i < n - 1$, so that by induction $I$ contains an
$M_1$-sequence $f_2, \ldots, f_n$.

For the proof of $(3) \Rightarrow (1)$ we do not need to assume that $A$ is Noetherian
or $M$ finite. Let $f_1, \ldots, f_n \in I$ be an $M$-sequence; we have the exact sequence
\[
  0 \to M \xrightarrow{f_1} M \to M_1 \to 0,
\]
and if $n > 1$ the inductive hypothesis $\Ext^i_A(N, M_1) = 0$ for $i < n - 1$, so that
\[
  0 \to \Ext^i_A(N, M) \xrightarrow{f_1} \Ext^i_A(N, M)
\]
is exact for $i < n$. But $\Ext^i_A(N, M)$ is annihilated by elements of $\ann(N)$.
Since $\Supp(N) = V(\ann(N)) \subset V(I)$, we have $I \subset \sqrt{\ann(N)}$, and a
sufficiently large power of $f_1$ annihilates $\Ext^i_A(N, M)$. Therefore,
$\Ext^i_A(N, M) = 0$ for $i < n$.
\end{proof}

Let $M$ and $I$ be as in the above theorem, and $a_1, \ldots, a_n$ an $M$-sequence
in $I$. For $1 \le i \le n$, set $M_i = M/(a_1, \ldots, a_i)M$; then it is easy to see
that $\Hom_A(A/I, M_i) \supset \Ext^1_A(A/I, M_{i-1}) \supset \cdots \supset \Ext^i_A(A/I, M)$. Therefore,
if $\Ext^i_A(A/I, M) = 0$ we can find another element $a_{n+1} \in I$ such that
$a_1, \ldots, a_{n+1}$ is an $M$-sequence. Hence if $a_1, \ldots, a_n$ is a maximal $M$-sequence
in $I$ we must have $\Ext^n_A(A/I, M) \neq 0$. We thus obtain the following theorem.

\begin{theorem}\label{thm:16.7}
Let $A$ be a Noetherian ring, $I$ an ideal of $A$ and $M$ a finite
$A$-module such that $M \neq IM$; then the length of a maximal $M$-sequence
in $I$ is a well-determined integer $n$, and $n$ is determined by
\[
  \Ext^i_A(A/I, M) = 0 \text{ for } i < n \quad \text{and} \quad \Ext^n_A(A/I, M) \neq 0.
\]
We write $n = \depth(I, M)$, and call $n$ the \emph{$I$-depth}\index{depth} of $M$. (If $M = IM$, the
$I$-depth is by convention $\infty$.) Theorem 7 takes the form
\[
  \depth(I, M) = \inf\{i \mid \Ext^i_A(A/I, M) \neq 0\}.
\]
In particular for a Noetherian local ring $(A, m, k)$, we call $\depth(m, M)$
simply the \emph{depth} of $M$, and write $\depth M$ or $\depth_A M$:
\[
  \depth M = \inf\{i \mid \Ext^i_A(k, M) \neq 0\}.
\]
From Theorem 6 we see that if $V(I) = V(I')$ then $\depth(I, M) =
\depth (I', M)$; this also follows easily from Theorem 1.

If $\ann (M) = \mathfrak{a}$ and we set $\bar{A} = A/\mathfrak{a}$ then $M$ is also an $\bar{A}$-module. Writing
$\bar{a}$ or $\bar{I}$ for the image of an element $a$ or an ideal $I$ of $A$ under the natural
homomorphism $A \to \bar{A}$ we clearly have that $a_1, \ldots, a_n$ is an $M$-sequence
if and only if $\bar{a}_1, \ldots, \bar{a}_n$ is. Thus $\depth (I, M) = \depth (\bar{I}, M)$, and if we set
$I + \mathfrak{a} = J$, then since $\bar{I} = \bar{J}$ we also have $\depth (I, M) = \depth (J, M)$.
\end{theorem}

We can also prove that the length of a maximal $M$-sequence is
well-determined by means of the Koszul complex.

\begin{theorem}\label{thm:16.8}
Let $A$ be a Noetherian ring, $I = (y_1, \ldots, y_n)$ an ideal of $A$,
and $M$ a finite $A$-module such that $M \neq IM$. If we set
\[
  q = \sup\{i \mid H_i(y, M) \neq 0\},
\]
then any maximal $M$-sequence in $I$ has length $n - q$.
\end{theorem}

\begin{proof}
Let $x_1, \ldots, x_s$ be a maximal $M$-sequence in $I$; we argue by induction
on $s$. If $s = 0$ then every element of $I$ is a zero-divisor of $M$, so that there
exists $P \in \Ass(M)$ containing $I$. By definition of $\Ass$, there exists $0 \neq \xi \in M$
such that $P = \ann(\xi)$, and hence $I\xi = 0$. Thus $\xi \in H_n(y, M)$ so that $q = n$,
and the assertion holds in this case.

If $s > 0$ we set $M_1 = M/x_1M$; then from the exact sequence
\[
  0 \to M \xrightarrow{x_1} M \to M_1 \to 0
\]
and from the fact that $IH_i(y, M) = 0$ (by Theorem 4), it follows that
\[
  0 \to H_i(y, M) \to H_i(y, M_1) \to H_{i-1}(y, M) \to 0
\]
is exact for every $i$. Thus $H_{q+1}(y, M_1) \neq 0$ and $H_i(y, M_1) = 0$ for $i > q + 1$;
but $x_2, \ldots, x_s$ is a maximal $M_1$-sequence in $I$, so that by induction we
have $q + 1 = n - (s - 1)$, and therefore $q = n - s$.
\end{proof}

In other words, $\depth(I, M)$ is the number of successive zero terms from
the left in the sequence
\[
  H_n(y, M), H_{n-1}(y, M), \ldots, H_0(y, M) = M/IM \neq 0.
\]
This fact is sometimes referred to as the `depth sensitivity' of the Koszul
complex.

\begin{corollary*}
In the situation of the theorem, $y_1, \ldots, y_n$ is an $M$-sequence if
and only if $\depth (I, M) = n$.
\end{corollary*}

\begin{proof}
$\depth(I, M) = n \iff H_i(y, M) = 0$ for all $i > 0 \iff y$ is an $M$-sequence.
\end{proof}

\subsection*{Grade}

A little before Auslander and Buchsbaum [2], Rees [5] introduced and
developed the theory of another notion related to regular sequences, that of
grade.\index{grade} Let $A$ be a Noetherian ring and $M \neq 0$ a finite $A$-module. Then Rees
made the definition
\[
  \grade M = \inf\{i \mid \Ext^i_A(M, A) \neq 0\}.
\]
For a proper ideal $J$ of $A$ we also call $\grade (A/J)$ the \emph{grade of the ideal}
$J$, and write $\grade J$. If we set $\mathfrak{a} = \ann(M)$ then since $\Supp(M) = V(\mathfrak{a})$,
Theorem 6 gives $\grade M = \depth(\mathfrak{a}, A)$. Moreover, if $g = \grade M$ then
$\Ext^g_A(M, A) \neq 0$, so that
\[
  \grade M \le \projdim M.
\]
If $I$ is an ideal then $\grade I (=\grade(A/I)) = \depth(I, A)$ is the length
of a maximal $A$-sequence in $I$, but in general if $a_1, \ldots, a_r$ is an $A$-sequence
then one sees easily from Theorem~\ref{thm:13.5} that $\htop(a_1, \ldots, a_r) = r$. Thus if
$a_1, \ldots, a_r$ is a maximal $A$-sequence in $I$, we have $r = \htop(a_1, \ldots, a_r) \le \htop I$.
Hence for an ideal $I$ we have $\grade I \le \htop I$.

\begin{theorem}\label{thm:16.9}
Let $A$ be a Noetherian ring, and $M, N$ finite $A$-modules;
suppose $M \neq 0$, $\grade M = k$ and $\projdim N = l \le k$. Then
\[
  \Ext^i_A(M, N) = 0 \quad \text{for } i < k - l.
\]
\end{theorem}

\begin{proof}
We use induction on $l$. If $l = 0$ then $N$ is a direct summand of some
free module $A^m$, so that we need only say what happens for $N = A$, but
then the assertion is just the definition of grade. If $l > 0$ we choose an
exact sequence
\[
  0 \to N_1 \to L_0 \to N \to 0
\]
with $L_0$ a finite free module; then $\projdim N_1 = l - 1$, so that by induction
\[
  \Ext^i_A(M, L_0) = 0 \quad \text{for } i < k \quad \text{and}
\]
\[
  \Ext^i_A(M, N_1) = 0 \quad \text{for } i < k - l;
\]
the assertion follows from this.
\end{proof}

% ============ 习题 ============
\section*{Exercises to \S\ref{sec:16}}
\emph{Prove the following propositions.}

\begin{exercise}\label{ex:16.1-ex}
Let $(A, m)$ be a Noetherian local ring, $M \neq 0$ a finite $A$-module, and
$a_1, \ldots, a_r \in m$ an $M$-sequence. Set $M' = M/(a_1, \ldots, a_r)M$. Then
\[
  \dim M' = \dim M - r.
\]
\end{exercise}

\begin{exercise}\label{ex:16.2-ex}
Let $A$ be a Noetherian ring, $\mathfrak{a}$ and $\mathfrak{b}$ ideals of $A$; then if $\grade \mathfrak{a}
> \projdim A/\mathfrak{b}$ we have $\mathfrak{b}:\mathfrak{a} = \mathfrak{b}$.
\end{exercise}

\begin{exercise}\label{ex:16.3-ex}
Let $A$ be a Noetherian ring. A proper ideal $I$ of $A$ is called a \emph{perfect ideal} if
$\grade I = \projdim A/I$. If $I$ is a perfect ideal of grade $k$ then all the prime
divisors of $I$ have grade $k$.
\end{exercise}

\begin{remark}
Quite generally, we have $\grade I (=\grade(A/I)) \le \projdim A/I$. If $A$ is a
regular local ring and $P \in \Spec A$ then as we will see in Theorems 19.1 and 19.2, $P$ is
perfect, that is $A/P$ is Cohen--Macaulay.
\end{remark}

\begin{exercise}\label{ex:16.4-ex}
Let $f\colon A \to B$ be a flat ring homomorphism, $M$ an $A$-module, and
$a_1, \ldots, a_n \in A$ an $M$-sequence; if $(M/(a_1, \ldots, a_n)M) \otimes_A B \neq 0$ then
$f(a_1), \ldots, f(a_n)$ is an $M \otimes_A B$-sequence.
\end{exercise}

\begin{exercise}\label{ex:16.5-ex}
Let $A$ be a Noetherian local ring, $M$ a finite $A$-module, and $P$ a prime ideal
of $A$; show that $\depth(P, M) \le \depth_{A_P} M_P$, and construct an example
where the inequality is strict.
\end{exercise}

\begin{exercise}\label{ex:16.6-ex}
Let $A$ be a ring and $a_1, \ldots, a_n \in A$ an $A$-quasi-regular sequence. If $A$
contains a field $k$ then $a_1, \ldots, a_n$ are algebraically independent over $k$.
\end{exercise}

\begin{exercise}\label{ex:16.7-ex}
Let $(A, m)$ and $(B, n)$ be Noetherian local rings, and suppose that $A \subset B$,
$n \cap A = m$ and that $mB$ is an $n$-primary ideal. Then for a finite $B$-module
$M$ we have
\[
  \depth_A M = \depth_B M.
\]
\end{exercise}

\begin{exercise}\label{ex:16.8-ex}
Let $A$ be a ring, $P_1, \ldots, P_r$ prime ideals, $I$ an ideal, and $x$ an element of $A$. If
$x \notin I + \bigcup_{i=1}^{r} P_i$ then there is a $y \in I$ such that $x + y \notin \bigcup_{i=1}^{r} P_i$
(E. Davis).
\end{exercise}

\begin{exercise}\label{ex:16.9-ex}
Use the previous question to show the following: let $A$ be a Noetherian
ring, and suppose that $I \neq A$ is an ideal generated by $n$ elements; then
$\grade I \le n$, and if $\grade I = n$ then $I$ can be generated by an $A$-sequence
([K], Th. 125).
\end{exercise}

\begin{exercise}\label{ex:16.10-ex}
Let $A$ be a Noetherian ring, and suppose that $P$ is a height $r > 0$ prime
ideal generated by $r$ elements $a_1, \ldots, a_r$.
\begin{enumerate}[(i)]
  \item Suppose either that $A$ is local, or that $A$ is $\N$-graded and the $a_i$ are
        homogeneous of positive degree. Then $A$ is an integral domain, and for
        $1 \le i \le r$ the ideal $(a_1, \ldots, a_i)$ is prime; hence $a_1, \ldots, a_r$ is an $A$-sequence.
  \item In general $a_1, \ldots, a_r$ does not have to be an $A$-sequence, but $P$ can in
        any case be generated by an $A$-sequence (E. Davis).
\end{enumerate}
\end{exercise}

% ============ §17 Cohen-Macaulay rings (原书 pp.133-138) ============
\section{Cohen--Macaulay rings}\label{sec:17}\index{Cohen-Macaulay rings, section 17}

\begin{theorem}[Ischebeck]\label{thm:17.1}
Let $(A, m)$ be a Noetherian local ring, $M$ and $N$
non-zero finite $A$-modules, and suppose that $\depth M = k$, $\dim N = r$. Then
\[
  \Ext^i_A(N, M) = 0 \quad \text{for } i < k - r.
\]
\end{theorem}

\begin{proof}
By induction on $r$; if $r = 0$ then $\Supp(N) = \{m\}$ and the assertion
holds by Theorem~\ref{thm:16.6}. Suppose $r > 0$. By Theorem~\ref{thm:6.4}, there exists a
chain
\[
  N = N_0 \supset N_1 \supset \cdots \supset N_t = (0) \quad \text{with } N_j/N_{j+1} \cong A/P_j
\]
of submodules $N_j$, where $P_j \in \Spec A$. It is easy to see that if $\Ext^i_A
(N_j/N_{j+1}, M) = 0$ for each $j$ then $\Ext^i_A(N, M) = 0$, and since $\dim N_j/N_{j+1} \le
\dim N = r$ it is enough to prove that $\Ext^i_A(N, M) = 0$ for $i < k - r$ in the
case $N = A/P$ with $P \in \Spec A$ and $\dim N = r$. Since $r > 0$ we can take
an element $x \in m - P$ and get the exact sequence
\[
  0 \to N \xrightarrow{x} N \to N' \to 0,
\]
where $N' = A/(P, x)$; then $\dim N' < r$ so that by induction we have
$\Ext^i_A(N', M) = 0$ for $i < k - r + 1$. Thus for $i < k - r$ we have an exact
sequence
\[
  0 \to \Ext^i_A(N, M) \xrightarrow{x} \Ext^i_A(N, M) \to \Ext^{i+1}_A(N', M) = 0.
\]
We have $x \in m$ so that by NAK, $\Ext^i_A(N, M) = 0$.
\end{proof}

\begin{theorem}\label{thm:17.2}
Let $A$ be a Noetherian local ring, $M$ a finite $A$-module,
and assume that $P \in \Ass(M)$; then $\dim(A/P) \ge \depth M$.
\end{theorem}

\begin{proof}
If $P \in \Ass(M)$ then $\Hom_A(A/P, M) \neq 0$, so that by the previous
theorem we cannot have $\dim A/P < \depth M$.
\end{proof}

\begin{definition}
Let $(A, m, k)$ be a Noetherian local ring, and $M$ a finite
$A$-module. We say that $M$ is a \emph{Cohen--Macaulay module}\index{Cohen-Macaulay module} (abbreviated to
CM module) if $M \neq 0$ and $\depth M = \dim M$, or if $M = 0$. If $A$ itself is a
CM module we say that $A$ is a \emph{CM ring}\index{Cohen-Macaulay ring} or a \emph{Macaulay ring}.
\end{definition}

\begin{theorem}\label{thm:17.3}
Let $A$ be a Noetherian local ring and $M$ a finite $A$-module.
\begin{enumerate}[(i)]
  \item If $M$ is a CM module then for any $P \in \Ass(M)$ we have
        $\dim(A/P) = \dim M = \depth M$. Hence $M$ has no embedded associated
        primes.
  \item If $a_1, \ldots, a_r \in m$ is an $M$-sequence and we set $M' = M/(a_1, \ldots, a_r)M$ then
        \[
          M \text{ is a CM module} \iff M' \text{ is a CM module}.
        \]
  \item If $M$ is a CM module then $M_P$ is a CM module over $A_P$ for
        every $P \in \Spec A$, and if $M_P \neq 0$ then
        \[
          \depth (P, M) = \depth_{A_P} M_P.
        \]
\end{enumerate}
\end{theorem}

\begin{proof}
(i) Quite generally, we have
\[
  \dim M = \sup\{\dim A/P \mid P \in \Ass M\}
\]
\[
  \ge \inf\{\dim A/P \mid P \in \Ass M\} \ge \depth M,
\]
so that this is clear.

(ii) By definition $\depth M' = \depth M - r$, and by Ex.~\ref{ex:16.1-ex}, $\dim M' =
\dim M - r$, so that this is clear.

(iii) It is enough to consider the case $M_P \neq 0$, when $P \supset \ann(M)$. Then
quite generally we have $\dim M_P \ge \depth M_P \ge \depth(P, M)$, so that we need
only show that
\[
  \dim M_P = \depth (P, M).
\]
We prove this by induction on $\depth(P, M)$. If $\depth(P, M) = 0$ then $P$ is
contained in an associated prime of $M$, but in view of $P \supset \ann(M)$ and
the fact that by (i) all the associated primes of $M$ are minimal, it follows
that $P$ is itself an associated prime of $M$; therefore $\dim M_P = 0$. If
$\depth(P, M) > 0$ then we can take an $M$-regular element $a \in P$, and set
$M' = M/aM$. Then
\[
  \depth (P, M') = \depth (P, M) - 1,
\]
$M'$ is a CM module with $M'_P \neq 0$, so that by induction $\dim M'_P =
\depth(P, M) - 1$; moreover, $a$ is $M_P$-regular as an element of $A_P$, and
so that using Ex.~\ref{ex:16.1-ex} once more, we have $\dim M_P =
\dim M'_P + 1$. Putting these together gives $\depth (P, M) = \dim M_P$.
\end{proof}

\begin{theorem}\label{thm:17.4}
Let $(A, m)$ be a CM local ring.
\begin{enumerate}[(i)]
  \item For a proper ideal $I$ of $A$ we have
        \[
          \htop I = \depth (I, A) = \grade I, \quad \text{and} \quad \htop I + \dim A/I = \dim A.
        \]
  \item For elements $a_1, \ldots, a_r \in m$ the following four conditions are
        equivalent:
        \begin{enumerate}[(1)]
          \item $a_1, \ldots, a_r$ is an $A$-sequence;
          \item $\htop(a_1, \ldots, a_i) = i$ for $1 \le i \le r$;
          \item $\htop(a_1, \ldots, a_r) = r$;
          \item $a_1, \ldots, a_r$ is part of a system of parameters of $A$.
        \end{enumerate}
  \item If $P \subset Q$ are prime ideals of $A$ then
        \[
          \htop Q - \htop P = \htop (Q/P).
        \]
\end{enumerate}
\end{theorem}

\begin{proof}
(iii) The implication $(1) \Rightarrow (2)$ follows from Theorem~\ref{thm:13.5}, together with
the fact that from the definition of $A$-sequence we have $0 < \htop(a_1) \le
\htop(a_1, a_2) \le \cdots \le r$.

$(2) \Rightarrow (3)$ is obvious.

$(3) \Rightarrow (4)$ If $\dim A = r$ this is obvious; if $\dim A > r$ then $m$ is not a minimal
prime divisor of $(a_1, \ldots, a_r)$, so that we can choose $a_{r+1} \in m$ not contained
in any minimal prime divisor of $(a_1, \ldots, a_r)$, and then $\htop(a_1, \ldots, a_{r+1}) = r + 1$.
Proceeding in the same way we arrive at a system of parameters of $A$.
(Note that so far we have not used the CM assumption.)

$(4) \Rightarrow (1)$ It is enough to show that any system of parameters $x_1, \ldots, x_n$
($n = \dim A$) is an $A$-sequence. If $P \in \Ass(A)$ then by Theorem 3, (i),
$\dim A/P = n$, so that $x_1 \notin P$. Thus $x_1$ is $A$-regular. Therefore if we set
$A' = A/x_1A$ we have by the previous theorem that $A'$ is an $(n - 1)$-
dimensional CM ring, and the images of $x_2, \ldots, x_n$ form a system of
parameters of $A'$. Thus by induction on $n$ we see that $x_1, \ldots, x_n$ is an
$A$-sequence.

(i) If $\htop I = r$ then we can take $a_1, \ldots, a_r \in I$ such that $\htop(a_1, \ldots, a_i) = i$ for
$1 \le i \le r$. Thus by (iii), $a_1, \ldots, a_r$ is an $A$-sequence. Thus $r \le \grade I$.
Conversely, if $b_1, \ldots, b_s$ is an $A$-sequence in $I$ then $\htop(b_1, \ldots, b_s) = s \le \htop I$, and
hence $r \ge \grade I$, so that equality must hold. For the second equality, let
$S$ be the set of minimal prime divisors of $I$, we have
\[
  \htop I = \inf\{\htop P \mid P \in S\}
\]
and
\[
  \dim(A/I) = \sup\{\dim A/P \mid P \in S\},
\]
so it is enough to show that $\htop P = \dim A - \dim A/P$ for every $P \in S$.
Let $\htop P = r$ and $\dim A = n$. By Theorem 3, (iii), $A_P$ is a CM ring
and $r = \depth(P, A)$. Now if we take an $A$-sequence $a_1, \ldots, a_r \in P$ then
by Theorem 3, (ii), $A/(a_1, \ldots, a_r)$ is an $(n - r)$-dimensional CM ring, and from
the fact that $\htop (a_1, \ldots, a_r) = r = \htop P$ we see that $P$ is a minimal prime
divisor of $(a_1, \ldots, a_r)$; thus by Theorem 3, (i), $\dim A/P = \dim A/(a_1, \ldots, a_r) =
n - r$.

(ii) Let $P \subset Q$ be prime ideals of $A$. Then since $A_Q$ is a CM ring, (i)
above gives $\dim A_Q = \htop QA_Q + \dim A_Q/QA_Q$; in other words $\htop P - \htop Q =
\htop (Q/P)$.
\end{proof}

If one system of parameters of a Noetherian local ring $A$ is an $A$-sequence
then $\depth A = \dim A$, so that $A$ is a CM ring, and therefore, by the above
theorem, every system of parameters of $A$ is an $A$-sequence.

\begin{theorem}\label{thm:17.5}
Let $A$ be a Noetherian local ring and $\hat{A}$ its completion;
then
\begin{enumerate}[(i)]
  \item $\depth A = \depth \hat{A}$;
  \item $A$ is CM $\iff \hat{A}$ is CM.
\end{enumerate}
\end{theorem}

\begin{proof}
(i) This comes for example from the fact that
\[
  \Ext^i_A(A/m, A) \otimes_A \hat{A} = \Ext^i_{\hat{A}}(\hat{A}/m\hat{A}, \hat{A}) \quad \text{for all } i.
\]
(ii) follows from (i) and the fact that
$\dim A = \dim \hat{A}$.
\end{proof}

\begin{definition}
A proper ideal $I$ in a Noetherian ring $A$ is said to be \emph{unmixed}\index{unmixed ideal}
if the heights of its prime divisors are all equal. We say that the \emph{unmixedness}
\emph{theorem} holds for $A$ if for every $r > 0$, every height $r$ ideal $I$ of $A$ generated
by $r$ elements is unmixed. This includes as the case $r = 0$ the statement
that $(0)$ is unmixed. By Theorem~\ref{thm:13.5}, if $I$ is an ideal satisfying the
hypotheses of this proposition, then all the minimal prime divisors of $I$
have height $r$, so that to say that $I$ is unmixed is to say that $I$ does not
have embedded prime divisors.
\end{definition}

A Noetherian ring $A$ is said to be a \emph{CM ring}\index{Cohen-Macaulay ring} if $A_m$ is a CM local
ring for every maximal ideal $m$ of $A$. By Theorem 3, (iii), a localisation
$S^{-1}A$ of a CM ring $A$ is again CM.

\begin{theorem}\label{thm:17.6}
A necessary and sufficient condition for a Noetherian ring
$A$ to be a CM ring is that the unmixedness theorem holds for $A$.
\end{theorem}

\begin{proof}
First suppose that $A$ is a CM ring and that $I = (a_1, \ldots, a_r)$ is an
ideal of $A$ with $\htop I = r$. We assume that $P$ is an embedded prime divisor
of $I$ and derive a contradiction. Localising at $P$ we can assume that $A$ is
a CM local ring; then by Theorem 4, (iii), $a_1, \ldots, a_r$ is an $A$-sequence, and
hence $A/I$ is also a CM local ring. But then $I$ does not have embedded
prime divisors, and this is a contradiction. Next we suppose that the
unmixedness theorem holds for $A$. If $P \in \Spec A$ with $\htop P = r$ then we can
choose $a_1, \ldots, a_r \in P$ such that
\[
  \htop(a_1, \ldots, a_i) = i \quad \text{for } 1 \le i \le r.
\]
Then by the unmixedness theorem, all the prime divisors of $(a_1, \ldots, a_i)$
have height $i$, and therefore do not contain $a_{i+1}$. Hence $a_{i+1}$ is an
$A/(a_1, \ldots, a_i)$-regular element; in other words, $a_1, \ldots, a_r$ is an $A$-sequence.
Therefore $\depth A_P = r = \dim A_P$, so that $A_P$ is a CM local ring; $P$ was
any element of $\Spec A$, so that $A$ is a CM ring.
\end{proof}

The unmixedness theorem for polynomial rings over a field was a
brilliant early result of Macaulay in 1916; for regular local rings, the
unmixedness theorem was proved by I. S. Cohen [1] in 1946. This explains
the term Cohen--Macaulay. Having come this far, we are now in a position
to give easy proofs of these two theorems.

\begin{theorem}\label{thm:17.7}
If $A$ is a CM ring then so is $A[X_1, \ldots, X_n]$.
\end{theorem}

\begin{proof}
We need only consider the case $n = 1$. Set $B = A[X]$ and let $P$ be
a maximal ideal of $B$. Set $P \cap A = m$; then $B_P$ is also a localisation of
$A_m[X]$, so that replacing $A$ by $A_m$ we have a local CM ring
$A$ with maximal ideal $m$, and we need only prove that $B_P$ is CM. Setting
$A/m = k$ we get
\[
  B/mB = k[X],
\]
so that $P/mB$ is a principal ideal of $k[X]$ generated by an irreducible
monic polynomial $\bar{\varphi}(X)$. If we let $f(X) \in A[X]$ be a monic polynomial of
$A[X]$ which reduces to $\varphi(X)$ modulo $mB$ then $P = (m, f)$. We choose a
system of parameters $a_1, \ldots, a_n$ for $A$, so that $a_1, \ldots, a_n, f$ is a system of
parameters of $B_P$. Since $B$ is flat over $A$ the $A$-sequence $a_1, \ldots, a_n$ is also
a $B$-sequence. We set $A/(a_1, \ldots, a_n) = A'$; then the image of $f$ in $A'[X]$ is
a monic polynomial, and therefore $A'[X]$-regular, so that $a_1, \ldots, a_n, f$ is
a $B$-sequence, and
\[
  \depth B_P \ge \depth (P, B) \ge n + 1 = \dim B_P.
\]
Therefore $B_P$ is a CM ring.
\end{proof}

\begin{remark}
If $A$ is a CM local ring, then a similar (if anything, rather easier)
method can be used to prove that $A[[X]]$ is also CM. The statement also
holds for a non-local CM ring, but the proof is a little more complicated,
and we leave it to \S23.
\end{remark}

\begin{theorem}\label{thm:17.8}
A regular local ring is a CM ring.
\end{theorem}

\begin{proof}
Let $(A, m)$ be an $n$-dimensional regular local ring, and $x_1, \ldots, x_n$ a
regular system of parameters. By Theorems 14.2 and 14.3, $(x_1), (x_1, x_2), \ldots,
(x_1, \ldots, x_n)$ is a strictly increasing chain of prime ideals; therefore $x_1, \ldots, x_n$
is an $A$-sequence.
\end{proof}

\begin{theorem}\label{thm:17.9}
Any quotient of a CM ring is universally catenary.
\end{theorem}

\begin{proof}
Clear from Theorems 7 and 4.
\end{proof}

\begin{theorem}\label{thm:17.10}
A necessary and sufficient condition for a Noetherian local
ring $(A, m, k)$ to be a regular ring is that $\gr_m(A)$ is isomorphic as a graded $k$-
algebra to a polynomial ring over $k$.
\end{theorem}

\begin{proof}
If $A$ is regular, let $x_1, \ldots, x_r$ be a regular system of parameters, that
is a minimal basis of $m$; then $x_1, \ldots, x_r$ is an $A$-sequence, so that by
Theorem~\ref{thm:16.2} (see also Theorem~\ref{thm:14.4} for another proof) $\gr_m(A) \cong
k[X_1, \ldots, X_r]$. Conversely, if $\gr_m(A) \cong k[X_1, \ldots, X_r]$, then comparing
the homogeneous components of degree $1$, we see that $m/m^2 \cong kX_1
+ \cdots + kX_r$. On the other hand, the homogeneous component of degree $n$
of $k[X_1, \ldots, X_r]$ is a vector space over $k$ of dimension $\binom{n+r-1}{r-1}$,
so that the Samuel function is
\[
  \chi_A(n) = l(A/m^{n+1}) = \sum_{i=0}^{n} l(m^i/m^{i+1}) = \binom{n+r}{r},
\]
and $\dim A = r$. Therefore $A$ is regular.
\end{proof}

We can also characterise CM local rings in terms of properties of
multiplicities. Let $A$ be a Noetherian local ring. An ideal of $A$ is said to
be a \emph{parameter ideal}\index{parameter ideal} if it can be generated by a system of parameters. By
Theorem~\ref{thm:14.10}, if $q$ is a parameter ideal then $l(A/q) \ge e(q)$. As we are about
to see, equality here is characteristic of CM rings.

\begin{theorem}\label{thm:17.11}
The following three conditions on a Noetherian local ring
$(A, m)$ are equivalent:
\begin{enumerate}[(1)]
  \item $A$ is a CM ring;
  \item $l(A/q) = e(q)$ for any parameter ideal $q$ of $A$;
  \item $l(A/q) = e(q)$ for some parameter ideal $q$ of $A$.
\end{enumerate}
\end{theorem}

\begin{proof}
$(1) \Rightarrow (2)$. If $x_1, \ldots, x_d$ is a system of parameters of $A$ and $q =
(x_1, \ldots, x_d)$ then by Theorem~\ref{thm:16.2}, $\gr_q(A) \cong (A/q)[X_1, \ldots, X_d]$, so that
as in the proof of the previous theorem, $\chi_A(n) = l(A/q)\binom{n+d}{d}$, so that
$e(q) = l(A/q)d!$.

$(2) \Rightarrow (3)$ is obvious.

$(3) \Rightarrow (1)$ Suppose that $q = (x_1, \ldots, x_d)$ is a parameter ideal satisfying
$e(q) = l(A/q)$. We set $B = (A/q)[X_1, \ldots, X_d]$; then there is a homogeneous
ideal $\mathfrak{b}$ of $B$ such that $\gr_q(A) \cong B/\mathfrak{b}$. We write $\varphi_B(n)$ and $\varphi_{\mathfrak{b}}(n)$ for the Hilbert
polynomials of $B$ and $\mathfrak{b}$ (see \S\ref{sec:13}); then
\[
  l(q^n/q^{n+1}) = \varphi_B(n) - \varphi_{\mathfrak{b}}(n),
\]
and for $n \gg 0$ we have $l(q^n/q^{n+1}) = \chi_A(n) - \chi_A(n-1)$. The left-hand side is a
polynomial in $n$ of degree $d - 1$, and the coefficient of $n^{d-1}$ is $e(q)/(d - 1)!$. By
hypothesis $e(q) = l(A/q)$, so that $\varphi_{\mathfrak{b}}(n)$ must be a polynomial in $n$ of degree at
most $d - 2$. However, if $\mathfrak{b} \neq (0)$ then we can take a non-zero homogeneous
element $f(X) \in \mathfrak{b}$. If $m' \subset q$ and we set $\bar{m} = m/q$ then in $B$ we have
$\bar{m}^\nu = (0)$, and therefore replacing $f$ by the product of $f$ with a suitable
element of $\bar{m}$, we can assume that $f \neq 0$ but $\bar{m}f = 0$. Then
if $f$ has degree $p$,
\[
  \varphi_{\mathfrak{b}}(n) \ge \binom{n-p+d-1}{d-1},
\]
the length of the homogeneous component of degree $n - p$ in $(A/m)[X_1, \ldots, X_d]$. This
contradicts $\deg \varphi_{\mathfrak{b}} < d - 1$. Hence $\mathfrak{b} = (0)$, and
\[
  \gr_q(A) = B = (A/q)[X_1, \ldots, X_d],
\]
so that by Theorem~\ref{thm:16.3}, $\{x_1, \ldots, x_d\}$ is an $A$-sequence. Therefore $A$ is a CM
ring.
\end{proof}

% ============ 习题 ============
\section*{Exercises to \S\ref{sec:17}}
\emph{Prove the following propositions.}

\begin{exercise}\label{ex:17.1-ex}
\begin{enumerate}[(a)]
  \item A zero-dimensional Noetherian ring is a CM ring.
  \item A one-dimensional ring is CM provided that it is reduced ($=$ no
        nilpotent elements); also, construct an example of a one-dimensional ring
        which is not CM.
\end{enumerate}
\end{exercise}

\begin{exercise}\label{ex:17.2-ex}
Let $k$ be a field, $x, y$ indeterminates over $k$, and set
\[
  A = k[x^3, x^2y, xy^2, y^3] \subset k[x, y] \quad \text{and} \quad P = (x^3, x^2y, xy^2, y^3)A.
\]
Is $R = A_P$ a CM ring? How about
$k[x^4, x^3y, xy^3, y^4]$?
\end{exercise}

\begin{exercise}\label{ex:17.3-ex}
A two-dimensional normal ring is CM.
\end{exercise}

\begin{exercise}\label{ex:17.4-ex}
Let $A$ be a CM ring, $a_1, \ldots, a_r$ an $A$-sequence, and set $J = (a_1, \ldots, a_r)$. Then
for every integer $\nu$ the ring $A/J^\nu$ is CM, and therefore $J^\nu$ is unmixed.
\end{exercise}

\begin{exercise}\label{ex:17.5-ex}
Let $A$ be a Noetherian local ring and $P \in \Spec A$. Then
\begin{enumerate}[(i)]
  \item $\depth A \le \depth (P, A) + \dim A/P$;
  \item call $\dim A - \depth A$ the \emph{codepth} of $A$. Then $\codepth A \ge
        \codepth A_P$.
\end{enumerate}
\end{exercise}

\begin{exercise}\label{ex:17.6-ex}
Let $A$ be a Noetherian ring, $P \in \Spec A$ and set $G = \gr_P(A)$. If $G$ is an
integral domain then $P^n = P^{(n)}$ for all $n > 0$. (This observation is due to
Robbiano. One sees from it that if $P$ is a prime ideal generated by an $A$-
sequence then $P^n = P^{(n)}$.)
\end{exercise}

% ============ §18 Gorenstein rings (原书 pp.139-152) ============
\section{Gorenstein rings}\label{sec:18}\index{Gorenstein rings, section 18}

\begin{lemma}\label{lem:18.1}
Let $A$ be a ring, $M$ an $A$-module, and $n > 0$ a given integer,
\[
  \injdim M \le n \iff \Ext^{n+1}_A(A/I, M) = 0 \text{ for all ideals } I.
\]
If $A$ is Noetherian, then we can replace `for all ideals' by `for all prime
ideals' in the right-hand condition.
\end{lemma}

\begin{proof}
($\Rightarrow$) This is clear on calculating the Ext by an injective resolution
of $M$.

($\Leftarrow$) If $n = 0$ then from the exact sequence $0 \to I \to A \to A/I \to 0$ and
from the fact that $\Ext^1_A(A/I, M) = 0$ we get that $\Hom(A, M) \to
\Hom(I, M) \to 0$ is exact. Since this holds for every $I$, Theorem B3 of
Appendix B implies that $M$ is injective. Suppose then that $n > 0$.
There exists an exact sequence
\[
  0 \to M \to Q^0 \to Q^1 \to \cdots \to Q^{n-1} \to C \to 0,
\]
with each $Q^i$ injective. (We can obtain this by taking an injective resolution
of $M$ up to $Q^{n-1}$ and setting $C$ for the cokernel of $Q^{n-2} \to Q^{n-1}$.)
One sees easily that $\Ext^{n+1}_A(A/I, M) \cong \Ext^1_A(A/I, C)$, so that by the argument
used in the $n = 0$ case, $C$ is injective, and so $\injdim M \le n$.

If $A$ is Noetherian then by Theorem~\ref{thm:6.4}, any finite $A$-module $N$ has a
chain $N = N_0 \supset N_1 \supset \cdots \supset N_{r+1} = 0$ of submodules such that $N_j/N_{j+1} \cong
A/P_j$, with $P_j \in \Spec A$. Using this, if $\Ext^i_A(A/P, M) = 0$ for all prime ideals
$P$ then we also have $\Ext^i_A(N, M) = 0$ for all finite $A$-modules $N$. Now we
just have to apply this with $i = n + 1$ and $N = A/I$.
\end{proof}

\begin{lemma}\label{lem:18.2}
Let $A$ be a ring, $M$ and $N$ two $A$-modules, and $x \in A$;
suppose that $x$ is both $A$-regular and $M$-regular, and that $xN = 0$. Set $B = A/xA$ and
$\bar{M} = M/xM$. Then
\begin{enumerate}[(i)]
  \item $\Hom_B(N, \bar{M}) = 0$, and $\Ext^{n+1}_A(N, M) \cong \Ext^n_B(N, \bar{M})$ for all $n \ge 0$;
  \item $\Ext^n_A(M, N) \cong \Ext^n_B(\bar{M}, N)$ for all $n \ge 0$;
  \item $\Tor^A_n(M, N) \cong \Tor^B_n(\bar{M}, N)$ for all $n \ge 0$.
\end{enumerate}
\end{lemma}

\begin{proof}
(i) The first formula is obvious. For the second, set $T^n(N) =
\Ext^{n+1}_A(N, M)$, and view $T^n$ as a contravariant functor from the category
of $B$-modules to that of Abelian groups. Then first of all, the exact sequence
\[
  0 \to M \xrightarrow{x} M \to \bar{M} \to 0
\]
gives $T^0(N) = \Hom_A(N, M) = \Hom_B(N, \bar{M})$. Moreover, since $x$ is
$A$-regular we have $\projdim_A B = 1$, and therefore $T^n(B) = 0$ for $n > 0$,
so that $T^n(L) = 0$ for $n > 0$ and every projective $B$-module $L$. Finally, for
any short exact sequence $0 \to N' \to N \to N'' \to 0$ of $B$-modules, there is
a long exact sequence
\[
  0 \to T^0(N'') \to T^0(N) \to T^0(N')
\]
\[
  \to T^1(N'') \to T^1(N) \to T^1(N') \to \cdots .
\]
This proves that $T^\bullet$ is the derived functor of $\Hom_B(-, \bar{M})$, and therefore
coincides with $\Ext^\bullet_B(-, \bar{M})$.

(ii) We first prove $\Tor^A_n(M, B) = 0$ for $n > 0$. For $n > 1$ this follows
from $\projdim_A B = 1$. For $n = 1$, consider the long exact sequence
$0 \to \Tor^A_1(M, B) \to M \xrightarrow{x} M \to M \otimes_A B \to 0$ associated with the short
exact sequence $0 \to A \xrightarrow{x} A \to B \to 0$. Since $x$ is $M$-regular we have
$\Tor^A_1(M, B) = 0$.

Now let $L_\bullet \to M \to 0$ be a free resolution of the $A$-module $M$. Then
$L_\bullet \otimes_A B \to M \otimes_A B \to 0$ is exact by what we have just proved, so that
$L_\bullet \otimes_A B$ is a free resolution of the $B$-module $M \otimes_A B = \bar{M}$. Then $\Ext^n_A(M, N)
= H^n(\Hom_A(L_\bullet, N)) = H^n(\Hom_B(L_\bullet \otimes_A B, N)) = \Ext^n_B(\bar{M}, N)$ by Formula 9
of Appendix A.

(iii) Using the same notation as above, we have $\Tor^A_n(M, N) =
H_n(L_\bullet \otimes_A N) = H_n((L_\bullet \otimes_A B) \otimes_B N) = \Tor^B_n(\bar{M}, N)$.
\end{proof}

\begin{lemma}\label{lem:18.3}
Let $(A, m, k)$ be a Noetherian local ring, $M$ a finite $A$-module, and
$P \in \Spec A$; suppose that $\htop (m/P) = 1$. Then
\[
  \Ext^{r+1}_A(k, M) = 0 \Rightarrow \Ext^r_A(A/P, M) = 0.
\]
\end{lemma}

\begin{proof}
Choose $x \in m - P$; then $0 \to A/P \xrightarrow{x} A/P \to A/(P + Ax) \to 0$ is
an exact sequence, and $P + Ax$ is an $m$-primary ideal, so that if we let
$N = A/(P + Ax)$, there exists a chain of submodules of $N$
\[
  N = N_0 \supset N_1 \supset \cdots \supset N_s = 0 \quad \text{such that } N_i/N_{i+1} \cong k.
\]
Hence from $\Ext^{r+1}_A(k, M) = 0$ we get $\Ext^{r+1}_A(A/(P + Ax), M) = 0$, and
\[
  \Ext^r_A(A/P, M) \xrightarrow{x} \Ext^r_A(A/P, M) \to 0
\]
is exact, so that by NAK $\Ext^r_A(A/P, M) = 0$.
\end{proof}

\begin{lemma}\label{lem:18.4}
Let $(A, m, k)$ be a Noetherian local ring, $M$ a finite $A$-module, and
$P \in \Spec A$; suppose that $\htop (m/P) = d$. Then
\[
  \Ext^{r+d}_A(k, M) = 0 \Rightarrow \Ext^r_{A_P}(\kappa(P), M_P) = 0.
\]
\end{lemma}

\begin{proof}
Let $m = P_0 \supset P_1 \supset \cdots \supset P_d = P$, with $P_i \in \Spec A$ and $\htop (P_i/P_{i+1}) = 1$;
then by the previous lemma, from $\Ext^{r+d}_A(k, M) = 0$ we get
\[
  \Ext^{r+d-1}_A(A/P_1, M) = 0,
\]
and localising at $P_1$ we get
\[
  \Ext^{r+d-1}_{A_{P_1}}(\kappa(P_1), M_{P_1}) = 0.
\]
Proceeding in the same way gives the result.
\end{proof}

\begin{theorem}\label{thm:18.1}
Let $(A, m, k)$ be an $n$-dimensional Noetherian local ring.
Then the following conditions are equivalent:
\begin{enumerate}[(1)]
  \item $\injdim A < \infty$;
  \item $\injdim A = n$;
  \item $\Ext^i_A(k, A) = 0$ for $i \neq n$ and $\cong k$ for $i = n$;
  \item $\Ext^i_A(k, A) = 0$ for some $i > n$;
  \item $\Ext^i_A(k, A) = 0$ for $i < n$ and $\cong k$ for $i = n$;
  \item[($4'$)] $A$ is a CM ring and $\Ext^n_A(k, A) \cong k$;
  \item[($5$)] $A$ is a CM ring, and every parameter ideal of $A$ is irreducible;
  \item[($5'$)] $A$ is a CM ring and there exists an irreducible parameter ideal.
\end{enumerate}
Recall that an ideal $I$ is \emph{irreducible} if $I = J \cap J'$ implies either $I = J$ or
$I = J'$ (see \S\ref{sec:6}).
\end{theorem}

\begin{definition}
A Noetherian local ring for which the above equivalent
conditions hold is said to be \emph{Gorenstein}.\index{Gorenstein ring}
\end{definition}

\begin{proof}[Proof of $(1) \Rightarrow (1')$]
Set $\injdim A = r$. If $P$ is a minimal prime ideal of $A$
such that $\htop (m/P) = \dim A = n$ then $PA_P \in \Ass(A_P)$, so that
$\Hom_{A_P}(\kappa(P), A_P) \neq 0$; hence, by Lemma~\hyperref[lem:18.4]{4}, $\Ext^n_A(k, A) \neq 0$, therefore $r \ge n$.
If $r = 0$ this means that $n = 0$, and we are done. If $r > 0$, set $\Ext^r_A(-, A) =
T$; then this is a right-exact contravariant functor, and by Lemma~\hyperref[lem:18.1]{1},
there is a prime ideal $P$ such that $T(A/P) \neq 0$. Now if $P \neq m$ and we take
$x \in m - P$, the exact sequence
\[
  0 \to A/P \xrightarrow{x} A/P
\]
leads to an exact sequence
\[
  T(A/P) \xrightarrow{x} T(A/P) \to 0;
\]
but then by NAK, $T(A/P) = 0$, which is a contradiction. Thus $P = m$, and
so $T(k) \neq 0$. We have $m \notin \Ass(A)$, since otherwise there would exist an
exact sequence $0 \to k \to A$, and hence an exact sequence
\[
  T(A) = \Ext^r_A(A, A) = 0 \to T(k) \to 0,
\]
which is a contradiction. Hence $m$ contains an $A$-regular element $x$. If we set
$B = A/xA$ then by Lemma~\hyperref[lem:18.2]{2}, $\Ext^n_B(N, B) = \Ext^{n+1}_A(N, A)$ for every $B$-
module $N$, so that $\injdim B = r - 1$. By induction on $r$ we have $r - 1 =
\dim B = n - 1$, and hence $r = n$.
\end{proof}

\begin{proof}[Proof of $(1') \Rightarrow (2)$]
When $n = 0$ we have $m \in \Ass(A)$, so there exists an exact
sequence $0 \to k \to A$, and since $\injdim A = 0$,
\[
  A = \Hom(A, A) \to \Hom(k, A) \to 0
\]
is exact. Therefore $\Hom(k, A)$ is generated by one element. But
$\Hom (k, A) \neq 0$, so that we must have $\Hom (k, A) \cong k$. By assumption, $A$ is
an injective module, so that $\Ext^i_A(k, A) = 0$ for $i > 0$; thus we are done in the
case $n = 0$. If $n > 0$ then, as we have seen above, $m$ contains an $A$-regular
element $x$, and if we set $B = A/xA$ then $\dim B = \injdim B = n - 1$, so that by
Lemma~\hyperref[lem:18.2]{2} and induction on $n$ we have
\[
  \Ext^i_A(k, A) = \Ext^{i-1}_B(k, B) =
  \begin{cases}
    0 & \text{if } 0 < i \neq n \\
    k & \text{if } i = n
  \end{cases}
\]
and
\[
  \Hom_A(k, A) = 0. \qquad\blacksquare
\]
\end{proof}

\begin{proof}[Proof of $(3) \Rightarrow (1)$]
We use induction on $n$. Assume that for some $i > n$ we have
$\Ext^i_A(k, A) = 0$. If $n = 0$ then $m$ is the unique prime ideal of $A$, so that by
Lemma~\hyperref[lem:18.1]{1}, $\injdim A \le i < \infty$. If $n > 0$ let $P$ be a prime ideal distinct
from $m$, and set $d = \htop (m/P)$ and $B = A_P$; then by Lemma~\hyperref[lem:18.4]{4} we have
$\Ext^{i-d}_B(\kappa(P), B) = 0$. Moreover, $\dim B \le n - d < i - d$, so that by induction
$\injdim B < \infty$. Thus for any finite $A$-module $M$ we have
\[
  (\Ext^i_A(M, A))_P = \Ext^i_B(M_P, B) = 0
\]
for $i > n \ge \dim B = \injdim B$). Therefore, setting $T(M) = \Ext^i_A(M, A)$,
we get $\Supp(T(M)) = \{m\}$, and since $T(M)$ is a finite $A$-module,
$T(M) \neq 0$ implies that $T(M) \otimes_A k \neq 0$. Using this, we now prove that $T(A/P) = 0$ for every prime
ideal $P$. If $T(A/P) \neq 0$ for some $P$, choose a maximal $P$ with this property. By
assumption $T(k) = 0$, so that $P \neq m$, so that we can take $x \in m - P$ and form
\[
  0 \to A/P \xrightarrow{x} A/P \to A/(P + Ax) \to 0.
\]
We write $A/(P + Ax) = M_0 \supset M_1 \supset \cdots \supset M_s = 0$ with $M_l/M_{l+1} \cong A/P_l$;
then $P_l$ is strictly bigger than $P$, so that $T(A/(P + Ax)) = 0$. Therefore
\[
  T(A/P) \xrightarrow{x} T(A/P) \to 0
\]
so that multiplication by $x$ in $T(A/P)$ is injective; but since
$l(T(A/P)) < \infty$, injective implies surjective. Hence by NAK, $T(A/P) = 0$,
which is a contradiction. Therefore $T(A/P) = 0$ for every $P \in \Spec A$, so that
$\injdim A < \infty$ by Lemma~\hyperref[lem:18.1]{1}.

Thus we have proved that $(1), (1'), (2)$ and $(3)$ are equivalent. Now we
prove the equivalence of $(2), (4), (4'), (5)$ and $(5')$.

$(4) \Rightarrow (2)$ is obvious. $(4) \Leftrightarrow (4')$ comes at once from the fact that $A$ is CM if and
only if $\Ext^i_A(k, A) = 0$ for all $i < n$ (the implication $(2) \Rightarrow (3)$ of
Theorem~\ref{thm:16.6}).

$(4') \Rightarrow (5)$. A system of parameters $x_1, \ldots, x_n$ in a CM ring $A$ is an $A$-
sequence, so that setting $B = A/\sum x_iA$, we have
\[
  \Hom_A(k, B) \cong \Ext^n_A(k, A) \cong k.
\]
Then $B$ is an Artinian ring, and any minimal non-zero ideal of $B$ is
isomorphic to $k$, so that the above formula says that $B$ has just one such
ideal, say $I_0$. If $I_1$ and $I_2$ are any non-zero ideals of $B$ then both of
them contain $I_0$, so that $I_1 \cap I_2 \neq (0)$. Lifting this up to $A$, this means
that $(x_1, \ldots, x_n)$ is an irreducible ideal.

$(5') \Rightarrow (2)$. If $A$ is CM we already have $\Ext^i_A(k, A) = 0$ for $i < n$. If $q$ is
an irreducible parameter ideal and we set $B = A/q$ then, in the same way as
above,
\[
  \Ext^{n+1}_A(k, A) \cong \Ext^1_B(k, B),
\]
so that it is enough to prove that in an Artinian ring $B$, $(0)$ is irreducible
implies that
\[
  \Hom_B(k, B) \cong k \quad \text{and} \quad \Ext^i_B(k, B) = 0 \text{ for } i > 0.
\]
The statement for Hom is easy: first of all, $B$ is Artinian, so that
$\Hom_B(k, B) \neq 0$; for non-zero $f, g \in \Hom_B(k, B)$ we must have $f(k) = g(k)$,
since otherwise $f(k) \cap g(k) = (0)$, which contradicts the irreducibility of $(0)$.
Hence $f(1) = g(\alpha)$ for some $\alpha \in k$, and $f = \alpha g$, so that $\Hom_B(k, B) = k$.

Now consider the $\Ext^i_B(k, B)$. Choose a chain of ideals $(0) = N_0 \subset
N_1 \subset \cdots \subset N_r = B$ such that $N_i/N_{i-1} \cong k$, and consider the exact sequences
\[
  0 \to N_1 \to N_2 \to k \to 0
\]
\[
  0 \to N_2 \to N_3 \to k \to 0
\]
\[
  \vdots
\]
\[
  0 \to N_{r-1} \to B \to k \to 0.
\]
From the long exact sequence
\[
  0 \to \Hom_B(k, B) \to \Hom_B(N_{i+1}, B) \to \Hom_B(N_i, B) \xrightarrow{\delta_i}
  \Ext^1_B(k, B) \to \cdots
\]
and an easy induction (using $N_1 \cong k$ and $\Hom_B(k, B) \cong k$), we get that
$l(\Hom_B(N_i, B)) \le i$, with equality holding if and only if $\delta_1, \ldots, \delta_{i-1}$ are all
zero. However,
\[
  l(\Hom_B(N_r, B)) = l(\Hom_B(B, B)) = l(B) = r,
\]
so that we must have $\delta_1 = \cdots = \delta_{r-1} = 0$. Then from
\[
  0 \to N_{r-1} \to B \to k \to 0
\]
we get the exact sequence
\[
  0 \to \Ext^1_B(k, B) \to \Ext^1_B(B, B) = 0,
\]
and therefore $\Ext^1_B(k, B) = 0$. Now from Lemma~\hyperref[lem:18.1]{1}, $B$ is an injective $B$-
module, so that $\Ext^i_B(k, B) = 0$ for all $i > 0$.
\end{proof}

\begin{lemma}\label{lem:18.5}
Let $A$ be a Noetherian ring, $S \subset A$ a multiplicative set, and $I$ an
injective $A$-module; then $I_S$ is an injective $A_S$-module.
\end{lemma}

\begin{proof}
Every ideal of $A_S$ is the localisation $\mathfrak{a}_S$ of an ideal $\mathfrak{a}$ of $A$. From
$0 \to \mathfrak{a} \to A$ we get the exact sequence $\Hom_A(A, I) \to \Hom_A(\mathfrak{a}, I) \to 0$,
and, since $\mathfrak{a}$ is finitely generated
\[
  \Hom_{A_S}(A_S, I_S) \to \Hom_{A_S}(\mathfrak{a}_S, I_S) \to 0
\]
is exact. This proves that $I_S$ is an injective $A_S$-module.
\end{proof}

\begin{theorem}\label{thm:18.2}
If $A$ is a Gorenstein local ring and $P \in \Spec A$ then $A_P$ is also
Gorenstein.
\end{theorem}

\begin{proof}
If $0 \to A \to I^0 \to I^1 \to \cdots \to I^n \to 0$ is an injective reso-
lution of $A$, then localising at $P$ and applying the previous lemma gives an
injective resolution of $A_P$, so that $\injdim A_P \le n$.
\end{proof}

\begin{definition}
A Noetherian ring $A$ is \emph{Gorenstein} if its localisation at every
maximal ideal is a Gorenstein local ring. (By the previous theorem, it then
follows that $A_P$ is Gorenstein for every $P \in \Spec A$.)
\end{definition}

\begin{theorem}\label{thm:18.3}
Let $A$ be a Noetherian local ring and $\hat{A}$ its completion.
Then $A$ is Gorenstein $\iff \hat{A}$ is Gorenstein.
\end{theorem}

\begin{proof}
We have $\dim A = \dim \hat{A}$, and since $\hat{A}$ is faithfully flat over $A$,
$\Ext^i_A(k, A) \otimes_A \hat{A} = \Ext^i_{\hat{A}}(k, \hat{A})$, so that we only need to use condition (3)
of Theorem 1.
\end{proof}

Closely related to the theory of Gorenstein rings is Matlis' theory
of injective modules over Noetherian rings. We now discuss the main
results of Matlis~\cite{Matl1}.

Let $A$ be a Noetherian ring, and $E$ an injective $A$-module. If $E$ is a
submodule of an $A$-module $M$ then since we can extend the identity map
$1_E$ to a linear map $f\colon M \to E$, we have $M = E \oplus F$ (with $F = \Ker f$).
An $A$-module $N$ is \emph{indecomposable}\index{indecomposable module} if $N$ cannot be written as a direct
sum of two submodules. We write $E_A(N)$ or $E(N)$ for the \emph{injective hull}\index{injective hull}
of an $A$-module $E$ (see Appendix B).

\begin{theorem}\label{thm:18.4}
Let $A$ be a Noetherian ring and $P, Q \in \Spec A$.
\begin{enumerate}[(i)]
  \item $E(A/P)$ is indecomposable.
  \item Any indecomposable injective $A$-module is of the form $E(A/P)$ for
        some $P \in \Spec A$.
  \item If $x \in A - P$ then multiplication by $x$ induces an automorphism of $E(A/P)$.
  \item $P \nsubseteq Q \iff E(A/P) \to E(A/Q)$ is injective.
  \item For any $\xi \in E(A/P)$ there exists a positive integer $\nu$ (depending on $\xi$)
        such that $P^\nu\xi = 0$.
  \item If $Q \subset P$ then $E(A/Q)$ is an $A_P$-module, and is an injective hull of
        $(A/Q)_P = A_P/QA_P$, that is
        \[
          E_A(A/Q) = E_{A_P}(A_P/QA_P).
        \]
\end{enumerate}
\end{theorem}

\begin{proof}
(i) If $I_1$ and $I_2$ are non-zero ideals of $A/P$ then $0 \neq I_1I_2 \subset I_1 \cap I_2$.
Since $E(A/P)$ is an essential extension of $A/P$ (see Appendix B), so that for
any two non-zero submodules $N_1, N_2$ of $E(A/P)$ we have $N_i \cap (A/P) \neq 0$, so
that
\[
  N_1 \cap N_2 \supset (N_1 \cap A/P) \cap (N_2 \cap A/P) \neq 0.
\]
(ii) Let $N \neq 0$ be an indecomposable injective $A$-module and choose
$P \in \Ass(N)$. Then $A/P$ can be embedded into $N$, and so $E(A/P)$ can also; but
an injective submodule is always a direct summand, and since $N$ is
indecomposable, $N = E(A/P)$.

(iii) Write $\varphi$ for multiplication by $x$ in $E(A/P)$; then $\Ker(\varphi) \cap (A/P) = 0$, so
that $\Ker(\varphi) = 0$, and $\Im \varphi$ is isomorphic to $E(A/P)$. Hence $\Im \varphi$ is injective,
and is therefore a direct summand of $E(A/P)$, so that by (i), $\Im \varphi = E(A/P)$.

(iv) If $P \nsubseteq Q$ and $x \in P - Q$ then multiplication by $x$ is injective in
$E(A/Q)$ but not in $E(A/P)$.

(v) By the proof of (ii) together with (iv), $\Ass(E(A/P)) = \{P\}$, so that the
submodule $A\xi \cong A/\ann(\xi)$ also has $\Ass(A\xi) = \{P\}$. Hence $\ann(\xi)$ is a
$P$-primary ideal.

(vi) By (iii), we can view $E(A/Q)$ as an $A_P$-module; hence it contains
$(A/Q)_P$. Since $E(A/Q)$ is an essential extension of $A/Q$ and $A/Q \subset (A/Q)_P \subset
E(A/Q)$, it is also an essential extension of $(A/Q)_P$. For $A_P$-modules $M$ and $N$,
any $A$-linear map $M \to N$ is also $A_P$-linear, and of course conversely,
so that for an $A_P$-module, being injective as an $A_P$-module is the same as
being injective as an $A$-module. Thus $E(A/Q)$ is an injective hull of the
$A_P$-module $(A/Q)_P$.
\end{proof}

\begin{example}\label{ex:18.1}
If $A$ is an integral domain and $K$ its field of fractions, $K = E(A)$
(prove this!).
\end{example}

\begin{example}\label{ex:18.2}
If $A$ is a DVR with uniformising element $x$ and field of fractions
$K$, and $k = A/xA$, then $E(k) = K/A$. Indeed, if $I$ is a non-zero ideal of $A$
we can write $I = x^rA$, and if $f\colon I \to K/A$ is a given map, let $f(x^r) = \alpha
\bmod A$ for some $\alpha \in K$; then $f$ can be extended to a map $\bar{f}\colon A \to K/A$
by setting $\bar{f}(1) = (\alpha/x^r) \bmod A$. Therefore $K/A$ is injective. We have
$(x^{-1}A)/A \cong A/xA = k$, and it is easy to see that $K/A$ is an essential extension
of $x^{-1}A/A$. Thus $K/A$ can be thought of as $E(k)$.
\end{example}

\begin{theorem}\label{thm:18.5}
We consider modules over a Noetherian ring $A$.
\begin{enumerate}[(i)]
  \item A direct sum of any number of injective modules is injective.
  \item Every injective module is a direct sum of indecomposable injective
        modules.
  \item The direct sum decomposition in (ii) is unique, in the sense that if
        \[
          M = \bigoplus M_i \quad \text{(with indecomposable } M_i\text{)}
        \]
        then for any $P \in \Spec A$, the sum $M(P)$ of all the $M_i$ isomorphic to $E(A/P)$
        depends only on $M$ and $P$, and not on the decomposition $M = \bigoplus M_i$.
        Moreover, the number of $M_i$ isomorphic to $E(A/P)$ is equal to
        \[
          \dim_{\kappa(P)} \Hom_{A_P}(\kappa(P), M_P), \quad \text{(where } \kappa(P) = A_P/PA_P),
        \]
        so that this also is independent of the decomposition.
\end{enumerate}
\end{theorem}

\begin{proof}
(i) Let $M_\lambda$ for $\lambda \in \Lambda$ be injective modules. It is enough to prove that
for an ideal $I$ of $A$, any linear map $\varphi\colon I \to \bigoplus M_\lambda$ can be extended to a linear
map $\bar{\varphi}\colon A \to \bigoplus M_\lambda$. Since $I$ is finitely generated, $\varphi(I)$ is contained
in a finite direct sum $M_{\lambda_1} \oplus \cdots \oplus M_{\lambda_n}$ of the $M_\lambda$. If $\varphi_\mu$ is the $\mu$th component
of $\varphi(\mathfrak{a})$, then $\varphi_\mu\colon I \to M_\mu$ extends
to $\bar{\varphi}_\mu\colon A \to M_\mu$, and we just set $\bar{\varphi}(1) = \bar{\varphi}_1(1) + \cdots + \bar{\varphi}_n(1)$.

(ii) Say that a family $\mathcal{F} = \{E_\lambda\}$ of indecomposable injective submodules
$E_\lambda$ of $M$ is \emph{free} if the sum in $M$ of the $E_\lambda$ is direct, that is if, for any finite
number $E_{\lambda_1}, \ldots, E_{\lambda_n}$ of them,
\[
  E_{\lambda_n} \cap (E_{\lambda_1} + \cdots + E_{\lambda_{n-1}}) = 0.
\]
Let $\mathfrak{M}$ be the set of all free families $\mathcal{F}$, ordered by inclusion. Then by
Zorn's lemma $\mathfrak{M}$ has a maximal element, say $\mathcal{F}_0$. Write $N = \sum_{E \in \mathcal{F}_0} E$; then
by (i), $N$ is injective, hence a direct summand of $M$, and $M = N \oplus N'$. If
$N' \neq 0$ then since it is a direct summand of $M$ it must be injective, and
for $P \in \Ass(N')$, the proof of Theorem 4, (ii), shows that $N'$ contains a
direct summand $E'$ isomorphic to $E(A/P)$. Thus $\mathcal{F}_0 \cup \{E'\}$ is a free family,
contradicting the maximality of $\mathcal{F}_0$. Hence $N' = 0$ and $M = N$.

(iii) If we can show that $M(P)$ has the property that every submodule
$E$ of $M$ isomorphic to $E(A/P)$ is contained in $M(P)$, then $M(P)$ is the
submodule of $M$ generated by all such $E$, and therefore is determined by
$M$ and $P$ only. To prove this, take any $\xi \in E$; we can write $\xi = \xi_1 + \cdots +
\xi_r$ with $\xi_i \in M(P_i)$, where $P_1, \ldots, P_r$ are distinct prime ideals and $P = P_1$.
Setting $\eta_1 = \xi - \xi_1$, and $\eta_i = \xi_i$ for $2 \le i \le r$ we have $\eta_1 + \cdots + \eta_r = 0$, with
$\eta_1 \in M(P_1) + E$ and $\eta_i \in M(P_i)$ for $i \ge 2$. We need only prove that in this case
each $\eta_i = 0$. Suppose that $P_1$ is minimal among $P_1, \ldots, P_r$; then for any $m$
we have $(P_1 \cdots P_r)^m \nsubseteq P_1$, so that taking $a \in (P_1 \cdots P_{r-1})^m - P_1$ and $m$
large enough, we get $a\eta_1 = \cdots = a\eta_{r-1} = 0$. Then also $\eta_r = 0$, but
multiplication by $a$ is an automorphism of $M(P_r)$, so that $\eta_r = 0$. By
induction on $r$ we get $\eta_i = 0$ for all $i$.

We now prove that if $M(P) = M_1 \oplus \cdots \oplus M_s$ with $M_i \cong E(A/P)$ then
\[
  s = \dim_{\kappa(P)} \Hom_{A_P}(\kappa(P), M_P).
\]
(We are writing this as if $s$ were finite, but, as one can see from the proof
below, the same works for any cardinal number.) By Theorem 4, (vi), both
sides of $M(P) = M_1 \oplus \cdots \oplus M_s$ are $A_P$-modules, and $M_i = E(\kappa(P))$.
Moreover, by Theorem 4, (v), $E(A/Q)_P = 0$ if $Q \nsubseteq P$, so that
\[
  M_P = \bigoplus_{Q \subset P} M(Q)_P = \bigoplus_{Q \subset P} M(Q).
\]
Hence we can replace $A$ by $A_P$, and assume that $A$ is a local ring with $P$
its maximal ideal; set $k = \kappa(P)$. If $Q \neq P$ then any $x \in P - Q$ gives an
automorphism of $M(Q)$, but $x \cdot k = 0$, so that $\Hom_A(k, M(Q)) = 0$. Hence
$\Hom_A(k, M) = \Hom_A(k, M(P))$, so that there is no loss of generality in
assuming that $M = M(P)$. For any $A$-module $N$, we can identify
$\Hom_A(k, N)$ with the submodule $\{\xi \in N \mid P\xi = 0\}$, but since $E(k)$ is an
essential extension of $k$ we must have $\dim_k \Hom_A(k, E(k)) = 1$, so that if
$M = M_1 \oplus \cdots \oplus M_s$ with $M_i \cong E(k)$ then $s = \dim_k \Hom_A(k, M)$.
\end{proof}

\begin{theorem}\label{thm:18.6}
Let $(A, m, k)$ be a Noetherian local ring, and $E = E_A(k)$ the
injective hull of $k$. For each $A$-module $M$ set $M' = \Hom_A(M, E)$.
\begin{enumerate}[(i)]
  \item If $M$ is an $A$-module and $0 \neq x \in M$, then there exists $\varphi \in M'$ such that
        $\varphi(x) \neq 0$. In other words the canonical map $\theta\colon M \to M''$ defined by $\theta(x)(\varphi)
        = \varphi(x)$ for $x \in M$ and $\varphi \in M'$ is injective.
  \item If $M$ is an $A$-module of finite length, then $l(M) = l(M')$ and the
        canonical map $M \to M''$ is an isomorphism.
  \item Let $\hat{A}$ be the completion of $A$; then $E$ is also an $\hat{A}$-module, and
        is an injective hull of $k$ as $\hat{A}$-module.
  \item $\Hom_A(E, E) = \Hom_{\hat{A}}(E, E) = \hat{A}$. In other words, each endomor-
        phism of the $A$-module $E$ is multiplication by a unique element of $\hat{A}$.
  \item $E$ is Artinian as an $A$-module and also as an $\hat{A}$-module. Assume
        now that $A$ is complete, and write $\mathcal{N}$ (resp.\ $\mathcal{A}$) for the category of
        Noetherian (respectively Artinian) $A$-modules. Then if $M \in \mathcal{A}$ we have
        $M' \in \mathcal{N}$ and $M \cong M''$; if $M \in \mathcal{N}$ we have $M' \in \mathcal{A}$ and $M \cong M''$.
\end{enumerate}
\end{theorem}

\begin{proof}
(i) Let $f\colon Ax \to E$ be the composite of the canonical maps
$Ax \cong A/\ann(x)$, $A/\ann(x) \to A/m = k$ and $k \to E$. Then $f(x) \neq 0$. Since $E$ is
injective we can extend $f$ to $\varphi\colon M \to E$.

(ii) If $l(M) = n < \infty$ then $M$ has a submodule $M_1$ of length $n - 1$, and
$0 \to M_1 \to M \to k \to 0$ is exact, so that $0 \to k' \to M' \to M_1' \to 0$ is exact.
However
\[
  k' = \Hom_A(k, E) = \Hom_A(k, k) \cong k,
\]
so that by induction on $n$ we get $l(M) = n = l(M')$. The canonical map
$M \to M''$ is injective by (i), and $l(M) = l(M') = l(M'')$, hence it must be an
isomorphism.

(iii) Each element of $E$ is annihilated by some power of $m$, so that the
canonical map $E \to E \otimes_A \hat{A}$ is surjective. However, since $\hat{A}$ is faithfully
flat over $A$ it is also injective, so that $E \cong E \otimes_A \hat{A}$, and we can view $E$ as
an $\hat{A}$-module. Let $F$ be the injective hull of $E$ as an $\hat{A}$-module. Then $F$
is also the $A$-injective hull of $k$, so that every element of $F$ is annihilated
by some power of $m\hat{A}$. As an $A$-module $F$ splits into a direct sum of $E$
and an $\hat{A}$-module $C$. If $x \in C$, and if $m^\nu Ax = 0$, then for each $\hat{a} \in \hat{A}$ we
can find $a \in A$ such that $\hat{a} \equiv a \bmod m^\nu\hat{A}$ and hence $\hat{a}x = ax \in C$.
Therefore $C$ is an $A$-module. But $F$ is indecomposable as an $\hat{A}$-module.
Hence $C = 0$ and $E = F$.

(iv) For $\nu > 0$ set $E_\nu = \{x \in E \mid m^\nu x = 0\}$. Then we have $(A/m^\nu)' \cong
\Hom_A(A/m^\nu, E) \cong E_\nu$, and $\Hom_A(E_\nu, E_\nu) = \Hom_A(E_\nu, E) = E'_\nu = (A/m^\nu)'' =
A/m^\nu$. Now $E_1 \subset E_2 \subset \cdots$ and $E = \bigcup_\nu E_\nu$ by Theorem 4, (v), hence $E =
\varinjlim E_\nu$. Therefore $\Hom_A(E, E) = \Hom_A(\varinjlim E_\nu, E) = \varprojlim \Hom_A(E_\nu, E) =
\varprojlim A/m^\nu = \hat{A}$.

(v) If an $A$-module $M$ is Artinian and $x \in M$, then $Ax \cong A/\ann(x)$ is also
Artinian and consequently $m^\nu \subset \ann(x)$ for some $\nu$. Therefore $M$ can be
viewed as an $\hat{A}$-module, and its $\hat{A}$-submodules are precisely its $A$-
submodules. It is also clear that if an $\hat{A}$-module $M$ is Artinian then we have
the same conclusion. Therefore to prove (v) we may assume that $A$ is
complete.

If $M$ is a submodule of $E$ set $M^\perp = \{a \in A \mid aM = 0\}$. If $I$ is an ideal of $A$ set
$I^\perp = \{x \in E \mid Ix = 0\}$. Then clearly $M^{\perp\perp} \supset M$. If $x \in E - M$ there exist
$\varphi \in (E/M)'$ satisfying $\varphi(x \bmod M) \neq 0$ by (i), and if we identify $E' =
\Hom_A(E, E)$ with $A$ then $(E/M)'$ is identified with $M^\perp$. Thus $\varphi(x \bmod M) =
ax$ for some $a \in M^\perp$, and $x \notin M^{\perp\perp}$. Therefore $M^{\perp\perp} = M$. Similarly, if
$a \in \hat{A} - I$ then there exists $\varphi \in (A/I)'$ such that $\varphi(a \bmod I) \neq 0$, and
$(A/I)'$ is identified with the submodule $I^\perp$ of $E = A'$. Thus, setting $x =
\varphi(1 \bmod I)$ we have $x \in I^\perp$ and $ax = \varphi(a \bmod I) \neq 0$. This proves $a \notin I^{\perp\perp}$, so
that $I = I^{\perp\perp}$. Thus $M \mapsto M^\perp$ is an order-reversing bijection from the set of
submodules of $E$ onto the set of ideals of $A$. Since $A$ is Noetherian, it follows
that $E$ is Artinian. By Theorem~\ref{thm:3.1} finite direct sums $E^n$ of $E$ are also
Artinian for all $n > 0$.

If $M \in \mathcal{N}$ then there is a surjection $A^n \to M$ for some $n$, and so
there is an injection $M' \hookrightarrow (A^n)' = E^n$. Hence $M'$ is Artinian. On the
other hand, if $M \in \mathcal{A}$ there is an injection $M \to E^n$ for some $n$. This can be
seen as follows: consider all linear maps $M \to E^n$, where $n$ is not fixed,
and take one $\varphi\colon M \to E^n$ whose kernel is minimal among the kernels of
those maps. Then, using (i) we can easily see that $\Ker(\varphi) = 0$. Now, from
$0 \to M \to E^n$ we have $(E^n)' = A^n \to M' \to 0$ exact, hence $M' \in \mathcal{N}$. Now the
assertion $M \cong M''$ for $M \in \mathcal{N}$ or $\mathcal{A}$ can easily be checked using (iv) if $M = A$
or $E$, and the general case follows from this and from (i).
\end{proof}

\begin{lemma}\label{lem:18.6}
Let $A$ be a Noetherian ring, $S \subset A$ a multiplicative set, $M$ an $A$-
module and $N \subset M$ a submodule. Assume that $M$ is an essential extension
of $N$; then $M_S$ is an essential extension of $N_S$.
\end{lemma}

\begin{proof}
For $\xi \in M$ we write $\xi_s = \xi/1 \in M_S$; then any element of $M_S$ can be written
$u^{-1}\xi_s$ (with $u$ a unit of $A$, and $\xi \in M$), so that it is enough to show that
for any non-zero $\xi_s$ we have $N_S \cap A_S\xi_s \neq 0$. Suppose that $\ann(\xi_s)$ is a
maximal element of the set of ideals $\{\ann(t\xi_s) \mid t \in S\}$; then if we set $\eta = \xi_s$,
we have $\eta_s = \eta$, and hence $\eta \neq 0$. Now let $\mathfrak{b} = \{b \in A \mid b\eta \in N\}$; by
assumption,
\[
  \mathfrak{b}\eta = A\eta \cap N \neq 0.
\]
Suppose that $\mathfrak{b} = (b_1, \ldots, b_r)$; if $b_1\eta_s = \cdots = b_r\eta_s = 0$ then there is a $t \in S$
such that $tb_i\eta = 0$ for all $i$. Then $t\mathfrak{b}\eta = 0$, but by choice of $\eta$ we have
$\ann(\eta) = \ann(t\eta)$, so that $\mathfrak{b}\eta = 0$, which is a contradiction. Thus $b_i\eta_s \neq 0$
for some $i$, and
\[
  b_i\eta_s \in A_S\eta_s \cap N_S = A_S\xi_s \cap N_S,
\]
as required.
\end{proof}

By Lemmas 5 and 6, if $M$ is an injective hull of $N$ then the $A_S$-module
$M_S$ is an injective hull of $N_S$. Hence if $0 \to M \to I^0 \to I^1 \to \cdots$ is
a minimal injective resolution of an $A$-module $M$, then $0 \to M_S \to I^0_S \to
I^1_S \to \cdots$ is a minimal injective resolution of the $A_S$-module $M_S$. The
$I^i$ are determined uniquely up to isomorphism by $M$. We can therefore
define $\mu_i(P, M)$ to be the number of summands isomorphic to $E(A/P)$
appearing in a decomposition of $I^i$ as a direct sum of indecomposable
modules. We can write symbolically
\[
  I^i = \bigoplus_{P \in \Spec A} \mu_i(P, M)E(A/P).
\]
From what we have just proved, for a multiplicative set $S \subset A$,
\[
  \mu_i(P, M) = \mu_i(PA_S, M_S) \quad \text{if } P \cap S = \emptyset.
\]

\begin{theorem}\label{thm:18.7}
Let $A$ be a Noetherian ring, $M$ an $A$-module, and
$P \in \Spec A$. Then
\[
  \mu_i(P, M) = \dim_{\kappa(P)} \Ext^i_{A_P}(\kappa(P), M_P) = \dim_{\kappa(P)}(\Ext^i_A(A/P, M))_P.
\]
In particular, if $M$ is a finite $A$-module then $\mu_i(P, M) < \infty$.
\end{theorem}

\begin{proof}
Replacing $A$ and $M$ by $A_P$ and $M_P$ we can assume that $(A, P, k)$ is
a local ring. Let $0 \to M \to I^0 \to I^1 \to \cdots$ be a minimal injective
resolution of $M$, so that $\Ext^i_A(k, M)$ is obtained as the homology of the
complex
\[
  \cdots \to \Hom_A(k, I^{i-1}) \to \Hom_A(k, I^i) \to \Hom_A(k, I^{i+1})
  \to \cdots
\]
We can identify $\Hom_A(k, I^i)$ with the submodule $T^i = \{x \in I^i \mid Px = 0\} \subset
I^i$. By construction of the minimal injective resolution, $I^i$ is an essential
extension of $d(I^{i-1})$, so that for $x \in T^i$ the submodule $Ax \cong k$ intersects
$d(I^{i-1})$, and $x \in d(I^{i-1})$. Therefore, $T^i \subset d(I^{i-1})$, and $dT^{i-1} = dT^i = 0$, so
that $\Ext^i_A(k, M) = T^i$. Also,
\[
  \dim_k T^i = \dim_k \Hom_A(k, I^i),
\]
and by Theorem 4, (iii), this is equal to $\mu_i(P, M)$.
\end{proof}

\begin{theorem}\label{thm:18.8}
A necessary and sufficient condition for a ring $A$ to be
Gorenstein is that a minimal injective resolution $0 \to A \to I^0 \to
I^1 \to \cdots$ of $A$ satisfies
\[
  I^i = \bigoplus_{\htop P = i} E(A/P),
\]
or, in other words, $\mu_i(P, A) = \delta_{i,\htop P}$ (the Kronecker $\delta$) for every $P \in \Spec A$.
\end{theorem}

\begin{proof}
By Theorem 7 and condition (2) of Theorem 1 we have
\[
  A_P \text{ is Gorenstein} \iff \mu_i(P, A) = \delta_{i,\htop P}. \qquad\blacksquare
\]
\end{proof}

\begin{theorem}\label{thm:18.9}
Let $(A, m)$ be a Noetherian local ring, and $M$ a finite
$A$-module. Then
\[
  \injdim M < \infty \iff \injdim M = \depth A.
\]
\end{theorem}

\begin{proof}
Suppose that $\injdim M = r < \infty$. If $P$ is a prime ideal distinct from
$m$, choose $x \in m - P$. Then
\[
  0 \to A/P \xrightarrow{x} A/P,
\]
together with the right-exactness of $\Ext^r_A(-, M)$ gives an exact sequence
\[
  \Ext^r_A(A/P, M) \xrightarrow{x} \Ext^r_A(A/P, M) \to 0,
\]
so that by NAK $\Ext^r_A(A/P, M) = 0$. Putting this together with Lemma~\hyperref[lem:18.1]{1},
we get $\Ext^r_A(k, M) \neq 0$. Set $t = \depth A$, and let $x_1, \ldots, x_t \in m$ be a
maximal $A$-sequence; then setting $A/(x_1, \ldots, x_t) = N$ we have $m \in \Ass(N)$.
Hence there exists an exact sequence $0 \to k \to N$, and we must have
$\Ext^t_A(N, M) \neq 0$. The Koszul complex $K(x_1, \ldots, x_t)$ is a projective resolution
of $N = A/(x_1, \ldots, x_t)$, so computing Ext by means of it we see that
\[
  \Ext^t_A(N, M) \cong M/(x_1, \ldots, x_t)M,
\]
and by NAK this is non-zero. Thus $\projdim N = t$, and from $\Ext^r_A(N, M) \neq
0$ we get $t \le r$, whereas from $\Ext^t_A(N, M) \neq 0$ we get $t \ge r$. Hence $t = r$.
\end{proof}

\begin{remark}[the Bass conjecture and the intersection theorem]
Let $(A, m, k)$ be
a Noetherian local ring of dimension $d$. H. Bass~\cite{Ba1} conjectured the
following:
\[
  (B) \quad \text{if there exists a finite } A\text{-module } M (\neq 0) \text{ of finite injective}
\]
\[
  \text{dimension, then } A \text{ is a CM ring.}
\]
According to Theorem 9 this is equivalent to asking that $\injdim M = d$.
The converse of the Bass conjecture is true. Indeed, if $A$ is CM, taking a
maximal $A$-sequence $x_1, \ldots, x_d$, and setting $B = A/(x_1, \ldots, x_d)$ and $E = E(k)$
we have $l_A(B) < \infty$. By Theorem 6, $M = \Hom_A(B, E)$ is also of finite length,
hence is finitely generated. We prove $\injdim_A M \le d$; the Koszul complex
$0 \to A \to A^d \to \cdots \to A^d \to A \to B \to 0$ with respect to $x_1, \ldots, x_d$ provides an $A$-
free resolution of $B$. Now applying the exact functor $\Hom_A(-, E)$ to this
gives the exact sequence $0 \to M \to E \to E^d \to \cdots \to E^d \to E \to 0$. This proves
$\injdim_A M \le d$.

$(B)$ is a special case of the following theorem.
\[
  (C) \quad \text{If } I^\bullet\colon 0 \to I^0 \to \cdots \to I^d \to 0 \text{ is a complex of injective modules such}
\]
\[
  \text{that } H^i(I^\bullet) \text{ is finitely generated for all } i \text{ and } I^\bullet \text{ is not exact, then }
\]
\[
  I^d \neq 0.
\]
Using the theory of dualizing complexes (see~\cite{Rob}) one can prove that $(C)$
is equivalent to the following

\emph{Intersection Theorem.} If $F_\bullet\colon 0 \to F_r \to \cdots \to F_0 \to 0$ is a complex of finitely
generated free modules such that $H_i(F_\bullet)$ has finite length for all $i$ and $F_\bullet$ is
not exact, then $F_r \neq 0$.

$(B)$ was proved by Peskine and Szpiro~\cite{PS1} in some important cases, and
by Hochster~\cite{H} in the equal characteristic case (i.e. when $A$ contains a
field) as a corollary of his existence theorem for the `big CM module', i.e. a
(not necessarily finite) $A$-module with $\depth = \dim A$, see~\cite{H} p. 10 and
p. 70. The intersection theorem was conjectured by Peskine--Szpiro~\cite{PS3} and
by P. Roberts independently. They pointed out that it was also a
consequence of Hochster's theorem. Finally, P. Roberts~\cite{Rob3} settled the
remaining unequal characteristic case of the intersection theorem by using
the advanced technique of algebraic geometry developed by W. Fulton
(\cite{Ful}). Therefore $(B)$, which was known as Bass's conjecture for 24 years,
is now a theorem. Some other conjectures listed in~\cite{H} are still open.
\end{remark}

% ============ 习题 ============
\section*{Exercises to \S\ref{sec:18}}
\emph{Prove the following propositions.}

\begin{exercise}\label{ex:18.1-ex}
Let $(A, m)$ be a Noetherian local ring, $x_1, \ldots, x_n$ an $A$-sequence, and set
$B = A/(x_1, \ldots, x_n)$; then $A$ is Gorenstein $\iff B$ is Gorenstein.
\end{exercise}

\begin{exercise}\label{ex:18.2-ex}
Use the result of Ex.~\ref{ex:18.1-ex} to give another proof of Theorem~\hyperref[thm:18.3]{3}.
\end{exercise}

\begin{exercise}\label{ex:18.3-ex}
If $A$ is Gorenstein then so is the polynomial ring $A[X]$.
\end{exercise}

\begin{exercise}\label{ex:18.4-ex}
Is the ring $R$ of Ex.~\ref{ex:17.2-ex} Gorenstein?
\end{exercise}

\begin{exercise}\label{ex:18.5-ex}
Let $(A, m, k)$ be a local ring; then $E = E_A(k)$ is a faithful $A$-module (that is
$0 \neq a \in A \Rightarrow aE \neq 0$).
\end{exercise}

\begin{exercise}\label{ex:18.6-ex}
Let $(A, m, k)$ be a complete Noetherian local ring and $M$ an $A$-module. If $M$
is a faithful $A$-module and is an essential extension of $k$ then $M \cong E_A(k)$.
\end{exercise}

\begin{exercise}\label{ex:18.7-ex}
Let $k$ be a field, $S = k[X_1, \ldots, X_n]$ and $P = (X_1, \ldots, X_n)$; set $A = S_P$,
$\hat{A} = k[[X_1, \ldots, X_n]]$ and $E = k[X_1^{-1}, \ldots, X_n^{-1}]$. We make $E$ into an
$A$-module by the following multiplication: if $X^\alpha = X_1^{\alpha_1} \cdots X_n^{\alpha_n}$ and $X^{-\beta}
= X_1^{-\beta_1} \cdots X_n^{-\beta_n}$, the product $X^\alpha X^{-\beta}$ is defined to be $X^{\alpha-\beta}$ if $\alpha_i \le \beta_i$ for all $i$
and $0$ otherwise. Then $E = E_A(S/P) = E_{\hat{A}}(k)$. (Use the preceding question;
see also Northcott~\cite{Nort8} for further results. The elements of this $A$-module
$E$ are called \emph{inverse polynomials}; they were defined and used by Macaulay
~\cite{Mac} as early as 1916.)
\end{exercise}

\begin{exercise}\label{ex:18.8-ex}
Let $k$ be a field and $t$ an indeterminate. Consider the subring $A =
k[t^3, t^5, t^7]$ of $k[t]$ and show that $A$ is a one-dimensional CM ring which
is not Gorenstein. How about $k[t^3, t^4, t^5]$ and $k[t^4, t^5, t^6]$?
\end{exercise}


% ============ Chapter 7. Regular rings (原书 pp.153-173) ============
\chapter{Regular rings}\label{ch:7}

Regular local rings have already been mentioned several times, and in
this chapter we are going to carry out a study of them using homological
algebra. Serre's Theorem~\ref{thm:19.2}, characterising regular local rings as Noetherian
local rings of finite global dimension, is the really essential result, and
from this one can deduce at once, for example, that a localisation of a
regular local ring is again regular (Theorem~\ref{thm:19.3}); this is a result which
ideal theory on its own was only able to prove with difficulty in special
cases. \S\ref{sec:20} on UFDs is centred around the theorem that a regular local
ring is a UFD, another important achievement of homological methods;
we only cover the basic topics. This section was written referring to the early
parts of Professor M.\ Narita's lectures at Tokyo Metropolitan University.
In \S\ref{sec:21} we give a simple discussion of the most elementary results on
complete intersection rings. This is an area where the homology theory of
rings plays an essential role, but we are only able to mention this in
passing.

% ============ §19 Regular rings (原书 pp.153-161) ============
\section{Regular rings}\label{sec:19}\index{regular local rings, section 19}\index{regular rings, section 19}\index{minimal resolution, section 19}\index{global dimension, section 19}

\textbf{Minimal free resolutions.} Let $(A, m, k)$ be a local ring, $M$ and $N$ finite
$A$-modules. An $A$-linear map $\varphi\colon M \to N$ induces a $k$-linear map
$\bar{\varphi}\colon M \otimes k \to N \otimes k$, which we denote $\bar{\varphi}$; then one sees easily that
\[
  \varphi \text{ is an isomorphism } \iff \bar{\varphi} \text{ is surjective and } \Ker \varphi = mM.
\]
In particular for free modules $M$ and $N$, if $\bar{\varphi}$ is an isomorphism then $\rank
M = \rank N$, and writing $\varphi$ as a matrix we have $\det \varphi \notin m$, so that
\[
  \varphi \text{ is an isomorphism } \iff \bar{\varphi} \text{ is an isomorphism}.
\]
Let $M$ be a finite $A$-module. An exact sequence
\[
  (*) \quad \cdots \xrightarrow{d_{i+1}} L_i \xrightarrow{d_i} L_{i-1} \to \cdots \to L_1 \xrightarrow{d_1} L_0 \xrightarrow{\varepsilon} M \to 0
\]
(or the complex $L_\bullet$) is called a \emph{minimal (free) resolution}\index{minimal resolution} of $M$ if it satisfies the
three conditions (1) each $L_i$ is a finite free $A$-module, (2) $\bar{d}_i = 0$, or in other
words $d_iL_i \subset mL_{i-1}$ for all $i$, and (3) $\bar{\varepsilon}\colon L_0 \otimes k \to M \otimes k$ is an isomorph-
ism. Breaking up $(*)$ into short exact sequences $0 \to K_1 \to L_0 \to M \to 0$,
$0 \to K_2 \to L_1 \to K_1 \to 0, \ldots$, we have $L_0 \otimes k \cong M \otimes k$,
$L_1 \otimes k \cong K_1 \otimes k, \ldots$. Any two minimal resolutions of $M$ are isomorphic as com-
plexes (prove this!).

\begin{example*}
Let $x_1, \ldots, x_n \in m$ be an $A$-sequence, and let $K_\bullet = K_\bullet(x_1, \ldots, x_n)$
be the Koszul complex
\[
  0 \to K_n \to K_{n-1} \to \cdots \to K_0 \to A/(x_1, \ldots, x_n) \to 0;
\]
then $K_\bullet$ is a minimal resolution of $A/(x_1, \ldots, x_n)$ over $A$.
\end{example*}

Let $(A, m, k)$ be a Noetherian local ring; then a finite $A$-module $M$
always has a minimal resolution. Construction: let $\{\omega_1, \ldots, \omega_p\}$ be a
minimal basis of $M$, let $L_0 = Ae_1 + \cdots + Ae_p$ be a free module, and define
$\varepsilon\colon L_0 \to M$ by $\varepsilon(e_i) = \omega_i$; taking $K_1$ to be the kernel of $\varepsilon$ we get
$0 \to K_1 \to L_0 \to M \to 0$ with $L_0 \otimes k \cong M \otimes k$. Now $K_1$ is again a
finite $A$-module, so that we need only proceed as before.

\begin{lemma}\label{lem:19.1}
Let $(A, m, k)$ be a local ring, and $M$ a finite $A$-module. Suppose
that $L_\bullet$ is a minimal resolution of $M$; then
\begin{enumerate}[(i)]
  \item $\dim_k \Tor^A_i(M, k) = \rank L_i$ for all $i$,
  \item $\projdim M = \sup \{i \mid \Tor^A_i(M, k) \neq 0\} \le \projdim_A k$,
  \item if $M \neq 0$ and $\projdim M = r < \infty$ then for any finite $A$-module $N \neq 0$
        we have $\Ext^r_A(M, N) \neq 0$.
\end{enumerate}
\end{lemma}

\begin{proof}
(i) We have $\Tor^A_i(M, k) = H_i(L_\bullet \otimes k)$, but from the definition of
minimal resolution, $\bar{d}_i = 0$, and hence $H_i(L_\bullet \otimes k) = L_i \otimes k$, and the dimen-
sion of this as a $k$-vector space is equal to $\rank_A L_i$.

(ii) follows from (i).

(iii) Since $L_{r+1} = 0$ and $L_r \neq 0$, $\Ext^r_A(M, N)$ is the cokernel of $d_r^*\colon
\Hom_A(L_r, N) \to \Hom_A(L_{r-1}, N)$, but since $L_i$ is free, $\Hom_A(L_i, N)$ is just
a direct sum of a number of copies of $N$; we can write $d_r\colon L_r \to L_{r-1}$ as a
matrix with entries in $m$, and then $d_r^*$ is given by the same matrix, so that
$\Im(d_r^*) \subset m\,\Hom_A(L_r, N)$, and by NAK $\Ext^r_A(M, N) \neq 0$.
\end{proof}

\begin{remark}
One sees from the above lemma that $\Tor_i(M, k) = 0$ implies that
$L_i = 0$, and therefore $\projdim M < i$, so that $\Tor_j(M, k) = 0$ for $j > i$. It is
conjectured that this holds in more generality, or more precisely:
\emph{Rigidity conjecture.}\index{rigidity conjecture, section 19} Let $R$ be a Noetherian ring, $M$ and $N$ finite $R$-modules;
suppose that $\projdim M < \infty$. Then $\Tor^R_i(M, N) = 0$ implies that
$\Tor^R_j(M, N) = 0$ for all $j > i$.
\end{remark}

This has been proved by Lichtenbaum~\cite{Li1} if $R$ is a regular ring, but
is unsolved in general.

The following theorem is not an application of Lemma~\hyperref[lem:19.1]{1}, but is proved
by a similar technique.

\begin{theorem}[Auslander and Buchsbaum]\label{thm:19.1}
Let $A$ be a Noetherian local
ring and $M \neq 0$ a finite $A$-module. Suppose that $\projdim M < \infty$; then
\[
  \projdim M + \depth M = \depth A.
\]
\end{theorem}

\begin{proof}
Set $\projdim M = h$; we work by induction on $h$. If $h = 0$ then $M$
is a free $A$-module, so that the assertion is trivial. If $h = 1$, let
\[
  (\varphi) \quad 0 \to A^m \xrightarrow{\varphi} A^n \to M \to 0
\]
be a minimal resolution of $M$. We can write $\varphi$ as an $m \times n$ matrix with
entries in $m$. From $(\varphi)$ we obtain the long exact sequence
\[
  \cdots \to \Ext^i_A(k, A^n) \xrightarrow{\varphi^*} \Ext^i_A(k, A^m) \to \Ext^{i+1}_A(k, M) \to \cdots,
\]
and writing out $\Ext^i_A(k, A^m) = \Ext^i_A(k, A)^m$ and $\Ext^i_A(k, A^n) = \Ext^i_A(k, A)^n$,
we can express $\varphi^*$ by the same matrix as $\varphi$. However, the entries of $\varphi$ are
elements of $m$, and therefore annihilate $\Ext^i_A(k, A)$, so that $\varphi^* = 0$, and we
have an exact sequence
\[
  0 \to \Ext^i_A(k, A)^n \to \Ext^{i+1}_A(k, M) \to \Ext^{i+1}_A(k, A)^m \to 0
\]
for every $i$. Since $\depth M = \inf \{i \mid \Ext^i_A(k, M) \neq 0\}$ we have $\depth M =
\depth A - 1$ and the theorem holds if $h = 1$. If $h > 1$ then taking any exact
sequence
\[
  0 \to M' \to A^n \to M \to 0,
\]
we have $\projdim M' = h - 1$, so that an easy induction completes the
proof.
\end{proof}

\begin{lemma}\label{lem:19.2}
Let $A$ be a ring and $n > 0$ a given integer. Then the following
conditions are equivalent.
\begin{enumerate}[(1)]
  \item $\projdim M \le n$ for every $A$-module $M$;
  \item $\projdim M \le n$ for every finite $A$-module $M$;
  \item $\injdim N \le n$ for every $A$-module $N$;
  \item $\Ext^{n+1}_A(M, N) = 0$ for all $A$-modules $M$ and $N$.
\end{enumerate}
\end{lemma}

\begin{proof}
$(1) \implies (2)$ is trivial.

$(2) \implies (3)$ For any ideal $I$, the $A$-module $A/I$ is finite, so that
$\Ext^{n+1}_A(A/I, N) = 0$, and hence by \S\ref{sec:18}, Lemma~\hyperref[lem:18.1]{1}, $\injdim N \le n$.

$(3) \implies (4)$ is trivial, and $(4) \implies (1)$ is well-known (see p.~280).
\end{proof}

We define the \emph{global dimension}\index{global dimension} of a ring $A$ by
\[
  \gldim A = \sup \{\projdim M \mid M \text{ is an $A$-module}\}.
\]
According to Lemma~\hyperref[lem:19.2]{2} above, this is also equal to the maximum projective
dimension of all finite $A$-modules. If $(A, m, k)$ is a Noetherian local ring then by
Lemma~\hyperref[lem:19.1]{1}, $\gldim A = \projdim_A k$.

We have defined regular local rings (see \S\ref{sec:14}) as Noetherian local rings
for which $\dim A = \embdim A$, and we have seen that they are integral
domains (Theorem~\ref{thm:14.3}) and CM rings (Theorem~\ref{thm:17.8}). A regular local
ring is Gorenstein (Theorem~\ref{thm:18.1}, (5')). A necessary and sufficient condition
for a Noetherian local ring $(A, m, k)$ to be regular is that $\gr_m(A)$ is a poly-
nomial ring over $k$ (Theorem~\ref{thm:17.10}). The following theorem gives another
important necessary and sufficient condition.

\begin{theorem}[Serre]\label{thm:19.2}
Let $A$ be a Noetherian local ring. Then
\[
  A \text{ is regular} \iff \gldim A = \dim A \iff \gldim A < \infty.
\]
\end{theorem}

\begin{proof}
(I) Suppose that $(A, m, k)$ is an $n$-dimensional regular local ring.
Let $x_1, \ldots, x_n$ be a regular system of parameters; then since this is an
$A$-sequence, the Koszul complex $K_\bullet(x_1, \ldots, x_n)$ is a minimal free resolution of
$A/(x_1, \ldots, x_n) = k$, and $K_n \neq 0$, $K_{n+1} = 0$, so that as we have already seen,
$\gldim A = \projdim k = n$.

(II) Conversely, suppose that $\gldim A = r < \infty$ and $\embdim A = s$. We
prove that $A$ is regular by induction on $s$; we can assume that $s > 0$, that
is $m \neq 0$. Then $m \notin \Ass(A)$: for if $0 \neq a \in A$ is such that $ma = 0$, consider
a minimal resolution
\[
  0 \to L_r \to L_{r-1} \to \cdots \to L_0 \to k \to 0
\]
of $k$ (with $r > 0$); then $L_r \subset mL_{r-1}$, but then $aL_r = 0$, which contradicts the
assumption that $L_r$ is a free module. Thus we can choose $x \in m$ not
contained in $m^2$ or in any associated prime of $A$. Then $x$ is $A$-regular,
hence also $m$-regular, so that if we set $B = A/xA$ then according to Lemma~\hyperref[lem:18.2]{2}
of \S\ref{sec:18}, $\Ext^q_A(m, N) = \Ext^q_B(m/xm, N)$ for all $B$-modules $N$, and hence
we obtain $\projdim_B m/xm \le r$.

Now we prove that the natural map $m/xm \to m/xA$ splits, so that
$m/xA$ is isomorphic to a direct summand of $m/xm$. Since $x \notin m^2$, we can take
a minimal basis $x_1 = x, x_2, \ldots, x_s$ of $m$ starting with $x$ (here $s = \embdim
A$). We set $b = (x_2, \ldots, x_s)$, so that by the minimal basis condition,
$b \cap xA \subset xm$, and therefore there exists a chain
\[
  m/xA = (b + xA)/xA \cong b/(b \cap xA) \to m/xm \to m/xA
\]
of natural maps, whose composite is the identity. This proves the above
claim. Now clearly,
\[
  \projdim_B m/xA \le \projdim_B m/xm \le r.
\]
Taking a minimal $B$-projective resolution of $m/xA$ and patching it together
with the exact sequence $0 \to m/xA \to B \to k \to 0$ gives a projective
resolution of $k$ of length $\le r + 1$, and hence $\gldim B = \projdim_B k \le r + 1$,
so that by induction, $B$ is a regular local ring. Since $x$ is not contained
in any associated prime of $A$ we have $\dim B = \dim A - 1$, and therefore
$A$ is regular.
\end{proof}

\begin{theorem}[Serre]\label{thm:19.3}
Let $A$ be a regular local ring and $P$ a prime ideal; then
$A_P$ is a regular local ring.
\end{theorem}

\begin{proof}
Since $\projdim_A A/P \le \gldim A < \infty$, as an $A$-module $A/P$ has a
projective resolution $L_\bullet$ of finite length. Then $L_\bullet \otimes_A A_P$ is a projective
resolution of $(A/P) \otimes_A A_P = A_P/PA_P = K(P)$ as an $A_P$-module, so that $K(P)$
has a projective resolution of finite length as an $A_P$-module, which means that $A_P$
has finite global dimension; thus by the previous theorem, $A_P$ is
regular.
\end{proof}

\begin{definition*}
A \emph{regular ring}\index{regular rings} is a Noetherian ring such that the localisation
at every prime is a regular local ring. By the previous theorem, it is
enough\index{regular local rings} for the localisation at every maximal ideal to be regular.
\end{definition*}

\begin{theorem}\label{thm:19.4}
A regular ring is normal.
\end{theorem}

\begin{proof}
The definition of normal is local, so that it is enough to show that a
regular local ring is normal. We show that the conditions of the corollary
to Theorem~\ref{thm:11.5} are satisfied. (a) The localisation at a height 1 prime ideal is a
DVR by the previous theorem and Theorem~\ref{thm:11.2}. (b) All the prime
divisors of a non-zero principal ideal have height 1 by Theorem~\ref{thm:17.8} (the
implication regular $\Rightarrow$ CM).
\end{proof}

\begin{theorem}\label{thm:19.5}
If $A$ is regular then so are $A[X]$ and $A[[X]]$.
\end{theorem}

\begin{proof}
For $A[X]$, let $P$ be a maximal ideal of $A[X]$ and set $P \cap A = m$.
Then $A[X]_P$ is a localisation of $A_m[X]$, so that replacing $A$ by $A_m$, we can assume
$A$ is a regular local ring. Then setting $A/m = k$ we have $A[X]/m[X] =
k[X]$, so that there is a monic polynomial $f(X)$ with coefficients in $A$ such that
$P = (m, f(X))$, and such that $f$ reduces to an irreducible polynomial in
$k[X]$ modulo $m$. Then by Theorem~\ref{thm:15.1}, we clearly have
\[
  \dim A[X]_P = \htop P = 1 + \htop m = 1 + \dim A;
\]
on the other hand $m$ is generated by $\dim A$ elements, so that $P = (m, f)$ is
generated by $\dim A + 1$ elements, and therefore $A[X]_P$ is regular.

For $A[[X]]$, set $B = A[[X]]$ and let $M$ be a maximal ideal of $B$; then $X \in M$
by Theorem~\ref{thm:8.2}, (i). Therefore $M \cap A = m$ is a maximal ideal of $A$. Now
although we cannot say that $B_M$ contains $A_m[[X]]$, the two have the same
completion, $(B_M)^\wedge = (A_m)^\wedge[[X]]$. A Noetherian local ring is regular if and
only if its completion is regular (since both the dimension and embedding
dimension remain the same on taking the completion). Thus if we replace
$A$ by $(A_m)^\wedge$, the maximal ideal of $B = A[[X]]$ is $M = (m, X)$, and $\htop M =
\htop m + 1$, so that $B$ is also regular.
\end{proof}

Next we discuss the properties of modules which have finite free
resolutions (the definition is given below).

\begin{lemma}[Schanuel]\label{lem:19.3}
Let $A$ be a ring and $M$ an $A$-module. Suppose that
\[
  0 \to K \to P \to M \to 0 \quad \text{and} \quad 0 \to K' \to P' \to M \to 0
\]
are exact sequences with $P$ and $P'$ projective. Then $K \oplus P' \cong K' \oplus P$.
\end{lemma}

% 注: 原书此处用 λ, λ' 与矩阵 (φ, ψ, φψ = 1) 的写法, 扫描版矩阵不可辨;
% 此处按标准 fibre product 证法转录, 结论与原书一致.
\begin{proof}
From the fact that $P$ and $P'$ are projective, there exist $\lambda\colon P \to P'$
and $\lambda'\colon P' \to P$, giving the diagram
\[
  \begin{array}{ccccccc}
    0 & \to & K & \to & P & \xrightarrow{\varepsilon} & M \to 0 \\
    & & & & \downarrow \lambda' & & \downarrow \lambda \\
    0 & \to & K' & \to & P' & \xrightarrow{\varepsilon'} & M \to 0
  \end{array}
\]
with $\varepsilon'\lambda = \varepsilon$ and $\varepsilon\lambda' = \varepsilon'$.
We add in harmless summands $P'$ and $P$ to the two exact rows, and line
up the middle terms:
\[
  0 \to K \oplus P' \to P \oplus P' \xrightarrow{(\varepsilon, 0)} M \to 0
  \quad \text{and} \quad
  0 \to P \oplus K' \to P \oplus P' \xrightarrow{(0, \varepsilon')} M \to 0.
\]
Consider the submodule $W = \{(x, y) \in P \oplus P' \mid \varepsilon(x) = \varepsilon'(y)\}$ of $P \oplus P'$,
the fibre product of $P$ and $P'$ over $M$. Projection onto the first factor gives
the exact sequence $0 \to K' \to W \to P \to 0$: indeed $W \to P$ is surjective since
$\varepsilon'$ is surjective, and its kernel is $\{(0, y) \mid \varepsilon'(y) = 0\} = K'$. Since $P$ is
projective, this sequence splits, so that $W \cong K' \oplus P$. Similarly, projection onto
the second factor gives the exact sequence $0 \to K \to W \to P' \to 0$, which splits, so
that $W \cong K \oplus P'$. Therefore $\varphi$ induces an isomorphism
$K \oplus P' \cong P \oplus K'$.
\end{proof}

\begin{lemma}[generalised Schanuel lemma]\label{lem:19.4}
Let $A$, $M$ be as above, and suppose
that
\[
  0 \to P_n \to \cdots \to P_1 \to P_0 \to M \to 0
  \quad \text{and} \quad
  0 \to Q_n \to \cdots \to Q_1 \to Q_0 \to M \to 0
\]
are exact sequences with $P_i$ and $Q_i$ projective for
$0 \le i \le n - 1$. Then
\[
  P_0 \oplus Q_1 \oplus P_2 \oplus \cdots \cong Q_0 \oplus P_1 \oplus Q_2 \oplus \cdots
\]
(the alternating sums of $P_i$ and $Q_i$).
\end{lemma}

\begin{proof}
Write $K$ for the kernel of $P_0 \to M$ and $K'$ for the kernel of
$Q_0 \to M$; then, by the previous lemma, $K \oplus Q_0 \cong P_0 \oplus K'$. Now add
in harmless summands $Q_n$ and $P_n$ to $0 \to P_n \to \cdots \to P_1 \to
K \to 0$ and $0 \to Q_n \to \cdots \to Q_1 \to K' \to 0$ respectively, to obtain
exact sequences
\[
  0 \to P_n \to \cdots \to P_1 \oplus Q_0 \to K \oplus Q_0 \to 0
  \quad \text{and} \quad
  0 \to Q_n \to \cdots \to Q_1 \oplus P_0 \to P_0 \oplus K' \to 0,
\]
and induction on $n$ gives the assertion.
\end{proof}

A \emph{finite free resolution} (or \emph{FFR}\index{FFR} for short) of an $A$-module $M$ is a
resolution
\[
  0 \to F_r \to \cdots \to F_1 \to F_0 \to M \to 0 \quad \text{(of finite length)}
\]
with each $F_i$ a finite free module. If $M$ has an FFR we set
\[
  \chi(M) = \sum_{i=0}^{r} (-1)^i \rank F_i,
\]
and call $\chi(M)$ the \emph{Euler number}\index{Euler number} of $M$. By Lemma~\hyperref[lem:19.4]{4},
$\chi(M)$ is independent of the choice of FFR. Moreover, since for any prime
ideal $P$,
\[
  0 \to (F_r)_P \to \cdots \to (F_1)_P \to (F_0)_P \to M_P \to 0
\]
is an FFR of the $A_P$-module $M_P$ we have $\chi(M) = \chi(M_P)$. If $M$ is itself
free then one sees easily from Lemma~\hyperref[lem:19.4]{4} that $\chi(M) = \rank M$.

\begin{theorem}\label{thm:19.6}
Let $(A, m)$ be a local ring, and suppose that for any finite
subset $E \subset m$ there exists $0 \neq y \in A$ such that $yE = 0$; then the only
$A$-modules having an FFR are the free modules.
\end{theorem}

\begin{remark}
If $A$ is Noetherian then the assumption on $m$ is equivalent to
$m \in \Ass(A)$, or $\depth A = 0$. In this case the theorem is a special case of
Theorem~\ref{thm:19.1}.
\end{remark}

\begin{proof}
Suppose that $0 \to F_r \to F_{r-1} \to \cdots \to F_0 \to M \to 0$ is an FFR
of $M$, and set $N = \Coker(F_r \to F_{r-1})$; if we prove that $N$ is free then
we can decrease $r$ by 1, so that we only need consider the case $0 \to
F_1 \to F_0 \to M \to 0$. Now let $0 \to L_1 \to L_0 \to M \to 0$ be a minimal
free resolution of $M$; then since $L_0$ and $F_0$ are finitely generated, by
Schanuel's lemma (or by Theorem~\ref{thm:2.6}), $L_1$ is also finite. Considering bases of
$F_1$ and $L_1$, we can write down a set of generators of $L_1$ as a submodule of
$F_0$ using only a finite number of elements of $m$. Then by assumption, there
exists $0 \neq y \in A$ such that $yL_1 = 0$. Since $L_1$ is a free module, we must
have $L_1 = 0$, so that $M \cong L_0$, and is free.
\end{proof}

\begin{theorem}\label{thm:19.7}
Let $A$ be any ring; if $M$ is an $A$-module having an FFR then
$\chi(M) \ge 0$.
\end{theorem}

\begin{proof}
Choose a minimal prime ideal $P$ of $A$; since $\chi(M) = \chi(M_P)$, we can
replace $A$ by $A_P$, and then $A$ is a local ring with maximal ideal $m$ equal to
$\nil(A)$. Then the assumption of the previous theorem is satisfied: for
$x_1, \ldots, x_r \in m$, we can assume by induction on $r$ that there is a $z \neq 0$
such that $zx_1 = \cdots = zx_{r-1}$; but $x_r$ is nilpotent, so that there is an $i \ge 0$
such that $zx_r^i \neq 0$ but $zx_r^{i+1} = 0$, and we can take $y = zx_r^i$. Therefore by
the previous theorem $M$ is a free module, and $\chi(M) = \rank M \ge 0$.
\end{proof}

\begin{theorem}[Auslander and Buchsbaum~\cite{AB2}]\label{thm:19.8}
Let $A$ be a Noetherian
ring and $M$ an $A$-module, and suppose that $M$ has an FFR. Then the
following three conditions are equivalent:
\begin{enumerate}[(1)]
  \item $\ann(M) \neq 0$;
  \item $\chi(M) = 0$;
  \item $\ann(M)$ contains an $A$-regular element.
\end{enumerate}
\end{theorem}

\begin{proof}
$(1) \implies (2)$ Suppose that $\chi(M) > 0$; then for any $P \in \Ass(A)$ we have
$\chi(M_P) > 0$, and hence $M_P \neq 0$. By Theorem~\hyperref[thm:19.6]{6}, $M_P$ is a free $A_P$-module, so
that setting $I = \ann(M)$ we have $I_P = \ann(M_P) = 0$. If we set $J = \ann(I)$
then this is equivalent to $J \not\subset P$. Since this holds for every $P \in \Ass(A)$
we see that $J$ contains an $A$-regular element, but then $J \cdot I = 0$ implies that
$I = 0$, a contradiction.

$(2) \implies (3)$ If $\chi(M) = 0$ then by Theorem~\hyperref[thm:19.6]{6}, $M_P = 0$ for every $P \in \Ass(A)$.
This means that $\ann(M) \not\subset P$, so that $\ann(M)$ contains an $A$-regular
element.

$(3) \implies (1)$ is obvious.
\end{proof}

\begin{theorem}[Vasconcelos~\cite{Vas1}]\label{thm:19.9}
Let $A$ be a Noetherian local ring, and
$I$ a proper ideal of $A$; assume that $\projdim I < \infty$. Then
\[
  I \text{ is generated by an $A$-sequence} \iff I/I^2 \text{ is a free module over } A/I.
\]
\end{theorem}

\begin{proof}
$(\implies)$ is already known (Theorem~\ref{thm:16.2}). In fact, $I^v/I^{v+1}$ is a free
$A/I$-module for $v = 1, 2, \ldots$

$(\Longleftarrow)$ We can assume that $I \neq 0$. Since $I$ has finite projective dimension
over $A$ so has $A/I$, and since $A$ is local, $A/I$ has an FFR. Now $\ann(A/I) = I$,
so that by the previous theorem $I$ is not contained in any associated prime
of $A$, and therefore we can choose an element $x \in I$ such that $x$ is not
contained in $m^2$ or in any associated prime of $A$. Then $x$ is $A$-regular,
and $\bar{x} = x \bmod I^2$ is a member of a basis of $I/I^2$ over $A/I$; let $\bar{x}$,
$\bar{y}_2, \ldots, \bar{y}_n \in I$ be such that their images form a basis of $I/I^2$. Then if we
set $B = A/xA$, we see by the same argument as in (II) of the proof of
Theorem~\hyperref[thm:19.2]{2} that $\projdim_B I/xI < \infty$, and that $I/xA$ is isomorphic to a direct
summand of $I/xI$. We now set $I^* = I/xA$, so that $\projdim_A I^* < \infty$. But
on the other hand one sees easily that $I^*/I^{*2}$ is a free module over $B/I^*$,
and an induction on the number of generators of $I$ completes the proof.
\end{proof}

\begin{remark}
In Lech~\cite{Le1}, a set $x_1, \ldots, x_n$ of elements of $A$ is defined to be
\emph{independent}\index{independent elements} if
\[
  \sum a_ix_i = 0 \Rightarrow a_i \in (x_1, \ldots, x_n) \quad \text{for all } i.
\]
If we set $I = (x_1, \ldots, x_n)$ then this condition is equivalent to saying that
the images of $x_1, \ldots, x_n$ in $I/I^2$ form a basis of $I/I^2$ over $A/I$. Then if $A$
and $I$ satisfy the hypotheses of the previous theorem, the theorem tells us
that $I = (y_1, \ldots, y_n)$ with $y_1, \ldots, y_n$ an $A$-sequence. Setting $x_i = \sum a_{ij}y_j$ we
see that the matrix $(a_{ij})$ is invertible when considered in $A/I$; this means that
the determinant of $(a_{ij})$ is not in the maximal ideal of $A$, and so $(a_{ij})$ itself
is invertible. Thus $x_1, \ldots, x_n$ is an $A$-quasi-regular sequence, hence an
$A$-sequence. In particular, we get the following corollary.
\end{remark}

\begin{corollary*}
Let $(A, m)$ be a regular local ring. Then if $x_1, \ldots, x_n \in m$ are
independent in the sense of Lech, they form an $A$-sequence.
\end{corollary*}

However, if we try to prove this corollary as it stands, the induction
does not go through. The key to success with Vasconcelos' theorem is to
strengthen the statement so that induction can be used effectively. Now
as Kaplansky has also pointed out, the main part of Theorem~\hyperref[thm:19.2]{2} (the
implication $\gldim A < \infty \Rightarrow$ regular) follows at once from Theorem~\hyperref[thm:19.9]{9},
because if $m$ is generated by an $A$-sequence then $\embdim A \le \depth A \le
\dim A$.

% ============ 习题 ============
\section*{Exercises to \S\ref{sec:19}}
\emph{Prove the following propositions.}

\begin{exercise}\label{ex:19.1-ex}
Let $k$ be a field and $R = R_0 + R_1 + R_2 + \cdots$ a Noetherian graded ring
with $R_0 = k$; set $m = R_1 + R_2 + \cdots$. Show that if $R_m$ is an $n$-dimensional
regular local ring then $R$ is a polynomial ring $R = k[y_1, \ldots, y_n]$ with $y_i$
homogeneous of positive degree.
\end{exercise}

\begin{exercise}\label{ex:19.2-ex}
Let $A$ be a ring and $M$ an $A$-module. Say that $M$ is \emph{stably free}\index{stably free} if there exist
finite free modules $F$ and $F'$ such that $M \oplus F = F'$. Obviously a stably free
$A$-module $M$ is a finite projective $A$-module, and has an FFR $0 \to
F \to F' \to M \to 0$. Prove that, conversely, a finite projective module
having an FFR is stably free.
\end{exercise}

\begin{exercise}\label{ex:19.3-ex}
Prove that if every finite projective module over a Noetherian ring $A$ is
stably free then every finite $A$-module of finite projective dimension has an
FFR.
\end{exercise}

\begin{exercise}\label{ex:19.4-ex}
Prove that if every finite module over a Noetherian ring $A$ has an FFR
then $A$ is regular.
\end{exercise}

% ============ §20 UFDs (原书 pp.161-169) ============
\section{UFDs}\label{sec:20}\index{UFD, section 20}\index{factorial ring, section 20}

This section treats UFDs, which we have already touched on in
\S\ref{sec:1}; note that the Bourbaki terminology for UFD is `factorial ring'. First of
all, we have the following criterion for Noetherian rings.

\begin{theorem}\label{thm:20.1}
A Noetherian integral domain $A$ is a UFD if and only if
every height 1 prime ideal is principal.
\end{theorem}

\begin{proof}
Proof of `only if'. Suppose that $A$ is a UFD and that $P$ is a height 1
prime ideal. Take any non-zero $a \in P$, and express $a$ as a product of prime
elements, $a = \prod \pi_i$. Then at least one of the $\pi_i$ belongs to $P$; if $\pi_i \in P$ then
$(\pi_i) \subset P$, but $(\pi_i)$ is a non-zero prime ideal and $\htop P = 1$, hence $P = (\pi_i)$.

Proof of `if'. Since $A$ is Noetherian, every element $a \in A$ which is neither
0 nor a unit can be written as a product of finitely many irreducibles.
Hence it will be enough to prove that an irreducible element $a$ is a prime
element. Let $P$ be a minimal prime divisor of $(a)$; then by the principal
ideal theorem (Theorem~\ref{thm:13.9}), $\htop P = 1$, so that by assumption we can
write $P = (b)$. Thus $a = bc$, and since $a$ is irreducible, $c$ is a unit, so that
$(a) = (b) = P$, and $a$ is a prime element.
\end{proof}

\begin{theorem}\label{thm:20.2}
Let $A$ be a Noetherian integral domain, $\Gamma$ a set of prime
elements of $A$, and let $S$ be the multiplicative set generated by $\Gamma$. If $A_S$
is a UFD then so is $A$.
\end{theorem}

\begin{proof}
Let $P$ be a height 1 prime ideal of $A$. If $P \cap S \neq 0$ then $P$ contains
an element $\pi \in \Gamma$, and since $\pi A$ is a non-zero prime ideal we have $P = \pi A$.
If $P \cap S = 0$ then $PA_S$ is a height 1 prime ideal of $A_S$, so that $PA_S = aA_S$
for some $a \in P$. Among all such $a$ choose one such that $aA$ is maximal;
then $a$ is not divisible by any $\pi \in \Gamma$. Now if $x \in P$ we have $sx = ay$ for
some $s \in S$ and $y \in A$. Let $s = \pi_1 \cdots \pi_r$ with $\pi_i \in \Gamma$; then $a \notin \pi_iA$, so
that $y \in \pi_iA$, and an induction on $r$ shows that $y \in sA$, so that $x \in aA$. Hence
$P = aA$.
\end{proof}

\begin{lemma}\label{lem:20.1}
Let $A$ be an integral domain, and $a$ an ideal of $A$ such that
$a \oplus A^n \cong A^{n+1}$; then $a$ is principal.
\end{lemma}

\begin{proof}
Fix the basis $e_0, \ldots, e_n$ of $A^{n+1}$, and viewing $a \oplus A^n \subset A \oplus A^n$, fix
$f_0, \ldots, f_n$ such that $f_0$ is a basis of $A$ and $f_1, \ldots, f_n$ a basis of $A^n$. Then
the isomorphism $\varphi\colon A^{n+1} \to a \oplus A^n$ can be given in the form $\varphi(e_i) =
\sum_{j=0}^{n} a_{ij}f_j$. Write $d_i$ for the $(i, 0)$th cofactor of the matrix $(a_{ij})$, and $d$
for the determinant, so that, since $\varphi$ is injective, $d \neq 0$, and $\sum_i a_{i0}d_i = d$,
$\sum_i a_{ij}d_i = 0$ if $j \neq 0$. Hence if we set $e_0' = \sum_i d_ie_i$, we have $\varphi(e_0') = df_0$.
Moreover, since the image of $\varphi$ includes $f_1, \ldots, f_n$, there exist $e_1', \ldots,
e_n' \in A^{n+1}$ such that $\varphi(e_j') = f_j$. Now define a matrix $(c_{jk})$ by $e_j' = \sum_{k=0}^{n} c_{jk}e_k$
for $j = 0, \ldots, n$ (so $c_{0k} = d_k$). Then we have
\[
  (c_{jk})(a_{ij}) = \begin{pmatrix}
    d & 0 & \cdots & 0 \\
    0 & 1 & \cdots & 0 \\
    \vdots & \vdots & \ddots & \vdots \\
    0 & 0 & \cdots & 1
  \end{pmatrix},
\]
so that by comparing the determinants of both sides we get $\det(c_{jk}) = 1$.
Therefore $e_0', \ldots, e_n'$ is another basis of $A^{n+1}$, and $\varphi(Ae_0') = dAf_0$,
so that $a = dA$.
\end{proof}

Let $K$ be the field of fractions of the integral domain $A$; for a finite
$A$-module $M$, the dimension of $M \otimes_A K$ as a vector space over $K$ is called
the \emph{rank}\index{rank of a module} of $M$. A torsion-free finite $A$-module of rank 1 is isomorphic to
an ideal of $A$. Lemma~\hyperref[lem:20.1]{1} can be formulated as saying that for an integral
domain $A$, a stably free rank 1 module is free (see Ex.~\ref{ex:19.2-ex}). The elementary
proof given above is taken from a lecture by M. Narita in 1971.

\begin{theorem}[Auslander and Buchsbaum~\cite{AB3}]\label{thm:20.3}
A regular local ring is a
UFD.
\end{theorem}

\begin{proof}
Let $(A, m)$ be a regular local ring: the proof works by induction
on $\dim A$. If $\dim A = 0$ then $A$ is a field and therefore (trivially) a UFD.
If $\dim A = 1$ then $A$ is a DVR, and therefore a UFD. We suppose that
$\dim A > 1$ and choose $x \in m - m^2$; then since $xA$ is a prime ideal, applying
Theorem~\hyperref[thm:20.2]{2} to $\Gamma = \{x\}$, we need only show that $A_x$ is a UFD (where
$A_x = A[x^{-1}]$ is as on p.~22). Let $P$ be a height 1 prime ideal of $A_x$, and
set $p = P \cap A$; we have $P = pA_x$. Since $A$ is a regular local ring, the
$A$-module $p$ has an FFR, so that the $A_x$-module $P$ has an FFR. For any
prime ideal $Q$ of $A_x$, the ring $(A_x)_Q = A_{Q \cap A}$ is a regular local ring of
dimension less than that of $A$, so by induction is a UFD. Thus $P_Q$ is free
as an $(A_x)_Q$-module, so that by Theorem~\ref{thm:7.12}, the $A_x$-module $P$ is projective;
hence by Ex.~\ref{ex:19.2-ex}, $P$ is stably free, and therefore by the previous lemma,
$P$ is a principal ideal of $A_x$.
\end{proof}

The above proof is due to Kaplansky. Instead of our Lemma~\hyperref[lem:20.1]{1}, he used
the following more general proposition, which he had previously proved:
if $A$ is an integral domain, and $I_i, J_i$ are ideals of $A$ for $1 \le i \le r$ such that
$\bigoplus_{i=1}^{r} I_i \cong \bigoplus_{i=1}^{r} J_i$, then $I_1 \cdots I_r \cong J_1 \cdots J_r$. This is an interesting property
of ideals, and we have given a proof in Appendix C.

\begin{theorem}\label{thm:20.4}
Let $A$ be a Noetherian integral domain. Then if any finite
$A$-module has an FFR, $A$ is a UFD.
\end{theorem}

\begin{proof}
By Ex.~\ref{ex:19.4-ex}, $A$ is a regular ring. Let $P$ be a height 1 prime ideal
of $A$. Then $A_m$ is a regular local ring for any $m \in \Spec A$, so by the previous
theorem, the ideal $P_m$ is principal, and is therefore a free $A_m$-module. Hence
by Theorem~\ref{thm:7.12}, $P$ is projective. Therefore by Ex.~\ref{ex:19.2-ex}, $P$ is stably free,
and so by Lemma~\hyperref[lem:20.1]{1} is principal.
\end{proof}

Let $A$ be an integral domain; for any two non-zero elements $a, b \in A$,
the notion of greatest common divisor (g.c.d.) and least common multiple
(1.c.m.) are defined as in the ring of integers. That is, $d$ is a g.c.d. of $a$ and
$b$ if $d$ divides both $a$ and $b$, and any element $x$ dividing both $a$ and $b$
divides $d$; and $e$ is an 1.c.m. of $a$ and $b$ if $e$ is divisible by both $a$ and $b$,
and any $y$ divisible by $a$ and $b$ is divisible by $e$; this condition is equivalent
to $(e) = (a) \cap (b)$.

\begin{lemma}\label{lem:20.2}
If an 1.c.m. of $a$ and $b$ exists then so does a g.c.d.
\end{lemma}

\begin{proof}
If $(a) \cap (b) = (e)$ then there exists $d$ such that $ab = ed$. From $e \in (a)$
we get $b \in (d)$ and similarly $a \in (d)$, so that $(a, b) \subset (d)$. Now if $x$ is a common
divisor of $a$ and $b$ then $a = xt$ and $b = xs$, so that $xst$ is a common multiple
of $a, b$, and is hence divisible by $e$. Then from $ed = ab = x \cdot xst$ we get that
$d$ is divisible by $x$. Therefore, $d$ is a g.c.d. of $a$ and $b$.
\end{proof}

\begin{remark}\label{rem:20.1}
If $A$ is a Noetherian integral domain which is not a UFD then
$A$ has an irreducible element $a$ which is not prime. If $xy \in (a)$ but $x \notin (a)$,
$y \notin (a)$ then the only common divisors of $a$ and $x$ are units, so that 1 is a
g.c.d. of $a$ and $x$. However, $xy \in (a) \cap (x)$, but $xy \notin (ax)$, so that $(a) \cap
(x) \neq (ax)$, and there does not exist any 1.c.m. of $a$ and $x$. Thus the converse
of Lemma~\hyperref[lem:20.2]{2} does not hold in general.
\end{remark}

\begin{remark}\label{rem:20.2}
If $A$ is a UFD then an intersection of an arbitrary collection
of principal ideals is again principal (we include (0)). Indeed, if $\bigcap_i a_iA \neq 0$,
then factorise each $a_i$ as a product of primes:
\[
  a_i = u_i \prod_p p^{n_{ip}},
\]
with $u_i$ units, and $p, \pi$ prime elements such that $pA \neq \pi A$ for $p \neq \pi$.
(We could even allow the $a_i$ to be elements of the field of fractions of $A$.)
Then $\bigcap_i a_iA = dA$, where $d = \prod_p p^{\min_i n_{ip}}$.
\end{remark}

\begin{theorem}\label{thm:20.5}
An integral domain $A$ is a UFD if and only if the ascending
chain condition holds for principal ideals, and any two elements of $A$ have
an 1.c.m.
\end{theorem}

\begin{proof}
The `only if' is already known, and we prove the `if'. From the first
condition it follows that every element which is neither 0 nor a unit can
be written as a product of a finite number of irreducible elements, so that
we need only prove that an irreducible element is prime. Let $a$ be an
irreducible element, and let $xy \in (a)$ and $x \notin (a)$. By assumption we can
write $(a) \cap (x) = (z)$; now 1 is a g.c.d. of $a$ and $x$, so that one sees from the
proof of Lemma~\hyperref[lem:20.2]{2} that $(z) = (ax)$, and then $xy \in (a) \cap (x) = (ax)$ implies
that $y \in (a)$. Therefore $(a)$ is prime.
\end{proof}

\begin{theorem}\label{thm:20.6}
Let $A$ be a regular ring and $u, v \in A$. Then $uA \cap vA$ is a
projective ideal.
\end{theorem}

\begin{proof}
$A_m$ is a UFD for every maximal ideal $m$, so that $(uA \cap vA)A_m =
uA_m \cap vA_m$ is a principal ideal, and hence a free module.
\end{proof}

\begin{theorem}\label{thm:20.7}
If $A$ is a UFD then a projective ideal is principal.
\end{theorem}

\begin{proof}
By Theorem~\ref{thm:11.3}, it is equivalent to say that a non-zero ideal $a$ is
projective or invertible. Hence if we set $K$ for the field of fractions of $A$,
then there exist $u_i \in K$ such that $u_ia \subset A$ and $a_i \in a$ such that $\sum u_ia_i = 1$.
We have $a \subset \bigcap u_i^{-1}A$, and conversely if $x \in \bigcap u_i^{-1}A$ then $x =
\sum (xu_i)a_i \in a$, and hence $a = \bigcap u_i^{-1}A$; now since $A$ is a UFD, the intersection
of principal fractional ideals is again principal.
\end{proof}

\begin{theorem}\label{thm:20.8}
If $A$ is a regular UFD then so is $A[X]$.
\end{theorem}

\begin{proof}
Set $B = A[X]$. By Theorem~\hyperref[thm:20.5]{5}, it is enough to prove that $uB \cap vB$
is principal for $u, v \in B$; set $a = uB \cap vB$. Then by Theorem~\hyperref[thm:20.6]{6} and
Theorem~\ref{thm:19.5}, $a$ is projective, so that
\[
  a \otimes_A A = a \otimes_A (B/XB) = a/Xa
\]
is projective as an $A$-module. Suppose that $a = X^rb$ with $b \nsubseteq XB$; then
$a/Xa \cong b/Xb$, so that $b$ is isomorphic to $a$, hence projective, and therefore
locally principal. $B$ is a regular ring, so that the prime divisors of $b$ all
have height 1. Since $XB$ is also a height 1 prime ideal and $b \nsubseteq XB$ we
have $b:XB = b$, hence $b \cap XB = Xb$. Therefore since we can view $b/Xb$
as $b/Xb = b/b \cap XB \subset B/XB = A$, by Theorem~\hyperref[thm:20.7]{7} it is principal, hence
$b = yB + Xb$ for some $y \in b$; then by NAK, $b = yB$, so that $a = X^ryB$.
\end{proof}

\begin{remark}
There are examples where $A$ is a UFD but $A[X]$ is not.
\end{remark}

It is easy to see that a UFD is a Krull ring. For any Krull ring $A$ we
can define the \emph{divisor class group}\index{divisor class group} of $A$, which should be thought of as a
measure of the extent to which $A$ fails to be a UFD. We can give the
definition in simple terms as follows: let $\mathscr{P}$ be the set of height 1 prime
ideals of the Krull ring $A$, and $D(A)$ the free Abelian group on $\mathscr{P}$. That
is, $D(A)$ consists of formal sums $\sum_{p \in \mathscr{P}} n_p \cdot p$ (with $n_p \in \mathbb{Z}$ and all but finitely
many $n_p = 0$), with addition defined by
\[
  \left(\sum n_pp\right) + \left(\sum n_p'p\right) = \sum (n_p + n_p')p.
\]
Let $K$ be the field of fractions of $A$, and $K^*$ the multiplicative group of
non-zero elements of $K$, and for $a \in K^*$ set $\divisor(a) = \sum_{p \in \mathscr{P}} v_p(a) \cdot p$, where $v_p$
is the normalised additive valuation of $K$ corresponding to $p$. Then
$\divisor(ab) = \divisor(a) + \divisor(b)$, so that $\divisor$ is a homomorphism from $K^*$ to $D(A)$. We
write $F(A)$ for the image of $K^*$; this is a subgroup of $D(A)$, so that
we can define $C(A) = D(A)/F(A)$ to be the divisor class group of $A$.

Now, if $A$ is a UFD then each $p \in \mathscr{P}$ is principal, and if $p = aA$ then for
an element of $D(A)$ we have $p = \divisor(a)$, so that $C(A) = 0$. Conversely, if
$C(A) = 0$ then each $p \in \mathscr{P}$ is a principal ideal, and putting this together
with the corollary of Theorem~\ref{thm:12.3}, one sees easily that $A$ is a UFD. Hence
\[
  A \text{ is a UFD} \iff C(A) = 0.
\]

Now let $A$ be any ring, and $M$ a finite projective $A$-module. For each
$P \in \Spec A$, the localisation $M_P$ is a free module over $A_P$, and we write
$n(P)$ for its rank. Then $n$ is a function on $\Spec A$, and is constant on every
connected component (since $n(P) = n(Q)$ if $P \supset Q$). This function $n$ is called
the \emph{rank} of $M$. If the rank is a constant $r$ over the whole of $\Spec A$ then
we say that $M$ is a \emph{projective module of rank $r$}\index{projective module of rank r}. We write $\Pit(A)$ for the
set of isomorphism classes of finite projective $A$-modules of rank 1; $\cl(M)$
denotes the isomorphism class of $M$. If $M$ and $N$ are finite projective rank
1 modules then so is $M \otimes_A N$; this is clear on taking localisations. Thus
we can define a sum in $\Pit(A)$ by setting
\[
  \cl(M) + \cl(N) = \cl(M \otimes N).
\]
We set $M^* = \Hom_A(M, A)$, and define $\varphi\colon M \otimes M^* \to A$ by
\[
  \varphi\left(\sum m_i \otimes f_i\right) = \sum f_i(m_i);
\]
then $\varphi$ is an isomorphism (taking localisations and using the corollary to
Theorem~\ref{thm:7.11} reduces to the case $M = A$, which is clear). Hence $\cl(M^*) =
-\cl(M)$, and $\Pit(A)$ becomes an Abelian group, called the \emph{Picard group}\index{Picard group}
of $A$. If $A$ is local then $\Pit(A) = 0$.

If $A$ is an integral domain with field of fractions $K$, then $M_{(0)} = M \otimes K$,
so that the rank we have just defined coincides with the earlier definition
(after Lemma~\hyperref[lem:20.1]{1}). If $M$ is a finite projective rank 1 module, then since $M$
is torsion-free we have $M \subset M_{(0)} \cong K$, so that $M$ is isomorphic as an
$A$-module to a fractional ideal; for fractional ideals, by Theorem~\ref{thm:11.3},
projective and invertible are equivalent conditions, so that for an integral
domain $A$, we can consider $\Pit(A)$ as a quotient of the group of invertible
fractional ideals under multiplication. A fractional ideal $I$ is isomorphic
to $A$ as an $A$-module precisely when $I$ is principal, so that
\[
  \Pit(A) = (\text{invertible fractional ideals})/(\text{principal fractional ideals}).
\]

Suppose in addition that $A$ is a Krull ring. Then we can view $\Pit(A)$
as a subgroup of $C(A)$. To prove this, for $p \in \mathscr{P}$ and $I$ a fractional ideal,
set
\[
  v_p(I) = \min \{v_p(x) \mid x \in I\};
\]
this is zero for all but finitely many $p \in \mathscr{P}$ (check this!), so that we can set
\[
  \divisor(I) = \sum_{p \in \mathscr{P}} v_p(I) \cdot p \in D(A).
\]
For a principal ideal $I = aA$ we have $\divisor(I) = \divisor(a)$. One sees easily that
$\divisor(II') = \divisor(I) + \divisor(I')$, and that $\divisor(A) = 0$, so that if $I$ is invertible,
$\divisor(I) = -\divisor(I^{-1})$.

For invertible $I$ we have $(I^{-1})^{-1} = I$: indeed, $I \subset (I^{-1})^{-1}$ from the
definition, and $I = I \cdot A \supset I(I^{-1}(I^{-1})^{-1}) = (I^{-1})^{-1}$, so $I = (I^{-1})^{-1}$. If $I$ is invertible and
$\divisor(I) = 0$ then $\divisor(I^{-1}) = 0$, so that $I \subset A$, $I^{-1} \subset A$; hence $A \subset (I^{-1})^{-1} = I$,
and $I = A$. It follows that if $I, I'$ are invertible, $\divisor(I) = \divisor(I')$ implies $I = I'$.
Thus we can view the group of invertible fractional ideals as a subgroup of
$D(A)$, and $\Pit(A)$ as a subgroup of $C(A)$.

If $A$ is a regular ring then as we have seen, $p \in \mathscr{P}$ is a locally free module,
and so is invertible. Clearly from the definition, $\divisor(p) = p$. Hence, in
the case of a regular ring, $D(A)$ is identified with the group of invertible
fractional ideals, and $C(A)$ coincides with $\Pit(A)$.

The notions of $D(A)$ and $\Pit(A)$ originally arise in algebraic geometry.
Let $V$ be an algebraic variety, supposed to be irreducible and normal. We
write $\mathscr{W}$ for the set of irreducible codimension 1 subvarieties of $V$, and
define the group of divisors $D(V)$ of $V$ to be the free Abelian group on $\mathscr{W}$;
a divisor (or Weil divisor) is an element of $D(V)$. Corresponding to a
rational function $f$ on $V$ and an element $W \in \mathscr{W}$, let $v_W(f)$ denote the
order of zero of $f$ along $W$, or minus the order of the pole if $f$ has a pole
along $W$. Write $\divisor(f) = \sum_{W \in \mathscr{W}} v_W(f) \cdot W$ for the divisor of $f$ on $V$ (or just
$\divisor(f)$). For $W \in \mathscr{W}$, the local ring $\mathcal{O}_W$ of $W$ on $V$ is a DVR of the function
field of $V$, and $v_W$ is the corresponding valuation. We say that two divisors
$M, N$ in $D(V)$ are linearly equivalent if their difference $M - N$ is the divisor of a
function, and write $M \sim N$. The quotient group of $D(V)$ by $\sim$, that is the
quotient by the subgroup of divisors of functions, is the divisor class group
of $V$ (up to linear equivalence), and we write $C(V)$ for this. (In addition
to linear equivalence one also considers other equivalence relations with
certain geometric significance (algebraic equivalence, numerical equivalence,
$\ldots$), and divisor class groups, quotients of $D(V)$ by the corresponding subgroups.)

A divisor $M$ on $V$ is said to be a \emph{Cartier divisor} if it is the divisor of a
function in a neighbourhood of every point of $V$. From a Cartier divisor $M$
one constructs a line bundle over $V$, and two Cartier divisors give rise to
isomorphic line bundles if and only if they are linearly equivalent. Cartier
divisors form a subgroup of $D(V)$, and their class group up to linear
equivalence is written $\Pit(V)$; this can also be considered as the group of
isomorphism classes of line bundles over $V$ (with group law defined by
tensor product). If $V$ is smooth then (by Theorem~\hyperref[thm:20.3]{3}) there is no distinction
between Cartier and Weil divisors, and $C(V) = \Pit(V)$.

The reader familiar with algebraic geometry will know that the divisor
class group and Picard group of a Krull ring are an exact translation of
the corresponding notions in algebraic geometry. If $V$ is an affine variety,
with coordinate ring $k[V] = A$ then $C(V) = C(A)$ and $\Pit(V) = \Pit(A)$. In
this case, to say that $A$ is a UFD expresses the fact that every codimension
1 subvariety of $V$ can be defined as the intersection of $V$ with a hypersurface.
If $V \subset \mathbb{P}^n$ is a projective algebraic variety, defined by a prime ideal
$I \subset k[X_0, \ldots, X_n]$, and we set $A = k[X]/I = k[t_0, \ldots, t_n]$ (with $t_i$ the
class of $X_i$) then $A$ is the so-called homogeneous coordinate ring of $V$. If
$A$ is integrally closed we say that $V$ is projectively normal (also
arithmetically normal). This condition is stronger than saying that $V$ is
normal (the local ring of any point of $V$ is normal). If $A$ is a UFD then
every codimension 1 subvariety of $V$ can be given as the intersection in
$\mathbb{P}^n$ of $V$ with a hypersurface. Let $m = (t_0, \ldots, t_n)$ be the homogeneous
maximal ideal of $A$, and write $R = A_m$ for the localisation. The above
statement holds if we just assume that $R$ is a UFD; see Ex.~\ref{ex:20.6-ex}. All the
information about $V$ is contained in the local ring $R$.

Thus $C(A)$, $\Pit(A)$ and the UFD condition are notions with important
geometrical meaning, and methods of algebraic geometry can also be used
in their study. For example, in this way Grothendieck~\cite{G5} was able to
prove the following theorem conjectured by Samuel: let $R$ be a regular
local ring, $P$ a prime ideal generated by an $R$-sequence, and set $A = R/P$;
if $A_p$ is a UFD for every $p \in \Spec A$ with $\htop p \le 3$ then $A$ is a UFD.

We do not have the space to discuss $C(A)$ and $\Pit(A)$ in detail, and we
just mention the following two theorems as examples:
\begin{enumerate}[(1)]
  \item If $A$ is a Krull ring then $C(A) \cong C(A[X])$.
        This generalises the well-known theorem (see Ex.~\ref{ex:20.2-ex}) that if $A$ is a UFD
        then so is $A[X]$.
  \item If $A$ is a regular ring then $C(A) \cong C(A[[X]])$.
        This generalises Theorem~\hyperref[thm:20.8]{8}.
\end{enumerate}
Finally we give an example. Let $k$ be a field of characteristic 0, and set
$A = k[X, Y, Z]/(Z^n - XY) = k[x, y, z]$ for some $n > 1$. Then $A/(z, x) \cong
k[X, Y, Z]/(X, Z) \cong k[Y]$, so that $p = (x, z)$ is a height 1 prime ideal of $A$.
In $D(A)$ we have $np = \divisor(x)$, and it can be proved that $C(A) \cong \mathbb{Z}/d$ (see
\cite{S2}, p.~58). The relation $xy = z^n$ shows that $A$ is not a UFD.

For those wishing to know more about UFDs, consult~\cite{K},~\cite{S2} and
\cite{F}.

% ============ 习题 ============
\section*{Exercises to \S\ref{sec:20}}
\emph{Prove the following propositions.}

\begin{exercise}\label{ex:20.1-ex}
(Gauss' lemma) Let $A$ be a UFD, and $f(X) = a_0 + a_1X + \cdots + a_nX^n \in
A[X]$; say that $f$ is \emph{primitive} if the g.c.d. of the coefficients $a_0, \ldots, a_n$
is 1. Then if $f(X)$ and $g(X)$ are primitive, so is $f(X)g(X)$.
\end{exercise}

\begin{exercise}\label{ex:20.2-ex}
If $A$ is a UFD so is $A[X]$ (use the previous question).
\end{exercise}

\begin{exercise}\label{ex:20.3-ex}
If $A$ is a UFD and $q_1, \ldots, q_r$ are height 1 primary ideals then $q_1 \cap \cdots \cap q_r$ is a
principal ideal.
\end{exercise}

\begin{exercise}\label{ex:20.4-ex}
Let $A$ be a Zariski ring (see \S\ref{sec:8}) and $\hat{A}$ the completion of $A$. Then if $A$ is
a UFD so is $\hat{A}$ (there are counter-examples to the converse).
\end{exercise}

\begin{exercise}\label{ex:20.5-ex}
Let $A$ be an integral domain. We say that $A$ is \emph{locally UFD} if $A_m$ is a UFD
for every maximal ideal $m$. If $A$ is a semilocal integral domain and $A$ is
locally UFD, then $A$ is a UFD.
\end{exercise}

\begin{exercise}\label{ex:20.6-ex}
Let $R = \bigoplus_{n \ge 0} R_n$ be a graded ring, and suppose that $R_0$ is a field. Set
$m = \bigoplus_{n > 0} R_n$. If $I$ is a homogeneous ideal of $R$ such that $IR_m$ is principal
then there is a homogeneous element $f \in I$ such that $I = fR$.
\end{exercise}

% ============ §21 Complete intersection rings (原书 pp.169-173) ============
\section{Complete intersection rings}\label{sec:21}\index{complete intersection rings, section 21}\index{c.i. ring, section 21}

Let $(A, m, k)$ be a Noetherian local ring; we choose a minimal basis
$x_1, \ldots, x_n$ of $m$, where $n = \embdim A$ is the embedding dimension (see
\S\ref{sec:14}). Set $E_\bullet = K_\bullet(x_1, \ldots, x_n)$ for the Koszul complex. The complex $E_\bullet$ is
determined by $A$ up to isomorphism. Indeed, if $x_1', \ldots, x_n'$ is another
minimal basis of $m$ then by Theorem~\ref{thm:2.3}, there is an invertible $n \times n$
matrix $(a_{ij})$ over $A$ such that $x_i = \sum a_{ij}x_j'$. It is proved in Appendix C that
$K_\bullet(x_1, \ldots, x_n)$ can be thought of as the exterior algebra $\bigwedge(Ae_1 + \cdots + Ae_n)$ with
differential defined by $d(e_i) = x_i$. Similarly,
$K_\bullet(x_1', \ldots, x_n') = \bigwedge(Ae_1' + \cdots + Ae_n')$ with $d(e_i') = x_i'$.
Now $f(e_j') = \sum a_{ij}e_i$ defines an isomorphism from the free $A$-module
$Ae_1' + \cdots + Ae_n'$ to $Ae_1 + \cdots + Ae_n$, which extends to an isomorphism $f$
of the exterior algebra; $f$ commutes with the differential $d$, since for a
generator $e_i'$ of $\bigwedge(Ae_1' + \cdots + Ae_n')$ we have $df(e_i') = \sum a_{ij}x_j = x_i = fd(e_i')$.
Therefore $f\colon K_\bullet(x_1', \ldots, x_n') \to K_\bullet(x_1, \ldots, x_n)$ is an isomorphism of complexes.

Since $mH_p(E_\bullet) = 0$ by Theorem~\ref{thm:16.4}, $H_p(E_\bullet)$ is a vector space over
$k$, and we set
\[
  \varepsilon_p = \dim_k H_p(E_\bullet) \quad \text{for } p = 0, 1, 2, \ldots;
\]
then these are invariants of the local ring $A$. In view of $H_0(E_\bullet) = A/(x) =
A/m = k$, we have $\varepsilon_0 = 1$. In this section we are concerned with $\varepsilon_1$. If $A$ is
regular then $x_1, \ldots, x_n$ is an $A$-sequence, so that $\varepsilon_1 = \cdots = \varepsilon_n = 0$, and
conversely by Theorem~\ref{thm:16.5}, $\varepsilon_n = 0$ implies that $A$ is regular.

Let us consider the case when $A$ can be expressed as a quotient of a
regular local ring $R$; let $A = R/a$, and write $n$ for the maximal ideal of $R$.
If $a \nsubseteq n^2$, we can take $x \in a - n^2$; then $R' = R/xR$ is again a regular local
ring, and $A = R'/a'$, so that we can write $A$ as a quotient of a ring of
dimension smaller than $R$. In this way we see that there exists an
expression $A = R/a$ of $A$ as a quotient of a regular local ring $(R, n)$ with
$a \subset n^2$. Then we have $m = n/a$ and $m/m^2 = n/(a + n^2) = n/n^2$, so that
$\dim R = \embdim A$. Conversely, equality here implies that $a \subset n^2$.

Let $(R, n)$ be a regular local ring and $A = R/a$ with $a \subset n^2$; choose a
regular system of parameters (that is a minimal basis of $n$) $t_1, \ldots, t_n$.
Then the images $x_i$ of $t_i$ in $A$ form a minimal basis $x_1, \ldots, x_n$ of $m$. Let
\[
  K_\bullet(t_1, \ldots, t_n)\colon 0 \to L_n \to L_{n-1} \to \cdots \to L_1 \to L_0 \to 0
\]
be the Koszul complex of $R$ and $t$. By Theorem~\ref{thm:16.5}, we know that this
becomes exact on adding $\cdots + L_0 \to k \to 0$ to the right-hand end, so that
$K_\bullet(t_1, \ldots, t_n)$ is a projective resolution of $k$ as an $R$-module. Taking the tensor
product with $A = R/a$, we get the complex $E_\bullet = K_\bullet(x_1, \ldots, x_n)$ of $A$-modules.
Thus we have
\[
  H_p(E_\bullet) = H_p(K_\bullet(t_1, \ldots, t_n) \otimes_R A) = \Tor_p^R(k, A) \quad \text{for all } p \ge 0.
\]
However, from the exact sequence of $R$-modules $0 \to a \to R \to A \to 0$ we
get the long exact sequence
\[
  0 = \Tor_1^R(k, R) \to \Tor_1^R(k, A) \to k \otimes_R a \to k \otimes_R R
      \to k \otimes_R A \to 0;
\]
at the right-hand end we have $k \otimes_R R \cong k \otimes_R A = k$, so that
\[
  \Tor_1^R(k, A) \cong k \otimes_R a = a/na.
\]
Quite generally, we write $\mu(M)$ for the minimum number of generators of
an $R$-module $M$. Then we see that
\[
  \mu(a) = \dim_k H_1(E_\bullet) = \varepsilon_1(A).
\]

\begin{theorem}\label{thm:21.1}
Let $(A, m, k)$ be a Noetherian local ring, and $\hat{A}$ its
completion.
\begin{enumerate}[(i)]
  \item $\varepsilon_p(A) = \varepsilon_p(\hat{A})$ for all $p \ge 0$.
  \item $\varepsilon_1(A) \ge \embdim A - \dim A$.
  \item If $R$ is a regular local ring, $a$ an ideal of $R$ and $A \cong R/a$, then
        $\mu(a) = \dim R - \embdim A + \varepsilon_1(A)$.
\end{enumerate}
\end{theorem}

\begin{proof}
(i) is clear from the fact that a minimal basis of $m$ is a minimal
basis of $m\hat{A}$, so that applying $\otimes_A \hat{A}$ to the complex $E_\bullet$ made from $A$
gives that made from $\hat{A}$. Then since $\hat{A}$ is $A$-flat, $H_p(E_\bullet) \otimes \hat{A} = H_p(E_\bullet \otimes \hat{A})$,
and $mH_p(E_\bullet) = 0$ gives $H_p(E_\bullet) \otimes \hat{A} = H_p(E_\bullet)$.

(ii) If $A$ is a quotient of a regular local ring, then as we have seen above,
there exists a regular ring $(R, n)$ such that $A = R/a$ with $a \subset n^2$, so that
$\varepsilon_1(A) = \mu(a) \ge \htop a = \dim R - \dim A = \embdim A - \dim A$, where the
equality for $\htop a$ comes from Theorem~\ref{thm:17.4}, (i). Now $A$ itself is not necessarily
a quotient of a regular local ring, but we will prove later (see \S29) that $\hat{A}$
always is, and we admit this in the section. Having said this, the two sides
of (ii) are unaltered on replacing $A$ by $\hat{A}$, and the inequality holds for $\hat{A}$.

(iii) Set $n = \rad(R)$. If $a \subset n^2$ then, as we have seen above, $\mu(a) = \varepsilon_1(A)$,
and $\dim R = \embdim A$, so that we are done. If $a \nsubseteq n^2$, take $x \in a - n^2$;
we pass to $R/xR$ and $a/xR$, each of $\dim R$ and $\mu(a)$ decreases by 1, so
induction completes the proof.
\end{proof}

\begin{definition*}
A Noetherian local ring $A$ is a \emph{complete intersection ring}
(abbreviated to \emph{c.i. ring}) if $\varepsilon_1(A) = \embdim A - \dim A$.
\end{definition*}

\begin{theorem}\label{thm:21.2}
Let $A$ be a Noetherian local ring.
\begin{enumerate}[(i)]
  \item $A$ is c.i. $\iff$ $\hat{A}$ is c.i.
  \item Let $A$ be a c.i. ring and $R$ a regular local ring such that $A = R/a$;
        then $a$ is generated by an $R$-sequence. Conversely, if $a$ is an ideal generated
        by an $R$-sequence then $R/a$ is a c.i. ring.
  \item A necessary and sufficient condition for $A$ to be a c.i. ring is that
        the completion $\hat{A}$ should be a quotient of a complete regular local ring
        by an ideal generated by an $R$-sequence.
\end{enumerate}
\end{theorem}

\begin{proof}
(i) is obvious.

(ii) By Theorem~\hyperref[thm:21.1]{1}, (iii), $\mu(a) = \dim R - \embdim A + \varepsilon_1(A)$, and by
Theorem~\ref{thm:17.4}, (i), $\htop a = \dim R - \dim A$, so that $A$ is a c.i. ring is equivalent
to $\htop a = \mu(a)$. But by Theorem~\ref{thm:17.4}, (iii), this is equivalent to $a$ being
generated by an $R$-sequence.

(iii) The sufficiency is clear from (i) and (ii). Necessity follows from the
fact that $\hat{A}$ is a quotient of a complete regular local ring (see \S29), together
with (i) and (ii).
\end{proof}

\begin{theorem}\label{thm:21.3}
A c.i. ring is Gorenstein.
\end{theorem}

\begin{proof}
If $A$ is c.i. then so is $\hat{A}$, and if $\hat{A}$ is Gorenstein then so is $A$, so that we
can assume that $A$ is complete. Then we can write $A = R/a$, where $R$ is a
regular local ring and $a$ is an ideal generated by a regular sequence. Since $R$
is Gorenstein, $A$ is also by Ex.~\ref{ex:18.1-ex}.
\end{proof}

Thus we have the following chain of implications for Noetherian local
rings:
\[
  \text{regular} \implies \text{c.i.} \implies \text{Gorenstein} \implies \text{CM}.
\]
Let $A$ be a c.i. ring, and $p$ a prime ideal of $A$. If $A$ is of the form
$R/(x_1, \ldots, x_r)$, where $R$ is regular and $x_1, \ldots, x_r$ is an $R$-sequence, then
since $A_p$ can be written $A_p = R_p/(x_1, \ldots, x_r)$, where $R_p$ is regular and
$x_1, \ldots, x_r$ is an $R_p$-sequence, it follows that $A_p$ is again a c.i. ring. The
question of deciding whether $A_p$ is still a c.i. ring even if $A$ is not a quotient
of a regular local ring remained unsolved for some time, but was answered
affirmatively by Avramov~\cite{Av1}, making use of Andr\'e's homology theory
[1, 2]. % 注: 原书此处引 [1, 2]; André 在文献表中仅一条论文条目, 故保持纯文本
This theory defines homology and cohomology groups
$H_n(A, B, M)$ and $H^n(A, B, M)$ for $n \ge 0$ associated with a ring $A$, an $A$-algebra $B$
and a $B$-module $M$. The definition is complicated, but in any case these are
$B$-modules having various nice functorial properties. If $A$ is a
Noetherian local ring with residue field $k$ then
\[
  A \text{ is regular} \iff H_2(A, k, k) = 0,
\]
and
\[
  A \text{ is c.i.} \iff H_3(A, k, k) = 0;
\]
for $n \ge 3$ the statements $H_3(A, k, k) = 0$ and $H_n(A, k, k) = 0$ are equivalent.
Thus Andr\'e homology is particularly relevant to the study of regular and
c.i. rings.

% ============ 习题 ============
\section*{Exercises to \S\ref{sec:21}}
\emph{Prove the following propositions.}

\begin{exercise}\label{ex:21.1-ex}
Let $R$ be a regular ring, $I$ an ideal of $R$, and let $A = R/I$; then the subset
$\{p \in \Spec A \mid A_p \text{ is c.i.}\}$ is open in $\Spec A$ (use Theorem~\ref{thm:19.9}).
\end{exercise}

\begin{exercise}\label{ex:21.2-ex}
Let $A$ be a Noetherian local ring with $\embdim A = \dim A + 1$; if $A$ is CM
then it is c.i.
\end{exercise}

\begin{exercise}\label{ex:21.3-ex}
Let $k$ be a field, and set $A = k[X, Y, Z]/(X^2 - Y^2, Y^2 - Z^2, XY, YZ, ZX)$;
then $A$ is Gorenstein but not c.i.
\end{exercise}


% ============ Chapter 8. Flatness revisited (原书 pp.173-189) ============
\chapter{Flatness revisited}\label{ch:8}

The main theme of this chapter is flatness over Noetherian rings. In
\S\ref{sec:22} we prove a number of theorems known as the `local flatness criterion'
(the main result is Theorem~\ref{thm:22.3}). Together with Theorem~\ref{thm:23.1} in the
following section, this is extremely useful in applications.

In \S\ref{sec:23} we consider a flat morphism $A \to B$ of Noetherian local rings, and
investigate the remarkable relationships holding between $A$, $B$ and the fibre
ring $F = B/m_AB$. Roughly speaking, good properties of $B$ are usually inherited
by $A$, and sometimes by $F$. Conversely, in order for $B$ to inherit good
properties of $A$ one also requires $F$ to be good.

In \S\ref{sec:24} we discuss the so-called generic freeness theorem in the
improved form due to Hochster and Roberts (Theorem~\ref{thm:24.1}), and
investigate, following the ideas of Nagata, the openness of loci of points at
which various properties hold, arising out of Theorem~\ref{thm:24.3}, which states
that the set of points of flatness is open.

% ============ §22 The local flatness criterion (原书 pp.173-178) ============
\section{The local flatness criterion}\label{sec:22}\index{local flatness criterion, section 22}\index{flatness, local criterion, section 22}

\begin{theorem}\label{thm:22.1}
Let $A$ be a ring, $B$ a Noetherian $A$-algebra, $M$ a finite $B$-module, and
$J$ an ideal of $B$ contained in $\rad(B)$; set $M_n = M/J^{n+1}M$ for $n \ge 0$.
If $M_n$ is flat over $A$ for every $n \ge 0$, then $M$ is also flat over $A$.
\end{theorem}

\begin{proof}
According to Theorem~\ref{thm:7.7}, we need only show that for a finitely
generated ideal $I$ of $A$, the standard map $u\colon I \otimes_A M \to M$ is
injective. Set $I \otimes_A M = M'$; then $M'$ is also a finite $B$-module, and
hence is separated for the $J$-adic topology. Let $x \in \Ker(u)$; we prove that
$x \in \cap_n J^nM' = 0$. For each $n$ we have $M'/J^{n+1}M' = (I \otimes_A M)
\otimes_B B/J^{n+1} = I \otimes_A M_n$, and the standard map $I \otimes_A M_n \to
M_n$ is injective, by the assumption that $M_n$ is flat. Then we deduce that
$x \in J^{n+1}M'$ from the commutative diagram
\[
\begin{tikzcd}
M' = I \otimes_A M \arrow[r] \arrow[d, "\bar{u}"'] & I \otimes_A M_n \arrow[d, "u_n"] \\
M \arrow[r] & M_n
\end{tikzcd}
\]
in which the horizontal maps are the natural ones. Indeed, the image of $x$ in
$I \otimes_A M_n$ is mapped to $0$ in $M_n$, hence is $0$, so that
$x \in J^{n+1}M'$. Since this holds for all $n$, $x \in \cap_n J^nM' = 0$ (the
intersection being zero by the Krull intersection theorem, because $J \subset
\rad(B)$ and $M'$ is a finite $B$-module). Thus $u$ is injective, and $M$ is
flat over $A$.
\end{proof}

\begin{theorem}\label{thm:22.2}
Let $A$ be a ring, $B$ a Noetherian $A$-algebra, and $M$ a finite $B$-module;
suppose that $b$ is an $M$-regular element of $\rad(B)$. Then if $M/bM$ is flat
over $A$, so is $M$.
\end{theorem}

\begin{proof}
For each $i \ge 0$ the sequence $0 \to M/b^iM \xrightarrow{b} M/b^{i+1}M \to
M/bM \to 0$ is exact, so that by Theorem~\ref{thm:7.9} and an induction on $i$,
every $M/b^iM$ is flat over $A$. Thus we can just apply the previous theorem.
\end{proof}

\begin{definition*}
Let $A$ be a ring and $I$ an ideal of $A$; an $A$-module $M$ is said to be
\emph{$I$-adically ideal-separated}\index{I-adically ideal-separated, section 22} if $a \otimes_A M$ is separated for the $I$-adic
topology for every finitely generated ideal $a$ of $A$.

For example, if $B$ is a Noetherian $A$-algebra and $IB \subset \rad(B)$ then a
finite $B$-module $M$ is $I$-adically ideal-separated as an $A$-module.
\end{definition*}

Let $A$ be a ring, $I$ an ideal of $A$ and $M$ an $A$-module. Set $A_n = A/I^{n+1}$,
$M_n = M/I^{n+1}M$ for $n \ge 0$ and $\gr(A) = \bigoplus_{n \ge 0} I^n/I^{n+1}$,
$\gr(M) = \bigoplus_{n \ge 0} I^nM/I^{n+1}M$. There exist standard maps
\[
  \gamma_n\colon (I^n/I^{n+1}) \otimes_{A_0} M_0 \to I^nM/I^{n+1}M \quad \text{for } n \ge 0,
\]
and we can put together the $\gamma_n$ into a morphism of $\gr(A)$-modules
\[
  \gamma\colon \gr(A) \otimes_{A_0} M_0 \to \gr(M).
\]

\begin{theorem}\label{thm:22.3}
In the above notation, suppose that one of the following two conditions is
satisfied:
\begin{enumerate}[(a)]
  \item $I$ is a nilpotent ideal;
or ($\beta$) $A$ is a Noetherian ring and $M$ is $I$-adically ideal-separated.
\end{enumerate}
Then the following conditions are equivalent.
\begin{enumerate}[(1)]
  \item $M$ is flat over $A$;
  \item $\Tor^A_1(N, M) = 0$ for every $A_0$-module $N$;
  \item $M_0$ is flat over $A_0$, and $I \otimes_A M = IM$;
  \item[$3'$] $M_0$ is flat over $A_0$, and $\Tor^A_1(A_0, M) = 0$;
  \item $M_0$ is flat over $A_0$, and $\gamma_n$ is an isomorphism for every $n \ge 0$;
  \item[$4'$] $M_0$ is flat over $A_0$, and $\gamma$ is an isomorphism;
  \item $M_n$ is flat over $A_n$ for every $n \ge 0$.
\end{enumerate}
In fact, the implications (1) $\Rightarrow$ (2) $\Leftarrow$ (3) $\Leftarrow$ (3$'$) $\Leftarrow$ (4) $\Leftarrow$ (5)
hold without any assumption on $M$.
\end{theorem}

\begin{proof}
First of all, let $M$ be arbitrary.

(1) $\Rightarrow$ (2) is trivial.

(2) $\Rightarrow$ (3) If $N$ is an $A_0$-module then we have
\[
  N \otimes_A M = (N \otimes_{A_0} A_0) \otimes_A M = N \otimes_{A_0} M_0,
\]
and hence for an exact sequence $0 \to N_1 \to N_2 \to N_3 \to 0$ of
$A_0$-modules we get an exact sequence
\[
  0 = \Tor^A_1(N_3, M) \to N_1 \otimes_{A_0} M_0 \to N_2 \otimes_{A_0} M_0 \to
      N_3 \otimes_{A_0} M_0 \to 0;
\]
therefore $M_0$ is flat over $A_0$. Also, from the exact sequence $0 \to I \to
A \to A_0 \to 0$ we get an exact sequence
\[
  0 = \Tor^A_1(A_0, M) \to I \otimes_A M \to M \to M_0 \to 0,
\]
so that $I \otimes_A M = IM$.

(3) $\Rightarrow$ (3$'$) is clear, and (3$'$) $\Rightarrow$ (2) follows as follows: if
$N$ is an $A_0$-module, we can choose an exact sequence of $A_0$-modules
$0 \to R \to F_0 \to N \to 0$ with $F_0$ a free $A_0$-module. From this we get
the exact sequence
\[
  \Tor^A_1(F_0, M) = 0 \to \Tor^A_1(N, M) \to R \otimes_{A_0} M_0 \to F_0 \otimes_{A_0} M_0,
\]
and since $M_0$ is flat over $A_0$, the final arrow is injective, so that
$\Tor^A_1(N, M) = 0$.

(3) $\Rightarrow$ (4) We have already seen that $\gamma_0$ is the identity. For $n = 1$,
from the exact sequence $0 \to I^2 \to I \to I/I^2 \to 0$ we get
$\Tor^A_1(I/I^2, M) = 0$, so that from $0 \to I^2 \to I \to I/I^2 \to 0$ the sequence
$0 \to I^2 \otimes_A M \to I \otimes_A M \to (I/I^2) \otimes_A M \to 0$
is exact. From $I \otimes_A M = IM$ we get $I^2 \otimes_A M = I^2M$ and
$(I/I^2) \otimes_A M \cong IM/I^2M$. Proceeding similarly, from
$0 \to I^{n+1} \to I^n \to I^n/I^{n+1} \to 0$ we get by induction
$I^{n+1} \otimes_A M = I^{n+1}M$ and $(I^n/I^{n+1}) \otimes_A M \cong I^nM/I^{n+1}M$.
(4$'$) is just a restatement of (4).

(4) $\Rightarrow$ (5) We fix an $n > 0$ and prove that $M_n$ is flat over $A_n$. For
$i \le n$ we have a commutative diagram
\[
\begin{tikzcd}
0 \arrow[r] &
(I^{i+1}/I^{n+1}) \otimes_A M \arrow[r] \arrow[d, "\alpha_{i+1}"] &
(I^i/I^{n+1}) \otimes_A M \arrow[r] \arrow[d, "\alpha_i"] &
(I^i/I^{i+1}) \otimes_A M \arrow[r] \arrow[d, "\gamma_i"] & 0 \\
0 \arrow[r] &
I^{i+1}M/I^{n+1}M \arrow[r] &
I^iM/I^{n+1}M \arrow[r] &
I^iM/I^{i+1}M \arrow[r] & 0
\end{tikzcd}
\]
with exact rows (the top row is exact since $\Tor^A_1(I^i/I^{i+1}, M) = 0$). By
assumption $\gamma_i$ is an isomorphism, and since $\alpha_{i+1}$ is an isomorphism,
by downwards induction on $i$ we see that $\alpha_n, \ldots, \alpha_1$ are
isomorphisms. In particular,
\[
  \alpha_1\colon (I/I^{n+1}) \otimes_A M \xrightarrow{\cong} IM/I^{n+1}M,
\]
and $(I/I^{n+1}) \otimes_A M = I \otimes_{A_n} M_n = IM_n$,
so that the conditions in (3) are satisfied by $A_n$, $M_n$ and $I/I^{n+1}$.
Therefore by (3) $\Rightarrow$ (2), we have $\Tor^{A_n}_1(N, M_n) = 0$ for every
$A_0$-module $N$. Now if $N$ is an $A_i$-module then $IN$ and $N/IN$ are both
$A_{i-1}$-modules, and
\[
  0 \to IN \to N \to N/IN \to 0
\]
is exact, so that by induction on $i$ we get finally that $\Tor^{A_n}_1(N, M_n) = 0$
for all $A_n$-modules $N$. Therefore $M_n$ is a flat $A_n$-module.

Next, assuming either (a) or ($\beta$) we prove (5) $\Rightarrow$ (1). In case (a) we have
$A = A_n$ and $M = M_n$ for large enough $n$, so that this is clear. In case
($\beta$), let $a$ be a finitely generated ideal of $A$; we prove that the standard map
$j\colon a \otimes_A M \to M$ is injective. By hypothesis we have
$\cap_n I^n(a \otimes_A M) = 0$, so that we need only prove that $\Ker(j) \subset
I^n(a \otimes_A M)$ for all $n > 0$. For a fixed $n$, by the Artin-Rees lemma,
$I^k \cap a \subset I^na$ for sufficiently large $k > n$. We now consider the
natural map
\[
  a \otimes_A M \to (a/I^ka) \otimes_A M = (a \otimes_A M)/I^k(a \otimes_A M).
\]
Since $M_{k-1}$ is flat over $A_{k-1} = A/I^k$, the map
\[
  (a/I^ka) \otimes_A M = (a/I^ka) \otimes_{A_{k-1}} M_{k-1} \to M_{k-1}
\]
is injective, so that from the commutative diagram
\[
\begin{tikzcd}
a \otimes_A M \arrow[r] \arrow[d, "j"'] & (a/I^ka) \otimes_A M \arrow[d] \\
M \arrow[r] & M_{k-1}
\end{tikzcd}
\]
we get $\Ker(j) \subset \Ker(f) \subset \Ker(gf) = I^n(a \otimes_A M)$. This is what
we needed to prove.
\end{proof}

This theorem is particularly effective when $A$ is a Noetherian local ring and
$I$ is the maximal ideal, since if $A_0$ is a field, $M_0$ is automatically flat
over $A_0$ in (3)---(4$'$). Also, in this case, requiring $M_n$ to be flat over $A_n$ in
(5) is the same as requiring it to be a free $A_n$-module, by
Theorem~\ref{thm:7.10}.

We now discuss some applications of the above theorem.

\begin{theorem}\label{thm:22.4}
Let $(A, m)$ and $(B, n)$ be Noetherian local rings, $\hat{A}$ and $\hat{B}$
their respective completions, and $A \to B$ a local homomorphism.
\begin{enumerate}[(i)]
  \item For $M$ a finite $B$-module, set $\hat{M} = M \otimes_B \hat{B}$; then
        \[
          M \text{ is flat over } A \iff \hat{M} \text{ is flat over } \hat{A} \iff \hat{M} \text{ is flat over } A.
        \]
  \item Writing $M^*$ for the $(mB)$-adic completion of $M$ we have
        \[
          M \text{ is flat over } A \iff M^* \text{ is flat over } \hat{A} \iff M^* \text{ is flat over } A.
        \]
\end{enumerate}
\end{theorem}

\begin{proof}
(i) The first equivalence comes from the transitivity law for flatness, together
with the fact that $\hat{B}$ is faithfully flat over $B$; the second, from the fact that
both sides are equivalent to $\hat{B}/m^n\hat{B}$ being flat over $A/m^n$ for all $n > 0$.

(ii) All three conditions are equivalent to $M/m^nM$ being flat over $A/m^n$ for
all $n$.
\end{proof}

\begin{theorem}\label{thm:22.5}
Let $(A, m, k)$ and $(B, n, k')$ be Noetherian local rings, $A \to B$ a local
homomorphism, and $u\colon M \to N$ a morphism of finite $B$-modules. Then if
$N$ is flat over $A$, the following two conditions are equivalent:
\begin{enumerate}[(1)]
  \item $u$ is injective and $N/u(M)$ is flat over $A$;
  \item $\bar{u}\colon M \otimes_A k \to N \otimes_A k$ is injective.
\end{enumerate}
\end{theorem}

\begin{proof}
(1) $\Rightarrow$ (2) is easy, so we only give the proof of (2) $\Rightarrow$ (1). Suppose that
$x \in M$ is such that $u(x) = 0$; then $\bar{u}(\bar{x}) = 0$, so that $\bar{x} = 0$, in
other words, $x \in mM$. Now assuming $x \in m^nM$, we will deduce $x \in m^{n+1}M$.
Let $\{a_1, \ldots, a_r\}$ be a minimal basis of the $A$-module $m^n$, and write
$x = \sum a_i y_i$ with $y_i \in M$; then $0 = \sum a_i u(y_i)$. Since $N$ is flat over
$A$, by Theorem~\ref{thm:7.6} there exist $c_{ij} \in A$ and $z_j \in N$ such that
\[
  \sum_i a_i c_{ij} = 0 \quad \text{for all } j, \quad \text{and} \quad u(y_i) = \sum_j c_{ij}z_j \quad \text{for all } i.
\]
By choice of $a_1, \ldots, a_r$, all the $c_{ij} \in m$, and hence $u(y_i) \in mN$ and
$\bar{u}(\bar{y}_i) = 0$, so that $\bar{y}_i = 0$, and $y_i \in mM$. Therefore
$x \in m^{n+1}M$. We have proved that $x \in \cap_n m^nM = 0$, and hence $u$ is
injective. Now from $0 \to M \xrightarrow{u} N \to N/u(M) \to 0$ we get
$\Tor^A_1(k, N/u(M)) = 0$, so that by Theorem~\hyperref[thm:22.3]{3},
$N/u(M)$ is flat over $A$.
\end{proof}

\begin{corollary*}
Let $A$, $B$ and $A \to B$ be as above, and $M$ a finite $B$-module; set
$\bar{B} = B \otimes_A k = B/mB$, and for $x_1, \ldots, x_n \in n$ write $\bar{x}_i$ for the
images in $\bar{B}$ of $x_i$. Then the following conditions are equivalent:
\begin{enumerate}[(1)]
  \item $x_1, \ldots, x_n$ is an $M$-sequence and $M_n = M/\sum x_iM$ is flat over $A$;
  \item $\bar{x}_1, \ldots, \bar{x}_n$ is an $M \otimes_A k$-sequence and $M$ is flat over $A$.
\end{enumerate}
\end{corollary*}

\begin{proof}
(2) $\Rightarrow$ (1) follows at once from the theorem. For (1) $\Rightarrow$ (2) we must prove
that $M_i = M/(x_1M + \cdots + x_iM)$ is flat for $i = 1, \ldots, n$; but if $M_i$ is flat
over $A$ then by Theorem~\hyperref[thm:22.2]{2}, so is $M_{i-1}$.
\end{proof}

\begin{theorem}\label{thm:22.6}
Let $A$ be a Noetherian ring, $B$ a Noetherian $A$-algebra, $M$ a finite
$B$-module, and $b \in B$ a given element. Suppose that $M$ is flat over $A$ and
that $b$ is $M/(P \cap A)$-regular for every maximal ideal $P$ of $B$; then $b$ is
$M$-regular and $M/bM$ is flat over $A$.
\end{theorem}

\begin{proof}
Write $K$ for the kernel of $M \xrightarrow{b} M$; then $K = \cap_P K_P = 0$ for all $P$.
Hence $b$ is $M$-regular if and only if $b$ is $M_P$-regular for all $P$. Moreover,
according to Theorem~\ref{thm:7.1}, $A$-flatness is also a local property in both $A$
and $B$, so that we can replace $B$ by $B_P$ (for a maximal ideal $P$ of $B$), $A$ by
$A_{(P \cap A)}$ and $M$ by $M_P$, and this case reduces to Theorem~\hyperref[thm:22.5]{5}.
\end{proof}

\begin{corollary*}
Let $A$ be a Noetherian ring, $B = A[X_1, \ldots, X_n]$ the polynomial ring over
$A$, and let $f \in B$. If the ideal of $A$ generated by the coefficients of $f$
contains $1$ then $f$ is a non-zero-divisor of $B$, and $B/fB$ is flat over $A$.
The same thing holds for the formal power series ring $B = A[[X_1, \ldots, X_n]]$.
\end{corollary*}

\begin{proof}
The polynomial ring is a free $A$-module, and therefore flat; the formal power
series ring is flat by Ex.~\ref{ex:7.4-ex}. Furthermore, for $\mathfrak{p} \in \Spec A$, if
$B = A[X_1, \ldots, X_n]$ then $B/\mathfrak{p}B = (A/\mathfrak{p})[X_1, \ldots, X_n]$, and in the formal
power series case we also have $B/\mathfrak{p}B = (A/\mathfrak{p})[[X_1, \ldots, X_n]]$ since
$\mathfrak{p}$ is finitely generated. In either case $B/\mathfrak{p}B$ is an integral domain,
so that the assertion follows directly from the theorem.
\end{proof}

\begin{remark}[\textbf{Flatness of a graded module}]\index{graded module, flatness of, section 22}
Let $G$ be an Abelian group, $R = \bigoplus_{g \in G} R_g$ a $G$-graded ring and
$M = \bigoplus_{g \in G} M_g$ a graded $R$-module, not necessarily finitely
generated.
\begin{enumerate}[(1)]
  \item The following three conditions are equivalent:
        \begin{enumerate}[(a)]
          \item $M$ is $R$-flat;
          \item If $\mathcal{F}\colon \cdots \to N \to N' \to N'' \to \cdots$ is an exact
                sequence of graded $R$-modules and $R$-linear maps preserving
                degrees, then $\mathcal{F} \otimes_R M$ is exact;
          \item $\Tor^R_1(M, R/H) = 0$ for every finitely generated homogeneous
                ideal $H$ of $R$.
        \end{enumerate}
        The proof is left to the reader as an exercise, or can be found in
        Herrmann and Orbanz~[3]. Using this criterion one can adapt the proof of
        Theorem~\hyperref[thm:22.3]{3} to prove the following graded version.
  \item Let $I$ be a (not necessarily homogeneous) ideal of $R$. Suppose that
        \begin{enumerate}[(i)]
          \item for every finitely generated homogeneous ideal $H$ of $R$, the
                $R$-module $H \otimes_R M$ is $I$-adically separated;
          \item $M_0 = M/IM$ is $R/I$-flat;
          \item $\Tor^R_1(M, R/I) = 0$.
        \end{enumerate}
        Then $M$ is $R$-flat.
  \item As an application one can prove the following:
\end{enumerate}
\begin{enumerate}[(3)]
  \item Let $A = \bigoplus_{n \ge 0} A_n$ and $B = \bigoplus_{n \ge 0} B_n$ be graded Noetherian
        rings. Assume that $A_0$, $B_0$ are local rings with maximal ideals $m$, $n$ and
        set $M = m + A_1 + A_2 + \cdots$, $N = n + B_1 + B_2 + \cdots$; let
        $f\colon A \to B$ be a ring homomorphism of degree $0$ such that $f(m) \subset n$.
        Then the following are equivalent:
        \begin{enumerate}[(a)]
          \item $B$ is $A$-flat;
          \item $B_N$ is $A_M$-flat;
          \item $B_0$ is $A_0$-flat.
        \end{enumerate}
\end{enumerate}
\end{remark}

% ============ 习题 ============
\section*{Exercises to \S\ref{sec:22}}
\emph{Prove the following propositions.}

\begin{exercise}\label{ex:22.1-ex}
(The Nagata flatness theorem, see [N1], p.~65). Let $(A, m, k)$ and $(B, n, k')$
be Noetherian local rings, and suppose that $A \subset B$ and that $mB$ is an
$n$-primary ideal. We say that the \emph{transition theorem} holds between $A$ and
$B$ if $l(A/q) \cdot l(B/mB) = l(B/qB)$ for every $m$-primary ideal $q$ of $A$. This
holds if and only if $B$ is flat over $A$.
\end{exercise}

\begin{exercise}\label{ex:22.2-ex}
Let $(A, m)$ be a Noetherian local ring, and $k \subset A$ a subfield. If
$x_1, \ldots, x_n \in m$ is an $A$-sequence then $x_1, \ldots, x_n$ are algebraically
independent over $k$, and $A$ is flat over $C = k[x_1, \ldots, x_n]$
(Hartshorne [2]).
\end{exercise}

\begin{exercise}\label{ex:22.3-ex}
Let $(A, m, k)$ be a Noetherian local ring, $B$ a Noetherian $A$-algebra, and $M$
a finite $B$-module. Suppose that $mB \subset \rad(B)$. If $x \in m$ is both
$A$-regular and $M$-regular, and if $M/xM$ is flat over $A/xA$ then $M$ is flat
over $A$.
\end{exercise}

\begin{exercise}\label{ex:22.4-ex}
Let $A$ be a Noetherian ring and $B$ a flat Noetherian $A$-algebra; if $I$ and $J$
are ideals of $A$ and $B$ such that $IB \subset J$ then the $J$-adic completion of
$B$ is flat over the $I$-adic completion of $A$.
\end{exercise}

% ============ §23 Flatness and fibres (原书 pp.178-185) ============
\section{Flatness and fibres}\label{sec:23}\index{flatness and fibres, section 23}\index{fibre ring, section 23}\index{normal ring, criterion of Serre, section 23}\index{regularity and flatness, section 23}

Let $(A, m)$ and $(B, n)$ be Noetherian local rings, and $\varphi\colon A \to B$ a local
homomorphism. We set $F = B \otimes_A k(m) = B/mB$ for the \emph{fibre ring}\index{fibre ring} of
$\varphi$ over $m$. If $B$ is flat over $A$ then according to Theorem~\ref{thm:15.1}, we
have
\[
  (*) \quad \dim B = \dim A + \dim F.
\]

As the following shows, under certain conditions the converse holds.

\begin{theorem}\label{thm:23.1}
Let $A$, $B$ and $F$ be as above. If $A$ is a regular local ring, $B$ is
Cohen-Macaulay, and $\dim B = \dim A + \dim F$ then $B$ is flat over $A$.
\end{theorem}

\begin{proof}
We use induction on $\dim A$. If $\dim A = 0$ then $A$ is a field, and we are
done. If $\dim A > 0$, take $x \in m - m^2$ and set $A' = A/xA$ and $B' = B/xB$. Then
\[
  \dim B' \le \dim A' + \dim F = \dim A - 1 + \dim F = \dim B - 1,
\]
and using a system of parameters of $B'$ one sees that $\dim B' \ge \dim B - 1$, so
that
\[
  \dim B' = \dim A' + \dim F = \dim B - 1.
\]
One sees easily from this that $x$ is $B$-regular and $B'$ is a CM ring. Hence by
induction $B'$ is flat over $A'$. Thus $\Tor^{A'}_1(A/m, B') = 0$; moreover, $x$ is
both $A$-regular and $B$-regular, so that $\Tor^A_1(A/m, B') = \Tor^A_1(A/m, B)$.
Therefore by Theorem~\hyperref[thm:22.3]{22.3}, $B$ is flat over $A$.
\end{proof}

We give a translation of the above theorem into algebraic geometry for the sake
of application. (The language is that of modern algebraic geometry, see for
example [Ha], Ch.~2.)

\begin{corollary*}
Let $k$ be a field, $X$ and $Y$ irreducible algebraic $k$-schemes, and
$f\colon Y \to X$ a morphism. Set $\dim X = n$, $\dim Y = m$, and suppose that the
following conditions hold: (1) $X$ is regular; (2) $Y$ is Cohen-Macaulay; (3) $f$
takes closed points of $Y$ into closed points of $X$ (this holds for example if
$f$ is proper); (4) for every closed point $x \in X$ the fibre $f^{-1}(x)$ is
$(m - n)$-dimensional (or empty). Then $f$ is flat.
\end{corollary*}

\begin{proof}
Let $y \in Y$ be a closed point, and set $x = f(y)$, $A = \mathcal{O}_{X,x}$, and
$B = \mathcal{O}_{Y,y}$. We have $\dim A = n$, $\dim B = m$, and since by
Theorem~\ref{thm:15.1} $\dim B/m_AB \le \dim B - \dim A = m - n$, we get
$\dim B/m_AB = m - n$ from (4). Therefore by the above theorem $B$ is flat over $A$,
and this is what was required to prove.
\end{proof}

\begin{theorem}\label{thm:23.2}
Let $\varphi\colon A \to B$ be a homomorphism of Noetherian rings, and let $E$ be an
$A$-module and $G$ a $B$-module. Suppose that $G$ is flat over $A$; then we have
the following:
\begin{enumerate}[(i)]
  \item if $p \in \Spec A$ and $G/pG \neq 0$ then $\Ass_A(G/pG) = \{p\}$; moreover, if
        $P \in \Ass_B(G/pG)$ then $P \cap A = p$;
  \item $\Ass_B(E \otimes_A G) = \bigcup_{p \in \Ass_A(E)} \Ass_B(G/pG)$.
\end{enumerate}
\end{theorem}

\begin{proof}
(i) $G/pG = G \otimes_A (A/p)$ is flat over $A/p$, and $A/p$ is an integral domain,
so that any non-zero element of $A/p$ is $G/pG$-regular (see Ex.~\ref{ex:7.5-ex}). In
other words, the elements of $A - p$ are $G/pG$-regular. This gives
$\Ass_A(G/pG) = \{p\}$. Also, if $P \in \Ass_B(G/pG)$ then there exists $t \in G/pG$
such that $\ann_B(t) = P$, and then $P \cap A = \ann_A(t) \in \Ass_A(G/pG) = \{p\}$.

(ii) If $p \in \Ass_A(E)$ then there is an exact sequence of the form
$0 \to A/p \to E$, and since $G$ is flat the sequence $0 \to G/pG \to E \otimes_A G$
is also exact; thus
\[
  \Ass_B(G/pG) \subset \Ass_B(E \otimes_A G).
\]
Conversely, if $P \in \Ass_B(E \otimes_A G)$ then there is an $e \in E \otimes_A G$ such
that $\ann_B(e) = P$. We write $e = \sum x_i \otimes y_i$ with $x_i \in E$ and $y_i \in G$,
and set $E' = \sum Ax_i$; then by flatness of $G$, we can view $E' \otimes_A G$ as a
submodule $E' \otimes_A G \subset E \otimes_A G$. Since $e \in E' \otimes_A G$ we have
$P \in \Ass_B(E' \otimes_A G)$. Now $E'$ is a finite $A$-module, so that we can
choose a shortest primary decomposition of $0$ in $E'$, say
$0 = Q_1 \cap \cdots \cap Q_r$. Since $E'$ can be embedded in $\bigoplus (E'/Q_i)$, if we set
$E_i = E'/Q_i$ then
\[
  \Ass_B(E' \otimes_A G) \subset \bigcup_i \Ass_B(E_i \otimes_A G),
\]
and therefore $P \in \Ass_B(E_i \otimes_A G)$ for some $i$. This $E_i$ is a finite
$A$-module having just one associated prime, say $p$. We have
$p \in \Ass_A(E') \subset \Ass_A(E)$. For large enough $\nu$ we get $p^\nu E_i = 0$, so
that $p^\nu(E_i \otimes_A G) = 0$, and thus $p \subset P \cap A$. Moreover, an
element of $A - p$ is $E_i$-regular, and hence also $E_i \otimes_A G$-regular, so that
finally $p = P \cap A$. Now choose a chain of submodules of $E_i$,
\[
  E_i = E_0 \supset E_1 \supset \cdots \supset E_r = 0
\]
such that $E_j/E_{j+1} \cong A/p_j$ with $p_j \in \Spec A$. Then also
\[
  E_i \otimes_A G \supset E_1 \otimes_A G \supset \cdots \supset E_r \otimes_A G = 0,
\]
with
\[
  (E_j \otimes_A G)/(E_{j+1} \otimes_A G) \cong (A/p_j) \otimes_A G = G/p_jG,
\]
so that $\Ass_B(E_i \otimes_A G) \subset \bigcup_j \Ass_B(G/p_jG)$. Therefore
$P \in \Ass_B(G/p_jG)$ for some $j$, but by (i), $P \cap A = p_j$, so that $p_j = p$
and $P \in \Ass_B(G/pG)$.
\end{proof}

\begin{theorem}\label{thm:23.3}
Let $(A, m, k)$ and $(B, n, k')$ be Noetherian local rings, and
$\varphi\colon A \to B$ a local homomorphism. Let $M$ be a finite $A$-module, $N$ a
finite $B$-module, and assume that $N$ is flat over $A$. Then
\[
  \depth_B(M \otimes_A N) = \depth_A M + \depth_B(N/mN).
\]
\end{theorem}

\begin{proof}
Let $x_1, \ldots, x_r \in m$ be a maximal $M$-sequence, and $y_1, \ldots, y_s \in n$ a
maximal $N/mN$-sequence. Writing $x_i'$ for the images of $x_i$ in $B$, let us
prove that $x_1', \ldots, x_r', y_1, \ldots, y_s$ is a maximal $M \otimes_A N$-sequence. Now
$x_1', \ldots, x_r'$ is an $M \otimes_A N$-sequence, and if we set
$M_1 = M/\sum x_iM$ then
\[
  m \in \Ass_A(M_1), \quad \text{and} \quad (M \otimes_A N) \Big/ \sum_{i=1}^r x_i(M \otimes_A N) = M_1 \otimes_A N.
\]
Moreover, by the corollary of Theorem~\hyperref[thm:22.5]{22.5}, $y_1$ is $N$-regular, and
$N_1 = N/y_1N$ is flat over $A$, so that from the exact sequence
$0 \to N \xrightarrow{y_1} N \to N_1 \to 0$ we get the exact sequence
$0 \to M_1 \otimes_A N \xrightarrow{y_1} M_1 \otimes_A N \to M_1 \otimes_A N_1 \to 0$.
Proceeding in the same way we see that $y_1, \ldots, y_s$ is an
$M_1 \otimes_A N$-sequence. After this we need only prove that the $B$-module
\[
  (M \otimes_A N) \Big/ \Big( \sum x_i(M \otimes_A N) + \sum y_j(M \otimes_A N) \Big) = M_1 \otimes_A N_s
\]
has depth $0$, that is $n \in \Ass_B(M_1 \otimes_A N_s)$; however,
$m \in \Ass_A(M_1)$ and $n \in \Ass_B(N_s/mN_s)$, so that this follows at once from
the previous theorem.
\end{proof}

\begin{corollary*}
Let $A \to B$ be a local homomorphism of Noetherian rings as in the theorem, and
set $F = B/mB$. Assume that $B$ is flat over $A$. Then
\begin{enumerate}[(i)]
  \item $\depth B = \depth A + \depth F$;
  \item $B$ is CM $\iff$ $A$ and $F$ are both CM.
\end{enumerate}
\end{corollary*}

\begin{proof}
(i) is the case $M = A$, $N = B$ of the theorem. From (i) and (*) we have
\[
  \dim B - \depth B = (\dim A - \depth A) + (\dim F - \depth F)
\]
and in view of $\dim A \ge \depth A$ and $\dim F \ge \depth F$, (ii) is clear.
\end{proof}

\begin{theorem}\label{thm:23.4}
Let $A \to B$ be a local homomorphism of Noetherian local rings, set $m = \rad(A)$
and $F = B/mB$. We assume that $B$ is flat over $A$; then
\[
  B \text{ is Gorenstein} \iff A \text{ and } F \text{ are both Gorenstein}.
\]
\end{theorem}

\begin{proof}
(K.~Watanabe [1]). By the corollary just proved, we can assume that $A$, $B$ and
$F$ are CM. Set $\dim A = r$ and $\dim F = s$, and let $\{x_1, \ldots, x_r\}$ be a
system of parameters of $A$, and $\{y_1, \ldots, y_s\}$ a subset of $B$ which reduces
to a system of parameters of $F$ modulo $mB$. Then as we have seen in the proof
of Theorem~\hyperref[thm:23.3]{3}, $(x_1, \ldots, x_r, y_1, \ldots, y_s)$ is a $B$-sequence, and
therefore a system of parameters of $B$, and $\bar{B} = B/(x, y)B$ is flat over
$\bar{A} = A/(x)A$. Thus replacing $A$ and $B$ by $\bar{A}$ and $\bar{B}$, we can reduce to
the case $\dim A = \dim B = 0$. Now in general, a zero-dimensional local ring
$(R, M)$ is Gorenstein if and only if $\Hom_R(R/M, R) = (0 : M)_R$ is isomorphic to
$R/M$. Now set
\[
  \rad(B) = n, \quad \rad(F) = n/mB = \bar{n} \quad \text{and} \quad (0 : m)_A = I.
\]
Then $I$ is of the form $I \cong (A/m)^t$ for some $t$, and
$(0 : mB)_B = IB \cong (A/m)^t \otimes_A B = F^t$. Furthermore, we have
$(0 : n)_B = (0 : \bar{n})_{IB} \cong ((0 : \bar{n})_F)^t$, and hence if we set
$(0 : \bar{n})_F \cong (F/\bar{n})^u = (B/n)^u$ then $(0 : n)_B \cong (B/n)^{tu}$.
Therefore
\[
  B \text{ is Gorenstein} \iff tu = 1 \iff t = u = 1 \iff A \text{ and } F \text{ are Gorenstein}.
\]
\end{proof}

\begin{theorem}\label{thm:23.5}
If $A$ is Gorenstein then so are $A[X]$ and $A[[X]]$.
\end{theorem}

\begin{proof}
We write $B$ for either of $A[X]$ or $A[[X]]$, so that $B$ is flat over $A$. For any
maximal ideal $M$ of $B$ we set $M \cap A = \mathfrak{p}$ and $A/\mathfrak{p}A = \kappa(\mathfrak{p})$. In case
$B = A[X]$, the local ring $B_M$ is a localisation of $B \otimes_A A_{\mathfrak{p}} =
A_{\mathfrak{p}}[X]$, and the fibre ring of $A_{\mathfrak{p}} \to B_M$ is a localisation of
$\kappa(\mathfrak{p})[X]$, hence regular. In case $B = A[[X]]$ then $X \in M$, and
$\mathfrak{p}$ a maximal ideal of $A$, so that $\kappa(\mathfrak{p}) = A/\mathfrak{p}$ and
\[
  B \otimes_A \kappa(\mathfrak{p}) = (A/\mathfrak{p})[[X]] = \kappa(\mathfrak{p})[[X]].
\]
This is a regular local ring, and is the fibre ring of $A_{\mathfrak{p}} \to B_M$. Thus in
either case $B_M$ is Gorenstein by the previous theorem.
\end{proof}

\begin{theorem}\label{thm:23.6}
Let $A$ be a Gorenstein ring containing a field $k$; then for any finitely
generated field extension $K$ of $k$, the ring $A \otimes_k K$ is Gorenstein.
\end{theorem}

\begin{proof}
We need only consider the case that $K$ is generated over $k$ by one element $x$.
If $x$ is transcendental over $k$ then $A \otimes_k K$ is isomorphic to a
localisation of $A \otimes_k k[X] = A[X]$, and since $A[X]$ is Gorenstein, so is
$A \otimes_k K$. If $x$ is algebraic over $k$ then since $K \cong k[X]/(f(X))$ with
$f(X) \in k[X]$ a manic polynomial, we have
\[
  A \otimes_k K = A[X]/(f(X));
\]
now $A[X]$ is Gorenstein and $f(X)$ is a non-zero-divisor of $A[X]$, so that we see
that $A \otimes_k K$ is also Gorenstein.
\end{proof}

\begin{remark}
Theorems 5 and 6 also hold on replacing Gorenstein by Cohen-Macaulay; the
proofs are exactly the same. For complete intersection rings the counterpart of
Theorem 4 also holds, so that the analogs of Theorems 5 and 6 follow; the proof
involves Andr\'e homology (Avramov [1]). As we see in the next theorem, a
slightly weaker form of the same result holds for regular rings.
\end{remark}

\begin{theorem}\label{thm:23.7}
Let $(A, m, k)$ and $(B, n, k')$ be Noetherian local rings, and $A \to B$ a local
homomorphism; set $F = B/mB$. We assume that $B$ is flat over $A$.
\begin{enumerate}[(i)]
  \item If $B$ is regular then so is $A$.
  \item If $A$ and $F$ are regular then so is $B$.
\end{enumerate}
\end{theorem}

\begin{proof}
(i) We have $\Tor^A_i(k, k) \otimes_A B = \Tor^B_i(B \otimes_A k, B \otimes_A k)$, and the right-hand
side is zero for $i > \dim B$. Since $B$ is faithfully flat over $A$, we have
$\Tor^A_i(k, k) = 0$ for $i \gg 0$, so that by \S19, Lemma~\ref{lem:19.1}, (i),
$\projdim_A k < \infty$ and since $\projdim k = \gldim A$, by
Theorem~\ref{thm:19.2}, $A$ is regular.

(ii) Set $r = \dim A$ and $s = \dim F$. Let $\{x_1, \ldots, x_r\}$ be a regular system of
parameters of $A$, and $\{y_1, \ldots, y_s\}$ a subset of $n$ which maps to a regular
system of parameters of $F$. Since $A \to B$ is injective, we can view $A$ as a
subring $A \subset B$. Then $\{x_1, \ldots, x_r, y_1, \ldots, y_s\}$ generates $n$ but
$\dim B = r + s$, so that $B$ is regular.
\end{proof}

\begin{remark}
In Theorem 7, even if $B$ is regular, $F$ need not be. For example, let $k$ be a
field, $x$ an indeterminate over $k$, and $B = k[x]_{(x)}$, $A = k[x^2]_{(x^2)}
\subset B$; then $F = B/x^2B = k[x]/(x^2)$ has a nilpotent element. By Theorem 1,
or directly, we see that $B$ is flat over $A$. (From a geometrical point of view,
this example corresponds to the projection of the plane curve $y = x^2$ onto the
$y$-axis, and, not surprisingly, the fibre over the origin is singular.)
\end{remark}

Consider the following conditions $(R_i)$ and $(S_i)$ for $i = 0, 1, 2, \ldots$ on a
Noetherian ring $A$:
\[
  (R_i) \quad A_P \text{ is regular for all } P \in \Spec A \text{ with } \htop P \le i;
\]
\[
  (S_i) \quad \depth A_P \ge \min(\htop P, i) \text{ for all } P \in \Spec A.
\]
$(S_0)$ always holds. $(S_1)$ says that all the associated primes of $A$ are minimal,
that is $A$ does not have embedded associated primes. $(R_0) + (S_1)$ is the
necessary and sufficient condition for $A$ to be reduced. $(S_i)$ for all $i \ge 0$
is just the definition of a CM ring.

For an integral domain $A$, $(S_1)$ is equivalent to the condition that every prime
divisor of a non-zero principal ideal has height $1$. The characterisation of
normal integral domain given in the corollary to Theorem~\ref{thm:11.5} can be
somewhat generalised as follows.

\begin{theorem}[Serre]\label{thm:23.8}
$(R_1) + (S_2)$ are necessary and sufficient conditions for a Noetherian ring $A$ to
be normal.
\end{theorem}

\begin{proof}
We defined a normal ring (see \S9), by the condition that the localisation at
every prime is an integrally closed domain. The conditions $(R_i)$ and $(S_i)$ are
also conditions on localisations, so that we can assume that $A$ is local.

\textit{Necessity}. This follows from Theorems~\ref{thm:11.2} and \ref{thm:11.5}.

\textit{Sufficiency}. Since $A$ satisfies $(R_0)$ and $(S_1)$ it is reduced, and the
shortest primary decomposition of $(0)$ is $(0) = P_1 \cap \cdots \cap P_r$, where $P_i$
are the minimal primes of $A$. Thus if we set $K$ for the total ring of fractions
of $A$, we have
\[
  K = K_1 \times \cdots \times K_r, \quad \text{with } K_i \text{ the field of fractions of } A/P_i.
\]
First of all we show that $A$ is integrally closed in $K$. Suppose that we have a
relation in $K$ of the form
\[
  (a/b)^n + c_1(a/b)^{n-1} + \cdots + c_n = 0,
\]
with $a, b, c_1, \ldots, c_n \in A$ and $b$ an $A$-regular element. This is equivalent to
a relation
\[
  a^n + \sum_{i=1}^n c_i a^{n-i}b^i = 0
\]
in $A$. Let $P \in \Spec A$ be such that $\htop P = 1$; then by $(R_1)$, $A_P$ is regular,
and therefore normal, so that $a_P \in b_P A_P$, where we write $a_P$, $b_P$ for the
images in $A_P$ of $a$, $b$. Now $b$ is $A$-regular, so that by $(S_1)$, all the prime
divisors of the principal ideal $bA$ have height $1$; thus if
$bA = \mathfrak{q}_1 \cap \cdots \cap \mathfrak{q}_t$ is a shortest primary decomposition and we set
$\mathfrak{p}_i$ for the prime divisor of $\mathfrak{q}_i$, then $a \in bA_{\mathfrak{p}_i} \cap A = \mathfrak{q}_i$ for all
$i$, and hence $a \in bA$, so that $a/b \in A$. Therefore $A$ is integrally closed in
$K$; in particular, the idempotents $e_i$ of $K$, which satisfy $e_i^2 - e_i = 0$, must
belong to $A$, so that from $1 = \sum e_i$ and $e_ie_j = 0$ for $i \neq j$ we get
\[
  A = Ae_1 \times \cdots \times Ae_r.
\]
Now since $A$ is supposed to be local, we must have $r = 1$, so that $A$ is an
integrally closed domain.
\end{proof}

\begin{theorem}\label{thm:23.9}
Let $(A, m)$ and $(B, n)$ be Noetherian local rings and $A \to B$ a local
homomorphism. Suppose that $B$ is flat over $A$, and that $i \ge 0$ is a given
integer. Then
\begin{enumerate}[(i)]
  \item if $B$ satisfies $(R_i)$, so does $A$;
  \item if both $A$ and the fibre ring $B \otimes_A \kappa(\mathfrak{p})$ over every prime ideal
        $\mathfrak{p}$ of $A$ satisfy $(R_i)$, so does $B$.
  \item The above two statements also hold with $(S_i)$ in place of $(R_i)$.
\end{enumerate}
\end{theorem}

\begin{proof}
(i) For $\mathfrak{p} \in \Spec A$, since $B$ is faithfully flat over $A$, there is a prime
ideal of $B$ lying over $\mathfrak{p}$; if we let $P$ be a minimal element among these then
$\htop(P/\mathfrak{p}B) = 0$, so that $\htop P = \htop \mathfrak{p}$. Hence
$\htop \mathfrak{p} \le i \Rightarrow B_P$ is regular, so that by
Theorem~\hyperref[thm:23.7]{7}, $A_{\mathfrak{p}}$ is regular. Also, by the corollary to
Theorem~\hyperref[thm:23.3]{3}, $\depth B_P = \depth A_{\mathfrak{p}}$, so that one sees easily that $(S_i)$ for
$B$ implies $(S_i)$ for $A$.

(ii) Let $P \in \Spec B$ and set $P \cap A = \mathfrak{p}$. If $\htop P \le i$ then we have
$\htop \mathfrak{p} \le i$ and $\htop(P/\mathfrak{p}B) \le i$, hence $A_{\mathfrak{p}}$ and
$B \otimes_A \kappa(\mathfrak{p})$ are both regular, so by Theorem~\hyperref[thm:23.7]{7}, (ii), $B_P$ is regular.
Hence $B$ satisfies $(R_i)$. Moreover, for $(S_i)$ we have
\begin{align*}
  \depth B_P &= \depth A_{\mathfrak{p}} + \depth B \otimes_A \kappa(\mathfrak{p}) \\
             &\ge \min(\htop \mathfrak{p}, i) + \min(\htop P/\mathfrak{p}B, i) \\
             &\ge \min(\htop \mathfrak{p} + \htop P/\mathfrak{p}B, i) = \min(\htop P, i).
\end{align*}
\end{proof}

\begin{corollary*}
Under the same assumptions as Theorem~\hyperref[thm:23.9]{9}, we have
\begin{enumerate}[(i)]
  \item if $B$ is normal (or reduced) then so is $A$;
  \item if both $A$ and the fibre rings of $A \to B$ are normal (or reduced) then
        so is $B$.
\end{enumerate}
\end{corollary*}

\begin{remark}
If $A$ and the closed fibre ring $F = B/mB$ only are normal, then $B$ does not
have to be; for instance, there are known examples of normal Noetherian rings
for which the completion is not normal.
\end{remark}

Finally, we would like to draw the reader's attention to the following obvious,
but useful, fact concerning the fibre ring. Let $\varphi'\colon A' \to B'$ be a ring
homomorphism and $I$ an ideal of $A'$; we set $A = A'/I$, $B = B'/IB'$, and write
$\varphi\colon A \to B$ for the map induced by $\varphi'$. If $\mathfrak{p}' \in \Spec A'$ is such that
$I \subset \mathfrak{p}'$, we set $\mathfrak{p} = \mathfrak{p}'/I$; then the fibre of $\varphi'$ over $\mathfrak{p}'$ coincides
with the fibre of $\varphi$ over $\mathfrak{p}$. To see this,
\[
  B' \otimes_{A'} \kappa(\mathfrak{p}') = B' \otimes_{A'} (A'/\mathfrak{p}')_{\mathfrak{p}'} = B \otimes_A (A/\mathfrak{p})_{\mathfrak{p}} = B \otimes_A \kappa(\mathfrak{p}).
\]
It follows from this that if all the fibre rings of $\varphi'$ have a good property,
the same is true of $\varphi$. For an example of this, see Ex.~\ref{ex:23.2-ex}.

% ============ 习题 ============
\section*{Exercises to \S\ref{sec:23}}
\emph{Prove the following propositions.}

\begin{exercise}\label{ex:23.1-ex}
If $A$ is a Gorenstein local ring then all the fibre rings of $A \to \hat{A}$ are again
Gorenstein; the same thing holds for Cohen-Macaulay.
\end{exercise}

\begin{exercise}\label{ex:23.2-ex}
If $A$ is a quotient of a CM local ring, and satisfies $(S_i)$, then the completion
$\hat{A}$ also satisfies $(S_i)$. In particular, if $A$ does not have embedded
associated primes then neither does $\hat{A}$.
\end{exercise}

\begin{exercise}\label{ex:23.3-ex}
Give another proof of Theorem~\hyperref[thm:23.4]{4} along the following lines:
\begin{enumerate}[(1)]
  \item Using $\Ext^i_A(A/m, A) \otimes_A \hat{B} = \Ext^i_B(F, B)$, show that $B$ Gorenstein
        implies $A$ Gorenstein.
  \item Assuming that $A$ is Gorenstein, prove that $F$ is Gorenstein if and only
        if $B$ is. Firstly reduce to the case $\dim A = 0$. Then prove that
        $\Ext^i_B(F, B) = 0$ for $i > 0$ and $\cong F$ for $i = 0$, and deduce that if
        $0 \to B \to I^{\bullet}$ is an injective resolution of $B$ as a $B$-module then
        $0 \to F \to \Hom_B(F, I^{\bullet})$ is an injective resolution of $F$ as an $F$-module,
        so that, writing $k$ for the residue field of $B$, we have
        $\Ext^i_B(k, B) = \Ext^i_F(k, F)$ for all $i$.
\end{enumerate}
\end{exercise}

% ============ §24 Generic freeness and open loci results (原书 pp.185-189) ============
\section{Generic freeness and open loci results}\label{sec:24}\index{generic freeness, section 24}\index{open loci, section 24}\index{Nagata criterion, section 24}\index{normal flatness, section 24}\index{flatness locus, section 24}

Let $A$ be a Noetherian integral domain, and $M$ a finite $A$-module. Then there
exists $0 \ne a \in A$ such that $M_a$ is a free $A_a$-module. This follows from
Theorem~\ref{thm:4.10}, or can be proved as follows: choose a filtration
\[
  M = M_0 \supset M_1 \supset \cdots \supset M_r = 0
\]
such that $M_{i-1}/M_i \cong A/p_i$, with $p_i \in \Spec A$; then if we take $a \ne 0$
contained in every non-zero $p_i$ we see that every $(M_{i-1}/M_i)_a$ is either zero
or isomorphic to $A_a$, so that $M_a$ is a free $A_a$-module.

For applications, we require a more general version of this, which does not
assume $M$ to be finite. We give below a theorem due to Hochster and Roberts
[1]. First we give the following lemma.

\begin{lemma*}
Let $B$ be a Noetherian ring, and $C$ a $B$-algebra generated over $B$ by a single
element $x$; let $E$ be a finite $C$-module, and $F \subset E$ a finite $B$-module
such that $CF = E$. Then $D = E/F$ has a filtration
\[
  0 = G_0 \subset G_1 \subset \cdots \subset G_i \subset G_{i+1} \subset \cdots \subset D
      \quad \text{with } D = \bigcup_{i=0}^{\infty} G_i
\]
such that the successive quotients $G_{i+1}/G_i$ are isomorphic to a finite number
of finite $B$-modules.
\end{lemma*}

\begin{proof}
Set
\[
  G_i' = F + xF + \cdots + x^iF \subset E, \quad G_i = G_i'/F,
\]
and
\[
  F_i = \{f \in F \mid x^{i+1}f \in G_i'\} \subset F.
\]
Then $0 \subset G_1 \subset \cdots \subset G_i \subset G_{i+1} \subset \cdots$ is a filtration of
$D$, and $G_{i+1}/G_i \cong F/F_i$; on the other hand, $F_1 \subset F_2 \subset \cdots
\subset F_i \subset \cdots$ is an increasing chain of $B$-submodules of $F$, so must
terminate.
\end{proof}

\begin{theorem}\label{thm:24.1}
Let $A$ be a Noetherian integral domain, $R$ a finitely generated $A$-algebra, and
$S$ a finitely generated $R$-algebra; we let $E$ be a finite $S$-module,
$M \subset E$ an $R$-submodule which is finite over $R$, and $N \subset R$ an
$A$-submodule which is finite over $A$, and set $D = E/(M + N)$. Then there
exists $0 \ne a \in A$ such that $D_a$ is a free $A_a$-module.
\end{theorem}

\begin{proof}
Write $A'$ for the image of $A$ in $R$, and suppose that
$R = A'[u_1, \ldots, u_h]$; similarly, write $R'$ for the image of $R$ in $S$, and
suppose that $S = R'[v_1, \ldots, v_k]$. We work by induction on $h + k$; if
$h = k = 0$ then $D$ is a finite $A$-module, and we have already dealt with this
case.

Write $R_j = A'[u_1, \ldots, u_j]$ for $0 \le j \le h$, and
$S_j = R'[v_1, \ldots, v_j]$ for $0 \le j \le k$.

Suppose first that $k > 0$; set $M + N = M' \subset E$. We have a filtration
\[
  S_0M' \subset S_1M' \subset \cdots \subset S_kM' = SM' \subset E,
\]
the successive quotients of which are $S_0M'$, $S_1M'/S_0M'$, $\ldots$,
$S_{k-1}M'/S_{k-2}M'$, $S_kM'/S_{k-1}M'$, $E/SM'$. We can apply the induction
hypothesis to each of these except the last two. By virtue of the lemma,
$S_kM'/S_{k-1}M'$ has a filtration with (up to isomorphism) just a finite number
of finite $S_{k-1}$-modules appearing as quotients, and so we can apply the
induction hypothesis again. For the final term, write $E' = E/SM'$, and let
$e_1, \ldots, e_n$ be a set of generators of $E'$ over $S$; write
$E_{k-1} = S_{k-1}e_1 + \cdots + S_{k-1}e_n$. Then $SE_{k-1} = E'$, so that the lemma
again gives a filtration of $E'$ with essentially finitely many finite
$S_{k-1}$-modules appearing as quotients, and we can apply the induction
hypothesis to this term also.

If $k = 0$ then $E$ is a finite $R$-module, and replacing $E$ by $E/M$ we can
assume that $M = 0$. The preceding proof then applies almost verbatim to this
case, with $R_j$ instead of $S_j$.
\end{proof}

\begin{theorem}[topological Nagata criterion]\label{thm:24.2}
Let $A$ be a Noetherian ring, and $U \subset \Spec A$ a subset. Then the following
two conditions are necessary and sufficient for $U \subset \Spec A$ to be open.
\begin{enumerate}[(1)]
  \item for $P, Q \in \Spec A$, $P \in U$ and $P \subset Q \Rightarrow Q \in U$;
  \item if $P \in U$ then $U$ contains a non-empty open subset of $V(P)$.
\end{enumerate}
\end{theorem}

\begin{proof}
Necessity is obvious, and we prove sufficiency. Let $V_1, \ldots, V_t$ be the
irreducible components of the closure of $U^c = \Spec A - U$, and let $p_i$ be
their generic points. If $p_i \in U$ then by (2) there is a proper closed subset
$W$ of $V_i$ such that $U^c \cap V_i \subset W$, and so
$U^c \subset W \cup (\bigcup_{j \neq i} V_j)$, which contradicts the definition of
$V_i$. Thus $p_i \notin U$, so that by (1), $V_i \subset U^c$ for all $i$ and therefore
$U^c$ is closed.
\end{proof}

\begin{theorem}\label{thm:24.3}
Let $A$ be a Noetherian ring, $B$ a finitely generated $A$-algebra, and $M$ a
finite $B$-module. Set $U = \{P \in \Spec B \mid M_P \text{ is flat over } A\}$;
then $U$ is open in $\Spec B$.
\end{theorem}

\begin{proof}
We verify the conditions (1) and (2) of Theorem~\hyperref[thm:24.2]{2}.

(1) If $P \subset Q$ are prime ideals of $B$ then for an $A$-module $N$ we have
$N \otimes_A M_Q = (N \otimes_A M_P) \otimes_{B_P} B_Q$, so that if $M_P$ is flat over $A$
then so is $M_Q$.

(2) Let $P \in U$ and $p = P \cap A$; set $\bar{A} = A/p$. Now if $Q \in V(P)$, we have
$pB_Q \subset \rad(B_Q)$, and hence by Theorem~\hyperref[thm:22.3]{22.3}, $M_Q$ is flat over $A$ if and
only if $M_Q \otimes_A \bar{A}$ is flat over $\bar{A}$ and $\Tor^A_1(M_Q, \bar{A}) = 0$.
Now $\Tor^A_1(M_P, \bar{A}) = 0$, and the left-hand side is equal to
$\Tor^A_1(M, \bar{A}) \otimes_B B_P$. By computing the Tor by means of a finite free
resolution of $\bar{A}$ over $A$, we see that $\Tor^A_1(M, \bar{A})$ is a finite $B$-module,
so that there is a neighbourhood $W$ of $P$ in $\Spec B$ such that
$\Tor^A_1(M_Q, \bar{A}) = 0$ for $Q \in W$. Moreover, by
Theorem~\hyperref[thm:24.1]{1}, there exists $a \in A - p$ such that $M_P/pM_P$ is a free
$\bar{A}_a$-module, so that if $Q \notin V(aB)$, then $M_Q \otimes_A \bar{A}$ is flat over
$\bar{A}$. Thus the open set $\{W \cap V(P)\} - V(aB)$ of $V(P)$ is contained in
$U$.
\end{proof}

\begin{remark}
If $A$ is Noetherian and $B$ is a finitely generated $A$-algebra which is flat over
$A$ then it is also known that the map $\Spec B \to \Spec A$ is open; see [M],
p.~48 or [G2], (2.4.6).
\end{remark}

Let $A$ be a ring, and $\mathcal{P}$ a property of local rings; we define a subset
$\mathcal{P}(A) \subset \Spec A$ by $\mathcal{P}(A) = \{p \in \Spec A \mid \mathcal{P} \text{ holds for } A_p\}$.
For example, if $\mathcal{P}$ = regular, complete intersection, Gorenstein or CM we
write $\Reg(A)$, $\mathrm{c.i.}(A)$, $\Gor(A)$ or $\CM(A)$ for these loci. The question as to
whether $\mathcal{P}(A)$ is open is an interesting and important question. For $\Reg(A)$
this is a classical question, but for the other properties the systematic study
was initiated by Grothendieck.

The following proposition is called the (ring-theoretic) \emph{Nagata criterion}
for the property $\mathcal{P}$, and we abbreviate this to (NC).

\textbf{(NC)}: Let $A$ be a Noetherian ring. If $\mathcal{P}(A/p)$ contains a non-empty open
subset of $\Spec(A/p)$ for every $p \in \Spec A$, then $\mathcal{P}(A)$ is open in
$\Spec A$.

The truth or otherwise of this proposition depends on $\mathcal{P}$; in the
remainder of this section we discuss some $\mathcal{P}$ for which (NC) holds. In
Theorem~\hyperref[thm:24.4]{24.4} and Ex.~\ref{ex:24.3-ex} we illustrate how (NC) can be applied to
prove openness results.

\begin{theorem}[Nagata]\label{thm:24.4}
(NC) holds for $\mathcal{P}$ = regular.
\end{theorem}

\begin{proof}
Let $U = \Reg(A)$. A localisation of a regular local ring is again regular, so
that $U$ satisfies condition (1) of Theorem~\hyperref[thm:24.2]{2}. We now check condition (2). If
$P \in U$ then $A_P$ is regular, so that we can take $x_1, \ldots, x_n \in P$ to form a
regular system of parameters of $A_P$ (where $n = \htop P$). Then there exists a
neighbourhood $W$ of $P$ in $\Spec A$ such that
\[
  PA_Q = (x_1, \ldots, x_n)A_Q
\]
for all $Q \in W$. (In fact, if $a \in A$ is an element not contained in $P$, but
contained in every other prime divisor of $(x_1, \ldots, x_n)$ then
$PA_Q = (x_1, \ldots, x_n)A_Q$.) Moreover, by the hypothesis in (NC) there exists a
neighbourhood $W'$ of $P$ in $V(P)$ such that $A_Q/PA_Q$ is regular for $Q \in W'$.
Then $A_Q$ is regular for $Q \in W' \cap W$, so that $W' \cap W \subset U$.
\end{proof}

\begin{theorem}\label{thm:24.5}
(NC) also holds for $\mathcal{P}$ = CM.
\end{theorem}

\begin{proof}
As with the previous proof, we reduce to checking condition (2) of
Theorem~\hyperref[thm:24.2]{2}. Let $P \in \CM(A)$. If we take $u \in A - P$ and replace $A$ by $A_u$,
then we are considering a neighbourhood of $P$ in $\Spec A$, so that we will refer
to this procedure as `passing to a smaller neighbourhood of $P$'. Since $A_P$ is
CM, if $\htop P = n$ we can choose an $A_P$-sequence $y_1, \ldots, y_n \in P$. One sees
easily that after passing to a smaller neighbourhood of $P$, we can assume that
\begin{enumerate}[(a)]
  \item $y_1, \ldots, y_n$ is an $A$-sequence; and
  \item $I = (y_1, \ldots, y_n)A$ is a $P$-primary ideal.
\end{enumerate}
Then for $Q \in V(P)$, it is equivalent to say that $A_Q$ is CM or that $A_Q/IA_Q$
is CM. Thus replacing $A$ by $A/I$ we can assume that $0$ is a $P$-primary ideal.
Then $P^r = 0$ for some $r > 0$. Now consider the filtration
$0 \subset P^{r-1} \subset \cdots \subset P \subset A$ of $A$. Each $P^i/P^{i+1}$ is a finite
$A/P$-module, but $A/P$ is an integral domain, so that passing to a smaller
neighbourhood of $P$ we can assume that $P^i/P^{i+1}$ is a free $A/P$-module for
$0 \le i < r$. It is then easy to see that if $x_1, \ldots, x_s \in A$ is an
$A/P$-sequence, it is also an $A$-sequence. However, according to the hypothesis
in (NC), passing to a smaller neighbourhood of $P$, we can assume that $A/P$ is a
CM ring. Then for $Q \in V(P)$ the ring $A_Q/PA_Q$ is CM, so that from what we have
said above,
\[
  \depth A_Q \ge \depth A_Q/PA_Q = \dim A_Q/PA_Q = \dim A_Q,
\]
and $A_Q$ is CM.
\end{proof}

Let $A$ be a Noetherian ring and $I$ an ideal of $A$; we set $B = A/I$ and write
$Y$ for the closed subset $V(I) \subset \Spec A$. Let $M$ be a finite $A$-module. We
say that $M$ is \emph{normally flat along $Y$}\index{normally flat} if the $B$-module
$\gr_I(M) = \bigoplus_{n \ge 0} I^nM/I^{n+1}M$ is flat over $B$. If $B$ is a local
ring, this is the same as saying that each $I^nM/I^{n+1}M$ is a free $B$-module.
Normal flatness is an important notion introduced by Hironaka, and it plays a
leading role in the problem of resolution of singularities; we have used it in
the above proof in the statement that if $P$ is nilpotent and $A$ is normally
flat along $V(P)$ then an $A/P$-sequence is an $A$-sequence. However, in this
book we do not have space to discuss the theory of normal flatness any further,
and we refer to Hironaka [1] and [G2], (6.10).

\begin{theorem}\label{thm:24.6}
(NC) holds for $\mathcal{P}$ = Gorenstein.
\end{theorem}

\begin{proof}
Once more we reduce to verifying condition (2) of Theorem~\hyperref[thm:24.2]{2}. Suppose that
$P \in \Gor(A)$; if $\htop P = n$ then since $A_P$ is CM, we can take
$x_1, \ldots, x_n \in P$ forming an $A_P$-sequence. Passing to a smaller
neighbourhood of $P$, we can assume that $x_1, \ldots, x_n$ is an $A$-sequence.
Moreover, replacing $A$ by $A/(x_1, \ldots, x_n)$ we can assume that $\htop P = 0$. In
addition, we assume that $P$ is the unique minimal prime ideal of $A$. Since $A_P$
is a zero-dimensional Gorenstein ring, we have
\[
  \Ext^i_A(A/P, A) \otimes_A A_P = \Ext^i_{A_P}(\kappa(P), A_P) = 0
\]
and
\[
  \Hom_A(A/P, A) \otimes_A A_P = \Hom_{A_P}(\kappa(P), A_P) = \kappa(P).
\]
Thus passing to a smaller neighbourhood, we can assume that
$\Ext^i_A(A/P, A) = 0$ and $\Hom_A(A/P, A) \cong A/P$. In addition, as in the proof of
the previous theorem, we can assume that $P^i/P^{i+1}$ is a free $A/P$-module for
$0 \le i < r$, where $P^r = 0$. Then using
\[
  0 \to P^i/P^{i+1} \to P/P^{i+1} \to P/P^i \to 0,
\]
we get by induction that $\Ext^i_A(P, A) = 0$; from this it follows that
$\Ext^i_A(A/P, A) = 0$, and in turn by induction that $\Ext^i_A(P, A) = 0$, so that
$\Ext^i_A(A/P, A) = 0$. Proceeding in the same way we see that
$\Ext^i_A(A/P, A) = 0$ for every $i > 0$. If we take an injective resolution
$0 \to A \to I^{\bullet}$ of $A$ as an $A$-module, and consider the complex obtained by
applying $\Hom_A(A/P, -)$ to it, from what we have just said we obtain an exact
sequence $0 \to A/P \to \Hom_A(A/P, I^{\bullet})$, and this is an injective resolution of
$A/P$ as an $A/P$-module. The same thing holds on replacing $A$ by $A_Q$ for
$Q \in V(P)$, and then setting $k = \kappa(Q)$, we get
$\Ext^i_{A_Q/PA_Q}(k, A_Q/PA_Q) = \Ext^i_{A_Q}(k, A_Q)$. Thus it is equivalent to say that
$A_Q$ is Gorenstein or that $A_Q/PA_Q$ is Gorenstein. Therefore from the hypothesis
in (NC) we have that $\Gor(A) \cap V(P)$ contains a neighbourhood of $P$ in $V(P)$.
\end{proof}

The above proof is due to Greco and Marinari [1]. Their paper also proves that
(NC) also holds for $\mathcal{P}$ = complete intersection.

% ============ 习题 ============
\section*{Exercises to \S\ref{sec:24}}
\emph{Prove the following propositions.}

\begin{exercise}\label{ex:24.1-ex}
Let $A$ be a Noetherian ring, and $I$ an ideal of $A$; assume that $I^r = 0$, and
that $I^i/I^{i+1}$ is a free $A/I$-module for $1 \le i \le r$. Then for
$x_1, \ldots, x_r \in A$, it is equivalent for $(x_1, \ldots, x_r)$ to be an
$A$-sequence or an $A/I$-sequence.
\end{exercise}

\begin{exercise}\label{ex:24.2-ex}
If $A$ is a quotient of a CM ring $R$ then $\CM(A)$ is open in $\Spec A$.
\end{exercise}

\begin{exercise}\label{ex:24.3-ex}
If $A$ is a quotient of a Gorenstein ring then $\Gor(A)$ is open in $\Spec A$.
\end{exercise}


% ============ Chapter 9. Derivations (原书 pp.190-212) ============
\chapter{Derivations}\label{ch:9}

This chapter can be read independently of the preceding ones; the main
themes are derivations of rings and modules of differentials. The results
of this chapter will be applied in the proof of the structure theorem for
complete local rings in the next chapter, but in addition derivations and
modules of differentials have an important influence on properties of rings,
for example via the connection with regularity.

In \S\ref{sec:25} we discuss the general theory of modules of differentials, and
also prove the Hochschild formula for derivations of rings in characteristic
$p$. \S\ref{sec:26} is pure field theory; Theorem~\ref{thm:26.8}, which states that a
$p$-basis of a separable extension is algebraically independent, is taken from
Matsumura [3]. The terminology 0-\'etale is due to Andr\'e, and corresponds
to `formally \'etale for the discrete topology' in EGA. In \S\ref{sec:27} we treat the
higher derivations of Hasse and F.~K.~Schmidt, concentrating on the
extension problem which they did not treat, in a version due to author.

% ============ §25 Derivations and differentials (原书 pp.190-198) ============
\section{Derivations and differentials}\label{sec:25}\index{derivation, section 25}\index{Kähler differentials, module of, section 25}\index{0-smooth, section 25}\index{0-unramified, section 25}\index{0-étale, section 25}\index{Hochschild formula, section 25}

Let $A$ be a ring and $M$ an $A$-module. A \emph{derivation} from $A$ to $M$ is a map
$D\colon A \to M$ satisfying $D(a + b) = Da + Db$ and $D(ab) = bDa + aDb$; the set of all
these is written $\Der(A, M)$. It becomes an $A$-module in a natural way, with
$D + D'$ and $aD$ defined by $(D + D')a = Da + D'a$ and $(aD)b = a(Db)$.

If $A$ is a $k$-algebra via a ring homomorphism $f\colon k \to A$, we say that $D$ is
a \emph{$k$-derivation}, or a derivation \emph{over $k$}, if $D \circ f = 0$; the set of all
$k$-derivations of $A$ into $M$ is written $\Der_k(A, M)$. It is an $A$-submodule of
$\Der(A, M)$. Since $1 \cdot 1 = 1$, for any $D \in \Der(A, M)$ we have
$D(1) = D(1) + D(1)$, so that $D(1) = 0$, and so viewing $A$ as $\mathbb{Z}$-algebra we have
$\Der(A, M) = \Der_{\mathbb{Z}}(A, M)$.

In the particular case $M = A$, we write $\Der_k(A)$ for $\Der_k(A, A)$. If $D$,
$D' \in \Der_k(A)$, we can compose $D$ and $D'$ as maps $A \to A$, and it is easy to see
that the bracket $[D, D'] = DD' - D'D$ is again an element of $\Der_k(A)$, and that
$\Der_k(A)$ becomes a Lie algebra with this bracket.

Quite generally, for $D \in \Der(A, M)$ and $a \in A$ one sees at once that
$D(a^n) = na^{n-1}Da$. Hence if $A$ is a ring of characteristic $p$ we have
$D(a^p) = 0$. Also, in general we have a Leibnitz formula for powers of $D$,
\[
  D^n(ab) = \sum_{i=0}^n \binom{n}{i} D^i a \cdot D^{n-i} b;
\]
if $A$ has characteristic $p$ then this reduces to $D^p(ab) = D^pa \cdot b + a \cdot
D^pb$, so that also $D^p \in \Der(A)$.

Let $k$ be a ring, $B$ a $k$-algebra, and $N$ an ideal of $B$ with $N^2 = 0$; set
$A = B/N$. The $B$-module $N$ can in fact be viewed as an $A$-module. In this
situation, we say that $B$ is an \emph{extension} of the $k$-algebra $A$ by the $A$-module
$N$; (note that $B$ does not contain $A$, so that this is a different usage of
extension). We write this extension as usual in the form of an exact sequence
\[
  0 \to N \to B \to A \to 0.
\]
We say that this extension is \emph{split}, or is the \emph{trivial extension}, if there exists
a $k$-algebra homomorphism $\sigma\colon A \to B$ such that $f\sigma = 1$ (the identity map
of $A$). Then we can identify $A$ and $\sigma(A)$, and we have $B = A \oplus N$ as a
$k$-module. Conversely, starting from any $k$-algebra $A$ and an $A$-module $N$, we
can make the direct sum $A \oplus N$ of $k$-modules into a trivial extension of $A$
by $N$, by defining the product
\[
  (a, x)(a', x') = (aa', ax' + a'x) \quad \text{for } a, a' \in A \text{ and } x, x' \in N.
\]
In this book, we will write $A \ltimes N$ for this algebra.

In general, given a commutative diagram in the category of $k$-algebras
\[
\begin{tikzcd}
B \arrow[r, "f"] & A \\
& C \arrow[ul, "g"'] \arrow[u, "h"']
\end{tikzcd}
\]
where we think of $f$ as being fixed, we say that $h$ is a \emph{lifting} of $g$ to $B$.
Write $N$ for the ideal $\Ker f$ of $B$. If $h'\colon C \to B$ is another lifting of $g$, then
$h - h'$ is a map from $C$ to $N$. If $N^2 = 0$ then $N$ is an $f(\bar{B})$-module, and
moreover, by means of $g\colon C \to f(B) \subset A$, we can consider $N$ as a $C$-module.
Then it is easy to see that $h - h'\colon C \to N$ is a $k$-derivation of $C$ to the
$C$-module $N$. Conversely, if $D \in \Der_k(C, N)$ then $h + D$ is another lifting of $g$
to $B$.

Let $k$ be a ring and $A$ a $k$-algebra, and write $\mathcal{M}_A$ for the category of
$A$-modules. We have a covariant functor $M \mapsto \Der_k(A, M)$ from $\mathcal{M}_A$ to
itself, which turns out to be a representable functor. In other words, there
exists an $A$-module $\Omega$ and a derivation $d \in \Der_k(A, \Omega)$ with the following
universal property: for any $A$-module $M$ and any $D \in \Der_k(A, M)$, there exists
a unique $A$-linear map $f\colon \Omega \to M$ such that $D = f \circ d$. We are now going
to prove this. Firstly, define $\mu\colon A \otimes_k A \to A$ by
\[
  \mu(x \otimes y) = xy;
\]
then $\mu$ is a homomorphism of $k$-algebras. Set
\[
  I = \Ker \mu, \quad \Omega_{A/k} = I/I^2 \quad \text{and} \quad B = (A \otimes_k A)/I^2;
\]
then $\mu$ induces $\bar{\mu}\colon B \to A$, and
\[
  0 \to \Omega_{A/k} \to B \xrightarrow{\bar{\mu}} A \to 0
\]
is an extension of the $k$-algebra $A$ by $\Omega_{A/k}$; this extension splits, and in
fact defining $\iota_i\colon A \to B$ for $i = 1, 2$ by
\[
  \iota_1(a) = a \otimes 1 \mod I^2 \quad \text{and} \quad \iota_2(a) = 1 \otimes a \mod I^2,
\]
we get two liftings of $1_A\colon A \to A$. Hence $d = \iota_1 - \iota_2$ is a derivation of
$A$ to $\Omega_{A/k}$. Now we prove that the pair $(\Omega_{A/k}, d)$ satisfies the
conditions for the above $(\Omega, d)$. If $D \in \Der_k(A, M)$ and we define
$\varphi\colon A \otimes_k A \to A \ltimes M$ by $\varphi(x \otimes y) = (xy, xDy)$ then
$\varphi$ is a homomorphism of $k$-algebras, and
\[
  \varphi\Big(\sum x_i \otimes y_i\Big) = 0 \Rightarrow \sum x_iy_i = 0
      \Rightarrow \varphi\Big(\sum x_i \otimes y_i\Big) = \Big(0, \sum x_iDy_i\Big),
\]
hence $\varphi$ maps $I$ into $M$. Now $M^2 = 0$, so that we finally get
$f\colon I/I^2 = \Omega_{A/k} \to M$. For $a \in A$ we have
\begin{align*}
  f(da) &= f(1 \otimes a - a \otimes 1 \mod I^2) = \varphi(1 \otimes a) - \varphi(a \otimes 1) \\
        &= Da - a \cdot D(1) = Da,
\end{align*}
so that $D = f \circ d$. Moreover, $\Omega_{A/k}$ has the $A$-module structure induced by
multiplication by $a \otimes 1$ in $A \otimes_k A$ (or multiplication by $1 \otimes a$; since
$a \otimes 1 - 1 \otimes a \in I$, they both come to the same thing); thus if
$\omega = \sum x_i \otimes y_i \mod I^2 \in \Omega_{A/k}$ then
$a\omega = \sum ax_i \otimes y_i \mod I^2$, and $f(a\omega) = \sum ax_iDy_i =
a f(\omega)$, so that $f$ is $A$-linear. We have
\[
  a \otimes a' = (a \otimes 1)(1 \otimes a' - a' \otimes 1) + aa' \otimes 1,
\]
so that if $\omega = \sum x_i \otimes y_i \in I$ then $\omega \mod I^2 = \sum x_i\,dy_i$.
Hence $\Omega_{A/k}$ is generated as an $A$-module by $\{da \mid a \in A\}$, so that the
uniqueness of a linear map $f\colon \Omega_{A/k} \to M$ satisfying $D = f \circ d$ is
obvious.

The $A$-module $\Omega_{A/k}$, which we have just obtained, is called the \emph{module of
differentials}\index{module of differentials} of $A$ over $k$, or the \emph{module of K\"ahler differentials}, and for
$a \in A$ the element $da \in \Omega_{A/k}$ is called the \emph{differential} of $a$. We can write
$d_{A/k}$ for $d$ to be more specific. From the definition, we see that
\[
  \Der_k(A, M) \cong \Hom_A(\Omega_{A/k}, M).
\]

\begin{example*}
If $A$ is generated as a $k$-algebra by a subset $U \subset A$ then $\Omega_{A/k}$ is
generated as an $A$-module by $\{da \mid a \in U\}$. Indeed, if $a \in A$ then there
exist $a_i \in U$ and a polynomial $f(X) \in k[X_1, \ldots, X_n]$ such that
$a = f(a_1, \ldots, a_n)$, and then from the definition of derivation we have
\[
  da = \sum_i f_i(a_1, \ldots, a_n)da_i, \quad \text{where } f_i = \partial f/\partial X_i.
\]
In particular, if $A = k[X_1, \ldots, X_n]$ then $\Omega_{A/k} = AdX_1 + \cdots + AdX_n$, and
$dX_1, \ldots, dX_n$ are linearly independent over $A$; this follows at once from the
fact that there are $D_i \in \Der_k(A)$ such that $D_iX_j = \delta_{ij}$.
\end{example*}

We say that a $k$-algebra $A$ is \emph{0-smooth}\index{0-smooth} (over $k$) if it has the following
property: for any $k$-algebra $C$, any ideal $N$ of $C$ satisfying $N^2 = 0$, and any
$k$-algebra homomorphism $\bar{u}\colon A \to C/N$, there exists a lifting $u\colon A \to C$
of $\bar{u}$ to $C$, as a $k$-algebra homomorphism. In terms of diagrams, given a
commutative diagram
\[
\begin{tikzcd}
A \arrow[r, "\bar{u}"] & C/N \\
k \arrow[u] \arrow[r] & C \arrow[u]
\end{tikzcd}
\]
there exists $v$ such that
\[
\begin{tikzcd}
A \arrow[r, "\bar{u}"] \arrow[rd, "v"'] & C/N \\
k \arrow[u] \arrow[r] & C \arrow[u]
\end{tikzcd}
\]
is commutative. Moreover, we say that $A$ is \emph{0-unramified}\index{0-unramified} over $k$ (or \emph{0-neat})
if there exists at most one such $v$. When $A$ is both 0-smooth and
0-unramified, that is when for given $\bar{u}$ there exists a unique $v$, we say that
$A$ is \emph{0-\'etale}\index{0-étale}. The condition for $A$ to be 0-unramified over $k$ is that
$\Omega_{A/k} = 0$: sufficiency is obvious, and if we recall that in the construction of
$\Omega_{A/k}$ we had $d = \iota_1 - \iota_2$, necessity is clear.

If $A$ is a ring and $S \subset A$ is a multiplicative set then the localisation $A_S$
is 0-\'etale over $A$. This follows from the fact (Ex.~\ref{ex:1.1-ex}) that if $x \in C$ is a unit
modulo a nilpotent ideal, then it is itself a unit. We leave the details to the
reader.

\begin{theorem}[First fundamental exact sequence]\label{thm:25.1}
A composite $k \xrightarrow{f} A \xrightarrow{g} B$ of ring homomorphisms leads to an
exact sequence of $B$-modules
\[
  (1) \quad \Omega_{A/k} \otimes_A B \xrightarrow{\alpha} \Omega_{B/k} \to \Omega_{B/A} \to 0,
\]
where the maps are given by $\alpha(d_{A/k}a \otimes b) = b\,d_{B/k}g(a)$ and
$\beta(d_{B/k}b) = d_{B/A}b$ for $a \in A$ and $b \in B$. If moreover $B$ is 0-smooth over
$A$ then the sequence
\[
  (2) \quad 0 \to \Omega_{A/k} \otimes_A B \to \Omega_{B/k} \to \Omega_{B/A} \to 0,
\]
obtained from (1) by adding $0 \to$ at the left, is a split exact sequence.
\end{theorem}

\begin{proof}
In order for a sequence $N' \xrightarrow{\alpha} N \xrightarrow{\beta} N''$ of $B$-modules to be
exact, it is sufficient that for every $B$-module $T$, the induced sequence
\[
  \Hom_B(N', T) \xleftarrow{\alpha^*} \Hom_B(N, T) \xleftarrow{\beta^*} \Hom_B(N'', T)
\]
is exact. Indeed, taking $T = N''$, we get $\alpha^*\beta^*(1_{N''}) = 0$, and therefore
$\beta\alpha = 0$; and taking $T = N/\Im \alpha$, we see easily that $\Ker \beta = \Im \alpha$. From
this, to prove that (1) above is exact, it is enough to show that for any $B$-module
$T$,
\[
  (3) \quad \Der_k(A, T) \xleftarrow{} \Der_k(B, T) \xleftarrow{} \Der_B(B, T) \to 0
\]
is exact, but this is obvious.

Now suppose that $B$ is 0-smooth over $A$. Choose $D \in \Der_k(A, T)$ and consider
the commutative diagram
\[
\begin{tikzcd}
A \arrow[r, "\sigma"] \arrow[d, "g"'] & B \ltimes T \arrow[d, "p"] \\
B \arrow[r, "1_B"'] & B
\end{tikzcd}
\quad \text{with } \sigma(a) = (ga, Da).
\]
Then by assumption, there exists $\kappa\colon B \to B \ltimes T$ which can be added to
the diagram, leaving it commutative. If we write $\kappa(b) = (b, D'b)$ then
$D'\colon B \to T$ is a derivation of $B$ such that $D = D' \circ g$, and $D'$ corresponds to a
$B$-linear map $\alpha'\colon \Omega_{B/k} \to T$. Now take $T$ to be $\Omega_{A/k} \otimes_A B$,
and define $D$ by $D(a) = d_{A/k}(a) \otimes 1$, so that $D = D' \circ g$ implies that
$\alpha'\alpha = 1_T$. Thus (2) is split.
\end{proof}

Now consider the case $k \to A \xrightarrow{g} B$ when $g$ is surjective; set
$\Ker g = m$, so $B = A/m$. Then in (1) of the previous theorem we of course have
$\Omega_{B/A} = 0$, and we want to determine $\Ker \alpha$.

\begin{theorem}[Second fundamental exact sequence]\label{thm:25.2}
In the above notation, we have an exact sequence
\[
  (4) \quad m/m^2 \xrightarrow{\delta} \Omega_{A/k} \otimes_A B \xrightarrow{\alpha} \Omega_{B/k} \to 0
\]
where $\delta$ is the $B$-linear map defined by $\delta(x \mod m^2) = d_{A/k}x \otimes 1$. If $B$ is
0-smooth over $k$ then
\[
  (5) \quad 0 \to m/m^2 \to \Omega_{A/k} \otimes_A B \to \Omega_{B/k} \to 0
\]
is a split exact sequence.
\end{theorem}

\begin{proof}
We once more take an arbitrary $B$-module $T$ and consider
\[
  (6) \quad \Hom_B(m/m^2, T) \xleftarrow{\delta^*} \Der_k(A, T) \xleftarrow{\alpha^*} \Der_k(B, T).
\]
For $D \in \Der_k(A, T)$, to say that $\delta^*(D) = 0$ is just to say that $D(m) = 0$, so
that $D$ can be considered as a derivation from $B = A/m$; hence (6) is exact.

If $B$ is 0-smooth over $k$ then the extension
\[
  0 \to m/m^2 \to A/m^2 \to B \to 0
\]
of the $k$-algebra $B$ by $m/m^2$ splits, that is there exists a homomorphism of
$k$-algebras $s\colon B \to A/m^2$ such that $\bar{g}s = 1_B$. Now $sg\colon A/m^2 \to A/m^2$ is a
homomorphism vanishing on $m/m^2$, and $g(1 - sg) = 0$, so that if we set
$D = 1 - sg$ then $D\colon A/m^2 \to m/m^2$ is a derivation. If $\psi \in \Hom_B(m/m^2, T)$
then the composite $D'$ of
\[
  A \to A/m^2 \xrightarrow{D} m/m^2 \xrightarrow{\psi} T
\]
is an element of $\Der_k(A, T)$ satisfying $\delta^*(D') = \psi$. Indeed, for $x \in m$, if we
let $\bar{x} = x \mod m^2$ then
\[
  D'(x) = \psi(D(\bar{x})) = \psi(\bar{x} - sg(\bar{x})) = \psi(\bar{x}).
\]
Therefore $\delta^*$ is surjective. If we set $T = m/m^2$ then we see that (5) is a split
exact sequence.
\end{proof}

\begin{example*}
Suppose that $B = k[X_1, \ldots, X_n]/(f_1, \ldots, f_m) = k[x_1, \ldots, x_n]$; then setting
$A = k[X_1, \ldots, X_n]$ and using the above theorem, we have
\[
  \Omega_{B/k} = (\Omega_{A/k} \otimes_A B)/R = F/R
\]
where $F$ is the free $B$-module with basis $dX_1, \ldots, dX_n$, and $R$ is the
submodule of $F$ generated by $df_i = \sum_j (\partial f_i/\partial X_j)dX_j$ for $1 \le i \le m$.
For example, if $k$ is a field of characteristic $\neq 2$ and
\[
  B = k[X, Y]/(X^2 + Y^2) = k[x, y],
\]
then $\Omega_{B/k} = Bdx + Bdy$, where the only relation between $dx$ and $dy$ is
$xdx + ydy = 0$. If $k$ has characteristic $2$ then $\Omega_{B/k}$ is the free $B$-module of
rank $2$ with basis $dx$, $dy$.
\end{example*}

\begin{theorem}\label{thm:25.3}
Suppose that a field $L$ is a separable algebraic extension of a subfield $K$; then
$L$ is 0-\'etale over $K$. Moreover, for any subfield $k \subset K$ we have
$\Omega_{L/k} = \Omega_{K/k} \otimes_K L$.
\end{theorem}

\begin{proof}
Suppose that $0 \to N \to C \to C/N \to 0$ is an extension of $K$-algebras with
$N^2 = 0$, and that $\bar{u}\colon L \to C/N$ is a given $K$-algebra homomorphism. If $L'$ is an
intermediate field $K \subset L' \subset L$ with $L'$ finite over $K$, then, as is well-known in
field theory, we can write $L' = K(\alpha)$; let $f(X)$ be the minimal polynomial of
$\alpha$ over $K$ so that $L' \cong K[X]/(f)$, and $f'(\alpha) \neq 0$. Thus to lift
$\bar{u}|_{L'}\colon L' \to C/N$ to $C$, we need only find an element $y \in C$ satisfying
$f(y) = 0$ and $y \mod N = \bar{u}(\alpha)$. Now choose some inverse image $Y \in C$ of
$\bar{u}(\alpha)$; then $f(Y) \mod N = \bar{u}(f(\alpha)) = 0$, so that $f(Y) \in N$. Moreover,
since $N^2 = 0$, so that for $\eta \in N$ we get
\[
  f(Y + \eta) = f(Y) + f'(Y)\eta;
\]
but $f'(\alpha)$ is a unit of $L$, so that $\bar{u}(f'(\alpha)) = f'(Y) \mod N$ is a unit of
$C/N$, and hence $f'(Y)$ is a unit of $C$ by Ex.~\ref{ex:1.1-ex}. Thus if we set
$\eta = -f(Y)/f'(Y)$ we have $\eta \in N$ and $f(Y + \eta) = 0$. The $K$-algebra
homomorphism $v\colon L' \to C$ obtained by taking $\alpha$ to $v(\alpha) = y + \eta$ is a
lifting of $\bar{u}|_{L'}$, and one can see by the construction that $v$ is unique. Thus for
every $\alpha \in L$ there is a uniquely determined lifting $v_\alpha\colon K(\alpha) \to C$ of
$\bar{u}|_{K(\alpha)}$, and we can define $v\colon L \to C$ by $v(\alpha) = v_\alpha(\alpha)$. In fact, for
$\alpha, \beta \in L$ there exists $\gamma \in L$ such that $K(\gamma)$ contains both $\alpha$ and
$\beta$, and then by uniqueness we have
\[
  v_\gamma|_{K(\alpha)} = v_\alpha \quad \text{and} \quad v_\gamma|_{K(\beta)} = v_\beta.
\]
The second half comes from $\Omega_{L/K} = 0$ and Theorem~\hyperref[thm:25.1]{1}.
\end{proof}

We turn now to derivations. As we have seen, if $A$ is a ring of characteristic $p$
then for $D \in \Der(A)$ we have $D^p \in \Der(A)$. What can we say if $i < p$?

\begin{theorem}\label{thm:25.4}
Let $K$ be a field of characteristic $p$, and let $0 \ne D \in \Der(K)$.
\begin{enumerate}[(i)]
  \item $1, D, D^2, \ldots, D^{p-1}$ are linearly independent over $K$;
  \item the only way in which $c_0 + c_1D + \cdots + c_{p-1}D^{p-1}$ with $c_i \in K$ can be
        a derivation is if $c_0 = c_1 = \cdots = c_{p-1} = 0$.
\end{enumerate}
\end{theorem}

\begin{proof}
For $a \in K$, write $a_L$ for the operation of multiplying by $a$; then the property
$D(ax) = D(a) \cdot x + a \cdot Dx$ of a derivation means that $D \circ a_L = D(a)_L + a_L \circ D$.
We can write the Leibnitz formula as
\[
  D^i \circ a_L = aD^i + i \cdot D(a)D^{i-1} + \cdots,
\]
our proof exploits this formula.

(i) For some $i < p$ suppose that $1, D, \ldots, D^{i-1}$ are linearly independent
over $K$, but that $1, D, \ldots, D^i$ are not. Then we can write
$D^i = c_{i-1}D^{i-1} + \cdots + c_0$, with $c_j \in K$. If we choose some $a \in K$ such that
$D(a) \ne 0$, then in view of $D^i \circ a_L = c_{i-1}D^{i-1} \circ a_L + \cdots$, we get
\[
  aD^i + i \cdot D(a)D^{i-1} + \cdots = c_{i-1}aD^{i-1} + \cdots,
\]
where $\cdots$ indicates a linear combination of $1, D, \ldots, D^{i-2}$. Subtracting $a$
times our original relation from this gives a relation of the form
\[
  i \cdot D(a)D^{i-1} = \cdots,
\]
and this contradicts the assumption that $1, D, \ldots, D^{i-1}$ are linearly
independent.

(ii) Suppose that $E = c_iD^i + \cdots + c_1D + c_0$ is a derivation of $K$, with $i < p$
and $c_i \ne 0$. Then $E(1) = c_0$, so that $c_0 = 0$. Now if $i > 1$ then take $a \in K$
such that $D(a) \ne 0$, and substitute both sides of $E \circ a_L = c_i a_L \circ D^i + \cdots$ in the
Leibnitz formula: we get
\[
  aE + E(a)_L = ac_iD^i + [i \cdot c_i \cdot D(a) + ac_{i-1}]D^{i-1} + \cdots,
\]
but then in view of the linear independence of $1, D, \ldots, D^{p-1}$, the coefficients
of $D^{i-1}$ on both sides must be equal; therefore $i \cdot c_i \cdot D(a) = 0$, which is a
contradiction.
\end{proof}

\begin{remarks}
(i) The theorem also holds if $\chark K = 0$.

(ii) If $K$ is not a field, this result does not necessarily hold. For example, let
$k$ be a field of characteristic $p$, and set $A = k[X]/(X^p) = k[x]$, with
$x^p = 0$; then every derivation of $k[X]$ will take the ideal $(X^p)$ into itself,
and therefore induces a derivation of $A$. In particular, the derivation
$X^{p-1} \cdot d/dX$ of $k[X]$ induces $D \in \Der_k(A)$ such that $D(x) = x^{p-1}$, but
$D^i(x) = i \cdot x^{i-1}x^{p-1} = 0$ if $i > 1$, and therefore for $p > 2$ we have $D^p = 0$.
\end{remarks}

\begin{theorem}[the Hochschild formula]\label{thm:25.5}
Let $A$ be a ring of characteristic $p$; then for $a \in A$ and $D \in \Der(A)$ we have
\[
  (aD)^p = a^pD^p + (aD)^{p-1}(a) \cdot D.
\]
\end{theorem}

\begin{proof}
Set $E = aD$. Then $E^2 = E \circ a_L \circ D = (aE + E(a))D = a^2D^2 + E(a)D$, and
proceeding by induction, we get a relation of the form
\[
  E^k = a^kD^k + \sum_{i=2}^{k-1} b_{k,i}D^i + E^{k-1}(a)D,
\]
where $b_{k,i}$ are elements of $A$ given by a purely formal computation, so that
\[
  b_{k,i} = f_{k,i}(a, D(a), D^2(a), \ldots, D^{k-1}(a)),
\]
where the $f_{k,i}$ are polynomials with coefficients in $\mathbb{Z}/(p)$ not depending on
$A$, on $a$ or on $D$. Now to prove our theorem we need only show that $f_{p,i} = 0$
for $1 < i < p$. Let $k$ be a field of characteristic $p$, and let $x_1, x_2, \ldots$ be a
countable number of indeterminates over $k$; set $K = k(x_1, x_2, \ldots)$. Define a
$k$-derivation $D$ of $K$ by $Dx_i = x_{i+1}$. (Since $\Omega_{K/k}$ is the free $K$-module with
basis $dx_1, dx_2, \ldots$, given any $f_i \in K$ there exists a unique $D \in \Der_k(K)$
such that $Dx_i = f_i$.) For this $D$, we set $E = x_1D$; then since
$E^p - x_1^pD^p = b_{p,p-1}D^{p-1} + \cdots + b_{p,2}D^2 + E^{p-1}(a) \cdot D$ is a derivation,
by the previous theorem we must have $b_{p,i} = 0$ for $1 < i < p$. Therefore
\[
  b_{p,i} = f_{p,i}(x_1, x_2, \ldots, x_{p-i+1}) = 0,
\]
and this proves that $f_{p,i} = 0$.
\end{proof}

This formula is known as the Hochschild formula, although it is also reported to
have been first proved by Serre. Be that as it may, it is an important fact that
$(aD)^p$ is a linear combination of $D^p$ and $D$.

% ============ 习题 ============
\section*{Exercises to \S\ref{sec:25}}
\emph{Prove the following propositions.}

\begin{exercise}\label{ex:25.1-ex}
Let $A$ be a ring, $a, b \in A$ and $D, D' \in \Der(A)$; then
\[
  [aD, bD'] = ab[D, D'] + aD(b)D' - bD'(a)D.
\]
Hence in order for an $A$-submodule $\mathfrak{g} \subset \Der(A)$ to be closed under
$[\, , \,]$, it is enough to have $\mathfrak{g} = \sum_{i \in I} AD_i$ with
$[D_i, D_j] \in \mathfrak{g}$ for all $i, j \in I$.
\end{exercise}

\begin{exercise}\label{ex:25.2-ex}
Let $A$ be a ring containing the rational field $\mathbb{Q}$. Suppose that $x \in A$ and
$D \in \Der(A)$ are such that $Dx = 1$ and $\bigcap_{n} x^nA = (0)$; then $x$ is a non-
zero-divisor of $A$.
\end{exercise}

\begin{exercise}\label{ex:25.3-ex}
Let $A$ be a ring, and $I$ an ideal of $A$; set $\hat{A}$ for the $I$-adic completion of
$A$. Then for $D \in \Der(A)$ we have $D(I^n) \subset I^{n-1}$ for all $n > 0$, so that $D$
is $I$-adically continuous, and hence induces a derivation of $\hat{A}$. Also for a
multiplicative set $S \subset A$, a derivation $D$ induces a derivation of $A_S$ by
means of $D(a/s) = (D(a) \cdot s - a \cdot D(s))/s^2$.
\end{exercise}

\begin{exercise}\label{ex:25.4-ex}
Let $k$ be a ring, $k'$ and $A$ two $k$-algebras, and set $A' = k' \otimes_k A$; let
$S \subset A$ be a multiplicative set. Then $\Omega_{A_S/k'} = \Omega_{A/k} \otimes_k k' =
\Omega_{A/k} \otimes_A A'$, and $\Omega_{A_S/k} = \Omega_{A/k} \otimes_A A_S$.
\end{exercise}

\begin{exercise}\label{ex:25.5-ex}
Let $A$ be a ring of characteristic $p$, and $x \in A$, $D \in \Der(A)$ elements such
that $D^p = 0$ and $Dx = 1$; set $A_0 = \{a \in A \mid Da = 0\}$. Then $A_0$ is a subring of
$A$, and $A = A_0[x] = A_0 + A_0x + \cdots + A_0x^{p-1}$, with $1, x, \ldots, x^{p-1}$ linearly
independent over $A_0$.
\end{exercise}

% ============ §26 Separability (原书 pp.198-207) ============
\section{Separability}\label{sec:26}\index{separable algebra, section 26}\index{separably generated extension, section 26}\index{linearly disjoint, section 26}\index{p-basis, section 26}\index{differential basis, section 26}\index{imperfection module, section 26}\index{Cartier equality, section 26}\index{0-étale, section 26}

Let $k$ be a field and $A$ a $k$-algebra. We say that $A$ is \emph{separable} over $k$ if for
every extension field $k'$ of $k$, the ring $A' = A \otimes_k k'$ is reduced, that is does
not contain nilpotents. From the definition, one sees at once the following:
\begin{enumerate}[(1)]
  \item a subalgebra of a separable $k$-algebra is separable;
  \item $A$ is separable over $k$ if and only if every finitely generated $k$-
        subalgebra of $A$ is separable over $k$;
  \item for $A$ to be separable over $k$ it is sufficient that $A \otimes_k k'$ is reduced
        for every finitely generated extension field $k'$ of $k$;
  \item if $A$ is separable over $k$ and $k'$ is an extension field of $k$ then
        $A \otimes_k k'$ is separable over $k'$.
\end{enumerate}

\begin{remark}
When $A$ is a finite $k$-algebra, the separability condition can be checked using
the discriminant. The \emph{trace} of an element $a$ of $A$, denoted by $\tr_{A/k}(a)$, is
defined to be the trace of the $k$-linear mapping $A \to A$ induced by
multiplication by $a$. Let $\omega_1, \ldots, \omega_n$ be a linear basis of $A$ over $k$. Then
$d = \det(\tr_{A/k}(\omega_i\omega_j))$ is called a \emph{discriminant} of $A$ over $k$. If we use
another basis $\omega'_1, \ldots, \omega'_n$, and if $\omega_i = \sum_j c_{ij}\omega_j$, then the
discriminant with respect to this basis is $\det(c_{ij})^2 \cdot d$. Thus $d = 0$ or
$d \ne 0$ is a property of $A$ independent of the choice of basis. Now we claim that
$A$ is separable if and only if $d \ne 0$.

If $k'$ is an extension field of $k$ and $A' = A \otimes_k k'$, then $\omega_1, \ldots, \omega_n$
is also a linear basis of $A'$ over $k'$, and so $d$ is also a discriminant of $A'$
over $k'$. If $A'$ is not reduced, let $N = \nil(A')$. Take a basis $\omega'_1, \ldots, \omega'_n$
of $A'$ such that $\omega'_1, \ldots, \omega'_r$ span $N$. Then $\omega'_i\omega'_j$ is nilpotent for
$i \le r$, hence its trace is zero. It follows that $\det(\tr(\omega'_i\omega'_j)) = 0$, and so
$d = 0$. Conversely, if $A$ is separable over $k$, take an algebraic closure $K$ of $k$.
Then $A \otimes_k K$ is reduced. Therefore we need only prove that if $k$ is
algebraically closed and $A$ is reduced then $d \ne 0$. Now $A$ is an Artinian reduced
ring, hence is a finite product of fields, each of which (as a finite extension of
$k$) is isomorphic to $k$. Thus $A = ke_1 + \cdots + ke_n$ with $e_ie_j = 0$ for $i \neq j$ and
$e_i^2 = e_i$. Hence $\tr(e_i) = 1$ and $d \ne 0$.
\end{remark}

In what follows we consider mainly the case when $A$ is a field. If $K$ is an
algebraic extension field of $k$, and is separable in the usual sense (that is
every element of $K$ is a root of a polynomial with coefficients in $k$ having no
multiple roots), then $K$ is separable over $k$ in our sense. To see this, by (2)
above we can assume that $K$ is finitely generated over $k$, and then according to
the well-known primitive element theorem of field theory, $K \cong k[X]/(f(X))$,
where $f \in k[X]$ is irreducible and with no multiple roots. Then if $k'$ is an
extension field of $k$ we have
\[
  K \otimes_k k' = k'[X]/(f(X)),
\]
and when we factorise $f$ into primes in $k'[X]$ we get $f = f_1 \cdots f_s$ with
$(f_i, f_j) = 1$ for $i \neq j$, so that by Theorem~\ref{thm:1.4},
\[
  k'[X]/(f) = k'[X]/(f_1) \times \cdots \times k'[X]/(f_s);
\]
since this is a direct product of fields, it is reduced.

We say that an extension field $K$ of $k$ is \emph{separably generated} over $k$ if $K$ has
a separating transcendence basis over $k$, that is a transcendence basis $T$ such
that $K$ is a separable algebraic extension of $k(T)$.

\begin{theorem}\label{thm:26.1}
A separably generated extension field is separable.
\end{theorem}

\begin{proof}
Let $k$ be a field and $K$ a separably generated extension of $k$, with $T$ a
separating transcendence basis of $K$. If $k'$ is any extension field of $k$ then
$k(T) \otimes_k k'$ is a ring of fractions of $k[T] \otimes_k k' = k'[T]$, so that it is an
integral domain with field of fractions $k'(T)$. Thus
$K \otimes_k k' = K \otimes_{k(T)} (k(T) \otimes_k k')$ is a subring of $K \otimes_k k'(T)$. Now
$K$ is a separable algebraic extension of $k(T)$, so that as we have seen above
$K \otimes_{k(T)} k'(T)$ is a direct product of fields, hence reduced; therefore so is
$K \otimes_k k'$.
\end{proof}

\begin{theorem}\label{thm:26.2}
Let $k$ be a field of characteristic $p$, and $K$ a finitely generated extension
field of $k$; then the following conditions are equivalent:
\begin{enumerate}[(1)]
  \item $K$ is separable over $k$;
  \item $K \otimes_k k^{1/p}$ is reduced;
  \item $K$ is separably generated over $k$.
\end{enumerate}
\end{theorem}

\begin{proof}
(1) $\Leftarrow$ (3) has just been proved, and (3) $\Rightarrow$ (2) is clear from the
proof of Theorem~\hyperref[thm:26.1]{1}, since $k^{1/p}$ is an extension field of $k$.

(1) $\Rightarrow$ (3) Let $K = k(x_1, \ldots, x_n)$; we can assume that $x_1, \ldots, x_r$ is a
transcendence basis for $K$ over $k$. Assume furthermore that $x_{r+1}, \ldots, x_q$ are
separable algebraic over $k(x_1, \ldots, x_r)$, and that $x_{q+1}$ is not; set $y = x_{q+1}$
and let $f(Y^p)$ be the minimal polynomial of $y$ over $k(x_1, \ldots, x_r)$. The
coefficients of $f(Y^p)$ are rational functions of $x_1, \ldots, x_r$, so that clearing
denominators we get an irreducible polynomial
$F(X_1, \ldots, X_r, Y^p) \in k[X_1, \ldots, X_r, Y]$, with $F(x, y^p) = 0$. Now if
$\partial F/\partial X_i = 0$ for $1 \le i \le r$ then $F(X, Y)$ is the $p$-th power of a polynomial
$G(X, Y)$ with coefficients in $k^{1/p}$, but then we would have
\[
  k[x_1, \ldots, x_r, y] \otimes_k k^{1/p} = (k[X, Y]/(F(X, Y))) \otimes_k k^{1/p}
      = k^{1/p}[X, Y]/(G(X, Y)^p);
\]
this is a subring of $K \otimes_k k^{1/p}$ containing nilpotent elements, and this
contradicts (2). Hence we can assume that $\partial F/\partial X_1 \ne 0$. Then $x_1$ is
separable algebraic over $k(x_2, \ldots, x_r, y)$, and hence so are $x_{r+1}, \ldots, x_q$.
Therefore exchanging $x_1$ and $y = x_{q+1}$, we find that $x_{r+1}, \ldots, x_{q+1}$ are
separably algebraic over $k(x_1, \ldots, x_r)$, so that by induction on $q$, we have
(3).
\end{proof}

\begin{remark}
As we have seen in the proof, if $K = k(x_1, \ldots, x_n)$ is separable over $k$ then we
can choose a separating transcendence basis from among $x_1, \ldots, x_n$.
\end{remark}

\begin{theorem}\label{thm:26.3}
If $k$ is a perfect field then every extension field $K$ of $k$ is separable over
$k$, and a $k$-algebra $A$ is separable if and only if it is reduced.
\end{theorem}

\begin{proof}
Recall that a field $k$ is perfect if every algebraic extension of $k$ is separable.
If $k$ has characteristic $0$ then every extension field $K$ is separably generated,
and therefore separable. In characteristic $p$, perfect implies $k = k^{1/p}$, so that
by the previous theorem, all subfields of $K$ finitely generated over $k$ are
separable, so that $K$ itself is separable over $k$. (Caution: $K$ may fail to be
separably generated over $k$; for a counter-example, let $x$ be an indeterminate
over $k$, and $K = k(x, x^{p^{-1}}, x^{p^{-2}}, \ldots)$.) Now we show that if $A$ is a
reduced $k$-algebra then $A$ is separable. We can assume that $A$ is finitely
generated over $k$. Then $A$ is a Noetherian ring, and the total ring of fractions
$K$ of $A$ is of the form $K = K_1 \times \cdots \times K_r$ by Ex.~\ref{ex:6.5-ex}. Each $K_i$ is
separable over $k$, so that $K$ is also separable, and hence so is its subring $A$.
\end{proof}

In general two subfields $K$, $K'$ of a given field $L$ are said to be \emph{linearly
disjoint}\index{linearly disjoint} over a common subfield $k$ if the following three equivalent
conditions are satisfied:
\begin{enumerate}[(a)]
  \item if $\alpha_1, \ldots, \alpha_n \in K$ are linearly independent over $k$ they are also
        linearly independent over $K'$;
  \item the same with $K$ and $K'$ interchanged;
  \item if we write $K[K']$ for the subring of $L$ generated by $K$ and $K'$, the
        natural map $K \otimes_k K' \to K[K']$ is an isomorphism.
\end{enumerate}

\textit{Proof of equivalence.} (a) $\Rightarrow$ (c) Let $\xi = \sum x_i \otimes y_i$ be an element of
the kernel of $K \otimes_k K' \to K[K']$. Suppose that $x_1, \ldots, x_r$ are linearly
independent over $k$, and that the remainder $x_{r+1}, \ldots, x_n$ are linear
combinations of them, and rewrite $\xi = \sum x_i \otimes y_i'$. The image in $K[K']$ of
$\xi$ is $\sum x_iy_i'$, but if this is $0$ then by (a) we have $y_i' = 0$ for all $i$, so that
$\xi = 0$. This proves (c).

(c) $\Rightarrow$ (a) is also easy; finally, since (c) is symmetric in $K$ and $K'$ we of
course also get (a) $\Leftrightarrow$ (b).

Let $k$ be a field of characteristic $p$, and $K$ an extension field of $k$. Inside
the algebraic closure of $K$, consider the subfields $k^{p^{-n}} =
\{\alpha \in \bar{K} \mid \alpha^{p^n} \in k\}$; these are purely inseparable extension fields of
$k$, and $\bigcup_n k^{p^{-n}}$ is the smallest perfect field containing $k$.

\begin{theorem}[S.~MacLane]\label{thm:26.4}
Let $k$ and $K$ be as above. We have
\begin{enumerate}[(i)]
  \item if $K$ is separable over $k$ then $K$ and $k^{p^{-n}}$ are linearly disjoint over $k$;
  \item if $K$ and $k^{p^{-n}}$ are linearly disjoint over $k$ for some $n > 0$ then $K$ is
        separable over $k$.
\end{enumerate}
\end{theorem}

\begin{proof}
(i) It is enough to show that $K$ and $k^{p^{-n}}$ are linearly disjoint for every $n$.
Suppose that $\alpha_1, \ldots, \alpha_m \in k^{p^{-n}}$ are linearly independent over $k$. If
$\sum \xi_i\alpha_i = 0$ with $\xi_i \in K$, then we have to show that $\xi_i = 0$ for all $i$. Since
$k$ and $K$ are as in the statement, we can assume that $\xi_1, \ldots, \xi_m$ are
linearly independent over $k$. If we set $k_1 = k(\alpha_1, \ldots, \alpha_m)$, so that $k_1$ is a
finite extension of $k$; for some sufficiently large $n$ we have $k^{p^{-n}} \supset k_1$, and
if we set $A = K \otimes_k k_1$, then $A$ is a reduced ring. However, $A$ is finite as a
$K$-module, so is a zero-dimensional ring, but the $p^n$-th power of any element of
$A$ is in $K$, so that we see that $A$ has only one prime ideal. Hence $A$ is a field,
and is isomorphic to $K[k_1]$. From this we get $\sum \xi_i \otimes \alpha_i = 0$, that is
$\xi_i = 0$ for all $i$.

(ii) If $K$ and $k^{p^{-n}}$ are linearly disjoint, then $k^{p^{-n}} \subset k^{p^{-m}}$, and hence
$K$ and $k^{p^{-1}}$ are also linearly disjoint over $k$, so that $K \otimes_k k^{p^{-1}}$ is a
field. If $K'$ is a subfield of $K$ which is finitely generated over $k$ then by
Theorem~\hyperref[thm:26.2]{2}, $K'$ is separable over $k$. Hence $K$ is also separable over $k$.
\end{proof}

Let $K$ be an extension field of a field $k$; then $\Omega_{K/k}$ is a vector space over
$K$, and is generated by $\{dx \mid x \in K\}$, so that there exists a subset $B \subset K$
such that $\{dx \mid x \in B\}$ forms a basis of the vector space $\Omega_{K/k}$. A subset
$B \subset K$ with this property is called a \emph{differential basis} of $K$ over $k$. The
following condition (*) is necessary and sufficient for a subset
$\{x_\lambda\}_{\lambda \in \Lambda} \subset K$ to form a differential basis for $K$ over $k$.

\textbf{(*)} if $y_\lambda \in K$ are specified for every $\lambda \in \Lambda$ in an arbitrary way, then
there exists a unique $D \in \Der_k(K)$ such that $D(x_\lambda) = y_\lambda$ for all $\lambda$.

For $x_1, \ldots, x_n \in K$, let us study the condition for $dx_1, \ldots, dx_n \in \Omega_{K/k}$ to
be linearly independent over $K$. If $k$ has characteristic $0$ then this is equivalent
to $x_1, \ldots, x_n$ being algebraically independent over $k$. Indeed, if there exists
$f(X) \in k[X_1, \ldots, X_n]$ such that $f(x_1, \ldots, x_n) = 0$, then choose such a
relation of smallest degree; suppose for instance that $X_1$ actually appears in $f$,
so that $f_1 = \partial f/\partial X_1$ is non-zero, but of smaller degree than $f$, and hence
$f_1(x) \ne 0$. Then $f(x) = 0$ gives
\[
  0 = df = \sum_i f_i(x)dx_i,
\]
so that $dx_1, \ldots, dx_n$ are linearly dependent. Conversely, if $x_1, \ldots, x_n$ are
algebraically independent over $k$ then there exists a transcendence basis $B$ of
$K/k$ containing these, so that there exists $k$-derivations $D_i$ of the purely
transcendental extension $k(B)$ satisfying
\[
  D_i(x_i) = 1 \quad \text{and} \quad D_i(x) = 0 \text{ for } x \in B, x \neq x_i,
\]
(namely, $\partial/\partial x_i$). Moreover, $K$ is a separable algebraic extension of $k(B)$, and
so by Theorem~\hyperref[thm:25.3]{25.3} is 0-\'etale, so that the derivations $D_i$ extend to
derivations from $K$ to $K$. Then since $D_i(x_j) = \delta_{ij}$, the differentials
$dx_1, \ldots, dx_n \in \Omega_{K/k}$ are clearly linearly independent. Thus in this case a
differential basis is the same thing as a transcendence basis.

Now consider a field $k$ of characteristic $p$. We say that elements
$x_1, \ldots, x_n \in K$ of an extension field $K$ are \emph{$p$-independent} over $k$ if
$[K^p(k, x_1, \ldots, x_n) : K^p(k)] = p^n$, and a subset $B \subset K$ is $p$-independent if any
finite subset of $B$ is $p$-independent. This condition means precisely that the
set
\[
  \Gamma_B = \left\{ x_1^{\alpha_1} \cdots x_n^{\alpha_n} \;\middle|\;
      \begin{array}{l}
        x_1, \ldots, x_n \text{ are distinct} \\
        \text{elements of } B \text{ and } 0 \le \alpha_i < p
      \end{array} \right\}
\]
is linearly independent over $K^p(k)$; the elements of $\Gamma_B$ are called the
\emph{$p$-monomials} of $B$. If $B$ is not $p$-independent, we say it is $p$-dependent.
The condition of $p$-independence is not just a property of $B$ and $k$, but also
depends on $K$. If $B \subset K$ is $p$-independent over $k$ and $K = K^p(k, B)$, we say
that $B$ is a \emph{$p$-basis} of $K/k$. If $C \subset K$ is $p$-independent over $k$ then one
can easily show by Zorn's lemma that there exists a $p$-basis of $K/k$ containing
$C$.

One sees easily that $B$ is a $p$-basis of $K/k$ is equivalent to $\Gamma_B$ being a
basis of $K$ over $K^p(k)$ in the sense of linear algebra. If this holds then any map
$D\colon B \to K$ has a unique extension to an element $D \in \Der_k(K)$. Indeed, for a
$p$-monomial of $B$ we set
\[
  D(x_1^{\alpha_1} \cdots x_n^{\alpha_n}) = \sum_{r=1}^n \alpha_r x_1^{\alpha_1} \cdots
      x_r^{\alpha_r - 1} \cdots x_n^{\alpha_n}D(x_r),
\]
and extend $D$ to $K$ as a $K^p(k)$-linear map; then $D$ is a $k$-derivation. Thus a
$p$-basis $B$ is a differential basis of $K/k$. Conversely, if $B'$ is a differential
basis of $K/k$ then $B'$ is $p$-independent over $k$; for if $x_1, \ldots, x_n \in B'$ are
$p$-dependent, we can assume that $x_1 \in K^p(k, x_2, \ldots, x_n)$, so that we can
write $x_1 = f(x_2, \ldots, x_n)$, where $f$ is a polynomial with coefficients in $K^p(k)$.
Then in $\Omega_{K/k}$ we get $dx_1 = \sum_{i \ge 2}(\partial f/\partial x_i)dx_i$, which
contradicts the linear independence of $dx_1, \ldots, dx_n$. Now if we take a $p$-basis
$B$ of $K/k$ containing $B'$, then since both $B$ and $B'$ are differential bases, we
have $B = B'$. We summarise the above as follows:

\begin{theorem}\label{thm:26.5}
The notion of differential basis coincides with transcendence basis in
characteristic $0$, and with $p$-basis in characteristic $p$.
\end{theorem}

Now we look at the relation between separability and differential bases. Letting
$\pi \subset k$ denote the prime subfield of $k$, we write $\Omega_k$ for $\Omega_{k/\pi}$.

\begin{theorem}\label{thm:26.6}
For a field extension $K/k$, the following conditions are equivalent:
\begin{enumerate}[(1)]
  \item $K/k$ is separable;
  \item for any subfield $k' \subset k$ the standard map $\Omega_{k/k'} \otimes_k K \to \Omega_{K/k'}$
        is injective;
  \item[$2'$] for any subfield $k' \subset k$ and any differential basis $B$ of $k/k'$, there
        exists a differential basis of $K/k'$ containing $B$;
  \item $\Omega_k \otimes_k K \to \Omega_K$ is injective;
  \item any derivation of $k$ to an arbitrary $K$-module $M$ extends to a
        derivation from $K$ to $M$.
\end{enumerate}
\end{theorem}

\begin{proof}
(2) and (2$'$) are clearly equivalent, and (2) $\Rightarrow$ (3) $\Rightarrow$ (4) are trivial. In
characteristic $0$, both (1) and (2$'$) hold, so that we need only consider the case
$p > 0$.

(1) $\Rightarrow$ (2$'$) Since $K$ and $k^{1/p}$ are linearly disjoint over $k$, we can apply the
isomorphism $x \mapsto x^p$ to all three of these to get $K^p$ and $k$ linearly disjoint
over $k^p$. Hence $K^p(k^p, k') = K^p(k')$ and $k$ are linearly disjoint over
$k^p(k')$ (see Ex.~\ref{ex:26.1-ex} below). If we choose a $p$-basis $B$ of $k$ over $k'$ then
the set $\Gamma_B$ of $p$-monomials of $B$ is linearly independent over $k^p(k')$, hence
also linearly independent over $K^p(k')$, and $B$ as a subset of $K$ is also
$p$-independent over $k'$. Therefore $B$ can be extended to a $p$-basis of $K/k'$.

(2$'$) $\Rightarrow$ (1) If we take a $p$-basis $B$ of $k$ over $\pi$ then the set $\Gamma_B$ of
$p$-monomials is a basis of $k$ over $k^p$. $\{dx \mid x \in B\}$ is a basis of $\Omega_k$ over
$k$, and by assumption is linearly independent in $\Omega_K$ over $K$, so that
$\Gamma_B$ is also linearly independent over $K^p$. Therefore
\[
  k \otimes_{k^p} K^p \cong k(K^p),
\]
and $k$ and $K^p$ are linearly disjoint over $k^p$, so that by
Theorem~\hyperref[thm:26.4]{4}, $K/k$ is separable.
\end{proof}

Let $k$ be a field of characteristic $p$, and $\pi \subset k$ the prime subfield; a $p$-basis
of $k$ over $\pi$ is called an \emph{absolute $p$-basis} of $k$. If $k_1 \subset k$ is any perfect field
contained in $k$ then $k^p(k_1) = k^p = k^p(\pi)$, so that an absolute $p$-basis of $k$ is
the same as a $p$-basis over any perfect subfield of $k$.

\begin{theorem}\label{thm:26.7}
Let $k$ be a field of characteristic $p$, and $K$ an extension field of $k$. If an
absolute $p$-basis of $k$ is also an absolute $p$-basis of $K$, then $K$ is 0-\'etale over
$k$, and conversely.
\end{theorem}

\begin{proof}
Consider a commutative diagram of ring homomorphisms
\[
\begin{tikzcd}
K \arrow[r, "\bar{u}"] & \bar{C} = C/N \\
k \arrow[u, "\subset"] \arrow[ru, "j"'] & C \arrow[u]
\end{tikzcd}
\]
where $C = C/N$ is as in the definition with $N$ an ideal of $C$ satisfying $N^2 = 0$,
and $g$ the natural map. For $\alpha \in K$, if we choose $a \in C$ such that
$\bar{u}(\alpha) = g(a)$, then $a^p$ is independent of the choice of $a$. For if
$g(a) = g(a')$ then we can write $a' = a + x$ with $x \in N$, and since $C$ is a $k$-algebra
via $j$ it is a ring of characteristic $p$, so that
\[
  a'^p = a^p + x^p = a^p.
\]
Now we define a map $v_0\colon K^p \to C$ by $v_0(\alpha^p) = a^p$ for $\alpha \in K$, where $a$ is
such that $\bar{u}(\alpha) = g(a)$; one checks easily that $v_0$ is a homomorphism, and
coincides with $j$ on $k^p$. So far we have not used the assumption on $K/k$. Now
since by assumption $K$ is separable over $k$, and $K = K^p[k]$, we can think of $K$ as
\[
  K = K^p \otimes_{k^p} k,
\]
and thus we can define $v\colon K \to C$ by letting $v$ be equal to $v_0$ on $K^p$, and
equal to $j$ on $k$; this is a lifting of $\bar{u}$ to $C$. Uniqueness of the lifting is
clear from the fact that $K^p(k) = K$, so that $\Omega_{K/k} = 0$.

Conversely, if we assume that $K/k$ is 0-\'etale, then first of all, from
0-unramified we have $\Omega_{K/k} = 0$, so that by 0-smoothness and
Theorem~\hyperref[thm:25.1]{25.1} we have $\Omega_K = \Omega_k \otimes_k K$. Thus an absolute $p$-basis of $k$ is
also an absolute $p$-basis of $K$.
\end{proof}

\begin{theorem}\label{thm:26.8}
Let $K/k$ be a separable extension of fields of characteristic $p$, and let $B$ be a
$p$-basis of $K/k$. Then $B$ is algebraically independent over $k$, and $K$ is 0-\'etale
over $k(B)$.
\end{theorem}

\begin{proof}
Suppose by contradiction that $b_1, \ldots, b_n \in B$ are algebraically dependent
over $k$. Suppose that $0 \ne f \in k[X_1, \ldots, X_n]$ is a polynomial of minimal degree
among all those with $f(b) = 0$, and set $\deg f = d$. Then write
\[
  f(X_1, \ldots, X_n) = \sum_{0 \le i_\nu < p} g_{i_1, \ldots, i_n}(X_1^p, \ldots, X_n^p)
      X_1^{i_1} \cdots X_n^{i_n}.
\]
Then since $b_1, \ldots, b_n$ are $p$-independent over $k$ and $f(b) = 0$, we have
$g_{i_1, \ldots, i_n}(b^p) = 0$ for all values of $i_1, \ldots, i_n$. However, since
\[
  d \ge \deg g_{i_1, \ldots, i_n}(X^p) + i_1 + \cdots + i_n,
\]
by choice of $f$ we must have $f(X) = g_{0, \ldots, 0}(X^p)$. Hence we can write $f$ in the
form $f(X) = h(X)^p$, with $h \in k^{1/p}[X_1, \ldots, X_n]$. However, since $K$ and $k^{1/p}$
are linearly disjoint over $k$, the monomials of degree $< d$ in $b_1, \ldots, b_n$, being
linearly independent over $k$, must also be linearly independent over $k^{1/p}$.
Thus $h(b) \ne 0$, but this contradicts $h(b)^p = f(b) = 0$. For the second half, see
the proof of the following theorem.
\end{proof}

\begin{remark}
Although $k(B)$ is purely transcendental over $k$, it does not follow that $K$ is
algebraic over $k(B)$; for a counter-example, let $K = k(x, x^{p^{-1}}, x^{p^{-2}}, \ldots)$,
with $x$ an indeterminate over $k$. In this case $B = \emptyset$.
\end{remark}

\begin{theorem}\label{thm:26.9}
If $K$ is a separable field extension of a field $k$, then $K$ is 0-smooth over $k$,
and conversely.
\end{theorem}

\begin{proof}
Let $B$ be a differential basis of $K/k$. If $K/k$ is separable then by
Theorems~\hyperref[thm:26.5]{5} and \hyperref[thm:26.8]{8}, $k(B)$ is purely transcendental over $k$. Then, as
one sees at once from the definition, $k(B)$ is 0-smooth over $k$. Moreover,
$K/k(B)$ is 0-\'etale. In characteristic $0$ this follows from
Theorem~\hyperref[thm:25.3]{25.3}; in characteristic $p$, by
Theorem~\hyperref[thm:26.6]{6} we have an exact sequence
$0 \to \Omega_k \otimes K \to \Omega_K \to \Omega_{K/k} \to 0$, so that putting together an absolute
$p$-basis of $k$ with $B$ we have an absolute $p$-basis of $K$, and this is clearly
also an absolute $p$-basis of $k(B)$. Hence $K/k(B)$ is 0-\'etale by
Theorem~\hyperref[thm:26.7]{7}. Therefore $K/k$ is 0-smooth.

Conversely, if $K/k$ is 0-smooth then by Theorem~\hyperref[thm:25.1]{25.1}, $\Omega_k \otimes K \to \Omega_K$
is injective, so that by Theorem~\hyperref[thm:26.6]{6}, $K/k$ is separable.
\end{proof}

\medskip
\noindent\textbf{Imperfection modules and the Cartier equality}

Quite generally, if $k \to A \to B$ are ring homomorphisms, we write $I_{B/A/k}$ for
the kernel of $\Omega_{A/k} \otimes_A B \to \Omega_{B/k}$, and call it the \emph{imperfection module} of
the $A$-algebra $B$ over $k$. If $k = \mathbb{Z}$ or $k = \mathbb{Z}/p$ (the prime field of characteristic
$p$) we omit $k$, and write $I_{B/A}$.

\begin{lemma}\label{lem:26.1}
Let $k \to K \to L \to L'$ be field homomorphisms. Then there exists a natural
exact sequence
\[
  0 \to I_{L/K/k} \otimes_K L' \to I_{L'/K/k} \to I_{L'/L/k}
      \to \Omega_{L/K} \otimes_L L' \to \Omega_{L'/K} \to \Omega_{L'/L} \to 0.
\]
\end{lemma}

\begin{proof}
We have a commutative diagram with exact rows
\[
\begin{tikzcd}
0 \arrow[r] & I_{L/K/k} \otimes_K L' \arrow[r] \arrow[d] &
\Omega_{K/k} \otimes_K L' \arrow[r] \arrow[d, "="] &
\Omega_{L/k} \otimes_L L' \arrow[r] \arrow[d, "\cong"] &
\Omega_{L/K} \otimes_L L' \arrow[r] \arrow[d] & 0 \\
0 \arrow[r] & I_{L'/K/k} \arrow[r] &
\Omega_{K/k} \otimes_K L' \arrow[r] &
\Omega_{L'/k} \arrow[r] &
\Omega_{L'/K} \arrow[r] & 0
\end{tikzcd}
\]
We abbreviate this as
\[
\begin{tikzcd}
0 \arrow[r] & X \arrow[r] \arrow[d, "i"] & A \arrow[r] \arrow[d, "="] & B \arrow[r] \arrow[d, "f_1"] & P \arrow[r] \arrow[d, "f_2"] & 0 \\
0 \arrow[r] & Y \arrow[r] & A \arrow[r] & C \arrow[r] & Q \arrow[r] & 0
\end{tikzcd}
\]
and from it construct
\[
\begin{tikzcd}
0 \arrow[r] & A/X \arrow[r] \arrow[d] & B \arrow[r] \arrow[d, "f_1"] & P \arrow[r] \arrow[d, "f_2"] & 0 \\
0 \arrow[r] & A/Y \arrow[r] & C \arrow[r] & Q \arrow[r] & 0
\end{tikzcd}
\]
applying the snake lemma gives the exact sequence
\[
  0 \to Y/X \to \Ker f_1 \to \Ker f_2 \to 0,
\]
from which we easily get
\[
  0 \to X \to Y \to \Ker f_1 \to P \to Q \to \Coker f_2 \to 0.
\]
This is just what we wanted to prove.
\end{proof}

\begin{theorem}[the Cartier equality]\label{thm:26.10}
Let $k$ be a perfect field, $K$ an extension of $k$ and $L$ a finitely generated
extension field of $K$; then
\[
  (*) \quad \dim_L \Omega_{L/K} = \trdeg_k L + \rank_L I_{L/K/k},
\]
(where $\rank_L$ denotes the dimension of a vector space over $L$).
\end{theorem}

\begin{proof}
Suppose that $k \subset K \subset L \subset L'$, with $L$ finitely generated over $K$ and $L'$
finitely generated over $L$. If the theorem holds for $k \subset L \subset L'$ and for
$k \subset K \subset L$ then both $I_{L'/L/k}$ and $I_{L/K/k} \otimes_K L$ are finite-dimensional over
$L$, so that by the lemma, $I_{L'/K/k}$ is also finite-dimensional, and
\begin{align*}
  \dim_{L'} \Omega_{L'/K} - \dim_{L'} \Omega_{L'/k}
      &= (\dim_{L'} \Omega_{L'/L} - \dim_{L'} \Omega_{L/k}) +
          (\dim_{L'} \Omega_{L/K} - \dim_L \Omega_{L/k}) \\
      &= \trdeg_k L' + \trdeg_K L = \trdeg_k L';
\end{align*}
thus the theorem also holds for $k \subset K \subset L'$. Now every finitely generated
field extension can be obtained by a succession of the following three kinds of
simple extensions:
\begin{enumerate}[(1)]
  \item $L = K(\alpha)$ where $\alpha$ is transcendental over $K$;
  \item $L = K(\alpha)$ where $\alpha$ is separable algebraic over $K$;
  \item $L = K(\alpha)$ where $\chark K = p$, and $\alpha^p = a \in K$, but $\alpha \notin K$.
\end{enumerate}
Hence we need only prove (*) in each of these special cases. (1) and (2) are
easy. For (3), if we write $L = K[X]/(X^p - a)$ we see that
\[
  \Omega_{L/k} = (\Omega_{K[X]/k} \otimes_{K[X]} L)/L\,da
      = (\Omega_{K/k} \otimes_K L \oplus L\,da)/L\,da,
\]
and $d\alpha \ne 0$. Furthermore, since $k$ is a perfect field, we have
$a \notin K^p = kK^p$, so that in $\Omega_{K/k}(= \Omega_K)$ we have $da \ne 0$,
$\rank \Omega_{L/K} = 1$ and $\rank I_{L/K/k} = 1$, so that (*) also holds in this case.
\end{proof}

\begin{remark}[\textbf{Harper's theorem}]
An ideal $I$ of a ring $R$ is called a \emph{differential ideal} if it maps into itself under
every derivation of $R$ to $R$. A ring $R$ is said to be \emph{differentiably simple} if it has
no non-trivial differential ideals. The following beautiful theorem is due to
L.~Harper, Jr.

\begin{theorem*}
A Noetherian ring $R$ of characteristic $p$ is differentiably simple if and only if
it has the form $R = k[T_1, \ldots, T_n]/(T_1^p, \ldots, T_n^p)$, where $k$ is a field of
characteristic $p$.
\end{theorem*}

The `if' part is easy. The proof of the `only if' part is not so easy and we refer
the reader to Harper [1] and Yuan [1]. Recently this theorem was used by
Kimura-Niitsuma [1] to prove the following theorem which had been known as
Kunz' Conjecture:

\begin{theorem*}
Let $R$ be a regular local ring of characteristic $p$, and let $S$ be a local subring
of $R$ containing $R^p$. Assume that $R$ is a finite $S$-module. Then $R$ has a
$p$-basis over $S$ if and only if $S$ is regular.
\end{theorem*}
\end{remark}

% ============ 习题 ============
\section*{Exercises to \S\ref{sec:26}}
\emph{Prove the following propositions.}

\begin{exercise}\label{ex:26.1-ex}
Let $L$ be a field and $k$, $k'$, $K$, $K'$ subfields of $L$; assume that $k \subset k' \subset K$
and $k \subset K'$, and that $K$ and $K'$ are linearly disjoint over $k$. Then we have
(i) $K \cap K' = k$, and (ii) $K$ and $k'(K')$ are linearly disjoint over $k'$.
\end{exercise}

\begin{exercise}\label{ex:26.2-ex}
Let $L$ be a separable extension of a field $K$; then $L((T_1, \ldots, T_n))$ is separable
over $K((T_1, \ldots, T_n))$. Here $L((T_1, \ldots, T_n))$ denotes the field of fractions of
$L[[T_1, \ldots, T_n]]$.
\end{exercise}

% ============ §27 Higher derivations (原书 pp.207-212) ============
\section{Higher derivations}\label{sec:27}\index{higher derivation, section 27}\index{iterative higher derivation, section 27}\index{p-basis, section 27}

Let $k \xrightarrow{f} A \xrightarrow{g} B$ be ring homomorphisms. Let $t$ be an indeterminate
over $B$, and set $B_m = B[t]/(t^{m+1})$ for $m = 0, 1, \ldots$, and $B_\infty = B[[t]]$. We can
view $B_m$ as a $k$-algebra in a natural way (for $m \le \infty$).

For $m < \infty$ we define a \emph{higher derivation} (over $k$) of length $m$ from $A$ to
$B$ to be a sequence $D = (D_0, D_1, \ldots, D_m)$ of $k$-linear maps
$D_i\colon A \to B$, satisfying the conditions
\[
  (*) \quad D_0 = g \quad \text{and} \quad D_i(xy) = \sum_{j=0}^i D_j(x)D_{i-j}(y)
      \quad \text{for } 1 \le i \le m \text{ and } x, y \in A.
\]
These conditions are equivalent to saying that the map $E_t\colon A \to B_m$, defined by
\[
  E_t(x) = \sum_{i=0}^m D_i(x)t^i,
\]
is a $k$-algebra homomorphism with $E_t(x) = g(x) \mod t$.

When $A = B$ and $g = 1$ (the identity map of $A$) then we speak simply of a higher
$k$-derivation of $A$. In what follows we consider mainly this case.

If $D = (D_0, D_1, D_2, \ldots)$ is a higher derivation then $D_1 \in \Der_k(A, B)$.
Furthermore, $D_i$ is $0$ on $f(k)$ for $i > 0$. In general we say that $D$ is \emph{trivial}
on $a \in A$ if $D_i(a) = 0$ for $i > 0$; this means precisely that $a$ goes over to a
constant under the homomorphism $E_t$ corresponding to $D$.

The theory of higher derivations was initiated by Hasse and F.~K.~Schmidt [1].
In view of this, in this book we write $\mathrm{HS}_k(A, m)$ for the set of all higher
$k$-derivations of length $m$ of $A$, and we also write $\mathrm{HS}_k(A)$ for
$\mathrm{HS}_k(A, \infty)$. When we are not concerned with $k$, we simplify this to
$\mathrm{HS}(A, m)$ and $\mathrm{HS}(A)$. These sets do not have a module structure like that
of $\Der_k(A)$, but they do have a group structure (generally non-Abelian), which we
now explain. The homomorphism $E_t\colon A \to A_m$ corresponding to
$D \in \mathrm{HS}_k(A, m)$ can be extended to an endomorphism of $A[[t]]$ by setting
\[
  E_t\Big(\sum a_\nu t^\nu\Big) = \sum E_t(a_\nu)t^\nu.
\]
Now $E_t$ is injective, since if $\xi = a_\nu t^\nu + a_{\nu+1}t^{\nu+1} + \cdots \in A[[t]]$ with
$a_\nu \ne 0$ then $E_t(\xi) \equiv a_\nu t^\nu \mod t^{\nu+1}$; also, by setting
\[
  \xi - E_t(a_\nu t^\nu) = b_{\nu+1}t^{\nu+1} + \cdots,
\]
\[
  \xi - E_t(a_\nu t^\nu + b_{\nu+1}t^{\nu+1}) = c_{\nu+2}t^{\nu+2} + \cdots,
\]
and proceeding in the same way, we see easily that $E_t$ is surjective. In other
words, $E_t$ is an automorphism of $A[[t]]$. Conversely, an automorphism $E$ of the
$k$-algebra $A[[t]]$ satisfying $E(a) = a \mod t$ corresponds to a higher derivation. Thus
if we identify $\mathrm{HS}_k(A, m)$ with a subgroup of the automorphism group
$\Aut_k(A[[t]])$ of $A[[t]]$, it acquires the group structure which we are looking for.
Let us start computing this structure. For $D = (D_0, D_1, \ldots)$, $D' = (D'_0, D'_1, \ldots)$,
set
\[
  D \cdot D' = (D_0'', D_1'', \ldots), \quad \text{and} \quad D^{-1} = (D_0^*, D_1^*, \ldots);
\]
then
\begin{align*}
  E_t(E'_t(a)) &= E_t(a + D'_1(a)t + D'_2(a)t^2 + \cdots) \\
               &= (a + D_1(a)t + D_2(a)t^2 + \cdots) \\
               &\quad - (D'_1(a) + D_1(D'_1(a))t + D_2(D'_1(a))t^2 + \cdots)t \\
               &\quad + (D'_2(a) + D_1(D'_2(a))t + \cdots)t^2 + \cdots \\
               &= a + (D_1 + D'_1)(a)t + (D_2 + D_1D'_1 + D'_2)(a)t^2 + \cdots,
\end{align*}
so that
\[
  D_i'' = \sum_{p+q=i} D_pD_q' \quad \text{for all } i;
\]
and the $D_i^*$ are obtained by solving $\sum_{p+q=i} D_pD_q^* = 0$ for $i > 0$, that is
\[
  D_0^* = D_1^* = 1, \quad D_2^* = -D_2, \quad D_3^* = D_2^2 - D_3,
\]
\[
  D_4^* = -D_4 + D_2D_3 + D_3D_2 - D_2^3, \ldots
\]
(where $D_i^*$ is viewed as the operator $a \mapsto D_i^*(a)$; the products are
compositions).

If $S \subset A$ and $T \subset B$ are multiplicative sets such that $g(S) \subset T$, then the
given homomorphism $g\colon A \to B$ induces a homomorphism $A_S \to B_T$. Now we show
that in a similar way, a higher derivation has a unique extension to the
localisations. To see this, let $D = (D_0, D_1, \ldots, D_m)$ be a higher derivation of
length $m$ from $A$ to $B$; if we compose the homomorphism $E_t\colon A \to B_m$
corresponding to $D$ with the localisation $B_m \to (B_T)_m$, then an element $x \in S$
maps to $g(x)_T + D_1(x)_T t + \cdots$, and this is a unit of $(B_T)_m$, since the
constant term $g(x)_T$ is a unit of $B_T$. This $E_t$ induces a homomorphism
$A_S \to (B_T)_m$, which provides a higher derivation $A_S \to B_T$.

Let $D = (D_0, D_1, \ldots, D_m)$ be a higher $k$-derivation of length $m < \infty$ from $A$
to $B$. Consider the problem of extending this to a higher derivation of length
$m + 1$. If $E_m\colon A \to B_m$ is the homomorphism corresponding to $D$, the problem
of extending $D$ is equivalent to that of lifting $E_m$ to a homomorphism
$A \to B_{m+1}$. The following theorem is then clear from this.

\begin{theorem}\label{thm:27.1}
If the ring $A$ is 0-smooth over a ring $k$, then a higher derivation of length
$m < \infty$ over $k$ from $A$ to an $A$-algebra $B$ can be extended to a derivation of
length $\infty$.
\end{theorem}

This theorem can be applied for example to the case of a field $k$ and a separable
extension field $A$ of $k$.

If $A$ is a ring of characteristic $p$ then an ordinary derivation of $A$ is zero on
the subring $A^p$, but higher derivations need not vanish on $A^p$. If
$D = (D_0, D_1, D_2, \ldots)$ is a higher derivation of length $m \ge p$ then from
$E_t(a^p) = E_t(a)^p = a^p + D_1(a)^p t^p + \cdots$, we get
\[
  D_p(a^p) = D_1(a)^p, \quad \text{and in general } D_{pn}(a^p) = D_n(a)^p.
\]
For example, it follows from this that if $k$ is a field of characteristic $p$, and
$K = k(\alpha)$ with $\alpha^p \in k$ but $\alpha \notin k$, then although there exists
$D \in \Der_k(K)$ such that $D(\alpha) = 1$, this $D$ cannot be extended to a higher
$k$-derivation of length $\ge p$. Since $K$ is separable over the prime subfield, $D$
can be extended to a higher derivation of length $\infty$ (over the prime subfield),
but this extension cannot be trivial on $k$.

We say that a higher derivation $D = (D_0, D_1, \ldots) \in \mathrm{HS}_k(A)$ is \emph{iterative} if
it satisfies the following conditions:
\[
  \binom{i+j}{i}D_{i+j} = D_iD_j \quad \text{for all } i, j.
\]
This condition is equivalent to asking that the following diagram is commutative:
\[
\begin{tikzcd}
A \arrow[r, "E_{t+u}"] \arrow[rd, "E_t"'] & A[[t, u]] \\
& A[[t]] \arrow[u, "\subset"]
\end{tikzcd}
\]
where $E_t(a) = \sum t^nD_n(a)$ and $E_t(\sum t^\nu a_\nu) = \sum t^\nu E_t(a_\nu)$, and the
right-hand vertical arrow is the inclusion map. Indeed,
\[
  E_t(E_u(a)) = E_t\Big(\sum u^nD_n(a)\Big) = \sum t^i u^j D_iD_j(a)
\]
and
\[
  E_{t+u}(a) = \sum (t + u)^nD_n(a) = \sum_{i,j} \binom{i+j}{i} t^i u^j D_{i+j}(a).
\]
If $A$ contains the rational field $\mathbb{Q}$, then one sees by induction on $n$ that an
iterative higher derivation satisfies $D_n = D_1^n/n!$, and is hence determined by
$D_1$ only. Conversely, for $D \in \Der_k(A)$, we see that
$(1, D, D^2/2!, D^3/3!, \ldots)$ is an iterative higher derivation. If $A$ has
characteristic $p$ then for an iterative $D = (D_0, D_1, \ldots)$ we have $D_i = D_1^i/i!$
for $i < p$, and $D_p = 0$. Thus one cannot hope to extend any derivation to an
iterative higher derivation, even if $A$ is a field.

\begin{theorem}\label{thm:27.2}
Let $k \to A \xrightarrow{g} B$ be ring homomorphisms, and suppose that $B$ is 0-\'etale over
$A$. Then given $D = (D_0, D_1, \ldots) \in \mathrm{HS}_k(A, B)$, there exists a unique
$D' = (D'_0, D'_1, \ldots) \in \mathrm{HS}_k(B)$ such that $D'_i(g(a)) = D_i(a)$ for all $i$.
Moreover, if $D^* = (D_0^*, D_1^*, \ldots) \in \mathrm{HS}_k(A)$ is such that $D_i = g \circ D_i^*$
for all $i$, and if $D^*$ is iterative, then $D'$ is also iterative.
\end{theorem}

\begin{proof}
There is no problem about $D'_0 = 1_B$. Now assume that
$(D'_0, \ldots, D'_m) \in \mathrm{HS}_k(B, m)$ has been constructed so that
$D'_i \circ g = D_i$ for $i < m$; then if we define
$h\colon A \to B_{m+1} = B[t]/(t^{m+2})$ by $h(a) = \sum_{i=0}^{m+1} t^iD_i(a)$,
and $u\colon B \to B_{m+1}$ by $u(b) = \sum_{i=0}^{m} t^iD'_i(b)$, we obtain the left-hand
commutative diagram
\[
\begin{tikzcd}
A \arrow[r, "h"] \arrow[d, "g"'] & B_{m+1} \arrow[d, "p"] \\
B \arrow[ru, "u"'] \arrow[r, "\bar{g}"'] & B
\end{tikzcd}
\]
where $p\colon B_{m+1} \to B$ is the natural map and $\bar{g}$ the induced map. Hence by the
0-\'etale assumption, there exists a unique $v\colon B \to B_{m+1}$ which makes the
right-hand diagram commutative. Repeating this we see that $D'$ exists and is
unique. If $D^*$ is iterative, and we consider the homomorphisms
$E_t\colon A \to A[[t]]$ and $E'_t\colon B \to B[[t]]$ corresponding respectively to $D^*$ and
$D'$, then we know $E_t \circ E_u = E_{t+u}$, and we need only prove that
$E'_t \circ E'_u = E'_{t+u}$. By induction on $m$, assume that
\[
  E'_u(E'_t(b)) = E'_{t+u}(b) \mod (t, u)^{m+1} \quad \text{for all } b \in B;
\]
then from the commutativity of
\[
\begin{tikzcd}
B \arrow[r, "E'_u"] \arrow[d, "g"'] & B[[t, u]]/(t, u)^{m+1} \arrow[d] \\
A \arrow[r, "E_t"] & A[[t, u]] \arrow[r] & B[[t, u]]/(t, u)^{m+2}
\end{tikzcd}
\]
from $E_{t+u} = E_t \circ E_u$ and from the assumption that $B$ is 0-\'etale over $A$, we get
\[
  E'_u(E'_t(b)) = E'_{t+u}(b) \mod (t, u)^{m+2} \quad \text{for all } b \in B.
\]
This proves that $E'_t \circ E'_u = E'_{t+u}$.
\end{proof}

\begin{remark}
The above $D'$ will be called an \emph{extension} of $D$ (or of $D^*$) to $B$ (even if $A$ is
not a subring of $B$).
\end{remark}

\begin{theorem}\label{thm:27.3}
\begin{enumerate}[(i)]
  \item Let $A$ be a ring of characteristic $p$, and suppose that $x \in A$,
        $D \in \Der(A)$ satisfy $Dx = 1$ and $D^p = 0$; set $A_0 = \{a \in A \mid Da = 0\}$.
        Then $A$ is a free module over $A_0$, with basis $1, x, \ldots, x^{p-1}$.
  \item Let $k$ be a field of characteristic $p$, and $K$ a separable extension of $k$;
        let $D \in \Der_k(K)$ be such that $D \ne 0$, $D^p = 0$. Set
        $K_0 = \{a \in K \mid Da = 0\}$. Then there exists an $x \in K$ such that $Dx = 1$,
        and a subset $B_0 \subset K_0$ such that $B = \{x\} \cup B_0$ is a $p$-basis for $K$ over
        $k$.
\end{enumerate}
\end{theorem}

\begin{proof}
(i) Suppose that $a_0 + a_1x + \cdots + a_ix^i = 0$ for some $i < p$ and $a_j \in A_0$.
Then applying $D^i$ we get $i!a_i = 0$, hence $a_i = 0$. Hence by downward induction
on $i$ we see that $1, x, \ldots, x^{p-1}$ are linearly independent over $A_0$. Now in
view of $D^p = 0$, for every $a \in A$ we have $D^{i+1}a = 0$ for some $0 \le i < p$. If
$i = 0$ we have $a \in A_0$. If $i > 0$ then $D^i(a - x^iD^ia/i!) = 0$, so that we see by
induction that
\[
  D^{i+1}a = 0 \Rightarrow a \in A_0 + A_0x + \cdots + A_0x^i;
\]
setting $i = p - 1$ in this gives $A = A_0 + A_0x + \cdots + A_0x^{p-1}$.

(ii) Since $D \ne 0$ we can find $z \in K$ such that $Dz \ne 0$. Now in view of
$D^pz = 0$, there is an $i$ such that $D^iz \ne 0$ but $D^{i+1}z = 0$. If we set $y = D^iz$
and $x = (D^{i-1}z)/y$ then $Dx = 1$, so that by (i) we have $K = K_0(x)$ and
$[K : K_0] = p$. Now if we had $x^p \in K_0^pk$, then $x \in K_0k^{1/p}$, and we could
write $x = \sum \omega_i\alpha_i$ with $\omega_1, \ldots, \omega_n \in K_0$ linearly independent over
$k$, and $\alpha_i \in k^{1/p}$. Now $k \subset K_0$ and $x \notin K_0$, so that
$x, \omega_1, \ldots, \omega_n$ are linearly independent over $k$, and hence by the assumption
that $K$ is separable over $k$, they are also linearly independent over $k^{1/p}$,
which contradicts $x = \sum \omega_i\alpha_i$. Thus $x^p \notin K_0^pk$. Hence we can
choose a $p$-basis $C$ of $K_0$ over $k$ such that $x^p \in C$; set $B_0 = C - \{x^p\}$. Then
if $y_1, \ldots, y_m$ are distinct elements of $B_0$, we have
$[K_0^pk(x^p, y_1, \ldots, y_m) : K_0^pk] = p^{m+1}$, and together with $K = K_0(x)$ this
gives $[K^pk(y_1, \ldots, y_m) : K^pk] = p^m$. Thus $B_0$ as a subset of $K$ is
$p$-independent over $k$. Since also $K_0 = K_0^pk(x^p, B_0)$, we have
$K = K_0(x) = K^pk(x, B_0)$, so that setting $B = B_0 \cup \{x\}$ we get a $p$-basis of $K$
over $k$.
\end{proof}

\begin{theorem}\label{thm:27.4}
Let $K$ be a field of characteristic $p$, and $k \subset K$ a subfield such that $K$ is
separable over $k$. Then a necessary and sufficient condition for $D \in \Der_k(K)$
to extend to an iterative element of $\mathrm{HS}_k(K)$ is that $D^p = 0$.
\end{theorem}

\begin{proof}
We have already seen necessity, and we prove sufficiency. We can assume that
$D \ne 0$; if $D^p = 0$ then we can choose $K_0$, $x$ and $B_0$ as in Theorem~\hyperref[thm:27.3]{3},
(ii). We set $K' = k(B_0)$; then $D$ is a $K'$-derivation, $K$ is 0-\'etale over
$K'(x)$, and $K'(x)$ is a purely transcendental extension of $K'$. We define a
homomorphism $E_t\colon K'(X) \to K'(X)[t]$ by setting $E_t(a) = a$ for $a \in K'$ and
$E_t(X) = X + t$; then
\[
  E_t(E_u(X)) = X + u + t = E_{t+u}(X),
\]
so that $E_t \circ E_u = E_{t+u}$ holds over the whole of $K'(X)$. Thus $E_t$ defines an
iterative higher derivation $\bar{D}$ of $K'(X)$ over $K'$. Since $K$ is 0-\'etale over
$K'(X)$, by Theorem~\hyperref[thm:27.2]{2} there is an extension of $\bar{D}$ to an iterative higher
derivation of $K$ over $K'$; the term of degree $1$ in $D$ is $D$, so that $D$ is an
extension of $D$ (or more precisely, of $(1, D)$).
\end{proof}

% ============ 习题 ============
\section*{Exercises to \S\ref{sec:27}}

\begin{exercise}\label{ex:27.1-ex}
Let $k$ be a ring, $A$ a $k$-algebra and $D \in \Der_k(A)$. Say that $D$ is
\emph{integrable} over $k$ if there exists an extension $D' \in \mathrm{HS}_k(A)$ with
$D' = (D_0, D_1, \ldots)$ and $D = D_1$ (then $D'$ is an \emph{integral} of $D$); set
$\mathrm{Ider}_k(A) = \{D \in \Der_k(A) \mid D \text{ is integrable over } k\}$. Then prove
that $\mathrm{Ider}_k(A) = \Der_k(A)$ is an $A$-submodule.
\end{exercise}

\begin{exercise}\label{ex:27.2-ex}
In the notation of this section, consider the construction of $E_t\colon A \to A[[t]]$
corresponding to $D$; then if $t' \in A[[t]]$ is any power series with no constant
term, we have $A[[t']] \subset A[[t]]$, so that $E_t\colon A \to A[[t']]$ can be composed
to a homomorphism $E_t\colon A \to A[[t]]$, to give a different higher derivation.
Thus for $D = (D_0, D_1, \ldots)$, the homomorphism $E_{t'}$ corresponds to the higher
derivation $D' = (D_0, 0, D_1, 0, D_2, \ldots)$; taking the product $D \cdot D'$ we get an
integral of $D$, different from $D$. Thus for given $D \in \mathrm{Ider}_k(A)$ there will in
general exist many integrals of $D$; verify that if $D^p = 0$ and we impose the
condition that the integral should be iterative, it is still not uniquely
determined.
\end{exercise}


% ============ Chapter 10. I-smoothness (原书 pp.213-245) ============
\chapter{$I$-smoothness}\label{ch:10}

$I$-smoothness is a notion which Grothendieck obtained by reformulating the
theory of simple (non-singular) points in algebraic geometry in terms of an
algebraic `infinitesimal analysis', which makes effective use of nilpotent
elements. The definition looks complicated at first sight, but it has various
alternative formulations, and is a natural and useful notion. In \S\ref{sec:28},
along with the general theory of $I$-smoothness following [G1] we prove the
existence of a coefficient field for a complete local ring of equal
characteristic, relating this to the author's idea of quasi-coefficient field
(Theorem~\ref{thm:28.3}), and discuss Faltings' very simple proof of the
equivalence of $m$-smoothness and geometric regularity for local rings. In
\S\ref{sec:29} we deduce the existence of a coefficient ring for a complete local ring
of unequal characteristic from Theorem~\ref{thm:28.10}, and prove some classical
theorems of Cohen on complete local rings; these results are of decisive
significance for the usefulness of taking completions. \S\ref{sec:30} is something of a
jumble of various theories, but is for the most part occupied with the
so-called Jacobian criterion for regularity. On this subject, we treat the
simple and powerful method obtained by the author's 1972 seminar in the case
of a ring containing a field of characteristic $0$; in the most difficult case of
power series rings in characteristic $p$, the only method currently available
is that of Nagata, and we explain this as simply as possible.

% ============ §28 I-smoothness (原书 pp.213-222) ============
\section{$I$-smoothness}\label{sec:28}\index{I-smooth, section 28}\index{I-unramified, section 28}\index{I-étale, section 28}\index{coefficient field, section 28}\index{quasi-coefficient field, section 28}\index{geometrically regular, section 28}

Let $A$ be a ring, $B$ an $A$-algebra, and $I$ an ideal of $B$; we consider $B$ with the
$I$-adic topology. We say that $B$ is \emph{$I$-smooth} over $A$ if given an $A$-algebra $C$,
an ideal $N$ of $C$ satisfying $N^2 = 0$, and an $A$-algebra homomorphism
$\bar{u}\colon B \to C/N$ which is continuous for the discrete topology of $C/N$ (that is,
such that $\bar{u}(I^\nu) = 0$ for some $\nu$), then there exists a lifting $u\colon B \to C$ of
$\bar{u}$ to $C$.
\[
\begin{tikzcd}
B \arrow[r, "\bar{u}"] & C/N \\
A \arrow[u] \arrow[r] & C \arrow[u]
\end{tikzcd}
\]
If $I = (0)$, that is if no continuity condition is imposed on $\bar{u}$, then this is the
definition of 0-smooth given in \S\ref{sec:25}. Write $f'\colon C \to C/N$ for the natural map;
then from $f'u(I^\nu) = \bar{u}(I^\nu) = 0$ we have $u(I^\nu) \subset N$, hence
$u(I^{2\nu}) \subset N^2 = 0$, so that $u\colon B \to C$ (assuming it exists) is continuous for
the discrete topology of $C$. From this, one sees that if $B$ is $I$-smooth over $A$, and
instead of the condition $N^2 = 0$ we assume that $C$ is an $N$-adically complete
ring, then a continuous homomorphism $\bar{u}\colon B \to C/N$ has a lifting $u\colon B \to C$,
and $u$ is continuous with respect to the $N$-adic topology of $C$; this is because
we can lift $\bar{u}$ successively to $B \to C/N^i$ for $i = 1, 2, \ldots$, and then $u$ is given
by $B \to \varprojlim C/N^i = C$.

We now return to the original assumption $N^2 = 0$; we say that $B$ is
\emph{$I$-unramified} (or $I$-neat) over $A$ if given $C$, $N$ and a continuous
homomorphism $\bar{u}\colon B \to C/N$, there exists at most one lifting of $\bar{u}$ to $C$. If
$B$ is both $I$-smooth and $I$-unramified over $A$, we say that $B$ is \emph{$I$-\'etale}.
These conditions become weaker if we replace $I$ by a larger ideal.

\begin{theorem}[Transitivity]\label{thm:28.1}
Let $A \to B \xrightarrow{g'} B'$ be ring homomorphisms, and suppose that $g'$ is
continuous for the $I$-adic topology of $B$ and the $I'$-adic topology of $B'$; if $B$ is
$I$-smooth over $A$, and $B'$ is $I'$-smooth over $B$ then $B'$ is $I'$-smooth over $A$.
The same thing holds with $I$-unramified in place of $I$-smooth.
\end{theorem}

\begin{proof}
Suppose that $\bar{u}$ is given in the diagram;
\[
\begin{tikzcd}
B' \arrow[r, "\bar{u}"] & C/N \\
B \arrow[u, "g'"] \arrow[ru, "\overline{ug'}"'] & C \arrow[u] \\
A \arrow[u] \arrow[ruu] & 
\end{tikzcd}
\]
then since $\bar{u}g'\colon B \to C/N$ is continuous, by the $I$-smoothness of $B$, there
exists a lifting $w\colon B \to C$. Next by the $I'$-smoothness of $B'$ over $B$, we can
lift $\bar{u}$ to $v\colon B' \to C$. Also if $B$ is $I$-unramified over $A$, and the map $v$ in
the diagram exists, then $w = vg'$ is unique, and if in addition $B'$ is
$I'$-unramified over $B$ then $v$ is unique.
\end{proof}

\begin{theorem}[Base-change]\label{thm:28.2}
Let $A$ be a ring, $B$ and $A'$ two $A$-algebras, and set $B' = B \otimes_A A'$. If $B$ is
$I$-smooth over $A$ then $B'$ is $IB'$-smooth over $A'$. The same thing holds for
$I$-unramified.
\end{theorem}

\begin{proof}
We have the diagram
\[
\begin{tikzcd}
B \arrow[r, "p"] & B' \arrow[r, "\bar{u}"] & C/N \\
A \arrow[u] \arrow[r] & A' \arrow[u] \arrow[ru, ""] & C \arrow[u]
\end{tikzcd}
\]
where $p$, $q$ are the natural homomorphisms. Then if $\bar{u}$ satisfies
$\bar{u}(I^\nu B') = 0$, there is a lifting $v\colon B \to C$ of $\bar{u}p$. Then if we define
$v'\colon B' = B \otimes_A A' \to C$, by $v' = v \otimes_A \mathrm{id}$, this is a lifting of $\bar{u}$ to $C$. For
unramified, this is clear from the fact that $v'$ is uniquely determined by $v$.
\end{proof}

\begin{example*}
Let $k$ be a ring, $(A, m)$ a local ring, $(\hat{A}, \hat{m})$ its completion, and
$k \to A$ a homomorphism. Then
\begin{enumerate}[(i)]
  \item $\hat{A}$ is $\hat{m}$-\'etale over $A$;
  \item $A$ is $m$-smooth (or $m$-unramified) over $k$ $\iff$ $\hat{A}$ is $\hat{m}$-smooth (or
        $\hat{m}$-unramified) over $k$.
\end{enumerate}
We get a proof at once from the fact that $A/m^\nu \cong \hat{A}/\hat{m}^\nu$ for all $\nu$.
\end{example*}

\begin{example*}
Let $A$ be any ring, and set $B = A[X_1, \ldots, X_n]$ and $I = \sum X_iB$; we give $B$
the $I$-adic topology. Then $B$ is $I$-smooth over $A$.
\end{example*}

\begin{remark}
The gap between $I$-smoothness and 0-smoothness has been studied by Tanimoto
[1], [2]. For instance, if $k$ is a field, then $k[X_1, \ldots, X_n]$ is 0-smooth over
$k$ only when $\chark k = p$ and $[k : k^p] < \infty$.
\end{remark}

Let $(A, m, K)$ be a local ring. If $A$ is of characteristic $p$, then also $\chark K = p$;
moreover, if $\chark K = 0$, then $\chark A = 0$, and $A$ contains the rational number
field $\mathbb{Q}$. In either of these cases $A$ is said to be \emph{equicharacteristic}, or a
local ring of \emph{equal characteristic}; this is equivalent to saying that $A$ contains a
field. If $A$ is not of equal characteristic, then either
\[
  \chark A = 0 \text{ and } \chark K = p,
\]
or
\[
  \chark A = p^n \text{ for some } n > 1 \text{ and } \chark K = p.
\]
In this case we say that $A$ is a local ring of \emph{unequal characteristic}.

Let $A$ be an equicharacteristic local ring and let $K'$ be a subfield of $A$. We say
that $K'$ is a \emph{coefficient field} of $A$ if $K'$ maps isomorphically to $K$ under the
natural map $A \to A/m = K$, or equivalently, if $A = K' + m$. Moreover, we say that
$K'$ is a \emph{quasi-coefficient field} of $A$ if $K$ is 0-\'etale over $K'$ (or rather, over
the image of $K'$ in $K$).

\begin{theorem}\label{thm:28.3}
Let $(A, m, K)$ be an equicharacteristic local ring. Then
\begin{enumerate}[(i)]
  \item $A$ has a quasi-coefficient field;
  \item if $A$ is complete, it has a coefficient field;
  \item if the residue field $K$ of $A$ is separable over a subfield $k \subset A$ then
        $A$ has a quasi-coefficient field $K'$ containing $k$;
  \item if $K'$ is a quasi-coefficient field of $A$, then there exists a unique
        coefficient field $K''$ of the completion $\hat{A}$ containing $K'$.
\end{enumerate}
\end{theorem}

\begin{proof}
(iii) Suppose that $B = \{\tau_1, \tau_2, \ldots\}$ is a differential basis of $K/k$, and for
each $\tau_i$, choose an inverse image $\xi_i \in A$. Then by
Theorem~\hyperref[thm:26.8]{26.8}, $\tau_1, \tau_2, \ldots$ are algebraically independent over $k$, so
that the subring $k[\xi_1, \xi_2, \ldots]$ of $A$ meets $m$ in $(0)$, and hence $A$ contains
the field $K' = k(\xi_1, \xi_2, \ldots)$. We identify $K'$ with its image $k(B)$ in $K$, so that
$K$ is clearly 0-\'etale over $K'$, and $K'$ is a quasi-coefficient field, as required.

(i) By assumption $A$ contains a field, so that it contains a perfect field (for
example, the prime subfield). We need only apply (iii) to this.

(iv) In the diagram
\[
\begin{tikzcd}
K \arrow[r, "\cong"] & \hat{A}/\hat{m} \\
K' \arrow[u, "\subset"] \arrow[ru, "\subset"'] & \hat{A} \arrow[u, "p"] \\
& A \arrow[u, "\subset"]
\end{tikzcd}
\]
there exists a unique lifting of the identity map $K \to \hat{A}/\hat{m}$ to $K \to \hat{A}$,
and its image is the required coefficient field.

(ii) follows from (i) and (iv).
\end{proof}

The next lemma will be made more precise in Theorem~\hyperref[thm:28.7]{7}.

\begin{lemma}\label{lem:28.1}
Suppose that $(A, m, K)$ is a Noetherian local ring containing a field $k$. If $A$ is
$m$-smooth over $k$ then $A$ is regular. The converse holds if the residue field $K$
is separable over $k$.
\end{lemma}

\begin{proof}
Take a perfect subfield $k_1 \subset k$; then $k$ is 0-smooth over $k_1$, so that by
transitivity, $A$ is also $m$-smooth over $k_1$, so that we can assume that $k$ is a
perfect field. Also, replacing $A$ by $\hat{A}$, we can assume that $A$ is complete.
Then $A$ has a coefficient field containing $k$; for ease of notation we write $K$
for this, and identify it with the residue field. If $\{x_1, \ldots, x_n\}$ is a minimal
basis of $m$ then as $K$-algebras we have
\[
  A/m^2 \cong K[X_1, \ldots, X_n]/(X_1, \ldots, X_n)^2.
\]
The composite
\[
  A \to A/m^2 \xrightarrow{\cong} K[X_1, \ldots, X_n]/(X)^2 \xrightarrow{\subset} K[X_1, \ldots, X_n]/(X)^2
\]
lifts to $A \to K[[X_1, \ldots, X_n]]$, and by Theorem~\hyperref[thm:8.4]{8.4}, this is surjective. Thus
$\dim A \ge \dim K[[X_1, \ldots, X_n]] = n$, and together with $\embdim A = n$ this
gives the regularity of $A$.

Conversely, if $A$ is regular and $K$ is separable over $k$, then $A$ has a
coefficient field $K$ containing $k$. Let $\{x_1, \ldots, x_n\}$ be a regular system of
parameters of $A$, and define a homomorphism of $K$-algebras
$\psi\colon K[[X_1, \ldots, X_n]] \to \hat{A}$ by $\psi(X_i) = x_i$; then once more by
Theorem~\hyperref[thm:8.4]{8.4}, $\psi$ is surjective, and comparing dimensions, we see that
\[
  K[[X_1, \ldots, X_n]] = \hat{A}.
\]
Therefore $\hat{A}$ is $\hat{m}$-smooth over $K$, and since $K$ is 0-smooth over $k$, we see
that $\hat{A}$ is $\hat{m}$-smooth over $k$, and therefore $A$ is $m$-smooth over $k$.
\end{proof}

Let $k \to A \xrightarrow{g} B$ be ring homomorphisms, and let $I$ be an ideal of $B$; we
consider $B$ in the $I$-adic topology. We say that $B$ is \emph{$I$-smooth over $A$ relative
to $k$} if the following condition holds: for any $A$-algebra $C$, and an ideal $N$ of
$C$ such that $N^2 = 0$, given an $A$-algebra homomorphism $\bar{u}\colon B \to C/N$
satisfying $\bar{u}(I^\nu) = 0$ for sufficiently large $\nu$, if $\bar{u}$ has a lifting
$v'\colon B \to C$ as a $k$-algebra homomorphism, it also has a lifting $v\colon B \to C$ as an
$A$-algebra homomorphism.

\begin{theorem}\label{thm:28.4}
Let $k \to A \xrightarrow{g} B$ and $I \subset B$ be as above; then the following conditions
are equivalent:
\begin{enumerate}[(1)]
  \item $B$ is $I$-smooth over $A$ relative to $k$;
  \item if $N$ is a $B$-module such that $I^\nu N = 0$ for sufficiently large $\nu$, then
        $\Der_k(B, N) \to \Der_k(A, N)$ is surjective;
  \item for every $\nu > 0$, the map $\Omega_{B/k} \otimes_B (B/I^\nu) \to
        \Omega_{A/k} \otimes_A (B/I^\nu)$ has a left inverse (that is, it maps injectively onto
        a direct summand).
\end{enumerate}
\end{theorem}

\begin{proof}
(1) $\Rightarrow$ (2) If $I^\nu N = 0$, set $C = (B/I^\nu) \ltimes N$, and let
$\bar{u}\colon B \to B/I^\nu = C/N$ be the natural map. Given $D \in \Der_k(A, N)$, define
$\lambda\colon A \to C$ by $\lambda(a) = (ug(a), D(a))$. If we consider $C$ as an $A$-algebra via
$\lambda$, then $b \mapsto (u(b), 0) \in C$ is a $k$-algebra homomorphism from $B$ to $C$
lifting $\bar{u}$, so that by assumption there exists a lifting $v'\colon B \to C$ of $\bar{u}$ as an
$A$-algebra homomorphism; then writing $v'$ in the form
\[
  v'(b) = (u(b), D'(b)), \quad \text{with } D' \in \Der_k(B, N),
\]
we have $v'g = \lambda$, so that $D'g = D$.

(2) $\Rightarrow$ (1) Suppose given a commutative diagram
\[
\begin{tikzcd}
B \arrow[r, "\bar{u}"] & C/N \\
A \arrow[u, "g"] \arrow[ru, ""] & C \arrow[u, "j"]
\end{tikzcd}
\]
with $j$ the natural map, and $\bar{u}(I^\nu) = 0$; if $v\colon B \to C$ is a $k$-algebra
homomorphism satisfying $jv = \bar{u}$ and $vg = \lambda$, then setting $D = \lambda - vg$, we
can view $D$ as an element of $\Der_k(A, N)$. By assumption there exists
$D' \in \Der_k(B, N)$ such that $D = D'g$. Using this, we set $v' = v + D'$; then
\[
  v'g = vg + D'g = \lambda - D + D = \lambda, \quad \text{and} \quad jv' = \bar{u}.
\]

(2) $\Leftrightarrow$ (3) comes from observing the general fact that for a ring $R$, a map
$\varphi\colon M \to M'$ of $R$-modules has a left inverse if and only if for every $R$-module
$N$ the induced map
\[
  \Hom_R(M', N) \to \Hom_R(M, N)
\]
is surjective.
\end{proof}

\begin{theorem}\label{thm:28.5}
Let $A$ be a ring, $B$ an $A$-algebra, and $I$ an ideal of $B$; set $\bar{B} = B/I$, and
assume that $B$ is $I$-smooth over $A$. Then $\Omega_{B/k} \otimes_B \bar{B}$ is projective as
a $\bar{B}$-module.
\end{theorem}

\begin{proof}
It is enough to show that for an exact sequence $L \to M \to 0$ of $\bar{B}$-modules, the
sequence
\[
  \Hom_{\bar{B}}(\Omega_{B/k} \otimes_B \bar{B}, L) \to \Hom_{\bar{B}}(\Omega_{B/k} \otimes_B \bar{B}, M) \to 0
\]
is exact, that is that
\[
  \Der_k(B, L) \to \Der_k(B, M) \to 0
\]
is exact. Set $C = \bar{B} \ltimes L$ and $N = \Ker \varphi$. If we view both $L$ and $N$ as
ideals of $C$, we have $L^2 = N^2 = 0$ and $C/N \cong \bar{B} \ltimes M$. Now for any
$D \in \Der_k(B, M)$ we have an $A$-algebra homomorphism $B \to C/N$ given by
\[
  b \mapsto (b, D(b)) \in \bar{B} \ltimes M,
\]
and lifting this to $B \to C$ is equivalent to lifting $D$ to an element of
$\Der_k(B, L)$.
\end{proof}

\begin{lemma}\label{lem:28.2}
Let $B$ be a ring and $I$ an ideal of $B$, and let $u\colon L \to M$ be a map of
$B$-modules; assume that $M$ is projective. Suppose also that one of the
following two conditions hold:
\begin{enumerate}[(a)]
  \item $I$ is nilpotent;
or ($\beta$) $L$ is a finite $B$-module and $I \subset \rad(B)$.
\end{enumerate}
Then
\[
  u \text{ has a left inverse} \iff \bar{u}\colon L/IL \to M/IM \text{ has a left inverse}.
\]
\end{lemma}

\begin{proof}
($\Leftarrow$) is trivial. To prove ($\Rightarrow$), suppose that $\bar{v}\colon M/IM \to L/IL$ is a
left inverse of $\bar{u}$. Since $M$ is projective, there is a map $w\colon M \to L$ such
that the diagram
\[
\begin{tikzcd}
M \arrow[r, "w"] \arrow[d] & L \arrow[d] \\
M/IM \arrow[r, "\bar{w}"'] & L/IL
\end{tikzcd}
\]
commutes. Set $\tau = uw$. Then $\tau$ induces the identity on $L/IL$, so that
$L = \tau(L) + IL$, so that by NAK, $L = \tau(L)$. Hence if $L$ is a finite $B$-module,
then by Theorem~\ref{thm:2.4}, $\tau$ is also injective. Furthermore, if $I^\nu = 0$ we do
the following: if $x \in \Ker \tau$ then $0 = \bar{\tau}(\bar{x}) = \bar{x} \mod IL$, so that $x \in IL$, and
we can write $x = \sum a_i y_i$ with $a_i \in I$ and $y_i \in L$. Then
\[
  0 = \tau(x) = \sum a_i \tau(y_i) \equiv \sum a_iy_i = x \mod I^2L,
\]
so that $x \in I^2L$, and proceeding in the same way we arrive at
$x \in I^\nu L = 0$. Hence also in this case $\tau$ is an automorphism of $L$, and
$\tau^{-1}u$ is the required left inverse of $u$.
\end{proof}

\begin{theorem}\label{thm:28.6}
Let $k \to A \to B$ be ring homomorphisms, $I$ an ideal of $B$, and suppose that $B$ is
$I$-smooth over $k$. Set $B_1 = B/I$. Then the following conditions are equivalent:
\begin{enumerate}[(1)]
  \item $B$ is $I$-smooth over $A$;
  \item $\Omega_{A/k} \otimes_A B_1 \to \Omega_{B/k} \otimes_B B_1$ has a $B_1$-linear left inverse.
\end{enumerate}
\end{theorem}

\begin{proof}
(1) $\Rightarrow$ (2) is contained in Theorem~\hyperref[thm:28.4]{4}. Conversely, suppose that (2)
holds. For any $\nu > 0$, set $B_\nu = B/I^\nu$; then since $I$-smoothness and
$I^\nu$-smoothness are the same, by Theorem~\hyperref[thm:28.5]{5}, $\Omega_{B/k} \otimes_B B_\nu$ is a
projective $B_\nu$-module. Now set $I_1 = I/I^\nu$; then $B_\nu/I_1 = B_1$, and
$I_1^2 = 0$, so that applying Lemma~\hyperref[lem:28.2]{2}, we see that
$\Omega_{A/k} \otimes_A B_\nu \to \Omega_{B/k} \otimes_B B_\nu$ has a left inverse. By
Theorem~\hyperref[thm:28.4]{4}, $B$ is $I$-smooth over $A$ relative to $k$, but since it is also
$I$-smooth over $k$, it is also $I$-smooth over $A$.
\end{proof}

\begin{corollary*}
Let $(A, m, K)$ be a regular local ring containing a field $k$; then
\[
  A \text{ is } m\text{-smooth over } k \iff \Omega_k \otimes_k K \to \Omega_A \otimes_A K \text{ is injective}.
\]
\end{corollary*}

\begin{proof}
Let $k_1 \subset k$ be the prime subfield. Then by Lemma~\hyperref[lem:28.1]{1}, $A$ is $m$-smooth
over $k_1$, so that we need only apply the theorem to $k_1 \to k \to A$.
\end{proof}

Let $A$ be a Noetherian local ring, and $k \subset A$ a subfield. We say that $A$ is
\emph{geometrically regular} over $k$ if $A \otimes_k k'$ is a regular ring for every finite
extension field $k'$ of $k$.

\begin{theorem}\label{thm:28.7}
Let $(A, m, K)$ be a Noetherian local ring, and $k \subset A$ a subfield; then
\[
  A \text{ is } m\text{-smooth over } k \iff A \text{ is geometrically regular over } k.
\]
\end{theorem}

\begin{proof}
($\Rightarrow$) Let $k'$ be a finite extension field of $k$. Then $A \otimes_k k' = A'$ is
$mA'$-smooth over $k'$ by base-change. Let $n$ be any maximal ideal of $A'$; then $A'$ is
a finite $A$-module, hence integral over $A$, so that $n \supset mA'$. Thus if we set
$A'' = A'_n$ and $m'' = nA''$ then $A' \to A''$ is continuous for the $mA'$-adic topology of
$A'$ and the $m''$-adic topology of $A''$, but as a localisation $A''$ is 0-\'etale over $A'$,
so that by Theorem~\hyperref[thm:28.1]{1}, $A''$ is $m''$-smooth over $k'$, and hence by
Lemma~\hyperref[lem:28.1]{1}, $A'' = A'_n$ is a regular local ring. This is what was required
to prove.

($\Leftarrow$) According to Lemma~\hyperref[lem:28.1]{1}, there is only a problem if $k$ is of
characteristic $p$. The proof we now give was discovered in 1977 by G.~Faltings
[1] while he was still a student.

By the corollary of the previous theorem, we need only prove that
$\Omega_k \otimes_k K \to \Omega_A \otimes_A K$ is injective. For this, we let
$x_1, \ldots, x_r \in k$ be $p$-independent elements, and prove that
$dx_1, \ldots, dx_r \in \Omega_A \otimes_A K$ are linearly independent over $K$. Write $\alpha_i$ for $p$-th
roots of the $x_i$, and set $k' = k(\alpha_1, \ldots, \alpha_r)$. Then
\[
  B = A \otimes_k k' = A[T_1, \ldots, T_r]/(T_1^p - x_1, \ldots, T_r^p - x_r)
\]
is a Noetherian local ring; write $n$ for its maximal ideal, and $L$ for the residue
field $L = B/n$. By Theorem~\hyperref[thm:25.2]{25.2} the sequence
\[
  0 \to n/n^2 \to \Omega_{B/k} \otimes_B L \to \Omega_{L/k} \to 0
\]
is exact. Similarly,
\[
  0 \to m/m^2 \to \Omega_{A/k} \otimes_A K \to \Omega_{K/k} \to 0
\]
is exact. Now consider the commutative diagram
\[
\begin{tikzcd}
0 \arrow[r] & n/n^2 \arrow[r] \arrow[d, "\varphi_1"'] & \Omega_{B/k} \otimes_B L \arrow[r] \arrow[d, "\varphi_2"'] & \Omega_{L/k} \arrow[r] \arrow[d, "\varphi_3"'] & 0 \\
0 \arrow[r] & (m/m^2) \otimes_K L \arrow[r] & \Omega_{A/k} \otimes_A L \arrow[r] & \Omega_{K/k} \otimes_K L \arrow[r] & 0
\end{tikzcd}
\]
where $\varphi_1$, $\varphi_2$ and $\varphi_3$ are the natural maps. Then by the snake lemma, we get a
long exact sequence of $L$-modules
\[
  0 \to \Ker \varphi_1 \to \Ker \varphi_2 \to \Ker \varphi_3
      \to \Coker \varphi_1 \to \Coker \varphi_2 \to \Coker \varphi_3 \to 0.
\]
By assumption $A$ and $B$ are regular local rings of the same dimension, so that
$\rank m/m^2 = \dim A = \rank n/n^2$, so $\rank \Ker \varphi_1$ and $\rank \Coker \varphi_1$ are
finite and equal; moreover, since $L$ is a finite extension field of $K$, both
$\rank \Ker \varphi_3$ and $\rank \Coker \varphi_3$ are finite and equal (the Cartier equality).
Therefore by the above exact sequence, we get
\[
  \rank \Ker \varphi_2 = \rank \Coker \varphi_2.
\]
However, $\Coker \varphi_2 = \Omega_{B/A} \otimes_B L$, and by Theorem~\hyperref[thm:25.2]{25.2},
\[
  \Omega_{B/A} = BdT_1 + \cdots + BdT_r,
\]
so that both of $\Ker \varphi_2$ and $\Coker \varphi_2$ have rank equal to $r$. Now if we set
$J = (T_1^p - x_1, \ldots, T_r^p - x_r)$ we have an exact sequence
\[
  J/J^2 \to \Omega_{A[T_1, \ldots, T_r]/k} \otimes_A B = \Omega_{A/k} \otimes_A B \oplus B dT_i \to \Omega_{B/k} \to 0,
\]
and since this remains exact after performing $\otimes_A L$, we see that $\Ker \varphi_2$ is
generated by $dx_1, \ldots, dx_r$. Therefore $dx_1, \ldots, dx_r \in \Omega_A \otimes_A L$ are
linearly independent over $L$, so that they must also be linearly independent
over $K$ as elements of $\Omega_A \otimes_A K$. This is what we had to prove.
\end{proof}

Let $B$ be an $A$-algebra, $I$ an ideal of $B$, and consider $B$ with the $I$-adic
topology. Let $N$ be a $B$-module such that $I^\nu N = 0$ for some $\nu > 0$; in what
follows we will say that a $B$-module with this property is \emph{discrete}. An
$A$-bilinear map $f\colon B \times B \to N$ will be called a \emph{continuous symmetric
2-cocycle} if it satisfies the three conditions
\begin{enumerate}[(a)]
  \item $xf(y, z) - f(xy, z) + f(x, yz) - f(x, y)z = 0$ for all $x, y, z \in B$,
  \item[($\beta$)] $f(x, y) = f(y, x)$,
  \item[($\gamma$)] there exists $\rho \ge \nu$ such that $f(x, y) = 0$ if either $x \in I^\rho$ or $y \in I^\rho$.
\end{enumerate}
If this holds, we set $f(1, 1) = \varepsilon$; then substituting $y = z = 1$ in (a) gives
$x\varepsilon = f(x, 1)$.

Define a product on the $A$-module $C = (B/I^\nu) \oplus N$ by
\[
  (x, v)(y, w) = (xy, -f(x, y) + xw + yv)
\]
for $x, y \in B$; then $C$ is a commutative ring with unit $(1, \varepsilon)$, and $N$ is an
ideal of $C$ satisfying $N^2 = 0$. If we define a map $A \to C$ by $a \mapsto (\bar{a}, a\varepsilon)$
then this is a ring homomorphism, and the diagram
\[
\begin{tikzcd}
B \arrow[r] & B/I^\nu = C/N \\
A \arrow[u] \arrow[r] & C \arrow[u]
\end{tikzcd}
\]
is commutative; then a necessary and sufficient condition for $\bar{u}$ to have a lifting
to $B \to C$ is that there should exist an $A$-linear map $g\colon B \to N$ such that
\[
  (\delta) \quad f(x, y) = xg(y) - g(xy) + g(x)y \quad \text{for all } x, y \in B.
\]
For if $g$ exists, then defining $v\colon B \to C$ by $v(x) = (\bar{x}, g(x))$ we find that $v$
is a lifting of $\bar{u}$, and conversely, if there is a lifting $v$ of $\bar{u}$ one checks easily
that one can find a $g$ as above.

We say that the 2-cocycle $f$ \emph{splits} if there exists $g$ satisfying ($\delta$). For any
$A$-linear map $g\colon B \to N$, we write $\delta g$ for the bilinear map
$B \times B \to N$ given by the right-hand side of ($\delta$); this satisfies (a) and ($\beta$),
and if $g$ is continuous (that is, $\exists \rho$ such that $g(I^\rho) = 0$), then it also
satisfies ($\gamma$).

\begin{theorem}\label{thm:28.8}
Let $A$ be a ring and $B$ an $A$-algebra with an $I$-adic topology.
\begin{enumerate}[(i)]
  \item If $B$ is $I$-smooth over $A$ then every continuous symmetric 2-cocycle
        $f\colon B \times B \to N$ with values in a discrete $B$-module $N$ splits.
  \item If $B/I^n$ is projective as an $A$-module for infinitely many $n$, and if every
        continuous symmetric 2-cocycle with values in a discrete $B$-module splits,
        then $B$ is $I$-smooth over $A$.
\end{enumerate}
\end{theorem}

\begin{proof}
(i) is what we have just said.

(ii) Suppose that we are given a commutative diagram
\[
\begin{tikzcd}
B \arrow[r, "\bar{u}"] & C/N \\
A \arrow[u] \arrow[r] & C \arrow[u]
\end{tikzcd}
\]
with $N^2 = 0$ and $\bar{u}(I^\nu) = 0$; then in view of $N^2 = 0$, the $C$-module $N$ can be
viewed as a $C/N$-module, and by means of $\bar{u}$ as a $B$-module; but then
$I^\nu N = 0$, so that $N$ is a discrete $B$-module. Take an integer $n > \nu$ such that
$B/I^n$ is projective as an $A$-module. Then $\bar{u}$ can be lifted as a map of $A$-modules
to $\lambda\colon B \to C$ such that $\lambda(I^n) = 0$. For $x, y \in B$ we set
\[
  f(x, y) = \lambda(xy) - \lambda(x)\lambda(y),
\]
and since $\lambda$ is a ring homomorphism modulo $N$, we have $f(x, y) \in N$. Now for
$\xi \in N$ and $x \in B$, by definition we have $\lambda(x) \cdot \xi = x \cdot \xi$ (both sides are
products evaluated in $C$), and using this one computes the left-hand side of (a) to
be zero. The symmetry ($\beta$) is obvious. Also $\lambda(I^n) = 0$, so that we also get
($\gamma$). Thus $f$ is a continuous symmetric 2-cocycle, and hence by assumption
there is an $A$-linear map $g\colon B \to N$ satisfying
\[
  f(x, y) = xg(y) - g(xy) + g(x)y.
\]
Now if we set $v = \lambda + g$, we have
\begin{align*}
  v(xy) &= \lambda(xy) + g(xy) \\
        &= \lambda(x)\lambda(y) + f(x, y) + g(xy) \\
        &= \lambda(x)\lambda(y) + xg(y) + g(x)y \\
        &= (\lambda(x) + g(x))(\lambda(y) + g(y)) = v(x)v(y),
\end{align*}
so that $v$ is an $A$-algebra homomorphism, and is a lifting of $\bar{u}$.
\end{proof}

\begin{theorem}[([G1], 19.7.1)]\label{thm:28.9}
Let $(A, m, k)$ and $(B, n, k')$ be Noetherian local rings, and $\varphi\colon A \to B$ a
local homomorphism; set $B_0 = B \otimes_A k = B/mB$ and $n_0 = n/mB$. Then the
following conditions are equivalent:
\begin{enumerate}[(1)]
  \item $B$ is $n$-smooth over $A$,
  \item $B$ is flat over $A$ and $B_0$ is $n_0$-smooth over $k$.
\end{enumerate}
\end{theorem}

\begin{proof}
This is an extremely important theorem, but the proof is long and difficult, and
we refer to [G1] for it. We content ourselves with proving the following
analogous theorem, which is all that we will need in what follows.
\end{proof}

\begin{theorem}\label{thm:28.10}
Let $(A, m, k)$ be a local ring, and $B$ a flat $A$-algebra; suppose that
$B_0 = B \otimes_A k$ is 0-smooth over $k$. Then $B$ is $mB$-smooth over $A$.
\end{theorem}

\begin{proof}
As one sees from the definition of $mB$-smoothness, it is enough to prove that
$B/m^\nu B$ is 0-smooth over $A/m^\nu$ for every $\nu > 0$. Since $B/m^\nu B$ is flat over
$A/m^\nu$, we can assume that $m$ is nilpotent. Then a flat $A$-module is free (by
Theorem~\ref{thm:7.10}), so that $B$ is a projective $A$-module, and hence by
Theorem~\hyperref[thm:28.8]{8}, we need only show that every symmetric 2-cocycle
$f\colon B \times B \to N$ with values in a $B$-module $N$ splits. First of all, in the case
that $N$ satisfies $mN = 0$, then since $f$ is $A$-bilinear $f$ is essentially a
2-cocycle over $B_0$, that is there is a map $f_0\colon B_0 \times B_0 \to N$ such that
$f(x, y) = f_0(\bar{x}, \bar{y})$.

Now $B_0$ is 0-smooth over $k$, so that by Theorem~\hyperref[thm:28.8]{8}, $f_0$ splits, that is
there is a map $g_0\colon B_0 \to N$ such that $f_0 = \delta g_0$. Thus setting
$g(x) = g_0(\bar{x})$ we get
\[
  f = \delta g.
\]
In the general case, write $\varphi$ for the natural map $N \to N/mN$, and consider
$\varphi \circ f$; then this splits, that is there exists $\bar{g}\colon B \to N/mN$ such that
$\varphi f = \delta \bar{g}$.

Now since $B$ is projective over $A$, we can lift $\bar{g}$ to an $A$-linear map
$g\colon B \to N$, and then $f - \delta g$ is a 2-cocycle with values in $mN$. Doing the same
thing once more, we find $h\colon B \to mN$ such that $f - \delta(g + h)$ is a 2-cocycle with
values in $m^2N$. Proceeding in the same way, since $m$ is nilpotent we finally see
that $f$ splits.
\end{proof}

% ============ 习题 ============
\section*{Exercises to \S\ref{sec:28}}
\emph{Prove the following propositions.}

\begin{exercise}\label{ex:28.1-ex}
Theorem~\hyperref[thm:28.10]{28.10} also holds on replacing smooth by unramified or \'etale.
\end{exercise}

\begin{exercise}\label{ex:28.2-ex}
Let $k$ be a non-perfect field of characteristic $p$, and $a \in k - k^p$; set
$A = k[X]_{(X^p - a)}$. Then the residue field $k(a^{1/p})$ of $A$ is inseparable over
$k$. This ring $A$ does not have a coefficient field containing $k$, but is
0-smooth over $k$.
\end{exercise}

% ============ §29 The structure theorems for complete local rings (原书 pp.223-230) ============
\section{The structure theorems for complete local rings}\label{sec:29}\index{coefficient ring, section 29}\index{p-ring, section 29}\index{quasi-coefficient ring, section 29}\index{Eisenstein extension, section 29}\index{ramified local ring, section 29}\index{unramified local ring, section 29}

By Theorem~\hyperref[thm:28.3]{28.3}, a complete local ring of equal characteristic $A$ has a
coefficient field. If $K$ is a coefficient field of $A$, and $x_1, \ldots, x_n$ are generators
of the maximal ideal, then any element of $A$ can be expanded as a power series
in $x_1, \ldots, x_n$ with coefficients in $K$, and therefore $A$ is a quotient of the
regular local ring $K[[X_1, \ldots, X_n]]$. We now want to extend all of this to the
case of unequal characteristic.

We say that a DVR of characteristic $0$ is a \emph{$p$-ring}\index{p-ring} if its maximal ideal is
generated by the prime number $p$. If $K$ is a given field of characteristic $p$, then
there exists a $p$-ring having $K$ as its residue field. This follows by applying the
next theorem to $A = \mathbb{Z}_{(p)}$.

\begin{theorem}\label{thm:29.1}
Let $(A, tA, k)$ be a DVR and $K$ an extension field of $k$; then there exists a DVR
$(B, tB, K)$ containing $A$.
\end{theorem}

\begin{proof}
Let $\{x_\lambda\}_{\lambda \in \Lambda}$ be a transcendence basis of $K$ over $k$, and set
$k_1 = k(\{x_\lambda\})$. We take indeterminates $\{X_\lambda\}_{\lambda \in \Lambda}$ over $A$ in
bijection with the $\{x_\lambda\}$, and set $A[\{X_\lambda\}] = A'$ and $A_1 = (A')_{tA'}$. Now
$A'$ is a free $A$-module, and hence separated for the $t$-adic topology; therefore,
so is $A_1$, and $A_1$ is a DVR with $A_1/tA_1 \cong k_1$. Hence, replacing $A$ and $k$ by $A_1$
and $k_1$, we can assume that $K$ is algebraic over $k$. Let $L$ be an algebraic
closure of the field of fractions of $A$, and let $\mathcal{F}$ be the set of all pairs
$(B, \varphi)$, where $B$ is an intermediate ring $A \subset B \subset L$ and
$\varphi\colon B \to K$ is an $A$-algebra homomorphism, satisfying the conditions
\[
  (*) \quad B \text{ is a DVR and } \Ker \varphi = \rad(B) \text{ is generated by } t.
\]
We introduce an order on $\mathcal{F}$ by defining $(B, \varphi) \le (C, \psi)$ if $B \subset C$ and
$\psi|_B = \varphi$. Let us show that $\mathcal{F}$ has a maximal element. Suppose that
$\{(B_i, \varphi_i)\}$ is a totally ordered subset of $\mathcal{F}$, and set
\[
  B_0 = \bigcup_i B_i, \quad \varphi_0 = \bigcup_i \varphi_i;
\]
then one sees easily that $B_0$ is a local ring with maximal ideal $tB_0$. If
$0 \ne x \in B_0$, then $x \in B_i$ for some $i$, and since $B_i$ is a DVR we can write
$x = t^\nu u$ with $u$ a unit of $B_i$ and some $\nu$. From this we get
$x \notin t^{\nu+1}B_0$, so that $B_0$ is $t$-adically separated. Hence $B_0$ is a DVR, and if
$\varphi_0\colon B_0 \to K$ is defined to be equal to $\varphi_i$ on $B_i$ then
$(B_0, \varphi_0) \in \mathcal{F}$. Hence by Zorn's lemma $\mathcal{F}$ has a maximal element;
suppose that $(B, \varphi)$ is one. Then if $\varphi(B) \ne K$, take $\alpha \in K - \varphi(B)$,
let $\bar{f}(X)$ be the minimal polynomial of $\alpha$ over $\varphi(B)$, and take a manic
$f(X) \in B[X]$ which is an inverse image of $\bar{f}(X)$. Then $f(X)$ is irreducible in
$B[X]$, and hence (by Ex.~\ref{ex:9.6-ex}) also irreducible over the field of fractions of
$B$. Let $\xi$ be a root of $f(X)$ in $L$, and set $B' = B[\xi]$; then
$B' = B[X]/(f)$. Therefore
\[
  B'/tB' = B[X]/(t, f) = \varphi(B)[X]/(\bar{f}) = \varphi(B)(\alpha)
\]
is a field; since $B'$ is integral over $B$, every maximal ideal of $B'$ contains
$tB'$, so that $B'$ is a local integral domain with maximal ideal $tB'$, and is
Noetherian because it is finite over $B$, therefore a DVR. This contradicts the
maximality of $B$, so that we must have $\varphi(B) = K$.
\end{proof}

\begin{remark}
Since $B$ is an integral domain containing $A$, it is flat over $A$ by
Ex.~\ref{ex:10.2-ex}. In EGA $0_{IV}$ (10.3.1), the following more general fact is proved:
let $(A, m, k)$ be a Noetherian local ring, and $K$ an extension field of $k$, then
there exists a Noetherian ring $B$ containing $A$ satisfying the three conditions
(1) $\rad(B) = mB$, (2) $B/mB$ is isomorphic over $k$ to $K$, and (3) $B$ is flat over $A$.
\end{remark}

\begin{theorem}\label{thm:29.2}
Let $(A, m, K)$ be a complete local ring, $(R, pR, k)$ a $p$-ring, and
$\varphi_0\colon k \to K$ a field homomorphism; then there exists a local homomorphism
$\varphi\colon R \to A$ which induces $\varphi_0$ on the residue fields.
\end{theorem}

\begin{proof}
Set $S = \mathbb{Z}_{(p)}$, and let $k_0 \subset k$ be the prime subfield. Since
$\varphi_0(k_0) \subset K$, the prime number $p$, viewed as an element of $A$, belongs to
$m$. Hence the standard homomorphism $\mathbb{Z} \to A$ extends to a local
homomorphism $S \to A$. Now $R \otimes_S k_0 = R/pR = k$ is a separable extension of
$k_0$, and hence 0-smooth over $k_0$; also $R$ is a torsion-free $S$-module, hence
flat over $S$, so that by Theorem~\hyperref[thm:28.10]{28.10}, $R$ is $pR$-smooth over $S$.
\[
\begin{tikzcd}
k_0 \arrow[r, "\subset"] \arrow[d, "\subset"'] & k \arrow[r, "\varphi_0"] & K = A/m \\
S \arrow[r, "\subset"] & R \arrow[ru, "\varphi"', dashed] & A \arrow[u, "p"] \\
& & \mathbb{Z} \arrow[u, "\subset"] \arrow[uu, "\subset"]
\end{tikzcd}
\]
Hence, as we discussed at the beginning of \S\ref{sec:28}, we can lift
$\varphi_0\colon k \to A/m = K$ successively to $R \to A/m^\nu$, and using the fact that
$A = \varprojlim A/m^\nu$, we get $\varphi\colon R \to A$ making the left-hand diagram
commute.
\end{proof}

\begin{theorem}\label{thm:29.3}
The complete $p$-ring is uniquely determined up to isomorphism by its residue
field.
\end{theorem}

\begin{proof}
Suppose that $R$ and $R'$ are both complete $p$-rings with residue field $k$; then by
the theorem there exists a local homomorphism $\varphi\colon R \to R'$ which induces
the identity map on the residue field. We have $R' = \varphi(R) + pR'$, and of course
$\varphi(p) = p$, so that by the completeness of $R$ we see that $\varphi$ is
surjective, and is also injective, since $p^\nu R$ is not contained in $\Ker \varphi$
for any $\nu$. Therefore $R \cong R'$.
\end{proof}

Let $(A, m, k)$ be a complete local ring of unequal characteristic, and let
$p = \chark k$. We say that a subring $A_0 \subset A$ is a \emph{coefficient ring}\index{coefficient ring}
of $A$ if $A_0$ is a complete Noetherian local ring with maximal ideal $pA_0$, and
\[
  A = A_0 + m, \quad \text{that is, } k = A/m \cong A_0/pA_0.
\]
By Theorem~\hyperref[thm:29.1]{1} applied to the residue field $k$ of $A$, there exists a $p$-ring
$S$ such that $S/pS = k$; write $R$ for the completion of $S$, so that $R$ is a
complete $p$-ring with residue field $k$. By Theorem~\hyperref[thm:29.2]{2}, there exists a local
homomorphism $\varphi\colon R \to A$ inducing an isomorphism on the residue fields. If
we set $\varphi(R) = A_0$, then this is clearly a coefficient ring of $A$. If $A$ has
characteristic $0$ then $\varphi$ is injective and $A_0 \cong R$. If $A$ has
characteristic $p^\nu$ then $A_0 \cong R/p^\nu R$. We summarise the above discussion in
the following theorem.

\begin{theorem}\label{thm:29.4}
If $(A, m, k)$ is a complete local ring and $p = \chark k$ then $A$ has a coefficient
ring $A_0$. If $A$ has characteristic $0$ then $A_0$ is a complete DVR, and if $A$ has
characteristic $p^\nu$ then $A_0$ is an Artinian local ring.
\end{theorem}

In what follows, in order to include the case of equal characteristic in our
discussion, we also consider coefficient fields as being coefficient rings.
From the previous theorem and Theorem~\hyperref[thm:28.3]{28.3}, we get the following
important result.

\begin{theorem}\label{thm:29.5}
\begin{enumerate}[(i)]
  \item If $(A, m)$ is a complete local ring and $m$ is finitely generated, then $A$
        is Noetherian.
  \item A Noetherian complete local ring is a quotient of a regular local ring;
        in particular it is universally catenary.
  \item If $A$ is a Noetherian complete local ring (and in the case of unequal
        characteristic, $A$ is an integral domain), then there exists a subring
        $A' \subset A$ with the following properties: $A'$ is a complete regular local
        ring with the same residue field as $A$, and $A$ is finitely generated as an
        $A'$-module.
\end{enumerate}
\end{theorem}

\begin{proof}
We choose a coefficient ring $A_0$ of $A$. If $m = (x_1, \ldots, x_n)$ then every element
of $A$ can be expanded as a power series in $x_1, \ldots, x_n$ with coefficients in
$A_0$, so that $A$ is a quotient of $A_0[[X_1, \ldots, X_n]]$ and hence Noetherian. Now
$A_0$ is a quotient of a $p$-ring $R$, so that $A$ is a quotient of
$R[[X_1, \ldots, X_n]]$, which is a regular local ring, hence a CM ring, and therefore
according to Theorem~\hyperref[thm:17.9]{17.9}, $A$ is universally catenary. To prove (iii), set
$\dim A = n$, and in the case of equal characteristic let $\{y_1, \ldots, y_n\}$ be any
system of parameters of $A$; if $A$ is an integral domain of characteristic $0$ and
$\chark k = p$, we can choose a system of parameters $\{y_1 = p, y_2, \ldots, y_n\}$ of $A$
starting with $p$. In either case $R = A_0$, so that we set $A' = R[y]$. Then $A'$ is the
image of $\varphi\colon R[Y] \to A$, where $R[Y]$ is the regular local ring
\[
  R[Y] = R[[Y_1, \ldots, Y_n]], \quad \text{or } R[[Y_2, \ldots, Y_n]] \text{ if } y_1 = p,
\]
and $\varphi$ is the $R$-algebra homomorphism defined by $\varphi(Y_i) = y_i$. Set
$m' = \sum y_iA'$. Since $A/m = A'/m'$, every $A$-module of finite length has the same
length when viewed as an $A'$-module. In particular $A/m'A$ is a finite module over
$A'/m'$, and $A$ is $m'$-adically separated, so that by Theorem~\hyperref[thm:8.4]{8.4} $A$ is a finite
$A'$-module. Therefore, we have
\[
  \dim A' = \dim A = n.
\]
$R[Y]$ is an $n$-dimensional integral domain, and if $\Ker \varphi \ne 0$ we would have
$\dim A' < n$, which is a contradiction. Therefore $\varphi$ is injective and
$A' \cong R[Y]$.
\end{proof}

\begin{remark}
In the case of unequal characteristic when $A$ is not an integral domain, (iii)
can fail to hold. If $A$ is of characteristic $p^m$ with $m > 1$, then every subring
of $A$ has the same characteristic $p^m$, so cannot be regular. Even if $A$ has
characteristic $0$, the following is a counter-example: let $R$ be a complete
$p$-ring and $A = R[X]/(pX)$; then $A$ is a complete one-dimensional Noetherian
local ring, but if a subring $A'$ as in (iii) were to exist, $A'$ would be a
one-dimensional regular local ring, hence a DVR, and since $A'$ has characteristic
$0$ and its residue field characteristic $p$, $A'/pA'$ would be an Artinian ring,
and hence also $A/pA$ would be Artinian. But $A/pA \cong k[X]$ is
one-dimensional, and this is a contradiction.

The proof of Theorem~\hyperref[thm:29.5]{5} shows that it is sufficient to assume that $p$ is not
in any minimal prime ideal of $A$.
\end{remark}

Corresponding to the definition of quasi-coefficient field of an
equicharacteristic local ring, let us define \emph{quasi-coefficient rings} in the case of
unequal characteristic. Let $(A, m, K)$ be a possibly non-complete local ring, and
suppose $\chark K = p$. A subring $S \subset A$ is said to be a quasi-coefficient ring if
it satisfies the following two conditions:
\begin{enumerate}[(1)]
  \item $S$ is a Noetherian local ring with maximal ideal $pS$;
  \item $K = A/m$ is 0-\'etale over $S/pS$.
\end{enumerate}
In view of (1) if $A$ has characteristic $0$ then $S$ is a DVR, and if $A$ has
characteristic $p^\nu$ then $S$ is an Artinian ring.

\begin{theorem}\label{thm:29.6}
Let $(A, m, K)$ be a local ring, and suppose $\chark K = p$. Let $C \subset A$ be a
subring, and assume that $C$ is a Noetherian local ring with maximal ideal $pC$,
and that $K = A/m$ is separable over $C/pC$. Then there exists a quasi-coefficient
ring $S$ of $A$ containing $C$; moreover, if $A$ is flat over $C$, then it is also
flat over $S$.
\end{theorem}

\begin{proof}
Let $\{B_\lambda\}_{\lambda \in \Lambda}$ be a $p$-basis of $K$ over $C/pC$, and choose an inverse
image $b_\lambda \in A$ for each $\beta_\lambda$. Setting $C[\{b_\lambda\}] = C'$, by
Theorem~\hyperref[thm:26.8]{26.8} we see that $C'/(m \cap C') = (C/pC)[\{\bar{X}_\lambda\}]$ is a
polynomial ring over $C/pC$. Hence if $f(\ldots, X_\lambda, \ldots)$ is a non-zero
polynomial with coefficients in $C$ which satisfies $f(b_\lambda) \in m$, then setting
$\rho$ for the highest common factor of the coefficients of $f$, we have
\[
  f(X) = \rho \cdot f_1(X) \quad \text{and} \quad f_1(b_\lambda) \not\equiv 0.
\]
Thus $f_1(b_\lambda) \notin m$, and we must have $\rho \in m$, so that we have shown that
$m \cap C' = pC'$. Setting $S = (C')_{m}$, we have $S \subset A$, $m \cap S = pS$ and
$S/pS = K'$; since $C'$ is $p$-adically separated, so is $S$, and hence all the
ideals of $S$ are of the form $(0)$ or $(p^\nu)$. Thus $S$ is Noetherian, and it
satisfies all the conditions for a quasi-coefficient ring of $A$. If $A$ is flat over
$C$, then
\[
  p^\nu C \otimes_C A \cong p^\nu A.
\]
and hence the composite
\[
  p^\nu C \otimes A = (p^\nu C \otimes_C S) \otimes_S A \to p^\nu S \otimes_S A \to p^\nu A
\]
is injective; but the first arrow is surjective, so that the second arrow
$p^\nu S \otimes_S A \to p^\nu A$ is injective. By Theorem~\hyperref[thm:7.7]{7.7}, this proves that $A$ is
flat over $S$.
\end{proof}

All the assumptions of this theorem are satisfied by taking $C$ to be the image
of $\mathbb{Z}_{(p)} \to A$, so that this proves that every local ring has a quasi-coefficient
ring (including the quasi-coefficient field of a local ring in the equal
characteristic case).

\begin{theorem}\label{thm:29.7}
Let $(A, m, K)$ be a local ring, and $\hat{A}$ its completion, and suppose
$\chark K = p$. Let $S$ be a quasi-coefficient ring of $A$, and write $\bar{S}$ for its
image in $\hat{A}$; then there exists a unique coefficient ring $A_0$ of $\hat{A}$
containing $\bar{S}$.
\end{theorem}

\begin{proof}
Since $\bar{S}$ is a quasi-coefficient ring of $\hat{A}$, we can assume that $A$ is a
complete local ring. If $A$ has characteristic $0$ then $S$ is a DVR, and by
Theorem~\hyperref[thm:29.1]{1} there exists a complete $p$-ring $R$ containing $S$ and with
residue field $K$. Now $K$ is 0-\'etale over $S/pS$, and $R$ is flat over $S$, so that
$R$ is $pR$-\'etale over $S$ (by Ex.~\ref{ex:28.1-ex}). Hence there exists a unique
$S$-algebra homomorphism $R \to A$ which induces the identity map on the residue
fields; we write $A_0$ for the image of $R$. Then $R \cong A_0$, and $A_0$ is a
coefficient ring of $A$. If $A$ had another coefficient ring $B$ containing $S$ then
$B$ would also be a complete $p$-ring, and for the same reason as above, there
exist unique $S$-algebra homomorphisms $B \to A_0$ and $B \to A_0$, so that we must
have $B = A_0$.

If $A$ has characteristic $p^\nu$ then $S$ is an Artinian ring, and therefore
complete, so that applying Theorem~\hyperref[thm:29.4]{4} to $S$ itself, we can write
$S = R_0/p^\nu R_0$ with $R_0$ a complete $p$-ring. Let $R$ be a complete $p$-ring
containing $R_0$ and with residue field $K$; then by Ex.~\ref{ex:28.1-ex}, $R$ is $pR$-\'etale
over $R_0$, so that there exists a unique $R_0$-algebra homomorphism $R \to A$
inducing the identity map on the residue field $K$. The image is a coefficient
ring of $A$ containing $S$; the uniqueness is proved as in the case of
characteristic $0$.
\end{proof}

Next we study the structure of complete regular local rings. Let $(A, m, K)$ be a
local ring of unequal characteristic, and suppose that $\chark K = p$; then $A$ is
said to be \emph{ramified}\index{ramified} if $p \in m^2$ and \emph{unramified} if $p \notin m^2$. We will
also say that $A$ is unramified in the case of equal characteristic.

\begin{theorem}\label{thm:29.8}
An unramified complete regular local ring is a formal power series ring over a
field or over a complete $p$-ring.
\end{theorem}

\begin{proof}
Let $R$ be a coefficient ring of $A$. In the case of equal characteristic, $R$ is a
field, and if $x_1, \ldots, x_n$ is a regular system of parameters of $A$ then
$A = R[x_1, \ldots, x_n] \cong R[[X_1, \ldots, X_n]]$ (see the proof of
Theorem~\hyperref[thm:29.5]{5}). In the case of unequal characteristic, $R$ is a complete $p$-ring,
and since $p \notin m^2$, we can choose a regular system of parameters
$\{p, x_2, \ldots, x_n\}$ of $A$ containing $p$. Then $A = R[x_2, \ldots, x_n] \cong
R[[X_2, \ldots, X_n]]$.
\end{proof}

In the ramified case, it is not necessarily the case that $A$ can be expressed as
a formal power series ring over a DVR. To give the structure theorem in this
case we need the notion of an \emph{Eisenstein extension}.

\begin{lemma}[\textbf{Eisenstein's irreducibility criterion}]\label{lem:29.1}
Let $A$ be a ring, and $f(X) = X^n + a_1X^{n-1} + \cdots + a_n$, with $a_i \in A$. If there
exists a prime ideal $P$ of $A$ such that $a_1, \ldots, a_n \in P$ but $a_n \notin P^2$ then $f$
is irreducible in $A[X]$. If in addition $A$ is an integrally closed domain then
the principal ideal $(f)$ is a prime ideal of $A[X]$.
\end{lemma}

\begin{proof}
If $f$ is reducible, we can write
\[
  f = (X^r + b_1X^{r-1} + \cdots + b_r)(X^s + c_1X^{s-1} + \cdots + c_s)
\]
with $0 < r < n$, $s = n - r$ and $b_i, c_j \in A$. Reducing the coefficients on either
side modulo $P$ we have
\[
  X^n = (X^r + \bar{b}_1X^{r-1} + \cdots + \bar{b}_r)(X^s + \bar{c}_1X^{s-1} + \cdots + \bar{c}_s)
\]
in $(A/P)[X]$, so that we must have $b_i, c_j \in P$ for all $i, j$, but then
$a_n = b_rc_s \in P^2$, which contradicts the assumption. If $A$ is an integrally closed
domain and $K$ is the field of fractions of $A$, then by Ex.~\ref{ex:9.6-ex}, $f$ remains
irreducible in $K[X]$. Also, $f$ is manic, so that we have
$f \cdot A[X] = f \cdot K[X] \cap A[X]$, and this is a prime ideal of $A[X]$.
\end{proof}

Let $(A, m)$ be a normal local ring; then an extension ring
\[
  B = A[x]/(f) = A[x]
\]
defined by an \emph{Eisenstein polynomial}\index{Eisenstein polynomial}
\[
  f = X^n + a_1X^{n-1} + \cdots + a_n \quad \text{with } a_i \in m \text{ for all } i,
      \text{ and } a_n \notin m^2
\]
is called an \emph{Eisenstein extension} of $A$. By the lemma, $B$ is an integral domain,
and is integral over $A$. We have $B/mB = (A/m)[X]/(X^n)$, so that $B$ has just one
maximal ideal $n = mB + xB$. Hence $B$ is a local ring, and its residue field
coincides with that of $A$.

\begin{theorem}\label{thm:29.9}
\begin{enumerate}[(i)]
  \item If $(A, m)$ is a regular local ring, then an Eisenstein extension of $A$ is
        again a regular local ring.
  \item If $A$ is a ramified complete regular local ring and $R$ is a coefficient
        ring of $A$ then there is a subring $A_1 \subset A$ with the following
        properties:
        \begin{enumerate}[(1)]
          \item $A_1$ is an unramified complete regular local ring containing $R$,
                and hence can be expressed as a formal power series ring over $R$;
          \item $A$ is an Eisenstein extension of $A_1$.
        \end{enumerate}
\end{enumerate}
\end{theorem}

\begin{proof}
(i) Let $B = A[x]$, and $x^n + a_1x^{n-1} + \cdots + a_n = 0$, with $a_i \in m$ and
$a_n \in m^2$ but $a_n \notin m^3$. Then there exists a regular system of parameters
$\{y_1, \ldots, y_n = a_n\}$ of $A$ with $a_n$ as an element. As we have seen above, the
maximal ideal of $B$ is $mB + xB$, but $a_n \in xB$, so that
$\{y_1, \ldots, y_{n-1}, x\}$ is a regular system of parameters of $B$.

(ii) Since $\htop pA = 1$, by a skilful choice of a regular system of parameters
$\{x_1, \ldots, x_d\}$ of $A$, we can arrange that $\{p, x_2, \ldots, x_d\}$ is a system of
parameters of $A$. If we set $A_1 = R[x_2, \ldots, x_d]$ then $A_1$ is a complete
unramified regular local ring, and $A$ is a finite module over $A_1$ (see the
proof of Theorem~\hyperref[thm:29.5]{5}). We set $m_1$ for the maximal ideal of $A_1$. Now
$A = m_1A + A_1[x_1]$, so that by Theorem~\hyperref[thm:8.4]{8.4} (or by NAK),
$A = A_1[x_1]$. Let
\[
  f(X) = X^n + a_1X^{n-1} + \cdots + a_n \quad \text{with } a_i \in A_1
\]
be the minimal polynomial of $x_1$ over $A_1$. Then $a_n = x_1 \cdot \ldots \in m$,
so that $a_n \notin m_1^2$. Therefore by Hensel's lemma (Theorem~\hyperref[thm:8.3]{8.3}), all the
$a_i \in m_1$. We are left to prove that $a_n \notin m_1^2$. Write $p = \sum_{i=1}^{d} b_ix_i$ with
$b_i \in A_1$, and express the $b_i$ in the form $b_i = q_i(x_1)$, with
$q_i(X) \in A_1[X]$; then $x_1$ is a root of
\[
  F(X) = q_1(X)X + \sum_{i=2}^{d} q_i(X)x_i - p
\]
so that $F(X)$ is divisible by $f(X)$. Hence the constant term $F(0)$ of $F$ is
divisible by $a_n$. However, $F(0) = \sum_{i=2}^{d} q_i(0)x_i - p$, and
$p, x_2, \ldots, x_d$ is a regular system of parameters, so that $F(0) \notin m^2$, hence
also $a_n \notin m_1^2$.
\end{proof}

% ============ 习题 ============
\section*{Exercises to \S\ref{sec:29}}
\emph{Prove the following propositions.}

\begin{exercise}\label{ex:29.1-ex}
Let $A$ be a complete $p$-ring, $y$ an indeterminate over $A$, and $B = A[[y]]$; let
$C = B[x]$ be the Eisenstein extension of $B$ given by $x^2 + yx + p = 0$. Then $C$
is a two-dimensional complete regular local ring, but is not a formal power
series ring over a DVR of characteristic $0$.
\end{exercise}

\begin{exercise}\label{ex:29.2-ex}
In Theorem~\hyperref[thm:29.2]{29.2}, if $k$ is a perfect field then $\varphi$ is uniquely determined
by $\varphi_0$.
\end{exercise}

% ============ §30 Connections with derivations (原书 pp.230-245) ============
\section{Connections with derivations}\label{sec:30}\index{Jacobian criterion, section 30}\index{weak Jacobian condition, section 30}\index{Nagata-Zariski-Lipman, section 30}\index{geometrically regular, section 30}

\begin{theorem}[\textbf{Nagata-Zariski-Lipman}]\label{thm:30.1}
Let $(A, m)$ be a complete Noetherian local ring with $\mathbb{Q} \subset A$. Suppose that
$x_1, \ldots, x_r \in m$ and $D_1, \ldots, D_r \in \Der(A)$ are elements satisfying
$\det(D_ix_j) \notin m$. Then
\begin{enumerate}[(i)]
  \item There is a subring $C \subset A$ such that
        \[
          A = C[x_1, \ldots, x_r] \cong C[[X_1, \ldots, X_r]].
        \]
        Therefore $x_1, \ldots, x_r$ are analytically independent over $C$, and $A$ is
        $I$-smooth over $C$, where $I = \sum Ax_i$, and therefore also $m$-smooth over
        $C$.
  \item If $\mathfrak{g} = \sum AD_i$ is a Lie algebra, (that is if $[D_i, D_j] \in \mathfrak{g}$ for all
        $i, j$) then we can take $C$ to be $\{a \in A \mid D_1a = \cdots = D_ra = 0\}$.
\end{enumerate}
\end{theorem}

\begin{proof}
Letting $(c_{ij})$ be the inverse matrix of $(D_ix_j)$, and setting
$D_i' = \sum c_{ij}D_j$, we have $D_i'x_j = \delta_{ij}$, so that we can assume that
$D_ix_j = \delta_{ij}$. Quite generally, for an element $t \in m$ and a derivation
$D \in \Der(A)$, we define a map $E(D, t)\colon A \to A$ by
\[
  E(D, t) = \sum_{n \ge 0} \frac{t^n}{n!}D^n;
\]
by our assumptions, $E(D, t)(a) = \sum (t^n/n!)D^n(a)$ is meaningful, and one sees
easily that $E(D, t)$ is a ring homomorphism. Now set
\[
  E_i = E(D_i, -x_i) \quad \text{and} \quad C_i = \Im(E_i);
\]
then $C_i$ is a subring of $A$, and by computation we see that
\[
  D_i\Big(\sum \frac{(-x_i)^n}{n!}D_i^n(a)\Big) = \sum \frac{(-x_i)^{n-1}}{(n-1)!}D_i^n(a)
      + \sum \frac{(-x_i)^n}{n!}D_i^{n+1}(a) = 0,
\]
so that $C_i \subset \{a \in A \mid D_ia = 0\}$. Conversely, if $D_ia = 0$ then
$E_i(a) = a$, so that $C_i = \{a \in A \mid D_ia = 0\}$. Also, since for any $a \in A$ we have
$E_i(a) \equiv a \mod x_1A$, we see that elements of $A$ can be expanded in power
series in $x_1$ with coefficients in $C_1$, so that $A = C_1[x_1]$. Now
$E_1(x_1) = x_1 - x_1 = 0$, so that $x_1A \subset \Ker E_1$, and conversely, if
$E_1(a) = a - x_1D_1a + \cdots = 0$ then $a \in x_1A$, and therefore $\Ker E_1 = x_1A$.
Also, $D_1 \circ D_1c = 0 \iff E_1(c) = c$, so that $C_1 \cap x_1A = 0$. Now we prove
that $x_1$ is analytically independent over $C_1$; by contradiction, suppose that
\[
  c_\nu x_1^\nu + c_{\nu+1}x_1^{\nu+1} + \cdots = 0 \quad \text{with } c_\nu \in C_1
      \text{ and } c_\nu \ne 0.
\]
Then by Ex.~\hyperref[ex:25.2-ex]{25.2}, since $x_1$ is not a zero-divisor in $A$, we have
$c_\nu \in x_1A$, which is a contradiction. Thus if $0 \ne \varphi(X) \in C_1[X]$ then
$\varphi(x_1) \ne 0$, as required.

If $r > 1$, then write $D_i'$ for the restriction to $C_1$ of $E_1 \circ D_i$; then
$D_i' \in \Der(C_1)$, and $x_j \in C_1$ with $D_i'x_j = \delta_{ij}$ for $2 \le i, j \le r$, so that by
induction we have
\[
  C_1 = C[x_2, \ldots, x_r] \cong C[[X_2, \ldots, X_r]].
\]
(i) follows from this.

If $\mathfrak{g}$ is a Lie algebra, then we first arrange as before that $D_ix_j = \delta_{ij}$,
and then set $[D_i, D_j] = \sum_\nu a_{ij\nu}D_\nu$ with $a_{ij\nu} \in A$; then
$[D_i, D_j]x_\nu = D_i(\delta_{j\nu}) - D_j(\delta_{i\nu}) = 0$, so that $a_{ij\nu} = 0$, hence
$[D_i, D_j] = 0$, and $D_1(D_i(C_1)) = D_i(D_1(C_1)) = D_i(0) = 0$. Therefore
$D_i(C_1) \subset C_1$ for $i > 1$, and then in the above notation $D_i' = D_i$ for $i > 1$.
Thus by induction we have $C = \{a \in A \mid D_1a = \cdots = D_ra = 0\}$.
\end{proof}

\begin{corollary*}
Let $(A, m)$ be a reduced $n$-dimensional local ring containing $\mathbb{Q}$, and suppose
that the completion $\hat{A}$ of $A$ is also reduced. If there exist $n$ elements
$D_1, \ldots, D_n \in \Der(A)$ and $x_1, \ldots, x_n \in m$ such that
$\det(D_ix_j) \notin m$ then $A$ is a regular local ring and $x_1, \ldots, x_n$ is a regular
system of parameters of $A$. Suppose in addition that $\mathfrak{g} = \sum AD_i$ is a Lie
algebra; then $C = \{a \in A \mid D_1a = \cdots = D_na = 0\}$ is a coefficient field of
$\hat{A}$.
\end{corollary*}

\begin{proof}
Consider $\hat{A}$, with each of the $D_i$ extended to $\hat{A}$. If
$[D_i, D_j] = \sum_\nu a_{ij\nu}D_\nu$ holds in $A$ then it also holds in $\hat{A}$, so that if
$\mathfrak{g}$ is a Lie algebra, so is $\hat{\mathfrak{g}} = \sum \hat{A}D_i$. By the theorem,
$\hat{A} = C[[x_1, \ldots, x_n]]$, and $C$ is isomorphic to $\hat{A}/\sum x_i\hat{A}$, so is a
zero-dimensional local ring. Now by assumption $C$ is also reduced, so that $C$
must be a field. Therefore $\hat{A}$ is a formal power series ring over a field, and
hence is regular, so that $A$ is also regular. The other assertion is also clear.
\end{proof}

\begin{remark}
If we view this corollary as a criterion for regularity, then the condition that
$\hat{A}$ should be reduced is rather a nuisance; however, as we will see later, for
a very wide variety of local rings, we have $A$ is reduced $\iff \hat{A}$ is reduced.
This is the case (corollary of Theorem~\hyperref[thm:32.6]{32.6}) if $A$ is a localisation of a ring
$B$ which is finitely generated over a field (such an $A$ is said to be
\emph{essentially of finite type} over $K$). Note also that if we start off with a
regular local ring $A$, then the corollary gives a concrete method of
constructing a coefficient field of $\hat{A}$.
\end{remark}

Next we consider rings which are finitely generated over a field, which are
important in algebraic geometry.

\begin{theorem}\label{thm:30.2}
Let $k$ be a field, and $A = k[x_1, \ldots, x_n]$ a finitely generated ring over $k$. If
$A_p$ is 0-smooth over $k$ for every $p \in \mathrm{m}$-$\Spec A$, then $A$ is 0-smooth over
$k$.
\end{theorem}

\begin{proof}
Write $k[X]$ for $k[X_1, \ldots, X_n]$, and let $I = \{f(X) \in k[X] \mid f(x) = 0\}$, so
that $A = k[X]/I$. Suppose that $I = (f_1, \ldots, f_s)$. Consider a commutative diagram
\[
\begin{tikzcd}
A \arrow[r, "\bar{\varphi}"] & C/N \\
k \arrow[u] \arrow[r] & C \arrow[u]
\end{tikzcd}
\]
where $C$ is a ring, and $N \subset C$ is an ideal satisfying $N^2 = 0$. To lift
$\bar{\varphi}$ to $A \to C$, we first of all choose $u_i \in C$ such that
$\bar{\varphi}(x_i) = \bar{u}_i$. If $f \in I$ then $f(u_1, \ldots, u_n) \in N$. Now if we can choose
$y_i \in N$ for $1 \le i \le n$ such that $f_j(u + y) = 0$ for all $j$, then the
homomorphism $A \to C$ defined by $X_i \mapsto u_i + y_i$ is a lifting of
$\bar{\varphi}$. We have
\[
  f_j(u + y) = f_j(u) + \sum_{i=1}^{n} (\partial f_j/\partial X_i)(u) \cdot y_i,
\]
so that we are looking for solutions in $N$ of the system of linear equations in
$y_1, \ldots, y_n$:
\[
  (*) \quad f_j(u) + \sum_{i=1}^{n} (\partial f_j/\partial X_i)(u) \cdot y_i = 0
      \quad \text{for } j = 1, \ldots, s.
\]
For each maximal ideal $p$, the local ring $A_p$ is 0-smooth over $k$, so that if we
set $S = \bar{\varphi}(A - p)$ and $\tilde{S} = \varphi^{-1}(S)$ then in the diagram
\[
\begin{tikzcd}
A_p \arrow[r, "\bar{\varphi}_p"] & C_{\tilde{S}}/N_{\tilde{S}} = C_S/N_S \\
k \arrow[u] \arrow[r] & C_{\tilde{S}} \arrow[u]
\end{tikzcd}
\]
there exists a $y^p\colon A_p \to C_{\tilde{S}}$ lifting $\bar{\varphi}_p$. From this we see that (*), as
a system of equations in $N_{\tilde{S}}$, has a solution in $N_{\tilde{S}}$. If we view $N$ as a
$C/N$-module then $N_{\tilde{S}} = N_S$. Thus the theorem reduces to the following
lemma.
\end{proof}

\begin{lemma}\label{lem:30.1}
Let $A$ and $B$ be rings, $\varphi\colon A \to B$ a ring homomorphism, and $N$ a
$B$-module; suppose that $b_{ij} \in B$ and $\beta_i \in N$. If the system of linear
equations
\[
  \sum_j b_{ij}y_j = \beta_i \quad (\text{for } i = 1, \ldots, s)
\]
has a solution in $N_{A-p}$ for every $p \in \mathrm{m}$-$\Spec A$, then it has a solution in $N$.
\end{lemma}

\begin{proof}
The assumption that there is a solution in $N_{A-p}$ means that there exist
$q_{jp} \in N$ for $1 \le j \le n$, and $t_p \in A - p$ such that
\[
  \sum_j b_{ij}q_{jp} - \varphi(t_p)\beta_i = 0 \quad \text{for } 1 \le i \le s.
\]
Now since $\sum_p t_pA = A$, there is a finite set $\{p_1, \ldots, p_\ell\} \subset
\mathrm{m}$-$\Spec A$ and $a_\nu \in A$ such that
\[
  \sum_{\nu=1}^{\ell} a_\nu t_{p_\nu} = 1.
\]
Hence if we set
\[
  q_j = \sum_\nu \varphi(a_\nu)q_{jp_\nu},
\]
we get $\sum_j b_{ij}q_j = \beta_i$ for $1 \le i \le s$.
\end{proof}

\begin{theorem}\label{thm:30.3}
Let $k$ be a field, $S = k[X_1, \ldots, X_n]$, and $I$, $P$ ideals of $S$ such that
$I \subset P \in \Spec S$. Set
\[
  S_P = R, \quad \rad(R) = PR = M, \quad \text{and} \quad R/IR = A, \quad \rad(A) = m,
\]
\[
  R/M = A/m = K,
\]
and suppose that $\htop IR = r$ and $I = (f_1(X), \ldots, f_t(X))$. Then the following
conditions are equivalent.
\begin{enumerate}[(1)]
  \item $\rank(\partial f_i/\partial X_j \bmod P)_{i=1, \ldots, t;\, j=1, \ldots, n} = r$;
  \item $A$ is 0-smooth over $k$;
  \item $A$ is $m$-smooth over $k$;
  \item $\Omega_{A/k}$ is a free $A$-module of rank $n - r$;
  \item $A$ is an integral domain, its field of fractions is separable over $k$, and
        $\Omega_{A/k}$ is a free $A$-module.
\end{enumerate}
\end{theorem}

\begin{proof}
Note that $R$ is a regular local ring.

(1) $\Rightarrow$ (2) By assumption, for a suitable choice of $r$ elements
$D_1, \ldots, D_r$ from $\partial/\partial X_1, \ldots, \partial/\partial X_n$, and of $r$ elements
$g_1, \ldots, g_r$ from $f_1, \ldots, f_t$, we have $\det(D_ig_j) \notin M$. We observe that
taking $R \ni f$ into $(\partial f/\partial X_1 \bmod M, \ldots, \partial f/\partial X_n \bmod M)
\in K^n$ induces a linear map $M/M^2 \to K^r$, so that the images of $g_1, \ldots, g_r$ in
$M/M^2$ are linearly independent. Therefore $\sum g_iR$ is a height $r$ prime ideal
contained in $IR$, and hence $\sum g_iR = IR$. Given a commutative diagram
\[
\begin{tikzcd}
A \arrow[r, "\bar{\varphi}"] & C/N \\
k \arrow[u] \arrow[r] & C \arrow[u]
\end{tikzcd}
\]
with $N^2 = 0$, write $\bar{X}_i \in A$ for the image of $X_i$, and choose $u_i \in C$
such that $\varphi(u_i) = \bar{\varphi}(\bar{X}_i) \in C/N$. Then there is a homomorphism
$R \to C$ defined by $X_i \mapsto u_i$, and this induces a lifting of $\bar{\varphi}$ to $C$ if and
only if $g_i(u) = 0$ for $1 \le i \le r$. Therefore, as in the proof of the previous
theorem, we need only solve the system of equations in unknowns
$y_1, \ldots, y_n \in N$:
\[
  (*) \quad g_i(u) + \sum_{j=1}^{n} (\partial g_i/\partial X_j)(u) \cdot y_j = 0
      \quad \text{for } 1 \le i \le r.
\]
However, we view $N$ as a $C/N$-module, then via $\bar{\varphi}$ as an $A$-module, so that we
can replace the coefficients $(\partial g_i/\partial X_j)(u)$ of this system by
$(\partial g_i/\partial X_j)(\bar{X}) \in A$. Now by assumption, one $r \times r$ minor of this
$r \times n$ matrix of coefficients is a unit of $A$, so that (*) can always be solved.

(2) $\Rightarrow$ (3) is trivial.

(3) $\Rightarrow$ (1) By \S\ref{sec:28}, Lemma~\hyperref[lem:28.1]{1}, $A$ is regular, so that $IR$ is generated by
$r$ elements forming a subset of a regular system of parameters of $R$, and the
image of the natural map $IR \to M \to M/M^2$ is an $r$-dimensional $K$-vector space.
Let $k_1 \subset k$ be the prime subfield; then by
Theorem~\hyperref[thm:26.9]{26.9} $K = R/M$ is 0-smooth over $k_1$, so that by
Theorem~\hyperref[thm:25.2]{25.2}, the sequence
\[
  0 \to M/M^2 \to \Omega_{R/k_1} \otimes_R K \to \Omega_{K/k_1} \to 0
\]
is exact. We can write $\Omega_{S/k_1} = (\Omega_{k/k_1} \otimes S) \oplus F$, where $F$ is the free
$S$-module with basis $dX_1, \ldots, dX_n$ (for example, by
Theorem~\hyperref[thm:25.1]{25.1}), and localising we get
\[
  \Omega_{R/k_1} = (\Omega_{k/k_1} \otimes R) \oplus F \otimes_S R.
\]
Hence $\Omega_{R/k_1} \otimes_R K = (\Omega_{k/k_1} \otimes_k K) \oplus (F \otimes K)$. However, from
\[
  0 \to I/I^2 \to \Omega_{S/k_1} \otimes_S A \to \Omega_{A/k_1} \to 0
\]
we get the exact sequence $(I/I^2) \otimes_A K \to \Omega_{R/k_1} \otimes_R K \to
\Omega_{A/k_1} \otimes_A K \to 0$, and if $A$ is $m$-smooth over $k$ then by the corollary to
Theorem~\hyperref[thm:28.6]{28.6}, $\Omega_{k/k_1} \otimes K \to \Omega_{A/k_1} \otimes K$ is injective, so that
$V = \Im((I/I^2) \otimes K \to \Omega_{R/k_1} \otimes K)$ maps isomorphically to its
projection $W \subset F \otimes K$ in the second factor of the decomposition
$\Omega_{R/k_1} \otimes K = (\Omega_{k/k_1} \otimes K) \oplus (F \otimes K)$. Now factor
$I/I^2 \otimes K \to \Omega_{R/k_1} \otimes K$ as the composite
$I/I^2 \otimes K \to M/M^2 \to \Omega_{R/k_1} \otimes K$; as we have seen above, the first arrow
has rank $r$, and the second is injective, so that $\rank V = r$. Now
$F \otimes K = KdX_1 + \cdots + KdX_n$, and if we write
$(\partial f_i/\partial X_j) \bmod P = a_{ij}$ then $W$ is spanned by $\sum_j a_{ij}dX_j$ for
$1 \le i \le t$. Therefore $\rank(a_{ij}) = r$, and this proves (1).

(2) $\Rightarrow$ (4) By Theorem~\hyperref[thm:25.2]{25.2},
\[
  0 \to IR/I^2R \to \Omega_{R/k} \otimes_R A \to \Omega_{A/k} \to 0
\]
is a split exact sequence, and since $\Omega_{R/k} \otimes_R A$ is a free $A$-module with basis
$dX_1, \ldots, dX_n$, the $A$-module $\Omega_{A/k}$ is projective; but $A$ is a local ring, so
that $\Omega_{A/k}$ is free. Also $A$ is a regular local ring, therefore an integral
domain, and if $L$ is its field of fractions then $L$ is 0-smooth over $A$, hence also
0-smooth over $k$, so that $L$ is separable over $k$.

(4) $\Rightarrow$ (5) By Theorem~\hyperref[thm:26.2]{26.2}, the field of fractions $L$ of $A$ is
separably generated over $k$, so that a separating transcendence basis of $L$ over
$k$ is a differential basis, and $\rank \Omega_{L/k} = \trdeg_k L = n - r$; but
$\Omega_{L/k} = \Omega_{A/k} \otimes_A L$ and hence $\rank \Omega_{A/k} = \rank \Omega_{L/k} = n - r$.

(5) $\Rightarrow$ (1) In the exact sequence
\[
  IR/I^2R \to \Omega_{R/k} \otimes_R A \to \Omega_{A/k} \to 0,
\]
set $E = \Im(IR/I^2R \to \Omega_{R/k} \otimes_R A)$. Then
$\Omega_{R/k} \otimes_R A \cong E \oplus A^{n-r}$, so that $E \cong A^r$, and therefore
$E \otimes K \cong K^r$. This gives (1).
\end{proof}

\begin{remark}
In the above proof, the equivalence of (1), (2), (4) and (5) was comparatively
easy. The proof of (3) $\Rightarrow$ (1) used the corollary to Theorem~\hyperref[thm:28.6]{28.6}, and so
is not very elementary.
\end{remark}

\begin{remark}
If $k$ is a perfect field, or more generally if the residue field $K$ is separable
over $k$, then $m$-smoothness is equivalent to $A$ being a regular local ring, so
that Theorem~\hyperref[thm:30.3]{3} gives a criterion for regularity. In the case of an imperfect
field $k$, if $A$ is 0-smooth over $k$ then so is the field of fractions $L$ of $A$, but
the residue field $K$ is not necessarily separable over $k$; for example,
$A = k[X]_{(X^p - a)}$ with $a \in k - k^p$. For the case of an imperfect field $k$, we
give a regularity criterion for $A$ in Theorem~\hyperref[thm:30.5]{5}.
\end{remark}

Quite generally, let $A$ be a ring and $P$ a prime ideal of $A$, and let
$D_1, \ldots, D_s \in \Der(A)$ and $f_1, \ldots, f_t \in A$; then we write
\[
  J(f_1, \ldots, f_t; D_1, \ldots, D_s)(P) = (D_if_j \bmod P).
\]
This is an $s \times t$ matrix with entries in the integral domain $A/P$.

\begin{theorem}\label{thm:30.4}
Let $R$ be a regular ring, $P \in \Spec R$, and let $I \subset P$ be an ideal of $R$;
suppose that $\htop IR_P = r$.
\begin{enumerate}[(i)]
  \item for any $D_1, \ldots, D_s \in \Der(R)$ and $f_1, \ldots, f_t \in I$ we have
        \[
          \rank J(f_1, \ldots, f_t; D_1, \ldots, D_s)(P) \le r;
        \]
  \item if $D_1, \ldots, D_r \in \Der(R)$ and $f_1, \ldots, f_r \in I$ are such that
        $\det(D_if_j) \notin P$ then $IR_P = (f_1, \ldots, f_r)R_P$ and $R_P/IR_P$ is regular.
\end{enumerate}
\end{theorem}

\begin{proof}
(i) If $Q$ is a prime ideal of $R$ with $I \subset Q \subset P$ and $\htop Q = r$ then
\[
  \rank J(f_1, \ldots; D_1, \ldots)(P) \le \rank J(f_1, \ldots; D_1, \ldots)(Q),
\]
and if we set $QR_Q = \mathfrak{n}$ then $R_Q$ is an $r$-dimensional regular local ring, so
that $m$ is generated by $r$ elements, $m = (g_1, \ldots, g_r)$. Working in $R_Q$, we can
write
\[
  f_j = \sum_{\nu=1}^{r} g_\nu a_{\nu j} \quad \text{with } a_{\nu j} \in R_Q,
\]
so that
\[
  D_if_j \equiv \sum_{\nu=1}^{r} (D_ig_\nu) \cdot a_{\nu j} \bmod Q,
\]
and therefore
\[
  \rank J(f_1, \ldots, f_t; D_1, \ldots, D_s)(P) \le
  \rank J(g_1, \ldots, g_r; D_1, \ldots, D_s)(Q) \le r.
\]
(ii) Set $M = PR_P$; then if $\det(D_if_j) \notin M$ one sees easily that the images in
$M/M^2$ of $f_1, \ldots, f_r$ are linearly independent over $R_P/M = \kappa(M)$, so that
$\sum f_iR_P$ is a prime ideal of height $r$, and therefore coincides with $IR_P$.
Also, $R_P/IR_P$ is regular.
\end{proof}

\begin{theorem}[\textbf{Zariski}]\label{thm:30.5}
Let $k$ be a field of characteristic $p$, $S = k[X_1, \ldots, X_n]$, and $I$ and $P$
ideals of $S$ with $I \subset P \in \Spec S$. Set $S_P = R$ and $R/IR = A$, and suppose
that $\htop IR = r$ and $I = (f_1, \ldots, f_t)$. Then the following three conditions are
equivalent.
\begin{enumerate}[(1)]
  \item $A$ is a regular local ring.
  \item For any $p$-basis $\{u_\gamma\}_{\gamma \in \Gamma}$ of $k$, define $D_\gamma \in \Der(S)$ by
        $D_\gamma(u_{\gamma'}) = \delta_{\gamma\gamma'}$ (the Kronecker $\delta$) and
        $D_\gamma X_i = 0$; then there are a finite number of elements
        $\alpha, \beta, \ldots, \gamma \in \Gamma$ such that
        \[
          \rank J(f_1, \ldots, f_t; D_\alpha, D_\beta, \ldots, D_\gamma,
              \partial/\partial X_1, \ldots, \partial/\partial X_n)(P) = r.
        \]
  \item There exists a subfield $k' \subset k$ with the following properties:
        $k^p \subset k' \subset k$, $[k : k'] < \infty$, and $\Omega_{A/k'}$ is a free $A$-module, with
        \[
          \rank \Omega_{A/k'} = n - r + \rank \Omega_{k/k'}.
        \]
\end{enumerate}
\end{theorem}

\begin{proof}
(2) $\Rightarrow$ (1) comes from Theorem~\hyperref[thm:30.4]{4}, (ii).

(1) $\Rightarrow$ (2), (3) If $A$ is regular then $IR$ is generated by $r$ elements, and
these form an $R$-sequence, so that $IR/I^2R$ is a free $A$-module of rank $r$. Set
$M = PR$, $m = M/IR$ and $K = R/M = A/m$; then the image of $IR \to M/M^2$ is an
$r$-dimensional $K$-vector space, so that the natural map $(IR/I^2R) \otimes_A K \to
M/M^2$ is injective. From the exact sequence
\[
  IR/I^2R \to \Omega_{R/k} \otimes_R A \to \Omega_{A/k} \to 0
\]
we get the exact sequence
\[
  (*) \quad (IR/I^2R) \otimes_A K \to \Omega_{R/k} \otimes_R K \to \Omega_{A/k} \otimes_A K \to 0.
\]
Now by Theorem~\hyperref[thm:25.2]{25.2},
\[
  0 \to M/M^2 \to \Omega_{R/k} \otimes_R K \to \Omega_{K/k} \to 0
\]
is an exact sequence. The first arrow in (*) is the composite
$(IR/I^2R) \to M/M^2 \to \Omega_{R/k} \otimes K$, and so is injective. Thus
\[
  (**) \quad 0 \to (IR/I^2R) \otimes K \to \Omega_{R/k} \otimes K \to \Omega_{A/k} \otimes K \to 0
\]
is exact. Let $\{du_\sigma\}_{\sigma \in \Sigma}$ be a basis of $\Omega_k$ over $k$; then $\Omega_S$ is the free
$S$-module with basis $\{du_\sigma\}_{\sigma \in \Sigma} \cup \{dX_1, \ldots, dX_n\}$, and
$\Omega_R = \Omega_S \otimes_S R$, $\Omega_R \otimes_R K = \Omega_k \otimes_k K \oplus
(KdX_1 \oplus \cdots \oplus KdX_n)$. Now reorder $f_1, \ldots, f_t$ so that $f_1, \ldots, f_r$ are
generators of $IR$; then $df_1, \ldots, df_r \in \Omega_R \otimes K$ can be expressed using
$dX_1, \ldots, dX_n$ together with finitely many elements $du_{\tau_1}, \ldots, du_{\tau_s}$, and this
gives (2). Now if we let $k'$ be the field obtained by adjoining to $k^p$ all the
elements $u_\sigma$ for $\sigma \in \Sigma$ other than the $du_{\tau_1}, \ldots, du_{\tau_s}$ just used, then
from (**) we see that
\[
  0 \to (IR/I^2R) \otimes K \to \Omega_{R/k'} \otimes K \to \Omega_{A/k'} \otimes K \to 0
\]
is also exact. In the exact sequence
\[
  (\dagger) \quad IR/I^2R \to \Omega_{R/k'} \otimes_R A \to \Omega_{A/k'} \to 0,
\]
the middle term is the free $A$-module with basis
$du_{\tau_1}, \ldots, du_{\tau_s}, dX_1, \ldots, dX_n$; now the generators $f_1, \ldots, f_r$ of $IR$ map
to $df_1, \ldots, df_r \in \Omega_{R/k'} \otimes_R A$, and since these are linearly independent over
$K$, they base a direct summand of $\Omega_{R/k'} \otimes_R A$. Thus $\Omega_{A/k'}$ is also a
free $A$-module, and $\rank \Omega_{A/k'} + r = \rank \Omega_{R/k'} = \rank \Omega_{k/k'} + n$.

(3) $\Rightarrow$ (1) In the exact sequence ($\dagger$) the two terms
$\Omega_{R/k'} \otimes_R A$ and $\Omega_{A/k'}$ are both free modules, and the difference between
their ranks is $r$, so that $IR/I^2R$ maps onto a direct summand of
$\Omega_{R/k'} \otimes_R A$, which is a free module of rank $r$. Thus we can choose
$f_1, \ldots, f_r \in IR$ such that $df_1 \otimes 1, \ldots, df_r \otimes 1$ base this direct summand,
and then by NAK we see that $Rdf_1 + \cdots + Rdf_r$ is a direct summand of
$\Omega_{R/k'}$. Hence there exist $D_1, \ldots, D_r \in \Der_{k'}(R)$ such that
$\det(D_if_j) \notin M$. Therefore, by Theorem~\hyperref[thm:30.4]{4}, $A$ is regular.
\end{proof}

\begin{corollary*}
Let $k$ be a field and $S = k[X_1, \ldots, X_n]$; let $I$ be an ideal of $S$, and set
$B = S/I$. Define $U = \{p \in \Spec B \mid B_p \text{ is 0-smooth over } k\}$ and
$\Reg(B) = \{p \in \Spec B \mid B_p \text{ is regular}\}$. Then both $U$ and $\Reg(B)$ are open
subsets of $\Spec B$.
\end{corollary*}

\begin{proof}
Set $V' = \Spec B$, and let $V_1, \ldots, V_s$ be the irreducible components of $V'$. To
say that $P \in V_i \cap V_j$ for $i \neq j$ means just that $B_p$ has at least two minimal
prime ideals, and such points cannot belong to $\Reg(B)$, (nor a fortiori to $U$).
Thus first of all we can throw out the closed subset $W = \bigcup_{i \neq j}(V_i \cap V_j)$,
and therefore we need only prove that $V_i \cap U$ and $V_i \cap \Reg(B)$ are open in
$V_i$ for each $i$; we fix $i$, and set $\dim V_i = n - r$. Then by
Theorems~\hyperref[thm:30.3]{3} and \hyperref[thm:30.4]{4}, if we let $\Delta_1, \ldots, \Delta_s$ denote the images in
$B$ of the $r \times r$ minors of the Jacobian matrix $(\partial f_i/\partial X_j)$, where
$I = (f_1, \ldots, f_t)$, then $V_i - U$ is the intersection of $V_i$ with the closed subset of
$V'$ defined by the ideal $(\Delta_1, \ldots, \Delta_s)B$ of $B$, and is thus closed in $V_i$.
Using Theorem~\hyperref[thm:30.5]{5}, we can argue similarly for $\Reg(B)$; the only difference
is that, instead of one Jacobian matrix, we have to consider the closed subsets
of $V'$ defined by the ideal of $B$ generated by all the $r \times r$ minors of the
infinitely many Jacobian matrices $J(f_1, \ldots, f_t; D_\alpha, D_\beta, \ldots, D_\gamma,
\partial/\partial X_1, \ldots, \partial/\partial X_n)$, where $\{D_\alpha, \ldots, D_\gamma\}$ runs through all finite
subsets of the set of derivations $\{D_\gamma\}_{\gamma \in \Gamma}$ appearing in
Theorem~\hyperref[thm:30.5]{5}, (2).
\end{proof}

Theorems~\hyperref[thm:30.3]{3} and \hyperref[thm:30.5]{5} contain the result known as the \emph{Jacobian
criterion for regularity} in polynomial rings. There is some purpose in trying to
extend this to more general rings. In cases when the module of differentials is
not finitely generated, then the above method cannot be used as it stands, so we
approach the problem using modules of derivations.

Quite generally, let $A$ be an integral domain with field of fractions $L$; then
for an $A$-module $M$, we write $\rank_A M$ for the dimension over $L$ of the vector
space $M \otimes_A L$.

\begin{theorem}[\textbf{M.~Nomura}]\label{thm:30.6}
Let $(R, m)$ be an equicharacteristic $n$-dimensional regular local ring, and $R^*$
the completion of $R$; suppose that $k$ is a quasi-coefficient field of $R$ and $K$ a
coefficient field of $R^*$ containing $k$. Let $x_1, \ldots, x_n$ be a regular system of
parameters of $R$.
\begin{enumerate}[(i)]
  \item $R^* = K[[x_1, \ldots, x_n]]$, and if we write $\partial/\partial x_i$ for the partial
        derivatives in this representation, then $\Der_k(R^*) = \Der_k(R^*)$ is the free
        $R^*$-module with basis $\partial/\partial x_i$ for $1 \le i \le n$.
  \item The following conditions are equivalent:
        \begin{enumerate}[(1)]
          \item $\partial/\partial x_i$ maps $R$ into $R$ for $1 \le i \le n$, so that they can be
                considered as elements of $\Der_k(R)$;
          \item there exist $D_1, \ldots, D_n \in \Der_k(R)$ and $a_1, \ldots, a_n \in R$ such
                that $D_ia_j = \delta_{ij}$;
          \item there exist $D_1, \ldots, D_n \in \Der_k(R)$ and $a_1, \ldots, a_n \in R$ such
                that $\det(D_ia_j) \notin m$;
          \item $\Der_k(R)$ is a free $R$-module of rank $n$;
          \item $\rank \Der_k(R) = n$.
        \end{enumerate}
\end{enumerate}
\end{theorem}

\begin{proof}
(i) Since $K$ is 0-\'etale over $k$, any derivation of $R^*$ which vanishes on $k$ also
vanishes on $K$. If $D \in \Der_k(R^*)$, set $Dx_i = y_i$; then for any
$f \in R^* = K[[x_1, \ldots, x_n]]$ we have $D(f) = \sum_{i=1}^{n}(\partial f/\partial x_i) \cdot y_i$,
and hence $D = \sum y_i \partial/\partial x_i$. Conversely, for any given $y_i$ we can
construct a derivation by this formula, so that $\Der_k(R^*)$ is the free $R^*$-module
with basis $\partial/\partial x_1, \ldots, \partial/\partial x_n$.

(ii) (1) $\Rightarrow$ (2) $\Rightarrow$ (3) and (4) $\Rightarrow$ (5) are trivial. If (3) holds then
$D_1, \ldots, D_n$ are linearly independent over $R$ and over $R^*$, so that by (i), any
$D \in \Der_k(R)$ can be written as a combination $D = \sum c_iD_i$ of the $D_i$ with
coefficients in the field of fractions of $R^*$, but since $Da_j = \sum c_iD_i(a_j)$ for
$j = 1, \ldots, n$, we have $c_i \in R$; therefore $D_1, \ldots, D_n$ form a basis of
$\Der_k(R)$, which proves (4).

(5) $\Rightarrow$ (1) If $D_1, \ldots, D_n$ are linearly independent over $R$ then there exist
$a_1, \ldots, a_n \in R$ such that $\det(D_ia_j) \ne 0$. Therefore $D_1, \ldots, D_n$ are also
linearly independent over $R^*$. Thus writing $L^*$ for the field of fractions of
$R^*$ we can write $\partial/\partial x_i = \sum_j c_{ij}D_j$, with $c_{ij} \in L^*$. From this we
get $\delta_{ij} = \sum_\nu c_{i\nu}D_\nu x_j$, so that the matrix $(c_{ij})$ is the inverse of
$(D_jx_i)$, and $c_{ij} \in L$, where $L$ is the field of fractions of $R$; therefore
$(\partial/\partial x_i)(R) \subset L \cap R^* = R$.
\end{proof}

\begin{lemma}\label{lem:30.2}
Let $R$ be a regular ring and $P \in \Spec R$ with $\htop P = r$; then the following
two conditions are equivalent:
\begin{enumerate}[(1)]
  \item there exist $D_1, \ldots, D_r \in \Der(R)$ and $f_1, \ldots, f_r \in P$ such that
        $\det(D_if_j) \notin P$;
  \item for all $Q \in \Spec R$ with $\htop Q = s \le r$ such that $Q \subset P$ and
        $R_P/QR_P$ is regular, there exist $D_1, \ldots, D_s \in \Der(R)$ and
        $g_1, \ldots, g_s \in Q$ such that $\det(D_ig_j) \notin P$.
\end{enumerate}
\end{lemma}

\begin{proof}
(2) contains (1) as the special case $P = Q$; conversely, suppose that (1) holds.
By Theorem~\hyperref[thm:30.4]{4} we see that $f_1, \ldots, f_r$ is a regular system of parameters
of $R_P$. If $R_P/QR_P$ is regular then we can take $g_1, \ldots, g_s \in Q$ forming a
minimal basis of $QR_P$, and then take $g_{s+1}, \ldots, g_r \in P$ such that
$g_1, \ldots, g_r$ is a regular system of parameters of $R_P$. Then from
$\det(D_if_j) \notin P$ we deduce that $\det(D_ig_j) \notin P$, or in other words
$\rank J(g_1, \ldots, g_r; D_1, \ldots, D_r)(P) = r$. Therefore
\[
  \rank J(g_1, \ldots, g_s; D_1, \ldots, D_s)(P) = s.
\]
\end{proof}

If the above condition (1) holds, we say that the \emph{weak Jacobian condition} (WJ)
holds at $P$. If (WJ) holds at every $P \in \Spec R$ then we say that (WJ) holds in $R$.
If this holds, then for any $P$, $Q \in \Spec R$ with $Q \subset P$, setting $\htop Q = s$ we
have
\[
  R_P/QR_P \text{ is regular} \iff
  \text{there exist } D_1, \ldots, D_s \in \Der(R) \text{ and }
  f_1, \ldots, f_s \in Q \text{ such that } \det(D_if_j) \notin P,
\]
(the implication ($\Leftarrow$) is given by Theorem~\hyperref[thm:30.4]{4}). This statement is
the Jacobian criterion for regularity. We can use $\Der_k(R)$ in place of
$\Der(R)$, and we then write (WJ)$_k$.

\begin{theorem}\label{thm:30.7}
Let $(A, m)$ be an $n$-dimensional Noetherian local integral domain containing
$\mathbb{Q}$, and let $k \subset A$ be a subfield such that $\trdeg_k(A/m) = r < \infty$. Then
$\Der_k(A)$ is isomorphic to a submodule of $A^{n+r}$, and is therefore a finite
$A$-module, with
\[
  \rank \Der_k(A) \le \dim A + \trdeg_k(A/m).
\]
\end{theorem}

\begin{proof}
We write $A^*$ for the completion of $A$; choose a quasi-coefficient field $k'$ of $A$
containing $k$, and let $K$ be a coefficient field of $A^*$ containing $k'$. Let
$u_1, \ldots, u_r$ be a transcendence basis of $k'$ over $k$, and $x_1, \ldots, x_n$ a system
of parameters of $A$. We define $\varphi\colon \Der_k(A) \to A^{n+r}$ by
$\varphi(D) = (Du_1, \ldots, Du_r, Dx_1, \ldots, Dx_n)$; then $\varphi$ is $A$-linear, so that we
need only prove that it is injective. Suppose then that $Du_i = Dx_j = 0$ for all $i$
and $j$; then $D$ has a continuous extension to $A^*$, and this vanishes on
$B = K[[x_1, \ldots, x_n]]$. Now we do not know whether $A^*$ is an integral domain,
but it is finite as a $B$-module, so that any $a \in A$ is integral over $B$, and if
$f(X) \in B[X]$ has $a$ as a root, and has minimal degree, then $f(a) = 0$,
$f'(a) \ne 0$. Then $0 = D(f(a)) = f'(a) \cdot Da$, and $Da \in A$, so that, since a
non-zero element of $A$ cannot be a zero-divisor in $A^*$, we must have $Da = 0$.
Hence $D = 0$.
\end{proof}

\begin{remark}
If $k$ is an imperfect field then there are counter-examples even if
$k = A/m$: suppose $\chark k = p$ and $a \in k - k^p$, and set
$A = k[X, Y]_{(X, Y)}/(X^p + aY^p)$; then $\dim A = 1$ but
$\rank \Der_k(A) = 2$.
\end{remark}

\begin{theorem}\label{thm:30.8}
Let $(R, m)$ be a regular local ring containing $\mathbb{Q}$ and $k$ a quasi-coefficient
field of $R$. Then the following three conditions are equivalent:
\begin{enumerate}[(1)]
  \item (WJ)$_k$ holds at $m$;
  \item $\rank \Der_k(R) = \dim R$;
  \item (WJ)$_k$ holds at every $P \in \Spec R$.
\end{enumerate}
Furthermore, if these conditions hold, then for any $P \in \Spec R$, every element
of $\Der_k(R/P)$ is induced by an element of $\Der_k(R)$ and
\[
  \rank \Der_k(R/P) = \dim R/P.
\]
\end{theorem}

\begin{proof}
(1) $\Leftrightarrow$ (2) is known from Theorem~\hyperref[thm:30.6]{6}, and (1) is contained in
(3).

(1) $\Rightarrow$ (3) Write $R^*$ for the completion of $R$, and let $K$ be the coefficient
field of $R^*$ containing $k$; then by Theorem~\hyperref[thm:30.6]{6}, if
$x_1, \ldots, x_n$ is a regular system of parameters of $R$ then the derivations
$\partial/\partial x_i$ of $R^* = K[[x_1, \ldots, x_n]]$ belong to $\Der_k(R)$ for $1 \le i \le n$, and
form a basis of it. Now let $P \in \Spec R$, and write $\sigma\colon R \to R/P$ for the
natural homomorphism; to say that $D' \in \Der_k(R/P)$ is induced by
$D \in \Der_k(R)$ means that there is a commutative diagram
\[
\begin{tikzcd}
R \arrow[r, "\sigma"] \arrow[d, "D"'] & R/P \arrow[d, "D'"] \\
R \arrow[r, "\sigma"'] & R/P
\end{tikzcd}
\]
Suppose then that $D'$ is given; then $D'\sigma \in \Der_k(R, R/P)$, and this has a
unique extension to an element of $\Der_k(R^*, R^*/PR^*)$, so that it is uniquely
determined by its values on $x_1, \ldots, x_n$. Therefore if we choose
$b_1, \ldots, b_n \in R$ such that $D'(\sigma(x_i)) = \sigma(b_i)$, and set
$D = \sum b_i \partial/\partial x_i$, then $D'$ is induced by $D$.

Now $\Der_k(R, R/P)$ is a free $R/P$-module with basis $\bar{\sigma} \circ \partial/\partial x_i$ for
$1 \le i \le n$, and $\Der_k(R/P)$ can be identified with the submodule
\[
  N = \{D \in \Der_k(R, R/P) \mid D(f) = 0 \text{ for all } f \in P\}.
\]
Therefore, if $P = (f_1, \ldots, f_t)$ and $\htop P = r$ then
\[
  \rank \Der_k(R/P) = n - \rank J(f_1, \ldots, f_t; \partial/\partial x_1, \ldots,
      \partial/\partial x_n)(P);
\]
according to Theorem~\hyperref[thm:30.4]{4}, the right-hand side is $\ge n - r$, and by
Theorem~\hyperref[thm:30.7]{7} the left-hand side is $\le \dim R/P = n - r$, so we see that
\[
  \rank J(f_1, \ldots, f_t; \partial/\partial x_1, \ldots, \partial/\partial x_n)(P) = r,
\]
\[
  \rank \Der_k(R/P) = \dim R/P.
\]
\end{proof}

This theorem has many applications. For example, in the ring $R$ of convergent
power series in $n$ variables over $k = \mathbb{R}$ or $\mathbb{C}$ ($R$ is denoted by
$k\{X_1, \ldots, X_n\}$ in [N1], and elsewhere by $k\{X_1, \ldots, X_n\}$) we have
$\partial X_i/\partial X_j = \delta_{ij}$, so that the Jacobian criterion for regularity holds in
$R$. In characteristic $0$, the assumptions of the theorem are satisfied by the
formal power series ring over any field, or more generally by the formal power
series ring $R[[Y_1, \ldots, Y_n]]$ over a regular local ring $R$ which satisfies the
assumptions of the theorem. Even if $R$ is not local, but is a regular ring
containing a field $k$ of characteristic $0$, and such that the residue field at
every maximal ideal is algebraic over $k$, then if (WJ)$_k$ holds at every maximal
ideal of $R$, it in fact holds at every prime ideal of $R$, so that the Jacobian
criterion for regularity holds for example in rings such as
$k[X_1, \ldots, X_n] \otimes k[[Y_1, \ldots, Y_m]]$. A further extension of this theorem can be
found in Matsumura [2].

Now we are going to prove the results analogous to Theorem~\hyperref[thm:30.5]{5} for formal
power series rings over a field of characteristic $p$. This is a difficult theorem
obtained by Nagata [6]. First of all we have to do some preparatory work.

Let $k$ be a field and $k' \subset k$ a subfield; we say that $k'$ is \emph{cofinite} in $k$ if
$[k : k'] < \infty$. We say that a family $\mathcal{F} = \{k_\alpha\}_{\alpha \in I}$ of subfields of $k$ is
a \emph{directed family} if for any $\alpha, \beta \in I$ there exists $\gamma \in I$ such that
$k_\alpha, k_\beta \subset k_\gamma$.

Let $K$ be a field of characteristic $p$, and $k \subset K$ a subfield. Then there
exists a directed family $\mathcal{F} = \{k_\alpha\}_{\alpha \in I}$ of intermediate fields
$k \subset k_\alpha \subset K$ cofinite in $K$ such that $\bigcap_\alpha k_\alpha = k(K^p)$. To construct
this, we let $B$ be a fixed $p$-basis of $K$ over $k$, and let $I$ be the set of all
finite subsets of $B$; then we only need take $k_\alpha = k(K^p, B - \alpha)$ for
$\alpha \in I$.

\begin{lemma}\label{lem:30.3}
Let $K$ be a field, $(k_\alpha)_{\alpha \in I}$ a directed family of subfields of $K$, and set
$k = \bigcap_\alpha k_\alpha$. Then if $V$ is a vector space over $K$, and
$v_1, \ldots, v_n \in V$ are linearly independent over $k$, there exists $\alpha \in I$ such that
$v_1, \ldots, v_n$ are also linearly independent over $k_\alpha$.
\end{lemma}

\begin{proof}
For each $\alpha \in I$, write $q(\alpha)$ for the number of linearly independent
elements over $k_\alpha$ among $v_1, \ldots, v_n$; let $\alpha$ be such that $q(\alpha)$ is maximal,
and set $q = q(\alpha)$. Now if $q < n$, and we assume that $v_1, \ldots, v_q$ are linearly
independent over $k_\alpha$, we have $v_n = \sum_{i=1}^{q} c_iv_i$ with $c_i \in k_\alpha$. Since
the $v_i$ are linearly independent over $k$, at least one of the $c_i$ does not belong
to $k$, so that we can assume $c_1 \notin k$. Hence there exists $\beta \in I$ such that
$c_1 \notin k_\beta$, and also $\gamma \in I$ such that $k_\alpha, k_\beta \subset k_\gamma$, so that
$v_1, \ldots, v_q$ and $v_n$ are linearly independent over $k_\gamma$; this contradicts the
maximality of $q$, so that $q = n$.
\end{proof}

\begin{lemma}\label{lem:30.4}
Let $k \subset K$ be fields of characteristic $p$, and let $\mathcal{F} = \{k_\alpha\}_{\alpha \in I}$
be a directed family of intermediate fields $k \subset k_\alpha \subset K$; then the
following conditions are equivalent:
\begin{enumerate}[(1)]
  \item $\bigcap_\alpha k_\alpha(K^p) = k(K^p)$;
  \item the natural map $\Omega_{K/k} \to \varprojlim \Omega_{K/k_\alpha}$ is injective;
  \item if a finite subset $\{u_1, \ldots, u_n\} \subset K$ is $p$-independent over $k$, then
        it is also $p$-independent over $k_\alpha$ for some $\alpha$;
  \item there exists a $p$-basis $B$ of $K$ over $k$ such that every finite subset of
        $B$ is $p$-independent over $k_\alpha$ for some $\alpha$.
\end{enumerate}
\end{lemma}

\begin{proof}
(1) $\Rightarrow$ (3) If $\{u_1, \ldots, u_n\}$ is $p$-independent over $k$ then the $p^n$ monomials
$u_1^{i_1} \cdots u_n^{i_n}$ for $0 \le i_\nu < p$ are linearly independent over $k(K^p)$, so that by
the previous lemma they are also linearly independent over $k_\alpha(K^p)$ for some
$\alpha$.

(3) $\Rightarrow$ (4) is trivial; any $p$-basis will do.

(4) $\Rightarrow$ (2) If $0 \ne \omega \in \Omega_{K/k}$ and $B$ is a $p$-basis then there is a unique
expression $\omega = \sum c_id_{K/k}b_i$, with $\{b_1, \ldots, b_n\} \subset B$ a finite set and
$c_i \in K$. If we take $\alpha$ such that $b_1, \ldots, b_n$ are $p$-independent over
$k_\alpha$, then $d_{K/k_\alpha}b_i$ for $1 \le i \le n$ are linearly independent as elements of
$\Omega_{K/k_\alpha}$, so that $\omega$ has non-zero image in $\varprojlim \Omega_{K/k_\alpha}$.

(2) $\Rightarrow$ (1) Let $a \in K$ be such that $a \notin k(K^p)$; then $d_{K/k}a \ne 0$, so that
$d_{K/k}a \ne 0$ in $\Omega_{K/k_\alpha}$ for some $k_\alpha$, in other words $a \notin k_\alpha(K^p)$.
\end{proof}

\begin{lemma}\label{lem:30.5}
Let $k \subset K$ and $\mathcal{F} = \{k_\alpha\}_{\alpha \in I}$ be as in the previous lemma, and
assume that $\bigcap_\alpha k_\alpha(K^p) = k(K^p)$. Then if $L$ is an extension field of $K$
which is either separable over $K$ or finitely generated over $K$, we again have
$\bigcap_\alpha k_\alpha(L^p) = k(L^p)$.
\end{lemma}

\begin{proof}
(i) If $L/K$ is separable, choose a $p$-basis $B$ of $K/k$ and a $p$-basis $C$ of
$L/K$; then in view of the exact sequence (Theorem~\hyperref[thm:25.1]{25.1})
\[
  0 \to \Omega_{K/k} \otimes_K L \to \Omega_{L/k} \to \Omega_{L/K} \to 0,
\]
$B \cup C$ is a $p$-basis of $L/k$. If $b_1, \ldots, b_m \in B$ and $c_1, \ldots, c_n \in C$ are
finitely many distinct elements, then by the previous lemma, $b_1, \ldots, b_m$ are
$p$-independent over some $k_\alpha$, and from this (replacing $k$ by $k_\alpha$ in the
above exact sequence) we see that $b_1, \ldots, b_m, c_1, \ldots, c_n$ are
$p$-independent over $k_\alpha$.

(ii) If $L/K$ is finitely generated, then since any finitely generated extension
can be obtained as a succession of elementary extensions of type
\begin{enumerate}[(a)]
  \item $L = K(x)$ with $x$ separable over $K$,
  or
  \item $L = K(x)$ with $x^p = a \in K$,
\end{enumerate}
we need only consider case (b). We further divide this into two subcases: in
the first, $d_{K/k}a = 0$, and then from Theorem~\hyperref[thm:25.2]{25.2} and the fact that
$L \cong K[X]/(X^p - a)$ we get $\Omega_{L/k} = (\Omega_{K/k} \otimes L) \oplus Ldx$; from this
one sees easily that $\Omega_{L/k} \to \varprojlim \Omega_{L/k_\alpha}$ is injective. In the second
subcase, $d_{K/k}a \ne 0$ and now we have
$\Omega_{L/k} \cong ((\Omega_{K/k} \otimes L)/L \cdot d_{K/k}a) \oplus Ldx$. Hence if we take a
$p$-basis of $K/k$ of the form $\{a\} \cup B'$ with $a \notin B'$, then $\{x\} \cup B'$ is a
$p$-basis of $L/k$. Now $L/k$ satisfies condition (4) of the previous lemma, since
for $b_1, \ldots, b_m \in B'$, if we choose $\alpha$ such that $\{a, b_1, \ldots, b_m\} \subset K$
is $p$-independent over $k_\alpha$, then $\{x, b_1, \ldots, b_m\} \subset L$ is $p$-independent
over $k_\alpha$.
\end{proof}

\begin{lemma}\label{lem:30.6}
Let $K$ be a field of characteristic $p$, and let $\{K_\alpha\}$ be a directed family of
cofinite subfields $K_\alpha \subset K$ such that $\bigcap_\alpha K_\alpha = K^p$. Then if $L$ is a
finite field extension of $K$, there exists $\alpha$ such that
\[
  \rank_L \Omega_{L/K_\alpha} = \rank_L \Omega_{L/K}
\]
for all subfields $K' \subset K_\alpha$ with $[K : K'] < \infty$.
\end{lemma}

\begin{proof}
Let $K = K_0 \subset K_1 \subset \cdots \subset K_t = L$ be a chain of intermediate
fields such that $K_i = K_{i-1}(x_i)$, with $x_i$ either separable algebraic over
$K_{i-1}$ or $x_i^p \in K_{i-1}$. Then by the previous lemma we have
$\bigcap_\alpha K_\alpha(L^p) = L^p$, so that we need only prove the lemma for $t = 1$;
hence suppose that $L = K(x)$. If $L$ is separable over $K$ then in view of
$\Omega_{L/K'} = \Omega_{K/K'} \otimes_K L$, the assertion is clear (any $\alpha$ will do). Thus
suppose that $x^p = a \in K$, but $a \notin K^p$; then there exists $\alpha$ such that
$a \notin K_\alpha$. If $K' \subset K_\alpha$, then computing by means of
$L = K[X]/(X^p - a)$ we see that
\[
  \Omega_{L/K'} = ((\Omega_{K/K'} \otimes_K L) \oplus Ldx)/L \cdot d_{K/K'}a,
\]
and if $\rank \Omega_{L/K'} < \infty$ then $\rank \Omega_{L/K'} = \rank \Omega_{L/K}$.
\end{proof}

\begin{theorem}[\textbf{Nagata}]\label{thm:30.9}
Let $k$ be a field of characteristic $p$, $S = k[Y_1, \ldots, Y_m]$ and $P \in \Spec S$;
suppose that $\mathcal{F} = \{k_\alpha\}_{\alpha \in I}$ is a directed family of cofinite subfields of
$k$ such that $\bigcap_\alpha k_\alpha = k^p$. Then there exists $\alpha \in I$ such that for
every intermediate field $k^p \subset k' \subset k_\alpha$ with $[k : k'] < \infty$ the
following formula holds:
\[
  \rank \Der_{k'}(S/P) = \dim(S/P) + \rank \Der_{k'}(k).
\]
\end{theorem}

\begin{proof}
Set $A = S/P$, let $L$ be the field of fractions of $A$, and $\dim A = n$. Choose a
system of parameters $x_1, \ldots, x_n$ of $A$, set $B = k[x_1, \ldots, x_n]$, let $K$ be the
field of fractions of $B$ and $m_B$ its maximal ideal. Then $A$ is a finite $B$-module,
and hence $[L : K] < \infty$. If $k'$ is an intermediate field $k^p \subset k' \subset k$
with $[k : k'] = p^\rho < \infty$, and if $u_1, \ldots, u_\rho$ is a $p$-basis of $k$ over $k'$,
then $\{u_1, \ldots, u_\rho, x_1, \ldots, x_n\}$ is a $p$-basis of $B$ over
$C' = k'[x_1^p, \ldots, x_n^p]$, in the sense that $B$ is the free $C'$-module with basis
the set of $p$-monomials in $u_1, \ldots, u_\rho, x_1, \ldots, x_n$. A derivation from $B$ to
$B$ is continuous in the $m_B$-adic topology, and any element of $\Der_{k'}(B)$ is $0$ on
$C'$. Therefore $\Der_{k'}(B) = \Der_{k'}(B) = \Hom_B(\Omega_{B/C'}, B)$, and $\Omega_{B/C'}$ is
the free $B$-module of rank $n + \rho$ with basis $du_1, \ldots, du_\rho, dx_1, \ldots, dx_n$.
Therefore $\Der_{k'}(B)$ is also a free $B$-module of rank $n + \rho$, and the theorem
holds in case $A = B$.

We write $F'$ for the field of fractions of $C'$, and for $k_\alpha \in \mathcal{F}$ we set
$C_\alpha = k_\alpha[x_1^p, \ldots, x_n^p]$ and write $K_\alpha$ for the field of fractions of
$C_\alpha$. As above we have
\[
  \Der_{k'}(A) = \Der_{k'}(A) = \Hom_A(\Omega_{A/C'}, A),
\]
and since $A$ is a finite $C'$-module, $\Omega_{A/C'}$ is a finite $A$-module. Hence
\[
  \Der_{k'}(A) \otimes_A L = \Hom_L(\Omega_{A/C'} \otimes_A L, L) = \Hom_L(\Omega_{L/F'}, L),
\]
so that $\rank_L \Der_{k'}(A) = \rank_L \Omega_{L/F'}$. Moreover,
\[
  n + \rank \Der_{k'}(k) = \rank \Omega_{L/F'} = \rank \Omega_{K/F'},
\]
so that the conclusion of the theorem can be rewritten
\[
  \rank \Omega_{L/F'} = \rank \Omega_{K/F'}.
\]
Any element of $K_\alpha$ can be written with its denominator in
$k^p[x_1^p, \ldots, x_n^p] = B^p$. However, $B$ is faithfully flat over $C_\alpha$, so that
$K_\alpha \cap B = C_\alpha$. From this one deduces easily that
$\bigcap_\alpha K_\alpha = K^p$. Thus the theorem follows from the previous lemma.
\end{proof}

\begin{theorem}[\textbf{Nagata's Jacobian criterion}]\label{thm:30.10}
Let $k$ be a field of characteristic $p$, $S = k[[X_1, \ldots, X_n]]$, and let $I$, $P$ be
ideals of $S$ such that $I \subset P \in \Spec S$; set $S_P = R$, $R/IR = A$, and suppose
that $\htop IR = r$ and $I = (f_1, \ldots, f_t)$. Choose a $p$-basis $\{u_\gamma\}_{\gamma \in \Gamma}$ of
$k$, and define $D_\gamma \in \Der(k)$ by $D_\gamma(u_{\gamma'}) = \delta_{\gamma\gamma'}$. We make
any element $D \in \Der(k)$ act on the coefficients of power series, thus
extending $D$ to an element of $\Der(S)$.

Then the following conditions are equivalent:
\begin{enumerate}[(1)]
  \item $A$ is a regular local ring;
  \item there exists a finite number of elements $\alpha, \beta, \ldots, \gamma \in \Gamma$
        such that
        \[
          \rank J(f_1, \ldots, f_t; D_\alpha, \ldots, D_\gamma,
              \partial/\partial X_1, \ldots, \partial/\partial X_n)(P) = r.
        \]
\end{enumerate}
\end{theorem}

\begin{proof}
(2) $\Rightarrow$ (1) follows from Theorem~\hyperref[thm:30.4]{4}.

(1) $\Rightarrow$ (2) As we see from the proof of Lemma~\hyperref[lem:30.2]{2}, we need only
prove (2) in the case $I = P$. We let $\mathcal{F}$ be the family of subfields of $k$
obtained by adjoining all but a finite number of $\{u_\gamma\}_{\gamma \in \Gamma}$ to $k^p$;
then the conditions of Theorem~\hyperref[thm:30.9]{9} are satisfied, and there exists
$k' \in \mathcal{F}$ such that
\[
  \rank \Der_{k'}(S/P) = n - \htop P + \rank \Der_{k'}(k).
\]
If we set $[k : k'] = p^s$ then by construction there exist
$\gamma_1, \ldots, \gamma_s \in \Gamma$ such that $k = k'(u_{\gamma_1}, \ldots, u_{\gamma_s})$; now set
$C = k'[[X_1, \ldots, X_n]]$, so that $\Omega_{S/C}$ is the free $S$-module with basis
$du_{\gamma_1}, \ldots, du_{\gamma_s}, dX_1, \ldots, dX_n$, and arguing as in the proof of
Theorem~\hyperref[thm:30.8]{8}, we see that the derivations of $S/P$ over $k'$ are all induced by
derivations of $S$ over $k'$, and that if $g_1, \ldots, g_m$ are generators of $P$, we have
\[
  \rank J(g_1, \ldots, g_m; D_{\gamma_1}, \ldots, D_{\gamma_s},
      \partial/\partial X_1, \ldots, \partial/\partial X_n)(P) = \htop P.
\]
\end{proof}

\begin{remark}
Conditions (1) and (2) in Theorem~\hyperref[thm:30.10]{10} are of the same form as the
corresponding conditions in Theorem~\hyperref[thm:30.5]{5}. Condition (3) of
Theorem~\hyperref[thm:30.5]{5} is not applicable as it stands to the present situation, since
$\Omega_{A/k}$ is in general not a finite $A$-module. However, in general for a module
$M$ over a local ring $(R, m)$, if we write (in temporary notation)
$\tilde{M} = M/\bigcap_\nu m^\nu M$ for the associated separated module, then for any
$R$-module $N$ which is separated ($N = \tilde{N}$) we have
$\Hom_R(\tilde{M}, N) = \Hom_R(M, N)$. Therefore in the situation of the above theorem
we have $\Der_k(S) = \Hom_S(\tilde{\Omega}_{S/k}, S)$ and $\Der_k(A) = \Hom_A(\tilde{\Omega}_{A/k}, A)$,
and moreover $\tilde{\Omega}_{S/k}$ is a free $S$-module of rank $\rho + \rank \Omega_{k/\pi}$.
From this, using the same argument as in the proof of Theorem~\hyperref[thm:30.5]{5}, we see
that replacing $\Omega_{A/k}$ by $\tilde{\Omega}_{A/k}$ in (3) of Theorem~\hyperref[thm:30.5]{5}, this
condition is equivalent to (1) and (2) of Theorem~\hyperref[thm:30.10]{10}. Verifying this is a
suitable exercise for the reader.
\end{remark}

\begin{corollary*}
Let $A$ be a complete Noetherian local ring; then $\Reg(A)$ is an open subset of
$\Spec A$.
\end{corollary*}

\begin{proof}
If $A$ is equicharacteristic, and $k$ is a coefficient field of $A$, then $A$ is of the
form $A = S/I$ with $S = k[[X_1, \ldots, X_n]]$, so that by Theorems~\hyperref[thm:30.8]{8} and
\hyperref[thm:30.10]{10}, we see as in the corollary of Theorem~\hyperref[thm:30.5]{5} that $\Reg(A)$ is open.

If $A$ is of unequal characteristic, then by Theorem~\hyperref[thm:24.4]{24.4}, it is enough to
prove, under the assumption that $A$ is an integral domain, that $\Reg(A)$
contains a non-empty open subset of $\Spec A$. Now by
Theorem~\hyperref[thm:29.5]{29.5}, $A$ contains a regular local ring $B$, and is a finite module
over $B$, so that if we let $L$ and $K$ be the fields of fractions of $A$ and $B$,
then $L$ is a finite extension of $K$, and is separable since $\chark K = 0$. Replacing
$A$ by $A_b$, and $B$ by $B_b$ for some suitable $0 \ne b \in B$ we can assume that $A$ is
a free $B$-module (although $A$ and $B$ are no longer local rings, $B$ remains
regular). Suppose that $\omega_1, \ldots, \omega_m$ are a basis of $A$ as a $B$-module, and
consider the discriminant
\[
  d = \det(\tr_{L/K}(\omega_i\omega_j)).
\]
Let us prove that if $P \in \Spec A$ is such that $d \notin P$ then $P \in \Reg(A)$. Set
$p = P \cap B$; then $A_p$ is flat over $B_p$, and $B_p$ is regular, so that we need only
prove that the fibre $A_p \otimes_B \kappa(p)$ is regular. Now $A = \sum B\omega_i$, so
that $A \otimes_B \kappa(p) = \sum \kappa(p)\bar{\omega}_i$, and in $\kappa(p)$ we have
$\det(\tr(\bar{\omega}_i\bar{\omega}_j)) = \bar{d} \ne 0$; hence $A \otimes_B \kappa(p)$ is reduced,
and is therefore a direct product of fields. Therefore $A_p \otimes_B \kappa(p)$ is a
field; this proves that $\Reg(A)$ contains a non-empty open subset of $\Spec A$.
\end{proof}

% ============ 习题 ============
\section*{Exercises to \S\ref{sec:30}}

\begin{exercise}\label{ex:30.1-ex}
Let $(A, m)$ be a complete Noetherian local ring, and $D = (D_0, D_1, \ldots) \in
\mathrm{HS}(A)$. Suppose that $x \in m$ satisfies $D_1x = 1$ and $D_ix = 0$ for $i > 0$,
that is $E_t(x) = x + t$, and define $\varphi = E_{-t}$ by
$\varphi(a) = \sum_i (-x)^iD_ia$. Then $\varphi$ is an endomorphism of $A$ to $A$, with
$\Ker \varphi = xA$; and if we set $C = \Im \varphi$ then $A = C[x] \cong C[[X]]$.
\end{exercise}

\begin{exercise}\label{ex:30.2-ex}
If the conditions of Theorem~\hyperref[thm:30.5]{5} hold, is it true that $A$ is 0-smooth over
the $k'$ appearing in (3)?
\end{exercise}

\begin{exercise}\label{ex:30.3-ex}
Are conditions (1) and (3) of Theorem~\hyperref[thm:30.3]{3} still equivalent if
$S = k[[X_1, \ldots, X_n]]$?
\end{exercise}

\begin{exercise}\label{ex:30.4-ex}
Let $R$ be a regular ring containing a field of characteristic $0$; if (WJ) holds in
$R$ then it also holds in the polynomial ring $R[X_1, \ldots, X_n]$.
\end{exercise}


% ============ Chapter 11. Applications of complete local rings (原书 pp.246-265) ============
\chapter{Applications of complete local rings}\label{ch:11}

It has become clear in the previous chapter that the completion of a local
ring has a number of good properties. In this chapter we give some
applications of this. \S\ref{sec:31} centres on the work of Ratliff, giving
characterisations of catenary and universally catenary rings; Ratliff is
practically the only current practitioner of the Krull and Nagata tradition,
obtaining deep results by a fluent command of the methods of classical
ideal theory, and there is something about his proofs which is to be
savoured. In \S\ref{sec:32} we discuss Grothendieck's theory of the formal fibre;
this book is already long enough, and we have only covered a part of the
theory of G-rings, referring to [G2] and [M] for more details. In \S\ref{sec:33} we
discuss some further important applications, again sending the reader to
appropriate references for the details.

% ============ §31 Chains of prime ideals (原书 pp.246-254) ============
\section{Chains of prime ideals}\label{sec:31}\index{chains of prime ideals, section 31}\index{catenary ring, section 31}\index{equidimensional ring, section 31}\index{formally equidimensional, section 31}\index{formally catenary, section 31}\index{universally catenary, section 31}

\begin{theorem}\label{thm:31.1}
Let $A$ be a Noetherian ring and $P \in \Spec(A)$. Then there are at most finitely
many prime ideals $P'$ of $A$ satisfying $P \subset P'$, $\htop(P'/P) = 1$ and
$\htop P' > \htop P + 1$. (Ratliff [3] in the semi-local case, and McAdam [3] in
the general case.)
\end{theorem}

\begin{proof}
Let $\htop P = n$ and take $a_1, \ldots, a_n \in P$ such that $\htop(a_1, \ldots, a_n) = n$. Set
$I = (a_1, \ldots, a_n)$ and let $P_1 = P, P_2, \ldots, P_r$ be the minimal prime divisors of $I$.
In general, if $\{Q_\lambda\}_\Lambda$ is an infinite set of prime ideals such that
$Q_\lambda \supset P$ and $\htop(Q_\lambda/P) = 1$, then $\bigcap_\Lambda Q_\lambda = P$. This is because
$\bigcap Q_\lambda$ is equal to its own radical, hence is a finite intersection of prime
ideals containing $P$, hence either
\[
  \bigcap_\Lambda Q_\lambda = P
\]
or
\[
  \bigcap_\Lambda Q_\lambda = Q_{\lambda_1} \cap \cdots \cap Q_{\lambda_k};
\]
but the second case cannot occur since $Q_{\lambda_j} \not\supset Q_{\lambda_j}$ for
$j \ne \{\lambda_1, \ldots, \lambda_k\}$. Therefore if there were an infinite number of prime
ideals $Q_\lambda$ such that $Q_\lambda \supset P$, $\htop(Q_\lambda/P) = 1$ and
$\htop Q_\lambda \ne n + 1$, then $\bigcap Q_\lambda = P$. Hence there would exist a
$Q_\lambda$, say $Q_0$, which does not contain $P_1, \ldots, P_r$. Then $P$ would be the only
prime ideal lying between $Q_0$ and $I$. Let $b \in Q_0 - P$. Then $Q_0$ would be a
minimal prime divisor of $I + bA = (a_1, \ldots, a_n, b)$, so that
$\htop Q_0 \le n + 1$, a contradiction.
\end{proof}

\begin{theorem}[\textbf{Ratliff's weak existence theorem}]\label{thm:31.2}
Let $A$ be a Noetherian ring, and $p$, $P \in \Spec A$ be such that $p \subset P$,
$\htop p = h$ and $\htop(P/p) = d > 1$; then there exist infinitely many $p' \in \Spec A$ with
the properties
\[
  p \subset p' \subset P, \quad \htop p' = h + 1 \quad \text{and} \quad \htop(P/p') = d - 1.
\]
\end{theorem}

\begin{proof}
We first observe that if $P \supset p_1 \supset p_2 \supset \cdots \supset p_d = p$ is a strictly
decreasing chain of prime ideals, and if $p_{d-2} \subset p' \subset p$ with
$\htop p' = h + 1$ then $\htop(P/p') = d - 1$. Now there exist infinitely many
$p' \in \Spec A$ such that $p_{d-2} \subset p' \subset p$ and $\htop(p'/p) = 1$: for if
$p_1', \ldots, p_k'$ are a finite set of these, let $a \in p_{d-2} - \bigcup p_i'$, and let
$p_{k+1}'$ be a minimal prime divisor of $p + aA$ contained in $p_{d-2}$; then
$\htop(p_{k+1}'/p) = 1$. By the previous theorem, all but finitely many of these
satisfy $\htop p' = \htop p + 1$.
\end{proof}

\begin{lemma}\label{lem:31.1}
Let $A$ be a Noetherian ring, and $P \in \Spec A$ with $\htop P = h > 1$; suppose that
$u \in P$ is such that $\htop(uA) = 1$. Then there exist infinitely many prime ideals
$Q \subset P$ such that
\[
  u \notin Q \quad \text{and} \quad \htop Q = h - 1.
\]
\end{lemma}

\begin{proof}
Suppose that $p_1, \ldots, p_t$ are the minimal prime ideals of $A$, and let
$P_1, \ldots, P_s$ be finitely many given height $1$ prime ideals not containing $u$.
Let $Q_1, \ldots, Q_m$ be the minimal prime divisors of $uA$, so that these are also
height $1$ prime ideals of $A$. Since $h > 1$, there exists $v \in P$ not contained in
any $p_i$, $P_j$ or $Q_k$. Then
\[
  \htop(u, v) = 2 \quad \text{and} \quad \htop(v) = 1.
\]
Now let $P_{s+1}, \ldots, P_{s'}$ be the minimal prime divisors of $(v)$; continuing in the
same way we can find infinitely many height $1$ prime ideals not containing $u$, so
that if $h = 2$ we are done. If $h > 2$ we set $\bar{A} = A/(u)$ and
$\bar{P} = P/(u)$, so that $\htop \bar{P} = h - 1$, and since the image $\bar{v}$ of $v$ in
$\bar{A}$ satisfies $\htop(\bar{v}) = 1$, by induction on $h$ we can find infinitely many prime
ideals $\bar{Q}$ of $\bar{A}$ satisfying
\[
  \bar{Q} \subset \bar{P} \quad \text{and} \quad \htop \bar{Q} = h - 2.
\]
The inverse image $Q$ of $\bar{Q}$ in $A$ does not contain $u$, and from
$Q/(u) = \bar{Q}$ we have $\htop Q = h - 1$.
\end{proof}

\begin{theorem}[\textbf{Ratliff's strong existence theorem}]\label{thm:31.3}
Let $A$ be a Noetherian integral domain, $p$, $P \in \Spec A$, and suppose that
$\htop p = h > 0$ and $\htop(P/p) = d$. Then for each $i$ with $0 \le i \le d$ the set
\[
  \{p' \in \Spec A \mid p' \subset P, \htop(P/p') = d - i \text{ and } \htop p' = h + i\}
\]
is infinite.
\end{theorem}

\begin{proof}
For $i > 0$ this follows at once from the weak existence theorem, so that we
consider the case $i = 0$.

\textit{Step 1.} Replacing $A$ by $A_p$ we can assume that $(A, P)$ is a local integral
domain. Choose $a_1, \ldots, a_h$ such that $\htop(a_1, \ldots, a_j) = j$ for $j = 1, \ldots, h$,
and set
\[
  \mathfrak{a} = (a_1, \ldots, a_h) \quad \text{and} \quad \mathfrak{b} = (a_1, \ldots, a_{h-1}).
\]
Then $p$ is a minimal prime divisor of $\mathfrak{a}$. Let
\[
  \mathfrak{a} = \mathfrak{a}_1 \cap \cdots \cap \mathfrak{a}_r \quad \text{and} \quad
  \mathfrak{b} = \mathfrak{b}_1 \cap \cdots \cap \mathfrak{b}_s
\]
be shortest primary decompositions of $\mathfrak{a}$ and $\mathfrak{b}$, and let $p_i$, $p_j'$ be the
prime divisors of $\mathfrak{a}_i$ and $\mathfrak{b}_j$, respectively; we can assume that
$p = p_1$. Suppose that $\mathfrak{b}_{t+1}, \ldots, \mathfrak{b}_s$ are all the $\mathfrak{b}_j$ not contained
in $p$. Then since
\[
  \mathfrak{a}_2 \cap \cdots \cap \mathfrak{a}_r \cap \mathfrak{b}_{t+1} \cap \cdots \cap
      \mathfrak{b}_s \not\subset p,
\]
we can choose an element $y \in P$ contained in the left-hand side and not
contained in $p$. Now $p_j' \subset p$ for $1 \le j \le t$, so that $y \notin p_1$, and hence
\[
  \mathfrak{a} : yA = \mathfrak{a}_1, \quad \text{and} \quad
  \mathfrak{b} : yA = \mathfrak{b}_1 \cap \cdots \cap \mathfrak{b}_t.
\]
Now set
\[
  B = A[x_1, \ldots, x_h], \quad \text{where } x_i = a_i/y,
\]
\[
  I = (x_1, \ldots, x_h)B \quad \text{and} \quad Q = PB + I.
\]

\textit{Step 2.} We prove that
\[
  B/I \cong A/\mathfrak{a}_1, \quad Q \in \Spec B \quad \text{and} \quad \htop Q = h + d.
\]
A general element of $B$ can be written in the form
\[
  a/y^\nu, \quad \text{with } a \in (\mathfrak{a} + yA)^\nu, \text{ for some } \nu \ge 0.
\]
Now $B = A + I$, so that $B/I \cong A/(I \cap A)$. Now if $\alpha \in I \cap A$, then there
exists $\nu > 0$ such that $y^\nu\alpha \in \mathfrak{a}(\mathfrak{a} + yA)^{\nu-1}$. Since
$\mathfrak{a} : yA = \mathfrak{a}_1$, we have $\alpha \in \mathfrak{a}_1$. Conversely,
$y\mathfrak{a}_1 \subset \mathfrak{a}$, giving $\mathfrak{a}_1 \subset I \cap A$, and hence $I \cap A = \mathfrak{a}_1$.
Therefore
\[
  B/I \cong A/\mathfrak{a}_1.
\]
Under this isomorphism, the prime ideals $p/\mathfrak{a}_1$ and $P/\mathfrak{a}_1$ of $A/\mathfrak{a}_1$
correspond to $(pB + I)/I$ and $(PB + I)/I$ respectively in $B/I$; hence, setting
\[
  \mathfrak{q} = pB + I = p + I \quad \text{and} \quad Q = PB + I = P + I,
\]
we have $\mathfrak{q}$, $Q \in \Spec B$, with $B/\mathfrak{q} = A/p$ and $B/Q = A/P$, and $Q$ is a
maximal ideal of $B$. Also, $B/I = A/\mathfrak{a}_1$, and since $I$ is generated by $h$
elements and $\mathfrak{a}_1$ is a $p$-primary ideal, we get
\[
  \htop Q = \dim B_Q \ge h + \htop(P/\mathfrak{a}_1) = h + \htop(P/p) = h + d.
\]
On the other hand, from $y \notin p$ we get $\mathfrak{q} = pA[y^{-1}] \cap B$: indeed, if
$a \in pA[y^{-1}] \cap B$, then we can write $a = c/y^\nu$ with
$c \in p \cap (\mathfrak{a} + y^\nu A)$; since $\mathfrak{a} \subset p$ and $y \notin p$, we have
\[
  p \cap (\mathfrak{a} + y^\nu A) = \mathfrak{a}(\mathfrak{a} + yA)^{\nu-1} + y^\nu p,
\]
and therefore $c \in I + p = \mathfrak{q}$. Also, $A[y^{-1}] = B[y^{-1}]$, so that
\[
  \htop \mathfrak{q} = \htop \mathfrak{q}A[y^{-1}] = \htop pA[y^{-1}] = \htop p = h.
\]
However, $Q/\mathfrak{q} = P/p$, so that
\[
  \htop(Q/\mathfrak{q}) = \htop(P/p) = d
\]
and therefore $\htop Q \ge d + h$. Putting this together with the previous
inequality, we see that $\htop Q = d + h$.

\textit{Step 3.} For $\nu = 1, 2, \ldots$, set
\[
  J_\nu = (x_1, \ldots, x_{h-1}, x_h - y^\nu)B,
\]
let $Q_\nu$ be a minimal prime divisor of $J_\nu$ satisfying
\[
  Q_\nu \subset Q \quad \text{and} \quad \htop(Q/Q_\nu) = \htop(Q/J_\nu),
\]
and set $P_\nu = Q_\nu \cap A$. We will complete the proof by showing that
$P_1, P_2, \ldots$ are all distinct, and that
\[
  \htop P_\nu = h, \quad \htop(P/P_\nu) = d \quad \text{for all } \nu.
\]
Now $B = A + J_\nu$, so that $B/Q_\nu \cong A/P_\nu$ and $Q/Q_\nu \cong P/P_\nu$, hence
$\htop(P/P_\nu) = \htop(Q/Q_\nu) = \htop(Q/J_\nu)$; moreover, $\htop Q = d + h$, and since
$J_\nu$ is generated by $h$ elements,
\[
  \htop(P/P_\nu) \ge d + h - h = d.
\]
Therefore we have
\[
  d \le \htop(P/P_\nu) = \htop(Q/Q_\nu) \le \htop Q - \htop Q_\nu = d + h - \htop Q_\nu,
\]
and so if we can prove that
\[
  (*) \quad \htop Q_\nu = \htop P_\nu \ge h
\]
then this will show simultaneously that $\htop(P/P_\nu) = d$ and $\htop P_\nu = h$.

We have already seen that $\mathfrak{q} = pA[y^{-1}] \cap B$, so that
$\mathfrak{q} \cap A = p$, and hence $y \notin \mathfrak{q}$. Also, $\mathfrak{a}_1$ is a $p$-primary ideal
with $B/I \cong A/\mathfrak{a}_1$, so that $I$ is a $\mathfrak{q}$-primary ideal, and from
$\htop \mathfrak{q} = h$ we get
\[
  \htop(I + yB) = \htop(x_1, \ldots, x_h, y)B = h + 1.
\]
In addition, we have $I + yB = J_\nu + yB$, and since $Q_\nu$ is a minimal prime
divisor of the ideal $J_\nu$, which is generated by $h$ elements,
$\htop Q_\nu \le h$, and hence $y \notin Q_\nu$, and
\[
  Q_\nu = Q_\nu A[y^{-1}] \cap B \quad \text{and} \quad P_\nu = Q_\nu A[y^{-1}] \cap A,
\]
so that, by \S\ref{sec:20}, Example 1,
\[
  \htop Q_\nu = \htop P_\nu.
\]
Furthermore
\[
  (\mathfrak{b} : yA)B \subset (x_1, \ldots, x_{h-1})B,
\]
\[
  (\mathfrak{b} : yA) + (a_h - y^{\nu+1})A \subset J_\nu \cap A \subset Q_\nu \cap A = P_\nu,
\]
and since all prime divisors of $(\mathfrak{b} : yA)$ are also prime divisors of
$\mathfrak{b}$, we have $\htop(\mathfrak{b} : yA) = h - 1$. Now $a_h \in \mathfrak{q}$ and
$y \notin p$, so that $a_h - y^{\nu+1} \notin p$, and since all minimal prime divisors of
$(\mathfrak{b} : yA)$ are contained in $p$ they do not contain $a_h - y^{\nu+1}$.
Therefore $\htop P_\nu \ge h$. This completes the proof of (*).

\textit{Step 4.} If $\nu < \mu$ then $(Q + Q_\mu)B_a$ contains $y^\nu = (y^\nu - y^\mu)/(1 - y^{\mu-\nu})$,
and therefore contains $(J_\nu + y^\nu B)B_a = (I + y^\nu B)B_a$, so that
\[
  \htop(Q + Q_\mu)B_a \ge h + 1, \quad \text{but } \htop Q_\nu \le h,
\]
and therefore $Q_\nu \ne Q_\mu$. From $Q_\nu A[y^{-1}] = P_\nu A[y^{-1}]$ we see that
$P_1, P_2, \ldots$ must also all be distinct.
\end{proof}

\begin{theorem}\label{thm:31.4}
A Noetherian local integral domain $(A, m)$ is catenary if and only if
\[
  \htop p + \coht p = \dim A \quad \text{for all } p \in \Spec A.
\]
\end{theorem}

\begin{proof}
`Only if' is trivial, and we prove `if'. Let $\dim A = n$; if $A$ is not catenary, then
there exist $p$, $P \in \Spec A$ such that
\[
  p \subset P, \quad \htop(P/p) = 1 \quad \text{but } \htop P > \htop p + 1.
\]
Set $\htop(m/P) = d$. Applying the strong existence theorem to $A/p$ we see that there
exist infinitely many $P_\lambda \in \Spec A$ such that
\[
  p \subset P_\lambda, \quad \htop(P_\lambda/p) = 1 \quad \text{and} \quad \htop(m/P_\lambda) = d.
\]
However, by assumption, $\htop(m/P_\lambda) + \htop P_\lambda = n$, so that
\[
  \htop P_\lambda = n - d = \htop P > \htop p + 1.
\]
But according to Theorem~\hyperref[thm:31.1]{1}, there are only finitely many such
$P_\lambda$, and we have a contradiction.
\end{proof}

If $A$ is a ring of finite Krull dimension, we say that $A$ is \emph{equidimensional} if
$\dim A/p = \dim A$ for every minimal prime $p$ of $A$.

\begin{lemma}\label{lem:31.2}
If an equidimensional local ring $(A, m)$ is catenary then
\[
  \htop p_2 = \htop p_1 + \htop(p_2/p_1) \quad \text{for all } p_1, p_2 \in \Spec A
      \text{ with } p_1 \subset p_2.
\]
\end{lemma}

\begin{proof}
If we choose a minimal prime ideal $p \subset p_1$ then
$\htop(p_2/p) = \htop(m/p) - \htop(m/p_2) = \dim A - \htop(m/p_2)$, and this is independent
of the choice of $p$, so that $\htop p_2 = \htop(p_2/p)$. Similarly,
$\htop p_1 = \htop(p_1/p)$, so that $\htop p_2 = \htop(p_2/p_1) + \htop p_1$.
\end{proof}

\begin{theorem}\label{thm:31.5}
Let $A$, $B$ be Noetherian local rings, and $A \to B$ a local homomorphism. If $B$ is
equidimensional and catenary and is flat over $A$ then $A$ is also
equidimensional and catenary, and $B/pB$ is equidimensional for every
$p \in \Spec A$.
\end{theorem}

\begin{proof}
Write $m$ and $\mathfrak{M}$ for the maximal ideals of $A$ and $B$. For any minimal prime ideal
$p_0$ of $A$ there exists a minimal prime ideal $P_0$ of $B$ lying over $p_0$; then
$\dim B/P_0 = \dim B$, so that $\dim B/p_0B = \dim B$, and then by
Theorem~\hyperref[thm:15.1]{15.1} we have
\[
  \htop(m/p_0) = \htop(\mathfrak{M}/p_0B) - \htop(\mathfrak{M}/mB) = \dim B - \htop(\mathfrak{M}/mB).
\]
This is independent of the choice of $p_0$, so that $A$ is equidimensional. If
$p \in \Spec A$ and $P$ is a minimal prime divisor of $pB$ then from the going-down
theorem (Theorem~\hyperref[thm:9.5]{9.5}) we see that $P \cap A = p$, so that by
Theorem~\hyperref[thm:15.1]{15.1}, $\htop P = \htop p$, and therefore
$\htop(\mathfrak{M}/P) = \htop \mathfrak{M} - \htop P = \htop \mathfrak{M} - \htop p$ is determined by $p$ only. That
is, $B/pB$ is equidimensional. Also, if $p' \in \Spec A$ is such that $p' \subset p$ and
$\htop(p/p') = 1$, we let $P'$ be a minimal prime divisor of $p'B$ contained in $P$;
then $B/p'B$ is also equidimensional and flat over $A/p'$, so that
$\htop(P/P') = \htop(P/p'B) = \htop(p/p') = 1$. However, $B$ is equidimensional and
catenary, so that $\htop(P/P') = \htop P - \htop P' = \htop p - \htop p'$, and therefore
$\htop p = \htop p' + 1$, and so $A$ is catenary.
\end{proof}

\begin{corollary*}
Let $A$ be a quotient of a regular local ring $R$. If $A$ is equidimensional then
so is its completion $A^*$.
\end{corollary*}

\begin{proof}
Let $P_0$ be a minimal prime ideal of $A^*$ and $p_0 = P_0 \cap A$; then writing
$\mathfrak{p} \subset R$ for the inverse image of $p_0$, we have
$R^*/\mathfrak{p}R^* = A^*/p_0A^*$. $R^*$ is an integral domain, and therefore
equidimensional, so we can apply the theorem to $R \to R^*$ and see that
$R^*/\mathfrak{p}R^*$ is equidimensional. Hence
\[
  \dim A^*/P_0 = \dim A^*/p_0A^* = \dim A/p_0 = \dim A.
\]
\end{proof}

\begin{definition*}
We say that a Noetherian local ring $A$ is \emph{formally equidimensional} (or
\emph{quasi-unmixed}) if its completion $A^*$ is equidimensional.
\end{definition*}

\begin{theorem}\label{thm:31.6}
Let $(A, m)$ be a formally equidimensional Noetherian local ring.
\begin{enumerate}[(i)]
  \item $A_p$ is formally equidimensional for every $p \in \Spec A$.
  \item If $I$ is an ideal of $A$, then
        \[
          A/I \text{ is equidimensional} \iff A/I \text{ is formally equidimensional}.
        \]
  \item If $B$ is a local ring which is essentially of finite type over $A$ (see
        p.~232) and if $B$ is equidimensional then it is also formally
        equidimensional.
  \item $A$ is universally catenary.
\end{enumerate}
\end{theorem}

\begin{proof}
(i) Let $P \in \Spec(A^*)$ be such that $P \cap A = p$, and set $B = (A^*)_P$.
Since $B$ is flat over $A_p$, by Theorem~\hyperref[thm:22.4]{22.4}, $B^*$ is flat over $(A_p)^*$. Now $B$ is
a quotient of a regular local ring, and is equidimensional, so that by the above
corollary, $B^*$ is also equidimensional. Hence by Theorem~\hyperref[thm:31.5]{5}, $(A_p)^*$ is also
equidimensional.

(ii) follows easily from Theorem~\hyperref[thm:31.5]{5}.

(iii) $B$ is a localisation of a quotient of $A[X_1, \ldots, X_n]$ for some $n$, so that
by (ii) we need only show that a localisation $B$ of $A[X_1, \ldots, X_n]$ is formally
equidimensional. Now $A^*[[X_1, \ldots, X_n]]$ is faithfully flat over
$A[X_1, \ldots, X_n]$, so that there is a local ring $C$ which is a localisation of
$A^*[[X_1, \ldots, X_n]]$, and a local homomorphism $B \to C$ such that $C$ is flat over
$B$, and hence $C^*$ is flat over $B^*$. By the remark after Theorem~\hyperref[thm:15.5]{15.5}, $C$ is
equidimensional, and is a quotient of a regular local ring, so that by the
corollary of Theorem~\hyperref[thm:31.5]{5}, $C^*$ is also equidimensional; hence by
Theorem~\hyperref[thm:31.5]{5}, $B^*$ is also equidimensional.

(iv) Any local integral domain essentially of finite type over $A$ is formally
equidimensional, and hence catenary, so that any integral domain which is
finitely generated over $A$ is catenary.
\end{proof}

We say that a Noetherian local ring $A$ is \emph{formally catenary} ([G2], (7.1.9)) if
$A/p$ is formally equidimensional for every $p \in \Spec A$. One sees easily from the
above theorem that formally catenary implies universally catenary. The converse
of this was proved by Ratliff [2]. Universally catenary is a property of finitely
generated $A$-algebras, and we have to deduce from this a property of the
completion, so that the proof is difficult. Before giving Ratliff's proof we make
the following observation.

Let $(R, m)$ be a Noetherian local integral domain, $K$ the field of fractions of $R$,
and $R'$ the integral closure of $R$ in $K$; let $S$ be an intermediate ring
$R \subset S \subset R'$ such that $S$ is a finite $R$-module. $S$ is a semilocal ring, and its
completion $S^*$ (with respect to the $m$-adic topology, which coincides with the
$\rad(S)$-adic topology) can be identified with $R^* \otimes_R S$. Now $R^*$ is flat over
$R$, so that $R \subset S \subset R' \subset K$ gives
$R^* \subset S^* \subset R^* \otimes_R R' \subset R^* \otimes_R K$. The ring $R^* \otimes_R K$ is the
localisation of $R^*$ with respect to $R - (0)$, so that writing $T$ for the total ring of
fractions of $R^*$, we can consider $R^* \otimes_R K \subset T$, and hence
$R^* \subset S^* \subset T$. This leads to the possibility that properties of $R^*$ will be
reflected in some $S$.

\begin{theorem}\label{thm:31.7}
For a Noetherian local ring $A$, the following conditions are equivalent:
\begin{enumerate}[(1)]
  \item $A$ is formally catenary,
  \item $A$ is universally catenary,
  \item $A[X]$ is catenary.
\end{enumerate}
\end{theorem}

\begin{proof}
(1) $\Rightarrow$ (2) was proved in Theorem~\hyperref[thm:31.6]{6} and (2) $\Rightarrow$ (3) is trivial. We
prove (3) $\Rightarrow$ (1). Suppose then that $A[X]$ is catenary; we will prove that $A^*$ is
equidimensional by assuming the contrary and deriving a contradiction. The
proof breaks up into several lemmas.

\begin{lemma}\label{lem:31.3}
Let $(R, m)$ be a catenary Noetherian local integral domain, and let $R^*$ be its
completion. Let $\dim R = n$, and suppose that there exists a minimal prime $Q$ of
$R^*$ such that
\[
  1 < \dim(R^*/Q) = d \le n.
\]
For $i = 1, 2, \ldots, d - 1$, write $\mathcal{O}_i$ for the set of $p \in \Spec R$ satisfying the
conditions
\begin{enumerate}[(1)]
  \item $\htop p = i$,
and (2) there exists a minimal prime divisor $P$ of $pR^*$ such that $Q \subset P$ and
$\dim(R^*/P) = d - i$.
\end{enumerate}
Then $\mathcal{O}_i$ is non-empty for each $i$.
\end{lemma}

\begin{proof}
We work by induction on $i$. If $0 \ne a \in m$ then any minimal prime divisor $P$ of
$aR^* + Q$ satisfies $\htop(P/Q) = 1$, $P \cap R \ne 0$ and contains $a$. Hence if we set
$\mathcal{M} = \{P \in \Spec(R^*) \mid Q \subset P, \htop(P/Q) = 1 \text{ and } P \cap R \ne 0\}$, then
$m = \bigcup_{P \in \mathcal{M}}(P \cap R)$. Now $\htop(mR^*/Q) = d > 1$, so that
$mR^* \notin \mathcal{M}$, and hence $m$ itself is not of the form $P \cap R$ for $P \in \mathcal{M}$;
therefore both $\mathcal{M}$ and $\{P \cap R \mid P \in \mathcal{M}\}$ are infinite sets. By
Theorem~\hyperref[thm:31.1]{1}, $\mathcal{M}' = \{P \in \mathcal{M} \mid \htop P = 1\}$ is also infinite; choose any
$P \in \mathcal{M}'$, and set $p = P \cap R$. Then $0 < \htop p = \htop pR^* \le \htop P = 1$, so that
$\htop p = 1$, and $P$ is a minimal prime divisor of $pR^*$. Since $R^*$ is catenary,
$\dim(R^*/P) = \dim(R^*/Q) - \htop(P/Q) = d - 1$; hence $p \in \mathcal{O}_1$ and the
assertion is true for $i = 1$.

If $i > 1$, take $p$ as above, and set $\bar{R} = R/p$ and $\bar{P} = P/pR^*$; then since
$\bar{R}$ is catenary, $\dim \bar{R} = n - 1$, and $\bar{P}$ is a minimal prime divisor of
$\bar{R}^* = R^*/pR^*$, with $\dim(\bar{R}^*/\bar{P}) = \dim(R^*/P) = d - 1 < n - 1$. Hence by
induction there exists a prime ideal $\bar{p}_i = p_i/p$ of $\bar{R}$ of height $i - 1$, and
a minimal prime divisor $\bar{P}_i$ of $\bar{p}_i\bar{R}^*$ such that $\bar{P} \subset \bar{P}_i$ and
$\dim(\bar{R}^*/\bar{P}_i) = (d - 1) - (i - 1) = d - i$. If $P_i = \bar{P}_i/pR^*$ then $P_i$ is a minimal
prime divisor of $p_iR^*$ containing $P$, and hence $Q$, and $R^*/P_i = \bar{R}^*/\bar{P}_i$ is
$(d - i)$-dimensional, so that from the fact that $R$ is catenary we get
$\htop p_i = \htop \bar{p}_i + \htop p = i$, and $p_i \in \mathcal{O}_i$.
\end{proof}

\begin{lemma}\label{lem:31.4}
Let $(R, m)$ be a Noetherian local integral domain, $R^*$ its completion, and let
$0 = \mathfrak{q}_1 \cap \cdots \cap \mathfrak{q}_r$ be a shortest primary decomposition of $0$ in $R^*$, with
$P_i = \sqrt{\mathfrak{q}_i}$ for $1 \le i \le r$. Suppose that $P_1$ satisfies
\[
  \htop P_1 = 0, \quad \coht P_1 = 1 < \dim R.
\]
Then there exist $b \in m$ and $\delta \in (\mathfrak{q}_2 \cap \cdots \cap \mathfrak{q}_r) - P_1$ with the
following properties:
\begin{enumerate}[(1)]
  \item $b - \delta \in \mathfrak{q}_1$,
  \item $(b, \delta)R^* = (b, c)R^*$,
and (3) $c/b \notin R$ but is integral over $R$.
\end{enumerate}
\end{lemma}

\begin{proof}
\textit{Step 1.} $P_1$ is a minimal prime ideal and $r > 1$, so that
$\mathfrak{q}_2 \cdots \mathfrak{q}_r \not\subset P_1$, and we can choose
$\delta' \in (\mathfrak{q}_2 \cap \cdots \cap \mathfrak{q}_r) - P_1$. Since $\coht P_1 = 1$, it follows that
$\mathfrak{q}_1 + \delta'R^*$ is $(mR^*)$-primary, and hence if we set
\[
  \mathfrak{a} = (\mathfrak{q}_1 + \delta'R^*) \cap R,
\]
then this is an $m$-primary ideal. Now if $0 \ne b \in \mathfrak{a}$ is any element, then in
$R^*$ we can write
\[
  b = \xi + \beta\delta' \quad \text{with } \xi \in \mathfrak{q}_1 \text{ and } \beta \in R^*.
\]
Since $b$ is not a zero-divisor in $R^*$ we have $\beta\delta' \notin P_1$, so that setting
$\delta = \beta\delta'$, we get (1).

\textit{Step 2.} We prove that $\delta \notin bR^*$. By contradiction, suppose that
$\delta = b\zeta$ with $\zeta \in R^*$; then $b - \delta = b(1 - \zeta) \in \mathfrak{q}_1$, and since
$b \notin P_1$ we have $1 - \zeta \in \mathfrak{q}_1 \subset mR^*$, so that $\zeta$ is a unit of $R^*$, and
$b\delta R^* \subset \mathfrak{q}_2$. This contradicts the fact that $b$ is not a zero-divisor of
$R^*$.

\textit{Step 3.} If $\mathfrak{b}$ is a non-zero ideal of $R$ such that $m$ is not a prime divisor of
$\mathfrak{b}$ then $\delta \notin \mathfrak{b}R^*$. Indeed, $\mathfrak{b} : m = \mathfrak{b}$, so that
$\mathfrak{b}R^* : mR^* = (\mathfrak{b} : m)R^* = \mathfrak{b}R^*$, and so $mR^*$ is not a prime divisor of
$\mathfrak{b}R^*$; if $P$ is any prime divisor of $\mathfrak{b}R^*$, then $P \ne mR^*$ and
$P \ne P_1$ (in view of $P_1 \cap R = 0$), so that $\coht P_1 = 1$ implies $P \not\supset P_1$,
hence $P \not\supset \mathfrak{q}_1$, and there exists $a \in \mathfrak{q}_1 - P$. If we write $\mathfrak{Q}$ for the
$P$-primary component of $\mathfrak{b}R^*$ then $a\delta = 0 \in \mathfrak{Q}$, so that
$\delta \in \mathfrak{Q}$. Thus finally $\delta \notin \mathfrak{b}R^*$.

\textit{Step 4.} By the previous two steps $m$ is a prime divisor of $bR$. Hence we can
write
\[
  bR = I \cap J \quad \text{with } I \text{ an } m\text{-primary ideal and } J : m = J,
\]
and then $bR^* = IR^* \cap JR^*$ with $\delta \in JR^*$ and $\delta \notin IR^*$. Moreover,
$(IR^* + JR^*)/IR^* \cong (I + J)/I$, so we can choose $c \in J$ such that
$\delta - c \in IR^*$. Then
\[
  \delta - c \in IR^* \cap JR^* = bR^*,
\]
so that $(b, c)R^* = (b, \delta)R^*$, and (2) is proved. If $c \in bR$ we would have
$(b, \delta)R^* = bR^*$, contradicting Step 2, so that $c/b \notin R$. On the other hand, we
have $b - \delta \in \mathfrak{q}_1$, so that $\delta(b - \delta) = 0$, that is $b\delta = \delta^2$, and
$c - \delta \in bR^*$. Set $c = \delta + b\eta$; then
$c^2 = \delta^2 + 2b\delta\eta + b^2\eta^2 \in b(\delta, b)R^* = b(c, b)R^*$, so that
\[
  c^2 \in (bc, b^2)R^* \cap R = (bc, b^2)R.
\]
From this, we get $c^2 = bcu + b^2v$ with $u, v \in R$, which proves that $c/b$ is integral
over $R$. Thus we have proved (3).
\end{proof}

\begin{lemma}\label{lem:31.5}
In the notation and assumptions of Lemma~\hyperref[lem:31.4]{4}, set $S = R[c/b]$; then $S$ has a
maximal ideal of height $1$.
\end{lemma}

\begin{proof}
Write $T$ for the total ring of fractions of $R^*$; then we can view $S^*$ as an
intermediate ring $R^* \subset S^* = R^*[c/b] \subset T$, and $T$ is the total ring of
fractions of $S^*$. In Lemma~\hyperref[lem:31.4]{4} we had $\Ass(R^*) = \{P_1, \ldots, P_r\}$, so that
setting $Q_i = P_iT \cap S^*$, we get
\[
  \Ass(S^*) = \{Q_1, \ldots, Q_r\} \quad \text{with } \htop Q_i = \htop P_i.
\]
Moreover, $S^*$ is integral over $R^*$, so that $S^*/Q_i$ is integral over $R^*/P_i$, and
hence also $\coht Q_i = \coht P_i$. Let $P^*$ be any maximal ideal of $S^*$ containing
$Q_1$. Then from $(b, c)R^* = (b, \delta)R^*$ we get $S^* = R^*[c/b] = R^*[\delta/b]$, and since
$\delta \in \mathfrak{q}_2 \cap \cdots \cap \mathfrak{q}_r$ and $\delta - b \in \mathfrak{q}_1$ we have
\[
  \delta/b \in Q_2 \cap \cdots \cap Q_r \quad \text{and} \quad \delta/b - 1 \in Q_1,
\]
so that $Q_i + Q_1 = S^*$ for all $i > 1$. Therefore $Q_1$ is the only minimal prime ideal
contained in $P^*$. However, $\coht Q_1 = 1$ and $P^* \cap R^* = mR^*$, so that
$\htop P^* = 1$. Setting $P = P^* \cap S$, we have $\htop P = \htop P^* = 1$, and $P$ is a maximal
ideal of $S$, since $P^*$ is a maximal ideal of $S^*$.
\end{proof}

Now we return to the proof of Theorem~\hyperref[thm:31.7]{7}. Let $A$ be a Noetherian local
integral domain, and suppose that $A^*$ is not equidimensional; then by
Lemma~\hyperref[lem:31.3]{3}, there exists a prime ideal $p$ of $A$ such that $(A/p)^* = A^*/pA^*$
has dimension $> 1$, and has a minimal prime ideal of coheight $1$. Set $R = A/p$;
then by Lemmas~\hyperref[lem:31.4]{4} and \hyperref[lem:31.5]{5} applied to $R$, there exists a subring $S$ of the
integral closure $R'$ of $R$ generated by one element, and having a maximal ideal
$P$ with $\htop P = 1 < \dim R$. Let $f\colon R[X] \to S$ be a surjective homomorphism of
$R$-algebras and let $\mathfrak{q}$ be its kernel. Set $Q = f^{-1}(P)$. Then $P = Q/\mathfrak{q}$ and
$Q \cap R = m$. Since $R \subset S$ we have $\mathfrak{q} \cap R = (0)$, hence
$\htop Q = \htop m + 1$ and $\htop \mathfrak{q} = 1$ by the remark after
Theorem~\hyperref[thm:15.5]{15.5}. Thus $\htop Q - \htop \mathfrak{q} = \htop m = \dim A/p > 1 = \htop(Q/\mathfrak{q})$,
so that $R[X]$ (and hence also $A[X]$) is not catenary.
\end{proof}

\begin{corollary*}
A Noetherian ring $A$ is universally catenary if and only if $A[X]$ is catenary.
\end{corollary*}

\begin{proof}
Suppose $A[X]$ is catenary and set $B = A[X_1, \ldots, X_n]$. In order to prove that
$B$ is catenary it suffices to prove that $B_P$ is catenary for every $P \in \Spec B$.
Let $p = P \cap A$. Then $B_P$ is a localisation of $A_p[X_1, \ldots, X_n]$. Since
$A_p[X_1]$ is catenary, $A_p[X_1, \ldots, X_n]$ is catenary by the theorem.
\end{proof}

\begin{corollary*}
A Noetherian ring of dimension $d$ is catenary if $d \le 2$ and is universally
catenary if $d \le 1$.
\end{corollary*}

\begin{proof}
The first assertion is obvious from the definitions and the second assertion
follows from the first because $\dim A[X] = d + 1$.
\end{proof}

% ============ §32 The formal fibre (原书 pp.255-260) ============
\section{The formal fibre}\label{sec:32}\index{formal fibre, section 32}\index{regular homomorphism, section 32}\index{G-ring, section 32}\index{excellent ring, section 32}\index{quasi-excellent ring, section 32}

Let $(A, m)$ be a Noetherian local ring and $A^*$ its completion. The fibre ring of
the natural homomorphism $A \to A^*$ over any $p \in \Spec A$ is called a \emph{formal
fibre} of $A$ (although strictly speaking we should distinguish between the fibre
and the fibre ring, we will not do so in what follows). If $I$ is an ideal of $A$
then $(A/I)^* = A^*/IA^*$, so that a formal fibre of $A/I$ is also a formal fibre of
$A$.

Let $A$ be a Noetherian ring and $k \subset A$ a subfield. We say that $A$ is
\emph{geometrically regular} over $k$ if $A \otimes_k k'$ is a regular ring for every finite
extension $k'$ of $k$ (see \S\ref{sec:28}). This is equivalent to saying that $A_p$ is
geometrically regular over $k$ for every maximal ideal $p$ of $A$.

We say that a homomorphism $\varphi\colon A \to B$ of Noetherian rings is \emph{regular} if
$\varphi$ is flat, and for every $p \in \Spec A$, the fibre $B \otimes_A \kappa(p)$ of $\varphi$ over $p$
is geometrically regular over the field $\kappa(p)$.

A Noetherian ring $A$ is said to be a \emph{G-ring} (here G stands for Grothendieck)
if $A_p \to (A_p)^*$ is regular for every prime ideal $p$ of $A$; this means that all the
formal fibres of all the local rings of $A$ are geometrically regular.

\begin{theorem}\label{thm:32.1}
Let $A \xrightarrow{\varphi} B \xrightarrow{\psi} C$ be homomorphisms of Noetherian rings;
then
\begin{enumerate}[(i)]
  \item if $\varphi$ and $\psi$ are regular then so is $\psi\varphi$;
  \item if $\psi\varphi$ is regular and $\psi$ is faithfully flat then $\varphi$ is also regular.
\end{enumerate}
\end{theorem}

\begin{proof}
(i) Clearly $\psi\varphi$ is flat. For $p \in \Spec A$, write $K = \kappa(p)$, and let $L$ be a
finite extension field of $K$. Set $B_L = B \otimes_A L$, and $C_L = C \otimes_A L$; then the
homomorphism $\psi_L\colon B_L \to C_L$ induced by $\psi$ is also regular. Indeed, if $P$ is a
prime ideal of $B_L$ then $C_L \otimes_{B_L} \kappa(P) = C \otimes_B (B \otimes_A L) = C \otimes_A
\kappa(P) = C_{\kappa(P)}$, and hence if $F$ is a finite extension of $\kappa(P)$ then
$C_L \otimes_{B_L} F = C \otimes_A F$; setting $P \cap B = Q$, since $B_L$ is a finite $B$-module we
have $[\kappa(P) : \kappa(Q)] < \infty$, and hence $[F : \kappa(Q)] < \infty$, so that
$C \otimes_A F$ is a regular ring. Now $\varphi$ is regular, so that $B_L$ is a regular ring,
and hence by Theorem~\hyperref[thm:23.7]{23.7}, (ii), $C_L$ is a regular ring.

(ii) The flatness of $\varphi$ is obvious. If we let $p$, $K$ and $L$ be as above, then
$C_L$ is a regular ring and is flat over $B_L$, so that by Theorem~\hyperref[thm:23.7]{23.7}, (i), $B_L$
is also regular.
\end{proof}

\begin{theorem}\label{thm:32.2}
Let $\varphi\colon A \to B$ be a homomorphism of Noetherian rings, and assume that
$\varphi$ is faithfully flat and regular.
\begin{enumerate}[(i)]
  \item $A$ is regular (or normal, reduced, CM, or Gorenstein) if and only if $B$ has
        the same property.
  \item If $B$ is a G-ring then so is $A$ (the converse is not true).
\end{enumerate}
\end{theorem}

\begin{proof}
(i) follows from Theorem~\hyperref[thm:23.7]{23.7}, the corollaries to
Theorems~\hyperref[thm:23.9]{23.9} and \hyperref[thm:23.3]{23.3}, and Theorem~\hyperref[thm:23.4]{23.4}.

(ii) Let $p \in \Spec A$, choose $P \in \Spec B$ lying over $p$, and consider the
commutative diagram
\[
\begin{tikzcd}
(A_p)^* \arrow[r, "f^*"] & (B_P)^* \\
A_p \arrow[u, "\alpha"] \arrow[r, "f"] & B_P \arrow[u, "\beta"]
\end{tikzcd}
\]
Here $f$ is the map induced by $\varphi$, and $f^*$ is the map induced by $f$; and the
vertical arrows are the natural maps. Now $f$ and $\beta$ are both regular, and $f^*$ is
faithfully flat, so that according to the previous theorem, $\alpha$ is also regular.

To construct an example where $A$ is a G-ring and $B$ is not, we let $A = k$ be a
perfect field, and $B$ a regular local ring containing $k$. Then $k \to B$ is certainly
faithfully flat and regular, and $k$ is a field, and so trivially a G-ring. However,
there are known examples in which $B$ is not a G-ring. (See the appendix to [N1];
a counter-example is provided by the ring $R$ in (E3.1) if $\chark k = p$, and by $R$ in
Example 7 if $\chark k = 0$. In (E3.1) the field $k$ is not perfect, but $R$ is
geometrically regular over $k$.)
\end{proof}

\begin{theorem}\label{thm:32.3}
A complete Noetherian local ring is a G-ring.
\end{theorem}

\begin{proof}
Let $A$ be a complete Noetherian local ring and $p \in \Spec A$; set $B = A_p$, and
let $B^*$ be the completion of $B$. We prove that $B \to B^*$ is a regular
homomorphism; that is, for any prime ideal $p'$ of $B$, we need to show that
$B^* \otimes_B \kappa(p')$ is geometrically regular over $\kappa(p')$. However, $A/p' \cap A$ is also a
complete local ring, so that we can replace $A$ by $A/p' \cap A$ and reduce to the
case $p' = (0)$. Thus assume that $A$ is an integral domain, and let $L$ be the
common field of fractions of $A$ and $B$; we must show that $B^* \otimes_B L$ is
geometrically regular over $L$.

The problem can be further reduced to the case when $A$ is a regular local ring.
In fact, by Theorem~\hyperref[thm:29.5]{29.5}, $A$ contains a complete regular local ring $R$ and
is a finite module over $R$. Set $p \cap R = \mathfrak{q}$, $R_\mathfrak{q} = S$ and
$B' = A_p = A \otimes_R S$; then $B'$ is a semilocal ring, and $B$ is a localisation of $B'$
at a maximal ideal, so that $B^*$ is one direct factor of $B'^* = B' \otimes_S S^*$. Write $K$
for the common field of fractions of $R$ and $S$.
\[
\begin{tikzcd}
B'^* = B' \otimes_S S^* \arrow[r] & B^* \\
B' = A \otimes_R S \arrow[u] \arrow[r] & B = A_p \arrow[u] \arrow[r] & L \\
R \arrow[u] \arrow[r] & S = R_\mathfrak{q} \arrow[u] \arrow[r] & K \arrow[u]
\end{tikzcd}
\]
Now $B^* \otimes_B L$ can be written as $B^* \otimes_S L$, and it is hence a direct factor
of $B'^* \otimes_S L = S^* \otimes_S L = (S^* \otimes_S K) \otimes_K L$, so that we need only show
that $S^* \otimes_S K$ is geometrically regular over $K$.

Now $R$, $S$ and $S^*$ are regular local rings, and $S^* \otimes_S K$ is a localisation of
$S^*$, so is a regular ring. Hence if $\chark K = 0$, there is nothing to prove. We
assume that $\chark K = p$ in what follows. Then $R$ has a coefficient field $k$, and
can be written $R = k[[X_1, \ldots, X_n]]$. Choose a directed family $\{k_\alpha\}$ of cofinite
subfields $k_\alpha \subset k$ such that $\bigcap_\alpha k_\alpha = k^p$, and set
$R_\alpha = k_\alpha[[X_1^p, \ldots, X_n^p]]$; write $K_\alpha$ for the field of fractions of
$R_\alpha$. Then one sees easily (compare the proof of Theorem~\hyperref[thm:30.9]{30.9}) that
$\bigcap_\alpha K_\alpha = K^p$.

We set $\mathfrak{q}_\alpha = \mathfrak{q} \cap R_\alpha$; then since $R_\alpha \supset R^p$, we see that
$\mathfrak{q}$ is the unique prime ideal of $R$ lying over $\mathfrak{q}_\alpha$. Hence if we let
$S_\alpha = (R_\alpha)_{\mathfrak{q}_\alpha}$ then $S = R_\mathfrak{q} = R \otimes_{R_\alpha} S_\alpha$, and $S$ is a
finite module over $S_\alpha$. Hence $S^* = S \otimes_{S_\alpha} S_\alpha^*$; let us prove that $S^*$
is 0-smooth over $S$ relative to $S_\alpha$ (see \S\ref{sec:28}). Suppose we are given a
commutative diagram of the form
\[
\begin{tikzcd}
S^* \arrow[r, "\bar{u}"] & C/N \\
S_\alpha \arrow[u] \arrow[r] & S \arrow[u] \arrow[r] & C \arrow[ul]
\end{tikzcd}
\]
where $C$ is a ring and $N$ is an ideal of $C$ with $N^2 = 0$. If there is a lifting
$u'\colon S^* \to C$ of $\bar{u}$ as a homomorphism of $S_\alpha$-algebras, set
$w = u'\big|_{S^*}$, and let $\bar{v} = u \otimes w\colon S^* = S \otimes_{S_\alpha} S_\alpha^* \to C$; then
one sees easily that $\bar{v}$ is a lifting of $\bar{u}$ over $S$. Hence $S^*$ is 0-smooth over
$S$ relative to $S_\alpha$. Now for $Q \in \Spec(S^*)$, let $Q \cap S = (0)$; then $(S^*)_Q$ is
a local ring of $S^* \otimes_S K$, and conversely, every local ring of $S^* \otimes_S K$ is of
this form. From the diagram
\[
\begin{tikzcd}
S \arrow[r] & K \arrow[r] & (S^*)_Q/N(S^*)_Q \\
S_\alpha \arrow[u] \arrow[r] & K_\alpha \arrow[u] \arrow[ru, ""'] & (S^*)_Q \arrow[u]
\end{tikzcd}
\]
one sees that $(S^*)_Q$ is 0-smooth over $K$ relative to $K_\alpha$. Set $E = (S^*)_Q$ and
$m = \rad(E)$; then $E$ is $m$-smooth over $K$ relative to $K_\alpha$, so that according
to Theorem~\hyperref[thm:28.4]{28.4},
\[
  \Omega_{K/K_\alpha} \otimes_K (E/m) \to \Omega_{E/K} \otimes_E (E/m)
\]
is injective for every $\alpha$. Moreover, since $\bigcap_\alpha K_\alpha = K^p$, by
\S\ref{sec:30}, Lemma~\hyperref[lem:30.4]{4},
\[
  \Omega_K \to \varprojlim \Omega_{K/K_\alpha}
\]
is injective, and hence
\[
  \Omega_K \otimes (E/m) \to \varprojlim (\Omega_{K/K_\alpha} \otimes (E/m))
\]
is also injective. Therefore, from the commutative diagram
\[
\begin{tikzcd}
\Omega_K \otimes (E/m) \arrow[r] & \Omega_E \otimes (E/m) \\
\varprojlim (\Omega_{K/K_\alpha} \otimes (E/m)) \arrow[u] \arrow[r] &
\varprojlim (\Omega_{E/K_\alpha} \otimes (E/m)) \arrow[u]
\end{tikzcd}
\]
we finally see that $\Omega_K \otimes (E/m) \to \Omega_E \otimes (E/m)$ is injective. Hence it
follows from the corollary to Theorem~\hyperref[thm:28.6]{28.6} that $E$ is $m$-smooth over $K$, and
thus is geometrically regular. Since $E$ is an arbitrary local ring of $S^* \otimes_S K$,
we see that $S^* \otimes_S K$ is geometrically regular over $K$.
\end{proof}

\begin{theorem}\label{thm:32.4}
Let $A$ be a Noetherian ring; if $A_m \to (A_m)^*$ is regular for every maximal ideal
$m$ of $A$, then $A$ is a G-ring.
\end{theorem}

\begin{proof}
Since $(A_m)^*$ is a G-ring, by Theorem~\hyperref[thm:32.2]{2}, $A_m$ is also a G-ring. For any
$p \in \Spec A$, if we let $m$ be a maximal ideal of $A$ containing $p$ then $A_p$ is a
localisation of the G-ring $A_m$, and hence $A_p \to (A_p)^*$ is regular.
\end{proof}

Theorem~\hyperref[thm:32.4]{4} makes it much easier to distinguish G-rings. For example, the
next theorem is based on Theorem~\hyperref[thm:32.4]{4}.

\begin{theorem}\label{thm:32.5}
Let $A$ be a Noetherian semilocal ring; then a sufficient condition for $A$ to be
a G-ring is that if $C$ is a finite $A$-algebra which is an integral domain, $m$ is
a maximal ideal of $C$ and we write $B = C_m$, then $(B^*)_Q$ is a regular local
ring for every $Q \in \Spec(B^*)$ such that $Q \cap B = (0)$.
\end{theorem}

\begin{remark}
It is easy to see that the condition is also necessary.
\end{remark}

\begin{proof}
By the previous theorem, we need only show that under the given condition,
$A \to A^*$ is regular. Let $p \in \Spec A$, and let $L$ be a finite extension of
$\kappa(p)$; we prove that $A^* \otimes_A L$ is regular. Suppose that
$L = \kappa(p)(\xi_1, \ldots, \xi_m)$; then multiplying each $\xi_i$ by an element of $A/p$ we
can assume that $\xi_i$ is integral over $A/p$, so that if we set
$C = (A/p)[\xi_1, \ldots, \xi_m]$, then $C$ is a finite $A$-module, and the field of
fractions of $C$ is $L$. Now $C^* = A^* \otimes_A C$, and if we write
$m_1, \ldots, m_s$ for the maximal ideals of $C$ and set $B_i = C_{m_i}$, then
$C^* = B_1^* \times \cdots \times B_s^*$. We can identify any local ring of
$A^* \otimes_A L = C^* \otimes_C L$ with the localisation $(B_i^*)_Q$ of one of the factors
$B_i^*$ at some prime ideal $Q$ of $B_i^*$ with $Q \cap B_i = (0)$, and by assumption
this is regular. Hence $A^* \otimes_A L$ is a regular ring.
\end{proof}

\begin{theorem}[\textbf{H.~Mizutani}]\label{thm:32.6}
Let $R$ be a regular ring. If the weak Jacobian condition (WJ) of \S\ref{sec:30} holds for
$R[X_1, \ldots, X_n]$ for every $n \ge 0$, then $R$ is a G-ring.
\end{theorem}

\begin{proof}
Since (WJ) is inherited by any localisation we can assume that $R$ is local. We
prove that the condition of Theorem~\hyperref[thm:32.5]{5} holds. Set
$R_n = R[X_1, \ldots, X_n]$. Any integral domain $C$ which is finite as an $R$-module
can be expressed as $C = R_n/Q$ with $Q \in \Spec(R_n)$ for some $n$. Let $m$ be a
maximal ideal of $C$, $M$ the maximal ideal of $R_n$ corresponding to $m$, and
$S = (R_n)_M$; then it is enough to show that $(S^*)_P/Q(S^*)_P$ is regular for every
$P \in \Spec(S^*)$ with $P \cap S = QS$. If $\htop Q = r$ then $\htop Q(S^*)_P = r$, and by
assumption there exist $D_1, \ldots, D_r \in \Der(R_n)$ and
$f_1, \ldots, f_r \in Q$ such that $\det(D_if_j) \notin Q$. Now $D_i$ has a natural extension
to $S$, and then to $S^*$, and since $P \cap R_n = Q$, we have
$\det(D_if_j) \notin P$, so that by Theorem~\hyperref[thm:30.4]{30.4}, $(S^*)_P/Q(S^*)_P$ is
regular.
\end{proof}

\begin{corollary*}
A ring which is finitely generated over a field, or a localisation of such a
ring, is a G-ring.
\end{corollary*}

\begin{proof}
It follows from the definition that a quotient or localisation of a G-ring is
again a G-ring, so that we need only show that for a field $k$, the ring
$k[X_1, \ldots, X_n]$ is a G-ring; but by Theorems~\hyperref[thm:30.3]{30.3} and \hyperref[thm:30.5]{30.5}, (WJ)
holds in $k[X_1, \ldots, X_{n+m}]$, so that $k[X_1, \ldots, X_n]$ is a G-ring by the
theorem.
\end{proof}

\begin{remark}
The local rings which appear in algebraic geometry are essentially of finite
type over a field, and therefore G-rings. Hence for these rings, properties such
as reduced and normal pass to the completion.
\end{remark}

\begin{remark}
If $R$ is a regular ring containing a field of characteristic $0$, and (WJ) holds
in $R$, then by Ex.~\hyperref[ex:30.4-ex]{30.4}, it also holds in $R[X_1, \ldots, X_n]$. Hence $R$ and
$R[X_1, \ldots, X_n]$ are G-rings. In particular by Theorem~\hyperref[thm:30.8]{30.8}, rings of
convergent power series over $\mathbb{R}$ and $\mathbb{C}$ are G-rings.
\end{remark}

A theorem proved by Grothendieck asserts that if $A$ is a G-ring, then so is
$A[X]$; the proof is very hard, and we omit it, referring only to [M], Theorem 77.
The analogous statement for $A[[X]]$ remained unsolved for a long time, but was
recently proved for a semilocal ring $A$ by C.~Rotthaus [3]. In the non-semilocal
case she proved in [4] that, if $A$ is a finite-dimensional excellent ring
containing the rational numbers, then $A[[X]]$ is excellent. On the other hand,
Nishimura [3] showed that there exists a G-ring $A$ such that $A[[X]]$ is not a
G-ring.

Nagata [8] studied the condition that $\Reg(A)$ is open in $\Spec A$; putting
together Nagata's work with his own theory of G-rings, Grothendieck gave the
definition of excellent ring in [G2].

\begin{definition*}
A Noetherian ring $A$ is \emph{excellent} if it satisfies the following three conditions:
\begin{enumerate}[(1)]
  \item $A$ is universally catenary;
  \item $A$ is a G-ring;
  \item $\Reg(B) \subset \Spec B$ is open for every finitely generated $A$-algebra $B$.
\end{enumerate}
A Noetherian ring satisfying (2) and (3) is said to be \emph{quasi-excellent}.
\end{definition*}

One can prove that the classes of rings satisfying each of (1), (2) and (3) are
closed under localisation, finitely generated extensions and passing to
quotients. It can also be proved that (2) implies (3) for semilocal Noetherian
rings. A complete Noetherian local ring is excellent, as are practically all
Noetherian rings in applications. For more information on excellent rings, see
[M], Ch.~13, or [G2].

R.~Y.~Sharp [S] defined the notion of an \emph{acceptable ring}, replacing condition
(2) by the condition that all formal fibres of all localisations of $A$ are
Gorenstein, and replacing $\Reg(B)$ by $\Gor(B)$ in (3), and showed that the
resulting theory is analogous to the theory of excellent rings (see also
Greco-Marinari [1], Sharp [6]).

Using his cohomology theory, M.~Andr\'e [1] proved the following theorem. Let
$A$, $B$ be Noetherian local rings, and $\varphi\colon A \to B$ a local homomorphism;
suppose that $A$ is quasi-excellent and $B$ is $m_B$-smooth over $A$, where
$m_B = \rad(B)$; then $\varphi$ is regular. This is an extremely strong theorem; the
result of Rotthaus mentioned above also makes use of this.

% ============ §33 Some other applications (原书 pp.261-265) ============
\section{Some other applications}\label{sec:33}\index{dimension of intersection, section 33}\index{integral closure, finiteness of, section 33}\index{Nagata ring, section 33}\index{Mori-Nagata integral closure theorem, section 33}

\medskip
\noindent\textbf{Dimension of intersection}

Let $k$ be a field and $V$ and $W$ irreducible algebraic varieties in affine
$n$-space over $k$. (Here one may either assume that $k$ is algebraically closed
and identify the varieties with the corresponding subsets of $k^n$, or take the
scheme-theoretic viewpoint.) Then it is well known that every irreducible
component of $V \cap W$ has dimension $\ge \dim V + \dim W - n$. Algebraically, this
is equivalent to the following theorem.

Let $P$ and $P'$ be two prime ideals in the polynomial ring $R = k[X_1, \ldots, X_n]$
over a field $k$, and let $Q$ be a minimal prime divisor of $P + P'$. Then
\[
  \dim R/Q \ge \dim(R/P) + \dim(R/P') - n,
\]
or equivalently,
\[
  (*) \quad \htop Q \le \htop P + \htop P'.
\]
The idea of the proof consists of transforming the intersection $V \cap W$ in $k^n$
into the intersection $\Delta \cap (V \times W)$ in $k^{2n}$, where $\Delta$ is the diagonal, and
availing oneself of the fact that $\Delta$ is defined by $n$ equations
$x_i - y_i = 0$, for $i = 1, \ldots, n$ (see [M], p.~93).

Now in algebraic geometry, the theorem remains true if one replaces affine
$n$-space by a non-singular (smooth) variety; namely, if $V$ and $W$ are
irreducible subvarieties of a non-singular variety $U$, then every irreducible
component of $V \cap W$ has dimension $\ge \dim V + \dim W - \dim U$. Algebraically,
the inequality (*) still holds if $R$ is an arbitrary regular local ring containing
a field. One can easily reduce to the case where $R$ is complete, and then by
I.~S.~Cohen's structure theorem $R$ is isomorphic to $k[[X_1, \ldots, X_n]]$, and so one
can apply the same diagonal trick. Thus it is clear that (*) holds in an
arbitrary regular local ring of equal characteristic. How about the unequal
characteristic case? If $R$ is unramified, then its completion $R^*$ is a formal
power series ring over a DVR by Theorem~\hyperref[thm:29.8]{29.8}, and a slight modification of the
argument used in the case of $k[X_1, \ldots, X_n]$ suffices. When $R$ is ramified, by
Theorem~\hyperref[thm:29.9]{29.9}, $R^*$ is of the form $D[X_1, \ldots, X_n]/(f)$, where $D$ is a complete
DVR and $f$ is an Eisenstein polynomial. Using this, and applying his deep
results on intersection multiplicity, J.-P.~Serre proved the inequality (*) for
general regular local rings $R$ in Chapter V of his book [Se]. We recommend this
excellent book to the reader.

\medskip
\noindent\textbf{Integral closure of a Noetherian integral domain}

Let $A$ be a Noetherian integral domain with field of fractions $K$, and let $A'$
denote the integral closure of $A$ in $K$ (the so-called \emph{derived normal ring} of
$A$). Is $A'$ a finite module over $A$? This is a difficult question, and the answer
is no in general. When $A$ is finitely generated over a field $k$ (the case
encountered in algebraic geometry) it is easy to prove finiteness. We need the
following two lemmas.

\begin{lemma}\label{lem:33.1}
Let $A$ be a Noetherian normal integral domain with field of fractions $K$;
suppose that $L$ is a finite separable extension of $K$, and let $A'$ be the
integral closure of $A$ in $L$. Then $A'$ is a finite $A$-module.
\end{lemma}

\begin{proof}
By enlarging $L$ if necessary we can assume that $L$ is a Galois extension of $K$.
Write $G = \{\sigma_i \mid 1 \le i \le n\}$ for the Galois group of $L/K$, where
$n = [L : K]$, and let $y_1, \ldots, y_n$ be elements of $A'$ which form a basis of $L$
over $K$. If $z \in A'$ and $z = \sum_j c_jy_j$ with $c_j \in K$, then
$\sigma_iz = \sum_j c_j\sigma_iy_j$ for $i = 1, \ldots, n$, and hence
$c_j = C_j/D$, where $D = \det(\sigma_iy_j)$ and $C_j \in A'$. Putting $d = D^2$ we see
$d \in K$. In fact it is easy to see that $d = \det(\tr_{L/K}(y_iy_j))$ is the
discriminant of the separable $K$-algebra $L$ (see p.~198). It follows that
$d \ne 0$ and $dc_j \in A' \cap K = A$ for all $j$. Therefore $A'$ is contained in the
finite $A$-module $\sum_j Ad^{-1}y_j$, so that $A'$ itself is finite over $A$.
\end{proof}

\begin{lemma}[\textbf{Normalisation theorem of E.~Noether}]\label{lem:33.2}
Let $A = k[x_1, \ldots, x_n]$ be a finitely generated algebra over a field $k$; then
there exist $y_1, \ldots, y_r \in A$ which are algebraically independent over $k$ such
that $A$ is integral over $k[y_1, \ldots, y_r]$.
\end{lemma}

\begin{proof}
Here we assume that $k$ is an infinite field, referring the reader to [M], (14.G)
or [N1], (14.4) for the general case. Suppose $x_1, \ldots, x_n$ are algebraically
dependent over $k$, and let $f(x_1, \ldots, x_n) = 0$ be a relation, where
$f(X_1, \ldots, X_n)$ is a non-zero polynomial with coefficients in $k$. Write $d$ for
the degree of $f$ and let $f_d(X_1, \ldots, X_n)$ be the homogeneous part of $f$ of
degree $d$. Take $c_1, \ldots, c_{n-1} \in k$ such that
$f_d(c_1, \ldots, c_{n-1}, 1) \ne 0$, and set $y_i = x_i - c_ix_n$ for $i = 1, \ldots, n - 1$. Then
\[
  0 = f(x_1, \ldots, x_n) = f(y_1 + c_1x_n, \ldots, y_{n-1} + c_{n-1}x_n, x_n)
      = f_d(c_1, \ldots, c_{n-1}, 1)x_n^d + g_1x_n^{d-1} + \cdots + g_d,
\]
with $g_i \in k[y_1, \ldots, y_{n-1}]$, so that $x_n$ is integral over
$k[y_1, \ldots, y_{n-1}]$. Then $x_i = y_i + c_ix_n$, for $i = 1, \ldots, n - 1$, are also
integral over $k[y_1, \ldots, y_{n-1}]$, hence $A$ is integral over
$k[y_1, \ldots, y_{n-1}]$. Thus the assertion is proved by induction on $n$.
\end{proof}

Now let $A$ be a finitely generated integral domain over $k$, with field of
fractions $K$ and derived normal ring $A'$. Take $y_1, \ldots, y_r \in A$ as in
Lemma~\hyperref[lem:33.2]{2}, so that $A'$ is also the integral closure of $k[y_1, \ldots, y_r]$ in $K$.
Set $K' = k(y_1, \ldots, y_r)$. Then $K$ is a finite algebraic extension of $K'$. If this
extension is separable then $A'$ is finite over $k[y_1, \ldots, y_r]$ by
Lemma~\hyperref[lem:33.1]{1}, hence is also finite over $A$. If $K$ is inseparable over $K'$, let
$p = \chark K$. Then there is a purely inseparable finite extension $K''$ of $K'$ such
that $K(K'')$ is separable over $K''$. Therefore it suffices to prove that the
integral closure of $k[y_1, \ldots, y_r]$ in $K''$ is finite over it. But $K''$ is
contained in a field $L$ which is obtained by adjoining to $K'$ the $q$-th roots of
a finite number of elements $a_1, \ldots, a_s$ of $k$ and also the $q$-th roots of
$y_1, \ldots, y_r$, where $q$ is a sufficiently high power of $p$. Then the integral
closure of $k[y_1, \ldots, y_r]$ in $L$ is $k'[y_1^{1/q}, \ldots, y_r^{1/q}]$, where
$k' = k(a_1^{1/q}, \ldots, a_s^{1/q})$, and it is clear that $k'[y_1^{1/q}, \ldots, y_r^{1/q}]$ is
finite over $k[y_1, \ldots, y_r]$. This completes the proof of finiteness of $A'$ over
$A$.

When $A$ is a complete Noetherian local domain one can prove the finiteness of
$A'$ along the same line as above. Using Theorem~\hyperref[thm:29.5]{29.5}, (iii), instead of
Lemma~\hyperref[lem:33.2]{2}, one reduces to proving the finiteness of the integral closure of a
complete regular local ring $A$ in a finite extension $L$ of the field of fractions
$K$ of $A$. If $\chark K = 0$ this is proved by Lemma~\hyperref[lem:33.1]{1}, so that we can assume
$\chark K = p > 0$. Then $A = k[[X_1, \ldots, X_n]]$ is a formal power series ring over a
field $k$. We can also assume, as in the above proof, that $L$ is purely
inseparable over $K$, so that there is a power $q = p^m$ of $p$ such that
$L \subset K^{1/q}$. But there is one problem. Since a formal power series has
infinitely many coefficients, it may not be possible to find a finite extension
$k_1$ of $k$ such that $L \subset k_1((Y))$, where $Y = (Y_1, \ldots, Y_n)$,
$Y_i = X_i^{1/q}$ and $k_1((Y))$ denotes the field of fractions of $k_1[[Y]]$. One can
overcome this difficulty either by an argument of Nagata in [N1, p.~113], or
(following J.~Tate) by induction on $n$ as follows:

We may assume that $x = X_1^{1/q} \in L$ ($1 \le i \le n$). Since $A$ is normal we have
$A' = \{f \in L \mid f^q \in A\}$. Set $P = X_1A$, $Q = Y_1A'$; then
$Q = (\xi \in L \mid \xi^q \in P)$, so that $Q$ is the only prime ideal lying over $P$. Now
$A_P$ and $A'_Q$ are DVRs by Theorem~\hyperref[thm:11.2]{11.2}, (3), and their fields of fractions
are $L$ and $K$ respectively. Let $K'$ and $K_0$ be their residue fields; then
$[K' : K_0] \le [L : K]$ by Ex.~\hyperref[ex:10.8-ex]{10.8}. Since $A'/Q$ is contained in the integral
closure of $A/P = k[[X_2, \ldots, X_n]]$ in $K'$, by the induction hypothesis, $A'/Q$ is
finite over $A/P$. Since
\[
  Q^i/Q^{i+1} = Y_1^iA'/Y_1^{i+1}A' \cong A'/Q \quad \text{and} \quad Q^q = PA'
\]
we see that $A'/PA'$ is finite over $A/P$. Moreover, $A'$ is separated in the $P$-adic
topology (which is the same as the $Y_1$-adic topology), because
$A' \subset k^{1/q}[[Y_1, \ldots, Y_n]]$. Since $A$ is $P$-adically complete, $A'$ is finite over
$A$ by Theorem~\hyperref[thm:8.4]{8.4}.

From this result it is easy to derive the following theorem: \emph{If $A$ is a Noetherian
local ring whose completion $A^*$ is reduced, then the integral closure $A'$ of $A$ in its
total ring of fractions is finite over $A$}. See [M] p.~237.

On the other hand, if $A$ is not reduced and if the maximal ideal $m$ contains a
regular element (that is, a non-zero-divisor), then $A'$ is not finite over $A$
(Krull [2]). In fact, if $x \ne 0$ is nilpotent, take a regular element $s$ such that
$x \notin sA$; (it is possible to find such an $s$ since $\bigcap_i m^i = (0)$). Then the
elements $x/s^j$ (for $j = 1, 2, 3, \ldots$) belong to $A'$. If $A'$ is finite over $A$ then
there must be some integer $n$ such that $s^n(x/s^j) \in A$ for all $j$. But then
$x \in \bigcap_{j \ge 0}s^jA = (0)$, a contradiction.

Suppose $(A, m)$ is a one-dimensional Noetherian local integral domain. Then $A^*$
is reduced if and only if $A'$ is finite over $A$ (Krull [2]). In fact, if $A'$ is finite
over $A$ then it is a semilocal ring, and for each maximal ideal $P$ of $A'$ the
local ring $(A')_P$ is a DVR. The $\rad(A')$-adic topology of $A'$ coincides with the
$m$-adic topology, and the completion $A'^*$ of $A'$ with respect to this topology is
a direct product of complete DVRs by Theorem~\hyperref[thm:8.15]{8.15}, hence is reduced. On
the other hand it coincides with $A^* \otimes_A A'$ by Theorem~\hyperref[thm:8.7]{8.7}. Since
$0 \to A \to A'$ is exact, $0 \to A^* \to A^* \otimes_A A'$ is also exact. Therefore $A^*$ is
reduced. The converse holds, as already mentioned, without the restriction on
dimension.

Rees [9] proved that a reduced Noetherian local ring $A$ has reduced completion
$A^*$ if and only if for every finite subset $T$ of the total ring of fractions $K$
of $A$, the integral closure of $A[t]$ in $K$ is finite over $A[t]$.

Akizuki [1] constructed the first example of a one-dimensional Noetherian local
integral domain with non-reduced completion (see also Larfeldt-Lech [1]). To
avoid such pathology, Nagata [N1] defined and studied the class of
pseudo-geometric rings, which were called `anneaux universellement japonais'
by Grothendieck ([G1], [G2]). These are now known as `Nagata rings' ([M],
[B9]). A Noetherian ring $A$ is called a \emph{Nagata ring}\index{Nagata ring} if for every prime
ideal $P$ of $A$ and for every finite extension field $L$ of the field of fractions
$\kappa(P)$ of $A/P$, the integral closure of $A/P$ in $L$ is finite over $A/P$. For the
basic properties of Nagata rings, see [M], \S31. An alternative definition is the
following: a Noetherian ring $A$ is a Nagata ring if (1) for every maximal ideal
$m$ the formal fibres of $A_m$ are geometrically reduced, and (2) for every finite
$A$-algebra $B$ which is an integral domain, the set $\Nor(B) = \{P \in \Spec B \mid B_P
\text{ is normal}\}$ is open in $\Spec B$. The equivalence of these two definitions
can easily be proved from [G2], (7.6.4) and (7.7.2).

Although the integral closure $A'$ of a Noetherian integral domain $A$ may fail to
be finite over $A$, it is a Krull ring by the theorem of Mori-Nagata mentioned in
\S\ref{sec:12}. Because of its importance we quote here the theorem in full.

\begin{theorem}[\textbf{Mori-Nagata integral closure theorem}]\label{thm:33.1}
Let $A$ be a Noetherian integral domain and let $A'$ be its derived normal ring.
Then (1) $A'$ is a Krull ring, and (2) for every prime ideal $P$ of $A$ there are
only finitely many prime ideals $P'$ of $A'$ lying over $P$, and for each such $P'$
the field of fractions $\kappa(P')$ of $A'/P'$ is finite over $\kappa(P)$.
\end{theorem}

For a proof, see [N1], (33.10). This proof depends on I.~S.~Cohen's structure
theorem. There are also more recent proofs which do not use the structure
theorem (Nishimura [2], Querr\'e [1], Kiyek [1]).

Note that $A'$ is Noetherian if $\dim A \le 2$. This follows easily from the above
theorem, Theorem~\hyperref[thm:11.7]{11.7} (Krull-Akizuki) and Theorem~\hyperref[thm:12.7]{12.7} (Mori-Nishimura).
When $\dim A = 3$, Nagata constructed a counter-example ([N1], p.~207).

Theorem~\hyperref[thm:28.9]{28.9}, which is due to Grothendieck and is not proved in this book,
was given a new proof by Radu [5]. This interesting proof depends heavily on
I.~S.~Cohen's structure theorem. The same remark applies also to Andr\'e's proof
[1] of the theorem mentioned at the end of \S\ref{sec:32}.

I.~S.~Cohen's structure theorem is also at the basis of the theories of canonical
modules ([HK]) and of dualising complexes (Sharp [3], [5]). Here, the fact that
a complete Noetherian local ring is a quotient of a Gorenstein ring is
important.



% ============ 附录 (原书 Appendices A--C, pp.266-286) ============
% ============ Appendix A. Tensor products, direct and inverse limits (原书 pp.266-273) ============
\chapter*{Appendix A}
\addcontentsline{toc}{chapter}{Appendix A: Tensor products, direct and inverse limits}
\markboth{Appendix A}{Tensor products, direct and inverse limits}
\renewcommand{\thetheorem}{A\arabic{theorem}}
\setcounter{theorem}{0}

\vspace*{-0.5cm}
\begin{center}
  {\Large\bfseries Tensor products, direct and inverse limits}
\end{center}
\vspace{0.3cm}

\medskip
\noindent\textbf{Tensor products}

Let $A$ be a ring, $L$, $M$ and $N$ three $A$-modules. We say that a map
$\varphi\colon M \times N \to L$ is \emph{bilinear} if fixing either of the entries it is $A$-linear
in the other, that is if
\[
  \varphi(x + x', y) = \varphi(x, y) + \varphi(x', y), \quad \varphi(ax, y) = a\varphi(x, y),
\]
\[
  \varphi(x, y + y') = \varphi(x, y) + \varphi(x, y'), \quad \varphi(x, ay) = a\varphi(x, y).
\]
Write $\mathcal{B}(M, N; L)$ or $\mathcal{B}_A(M, N; L)$ for the set of all bilinear maps from
$M \times N$ to $L$; as with $\Hom(M, L)$, this has an $A$-module structure (since we are
assuming that $A$ is commutative).

If $g\colon L \to L'$ is an $A$-linear map and $\varphi \in \mathcal{B}(M, N; L)$ then
$g\varphi \in \mathcal{B}(M, N; L')$. Bearing this in mind, for given $M$ and $N$, consider a
bilinear map $\theta\colon M \times N \to L_0$ having the following property, where we write
$x \otimes y$ instead of $\theta(x, y)$: for any $A$-module $L$ and any $\varphi \in \mathcal{B}(M, N; L)$
there exists a unique $A$-linear map $g\colon L_0 \to L$ satisfying
\[
  \varphi(x \otimes y) = g(x \otimes y).
\]
If this holds we say that $L_0$ is the \emph{tensor product} of $M$ and $N$ over $A$, and
write $L_0 = M \otimes_A N$; we sometimes omit $A$ and write $M \otimes N$. As usual with
this kind of definition, $M \otimes_A N$, assuming it exists, is uniquely determined
up to isomorphism. To prove existence, write $F$ for the free $A$-module with
basis the set $M \times N$, and let $R \subset F$ be the submodule generated by all
elements of the form
\[
  (x + x', y) - (x, y) - (x', y), \quad (ax, y) - a(x, y)
\]
\[
  (x, y + y') - (x, y) - (x, y'), \quad (x, ay) - a(x, y).
\]
Then set $L_0 = F/R$ and write $x \otimes y$ for the image in $L_0$ of $(x, y) \in F$. It is now
easy to check that $L_0$ and $\otimes$ satisfy the above condition.

Note that the general element of $M \otimes_A N$ is a sum of the form $\sum x_i \otimes y_i$,
and cannot necessarily be written $x \otimes y$.

For $A$-modules $M$, $N$ and $L$ the definition of tensor product gives:

\textbf{Formula 1.} $\Hom_A(M \otimes_A N, L) \cong \mathcal{B}(M, N; L)$.

The canonical isomorphism is obtained by taking an element $\varphi$ of the
right-hand side to the element $g$ of the left-hand side satisfying
$g(x \otimes y) = \varphi(x, y)$.

We can define multilinear maps from an $r$-fold product of $A$-modules
$M_1, \ldots, M_r$ to an $A$-module $L$ just as in the bilinear case, and get modules
$\mathcal{B}(M_1, \ldots, M_r; L)$ and $M_1 \otimes_A \cdots \otimes_A M_r$; the following `associative
law' then holds:

\textbf{Formula 2.} $(M \otimes_A M') \otimes_A M'' = M \otimes_A (M' \otimes_A M'')$.

For example, for the first equality it is enough to check that the trilinear map
$M \times M' \times M'' \to (M \otimes_A M') \otimes_A M''$ given by
$(x, y, z) \mapsto (x \otimes y) \otimes z$ has the universal property for trilinear maps, and this
is easy. The following Formulas 3, 4 and 5 are also easy:

\textbf{Formula 3.} $M \otimes_A N \cong N \otimes_A M$ \quad (by $x \otimes y \mapsto y \otimes x$).

\textbf{Formula 4.} $M \otimes_A A = M$.

\textbf{Formula 5.} $(\bigoplus_\lambda M_\lambda) \otimes_A N = \bigoplus_\lambda(M_\lambda \otimes_A N)$.

If $f\colon M \to M'$ and $g\colon N \to N'$ are both $A$-linear then
$(x, y) \mapsto f(x) \otimes g(y)$ is a bilinear map from $M \times N$ to $M' \otimes_A N'$, and so it
defines a linear map $M \otimes_A N \to M' \otimes_A N'$, which we denote $f \otimes g$. From
the definition we have:

\textbf{Formula 6.} $(f \otimes g)(\sum x_i \otimes y_i) = \sum f(x_i) \otimes g(y_i)$.

In particular, if both $f$ and $g$ are surjective then we see from this that
$f \otimes g$ is surjective; its kernel is generated by
$\{x \otimes y \mid f(x) = 0 \text{ or } g(y) = 0\}$. Indeed, let $T \subset M \otimes N$ be the submodule
generated by this set; then $T \subset \Ker(f \otimes g)$ so that $f \otimes g$ induces a linear
map $\alpha\colon (M \otimes N)/T \to M' \otimes N'$; furthermore, we can define a bilinear map
$M' \times N' \to (M \otimes N)/T$ by
\[
  (x', y') \mapsto (x \otimes y \bmod T), \quad \text{where } f(x) = x',\ g(y) = y',
\]
since a different choice of inverse images $x$ and $y$ leads to a difference
belonging to $T$. This defines a linear map $\beta\colon M' \otimes N' \to (M \otimes N)/T$,
which is obviously an inverse of $\alpha$. We summarise the above (writing $1$ for
the identity maps):

\textbf{Formula 7.} Suppose given exact sequences
\[
  0 \to K \to M \xrightarrow{f} M' \to 0 \quad \text{and} \quad 0 \to L \xrightarrow{j} N \to N' \to 0,
\]
then $M' \otimes N' \cong (M \otimes N)/T$, where
\[
  T = (i \otimes 1)(K \otimes N) + (1 \otimes j)(M \otimes L).
\]

\textbf{Formula 8 (right-exactness of the tensor product).} If
\[
  M_1 \xrightarrow{f} M_2 \xrightarrow{g} M_3 \to 0
\]
is an exact sequence then so is
\[
  M_1 \otimes_A N \xrightarrow{f \otimes 1} M_2 \otimes_A N \xrightarrow{g \otimes 1} M_3 \otimes_A N \to 0.
\]

In general, even if $f\colon M \to M'$ is injective, $f \otimes 1\colon M \otimes N \to M' \otimes N$ need
not be. (Counter-example: let $A = \mathbb{Z}$ and $N = \mathbb{Z}/n\mathbb{Z}$ for some $n > 1$. Let
$f\colon \mathbb{Z} \to \mathbb{Z}$ be multiplication by $n$; then $\mathbb{Z} \otimes_A N \cong N \ne 0$, but
$f \otimes 1\colon N \to N$ is the zero map, and so is not injective.) However, if $\Im f$ is a
direct summand of $M'$ (in which case we say that the exact sequence
$0 \to M \to M' \to M'/M \to 0$ \emph{splits}), then there is a map $g\colon M' \to M$ such that
$gf = 1$. Then
\[
  (g \otimes 1)(f \otimes 1) = gf \otimes 1 = 1 \otimes 1
\]
is the identity map of $M \otimes N$, and hence $f \otimes 1$ is injective, and the sequence
$0 \to M \otimes N \to M' \otimes N \to (M'/M) \otimes N \to 0$ is split. In particular if $A$ is a
field then any submodule is a direct summand, so that the operation
$\otimes N$ takes exact sequences into exact sequences; in other word, $\otimes N$ is
an exact functor. For an arbitrary ring $A$, an $A$-module $N$ is said to be
\emph{flat} if $\otimes N$ is an exact functor. For more on this see \S\ref{sec:7}.

\medskip
\noindent\textbf{Change of coefficient ring}

Let $A$ and $B$ be rings, and $P$ a two-sided $A$-$B$-module; that is, for $a \in A$,
$b \in B$ and $x \in P$ the products $ax$ and $xb$ are defined, and in addition to the
usual conditions for $A$-modules and $B$-modules we assume that
\[
  (ax)b = a(xb).
\]
Then multiplication by an element $b \in B$ induces an $A$-linear map of $P$ to
itself, which we continue to denote by $b$. This determines a map
$1 \otimes b\colon M \otimes_A P \to M \otimes_A P$ for any $A$-module $M$, and by definition we take
this to be scalar multiplication by $b$ in $M \otimes_A P$; that is, we set
\[
  (\sum y_i \otimes x_i)b = \sum y_i \otimes x_ib \quad \text{for } y_i \in M \text{ and } x_i \in P.
\]
If $N$ is a $B$-module, then for $\varphi \in \Hom_B(P, N)$ we define the product
$\varphi a$ of $\varphi$ and $a \in A$ by
\[
  (\varphi a)(x) = \varphi(ax) \quad \text{for } x \in P;
\]
we have $\varphi a \in \Hom_B(P, N)$, and this makes $\Hom_B(P, N)$ into an $A$-module.

\textbf{Formula 9.} $\Hom_B(M, \Hom_B(P, N)) \cong \Hom_A(M \otimes_A P, N)$.

\textbf{Formula 10.} $(M \otimes_A P) \otimes_B N \cong M \otimes_A (P \otimes_B N)$.

Both of these are easy to prove, and we leave them to the reader.
Formula 10 generalises Formula 2.

Given a ring homomorphism $\lambda\colon A \to B$, we can think of $B$ as a two-sided
$A$-$B$-module by setting $ab = \lambda(a)b$; then for any $A$-module $M$, $M \otimes_A B$ is a
$B$-module, called the \emph{extension of scalars} in $M$ from $A$ to $B$, and written
$M_{(B)}$. For $A$-modules $M$ and $M'$ the following formula holds, so that
tensor product commutes with change of scalars.

\textbf{Formula 11.} $(M \otimes_A B) \otimes_B (M' \otimes_A B) = (M \otimes_A M') \otimes_A B$.

Indeed, using Formulas 10, 4 and 2, the left-hand side is equal to
$M \otimes_A (B \otimes_B (M' \otimes_A B)) = M \otimes_A (M' \otimes_A B) = (M \otimes_A M') \otimes_A B$.

\medskip
\noindent\textbf{Tensor product of $A$-algebras}

Given a ring homomorphism $\lambda\colon A \to B$ we say that $B$ is an \emph{$A$-algebra}. Let
$B'$ be another $A$-algebra defined by $\lambda'\colon A \to B'$. We say that a map
$f\colon B \to B'$ is a homomorphism of $A$-algebras if it is a ring homomorphism
satisfying $\lambda' = f\lambda$. If $B$ and $C$ are $A$-algebras, then we can take the tensor
product $B \otimes_A C$ of $B$ and $C$ as $A$-modules and this is again an $A$-algebra.
That is, we define the product by
\[
  \Big(\sum b_i \otimes c_i\Big)\Big(\sum b_j' \otimes c_j'\Big) = \sum b_ib_j' \otimes c_ic_j',
\]
and the ring homomorphism $A \to B \otimes_A C$ by $a \mapsto a \otimes 1 (= 1 \otimes a)$. The
fact that the above product is well-defined can easily be seen using the
bilinearity of $bb' \otimes cc'$ with respect to both $(b, c)$ and $(b', c')$. The algebra
$B \otimes_A C$ contains $B \otimes 1$ (short for the subset $\{b \otimes 1 \mid b \in B\} \subset B \otimes C$)
and $1 \otimes C$ as subalgebras, and is generated by these. Note that $B \otimes 1$ is not
necessarily isomorphic to $B$.

\begin{example*}
If $\mathfrak{a}$ is an ideal of $A$ and $C = A/\mathfrak{a}$, then
$B \otimes_A C = B \otimes_A (A/\mathfrak{a}) = B/\mathfrak{a}B$, and the above $B \otimes 1$ is also equal to
$B/\mathfrak{a}B$.
\end{example*}

\begin{example*}
If $B$ is an $A$-algebra and $A[X]$ is the polynomial ring over $A$ then
$B \otimes_A A[X]$ can be identified with $B[X]$. Indeed, $A[X]$ is a free $A$-module
with basis $\{X^\nu \mid \nu = 0, 1, 2, \ldots\}$, so that $B \otimes_A A[X]$ is also the free $B$-module
with basis $\{X^\nu\}$, and is isomorphic to $B[X]$ both as an $A$-module and as a
ring. Similarly for the polynomial ring in several variables.
\end{example*}

\medskip
\noindent\textbf{Direct limits}

A \emph{directed set} is a partially ordered set $\Lambda$ such that for any
$\lambda, \mu \in \Lambda$ there exists $\nu \in \Lambda$ with $\lambda \le \nu$ and $\mu \le \nu$. For example, a
totally ordered set is directed; the set of finite subsets of a set $S$, ordered by
inclusion, is a directed set which is not totally ordered.

Suppose that for each element $\lambda$ of a directed set $\Lambda$ we are given a set $M_\lambda$,
and whenever $\lambda \le \mu$ we are given a map $f_{\mu\lambda}\colon M_\lambda \to M_\mu$ satisfying the
conditions
\[
  f_{\lambda\lambda} = 1, \quad \text{and} \quad f_{\nu\mu} \circ f_{\mu\lambda} = f_{\nu\lambda}
      \quad \text{for } \lambda \le \mu \le \nu;
\]
we express all this data as $(M_\lambda; f_{\mu\lambda})$, and refer to it as a \emph{direct system}
over $\Lambda$ (or indexed by $\Lambda$). If each $M_\lambda$ is an $A$-module, and each $f_{\mu\lambda}$
$A$-linear we speak of a direct system of $A$-modules; if each $M_\lambda$ is a ring,
and each $f_{\mu\lambda}$ a ring homomorphism, a direct system of rings. More generally,
we can define direct systems in any category.

Given two direct systems $\mathcal{F} = \{M_\lambda; f_{\mu\lambda}\}$ and
$\mathcal{F}' = \{M'_\lambda; f'_{\mu\lambda}\}$ indexed by the same set, a morphism
$\varphi\colon \mathcal{F} \to \mathcal{F}'$ is a system of maps $\{\varphi_\lambda\colon M_\lambda \to M'_\lambda\}$ such that
\[
  f'_{\mu\lambda} \circ \varphi_\lambda = \varphi_\mu \circ f_{\mu\lambda} \quad \text{for } \lambda \le \mu.
\]
By a map from $\mathcal{F}$ to a set $X$ we mean a system $\{\varphi_\lambda\}$ of maps
$\varphi_\lambda\colon M_\lambda \to X$ satisfying $\varphi_\lambda = \varphi_\mu \circ f_{\mu\lambda}$ for $\lambda \le \mu$. Now if a
map $\psi\colon \mathcal{F} \to M_\infty$ from $\mathcal{F}$ to a set $M_\infty$ has the universal property for maps
from $\mathcal{F}$ to sets, that is, if for any map $\varphi\colon \mathcal{F} \to X$ there exists a unique
map $h\colon M_\infty \to X$ such that $\varphi_\lambda = h \circ \psi_\lambda$ for all $\lambda \in \Lambda$, then $M_\infty$ is called
the \emph{direct limit} of $\mathcal{F}$, or simply the limit of $\mathcal{F}$, and we write
$M_\infty = \varinjlim M_\lambda$, or $M_\infty = \lim M_\lambda$. As one sees easily from the definition, a map
$\varphi\colon \mathcal{F} \to \mathcal{F}'$ induces a map $\lim M_\lambda \to \lim M'_\lambda$, which in this book we write
$\varphi_\infty$ or $\lim \varphi$.

The limit of a direct system $\mathcal{F} = \{M_\lambda; f_{\mu\lambda}\}$ always exists. In order to
construct it we do the following: take the disjoint union $\bigsqcup_\lambda M_\lambda$ of the
$M_\lambda$, and define a relation $\sim$ by
\[
  x \sim y \quad (x \in M_\lambda, y \in M_\mu) \iff \text{there exists a } \nu
      \text{ with } \lambda \le \nu, \mu \le \nu \text{ and } f_{\nu\lambda}(x) = f_{\nu\mu}(y).
\]
Then it is easy to see that $\sim$ is an equivalence relation. We write $M_\infty$ for the
quotient set $(\bigsqcup_\lambda M_\lambda)/\sim$, that is the set of equivalence classes under $\sim$;
then one sees easily that $M_\infty$ satisfies the conditions for a direct limit. We write
$\lim x \in M_\infty$ for the equivalence class of $x \in M_\lambda$. If $\mathcal{F}$ is a direct system of modules
then $M_\infty$ can be given a natural structure of $A$-module, and
$x \mapsto \lim x$ is an $A$-linear map from $M_\lambda$ to $M_\infty$. Similarly for direct systems
or rings.

The above is general theory. In this book the following two theorems are of
particular importance.

\begin{theorem}\label{thm:A1}
Let $A$ be a ring, $N$ an $A$-module, and let $\mathcal{F} = \{M_\lambda; f_{\mu\lambda}\}$ be a direct
system of $A$-modules. Then
\[
  \varinjlim (M_\lambda \otimes_A N) = (\varinjlim M_\lambda) \otimes_A N.
\]
(In other words, tensor product commutes with direct limits.)
\end{theorem}

\begin{proof}
Set $\lim M_\lambda = M_\infty$ and $\lim(M_\lambda \otimes N) = L_\infty$. We write
$\varphi_\lambda\colon M_\lambda \to M_\infty$ for the $A$-linear map given by $x \mapsto \lim x$, so that
$\{\varphi_\lambda \otimes 1\}$ is a map from the direct system $(M_\lambda \otimes N; f_{\mu\lambda} \otimes 1)$ to the
$A$-module $M_\infty \otimes N$; this determines a unique $A$-linear map
$h\colon L_\infty \to M_\infty \otimes N$. For $x \in M_\lambda$ and $y \in N$ we have
$h(\lim(x \otimes y)) = (\lim x) \otimes y$. On the other hand, fixing $y \in N$ we can define
$g_{\lambda, y}\colon M_\lambda \to L_\infty$ by $g_{\lambda, y}(x) = \lim(x \otimes y)$, and in the limit we get
\[
  g_y\colon M_\infty \to L_\infty.
\]
If $x_\infty \in M_\infty$ we can write $x_\infty = \lim x$ for some $\lambda$ and some $x \in M_\lambda$. Then
$g_y(x_\infty) = g_{\lambda, y}(x) = \lim(x \otimes y)$. From this we can see that $g_y(x_\infty)$ is
bilinear in $x_\infty$ and in $y$, and so defines an $A$-linear map
$g\colon M_\infty \otimes N \to L_\infty$ such that $g(x_\infty \otimes y) = g_y(x_\infty)$. Now it is easy to see
that $g$ and $h$ are inverse maps, so that $M_\infty \otimes N \cong L_\infty$.
\end{proof}

\begin{theorem}\label{thm:A2}
Suppose that we have three direct systems of $A$-modules indexed by the same
set $\Lambda$, $\mathcal{F}' = \{M'_\lambda; f'_{\mu\lambda}\}$, $\mathcal{F} = \{M_\lambda; f_{\mu\lambda}\}$ and
$\mathcal{F}'' = \{M''_\lambda; f''_{\mu\lambda}\}$, and maps $\{\varphi_\lambda\}\colon \mathcal{F}' \to \mathcal{F}$ and
$\{\psi_\lambda\}\colon \mathcal{F} \to \mathcal{F}''$ such that for every $\lambda$,
\[
  M'_\lambda \xrightarrow{\varphi_\lambda} M_\lambda \xrightarrow{\psi_\lambda} M''_\lambda
\]
is an exact sequence; then the sequence obtained in the limit
\[
  \varinjlim M'_\lambda \xrightarrow{\varphi_\infty} \varinjlim M_\lambda \xrightarrow{\psi_\infty} \varinjlim M''_\lambda
\]
is also exact. (In other words, direct limit is an exact functor.)
\end{theorem}

\begin{proof}
Write $M_\infty$ for the limit of $\mathcal{F}$, and let $y_\infty \in M_\infty$ be such that
$\psi_\infty(y_\infty) = 0$. For some $\lambda$ and $y \in M_\lambda$ we can write $y_\infty = \lim y$, and then
$0 = \psi_\infty(\lim y) = \lim \psi_\lambda(y)$, so that for some $\mu \ge \lambda$ we have
$\psi_\mu(f_{\mu\lambda}(y)) = 0$; the left-hand side here is equal to
$\psi_\mu(f_{\mu\lambda}(y))$, so that by assumption there is $x \in M'_\mu$ such that
$f_{\mu\lambda}(y) = \varphi_\mu(x)$. Thus $y_\infty = \lim f_{\mu\lambda}(y) = \lim \varphi_\mu(x) =
\varphi_\infty(\lim x) \in \Im(\varphi_\infty)$. Also, $\psi_\infty \circ \varphi_\infty = 0$ is obvious.
\end{proof}

Given an $A$-module $M$, write $\{M_\lambda\}_{\lambda \in \Lambda}$ for the collection of all finitely
generated submodules of $M$. We define a partial order on $\Lambda$ by letting
$\lambda \le \mu$ if $M_\lambda \subset M_\mu$; this makes $\Lambda$ into a directed set, and we write
$f_{\mu\lambda}\colon M_\lambda \to M_\mu$ for the natural inclusion. Then $\{M_\lambda; f_{\mu\lambda}\}$ is a
direct system of $A$-modules, the limit of which is the original $M$, that is
$M = \varinjlim M_\lambda$. Hence any $A$-module can be expressed as a direct limit of
finitely generated $A$-modules.

In a similar way, given any ring $A$ and a subring $A_0 \subset A$, we can express
$A$ as the direct limit of subrings which are finitely generated over $A_0$ as
rings. If we take $A_0$ to be the minimal subring of $A$ (that is, the image in $A$
of $\mathbb{Z}$), then a ring which is finitely generated over $A_0$ is Noetherian, and
hence every ring is a direct limit of Noetherian rings.

\medskip
\noindent\textbf{Inverse limits}

Inverse systems and inverse limits are defined as the dual notions to direct
systems and direct limits, that is by reversing all the arrows in the
definitions. That is, we take a directed set $\Lambda$ as indexing set; an inverse
system of sets is the data of a set $M_\lambda$ for each $\lambda \in \Lambda$, and of a map
$f_{\lambda\mu}\colon M_\mu \to M_\lambda$ whenever $\lambda \le \mu$, such that
\[
  f_{\lambda\lambda} = 1, \quad \text{and} \quad f_{\lambda\mu} \circ f_{\mu\nu} = f_{\lambda\nu}
      \quad \text{for } \lambda \le \mu \le \nu;
\]
we write this as $(M_\lambda; f_{\lambda\mu})$. A morphism between two inverse systems with
the same indexing set, and a map from a set $N$ to an inverse system
$\mathcal{F} = \{M_\lambda; f_{\lambda\mu}\}$ are defined dually to the case of direct systems. We say
that $M_\infty$ is an \emph{inverse limit}, or projective limit of $\mathcal{F}$, and write
$M_\infty = \varprojlim M_\lambda$, if there is a map $\varphi = \{\varphi_\lambda\}\colon M_\infty \to \mathcal{F}$ which has the
property that for any set $X$, and any map $\psi = \{\psi_\lambda\}\colon X \to \mathcal{F}$, there exists a
unique map $h\colon X \to M_\infty$ such that $\psi_\lambda = \varphi_\lambda h$ for all $\lambda$. To prove the
existence of $\varprojlim M_\lambda$ we only have to let $M_\infty$ be the following subset of the
direct product $\prod_\lambda M_\lambda$:
\[
  M_\infty = \{(x_\lambda)_{\lambda \in \Lambda} \mid \lambda \le \mu \Rightarrow x_\lambda =
      f_{\lambda\mu}(x_\mu)\}.
\]
If each $M_\lambda$ is a module and each $\varphi_{\lambda\mu}$ is a linear map then this $M_\infty$ is a
submodule of the direct product module, and is the inverse limit of modules.
In a similar way, the inverse limit of an inverse system of rings is again a
ring.

\begin{example*}
Let $\Lambda = \{1, 2, 3, \ldots\}$ and let $p$ be a prime number. Consider the inverse
system of rings
\[
  \mathbb{Z}/(p) \leftarrow \mathbb{Z}/(p^2) \leftarrow \mathbb{Z}/(p^3) \leftarrow \cdots,
\]
where each arrow is the natural homomorphism. The inverse limit
$\varprojlim \mathbb{Z}/(p^n)$ is known as the ring of \emph{$p$-adic integers}. Its elements are of
the form
\[
  (a_1, a_2, a_3, \ldots), \quad \text{with } a_i \in \mathbb{Z}/(p^i) \text{ and } a_i = a_{i-1} \bmod p^{i-1};
\]
addition and multiplication is carried out term-by-term:
\[
  (a_1, a_2, \ldots) + (b_1, b_2, \ldots) = (a_1 + b_1, a_2 + b_2, \ldots)
\]
\[
  (a_1, a_2, \ldots) \cdot (b_1, b_2, \ldots) = (a_1b_1, a_2b_2, \ldots).
\]
More generally, if $A$ is any ring and $I$ an ideal of $A$, the inverse limit
$\varprojlim A/I^n$ of the inverse system of rings $A/I \leftarrow A/I^2 \leftarrow \cdots$ is called the
\emph{$I$-adic completion} of $A$ (see \S\ref{sec:8}).
\end{example*}

Taking the inverse limit of an inverse system of modules is a left-exact
functor, but is not an exact functor, so that the analog of
Theorem~\hyperref[thm:A2]{A2} for inverse systems does not hold.

\begin{example*}
Consider the diagram
\[
\begin{tikzcd}
0 \arrow[r] & \mathbb{Z} \arrow[r, "n"] \arrow[d] & \mathbb{Z} \arrow[r] \arrow[d] & \mathbb{Z}/(n) \arrow[r] \arrow[d] & 0 \\
0 \arrow[r] & \mathbb{Z} \arrow[r, "p"] \arrow[d] & \mathbb{Z} \arrow[r] \arrow[d] & \mathbb{Z}/(n) \arrow[r] \arrow[d] & 0 \\
0 \arrow[r] & \mathbb{Z} \arrow[r, "n"] & \mathbb{Z} \arrow[r] & \mathbb{Z}/(n) \arrow[r] & 0
\end{tikzcd}
\]
here $p$ and $n$ are coprime integers, and the arrows marked $\times$ are multiplication
by $n$. The rows are exact sequences, and each column defines an inverse
system. The left-hand and middle columns have $0$ as their inverse limits, but
since every arrow in the right-hand column is an isomorphism, the inverse
limit is isomorphic to $\mathbb{Z}/n$. Thus going to the inverse limit, we find that
$0 \to 0 \to \mathbb{Z}/(n)$ is exact, but the second arrow is not surjective.
\end{example*}

% ============ 习题 ============
\section*{Exercises to Appendix A}
\emph{Prove the following propositions.}

\noindent\textbf{A.1.}\quad Let $M$ and $N$ be $A$-modules. If the natural map
$M' \otimes_A N \to M \otimes_A N$ is injective for every finitely generated submodule
$M' \subset M$ then the same thing holds for every submodule $M' \subset M$.

\medskip
\noindent\textbf{A.2.}\quad Let $A$ be a ring, and $B$, $C$, $D$ (commutative) $A$-algebras. Then to
give a homomorphism of $A$-algebras from $B \otimes_A C$ to $D$ is the same thing as to
give a pair of homomorphisms of $A$-algebras $B \to D$ and $C \to D$; in other words,
$B \otimes_A C$ is the category-theoretical direct product of $B$ and $C$ in the category of
$A$-algebras.

% ============ Appendix B. Some homological algebra (原书 pp.274-282) ============
\chapter*{Appendix B}
\addcontentsline{toc}{chapter}{Appendix B: Some homological algebra}
\markboth{Appendix B}{Some homological algebra}
\renewcommand{\thetheorem}{B\arabic{theorem}}
\setcounter{theorem}{0}

\vspace*{-0.5cm}
\begin{center}
  {\Large\bfseries Some homological algebra}
\end{center}
\vspace{0.3cm}

Let $A$ be a ring; by a map from an $A$-module into another we mean an
$A$-linear map.

\medskip
\noindent\textbf{Complexes}

By a complex we mean a sequence
\[
  \cdots \to K_{n+1} \xrightarrow{d_{n+1}} K_n \xrightarrow{d_n} K_{n-1} \to \cdots
\]
of $A$-modules and $A$-linear maps such that $d_{n-1} \circ d_n = 0$ for every $n$.
This complex is written $K_\bullet$. Since $\Im(d_{n+1}) \subset \Ker(d_n)$ we can define
$H_n(K_\bullet) = \Ker(d_n)/\Im(d_{n+1})$ to be the \emph{homology} of $K_\bullet$ in dimension $n$. To
say that $H_n(K_\bullet) = 0$ for all $n$ is to say that $K_\bullet$ is exact. We also consider
complexes in which the indices go the other way, $\cdots \to K^n \to K^{n+1} \to \cdots$,
and for these we write $K^\bullet$ for the complex, and
$H^n(K^\bullet) = \Ker(d_n)/\Im(d_{n-1})$ for the \emph{cohomology} of $K^\bullet$ in dimension $n$.
From now on we often omit the indices, writing $d$ for $d_n$. We call $d$ the
\emph{differential} of the complex $K_\bullet$.

A morphism $f\colon K_\bullet \to K'_\bullet$ of complexes is a family $f = (f_n)_{n \in \mathbb{Z}}$ of
$A$-linear maps $f_n\colon K_n \to K'_n$ satisfying $d' \circ f_n = f_{n-1} \circ d$, or in other
words a commutative diagram
\[
\begin{tikzcd}
\cdots \arrow[r] & K_n \arrow[r] \arrow[d, "f_n"] & K_{n-1} \arrow[r] \arrow[d, "f_{n-1}"] & K_{n-2} \arrow[r] \arrow[d, "f_{n-2}"] & \cdots \\
\cdots \arrow[r] & K'_n \arrow[r] & K'_{n-1} \arrow[r] & K'_{n-2} \arrow[r] & \cdots
\end{tikzcd}
\]
In an obvious way, $f$ induces a linear map $H_n(K_\bullet) \to H_n(K'_\bullet)$ between the
homology modules in each dimension; this is often written $f_{n*}$, or simply $f$
if there is no fear of confusion. If $f, g\colon K_\bullet \to K'_\bullet$ are two morphisms we say
that $f$ and $g$ are \emph{homotopic} (denoted $f \sim g$) if for each $n$ there is a linear map
$h_n\colon K_n \to K'_{n+1}$ such that
\[
  f_n - g_n = d'h_n + h_{n-1}d.
\]
If this happens then $f$ and $g$ induce the same map $H_n(K_\bullet) \to H_n(K'_\bullet)$ on
homology. Two complexes $K_\bullet$ and $K'_\bullet$ are said to be \emph{homotopy equivalent} if
there exist morphisms $f\colon K_\bullet \to K'_\bullet$ and $g\colon K'_\bullet \to K_\bullet$ such that $gf \sim 1$,
and $fg \sim 1_{K'}$, where $1$ denotes the identity map $K_\bullet \to K_\bullet$. Homotopy
equivalent complexes have the same homology.

A sequence of complexes
\[
  0 \to K'_\bullet \xrightarrow{f} K_\bullet \xrightarrow{g} K''_\bullet \to 0
\]
is said to be \emph{exact} if
\[
  0 \to K'_n \xrightarrow{f_n} K_n \xrightarrow{g_n} K''_n \to 0
\]
is exact for every $n$. In this case, a connecting homomorphism
$\delta_n\colon H_n(K''_\bullet) \to H_{n-1}(K'_\bullet)$ is defined as follows: for $\zeta \in H_n(K''_\bullet)$, choose
$x \in \Ker d''_n$ representing $\zeta$, and take $y \in K_n$ such that $g(y) = x$; then since
$g(dy) = 0$ there is a well-determined $z \in K'_{n-1}$ for which $f_{n-1}(z) = dy$, and
$dz = 0$. The class $\xi \in H_{n-1}(K')$ represented by $z$ can easily be seen to depend
only on $\zeta$, and $\delta_n$ is defined by $\delta_n(\zeta) = \xi$. The following sequence is then
exact:
\[
  \cdots \to H_n(K'_\bullet) \xrightarrow{f} H_n(K_\bullet) \xrightarrow{g} H_n(K''_\bullet)
      \xrightarrow{\delta_n} H_{n-1}(K'_\bullet) \to \cdots
\]
The proof does not require anything new, and is well-known, so that we omit
it; this should be thought of as a fundamental theorem of homology theory.
The above sequence is called the \emph{homology long exact sequence} of the short exact
sequence $0 \to K'_\bullet \to K_\bullet \to K''_\bullet \to 0$.

\medskip
\noindent\textbf{Double complexes}

A double complex of $A$-modules is a doubly indexed family
$K_{\bullet\bullet} = (K_{pq})_{p, q \in \mathbb{Z}}$ of $A$-modules, with two sets of $A$-linear maps
$d'\colon K_{pq} \to K_{p-1,q}$ and $d''\colon K_{pq} \to K_{p,q-1}$ for which $d'd' = 0$, $d''d'' = 0$ and
$d'd'' = d''d'$. Given a double complex $K_{\bullet\bullet}$, if we set
\[
  K_n = \bigoplus_{p + q = n} K_{pq}
\]
and define $d_n\colon K_n \to K_{n-1}$ by
\[
  dx = d'x + (-1)^p d''x \quad \text{if } x \in K_{pq},
\]
then since $dd = 0$, the $\{K_n\}$ form an ordinary complex with differential $d$. The
homology of this complex is called the homology of $K_{\bullet\bullet}$, and written
$H_\bullet(K_{\bullet\bullet})$, or simply $H(K)$.

To treat homology and cohomology in a unified manner, we fix the following
convention on raising and lowering indices: $K_{pq} = K^{-p,-q} = K^{p,q} = K^{-p,-q}$.
For example, given a double complex $(K^{pq})$ with
$d'\colon K^{pq} \to K^{p+1,q}$ and $d''\colon K^{pq} \to K^{p,q+1}$, we think of $K^{pq}$ as
$K_{-p,-q}$, and set
\[
  K^n = \bigoplus_{p + q = n} K^{pq}.
\]

The basic technique for studying the homology of double complexes is spectral
sequences, but we leave this to specialist texts, and only consider here the
extreme cases which we will use later.

We can fix the first index $p$ in $K_{\bullet\bullet}$ getting a complex $K_{p\bullet}$:
\[
  \cdots \to K_{p,q+1} \xrightarrow{d''} K_{pq} \xrightarrow{d''} K_{p,q-1} \to \cdots
\]
Similarly, for fixed $q$, $d'$ defines a complex $K_{\bullet q}$.

Now suppose that $K_{\bullet\bullet}$ satisfies the condition $K_{pq} = 0$ if $p$ or $q < 0$ (a
first quadrant double complex). We set $H_0(K_{p\bullet}) = K_{p0}/d''K_{p1} = X_p$; then
$d'$ induces a map $X_p \to X_{p-1}$, making the $X_p$ into a complex $X_\bullet$. Similarly,
$d''$ makes the $H_q(K_{\bullet q}) = Y_q$ into a complex $Y_\bullet$. In this notation we have
the following theorem.

\begin{theorem}\label{thm:B1}
Suppose that the double complex $K_{\bullet\bullet}$ satisfies the conditions
\[
  K_{pq} = 0 \quad \text{for } p \text{ or } q < 0,
\]
and
\[
  H_q(K_{p\bullet}) = 0 \quad \text{for } q > 0 \text{ and all } p.
\]
Then in the above notation we have
\[
  H_n(K) \cong H_n(X_\bullet) \quad \text{for all } n.
\]
If in addition we have $H_p(K_{\bullet q}) = 0$ for $p > 0$ and all $q$ then
\[
  H_n(X_\bullet) = H_n(K) = H_n(Y_\bullet).
\]
\end{theorem}

\begin{proof}[Sketch proof]
Write $a_{ij}$ to denote an element of $K_{ij}$. We define a map
$\Phi\colon K_n \to X_n$ by taking $a = \sum_{i=0}^{n} a_{n-i,i} \in K_n$ into $\varphi(a_{n0}) \in X_n$, where
$\varphi\colon K_{n0} \to X_n$ denotes the canonical map. Then $\Phi$ is a morphism of
complexes, and we prove that it induces an isomorphism on homology.

Let $x \in X_n$ with $d'x = 0$. We can take $a_{n0}$ such that $x = \varphi(a_{n0})$, and
then since $\varphi(d'a_{n0}) = d'x = 0$ there exists $a_{n-1,1}$ such that
$d'a_{n0} = d''a_{n-1,1}$. In turn, since
$d''(d'a_{n-1,1}) = d'(d''a_{n-1,1}) = d'd'a_{n0} = 0$ and since $H_1(K_{n-2,\bullet}) = 0$ there
exists $a_{n-2,2}$ such that $da_{n-1,1} = d''a_{n-2,2}$; then proceeding as before, we can
choose $a_{n-i,i}$ for $0 \le i \le n$ such that $d'a_{n-i,i} = d''a_{n-i-1,i+1}$ for
$0 \le i < n$. Then for a suitable choice of signs,
$a = \sum_i a_{n-i,i} \in K_n$ satisfies $da = 0$ and $\Phi(a) = x$, and this proves that
$\Phi$ induces a surjection $H_n(K) \to H_n(X_\bullet)$. The proof that this is also
injective is similar. The second part follows by symmetry from the first.
\end{proof}

Dually, we have the following theorem for cohomology.

\begin{theorem}\label{thm:B2}
Suppose that the double complex $K^{\bullet\bullet}$ satisfies
\[
  K^{pq} = 0 \quad \text{for } p \text{ or } q < 0
\]
and
\[
  H^q(K^{p\bullet}) = 0 \quad \text{for } q > 0 \text{ and all } p.
\]
Then making $X^p = \Ker(d''\colon K^{p0} \to K^{p1})$ into a complex $X^\bullet$ by means of
$d'$, we have
\[
  H^n(K) \cong H^n(X^\bullet),
\]
If in addition $H^p(K^{\bullet q}) = 0$ for $p > 0$ and all $q$ then
\[
  H^n(Y^\bullet) \cong H^n(K) \cong H^n(X^\bullet),
\]
where $Y^\bullet$ is the complex made from $Y^q = \Ker(d'\colon K^{0q} \to K^{1q})$.
\end{theorem}

\begin{proof}
We leave the proof as a suitable exercise for the reader.
\end{proof}

\medskip
\noindent\textbf{Projective and injective modules}

An $A$-module $P$ is said to be \emph{projective} if it satisfies the following condition:
for any surjection $f\colon M \to N$ and any map $g\colon P \to N$, there exists a lifting
$h\colon P \to M$ such that $g = fh$. A free module is projective, and we can
characterise projective modules as direct summands of free modules. Indeed, if
we express a projective module $P$ as a quotient $P = F/G$ of a free module $F$
then the identity map $P \to P$ has a lifting such that $P \to F \to P$ is the identity
map, and then $F \cong P \oplus G$. Reversing the arrows and replacing surjection by
injection in the definition of projective module, we get the definition of
\emph{injective} module. There is no dual notion to that of a free module, so that
injective modules do not have any very simple characterisation, but we can
easily prove the following theorem using Zorn's lemma.

\begin{theorem}\label{thm:B3}
A necessary and sufficient condition for an $A$-module $I$ to be injective is that
for any ideal $\mathfrak{a}$ of $A$, and any map $\varphi\colon \mathfrak{a} \to I$, it is possible to extend
$\varphi$ to a map from the whole of $A$ to $I$.
\end{theorem}

Any $A$-module can be written as a quotient of a projective module (take for
example a free module). Dually to this, any module can be embedded into an
injective module; the proof of this is a little tricky, and we leave it to more
specialist textbooks. Given a module $M$, consider a surjection
$P_0 \to M$ from a projective module $P_0$ to $M$; letting $K_1$ be the kernel, we get an
exact sequence $0 \to K_1 \to P_0 \to M \to 0$. In the same way, we construct for $K_1$
an exact sequence $0 \to K_2 \to P_1 \to K_1 \to 0$ with $P_1$ projective, and proceeding as
before we get exact sequences $0 \to K_i \to P_{i-1} \to K_{i-1} \to 0$ for $i = 1, 2, \ldots$ with
$P_i$ projective. The resulting complex
\[
  P_\bullet\colon \cdots \to P_n \to P_{n-1} \to \cdots \to P_1 \to P_0 \to 0
\]
is called a \emph{projective resolution} of $M$. Since by construction this becomes an exact
sequence on replacing the final $P_0 \to 0$ by $P_0 \to M \to 0$, we have $H_n(P_\bullet) = 0$
for $n > 0$ and $H_0(P_\bullet) = M$. In the case that $A$ is Noetherian and $M$ is finite,
we can take $P_0$ to be a free module of finite rank, and then $K_1$ is again
finitely generated. Proceeding in the same way, we see that $M$ has a projective
resolution in which each $P_n$ is a finite free module.

Dually, for any $A$-module $M$ there exists an exact sequence of the form
$0 \to M \to Q^0 \to Q^1 \to \cdots$ with each $Q^n$ an injective module. The complex
$Q^\bullet\colon 0 \to Q^0 \to Q^1 \to \cdots$ is called an \emph{injective resolution} of $M$; it
satisfies $H^0(Q^\bullet) = M$ and $H^n(Q^\bullet) = 0$ for $n > 0$.

If $f\colon M \to N$ is a map of $A$-modules and $P_\bullet$, $P'_\bullet$ are projective
resolutions of $M$ and $N$ then there exists a morphism of complexes
$\varphi\colon P_\bullet \to P'_\bullet$ for which $f\varepsilon = \varepsilon'\varphi_0$, that is, a commutative diagram
\[
\begin{tikzcd}
\cdots \arrow[r] & P_1 \arrow[r] \arrow[d, "\varphi_1"] & P_0 \arrow[r] \arrow[d, "\varphi_0"] & M \arrow[r] \arrow[d, "f"] & 0 \\
\cdots \arrow[r] & P'_1 \arrow[r] & P'_0 \arrow[r] & N \arrow[r] & 0
\end{tikzcd}
\]
The existence of $\varphi_0, \varphi_1, \ldots$ can easily be proved successively, using the fact
that the $P_n$ are projective and the exactness of the lower sequence. Up to
homotopy, this $\varphi$ is unique, that is if $\varphi$ and $\psi$ both have the given property
then $\varphi \sim \psi$. We leave the proof of this to the reader. From this it follows
that any two projective resolutions of $M$ are homotopy equivalent. Exactly the
same thing holds for injective resolutions.

\medskip
\noindent\textbf{The Tor functors}

Let $M$ and $N$ be $A$-modules and $P_\bullet$, $Q_\bullet$ projective resolutions of $M$ and
$N$, respectively. We write $P_\bullet \otimes N$ for the complex obtained by tensoring
$P_\bullet$ through with $N$:
\[
  P_\bullet \otimes N\colon \cdots \to P_1 \otimes N \to P_0 \otimes N \to 0.
\]
The complex $M \otimes Q_\bullet$ is constructed similarly. Moreover, we can define a
double complex $K_{\bullet\bullet}$ by $K_{pq} = P_p \otimes_A Q_q$, with the obvious definitions of
$d'$, $d''$. Each $P_p$ is a direct summand of a free module, and is therefore flat
(that is performing $\otimes P_p$ takes exact sequences into exact sequences). Thus
$H_q(K_{p\bullet}) = H_q(P_p \otimes Q_\bullet) = 0$ for $q > 0$, and
$H_0(K_{p\bullet}) = H_0(P_p \otimes Q_\bullet) = P_p \otimes N$. In exactly the same way,
$H_p(K_{\bullet q}) = 0$ for $p > 0$ and $H_0(K_{\bullet q}) = M \otimes Q_q$, and therefore by
Theorem~\hyperref[thm:B1]{B1} $H_n(P_\bullet \otimes N) \cong H_n(K_{\bullet\bullet}) \cong H_n(M \otimes Q_\bullet)$. This
module (defined uniquely up to isomorphism) is written $\Tor^A_n(M, N)$; it is
independent of the choice of the projective resolutions of $M$ and $N$ chosen,
since if $P_\bullet$ and $P'_\bullet$ are two projective resolutions, we have $P_\bullet \sim P'_\bullet$, and
therefore $P_\bullet \otimes N \sim P'_\bullet \otimes N$.

The Tor functors have the following properties, (all of which can be proved
directly from the definition):
\begin{enumerate}[(1)]
  \item $\Tor^A_0(M, N) = M \otimes_A N$;
  \item if $M$ is flat then $\Tor^A_n(M, N) = 0$ for any $N$ and $n > 0$;
  \item $\Tor^A_n(M, N) \cong \Tor^A_n(N, M)$;
  \item $\Tor^A_n(M, N)$ is a covariant functor in both of its entries, and each short
        exact sequence $0 \to M' \to M \to M'' \to 0$ leads to a long exact sequence
        \[
          \cdots \to \Tor^A_n(M', N) \to \Tor^A_n(M, N) \to \Tor^A_n(M'', N)
          \to \Tor^A_{n-1}(M', N) \to \cdots \to \Tor^A_0(M'', N)
          \to M' \otimes N \to M \otimes N \to M'' \otimes N \to 0.
        \]
  \item If $(N_\lambda, f_{\mu\lambda})$ is a direct system of $A$-modules then
        \[
          \Tor^A_n(M, \varinjlim N_\lambda) = \varinjlim \Tor^A_n(M, N_\lambda).
        \]
\end{enumerate}

\medskip
\noindent\textbf{The Ext functors}

Let $M$ and $N$ be $A$-modules. The functor $\Hom_A(M, -)$ is left-exact, that is it
takes an exact sequence $0 \to N' \to N \to N'' \to 0$ into an exact sequence
$0 \to \Hom_A(M, N') \to \Hom_A(M, N) \to \Hom_A(M, N'')$; and $M$ is projective if and
only if $\Hom_A(M, -)$ is exact. In addition, $\Hom_A(-, N)$ is left-exact, in the sense
that it takes an exact sequence $0 \to M' \to M \to M'' \to 0$ into an exact sequence
$0 \to \Hom_A(M'', N) \to \Hom_A(M, N) \to \Hom_A(M', N)$, and $N$ is injective if and
only if $\Hom_A(-, N)$ is exact.

Choose a projective resolution $P_\bullet$ of $M$ and an injective resolution $Q_\bullet$ of $N$;
we define a double complex $K^{\bullet\bullet}$ by $K^{pq} = \Hom_A(P_p, Q^q)$, and construct
the two complexes
\[
  \Hom_A(M, Q^\bullet)\colon 0 \to \Hom_A(M, Q^0) \to \Hom_A(M, Q^1) \to \cdots
\]
and
\[
  \Hom_A(P_\bullet, N)\colon 0 \to \Hom_A(P_0, N) \to \Hom_A(P_1, N) \to \cdots.
\]
Then by Theorem~\hyperref[thm:B2]{B2} we get
\[
  H^n(\Hom_A(M, Q^\bullet)) \cong H^n(K^{\bullet\bullet}) \cong H^n(\Hom_A(P_\bullet, N)).
\]
Identifying these three, we write $\Ext^n_A(M, N)$. As with Tor, this does not depend
on the choice of $P_\bullet$ and $Q^\bullet$.

The main properties of the Ext functors are as follows:
\begin{enumerate}[(1)]
  \item $\Ext^0_A(M, N) = \Hom_A(M, N)$,
  \item If $M$ is projective, or if $N$ is injective, then $\Ext^n_A(M, N) = 0$ for $n > 0$;
  \item $\Ext^n_A(M, N)$ is a contravariant functor in $M$ and a covariant functor
        in $N$. A short exact sequence $0 \to M' \to M \to M'' \to 0$ gives rise to
        a long exact sequence
        \[
          0 \to \Hom_A(M'', N) \to \Hom_A(M, N) \to \Hom_A(M', N)
          \to \Ext^1_A(M'', N) \to \Ext^1_A(M, N) \to \Ext^1_A(M', N)
          \to \Ext^2_A(M'', N) \to \cdots,
        \]
        and a short exact sequence $0 \to N' \to N \to N'' \to 0$ gives rise to a long
        exact sequence
        \[
          0 \to \Hom_A(M, N') \to \Hom_A(M, N) \to \Hom_A(M, N'')
          \to \Ext^1_A(M, N') \to \Ext^1_A(M, N) \to \Ext^1_A(M, N'')
          \to \Ext^2_A(M, N') \to \cdots;
        \]
  \item $M$ is projective $\iff \Ext^1_A(M, N) = 0$ for all $N$,
        and
        \[
          N \text{ is injective } \iff \Ext^1_A(M, N) = 0 \text{ for all } M.
        \]
\end{enumerate}

\medskip
\noindent\textbf{Projective and injective dimensions}

If $M$ is an $A$-module for which there exists a projective resolution $P_\bullet$ with
$P_n = 0$ for $n > d$, but such that $P_d \ne 0$ for any choice of projective resolution
$P_\bullet$, then we say that $M$ has \emph{projective dimension} $d$, and write
$\projdim M = d$. If there is no such $d$ then we write $\projdim M = \infty$. The
\emph{injective dimension} $\injdim M$ is defined in the same way using injective
resolutions. Clearly $\projdim M = 0$ if and only if $M$ is projective, and
$\injdim M = 0$ if and only if $M$ is injective.

For a projective resolution $P_\bullet$ of $M$ and some $i > 0$, let $K_i$ denote the image of
$P_i \to P_{i-1}$; then $\cdots \to P_n \to P_{n-1} \to \cdots \to P_i \to 0$ is a projective
resolution of $K_i$, so that for $n > i$ we have $\Ext^n_A(M, N) = \Ext^{n-i}_A(K_i, N)$.
Now if $\Ext^{d+1}_A(M, N) = 0$ for all $N$, we have $\Ext^1_A(K_d, N) = 0$ for all $N$, and
hence $K_d$ is projective, so that $0 \to K_d \to P_{d-1} \to \cdots \to P_0 \to 0$ is also
a projective resolution of $P_0$, and $\projdim M \le d$. Conversely, if
$\projdim M \le d$ then obviously $\Ext^n_A(M, N) = 0$ for $n > d$.

Similarly, $\injdim N \le d \iff \Ext^{d+1}_A(M, N) = 0$ for all $M$.

\medskip
\noindent\textbf{Derived functors}

As we have just seen, the definition of functors like Tor and Ext can be given
using a resolution of just one entry. For instance, write $T$ for the functor
$\Hom_A(-, N)$, and let $P_\bullet$ be a projective resolution of a given module $M$;
construct the complex $T(P_\bullet)\colon \cdots \to T(P_n) \to T(P_{n-1}) \to \cdots \to T(P_0) \to 0$,
and take the cohomology $H^n(T(P_\bullet))$. Setting $R^nT(M) = H^n(T(P_\bullet))$ defines
a functor $R^nT$ in $M$, which we call the \emph{right derived functor} of the left exact
functor $T$. In the present case we have $R^nT = \Ext^n_A(-, N)$, but we can in
general define the right derived functor of a left exact contravariant functor
in the same way.

The right derived functor is uniquely determined by the following three
properties: (1) $R^0T = T$, (2) if $M$ is projective then $R^nT(M) = 0$ for all $n > 0$,
and (3) a short exact sequence $0 \to M' \to M \to M'' \to 0$ gives rise to a `natural'
long exact sequence
\[
  0 \to T(M'') \to T(M) \to T(M')
  \to R^1T(M'') \to R^1T(M) \to R^1T(M')
  \to R^2T(M'') \to \cdots
\]
(For the meaning of `natural' see a textbook on homological algebra.) For a
left exact covariant functor, we have to replace `projective' by `injective' in
the above. For right exact functors we can define a \emph{left derived functor} by
taking a projective resolution in the covariant case and an injective resolution
in the contravariant case.

\medskip
\noindent\textbf{Injective hull}

Let $L$ be an $A$-module and $M \subset L$ a submodule; we say that $L$ is an
\emph{essential extension} of $M$ if $N \cap M \ne 0$ for every non-zero submodule $N \subset L$,
or equivalently if
\[
  0 \ne x \in L \Rightarrow \text{there exists } a \in A \text{ such that } 0 \ne ax \in M.
\]

\begin{theorem}\label{thm:B4}
An $A$-module $M$ is injective if and only if it has no essential extensions
except $M$ itself.
\end{theorem}

We leave the proof to the reader. Now suppose that $M$ is a given $A$-module,
and choose an injective module $I$ with $M \subset I$. If we let $E$ be a maximal
element among all essential extensions of $M$ in $I$ then by the above theorem
$E$ is injective. An injective module $E$ such that $M \subset E$ is an essential
extension is called an \emph{injective hull} of $M$, and written $E(M)$ or $E_A(M)$; this
notion plays an important role in \S\ref{sec:18}. If $E$ and $E'$ are injective hulls of
$M$ then it is easy to see that there is an isomorphism $\varphi\colon E \to E'$ which
fixes the elements of $M$, although $\varphi$ itself is not necessarily unique.

Let $M$ be an $A$-module. Take an injective hull $I^0$ of $M$, and set
$K^1 = I^0/M$. Take an injective hull $I^1$ of $K^1$, and set $K^2 = I^1/K^1$.
Proceeding in the same way we obtain an injective resolution
$0 \to I^0 \to I^1 \to I^2 \to \cdots$ of $M$, which is called a \emph{minimal injective
resolution} of $M$.

The following two propositions are both famous and useful; the proofs are easy.

\medskip
\noindent\textbf{The five lemma.} Let
\[
\begin{tikzcd}
A \arrow[r] \arrow[d, "f_1"] & B \arrow[r] \arrow[d, "f_2"] & C \arrow[r] \arrow[d, "f_3"] & D \arrow[r] \arrow[d, "f_4"] & E \arrow[d, "f_5"] \\
A' \arrow[r] & B' \arrow[r] & C' \arrow[r] & D' \arrow[r] & E'
\end{tikzcd}
\]
be a commutative diagram of modules with exact rows. Then
\begin{enumerate}[(1)]
  \item $f_2$ surjective, and $f_1$ and $f_4$ injective $\Rightarrow f_3$ is injective;
  \item $f_4$ injective, and $f_2$ and $f_5$ surjective $\Rightarrow f_3$ is surjective.
\end{enumerate}

\medskip
\noindent\textbf{The snake lemma.} Let
\[
\begin{tikzcd}
A \arrow[r] \arrow[d, "\alpha"] & B \arrow[r] \arrow[d, "\beta"] & C \arrow[r] \arrow[d, "\gamma"] & 0 \\
0 \arrow[r] & A' \arrow[r] & B' \arrow[r] & C'
\end{tikzcd}
\]
be a commutative diagram of modules with exact rows. Then there exists an
exact sequence of the form
\[
  \Ker(\alpha) \to \Ker(\beta) \to \Ker(\gamma) \xrightarrow{\delta}
  \Coker(\alpha) \to \Coker(\beta) \to \Coker(\gamma).
\]

\medskip
\noindent\textbf{Tensor product of complexes}

Given two complexes of $A$-modules $K_\bullet$ and $L_\bullet$ the tensor product
$K_\bullet \otimes_A L_\bullet$ is defined as follows: firstly, set
\[
  (K \otimes L)_n = \bigoplus_{p + q = n} K_p \otimes_A L_q,
\]
and define the differential $d$ by setting
\[
  d(x \otimes y) = dx \otimes y + (-1)^p x \otimes dy
\]
for $x \in K_p$ and $y \in L_q$. In other words, $K \otimes L$ is the (single) complex obtained
from the double complex $W_{\bullet\bullet}$, where $W_{pq} = K_p \otimes L_q$.

There is an isomorphism of complexes $K \otimes L \cong L \otimes K$ obtained by sending
$x \otimes y$ into $(-1)^{pq}y \otimes x$ for $x \otimes y \in K_p \otimes L_q$. For a third complex of
$A$-modules $M$, the associative law holds:
\[
  (K \otimes L) \otimes M = K \otimes (L \otimes M).
\]
Hence the tensor product $K^{(1)} \otimes \cdots \otimes K^{(r)}$ of a finite number of complexes
can be defined by induction. This is used in \S\ref{sec:16}.

The information on homological algebra given above should be adequate for
the purpose of reading this book. However, a student intending to become a
specialist in algebra or geometry will require rather more detailed knowledge,
including the theory of spectral sequences. We mention here just three
representative references, two books by the originators of homological algebra
and category theory:
\begin{quote}
  H.~Cartan and S.~Eilenberg, \emph{Homological Algebra}, Princeton, 1956,\\
  S.~MacLane, \emph{Homology}, Springer, 1963,
\end{quote}
together with A.~Grothendieck's paper
\begin{quote}
  \emph{Sur quelques points d'alg\`ebre homologique}, T\^ohoku Math.~J.~9 (1957), 119--221.
\end{quote}

% ============ Appendix C. The exterior algebra (原书 pp.283-286) ============
\chapter*{Appendix C}
\addcontentsline{toc}{chapter}{Appendix C: The exterior algebra}
\markboth{Appendix C}{The exterior algebra}
\renewcommand{\thetheorem}{C\arabic{theorem}}
\setcounter{theorem}{0}

\vspace*{-0.5cm}
\begin{center}
  {\Large\bfseries The exterior algebra}
\end{center}
\vspace{0.3cm}

\begin{enumerate}[(1)]
  \item Let $M$ and $N$ be modules over a ring $A$. An $r$-multilinear map
        $\varphi\colon M^r = M \times \cdots \times M \to N$ from the direct product of $r$ copies of
        $M$ is said to be \emph{alternating} if $\varphi(x_1, \ldots, x_r) = 0$ whenever any of the
        elements $x_1, \ldots, x_r$ appears more than once. If $\varphi$ is alternating then for
        any $x_1, \ldots, x_r \in M$ we have
        \[
          \varphi(x_1, \ldots, x_{i-1}, x_i + x_j, x_{i+1}, \ldots, x_{j-1}, x_i + x_j,
              x_{j+1}, \ldots) = 0,
        \]
        and expanding out the left-hand side gives
        \[
          \varphi(x_1, \ldots, x_i, \ldots, x_j, \ldots) + \varphi(x_1, \ldots, x_j, \ldots,
              x_i, \ldots) = 0.
        \]
        In other words, on interchanging two of its entries, $\varphi$ changes sign.

        The \emph{$r$th exterior product} of $M$ is defined as the module $N_r$ having a
        universal alternating $r$-fold multilinear map $f_0\colon M^r \to N_r$, that is a map
        satisfying the property that every alternating $r$-fold multilinear map
        $f\colon M^r \to N$ factorises as $f = h \circ f_0$ for a unique $A$-linear map
        $h\colon N_r \to N$. We write $N_r = \bigwedge^r M$, and use
        $x_1 \wedge \cdots \wedge x_r$ to denote $f_0(x_1, \ldots, x_r)$. To prove the existence of the
        exterior product, let $N_r$ be the quotient of the $r$-fold tensor product
        $M \otimes \cdots \otimes M$ by the submodule generated by elements of the form
        $x_1 \otimes \cdots \otimes x \otimes \cdots \otimes x \otimes \cdots \otimes x_r$. Then $N_r$ satisfies the above
        condition. The fact that the exterior product is uniquely determined up to
        isomorphism is obvious from the definition.

  \item If $M$ is a free $A$-module of rank $n$, with basis $e_1, \ldots, e_n$, then
        $\bigwedge^r M$ is zero if $r > n$, and if $r \le n$ is the free $A$-module of rank
        $\binom{n}{r}$ with basis $(e_{i_1} \wedge \cdots \wedge e_{i_r} \mid 1 \le i_1 < \cdots < i_r \le n)$.
        (If $r > n$ this is easy; if $r \le n$ then the elements given above obviously
        generate $\bigwedge^r M$, and the fact that they are linearly independent can
        also be proved by reducing to the theory of determinants.)

  \item However, if $I \subset A$ is an ideal, then $\bigwedge^r A = 0$, but nevertheless
        $\bigwedge^r I$ is not necessarily $0$. For example, let $k$ be a field, $x$ and $y$
        indeterminates, and $A = k[x, y]$; if $I = xA + yA$ then $\bigwedge^2 I \ne 0$. Indeed,
        we can define $\varphi\colon I \times I \to k = A/I$ by
        \[
          \varphi(ax + by, a'x + b'y) = ab' - a'b,
        \]
        and it is then easy to check that $\varphi$ is alternating and bilinear, with
        $\varphi(x, y) = 1$, so that $\varphi \ne 0$, and we must have $\bigwedge^2 I \ne 0$.

  \item The operation of taking the exterior product commutes with extensions
        of scalars. That is, let $B$ be an $A$-algebra and $M$ an $A$-module, and set
        $M \otimes_A B = M_B$. Then $(\bigwedge^r M) \otimes_A B = \bigwedge^r M_B$, where of course
        $\bigwedge^r$ on the right-hand side refers to the exterior product of $B$-modules.
        For the proof, according to Appendix A, Formula~11, we have
        \[
          M_B \otimes_B \cdots \otimes_B M_B = (M \otimes_A \cdots \otimes_A M) \otimes_A B,
        \]
        so that if we let $f_0$ be the composite
        \[
          (M_B)^r = \underbrace{M \otimes_A \cdots \otimes_A M}_{r} \otimes_A B \to
              \Big(\bigwedge^r M\Big) \otimes_A B,
        \]
        then $f_0$ is an alternating $r$-fold $B$-multilinear map. Write
        $v\colon M \to M_B$ for the natural map $x \mapsto x \otimes 1$. Let $N$ be a $B$-module and
        $\varphi\colon (M_B)^r \to N$ be an alternating $r$-fold $B$-multilinear map. Then
        $\varphi$ induces an alternating $r$-fold $A$-multilinear map $\varphi\colon M^r \to N$,
        and therefore an $A$-linear map $\bigwedge^r M \to N$, and finally a $B$-linear map
        $(\bigwedge^r M) \otimes_A B \to N$ which we denote $h$. Then on $v(M^r)$ the two maps
        $\varphi$ and $h \circ f_0$ coincide; but $M_B$ is generated over $B$ by $v(M)$, so that
        $\varphi = h \circ f_0$. Thus $f_0$ has the universal property, and we can think of
        $(\bigwedge^r M) \otimes_A B$ as $\bigwedge^r M_B$.
\end{enumerate}

\begin{theorem}\label{thm:C1}
Let $A$ be an integral domain with field of fractions $K$, and let
$I_1, \ldots, I_r$ be ideals of $A$. Set $M = I_1 \otimes \cdots \otimes I_r$, and let $T$ be the torsion
submodule of $\bigwedge^r M$; then $(\bigwedge^r M)/T \cong I_1 \cdots I_r$. Therefore if
$J_1, \ldots, J_r$ are ideals of $A$ such that $I_1 \otimes \cdots \otimes I_r \cong J_1 \otimes \cdots \otimes J_r$
we have $I_1 \cdots I_r = J_1 \cdots J_r$.
\end{theorem}

\begin{proof}
We have $M_K \cong K \otimes \cdots \otimes K$ (the direct sum of $r$ copies of $K$), so that
$(\bigwedge^r M) \otimes K = \bigwedge^r M_K \cong K$. The kernel of the natural map
$\bigwedge^r M \to (\bigwedge^r M) \otimes K \cong K$ is obviously $T$ (since tensoring with $K$ is
the same thing as the localisation with respect to the zero ideal of $A$, see
\S\ref{sec:4}). In addition, the image is $I_1 \cdots I_r$. Indeed, viewing each $I_i$ as a
submodule of $K$, we can assume that the map is
\[
  \bigwedge^r (I_1 \otimes \cdots \otimes I_r) \to K,
\]
and since $\xi_i = \sum_{j=1}^{r} a_{ij}e_j \in K^r$ for $\xi_i \in I_i \otimes K$ we have
$\xi_1 \wedge \cdots \wedge \xi_r = \det(a_{ij})e_1 \wedge \cdots \wedge e_r$; it is clear from this that the above
map has image $I_1 \cdots I_r e_1 \wedge \cdots \wedge e_r$.
\end{proof}

If $A$ is a Dedekind ring then it is known that $I_1 \otimes \cdots \otimes I_r \cong A^{r-1} \otimes
I_1 \cdots I_r$ (see for example [B7], \S4, Prop.~24).

\begin{theorem}\label{thm:C2}
Let $A$ be a ring and $M$, $N$ $A$-modules. Then
\[
  \bigwedge^r (M \oplus N) = \bigoplus_{s + t = r} \Big[\Big(\bigwedge^s M\Big) \otimes_A
      \Big(\bigwedge^t N\Big)\Big].
\]
\end{theorem}

\begin{proof}
$\bigotimes^r(M \oplus N)$ can be written as a direct sum of all possible $r$-fold tensor
products of copies of $M$ and $N$ (with $2^r$ summands). Write $L_{s,t}$ for the
submodule which is the direct sum of all tensor products involving $s$ copies of
$M$ and $t$ copies of $N$ (with $s + t = r$). Thus
\[
  \bigotimes^r(M \oplus N) = \bigoplus_{s + t = r} L_{s,t}.
\]
For example, when $r = 2$ we have $L_{2,0} = M \otimes M$,
$L_{1,1} = (M \otimes N) \oplus (N \otimes M)$ and $L_{0,2} = N \otimes N$. Now let $Q$ be the
submodule of $\bigotimes^r(M \oplus N)$ generated by all elements of the form
$\cdots \otimes x \otimes \cdots \otimes x \otimes \cdots$; we have
$Q = \bigoplus[Q \cap L_{s,t}]$. We see at once that $Q \cap L_{s,t}$ is the submodule generated by
all elements of the forms $\cdots \otimes y \otimes \cdots \otimes y \otimes \cdots$ (with either $y \in M$
or $y \in N$), and $x \otimes y \otimes \xi \otimes y + \xi \otimes y \otimes x \otimes y$ (with $x \in M$ and
$y \in N$). Thus one sees easily that
\[
  \bigwedge^r (M \oplus N) = \Big(\bigotimes^r(M \oplus N)\Big)/Q = \bigoplus_{s + t = r}
      (L_{s,t}/Q \cap L_{s,t}),
\]
and
\[
  L_{s,t}/L_{s,t} \cap Q \cong \Big(\bigwedge^s M\Big) \otimes \Big(\bigwedge^t N\Big).
\]
\end{proof}

\begin{enumerate}[(5)]
  \item Let $A$ be a commutative ring. We say that a (possibly non-commutative)
        $A$-algebra $E$ is a \emph{skew-commutative graded algebra} if it has a direct sum
        decomposition $E = \bigoplus_{n \ge 0} E_n$ as an $A$-module, such that
        \begin{enumerate}[(i)]
          \item $E_pE_q = E_{p+q}$;
          \item $xy = (-1)^{pq}yx$ for $x \in E_p$ and $y \in E_q$;
          \item $x^2 = 0$ for $x \in E_{2k+1}$.
        \end{enumerate}
        For such an algebra $E$, a \emph{skew-derivation} is an $A$-linear map $d\colon E \to E$
        such that
        \begin{enumerate}[(a)]
          \item $d(E_n) \subset E_{n-1}$;
          \item[($\beta$)] $d(xy) = (dx)y + (-1)^px(dy)$ for $x \in E_p$, $y \in E_q$.
        \end{enumerate}
  \item Let $A$ be a ring and $M$ an $A$-module. We show how to define an
        $A$-bilinear map $\psi\colon (\bigwedge^p M) \times (\bigwedge^q M) \to \bigwedge^{p+q} M$. If we
        define $\varphi\colon M^p \times M^q \to \bigwedge^{p+q} M$ by
        $\varphi(x_1, \ldots, x_p, y_1, \ldots, y_q) = x_1 \wedge \cdots \wedge x_p \wedge y_1 \wedge \cdots \wedge y_q$,
        then for fixed $y_1, \ldots, y_q$ this is an alternating $p$-multilinear map from
        $M^p$ to $\bigwedge^{p+q} M$, so that there is a map $\Phi\colon (\bigwedge^p M) \times M^q \to
        \bigwedge^{p+q} M$ such that $\Phi(\xi, y_1, \ldots, y_q)$ is linear in $\xi$ and satisfies
        $\Phi(\xi, y_1, \ldots, y_q) = x_1 \wedge \cdots \wedge x_p \wedge y_1 \wedge \cdots \wedge y_q$ if
        $\xi = x_1 \wedge \cdots \wedge x_p$. Now for fixed $\xi$, $\Phi$ is alternating and multilinear in
        $y_1, \ldots, y_q$, defining a bilinear map $\psi\colon (\bigwedge^p M) \times (\bigwedge^q M) \to
        \bigwedge^{p+q} M$ such that $\psi(\xi, y_1 \wedge \cdots \wedge y_q) = \Phi(\xi, y_1, \ldots, y_q)$. For
        $\xi \in \bigwedge^p M$ and $\eta \in \bigwedge^q M$ we write $\xi \wedge \eta$ for $\psi(\xi, \eta)$; if
        $\xi = \sum_\alpha x_\alpha^{(1)} \wedge \cdots \wedge x_\alpha^{(p)}$ and
        $\eta = \sum_\beta y_\beta^{(1)} \wedge \cdots \wedge y_\beta^{(q)}$ then
        $\xi \wedge \eta = \sum_{\alpha, \beta} x_\alpha^{(1)} \wedge \cdots \wedge x_\alpha^{(p)}
        \wedge y_\beta^{(1)} \wedge \cdots \wedge y_\beta^{(q)}$. (It might be tempting to make the
        definition directly in terms of this formula, but the expression of $\xi$ and
        $\eta$ in the above form is non-unique, so that this requires an awkward
        proof.) The multiplication $\wedge$ satisfies the associative law
        $\xi \wedge (\eta \wedge \zeta) = (\xi \wedge \eta) \wedge \zeta$, so that we can use it to define a
        product on $\bigoplus_{n \ge 0} \bigwedge^n M$ (we set $\bigwedge^0 M = A$), which becomes an
        $A$-algebra, and it is easy to see that this satisfies the conditions for a
        skew-commutative algebra. We write $\bigwedge M$ for this algebra, and call it
        the \emph{exterior algebra} of $M$.

        Given any linear map $\alpha\colon M \to A$, there exists a unique skew-derivation
        $d$ of $\bigwedge M$ such that $d$ coincides with $\alpha$ on
        $\bigwedge^1 M = M$. The uniqueness is clear from the fact that $\bigwedge M$ is
        generated as an $A$-algebra by $M$: we must have
        \[
          d(x_1 \wedge \cdots \wedge x_p) = \sum_{i=1}^{p} (-1)^{i-1}\alpha(x_i)x_1
              \wedge \cdots \wedge \widehat{x_i} \wedge \cdots \wedge x_p.
        \]
        Conversely, the existence follows easily from the fact that the right-hand
        side defines an alternating $p$-fold multilinear map of $x_1, \ldots, x_p$.

        In particular, let $M$ be a free $A$-module of rank $n$ with basis
        $e_1, \ldots, e_n$, so that $M = Ae_1 \oplus \cdots \oplus Ae_n$. Then taking arbitrary
        elements $c_1, \ldots, c_n \in A$, we can define $\alpha\colon M \to A$ by
        $\alpha(e_i) = c_i$, and the skew-derivation of $\bigwedge M$ takes the form
        \[
          d(e_{i_1} \wedge \cdots \wedge e_{i_p}) = \sum_{j=1}^{p} (-1)^{j-1}c_{i_j}
              e_{i_1} \wedge \cdots \wedge \widehat{e_{i_j}} \wedge \cdots \wedge e_{i_p}.
        \]
        This can be identified with the differential operator of the Koszul complex
        $K_{\bullet}^{c_1, \ldots, c_n}$ discussed in \S\ref{sec:16}. Thus the Koszul complex can be
        thought of as the exterior algebra $\bigwedge(Ae_1 \oplus \cdots \oplus Ae_n)$ with the
        skew-derivation defined by $d(e_i) = c_i$.
\end{enumerate}

% ============ 习题 ============
\section*{Exercise to Appendix C}

\noindent\textbf{C.1.}\quad Let $(A, m, k)$ be a local ring and $M$ a finitely generated
$A$-module. Prove that $\min\{r \mid \bigwedge^r M \ne 0\}$ is equal to the minimal number of
generators of $M$.


% ============ 习题解答 ============
\backmatter
\chapter*{Solutions and hints for exercises}
\addcontentsline{toc}{chapter}{Solutions and hints for exercises}
% ============ Solutions and hints for exercises (Part 1: §1-§3) ============
% 原书 pp.287-288; 原书对 §1 的 Ex.1.4 与 §3 的 Ex.3.3 未给解答, 编号按原书跳号.

\noindent\textit{(Please be sure to try each exercise on your own before looking at
the solution)}

\section*{\S\ref{sec:1}}

\noindent\textbf{1.1.}\quad If $ab = 1 - x$ with $x^n = 0$ then $ab(1 + x + \cdots + x^{n-1}) = 1$.

\noindent\textbf{1.2.}\quad Set $e_i = (0, \ldots, 1, \ldots, 0) \in A_1 \times \cdots \times A_n$ (with $1$ in the $i$th place); then since
$e_ie_j = 0$ for $i \neq j$, any prime ideal $p$ of $A_1 \times \cdots \times A_n$ must contain all but one
of the $e_i$.

\noindent\textbf{1.3.}\quad (a) Use the fact that $\rad(A) = \{x \in A \mid 1 + ax \text{ is a unit of } A, \forall a \in A\}$.
Counter-example: $A = \Z$, $B = \Z/(4)$; then $\rad(A) = (0)$, $\rad(B) = 2B$.

\smallskip
\noindent\hspace*{1.5em}(b) Let $m_1, \ldots, m_r$ be the maximal ideals of $A$ and $I = \Ker f$. Suppose
$I \subset m_i$ for $1 \le i \le s$ and $I \nsubseteq m_i$ for $s < i \le r$; then the maximal ideals of $B$
are $f(m_i)$ for $1 \le i \le s$, and $f(m_i) = B$ for $i > s$. Now $\rad(A) = m_1 \cdots m_r$,
hence $f(\rad(A)) = f(m_1) \cdots f(m_r) = f(m_1) \cdots f(m_s) = \rad(B)$.

\noindent\textbf{1.5.}\quad The first half is easy; for the second, use Zorn's lemma.

\noindent\textbf{1.6.}\quad We can assume that there are no inclusions among $P_1, \ldots, P_r$. When $r = 2$,
take $x \in I - P_1$, $y \in I - P_2$; then one of $x$, $y$, $x + y$ is not in $P_1$ or $P_2$. When
$r > 2$, we can take $x \in I - (P_1 \cup \cdots \cup P_{r-1})$ by induction. Also, since $P_r$ is
prime, $P_r \not\supset IP_1 \cdots P_{r-1}$, so take $y \in IP_1 \cdots P_{r-1} - P_r$; then either $x$ or $x + y$
satisfies the condition.

\section*{\S\ref{sec:2}}

\noindent\textbf{2.1.}\quad By NAK there is an $e \in I$ such that $(1 - e)I = 0$. One sees easily that then
$I = Ie = Ae$ and $e^2 = e$.

\noindent\textbf{2.2.}\quad If $x \in \ann(M/IM)$ then $xM \subset IM$, so that by Theorem~\hyperref[thm:2.1]{1} there exists $y \in I$
such that $(x^n + y)M = 0$.

\noindent\textbf{2.3.}\quad $(M + N)/N \cong M/(M \cap N)$ shows that $M$ is finite, and similarly for $N$.

\noindent\textbf{2.4.}\quad (a) If $M \cong A^n$ and $P$ is a maximal ideal of $A$ with $k = A/P$ then $M/PM = k^n$;
for a field the result is well-known.

\smallskip
\noindent\hspace*{1.5em}(b) The first part is easy by the theory of determinants; the second half
comes from the fact that $A^n$ has $n$ linearly independent elements, but any
$n + 1$ elements are linearly dependent.

\smallskip
\noindent\hspace*{1.5em}(c) Use Theorem~\hyperref[thm:2.3]{3}, (iii).

\noindent\textbf{2.5.}\quad (a) If $F$ and $F'$ are free modules and $\alpha\colon F \to L$, $\beta\colon F' \to N$ are surjections
then there is a map $\mu$ making
\[
  \begin{tikzcd}[column sep=large, row sep=large]
    0 \arrow[r] & F \arrow[r] \arrow[d, "\alpha"] & F \oplus F' \arrow[r] \arrow[d, "\mu"] & F' \arrow[r] \arrow[d, "\beta"] & 0 \\
    0 \arrow[r] & L \arrow[r] & M \arrow[r] & N \arrow[r] & 0
  \end{tikzcd}
\]
commute. The assertion follows from this and the snake lemma
(Appendix B).

\smallskip
\noindent\hspace*{1.5em}(b) can be proved similarly.

\section*{\S\ref{sec:3}}

\noindent\textbf{3.1.}\quad Use the fact that $A$ is isomorphic to a submodule of $(A/I_1) \oplus \cdots \oplus (A/I_n)$.

\noindent\textbf{3.2.}\quad Use the previous question.

\noindent\textbf{3.4.}\quad If $II^{-1} = A$ then there exist $x_i \in I$ and $y_i \in I^{-1}$ such that $\sum x_iy_i = 1$; then it
follows easily that $I = \sum x_iA$.

\noindent\textbf{3.5.}\quad If $J$ is a fractional ideal generated by $a_1/b_1, \ldots, a_n/b_n$ with $a_i$ and
$b_i$ coprime, then $J^{-1}$ is the principal fractional ideal generated by $u/v$, where
$v = \lcm(b_1, \ldots, b_n)$ and $u = \hcf(a_1, \ldots, a_n)$.

\noindent\textbf{3.6.}\quad $\Ker(\varphi^n)$, $n = 1, 2, \ldots$, is an ascending chain of ideals of $A$.

\noindent\textbf{3.7.}\quad Choose an ideal $I$ of $A$ which is not finitely generated, and set $M = A/I$;
then by Theorem~\ref{thm:2.6}, $M$ cannot be of finite presentation.

% ============ Solutions and hints for exercises (Part 2: §4-§6) ============
% 原书 pp.288-289; 原书对 §4 的 Ex.4.1-4.4,4.7-4.9,4.12 与 §5 之外及 §6 的 Ex.6.5-6.7、附录 6.8-6.9 未给解答, 编号按原书跳号.

\section*{\S\ref{sec:4}}

\noindent\textbf{4.5.}\quad Write $V(I_\lambda)$ for the complement of $U_\lambda$, where $I_\lambda$ is an ideal of $A$. Then
$\bigcap V(I_\lambda) = V(\sum I_\lambda) = \emptyset$, so $1 \in \sum I_\lambda$, and therefore a finite sum of $I_\lambda$ also
contains $1$.

\noindent\textbf{4.6.}\quad If $\Spec A = V(I_1) \cup V(I_2)$ with $V(I_1) \cap V(I_2) = \emptyset$ then $I_1 + I_2 = A$ and
$I_1I_2 \subset \nil(A)$. So $1 = e_1 + e_2$ with $e_i \in I_i$ for $i = 1, 2$ and $(e_1e_2)^n = 0$. So $1 =
(e_1 + e_2)^{2n} = e_1^{2n}x_1 + e_2^{2n}x_2$ with $x_i \in A$. So $e = e_1^{2n}x_1$ satisfies $e(1 - e) = 0$.

\noindent\textbf{4.10.}\quad For $p \in \Spec A$, if $V(p) = V(a) \cup V(b)$ then $p \in V(p)$ gives $p \supset a$ or $p \supset b$, and
hence either $V(p) = V(a)$ or $V(p) = V(b)$. Conversely, if $V(I)$ is irreducible,
then for $x, y \in A$ with $xy \in \sqrt{I}$, from $V = V(I + Ax) \cup V(I + Ay)$ we have,
say, $V = V(I + Ax)$, and $x \in \sqrt{I}$, this proves $\sqrt{I} \in \Spec A$.

\noindent\textbf{4.11.}\quad If there is a closed subset which cannot be so expressed, let $V$ be a minimal
one. Then $V$ must be reducible, but if $V = V_1 \cup V_2$ with $V_i \neq V$ then, by
minimality, each of $V_1$ and $V_2$ is a union of a finite number of irreducible
closed sets, hence also $V$, a contradiction.

\section*{\S\ref{sec:5}}

\noindent\textbf{5.1.}\quad Set $k[X_1, \ldots, X_n]/p = k[x_1, \ldots, x_n]$; then by Theorem~\hyperref[thm:5.6]{6}, $\coht p =
\trdeg_k k(x)$. Suppose this is $r$, and that $x_1, \ldots, x_r$ is a transcendence basis
of $k(x)$ over $k$, and set $K = k(X_1, \ldots, X_r)$; then $k[X_1, \ldots, X_n]_p$ is the
localisation of $K[X_{r+1}, \ldots, X_n]$ at a prime ideal $P$, with $\htop p = \htop P$. This
reduces us to proving that if $r = 0$ then $\htop p = n$. In this case $x_1, \ldots, x_n$ are
algebraic over $k$, and letting $p_i$ be the kernel of $k[X_1, \ldots, X_n] \to
k[x_1, \ldots, x_i, X_{i+1}, \ldots, X_n]$ we get a strictly increasing chain $0 \subset p_1 \subset p_2
\subset \cdots \subset p_n = p$, giving $\htop p \ge n$, but by the corollary to Theorem~\hyperref[thm:5.6]{6},
$\htop p \le n$.

\noindent\textbf{5.2.}\quad If $A$ is a zero-dimensional Noetherian ring then all prime ideals are
minimal, and by Ex.~\ref{ex:4.12-ex}, there are only finitely many of these. Let these be
$p_1, \ldots, p_n$; then since $p_1 \cdots p_n = \nil(A)$ there is an $n$ such that $(p_1 \cdots p_n)^n = 0$.
For any ideal $I$ and any $i$, the module $I/Ip_i$ is a finite-dimensional vector
space over $A/p_i$, so that $l(I/Ip_i) < \infty$. It follows easily that $l(A) < \infty$, so that
$A$ is Artinian.

\section*{\S\ref{sec:6}}

\noindent\textbf{6.1.}\quad $\Ass M = \{(0), (3)\}$. ($\supset$ is obvious, $\subset$ from Theorem~\hyperref[thm:6.3]{3}.)

\noindent\textbf{6.2.}\quad No. Let $M$ be as in the previous question, $M_1 = \{(a, 0) \mid a \in \Z\}$ and
$M_2 = \{(0, a) \mid a \in \Z\}$; then $M = M_1 + M_2$, but each $M_i \cong \Z$.

\noindent\textbf{6.3.}\quad Since $xA/x^nA \cong A/x^{n-1}A$, there is an exact sequence
\[
  0 \to A/x^{n-1}A \to A/x^nA \to A/xA \to 0.
\]

\noindent\textbf{6.4.}\quad Use a primary decomposition of $I$.

% ============ Solutions and hints for exercises (Part 3: §7-§9) ============
% 原书 pp.289-291; 原书对 §7 的 Ex.7.1,7.6,7.7、§8 的 Ex.8.5,8.6、§9 的 Ex.9.5,9.7,9.9,9.11 未给解答, 编号按原书跳号.

\section*{\S\ref{sec:7}}

\noindent\textbf{7.2.}\quad For $b \in B$ write $b = y/x$ with $x, y \in A$. Then $y = bx \in Bx \cap A = xA$
(by Theorem~\ref{thm:7.5}, (ii)), so $b \in A$.

\noindent\textbf{7.3.}\quad Write $N \subset M$ for the $A$-submodule generated by $\{m_\lambda\}$; then $B \otimes
(M/N) = 0$, so $M/N = 0$.

\noindent\textbf{7.4.}\quad Set $M = \prod M_\lambda$. It is enough to show that $I \otimes M \to IM$ is injective
for an ideal $I = (a_1, \ldots, a_n)A$ of $A$ (by Theorem~\hyperref[thm:7.6]{6}). Define $f\colon A^n \to A$ by
$f(x_1, \ldots, x_n) = \sum a_ix_i$, and set $K = \Ker f$. Then $0 \to K \to A^n \to A$ is
exact, hence also $0 \to K \otimes M \to M^n \to M$, and if $\sum a_i \otimes \xi_i \in I \otimes M$ satisfies
$\sum a_i\xi_i = 0$ then writing $\xi_{i\lambda}$ for the $\lambda$-th coordinate of $\xi_i \in M$ we have
$\sum a_i\xi_{i\lambda} = 0$ for all $\lambda$, and hence $(\xi_{1\lambda}, \ldots, \xi_{n\lambda}) \in K \otimes M$.
Now since $A$ is Noetherian we can write $K = A\beta_1 + \cdots + A\beta_r$ with $\beta_j
= (b_{1j}, \ldots, b_{nj}) \in K \subset A^n$ for $1 \le j \le r$. Thus we can write $\xi_{i\lambda} = \sum_j b_{ij}\eta_{j\lambda}$
with $\eta_{j\lambda} \in M_\lambda$. Since $\sum_i a_ib_{ij} = 0$, setting $\eta_j = (\eta_{j\lambda})_\lambda \in M$, we get
$\xi_i = \sum_{j=1}^{r} b_{ij}\eta_j$, and $\sum_i a_i \otimes \xi_i = \sum_{i,j} a_ib_{ij} \otimes \eta_j = 0$.

\noindent\textbf{7.5.}\quad Tensor the exact sequence $0 \to A \xrightarrow{a} A$ with $N$ to get the exact sequence
$0 \to N \xrightarrow{a} N$.

\noindent\textbf{7.8.}\quad Tensor product does not commute with infinite direct products. If $\{p_i\}$ is
an infinite set of prime numbers then $\bigcap_i p_i\Z = (0)$, but $\bigcap_i (p_i\Z \otimes \Q) \neq 0$.

\noindent\textbf{7.9.}\quad If $I$ is an ideal of $A$ then $IB \cap A = I$, so that given a chain $I_1 \subset I_2 \subset \cdots$ of
ideals of $A$, $I_1B \subset I_2B \subset \cdots$ eventually terminates, hence so does the given
chain.

\section*{\S\ref{sec:8}}

\noindent\textbf{8.1.}\quad $(I+J)^{2n} \subset I^n + J^n$, so that given $x_1, x_2, \ldots$ such that $x_{n+1} - x_n \in (I+J)^{2n}$
we can write $x_{n+1} - x_n = u_n + v_n$ with $u_n \in I^n$ and $v_n \in J^n$. Thus $\{x_n\}$
converges to $x_1 + \sum u_i + \sum v_i$. Also, $I, J \subset \rad(A)$, so that $I + J \subset \rad(A)$, and
therefore $\bigcap_n (I + J)^n = (0)$.

\noindent\textbf{8.2.}\quad If $\{x_n\}$ satisfies $x_{n+1} - x_n \in J^n$ then there is a limit $x$ for the $I$-adic topology.
For any $i$, taking $m$ large enough we get $x_m - x \in I^i$, so that $x_n - x = x_n - x_m + x_m - x
\in J^n + I^i$, and since by Theorem~\hyperref[thm:8.10]{10}, (i) we have $\bigcap_i (J^n + I^i) = J^n$, we get
$x_n - x \in J^n$, so that $x$ is also a $J$-adic limit.

\noindent\textbf{8.3.}\quad Let $a\hat{A} = x\hat{A}$ and $x = \sum a_i\hat{x}_i$ with $a_i \in a$ and $\hat{x}_i \in \hat{A}$. Let $I$ be an ideal of
definition of $A$. Take $x_i \in A$ such that $x_i - \hat{x}_i \in I\hat{A}$ and set $a = \sum a_ix_i$. Then
$a \in a$ and $a\hat{A} \subset a\hat{A} + I\hat{A}$, so that by NAK, $x\hat{A} = a\hat{A}$, so $a = x\hat{A} \cap A = aA$.

\noindent\textbf{8.4.}\quad $e = (e - X)(e + eX + eX^2 + \cdots)$.

\noindent\textbf{8.7.}\quad $A/m^n$ is Artinian, so that there exists $t(n)$ such that $a_{t(n)} + m^n = a_j + m^n$ for
$j > t(n)$. We can assume that $t(n) < t(n+1) < \cdots$. Supposing that $a_{t(n)} \nsubseteq m^r$ for some
$r$, then we take $a_r \in a_{t(r)} - m^r$, then $a_{r+1} \in a_{t(r+1)}$ such that $a_{r+1} - a_r \in m^r$, and
proceed in the same way taking $a_i \in a_{t(i)}$ such that $a_i - a_{i-1} \in m^{i-1}$ for $i \ge r$. Then
$\lim a_i$ belongs to $\bigcap_\nu a_\nu$, but not to $m^r$, which is a contradiction.

\noindent\textbf{8.8.}\quad This can be done by following the proof of Theorem~\hyperref[thm:8.5]{5}, and replacing the
use of homogeneous polynomials by multihomogeneous polynomials.

\noindent\textbf{8.9.}\quad We can assume that $A$ is a Noetherian local ring with maximal ideal $P$.
There is an $x \neq 0$ such that $xP = 0$, and then since $\bigcap_n P^n = (0)$ there exists $c$
such that $x \notin P^c$. Then if $I \subset P^c$, $I:x = P$.

\noindent\textbf{8.10.}\quad Let $m = (X, Y) \subset k[X, Y]$ and set $A = k[X, Y]_m$; let $q(X) = \sum c_iX^i \in k[[X]]$ be
transcendental over $k(X)$ and set $a_\nu = (X^{\nu+1}, Y - 1 - c_1X - \cdots - c_\nu X^\nu)$.

\section*{\S\ref{sec:9}}

\noindent\textbf{9.1.}\quad $B_p$ is integral over $A_p$, so that any maximal ideal of $B_p$ lies over $pA_p$ and
therefore coincides with $PB_p$. Hence $B_p$ is a local ring, and the elements of $B - P$
are units of $B_p$.

\noindent\textbf{9.2.}\quad $\dim A \le \dim B$ from the going-up theorem, and $\dim B \le \dim A$ from Theorem~\hyperref[thm:9.3]{3}, (ii).

\noindent\textbf{9.3.}\quad Replacing $A$ and $B$ by $A_p$ and $B_p$ we can assume $p$ is maximal; then set
$k = A/p$, so that $B/pB$ is a finite $k$-module, hence an Artinian ring.

\noindent\textbf{9.4.}\quad If $ax^n \in A$ for all $n$ then $A[x]$ is a submodule of the finite $A$-module $a^{-1}A$; if
$A$ is Noetherian then $A[x]$ is also a finite $A$-module.

\noindent\textbf{9.6.}\quad Suppose $f = gh$ with $g, h \in K[X]$ monic. Roots of $g$ are roots of $f$, hence
integral over $A$, and expressing the coefficients of $g$ in terms of the roots, we
have that the coefficients of $g$ are integral over $A$; since $A$ is integrally closed,
$g \in A[X]$, and similarly for $h$.

\noindent\textbf{9.8.}\quad By Theorem~\hyperref[thm:9.3]{3}, (ii).

\noindent\textbf{9.10.}\quad $L[X]$ is a free module over $K[X]$, hence flat; and if $L$ is algebraic over $K$
then $L[X]$ is integral over $K[X]$. The first part follows from the previous
two questions, together with Theorem~\hyperref[thm:9.5]{5}. If $f, g$ have a common factor $\alpha(X)$
in $L[X]$ then set $P = (\alpha(X))$, so that $\htop P = 1$ (it can easily be seen that a
non-zero principal prime ideal in a Noetherian integral domain has height
$1$). Hence $\htop p \le 1$. But $f, g \in p$, so that $\htop p = 1$. There is an irreducible
divisor $h$ of $f$ in $p$, and $p = (h)$, so $h \mid g$.

% ============ Solutions and hints for exercises (Part 4: §10-§12) ============
% 原书 pp.290-292; 原书对 §10 的 Ex.10.1,10.6,10.7,10.10,10.11、§11 的 Ex.11.3,11.4,11.8-11.10、
% §12 的 Ex.12.3 未给解答, 编号按原书跳号.

\section*{\S\ref{sec:10}}

\noindent\textbf{10.2.}\quad Use the previous question and Theorem~\ref{thm:7.7}.

\noindent\textbf{10.3.}\quad In the proof of Theorem~\hyperref[thm:10.4]{4} we can choose $p$ to contain $(m, y)A[y]$.

\noindent\textbf{10.4.}\quad Let $0 \subset p_1 \subset p_2$ be a strictly increasing chain of prime ideals of $R$ and let
$0 \neq b \in p_1$, $a \in p_2 - p_1$; thus $ba^{-n} \in R$ for all $n \ge 0$. Take $F = \sum u_iX^i$ to be a
root of $f^2 + af + X = 0$. Then $u_1 = -a^{-1}$, and for all $i$ we have
$u_i \in a^{-2i+1}R$, so that $bF(X) \in R[[X]]$ but $F(X) \notin R[[X]]$.

\noindent\textbf{10.5.}\quad The first part comes from Theorem~\hyperref[thm:10.1]{1}. For the second part, by \S\ref{sec:9}, Lemma~\hyperref[lem:9.1]{1}
the integral closure of $R$ in $K$ is not the whole of $K$, and therefore coincides
with $R$, so $R$ is integrally closed; on the other hand, for $x \in K - R$ we have
$R[x] = K$, hence $x^{-1} \in R[x]$, so that $x^{-1}$ is integral over $R$ and hence in $R$.
Thus $R$ is a valuation ring. If $\dim R > 1$ then there is a prime ideal $p$ of $R$
distinct from $(0)$ and from the maximal ideal, and thus $R_p$ is intermediate
between $R$ and $K$.

\noindent\textbf{10.8.}\quad Let $v\colon L^* \to G$ be the additive valuation corresponding to $S$, and choose
$x_1, \ldots, x_e \in L$ such that $v(x_1), \ldots, v(x_e)$ represent the different cosets of $G'$ in
$G$, and $y_1, \ldots, y_f \in S$ such that their images in $k$ are linearly independent
over $k'$. It follows easily from the previous question that the $ef$ elements
$x_iy_j$ are linearly independent over $K$.

\noindent\textbf{10.9.}\quad If $S \subset S_1$ then the residue field $k_1$ of $S_1$ contains a valuation ring $A \neq k_1$
such that $S$ is the composite of $S_1$ and $A$. We have $k \subset A \subset k_1$, but by the
previous question $k_1$ is an algebraic extension field of $k$, hence integral
over $A$. But $A$ is integrally closed, therefore $A = k_1$, a contradiction.

\section*{\S\ref{sec:11}}

\noindent\textbf{11.1.}\quad Let $B$ be a valuation ring of $\bar{K}$ dominating $A$, and $G$ its value group. Then
for $\alpha \in m_B$ we have $v(\alpha)/n \in G$ for all $n \ge 1$, so that $G$ has no minimal element. Also it is
easy to see that some multiple of $v(\alpha)$ belongs to $v(K^*)$, so that $G$ is
Archimedean.

\noindent\textbf{11.2.}\quad If $B$ is a valuation ring of $L$ dominating $A$ and $G$ its value group, set
$H = v(K^*)$ and $e = [G:H]$. Then $x \in G \Rightarrow e \cdot x \in H$. Hence $G$ is isomorphic to a
subgroup of $H$, and $G \cong \Z$.

\noindent\textbf{11.5.}\quad Just use Forster's theorem (5.7).

\noindent\textbf{11.6.}\quad By Ex.~\ref{ex:9.7-ex}, $A = \Z[\sqrt{10}] = \Z[X]/(X^2 - 10)$. Then $A/3A = \Z[X]/(3,
X^2 - 1) = (\Z/3\Z)[X]/(X - 1)(X + 1)$, so that $P = (3, \sqrt{10} - 1)$ is a prime
ideal of $A$. This is not principal, since if $P = (\alpha)$ with $\alpha = a + b\sqrt{10}$, then
one and sees easily that the norm $N(\alpha) = a^2 - 10b^2$ would have to be $\pm 3$,
but this is impossible since the congruence $a^2 \equiv \pm 3 \bmod 5$ has no solution.

\noindent\textbf{11.7.}\quad Let $P_1, \ldots, P_r$ be the maximal ideals of $A$; choose an element $\alpha \in P_1$ such
that $\alpha \notin P_2 \cup P_3 \cup \cdots \cup P_r$. Then $\alpha A = P_1$, and similarly each of the prime
ideals is principal. Thus by Theorem~\hyperref[thm:11.6]{6} any ideal is principal. (Of course
this also follows from Theorem~\ref{thm:5.8}.)

\section*{\S\ref{sec:12}}

\noindent\textbf{12.1.}\quad Suppose that $L$ is normal over $K$, and let $G = \Aut_K(L)$. Let $S_1$ and $S$ be
valuation rings of $L$ dominating $R$, and let $S_1, \ldots, S_l$ be the conjugates
of $S_1$ by elements of $G$, and $A = S_1 \cap \cdots \cap S_l \cap S$. If $S \neq S_i$ for any $i$ then by
Ex.~\ref{ex:10.9-ex}, there are no inclusions among $S_1, \ldots, S_l, S$, and we can apply
Theorem~\hyperref[thm:12.2]{2}. Letting $n, n_i$ be the maximal ideals of $S, S_i$, and setting
$n \cap A = p$, $n_i \cap A = p_i$ we have $p_1 \cdots p_l \not\subset p$, so that we can choose
$x \in p_1 \cap \cdots \cap p_l$ with $x \notin p$. Then $x \notin n$, but since $x \in n_i$ for all $i$, all the
conjugates of $x$ over $K$ belong to $n_1$, and the coefficients of the minimal
polynomial of $x$ over $K$ belong to $n_1 \cap K = \rad(R)$. Thus it follows easily
that $x \in n$, a contradiction.

\noindent\textbf{12.2.}\quad By Ex.~\ref{ex:10.3-ex}, $\bar{R}$ is the intersection of all valuation rings of $L$ dominating $R$.
Ex.~\ref{ex:10.9-ex} can easily be extended to the infinite case, so that the second part
follows from the first. For the first part, reduce to the finite case and use the
previous question and Theorem~\hyperref[thm:12.2]{2}.

\noindent\textbf{12.4.}\quad Let $\mathcal{P}$ be the set of height $1$ prime ideals of $A$, and for $p \in \mathcal{P}$ set $I_p = a_pA_p$.
Then $x \in \tilde{I} \iff xI^{-1} \subset A_p$ for all $p \in \mathcal{P} \iff x \in a_pA_p$ for all $p \in \mathcal{P}$. Hence $\tilde{I}$ is the
intersection of $I_p \cap A$ taken over the finitely many $p \in \mathcal{P}$ such that $I_p \neq A_p$.

% ============ Solutions and hints for exercises (Part 5: §13-§15) ============
% 原书 pp.292-293; 原书对 §13 的 Ex.13.1,13.2,13.5、§14 的 Ex.14.1-14.6、§15 的 Ex.15.3 未给解答, 编号按原书跳号.

\section*{\S\ref{sec:13}}

\noindent\textbf{13.3.}\quad Let $P \in \Ass(A)$ be such that $\htop P \ge 1$, and let $p_1, \ldots, p_r$ be the prime
divisors of $(a)$. If $P \neq p_i$ for all $i$, then there exists $x \in P$ such that $(a):x = (a)$.
This is a contradiction, since $x$ is a zero-divisor, but if $xy = 0$ then $y \in \sqrt{(a)A}$.

\noindent\textbf{13.4.}\quad (i) is easy. (ii) The homogeneous elements of $P$ are nilpotent mod $Q^*$,
hence so are all elements of $P$. Now we show that if $f \notin P$, $g \notin Q^*$ then
$fg \notin Q^*$. Let $f = f_1 + \cdots + f_r$, $g = g_1 + \cdots + g_s$ with $f_i$ and $g_j$ homogeneous,
and $\deg f_1 < \deg f_2 < \cdots$, $\deg g_1 < \deg g_2 < \cdots$; we work by induction on
$r + s$. If $r = s = 1$ there is nothing to prove. Also, since we can assume that
$g_1 \notin Q^*$ we have $g_1 \notin Q$. If $f_1 \notin P$ then $f_1g_1 \notin Q^*$. Next, suppose $f_1 \in P$. If
$f_1g_1 \in Q^*$ then $fg \in Q^*$, since $(f_1 + \cdots + f_r)g \in Q^*$. If $f_1g_1 \notin Q^*$ then
$f_1^tg_1 \in Q^*$ and $f_1^{t-1}g_1 \notin Q^*$ for some $t > 1$ (since $f_1^n \in Q^*$ for $n \gg 0$).
Replacing $g$ by $f_1^tg$ reduces to the case $f_1g_1 \in Q^*$, so that $f_1^tg_1 \notin Q^*$.

\noindent\textbf{13.6.}\quad First half: let $S$ be the multiplicative set made up of homogeneous elements
of $R$ not in $P$; then $R_S/P^*R_S$ can be viewed as the localisation of $R/P^*$ with
respect to all non-zero homogeneous elements, and by the previous question
this is $\cong K[X, X^{-1}]$, which is a one-dimensional ring. Second half: proof by
induction on $\htop P = n$; take a prime ideal $Q \subset P$ with $\htop Q = n - 1$. If $Q \neq P^*$
then $Q$ is inhomogeneous, $Q^* \subset P^*$ and $\htop(P/Q^*) \ge 2$, so by the first half,
$P^* \neq Q^*$, hence $\htop P^* \ge \htop Q^* + 1 = n - 1$.

\section*{\S\ref{sec:14}}

\noindent\textbf{14.7.}\quad Let $p \in \Spec A$, $f, g \notin p$, and let $l = \dim A/p$. If $l = 1$ then $pA_f$ is a maximal
ideal. Suppose that $l > 1$. We can choose $x_2, \ldots, x_r \in m$ such that
$\htop(p, f, x_2, \ldots, x_i)/p = \htop(p, g, x_2, \ldots, x_i)/p = i$ for $2 \le i \le r$. Then any mini-
mal prime divisor $P$ of $(p, x_2, \ldots, x_r)$ satisfies $\dim A/P = l$, $f \notin P$ and $g \notin P$,
so $g \notin PA_f \in \text{m-Spec}(A_f)$.

\section*{\S\ref{sec:15}}

\noindent\textbf{15.1.}\quad Set $Z = X/Y$ so that $X = YZ$, $B = k[Y, Z] \supset A = k[YZ, Y]$, and $pB =
YB$. So $B/pB \cong k[Z]$ and $\dim B_P/pB_P = 1$. Now let $p' = (X - \alpha Y)A$ for
$0 \neq \alpha \in k$; then any height $1$ prime ideal of $B$ containing $X - \alpha Y =
Y(Z - \alpha)$ must be $YB$ or $(Z - \alpha)B$, but since $YB \cap A \neq p'$ and $(Z - \alpha)B \cap A \neq p'$
there does not exist any prime ideal of $B$ contained in $P$ and lying over $p'$.

\noindent\textbf{15.2.}\quad No. Set $f = XY - 1$. Then $fB$ is a prime ideal of $B$, and $fB \cap A = (0)$.
Since $fB + XB = B$ there does not exist any prime ideal of $B$ containing
$fB$ and lying over $XA$. Note that the fibres of $A \to B$ are all one-
dimensional.

% ============ Solutions and hints for exercises (Part 6: §16-§18) ============
% 原书 pp.293-294; 原书对 §16 的 Ex.16.4,16.6,16.8,16.9、§17 的 Ex.17.6、§18 的 Ex.18.2,18.7,18.8 未给解答, 编号按原书跳号.

\section*{\S\ref{sec:16}}

\noindent\textbf{16.1.}\quad $\le$ is easy; for $\ge$ consider a system of parameters of $M'$.

\noindent\textbf{16.2.}\quad $\Hom_A(A/\mathfrak{a}, A/\mathfrak{b}) = 0$ by Theorem~\hyperref[thm:16.9]{9}.

\noindent\textbf{16.3.}\quad For $P \in \Ass(A/I)$, $\grade P \ge k$ is clear from $P \supset I$. If $\grade P > k$ then by
Ex.~\ref{ex:16.2-ex}, $I:P = I$, which is a contradiction.

\noindent\textbf{16.5.}\quad Suppose that $(A, m)$ is local; then if $m \in \Ass(A)$ we have $\depth A = 0$, but if
$\htop P > 0$ and $P \notin \Ass(A)$ then $\depth A_P > 0$, so that $M = A$ gives a counter-
example. For example, $A = k[X, Y, Z]_{(X,Y,Z)}/(X, Y, Z)^2 \cap (Z)$ and $P = (x, z)A$
satisfy these conditions.

\noindent\textbf{16.7.}\quad Let $\xi = (x_1, \ldots, x_r)$ be a maximal $M$-sequence in $m$, and set $M' = M/\sum x_iM$,
then there exists $0 \neq \xi' \in M'$ such that $m\xi' = 0$. Thus $n\xi' = 0$, but $n^\nu \subset mB$
so $n^\nu\xi' = 0$ and $n \in \Ass_B(M')$, therefore $\xi$ is also a maximal $M$-sequence in $n$.

\noindent\textbf{16.10.}\quad (i) For $r = 1$ the proof is similar to Theorem~\ref{thm:14.3}. If $r > 1$, applying the case
$r = 1$ gives that $(a_1, \ldots, a_{r-1})$ is prime, and we can then use induction. (ii) For
any $Q \in \Ass(A)$ we have $QA_P \in \Ass(A_P)$; if we had $P \subset Q$ then by (i), $A_Q$
is an integral domain, a contradiction. Hence $P \nsubseteq Q$. Therefore using
Ex.~\ref{ex:16.8-ex}, we see that $P$ can be generated by an $A$-sequence. For a counter-
example let $A = k[x, y, z] = k[X, Y, Z]/(X(1 - YZ))$; $P = (x, y, z) = (y, z) = (y - y^2z, z)$
is a prime ideal of height $2$, but $y - y^2z$ is a zero-divisor in $A$.

\section*{\S\ref{sec:17}}

\noindent\textbf{17.1.}\quad (b) Let $k$ be a field; then $A = k[X, Y]/(XY, Y^2)$ is a one-dimensional
ring which is not CM.

\noindent\textbf{17.2.}\quad $x^3, y^3$ is an $A$-sequence, hence also an $R$-sequence, so that $R$ is CM. The
ring $k[x^4, x^3y, xy^3, y^4]$ is not CM.

\noindent\textbf{17.3.}\quad By localisation we need only consider the case of an integral domain, and
it then follows from Theorem~\ref{thm:11.5}, (i).

\noindent\textbf{17.4.}\quad We can assume that $A$ is a local ring. Since $A/J$ is CM we have $\dim
A/J = \depth A/J = r$, and if we set $k$ for the residue class field of $A$ then
$\Ext^i_A(k, A/J) = 0$ for $i < r$. Using the exact sequence $0 \to
J^\nu/J^{\nu+1} \to A/J^{\nu+1} \to A/J^\nu \to 0$ and the fact that $J^\nu/J^{\nu+1}$ is isomorphic
to a direct sum of a number of copies of $A/J$ we get by induction that
$\Ext^i_A(k, A/J^\nu) = 0$ for $i < r$.

\noindent\textbf{17.5.}\quad (i) Let $x_1, \ldots, x_s$ be a maximal $A$-sequence in $P$, and extend to a maximal $A$-
sequence in $m$, $x_1, \ldots, x_s, y_1, \ldots, y_t$. There exists $Q \in \Ass_A(A/(x_1, \ldots, x_s))$
containing $P$, so that by Theorem~\hyperref[thm:17.2]{2}, $\dim A/P \ge \dim A/Q \ge \depth
A/(x_1, \ldots, x_s) = s$. (ii) $\dim A - \htop P \ge \dim (A/P) \ge \depth A - \depth (P, A) \ge \depth A -
\depth A_P$.

\section*{\S\ref{sec:18}}

\noindent\textbf{18.1.}\quad Using Ex.~\ref{ex:16.1-ex}, we see that $A$ is CM $\iff B$ is CM. Assuming CM, we need
only use condition ($5'$) of Theorem~\hyperref[thm:18.1]{1} as a criterion.

\noindent\textbf{18.3.}\quad Given a prime ideal $P$ of $A[X]$, by localising $A$ at $P \cap A$ and factoring out
by a system of parameters we reduce to proving that if $(A, m, k)$ is a zero-
dimensional Gorenstein local ring, and $P$ a prime ideal of $A[X]$ such that
$P \cap A = m$ then $B = A[X]_P$ is Gorenstein. Then $P$ is generated by $m$ together
with a monic polynomial $f(X)$, and the image of $f$ in $k[X]$ is irreducible.
Since $f$ is $B$-regular, if we set $C = B/(f)$ then the maximal ideal of $C$ is
$mC$, and $C \cong A[X]/(f)$; this is a free $A$-module of finite rank, so that
$\Hom_A(C/mC, C) = \Hom_A(k, A) \otimes_A C$ (by Ex.~\ref{ex:7.7-ex}) $\cong C/mC$. So $C$ is Goren-
stein, therefore $B$ also.

\noindent\textbf{18.4.}\quad $R/(x^3, y^3) = k[x^3, y^3, x^2y, xy^2]/(x^3, y^3) \cong k[U, V]/(U^2, V^2, UV)$. In this
ring $(0)$ is not irreducible, so that $R$ is not Gorenstein.

\noindent\textbf{18.5.}\quad For $0 \neq a \in A$, the ideal $aA$ is $\cong A/I$ with $I \neq A$, so there exists a non-zero
map $\varphi\colon aA \to k$; viewed as a map $aA \to E$, this extends to $A \to E$, so that
$0 \neq \Im \varphi \subset aE$.

\noindent\textbf{18.6.}\quad We can consider $M$ as a submodule of $E = E_A(k)$. By faithfulness,
$A \subset \Hom_A(M, M) \subset \Hom_A(M, E)$. But $0 \to \Hom_A(E/M, E) \to
\Hom_A(E, E) = A \to \Hom_A(M, E) \to 0$ is exact, hence $E/M = 0$.

% ============ Solutions and hints for exercises (Chapter 7: §19-§21) ============
% 原书 pp.294-295; 原书对 §19 的 Ex.19.1,19.2、§20 的 Ex.20.1,20.2,20.3,20.6 未给解答, 编号按原书跳号.

\section*{\S\ref{sec:19}}

\noindent\textbf{19.3.}\quad If $0 \to P_r \to \cdots \to P_0 \to M \to 0$ is an exact sequence and each $P_i$ is finite
and projective, and if $P_r \oplus A^m = A^n$, then $0 \to P_r \oplus A^m \to P_{r-1} \oplus A^m \to \cdots \to P_0 \oplus A^m \to M \to 0$
is again exact, with $P_r \oplus A^m$ free. Proceeding in the same
way, adding a free module to $P_i$ at each stage, we get an FFR $0 \to L_{r+1}
\to L_r \to \cdots \to L_0 \to M \to 0$.

\noindent\textbf{19.4.}\quad For every maximal ideal $m$ of $A$, since the $A$-module $A/m$ has an FFR, also
the $A_m$-module $A_m/mA_m$ has an FFR, so that the projective dimension over
$A_m$ is finite. Thus $A_m$ is a regular local ring.

\section*{\S\ref{sec:20}}

\noindent\textbf{20.4.}\quad Use Theorem~\hyperref[thm:20.5]{5} and Ex.~\ref{ex:8.3-ex}.

\noindent\textbf{20.5.}\quad This follows from Theorem~\hyperref[thm:20.1]{1}. If an ideal $I$ is locally principal it is finitely
generated, and principal by Theorem~\ref{thm:5.8}.

\section*{\S\ref{sec:21}}

\noindent\textbf{21.1.}\quad For $P \in \Spec R$ with $P \supset I$ set $P/I = p$; then $A_p$ is c.i. $\iff I_p$ is generated by an
$R_p$-sequence $\iff (I/I^2)_p$ is free over $A_p$. Now $I/I^2$ is a finite $A$-module, so that
by Theorem~\ref{thm:4.10}, $\{p \in \Spec A \mid (I/I^2)_p \text{ is free}\}$ is an open subset of $\Spec A$.

\noindent\textbf{21.2.}\quad We can assume that $A$ is complete. Then $A = R/I$ with $R$ regular and
$\dim R = \dim A + 1$. Now $\htop I = 1$ and $A$ is a CM ring, so that all the prime
divisors of $I$ have height 1. Since $R$ is a UFD, $I$ is principal.

\noindent\textbf{21.3.}\quad $A = k[x, y, z] = k + kx + ky + kz + kx^2$, with $x^2 = y^2 = z^2$ and $xy = yz =
zx = 0$. Therefore $0:m = kx^2$, and $A$ is a zero-dimensional Gorenstein
ring. Set $I = (X^2 - Y^2, Y^2 - Z^2, XY, YZ, ZX)$ and $M = (X, Y, Z)$; then
$I \to M^2/M^3$ has five-dimensional image, so that at least five elements
are needed to generate $I$.

% ============ Solutions and hints (Part 8: §22, §24) ============
% 原书 pp.295-296; 原书对 §22 的 Ex.22.1, 22.4、§23 的 Ex.23.1-23.3 未给解答, 编号按原书跳号.

\section*{\S\ref{sec:22}}

\noindent\textbf{22.2.}\quad Algebraic independence comes from Theorem~\hyperref[thm:16.2]{16.2}, (i).
To prove flatness, setting $I = \sum x_iC$ and using Theorem~\hyperref[thm:22.3]{3}, we reduce to
proving that $\Tor^C_1(k, A) = 0$. Since $x$ is a $C$-sequence, the Koszul complex
$L_\bullet = K_\bullet(x, C)$ constructed from $C$ and $x$ is a free resolution of the
$C$-module $k = C/I$, and $\Tor^C_1(k, A) = H_1(L_\bullet \otimes_C A)$. However,
$L_\bullet \otimes_C A$ is just the Koszul complex constructed from $A$ and $\bar{x}$, and since
$\bar{x}$ is an $A$-sequence, $H_1(L_\bullet \otimes A) = 0$.

\noindent\textbf{22.3.}\quad We need only show that $\Tor^A_1(k, M) = 0$. By Lemma~\hyperref[lem:18.2]{2} of
\S\ref{sec:18}, $\Tor^A_1(k, M) = \Tor^{A/xA}_1(k, M/xM)$, but by assumption the right-hand
side is $0$.

\section*{\S\ref{sec:24}}

\noindent\textbf{24.1.}\quad Use $0 \to I^i/I^{i+1} \to A/I^{i+1} \to A/I^i \to 0$ to deduce that
$\Ass(A/I^i) = \Ass(A/I)$ for all $i$.

\noindent\textbf{24.2.}\quad By Theorem~\hyperref[thm:24.5]{5}, it is enough to show that for a prime ideal $p$ of
$A$, $\CM(A/p)$ contains a non-empty open. Let $P$ be the inverse image of $p$ in $R$,
so that $A/p = R/P$. If $x_1, \ldots, x_r \in P$ are chosen to form a system of
parameters of $R$, then since $R_P$ is CM, they form an $R_P$-sequence. Thus passing
to a smaller neighbourhood of $P$, we can assume (i) $P$ is the unique minimal
prime divisor of $(x) = (x_1, \ldots, x_r)R$, and (ii) $\bar{x}$ is an $R$-sequence. Now
replacing $R$ by $R/(x)$, we can assume that $P$ is nilpotent; moreover, we can
take $P^i/P^{i+1}$ to be free $R/P$-modules. Now using the previous question it
follows easily that $R$ is CM implies $R/P$ is CM.

\noindent\textbf{24.3.}\quad After a preliminary reduction as in the previous question, use the
proof of Theorem~\hyperref[thm:24.6]{6}.

% ============ Solutions and hints (Part 9: §25, §26) ============
% 原书 p.296; 原书对 §25 的 Ex.25.1, 25.3 未给解答, 编号按原书跳号.

\section*{\S\ref{sec:25}}

\noindent\textbf{25.2.}\quad If $xy = 0$ we prove that $y \in \bigcap_n x^nA$; suppose that $y \in x^nA$,
and set $y = x^nz$; then $0 = D(x^{n+1}z) = (n + 1)x^nz + x^{n+1}Dz$, and
$y \in x^{n+1}A$.

\noindent\textbf{25.4.}\quad $0 \to I \to A \otimes_k A \to A \to 0$ is a split exact sequence, so that
$0 \to I \otimes_k k' \to A \otimes_k A \otimes_k k' = A' \otimes_{k'} A' \to A' \to 0$ is also
exact, and hence $\Omega_{A'/k'} = (I/I^2) \otimes_k k' = \Omega_{A/k} \otimes_k k'$. For $A_S$, use the
fact that $A_S$ is 0-\'etale over $A$ and Theorem~\hyperref[thm:25.1]{1}.

\noindent\textbf{25.5.}\quad See Theorem~\hyperref[thm:27.3]{27.3}.

\section*{\S\ref{sec:26}}

\noindent\textbf{26.1.}\quad (i) Let $\alpha \in K \cap K'$; then $1, \alpha \in K$ are linearly dependent over
$K'$, hence also over $k$, so $\alpha \in k$. (ii) Assume that $u_1, \ldots, u_n \in K$ are
linearly independent over $k'$, and that $\sum \xi_iu_i = 0$ with
$\xi_i \in k'(K')$; we show that $\xi_i = 0$. Clearing denominators, we can assume
that $\xi_i \in k'[K']$. Choosing a basis $\{\omega_j\}$ of $K'$ over $k$ we can write
$\xi_i = \sum_j c_{ij}\omega_j$ with $c_{ij} \in k'$. Then since
$\sum_{i,j}c_{ij}\omega_ju_i = 0$ we get $\sum_i c_{ij}u_i = 0$, therefore
$c_{ij} = 0$ for all $i, j$.

\noindent\textbf{26.2.}\quad It is enough to show that $K((T))$ and $L^p((T^p))$ are linearly
disjoint over $K^p((T^p))$. Assume that $\omega_1(T), \ldots, \omega_r(T) \in K((T))$ are
linearly independent over $K^p((T^p))$, and that $\sum \varphi_i\omega_i = 0$ with
$\varphi_i \in L^p((T^p))$; we show that $\varphi_i = 0$ for all $i$. Clearing
denominators, we can assume that $\omega_i \in K[T]$ and
$\varphi_i \in L^p[T^p]$. Letting $\{\eta_\lambda\}$ be a basis of $L$ over $K$ we can in a
unique way write $\varphi_i = \sum_\lambda \varphi_{i\lambda}(T^p)$ with
$\varphi_{i\lambda}(T^p) \in K^p[T^p]$. Here $\sum_{\lambda,i}\varphi_{i\lambda}\omega_i$ is in
general an infinite sum, but only a finite number of terms appear in the sum
for the coefficient of some monomial in the $T$'s, so that the sum is
meaningful. Then $\sum_\lambda\big(\sum_i \varphi_{i\lambda}(T^p)\omega_i(T)\big) = 0$, so that
$\sum_i \varphi_{i\lambda}(T^p)\omega_i = 0$ for all $\lambda$, so
$\varphi_{i\lambda} = 0$ for all $i, \lambda$.

% ============ Solutions and hints (Part 10: §28, §29, §30) ============
% 原书 pp.296-297; 原书对 §27 的 Ex.27.1, 27.2、§31-§33 的习题未给解答, 编号按原书跳号.

\section*{\S\ref{sec:28}}

\noindent\textbf{28.1.}\quad Let $N$ be a $B$-module satisfying $m^\nu N = 0$, and $D\colon B \to N$ a
derivation over $A$; then $D$ induces a derivation $\bar{D}\colon B_0 = B/mB \to N/mN$.
If $B_0$ is 0-unramified over $k$ then $\bar{D} = 0$, so that $D(B) \subset mN$.
Proceeding as before gives $D(B) \subset m^2N, \ldots$ so that $D = 0$. The statement
about \'etale is just putting together those for smooth and unramified.

\noindent\textbf{28.2.}\quad Since some $D \in \Der(k)$ with $Da \ne 0$ can be extended to a derivation
of $A$, $a \notin A^p$. If $k'$ were a coefficient field containing $k$ then we would
have to have $a \in k'^p \subset A^p$. Also, $A$ is 0-smooth over $k$ because
$k[X]$ is.

\section*{\S\ref{sec:29}}

\noindent\textbf{29.1.}\quad Suppose that $C = R[[y]]$ with $R$ a DVR; if we let $u$ be a
uniformising element of $R$ then $pR$ is a power of $uR$, so that $pC$ has the
single prime divisor $uC$. However, in fact, in our case
$C/pC = (B/pB)[X]/(X(X + y))$, so that $pC$ has the two prime divisors $(p, x)$
and $(p, x + y)$.

\noindent\textbf{29.2.}\quad $R$ is $pR$-\'etale over $\mathbb{Z}_{(p)}$ (see Ex.~\hyperref[ex:28.1-ex]{28.1}).

\section*{\S\ref{sec:30}}

\noindent\textbf{30.1.}\quad $\varphi$ is the composite of $E_t\colon A \to A[[t]]$ and of the map
$A[[t]] \to A$ obtained by substituting $-x$ for $t$. Since $\varphi(x) = 0$ we have
$xA \subset \Ker \varphi$. For any $a \in A$ we can write $\varphi(a) = a + xb$, so that
$\varphi(\varphi(a)) = \varphi(a)$, and therefore $C \cap xA = (0)$ and $A = C + xA$,
therefore $A = C[x]$. Now $E_t(x) = x + t$, and it is easy to see that this is a
non-zero-divisor of $A[[t]]$, so that $x$ is a non-zero-divisor of $A$. If
$c_\nu x^\nu + c_{\nu+1}x^{\nu+1} + \cdots = 0$ with $c_i \in C$ then dividing by $x^\nu$ we get
$c_\nu \in C \cap xA = (0)$, so that $x$ is analytically independent over $C$.

\noindent\textbf{30.2.}\quad No. If $A$ is 0-smooth over $k'$ then so is the field of fractions $L$
of $A$, so $L/k'$ is separable, hence also $k/k'$, and this is not the case.

\noindent\textbf{30.3.}\quad (1) $\Rightarrow$ (3) is easy using Theorem~\hyperref[thm:28.7]{28.7}. If
$[k : k^p] = \infty$ then there are counter-examples to (3) $\Rightarrow$ (1) (see [C1],
(22.7.7)).


% ============ 参考文献 ============
\chapter*{References}
\addcontentsline{toc}{chapter}{References}
% ============ References (原书 pp.298-315) ============
% 原书文献分两部分: BOOKS (作者缩写键, 如 [M][K][AM][B1-9]) 与 RESEARCH PAPERS
% (按作者字母序, 每位作者内部编号 [1][2]...). 共 435 条 (BOOKS 55 + PAPERS 380),
% 2026-08-08 从原书全部提取, 键经 400dpi 页面图像核验.
% 内部键规则: 书用原键去空格 (如 M, K, AM, G6, N1); 论文用 作者码+编号
% (如 Kap2=Kaplansky 论文2, Kr3=Krull 论文3, Nort8=Northcott 论文8, Coh3=Cohen 论文3).
% 注: 正文中 [B5]/[B6,7]/[B7] 为 Bourbaki 分卷引用 (总条目为 [B1-9]), 保持纯文本不 \\cite.

\section*{Books}
\begin{thebibliography}{XXX}
\bibitem[AB]{AB} M. Auslander and M. Bridger, \textit{Stable module theory}. Memoirs of the AMS 94 (1969).
\bibitem[AM]{AM} M.F. Atiyah and I.G. Macdonald, \textit{Introduction to commutative algebra}. Addison-Wesley, 1969.
\bibitem[An1]{An1} M. Andr\'e, \textit{M\'ethodes simpliciales en alg\`ebre homologique et alg\`ebre commutative}. LNM 32, Springer, 1967.
\bibitem[An2]{An2} M. Andr\'e, \textit{Homologie des alg\`ebres commutatives}. Springer, 1974.
\bibitem[B1-9]{B1-9} N. Bourbaki, \textit{Alg\`ebre commutative}, Chap. 1--Chap. 9. Hermann, 1961--83.
\bibitem[Con]{Con} \textit{Conference on commutative algebra}, Kansas 1972, LNM 311. Springer, 1973.
\bibitem[Diff]{Diff} R. Berger, R. Kiehl, E. Kunz, and H.-J. Nastold, \textit{Differentialrechnung in der analytischen Geometrie}. LNM 38, Springer, 1967.
\bibitem[EG]{EG} E. G. Evans and P. Griffith, \textit{Syzygies}. LMS Lecture Note Series 106, Cambridge University Press, 1985.
\bibitem[F]{F} R. Fossum, \textit{The divisor class group of a Krull domain}. Springer Ergebnisse 74, 1973.
\bibitem[Ful]{Ful} W. Fulton, \textit{Intersection Theory}. Ergebnisse d. Math. 3. Folge Vol. 2, Springer, 1984.
\bibitem[G1]{G1} A. Grothendieck, \textit{\'El\'ements de g\'eom\'etrie alg\'ebrique}, IV. Publ. Math. IHES 20, 1964.
\bibitem[G2]{G2} A. Grothendieck, \textit{\'El\'ements de g\'eom\'etrie alg\'ebrique}, IV. Publ. Math. IHES 24, 1965.
\bibitem[G3]{G3} A. Grothendieck, \textit{\'El\'ements de g\'eom\'etrie alg\'ebrique}, IV. Publ. Math. IHES 28, 1966.
\bibitem[G4]{G4} A. Grothendieck, \textit{\'El\'ements de g\'eom\'etrie alg\'ebrique}, IV. Publ. Math. IHES 32, 1967.
\bibitem[G5]{G5} A. Grothendieck et al., \textit{Cohomologie locale des faisceaux coh\'erents et th\'eor\`emes de Lefschetz locaux et globaux} (SGA 2). North-Holland, 1968.
\bibitem[G6]{G6} A. Grothendieck, \textit{Local cohomology}. LNM 41, Springer, 1967.
\bibitem[GL]{GL} T. H. Gulliksen and G. Levin, \textit{Homology of local rings}. Queen's papers in pure and appl. math., 20, Queen's Univ. Kingston, 1969.
\bibitem[GS]{GS} S. Greco and P. Salmon, \textit{Topics in m-adic topologies}. Springer Ergebnisse 58, 1971.
\bibitem[H]{H} M. Hochster, \textit{Topics in the homological theory of modules over commutative rings}. Regional conference series 24, AMS 1975.
\bibitem[Ha]{Ha} R. Hartshorne, \textit{Algebraic geometry}. Springer, 1977.
\bibitem[HK]{HK} J. Herzog and E. Kunz, \textit{Der kanonische Modul eines Cohen-Macaulay Rings}. LNM 238, Springer, 1971.
\bibitem[HSV]{HSV} M. Hermann, R. Schmidt and W. Vogel, \textit{Theorie der normalen Flachheit}. Teubner Texte zur Math., Leipzig, 1977.
\bibitem[J]{J} N. Jacobson, \textit{Lectures in abstract algebra, Vol. 3. Theory of fields and Galois theory}. Van Nostrand, 1964.
\bibitem[K]{K} I. Kaplansky, \textit{Commutative rings}. Allyn and Bacon, 1970.
\bibitem[Kunz1]{Kunz1} E. Kunz, \textit{Introduction to commutative algebra and algebraic geometry}. Birkh\"auser, 1985.
\bibitem[Kunz2]{Kunz2} E. Kunz, \textit{K\"ahler Differentials}. Vieweg Advanced Lectures in Math., Vieweg \& Sohn, Braunschweig, 1986.
\bibitem[KPR]{KPR} H. Kurke, G. Pfister and M. Roczen, \textit{Henselsche Ringe und algebraische Geometrie}. VEB Deutscher Verlag der Wissensch., Berlin, 1975.
\bibitem[KPPRM]{KPPRM} H. Kurke, G. Pfister, D. Popescu, M. Roczen and T. Mostowski, \textit{Die Approximationseigenschaft lokaler Ringe}. LNM 634, Springer, 1978.
\bibitem[Kr]{Kr} W. Krull, \textit{Idealtheorie}. Springer Ergebnisse, 1935.
\bibitem[La]{La} T.Y. Lam, \textit{Serre's conjecture}. LNM 635, Springer, 1978.
\bibitem[M]{M} H. Matsumura, \textit{Commutative algebra}. Benjamin, 1970 (second edn. 1980).
\bibitem[McA]{McA} S. McAdam, \textit{Asymptotic prime divisors}. LNM 1023, Springer, 1983.
\bibitem[Mac]{Mac} F.S. Macaulay. \textit{Algebraic theory of modular systems}. Cambridge Tracts 19, Cambridge University Press, 1916.
\bibitem[N1]{N1} M. Nagata, \textit{Local rings}. Interscience, 1962.
\bibitem[N2]{N2} M. Nagata, \textit{Lectures on the fourteenth problem of Hilbert}. Tata Inst., 1965.
\bibitem[N3]{N3} M. Nagata, \textit{Polynomial rings and affine spaces}. Regional conference series 37, AMS 1978.
\bibitem[Nor1]{Nor1} D.G. Northcott, \textit{Ideal theory}. Cambridge Tracts 42, Cambridge University Press, 1952.
\bibitem[Nor2]{Nor2} D.G. Northcott, \textit{An introduction to homological algebra}. Cambridge University Press, 1960.
\bibitem[Nor3]{Nor3} D.G. Northcott, \textit{Lessons on rings, modules and multiplicities}. Cambridge University Press, 1968.
\bibitem[Nor4]{Nor4} D.G. Northcott, \textit{Finite free resolutions}. Cambridge Tracts 71, Cambridge University Press, 1976.
\bibitem[R]{R} M. Raynaud, \textit{Anneaux locaux hens\'eliens}. LNM 169, Springer, 1970.
\bibitem[Rob]{Rob} P. Roberts, \textit{Homological invariants of modules over commutative rings}. Presses de l'Univ. de Montr\'eal, 1980.
\bibitem[Rot]{Rot} J. Rotman, \textit{Notes on homological algebra}. Van Nostrand Math Studies 26, 1968.
\bibitem[S1]{S1} P. Samuel, \textit{Alg\`ebre locale}. M\'em. Sci. Math. 123, Gauthier-Villars, 1953.
\bibitem[S2]{S2} P. Samuel, \textit{Lectures on unique factorisation domains}. Tata Inst., 1964.
\bibitem[S3]{S3} P. Samuel, \textit{M\'ethodes d'alg\`ebre abstraite en g\'eom\'etrie alg\'ebrique}. Springer Ergebnisse, 1955.
\bibitem[Sal]{Sal} J. Sally, \textit{Numbers of generators of ideals in local rings}. Lecture Notes in Pure and Appl. Math. 35, Marcel Dekker, 1978.
\bibitem[Sch]{Sch} P. Schenzel, \textit{Dualisierende Komplexe in der lokalen Algebra und Buchsbaum-Ringe}. LNM 907, Springer 1982.
\bibitem[Sh]{Sh} I.R. Shafarevich, \textit{Basic algebraic geometry}. Springer, 1974.
\bibitem[Se]{Se} J.-P. Serre, \textit{Alg\`ebre locale, multiplicit\'es}. LNM 11, Springer, 1965.
\bibitem[St]{St} R.P. Stanley, \textit{Combinatorics and commutative algebra}. Birkh\"auser, 1983.
\bibitem[Su]{Su} S. Suzuki, \textit{Differentials of commutative rings}. Queen's papers in pure and appl. math., 29, Queen's Univ., Kingston, 1971.
\bibitem[SV]{SV} J. St\"uckrad and W. Vogel, \textit{Buchsbaum Rings and Applications}. VEB Deutscher Verlag d. Wissenschaften, Berlin 1986.
\bibitem[V]{V} W. Vasconcelos, \textit{Divisor theory in module categories}. North-Holland Math. Studies 14, North-Holland, 1974.
\bibitem[ZS]{ZS} O. Zariski and P. Samuel, \textit{Commutative algebra}, Vols. I and II. Van Nostrand 1958, 1960 (new edn. Springer, 1975).
\end{thebibliography}

\section*{Research papers}
\begin{thebibliography}{XXX}
\bibitem[1]{Ab1} S. Abhyankar, On the valuations centered in a local domain. Amer. J. Math. 78 (1956), 321--48.
\bibitem[2]{Ab2} S. Abhyankar, Local rings of high embedding dimension. Amer. J. Math. 89 (1967), 1073--7.
\bibitem[1]{AHE1} S. Abhyankar, W. Heinzer and P. Eakin, On the uniqueness of the coefficient ring in a polynomial ring. J. Alg. 23 (1972), 310--42.
\bibitem[1]{Ak1} Y. Akizuki, Einige Bemerkungen \"uber prim\"are Integrit\"atsbereiche mit Teilerkettensatz. Proc. Phys.-Math. Soc. Japan 17 (1935), 327--36.
\bibitem[2]{Ak2} Y. Akizuki, Teilerkettensatz und Vielfachensatz. Proc. Phys.-Math. Soc. Japan 17 (1935), 337--45.
\bibitem[3]{Ak3} Y. Akizuki, Zur Idealtheorie der einartigen Ringbereiche mit dem Teilerkettensatz. Japan. J. Math. 14 (1938), 85--102; II. Japan. J. Math. 14 (1938), 177--87; III. Japan. J. Math. 15 (1939), 1--11.
\bibitem[1]{And1} M. Andr\'e, Localisation de la lissit\'e formelle. manuscripta math. 13 (1974), 297--307.
\bibitem[1]{Ao1} Y. Aoyama, Some basic results on canonical modules. J. Math. Kyoto Univ. 23 (1983), 85--94.
\bibitem[1]{Au1} M. Auslander, Modules over unramified regular local rings. Illinois J. Math. 5 (1961), 631--47.
\bibitem[2]{Au2} M. Auslander, On the purity of the branch locus. Amer. J. Math. 84 (1962), 116--25.
\bibitem[1]{AB1} M. Auslander and D. Buchsbaum, Homological dimension in local rings. Trans. AMS 85 (1957), 390--405.
\bibitem[2]{AB2} M. Auslander and D. Buchsbaum, Codimension and multiplicity. Ann. Math. 68 (1958), 625--57.
\bibitem[3]{AB3} M. Auslander and D. Buchsbaum, Unique factorization in regular local rings. Proc. Natl. Acad. Sci. USA 45 (1959), 733--4.
\bibitem[1]{Av1} L. Avramov, Flat morphisms of complete intersections. Soviet Math. Dokl. 16 (1975), 1413--17.
\bibitem[2]{Av2} L. Avramov, Complete intersections and symmetric algebras. J. Alg. 73 (1981), 248--63.
\bibitem[3]{Av3} L. Avramov, Homology of local flat extensions and complete intersection defects. Math. Ann. 228 (1977), 27--37.
\bibitem[1]{Ba1} H. Bass, On the ubiquity of Gorenstein rings. Math. Z. 82 (1963), 8--28.
\bibitem[1]{Bo1} M. Boratynski, A note on set-theoretical complete intersection ideals. J. Alg. 54 (1978), 1--5.
\bibitem[2]{Bo2} M. Boratynski, Generating ideals up to radical and system of parameters of graded rings. J. Alg. 78 (1982), 20--4.
\bibitem[1]{BI1} A. Brezuleanu and C. Ionescu, On the localization theorems and completion of P-rings. Rev. Roum. Math. Pures Appl. 29 (1984), 371--80.
\bibitem[1]{BR1} A. Brezuleanu and N. Radu, Geometric regularity and formal smoothness. Rend. Accad. Nat. Lincei 51 (1971), 326--7.
\bibitem[2]{BR2} A. Brezuleanu and N. Radu, On the localisation of formal smoothness. Rev. Roum. Math. Pures Appl. 20 (1975), 189--200.
\bibitem[3]{BR3} A. Brezuleanu and N. Radu, Excellent rings and good separation of the module of differentials. Rev. Roum. Math. Pures Appl. 23 (1978), 1455--70.
\bibitem[1]{BRT1} A. Brezuleanu and C. Rotthaus, Eine Bemerkung \"uber Ringe mit geometrisch normalen formalen Fasern. Archiv d. Math. 39 (1982), 19--27.
\bibitem[1]{Br1} M. Brodmann, Rees rings and form rings of almost complete intersections. Nagoya Math. J. 88 (1982), 1--16.
\bibitem[2]{Br2} M. Brodmann, Asymptotic stability of Ass($M/I^nM$). Proc. AMS 74 (1979), 16--18.
\bibitem[1]{BRo1} M. Brodmann and C. Rotthaus, \"Uber den regul\"aren Ort in ausgezeichneten Ringen. Math. Z. 175 (1980), 81--7.
\bibitem[2]{BRo2} M. Brodmann and C. Rotthaus, Local domains with bad sets of formal prime divisors. J. Alg. 75 (1982), 386--94.
\bibitem[1]{Bru1} W. Bruns, `Jede' endliche freie Aufl\"osung ist freie Aufl\"osung eines von drei Elementen erzeugten Ideals. J. Alg. 39 (1976), 429--39.
\bibitem[2]{Bru2} W. Bruns, Zur Konstruktion basischer Elemente. Math. Z. 172 (1980), 63--75.
\bibitem[3]{Bru3} W. Bruns, The Eisenbud-Evans generalized principal ideal theorem and determinantal ideals. Proc. AMS 83 (1981), 19--24.
\bibitem[1]{BE1} D. Buchsbaum and D. Eisenbud, What makes a complex exact? J. Alg. 25 (1973), 259--68.
\bibitem[2]{BE2} D. Buchsbaum and D. Eisenbud, Some structure theorems for finite free resolutions. Adv. in Math. 12 (1974), 84--139.
\bibitem[3]{BE3} D. Buchsbaum and D. Eisenbud, What annihilates a module? J. Alg. 47 (1977), 231--43.
\bibitem[1]{Bu1} L. Burch, On ideals of finite homological dimension in local rings. Proc. Camb. Phil. Soc. 64 (1968), 941--52.
\bibitem[1]{Ca1} A. Caruth, A short proof of the principal ideal theorem. Quart. J. Math. Oxford 31 (1980), 401.
\bibitem[1]{Ch1} S. Chase, Direct product of modules. Trans. AMS 97 (1960), 457--73.
\bibitem[1]{Chv1} C. Chevalley, On the theory of local rings. Ann. Math. 44 (1943), 690--708.
\bibitem[2]{Chv2} C. Chevalley, Intersection of algebraic and algebroid varieties. Trans. AMS 57 (1945), 1--85.
\bibitem[1]{Chi1} L. Chiantini, On form ideals and Artin-Rees condition. manuscripta math. 36 (1981), 125--45.
\bibitem[1]{Chir1} G. Chiriacescu, On the theory of Japanese rings. Rev. Roum. Math. Pures Appl. 27 (1982), 945--8.
\bibitem[1]{Coh1} I.S. Cohen, On the structure and ideal theory of complete local rings. Trans. AMS 59 (1946), 54--106.
\bibitem[2]{Coh2} I.S. Cohen, Commutative rings with restricted minimum condition. Duke Math. J. 17 (1950), 27--42.
\bibitem[3]{Coh3} I.S. Cohen, Lengths of prime ideal chains. Amer. J. Math. 76 (1954), 654--68.
\bibitem[1]{CS1} I.S. Cohen and A. Seidenberg, Prime ideals and integral dependence. Bull. AMS 52 (1946), 252--61.
\bibitem[1]{CN1} C. Cowsik and M.V. Nori, Affine curves in characteristic p are set-theoretic complete intersection. Inv. math. 45 (1978), 111--14.
\bibitem[1]{Da1} E. Davis, Ideals of the principal class, R-sequences and a certain monoidal transformation. Pacific J. Math. 20 (1967), 197--205.
\bibitem[1]{DEP1} C. DeConcini, D. Eisenbud and C. Procesi, Young diagrams and determinantal varieties. Inv. math. 56 (1980), 129--65.
\bibitem[2]{DEP2} C. DeConcini, D. Eisenbud and C. Procesi, Hodge algebras. Ast\'erisque 91 (1982), 1--87.
\bibitem[1]{DP1} C. DeConcini and C. Procesi, A characteristic-free approach to invariant theory. Adv. in Math. 21 (1976), 330--54.
\bibitem[1]{Die1} J. Dieudonn\'e, On regular sequences. Nagoya Math. J. 27 (1966), 355--6.
\bibitem[1]{Du1} S. Dutta, Symbolic powers, intersection multiplicity, and asymptotic behaviour of Tor. J. London Math. Soc. 28 (1983), 261--81.
\bibitem[2]{Du2} S. Dutta, Frobenius and multiplicities. J. Alg. 85 (1983), 424--48.
\bibitem[1]{DHM1} S. Dutta, M. Hochster and J.E. McLaughlin, Modules of finite projective dimension with negative intersection multiplicities. Inv. math. 79 (1985), 253--91.
\bibitem[1]{Ea1} J.A. Eagon, Examples of Cohen-Macaulay rings which are not Gorenstein. Math. Z. 109 (1969), 109--11.
\bibitem[1]{EH1} J.A. Eagon and M. Hochster, R-sequences and indeterminates. Quart. J. Math. Oxford 25 (1974), 61--71.
\bibitem[1]{EN1} J.A. Eagon and D.G. Northcott, Ideals defined by matrices and a certain complex associated with them. Proc. Roy. Soc. London 269 (1962), 188--204.
\bibitem[2]{EN2} J.A. Eagon and D.G. Northcott, Generically acyclic complexes and generically perfect ideals. Proc. Roy. Soc. London 299 (1967), 147--72.
\bibitem[1]{Eak1} P. Eakin, The converse to a well known theorem on Noetherian rings. Math. Ann. 177 (1968), 278--82.
\bibitem[1]{ES1} B. Eckmann and A. Schopf, \"Uber injektiven Moduln. Archiv d. Math. 4 (1953), 75--8.
\bibitem[1]{Ei1} D. Eisenbud, Some directions of recent progress in commutative algebra. Proc. Symposia in Pure Math. 29 (1975), 111--28.
\bibitem[2]{Ei2} D. Eisenbud, Introduction to algebras with straightening laws. In Ring Theory and Algebra III, pp. 243--68. Dekker, 1980.
\bibitem[1]{EB1} D. Eisenbud and D. Buchsbaum, Algebra structures for finite free resolutions, and some structure theorems for ideals of codimension 3. Amer. J. Math. 99 (1977), 447--85.
\bibitem[2]{EB2} D. Eisenbud and D. Buchsbaum, Gorenstein ideals of height 3. In Seminar Eisenbud/Singh/Vogel, Teubner Texte zur Math. Bd. 48, 30--48. Teubner, Leipzig, 1980.
\bibitem[1]{EE1} D. Eisenbud and G. Evans, Every algebraic set in n-space is the intersection of n hypersurfaces. Inv. math. 19 (1973), 107--12.
\bibitem[2]{EE2} D. Eisenbud and G. Evans, Generating modules efficiently: Theorems from algebraic K-theory. J. Alg. 27 (1973), 278--305.
\bibitem[3]{EE3} D. Eisenbud and G. Evans, A generalized principal ideal theorem. Nagoya Math. J. 62 (1976), 41--53.
\bibitem[1]{EG1} G. Evans and P. Griffith, The syzygy problem. Ann. Math. 114 (1981), 323--33.
\bibitem[2]{EG2} G. Evans and P. Griffith, Order ideals of minimal generators. Proc. AMS 86 (1982), 375--8.
\bibitem[1]{Fa1} G. Faltings, Ein einfacher Beweis, da\ss{} geometrische Regularit\"at formale Glattheit impliziert. Arch. d. Math. 30 (1978), 284--5.
\bibitem[2]{Fa2} G. Faltings, \"Uber Annulatoren lokaler Kohomologiegruppen. Arch. d. Math. 30 (1978), 473--6.
\bibitem[3]{Fa3} G. Faltings, Zur Existenz dualisierender Komplexe. Math. Z. 162 (1978), 75--86.
\bibitem[4]{Fa4} G. Faltings, \"Uber Macaulaytizierung. Math. Ann. 238 (1978), 175--92.
\bibitem[1]{Fe1} D. Ferrand, Suite r\'eguli\`ere et intersection compl\`ete. C.R. Acad. Paris 264 (1967), 427--8.
\bibitem[2]{Fe2} D. Ferrand, Trivialisation des modules projectifs. La m\'ethode de Kronecker. J. Pure Appl. Alg. 24 (1982), 261--4.
\bibitem[1]{FR1} D. Ferrand and M. Raynaud, Fibres formelles d'un anneau local noeth\'erien. Ann. Sci. \'Ec. Norm. Sup. 3 (1970), 293--311.
\bibitem[1]{Fi1} M. Fiorentini, On relative regular sequences. J. Alg. 18 (1971), 384--9.
\bibitem[1]{Fit1} H. Fitting, Die Determinantenideale eines Moduls. Jahresberichte Deut. Math. Verein. 46 (1936), 195--228.
\bibitem[1]{Fl1} H. Flenner, Die S\"atze von Bertini f\"ur lokale Ringe. Math. Ann. 229 (1977), 97--111.
\bibitem[1]{Fo1} E. Formanek, Faithful Noetherian modules. Proc. AMS 41 (1973), 381--3.
\bibitem[1]{For1} O. Forster, \"Uber die Anzahl der Erzeugenden eines Ideals in einem Noetherschen Ring. Math. Z. 84 (1964), 80--7.
\bibitem[1]{Fos1} R. Fossum, Duality over Gorenstein rings. Math. Scand. 26 (1970), 165--76.
\bibitem[2]{Fos2} R. Fossum, The structure of indecomposable injective modules. Math. Scand. 36 (1975), 291--312.
\bibitem[3]{Fos3} R. Fossum, Commutative extensions by canonical modules are Gorenstein rings. Proc. AMS 40 (1973), 395--400.
\bibitem[1]{FF1} R. Fossum and H.-B. Foxby, The category of graded modules. Math. Scand. 35 (1974), 288--300.
\bibitem[1]{FFGR1} R. Fossum, H.-B. Foxby, P. Griffith and I. Reiten, Minimal injective resolutions with applications to dualizing modules and Gorenstein modules. Publ. Math. IHES 45 (1976), 193--215.
\bibitem[1]{Fox1} H.-B. Foxby, On the $\mu_i$ in a minimal injective resolution. Math. Scand. 29 (1971), 175--186; II: Math. Scand. 41 (1977), 19--44.
\bibitem[2]{Fox2} H.-B. Foxby, Injective modules under flat base change. Proc. AMS 50 (1975), 23--7.
\bibitem[3]{Fox3} H.-B. Foxby, Isomorphisms between complexes with applications to the homological theory of modules. Math. Scand. 40 (1977), 5--19.
\bibitem[4]{Fox4} H.-B. Foxby, Bounded complexes of flat modules. J. Pure Appl. Alg. 15 (1979), 149--72.
\bibitem[5]{Fox5} H.-B. Foxby, A homological theory of complexes of modules. Preprint series No. 19a, 19b., Copenhagen Univ. Math. Inst., 1981.
\bibitem[1]{FT1} H.-B. Foxby and A. Thorup, Minimal injective resolutions under flat base change. Proc. AMS 67 (1977), 27--31.
\bibitem[1]{Go1} O. Goldman, Hilbert rings and the Hilbert Nullstellensatz. Math. Z. 54 (1951), 136--40.
\bibitem[1]{Got1} S. Goto, On Buchsbaum rings. J. Alg. 67 (1980), 272--9.
\bibitem[2]{Got2} S. Goto, On Buchsbaum rings obtained by gluing. Nagoya Math. J. 83 (1981), 123--35.
\bibitem[3]{Got3} S. Goto, Buchsbaum rings of multiplicity 2. J. Alg. 74 (1982), 494--508.
\bibitem[4]{Got4} S. Goto, On the associated graded rings of parameter ideals in Buchsbaum rings. J. Alg. 85 (1983), 490--534.
\bibitem[5]{Got5} S. Goto, A note on quasi-Buchsbaum rings. Proc. AMS 90 (1984), 511--16.
\bibitem[1]{GS1} S. Goto and Y. Shimoda, On the Rees algebra of Cohen-Macaulay local rings. In R. Draper (ed.), Commutative algebra, analytical methods, pp. 201--231. Dekker, 1982.
\bibitem[1]{GSu1} S. Goto and N. Suzuki, Index of reducibility of parameter ideals in a local ring. J. Alg. 87 (1984), 53--88.
\bibitem[1]{GSW1} S. Goto, N. Suzuki and K. Watanabe, On affine semigroup rings. Japan. J. Math. 2 (1976), 1--12.
\bibitem[1]{GW1} S. Goto and K. Watanabe, On graded rings, I, II. I J. Math. Soc. Japan 30 (1978), 179--213; II (P-graded rings) Tokyo J. Math. 1 (1978), 237--61.
\bibitem[2]{GW2} S. Goto and K. Watanabe, The structure of one-dimensional F-pure rings. J. Alg. 49 (1977), 415--31.
\bibitem[1]{Gr1} S. Greco, Henselization of a ring with respect to an ideal. Trans. AMS 144 (1969), 43--65.
\bibitem[2]{Gr2} S. Greco, Two theorems on excellent rings. Nagoya Math. J. 60 (1976), 139--46.
\bibitem[3]{Gr3} S. Greco, A note on universally catenary rings. Nagoya Math. J. 87 (1982), 95--100.
\bibitem[1]{GM1} S. Greco and M. Marinari, Nagata's criterion and openness of loci for Gorenstein and complete intersections. Math. Z. 160 (1978), 207--16.
\bibitem[1]{GrS1} S. Greco and N. Sankaran, On the separable and algebraic closedness of a Hensel couple in its completion. J. Alg. 39 (1976), 335--48.
\bibitem[1]{GT1} S. Greco and C. Traverso, On seminormal schemes. Compositio Math. 40 (1980), 325--65.
\bibitem[1]{GR1} L. Gruson and M. Raynaud, Crit\`eres de platitude et de projectivit\'e. Inv. math. 13 (1971), 1--89.
\bibitem[1]{Har1} R. Hartshorne, Complete intersections and connectedness. Amer. J. Math. 84 (1962), 497--508.
\bibitem[2]{Har2} R. Hartshorne, A property of A-sequences. Bull. Soc. Math. France 94 (1966), 61--6.
\bibitem[1]{Hap1} L. Harper, On differentiably simple algebras. Trans. AMS 100 (1961), 63--72.
\bibitem[1]{HS1} H. Hasse and F.K. Schmidt, Noch eine Begr\"undung der Theorie der h\"oheren Differential-quotienten in einem algebraischen Funktionenk\"orper einer Unbestimmten. J. reine angew. Math. 177 (1937), 215--37.
\bibitem[1]{He1} W. Heinzer, Some remarks on complete integral closure. J. Australian Math. 9 (1969), 310--14.
\bibitem[2]{He2} W. Heinzer, Quotient overrings of an integral domain. Mathematika 17 (1970), 139--48.
\bibitem[1]{HO1} W. Heinzer and J. Ohm, Locally Noetherian commutative rings. Trans. AMS 158 (1971), 273--84.
\bibitem[2]{HO2} W. Heinzer and J. Ohm, Noetherian intersections of integral domains. Trans. AMS 167 (1972), 291--308.
\bibitem[1]{HI1} M. Herrmann and S. Ikeda, Remarks on lifting of Cohen-Macaulay property. Nagoya Math. J. 92 (1983), 121--32.
\bibitem[2]{HI2} M. Herrmann and S. Ikeda, On the Gorenstein property of Rees algebras. manuscripta math. 59 (1987), 471--90.
\bibitem[1]{HOr1} M. Herrmann and U. Orbanz, Faserdimension von Aufblasungen lokaler Ringe und \"Aquimultiplizit\"at. J. Math. Kyoto Univ. 20 (1980), 651--9.
\bibitem[2]{HOr2} M. Herrmann and U. Orbanz, On equimultiplicity. Math. Proc. Camb. Phil. Soc. 91 (1982), 207--13.
\bibitem[3]{HOr3} M. Herrmann and U. Orbanz, Two notes on flatness. manuscripta math. 40 (1982), 109--33.
\bibitem[1]{HSc1} M. Herrmann and R. Schmidt, Zur Transitivit\"at der normalen Flachheit. Inv. math. 28 (1975), 129--36.
\bibitem[1]{BHe1} B. Herzog, On a relation between the Hilbert functions belonging to a local homomorphism. J. London Math. Soc. 25 (1982), 458--66.
\bibitem[1]{JHe1} J. Herzog, Generators and relations of Abelian semigroups and semigroup rings. manuscripta math. 3 (1970), 175--93.
\bibitem[2]{JHe2} J. Herzog, Ein Cohen-Macaulay-Kriterium mit Anwendungen auf den Konormalenmodul und den Differentialmodul. Math. Z. 163 (1978), 149--62.
\bibitem[3]{JHe3} J. Herzog, Deformationen von Cohen-Macaulay Algebren. J. reine angew. Math. 318 (1980), 83--105.
\bibitem[1]{HK1} J. Herzog and E. Kunz, Die Werthalbgruppe eines lokalen Rings der Dimension 1. S.-B. Heidelberger Akad. Wiss. (1971), 27--67.
\bibitem[1]{HSV1} J. Herzog, A. Simis and W. Vasconcelos, Approximation complexes of blowing-up rings. J. Alg. 74 (1982), 466--93; II. J. Alg. 82 (1983), 53--83.
\bibitem[1]{Hib1} T. Hibi, Every affine graded ring has a Hodge algebra structure. Rend. Sem. Mat. Univers. Politecn. Torino, 44-2 (1986), 27--66.
\bibitem[2]{Hib2} T. Hibi, Distributive lattices, affine semigroup rings and algebras with straightening laws. In Commutative Algebra and Combinatorics, Advanced Studies in Pure Math. 11, 1987, pp. 93--109. North-Holland, Amsterdam.
\bibitem[1]{Hil1} D. Hilbert, \"Uber die Theorie der algebraischen Formen. Math. Ann. 36 (1890), 471--534.
\bibitem[2]{Hil2} D. Hilbert, \"Uber die vollen Invariantensysteme. Math. Ann. 42 (1893), 313--73.
\bibitem[1]{Hir1} H. Hironaka, Resolution of singularities of an algebraic variety over a field of characteristic zero. Ann. Math. 79 (1964), 109--326.
\bibitem[1]{Hoc1} M. Hochster, Rings of invariants of tori, Cohen-Macaulay rings generated by monomials, and polytopes. Ann. Math. 96 (1972), 318--37.
\bibitem[2]{Hoc2} M. Hochster, Non-openness of loci in Noetherian rings. Duke Math. J. 40 (1973), 215--19.
\bibitem[3]{Hoc3} M. Hochster, Contracted ideals from integral extensions of regular rings. Nagoya Math. J. 51 (1973), 25--43.
\bibitem[4]{Hoc4} M. Hochster, Constraints on systems of parameters. In Ring Theory, Proc. Oklahoma Conf., pp. 121--61, Dekker, 1974.
\bibitem[5]{Hoc5} M. Hochster, Cohen-Macaulay rings, combinatorics, and simplicial complexes. In Ring Theory II, Proc. Ohio Conf., pp. 171--223. Dekker, 1977.
\bibitem[6]{Hoc6} M. Hochster, The Zariski-Lipman conjecture for homogeneous complete intersections. Proc. AMS 49 (1975), 261--2.
\bibitem[7]{Hoc7} M. Hochster, Cyclic purity versus purity in excellent Noetherian rings. Trans. AMS 231 (1977), 463--88.
\bibitem[8]{Hoc8} M. Hochster, Some applications of the Frobenius in characteristic 0. Bull. AMS 84 (1978), 886--912.
\bibitem[9]{Hoc9} M. Hochster, Canonical elements in local cohomology modules and the direct summand conjecture. J. Alg. 84 (1983), 503--53.
\bibitem[1]{HE1} M. Hochster and J. Eagon, Cohen-Macaulay rings, invariant theory, and generic perfection of determinantal loci. Amer. J. Math. 93 (1971), 1020--58.
\bibitem[1]{HH1} M. Hochster and C. Huneke, Tightly closed ideals. Bull. AMS 18 (1988), 45--8.
\bibitem[1]{HR1} M. Hochster and L.J. Ratliff, Five theorems on Macaulay rings. Pacific J. Math. 44 (1973), 147--72.
\bibitem[1]{HRo1} M. Hochster and J.L. Roberts, Rings of invariants of reductive groups acting on regular rings are Cohen Macaulay. Adv. in Math. 13 (1974), 115--75.
\bibitem[1]{Hop1} C. Hopkins, Rings with minimum condition for left ideals. Ann. Math. 40 (1939), 712--30.
\bibitem[1]{Hu1} C. Huneke, The theory of d-sequences and powers of ideals. Adv. in Math. 46 (1982), 249--99.
\bibitem[2]{Hu2} C. Huneke, On the symmetric and Rees algebra of an ideal generated by a d-sequence. J. Alg. 62 (1980), 268--75.
\bibitem[3]{Hu3} C. Huneke, A remark concerning multiplicities. Proc. AMS 85 (1982), 331--2.
\bibitem[4]{Hu4} C. Huneke, Linkage and the Koszul homology of ideals. Amer. J. Math. 104 (1982), 1043--62.
\bibitem[5]{Hu5} C. Huneke, Numerical invariants of liaison classes. Inv. Math. 75 (1984), 301--25.
\bibitem[1]{HU1} C. Huneke and B. Ulrich, Divisor class groups and deformations. Amer. J. Math. 107 (1985), 1265--303.
\bibitem[2]{HU2} C. Huneke and B. Ulrich, The structure of linkage. Ann. Math. 126 (1987), 277--334.
\bibitem[1]{Ik1} S. Ikeda, The Cohen-Macaulayness of the Rees algebras of local rings. Nagoya Math. J. 89 (1983), 47--63.
\bibitem[2]{Ik2} S. Ikeda, On the Gorensteinness of Rees algebras over local rings. Nagoya Math. J. 102 (1986).
\bibitem[1]{Is1} F. Ischebeck, Eine Dualit\"at zwischen den Funktoren Ext und Tor. J. Alg. 11 (1969), 510--31.
\bibitem[1]{Iv1} B. Iversen, Amplitude inequalities for complexes. Ann. Sci. \'Ec. Norm. Sup. 10 (1977), 547--58.
\bibitem[1]{Jo1} P. Jothilingam, A note on grade. Nagoya Math. J. 59 (1975), 149--52.
\bibitem[1]{KW1} V. Kac and K. Watanabe, Finite linear groups whose ring of invariants is a complete intersection. Bull. AMS 6 (1982), 221--3.
\bibitem[1]{Kap1} I. Kaplansky, Modules over Dedekind rings and valuation rings. Trans. AMS 72 (1952), 327--40.
\bibitem[2]{Kap2} I. Kaplansky, Projective modules. Ann. Math. 68 (1958), 372--7.
\bibitem[1]{Ke1} M. Kersken, Cousinkomplex und Nennersysteme. Math. Z. 182 (1983), 389--402.
\bibitem[1]{KN1} T. Kimura and H. Niitsuma, On Kunz's conjecture. J. Math. Soc. Japan 34 (1982), 371--8.
\bibitem[1]{Ki1} D. Kirby, Coprimary decomposition of Artinian modules. J. London Math. Soc. 6 (1973), 571--6.
\bibitem[1]{Kiy1} K.H. Kiyek, Anwendung von Ideal-Transformationen. manuscripta math. 34 (1981), 327--53.
\bibitem[1]{Kr1} W. Krull, Primidealketten in allgemeinen Ringbereichen. S.-B. Heidelberger Akad. Wiss. 7 (1928).
\bibitem[2]{Kr2} W. Krull, Ein Satz \"uber prim\"are Integrit\"atsbereiche. Math. Ann. 103 (1930), 450--65.
\bibitem[3]{Kr3} W. Krull, Allgemeine Bewertungstheorie. J. reine angew. Math. 167 (1931), 160--96.
\bibitem[4]{Kr4} W. Krull, \"Uber die Zerlegung der Hauptideale in allgemeinen Ringen. Math. Ann. 105 (1931), 1--14.
\bibitem[5]{Kr5} W. Krull, Beitr\"age zur Arithmetik kommutativer Integrit\"atsbereiche. I. Math. Z. 41 (1936), 545--77; II. Math. Z. 41 (1936), 665--79; III. Math. Z. 42 (1937), 745--66; IV. Math. Z. 42 (1937), 767--73; V. Math. Z. 43 (1938), 768--82; VI. Math. Z. 45 (1939), 1--19; VII. Math. Z. 45 (1940), 319--34; VIII. Math. Z. 48 (1942/43), 533--52.
\bibitem[6]{Kr6} W. Krull, Dimensionstheorie in Stellenringen. J. reine angew. Math. 179 (1938), 204--26.
\bibitem[7]{Kr7} W. Krull, Jacobsonsche Ringe, Hilbertscher Nullstellensatz, Dimensionstheorie. Math. Z. 54 (1951), 354--87.
\bibitem[1]{Ku1} E. Kunz, Characterizations of regular local rings of characteristic p. Amer. J. Math. 91 (1969), 772--84.
\bibitem[2]{Ku2} E. Kunz, On Noetherian rings of characteristic p. Amer. J. Math. 98 (1976), 999--1013.
\bibitem[3]{Ku3} E. Kunz, Almost complete intersections are not Gorenstein rings. J. Alg. 28 (1974), 111--15.
\bibitem[1]{KM1} A. Kustin and M. Miller, Deformation and linkage of Gorenstein algebras. Trans. AMS 284 (1984), 501--23.
\bibitem[2]{KM2} A. Kustin and M. Miller, Tight double linkage of Gorenstein algebras. J. Alg. 95 (1985), 384--97.
\bibitem[1]{KMU1} A. Kustin, M. Miller and B. Ulrich, Linkage theory for algebras with pure resolutions. J. Alg. 102 (1986), 199--228.
\bibitem[1]{LR1} K. Langmann and C. Rotthaus, \"Uber den regul\"aren Ort bei japanischen lokalen Ringen. manuscripta math. 27 (1979), 19--30.
\bibitem[1]{Las1} A. Lascoux, Syzygies des vari\'et\'es d\'eterminantales. Adv. in Math. 30 (1978), 202--37.
\bibitem[1]{LL1} T. Larfeldt and C. Lech, Analytic ramifications and flat couples of local rings. Acta math. 146 (1981), 201--8.
\bibitem[1]{Laz1} D. Lazard, Autour de la platitude. Bull. Soc. Math. France 97 (1969), 81--128.
\bibitem[1]{Le1} C. Lech, Inequalities related to certain couples of local rings. Acta math. 112 (1964), 69--89.
\bibitem[1]{LV1} G. Levin and W. Vasconcelos, Homological dimensions and Macaulay rings. Pacific J. Math. 25 (1968), 315--23.
\bibitem[1]{Li1} S. Lichtenbaum, On the vanishing of Tor in regular local rings. Illinois J. Math. 10 (1966), 220--6.
\bibitem[1]{Lip1} J. Lipman, Free derivation modules on algebraic varieties. Amer. J. Math. 87 (1965), 874--98.
\bibitem[2]{Lip2} J. Lipman, On the Jacobian ideal of the module of differentials. Proc. AMS 21 (1969), 422--6.
\bibitem[1]{MacD1} I.G. MacDonald, Secondary representation of modules over a commutative ring. Symposia Mathematica XI (1973), 23--43.
\bibitem[1]{MS1} I.G. MacDonald and R.Y. Sharp, An elementary proof of the non-vanishing of certain local cohomology modules. Quart. J. Math. Oxford 23 (1972), 197--204.
\bibitem[1]{MacR1} R. MacRae, On the homological dimension of certain ideals. Proc. AMS 14 (1963), 746--50.
\bibitem[2]{MacR2} R. MacRae, On an application of the Fitting invariants. J. Alg. 2 (1965), 153--69.
\bibitem[1]{McA1} S. McAdam, Finite coverings by ideals. In Ring Theory, Proc. Oklahoma Conf., pp. 163--71. Dekker, 1974.
\bibitem[2]{McA2} S. McAdam, Going down. Duke J. Math. 39 (1972), 633--6.
\bibitem[3]{McA3} S. McAdam, Saturated chains in Noetherian rings. Indiana Univ. Math. J. 23 (1974), 719--28.
\bibitem[4]{McA4} S. McAdam, Two applications of asymptotic prime divisors. Proc. AMS 84 (1982), 179--80.
\bibitem[1]{ME1} S. McAdam and P. Eakin, The asymptotic Ass. J. Alg. 61 (1979), 71--81.
\bibitem[1]{Mar1} J. Marot, Sur les anneaux universellement japonais. C. R. Acad. Paris, 277 (1973), 1029--31; C.R. Acad. Paris 278 (1974), 1169--72.
\bibitem[1]{MV1} C. Massaza and P. Valabrega, Sul apertura di luoghi in uno schema localmente noetheriano. Boll. UMI 14-A (1977), 564--74.
\bibitem[1]{Mat1} J. Matijevic, Maximal ideal transforms of Noetherian rings. Proc. AMS 54 (1976), 49--52.
\bibitem[1]{MR1} J. Matijevic and P. Roberts, A conjecture of Nagata on graded Cohen-Macaulay rings. J. Math. Kyoto Univ. 14 (1974), 125--8.
\bibitem[1]{Matl1} E. Matlis, Injective modules over Noetherian rings. Pacific J. Math. 8 (1958), 511--28.
\bibitem[2]{Matl2} E. Matlis, Reflexive domains. J. Alg. 8 (1968), 1--33.
\bibitem[3]{Matl3} E. Matlis, The Koszul complex and duality. Comm. Alg. 1 (1974), 87--144.
\bibitem[4]{Matl4} E. Matlis, The higher properties of R-sequences. J. Alg. 50 (1978), 77--112.
\bibitem[1]{Mats1} H. Matsumura, Formal power series rings over polynomial rings, I. In Number Theory, Algebraic Geometry and Commutative Algebra in honor of Y. Akizuki, pp. 511--20. Kinokuniya, Tokyo, 1973.
\bibitem[2]{Mats2} H. Matsumura, Noetherian rings with many derivations. In Contributions to Algebra (dedicated to E. Kolchin), pp. 279--94. Academic Press, New York, 1977.
\bibitem[3]{Mats3} H. Matsumura, Quasi-coefficient rings. Nagoya Math. J. 68 (1977), 123--30.
\bibitem[4]{Mats4} H. Matsumura, Integrable derivations. Nagoya Math. J. 87 (1982), 227--45.
\bibitem[1]{MK1} N. Mohan Kumar, Complete intersections. J. Math. Kyoto Univ. 17 (1977), 533--9.
\bibitem[2]{MK2} N. Mohan Kumar, On two conjectures about polynomial rings. Inv. math. 46 (1978), 225--36.
\bibitem[1]{Mo1} Y. Mori, On the integral closure of an integral domain, I, II. I Mem. Coll. Sci. Univ. Kyoto 27 (1953), 249--56; II Bull. Kyoto Gakugei Univ. 7 (1955), 19--30.
\bibitem[1]{Mur1} P. Murthy, A note on factorial rings. Arch. d. Math. 15 (1964), 418--21.
\bibitem[1]{Nagr1} K.R. Nagarajan, Groups acting on Noetherian rings. Nieuw Archief voor Wiskunde 16 (1968), 25--29.
\bibitem[1]{Nag1} M. Nagata, On the theory of Henselian rings, I, II, III. I Nagoya Math. J. 5 (1953), 45--57; II Nagoya Math. J. 7 (1954), 1--19; III Mem. Coll. Sci. Univ. Kyoto 32 (1959/60), 93--101.
\bibitem[2]{Nag2} M. Nagata, Basic theorems on general commutative rings. Mem. Coll. Sci. Univ. Kyoto 29 (1955), 59--77.
\bibitem[3]{Nag3} M. Nagata, On the derived normal rings of Noetherian integral domains. Mem. Coll. Sci. Univ. Kyoto 29 (1955), 293--303.
\bibitem[4]{Nag4} M. Nagata, The theory of multiplicity in general local rings. In Proc. Intern. Symp. Tokyo-Nikko 1955, Science Council of Japan, 1956, pp. 191--226.
\bibitem[5]{Nag5} M. Nagata, On the chain problem of prime ideals. Nagoya Math. J. 10 (1956), 51--64.
\bibitem[6]{Nag6} M. Nagata, A Jacobian criterion of simple points. Illinois J. Math. 1 (1957), 427--32.
\bibitem[7]{Nag7} M. Nagata, On the purity of branch loci in regular local rings. Illinois J. Math. 3 (1959), 328--33.
\bibitem[8]{Nag8} M. Nagata, On the closedness of singular loci. Publ. Math. IHES 2 (1959), 29--36.
\bibitem[9]{Nag9} M. Nagata, A type of subrings of a Noetherian ring. J. Math. Kyoto Univ. 8 (1968), 465--7.
\bibitem[10]{Nag10} M. Nagata, Flatness of an extension of a commutative ring. J. Math. Kyoto Univ. 9 (1969), 439--48.
\bibitem[1]{Nak1} Y. Nakai, Theory of differentials in commutative rings. J. Math. Soc. Japan 13 (1961), 68--84.
\bibitem[2]{Nak2} Y. Nakai, High order derivations I. Osaka J. Math. 7 (1970), 1--27.
\bibitem[3]{Nak3} Y. Nakai, On a ring with plenty of high order derivations. J. Math. Kyoto Univ. 13 (1972), 159--64.
\bibitem[4]{Nak4} Y. Nakai, On locally finite iterative higher derivations. Osaka J. Math. 15 (1978), 655--62.
\bibitem[1]{NKI1} Y. Nakai, K. Kosaki and Y. Ishibashi, High order derivations II. J. Sci. Hiroshima Univ. Ser. A-I 34 (1970), 17--27.
\bibitem[1]{No1} E. Noether, Idealtheorie in Ringbereichen. Math. Ann. 83 (1921), 24--66.
\bibitem[2]{No2} E. Noether, Abstrakter Aufbau der Idealtheorie in algebraischen Zahl- und Funktionenk\"orpern. Math. Ann. 96 (1926), 26--61.
\bibitem[1]{Nish1} J. Nishimura, Note on Krull domain. J. Math. Kyoto Univ. 15 (1975), 397--400.
\bibitem[2]{Nish2} J. Nishimura, Note on integral closures of a Noetherian integral domain. J. Math. Kyoto Univ. 16 (1976), 117--22.
\bibitem[3]{Nish3} J. Nishimura, On ideal-adic completion of Noetherian rings. J. Math. Kyoto Univ. 21 (1981), 153--69.
\bibitem[1]{Nom1} M. Nomura, Formal power series rings over polynomial rings, II. In Number Theory, Algebraic Geometry and Commutative Algebra in honor of Y. Akizuki, pp. 521--8. Kinokuniya, Tokyo, 1973.
\bibitem[1]{Nort1} D.G. Northcott, On unmixed ideals in regular local rings. Proc. London Math. Soc. 3 (1953), 20--8.
\bibitem[2]{Nort2} D.G. Northcott, A note on classical ideal theory. Proc. Camb. Phil. Soc. 51 (1955), 766--7.
\bibitem[3]{Nort3} D.G. Northcott, On homogeneous ideals. Proc. Glasgow Math. Assoc. 2 (1955), 105--11.
\bibitem[4]{Nort4} D.G. Northcott, On irreducible ideals in local rings. J. London Math. Soc. 32 (1957), 82--8.
\bibitem[5]{Nort5} D.G. Northcott, Syzygies and specializations. Proc. London Math. Soc. 15 (1965), 1--25.
\bibitem[6]{Nort6} D.G. Northcott, Hilbert functions and the Koszul complex. Bull. London Math. Soc. 2 (1970), 69--72.
\bibitem[7]{Nort7} D.G. Northcott, Grade sensitivity and generic perfection. Proc. London Math. Soc. 20 (1970), 597--618.
\bibitem[8]{Nort8} D.G. Northcott, Injective envelopes and inverse polynomials. J. London Math. Soc. 8 (1974), 290--296.
\bibitem[9]{Nort9} D.G. Northcott, Generalized Koszul complexes and Artinian modules. Quart. J. Math. Oxford 23 (1972), 289--297.
\bibitem[1]{NR1} D.G. Northcott and D. Rees, Reductions of ideals in local rings. Proc. Camb. Phil. Soc. 50 (1954), 145--58.
\bibitem[2]{NR2} D.G. Northcott and D. Rees, A note on reductions of ideals with an application to the generalized Hilbert function. Proc. Camb. Phil. Soc. 50 (1954), 353--9.
\bibitem[3]{NR3} D.G. Northcott and D. Rees, Principal systems. Quart. J. Math. Oxford 8 (1957), 119--27.
\bibitem[4]{NR4} D.G. Northcott and D. Rees, Extensions and simplifications of the theory of regular local rings. J. London Math. Soc. 32 (1957), 367--74.
\bibitem[1]{OC1} L. O'Carroll, Generalized fractions, determinantal maps, and top cohomology modules. J. Pure Appl. Alg. 32 (1984), 59--70.
\bibitem[2]{OC2} L. O'Carroll, On the generalized fractions of Sharp and Zakeri. J. London Math. Soc. 28 (1983), 417--27.
\bibitem[1]{Og1} T. Ogoma, Non-catenary pseudo-geometric normal rings. Japan. J. Math. 6 (1980), 147--63.
\bibitem[2]{Og2} T. Ogoma, Cohen-Macaulay factorial domain is not necessarily Gorenstein. Mem. Fac. Sci. Kochi Univ. (Math.) 3 (1982), 65--74.
\bibitem[3]{Og3} T. Ogoma, Some examples of rings with curious formal fibres. Mem. Fac. Sci. Kochi Univ. Math. 1 (1980), 17--22.
\bibitem[4]{Og4} T. Ogoma, Existence of dualizing complexes. J. Math. Kyoto Univ. 24 (1984), 27--48.
\bibitem[1]{Oh1} J. Ohm, Semi-valuations and groups of divisibility. Can. J. Math. 21 (1969), 576--91.
\bibitem[1]{Or1} U. Orbanz, H\"ohere Derivationen und Regularit\"at. J. reine angew. Math. 262/263 (1973), 194--204.
\bibitem[1]{Pe1} C. Peskine, Une g\'en\'eralisation du `main theorem' de Zariski. Bull. Sci. Math. 90 (1966), 119--27.
\bibitem[1]{PS1} C. Peskine and L. Szpiro, Dimension projective finie et cohomologie locale. Publ. Math. IHES 42 (1973), 47--119.
\bibitem[2]{PS2} C. Peskine and L. Szpiro, Liaison des vari\'et\'es alg\'ebriques. I. Inv. math. 26 (1974), 271--302.
\bibitem[3]{PS3} C. Peskine and L. Szpiro, Syzygies et multiplicit\'es. C. R. Acad. Sci. Paris, Ser. A-B. 278 (1974), 1421--4.
\bibitem[1]{Pf1} G. Pfister, Schlechte Henselsche Ringe. Bull. Acad. Polon. Sci. 25 (1977), 1083--8.
\bibitem[1]{PP1} G. Pfister and D. Popescu, Die strenge Approximationseigenschaft lokaler Ringe. Inv. math. 30 (1975), 145--74.
\bibitem[2]{PP2} G. Pfister and D. Popescu, Die Approximationseigenschaft von Primidealen. Bull. Acad. Polon. Sci. 27 (1979), 771--9.
\bibitem[1]{Po1} D. Popescu, General N\'eron desingularization. Nagoya Math. J. 100 (1985), 97--126.
\bibitem[2]{Po2} D. Popescu, General N\'eron desingularization and approximation. Nagoya Math. J. 104 (1986), 85--115.
\bibitem[1]{Qu1} J. Querre, Sur un th\'eor\`eme de Mori-Nagata. C.R. Acad. Paris 285 (1977), 323--4.
\bibitem[2]{Qu2} J. Querre, Sur une propri\'et\'e des anneaux de Krull. Bull. Sci. Math. France, 95 (1971), 341--54.
\bibitem[1]{Qui1} D. Quillen, Projective modules over polynomial rings. Inv. math. 36 (1976), 167--71.
\bibitem[1]{Rab1} S. Rabinowitch, Zum Hilbertschen Nullstellensatz. Math. Ann. 102 (1929), 520.
\bibitem[1]{Ra1} N. Radu, Un crit\`ere diff\'erentiel de lissit\'e formelle. C.R. Acad. Paris 272 (1971), 1166--8.
\bibitem[2]{Ra2} N. Radu, Une caract\'erisation des alg\`ebres noeth\'eriennes r\'eguli\`eres sur un corps de caract\'eristique z\'ero. C.R. Acad. Paris 270 (1970), 851--3.
\bibitem[3]{Ra3} N. Radu, Sur les alg\`ebres dont le module des diff\'erentielles est plat. Rev. Roum. Math. Pures et Appl. 21 (1976), 933--9.
\bibitem[4]{Ra4} N. Radu, Sur un crit\`ere de lissit\'e formelle. Bull. Math. Soc. Sci. Roum. 21 (1977), 133--5.
\bibitem[5]{Ra5} N. Radu, Sur la structure des alg\`ebres locales noeth\'eriennes formellement lisses. Analele Univ. Bucure\c{s}ti, Mathematics, 29 (1980), 81--4.
\bibitem[1]{Rat1} L.J. Ratliff, On quasi-unmixed local domains, the altitude formula, and the chain condition for prime ideals. I Amer. J. Math. 91 (1969), 508--28; II Amer. J. Math. 92 (1970), 99--144.
\bibitem[2]{Rat2} L.J. Ratliff, Characterizations of catenary rings. Amer. J. Math. 93 (1971), 1070--108.
\bibitem[3]{Rat3} L.J. Ratliff, Catenary rings and the altitude formula. Amer. J. Math. 94 (1972), 458--66.
\bibitem[1]{Re1} D. Rees, Valuations associated with a local ring. I Proc. London Math. Soc. 5 (1955), 107--28; II J. London Math. Soc. 31 (1956), 228--35.
\bibitem[2]{Re2} D. Rees, Valuations associated with ideals. I Proc. London Math. Soc. 6 (1956), 161--74; II J. London Math. Soc. 31 (1956), 221--9.
\bibitem[3]{Re3} D. Rees, Two classical theorems of ideal theory. Proc. Camb. Phil. Soc. 52 (1956), 155--7.
\bibitem[4]{Re4} D. Rees, A theorem of homological algebra. Proc. Camb. Phil. Soc. 52 (1956), 605--10.
\bibitem[5]{Re5} D. Rees, The grade of an ideal or module. Proc. Camb. Phil. Soc. 53 (1957), 28--42.
\bibitem[6]{Re6} D. Rees, A note on form rings and ideals. Mathematika 4 (1957), 51--60.
\bibitem[7]{Re7} D. Rees, A note on general ideals. J. London Math. Soc. 32 (1957), 181--6.
\bibitem[8]{Re8} D. Rees, On a problem of Zariski. Illinois J. Math. 2 (1958), 145--59.
\bibitem[9]{Re9} D. Rees, A note on analytically unramified local rings. J. London Math. Soc. 36 (1961), 24--8.
\bibitem[10]{Re10} D. Rees, Rings associated with ideals and analytic spread. Proc. Camb. Phil. Soc. 89 (1981), 423--32.
\bibitem[1]{RS1} D. Rees and R.Y. Sharp, On a theorem of B. Teissier on multiplicities of ideals in local rings. J. London Math. Soc. 18 (1978), 449--63.
\bibitem[1]{Rei1} G.A. Reisner, Cohen-Macaulay quotients of polynomial rings. Adv. in Math. 21 (1976), 30--49.
\bibitem[1]{RM1} G. Restuccia and H. Matsumura, Integrable derivations in rings of unequal characteristic. Nagoya Math. J. 93 (1984), 173--8.
\bibitem[1]{RV1} L. Robbiano and G. Valla, On normal flatness and normal torsion-freeness. J. Alg. 43 (1976), 552--60.
\bibitem[1]{Rob1} P. Roberts, Two applications of dualizing complexes over local rings. Ann. Sci. \'Ec. Norm. Sup. 9 (1976), 103--6.
\bibitem[2]{Rob2} P. Roberts, The vanishing of intersection multiplicities of perfect complexes. Bull. AMS 13 (1985), 127--30.
\bibitem[3]{Rob3} P. Roberts, Le th\'eor\`eme d'intersection. C. R. Acad. Sci. Paris, 304 (1987), 177--80.
\bibitem[1]{Roth1} C. Rotthaus, Nicht ausgezeichnete, universell japanische Ringe. Math. Z. 152 (1977), 107--25.
\bibitem[2]{Roth2} C. Rotthaus, Universell japanische Ringe mit nicht offenem regul\"arem Ort. Nagoya Math. J. 74 (1979), 123--35.
\bibitem[3]{Roth3} C. Rotthaus, Komplettierung semilokaler quasiausgezeichneter Ringe. Nagoya Math. J. 76 (1979), 173--80.
\bibitem[4]{Roth4} C. Rotthaus, Zur Komplettierung ausgezeichneter Ringe. Math. Ann. 253 (1980), 213--26.
\bibitem[5]{Roth5} C. Rotthaus, Potenzreihenerweiterung und formale Fasern in lokalen Ringen mit Approximationseigenschaft. manuscripta math. 42 (1983), 53--65.
\bibitem[1]{Sally1} J. Sally, On the number of generators of powers of an ideal. Proc. AMS 53 (1975), 24--6.
\bibitem[2]{Sally2} J. Sally, Bounds for numbers of generators of Cohen-Macaulay ideals. Pacific J. Math. 63 (1976), 517--20.
\bibitem[3]{Sally3} J. Sally, On the associated graded ring of a local Cohen-Macaulay ring. J. Math. Kyoto Univ. 17 (1977), 19--21.
\bibitem[4]{Sally4} J. Sally, Cohen-Macaulay local rings of maximal embedding dimension. J. Alg. 56 (1979), 168--83.
\bibitem[5]{Sally5} J. Sally, Reductions, local cohomology and Hilbert functions of local rings. In Commutative Algebra, Durham 1981, pp. 231--41; Cambridge University Press 1982.
\bibitem[1]{Sam1} P. Samuel, La notion de multiplicit\'e en alg\`ebre et en g\'eom\'etrie alg\'ebrique. J. math. pure et appl. 30 (1951), 159--274.
\bibitem[2]{Sam2} P. Samuel, On unique factorization domains. Illinois J. Math. 5 (1961), 1--17.
\bibitem[3]{Sam3} P. Samuel, Sur les anneaux factoriels. Bull. Soc. Math. France 89 (1961), 155--73.
\bibitem[1]{SS1} G. Scheja and U. Storch, Differentielle Eigenschaften der Lokalisierungen analytischer Algebren. Math. Ann. 197 (1972), 137--70.
\bibitem[2]{SS2} G. Scheja and U. Storch, \"Uber Spurfunktionen bei vollst\"andigen Durchschnitten. J. reine angew. Math. 278/279 (1975), 174--90.
\bibitem[1]{Sche1} P. Schenzel, Applications of dualizing complexes to Buchsbaum rings. Adv. in Math. 44 (1982), 61--77.
\bibitem[2]{Sche2} P. Schenzel, Multiplizit\"aten in verallgemeinerten Cohen-Macaulay-Moduln. Math. Nachr. 88 (1979), 295--306.
\bibitem[1]{STC1} P. Schenzel, Trung, Ngo Viet and Cuong, Nguyen Tu, Verallgemeinerte Cohen-Macaulay-Moduln. Math. Nachr. 85 (1978), 57--73.
\bibitem[1]{Sd1} A. Seidenberg, Derivations and integral closure. Pacific J. Math. 16 (1966), 167--73.
\bibitem[1]{Ser1} J.-P. Serre, Sur la dimension homologique des anneaux et des modules noeth\'eriens. Proc. Intern. Symp., Tokyo-Nikko, 1955, Science Council of Japan, 1956, 175--89.
\bibitem[1]{Sey1} H. Seydi, Anneaux hens\'eliens et conditions de cha\^ines. I Bull. Soc. Math. France 98 (1970), 9--31; II C.R. Acad. Paris 270 (1970), 696--8; III Bull. Soc. Math. France 98 (1970), 329--36.
\bibitem[2]{Sey2} H. Seydi, Sur deux th\'eor\`emes d'alg\`ebre commutative. C.R. Acad. Paris 271 (1970), 1105--8.
\bibitem[3]{Sey3} H. Seydi, Sur deux th\'eor\`emes d'alg\`ebre commutative. C.R. Acad. Paris 271 (1970), 1169--72.
\bibitem[4]{Sey4} H. Seydi, Un th\'eor\`eme de descente effective universelle et une application. C.R. Acad. Paris 270 (1970), 801--3.
\bibitem[1]{Shar1} R.Y. Sharp, The Cousin complex for a module over a commutative Noetherian ring. Math. Z. 112 (1969), 340--56.
\bibitem[2]{Shar2} R.Y. Sharp, Gorenstein modules. Math. Z. 115 (1970), 117--39.
\bibitem[3]{Shar3} R.Y. Sharp, Local cohomology theory in commutative algebra. Quart. J. Math. Oxford 21 (1970), 425--34.
\bibitem[4]{Shar4} R.Y. Sharp, Dualizing complexes for commutative Noetherian rings. Math. Proc. Camb. Phil. Soc. 78 (1975), 369--86.
\bibitem[5]{Shar5} R.Y. Sharp, Acceptable rings and homomorphic images of Gorenstein rings. J. Alg. 44 (1977), 246--61.
\bibitem[6]{Shar6} R.Y. Sharp, A commutative Noetherian ring which possesses a dualizing complex is acceptable. Math. Proc. Camb. Phil. Soc. 82 (1977), 197--213.
\bibitem[7]{Shar7} R.Y. Sharp, On the attached prime ideals of certain Artinian local cohomology modules. Proc. Edinburgh Math. Soc. 24 (2) (1981), 9--14.
\bibitem[8]{Shar8} R.Y. Sharp, Secondary representations for injective modules over commutative Noetherian rings. Proc. Edinburgh Math. Soc. 20 (2) (1976), 143--51.
\bibitem[1]{SharV1} R.Y. Sharp and P. V\'amos, The dimension of the tensor product of a finite number of field extensions. J. Pure Appl. Alg. 10 (1977), 249--52.
\bibitem[1]{SZ1} R.Y. Sharp and H. Zakeri, Modules of generalized fractions. Mathematika 29 (1982), 32--41.
\bibitem[2]{SZ2} R.Y. Sharp and H. Zakeri, Modules of generalized fractions and balanced big Cohen-Macaulay modules. In R.Y. Sharp (ed.), Commutative Algebra, Durham 1981, pp. 61--82; Cambridge University Press, 1982.
\bibitem[1]{Shi1} Y. Shimoda, A note on Rees algebras of two-dimensional local domains. J. Math. Kyoto Univ. 19 (1979), 327--33.
\bibitem[1]{Si1} B. Singh, On two numerical criteria for normal flatness in codimension one. Math. Nachr. 83 (1978), 241--6.
\bibitem[2]{Si2} B. Singh, Maximally differential prime ideals in a complete local ring. J. Alg. 82 (1983), 331--9.
\bibitem[1]{Stan1} R. Stanley, Relative invariants of finite groups generated by pseudoreflections. J. Alg. 49 (1977), 134--48.
\bibitem[2]{Stan2} R. Stanley, Hilbert functions of graded algebras. Adv. in Math. 28 (1978), 57--83.
\bibitem[1]{Sto1} U. Storch, Bemerkung zu einem Satz von M. Kneser. Arch. d. Math. 23 (1972), 403--4.
\bibitem[1]{Stu1} J. St\"uckrad, \"Uber die kohomologische Charakterisierung von Buchsbaum-Moduln. Math. Nachr. 95 (1980), 265--72.
\bibitem[1]{StV1} J. St\"uckrad and W. Vogel, Toward a theory of Buchsbaum singularities. Amer. J. Math. 100 (1978), 121--46.
\bibitem[1]{SuN1} N. Suzuki, On the generalized local cohomology and its duality. J. Math. Kyoto Univ. 18 (1978), 71--85.
\bibitem[2]{SuN2} N. Suzuki, On the Koszul complex generated by a system of parameters for a Buchsbaum module. Bull. Dept. General Education of Shizuoka College of Pharmacy 8 (1979), 27--35.
\bibitem[3]{SuN3} N. Suzuki, Canonical duality for Buchsbaum modules -- an application of Goto's lemma on Buchsbaum modules. Bull. Dept. General Education of Shizuoka College of Pharmacy 13 (1984), 47--60.
\bibitem[4]{SuN4} N. Suzuki, On a basic theorem for quasi-Buchsbaum modules. Bull. Dept. General Education of Shizuoka College of Pharmacy 11 (1982), 33--40.
\bibitem[1]{SuS1} S. Suzuki, Some results on Hausdorff m-adic modules and m-adic differentials. J. Math. Kyoto Univ. 2 (1963), 157--82.
\bibitem[2]{SuS2} S. Suzuki, On torsion of the module of differentials of a locality which is a complete intersection. J. Math. Kyoto Univ. 4 (1966), 471--5.
\bibitem[3]{SuS3} S. Suzuki, On N\"agler's numbers of discrete valuation rings. J. Math. Kyoto Univ. 11 (1971), 373--5.
\bibitem[1]{Sw1} R. Swan, The number of generators of a module. Math. Z. 102 (1967), 318--22.
\bibitem[2]{Sw2} R. Swan, On seminormality. J. Alg. 67 (1980), 210--29.
\bibitem[1]{Ta1} H. Tanimoto, Some characterisations of smoothness. J. Math. Kyoto Univ. 23 (1983), 695--706.
\bibitem[2]{Ta2} H. Tanimoto, Smoothness of Noetherian rings. Nagoya Math. J. 95 (1984), 163--79.
\bibitem[1]{Tat1} J. Tate, Homology of Noetherian rings and local rings. Illinois J. Math. 1 (1957), 246--61.
\bibitem[1]{Te1} B. Teissier, Cycles \'evanescents, sections planes et conditions de Whitney. Ast\'erisque 7/8 (1973), 285--362.
\bibitem[2]{Te2} B. Teissier, Sur une \'egalit\'e \`a la Minkowski pour les multiplicit\'es. Ann. Math. 106 (1977), 38--44.
\bibitem[3]{Te3} B. Teissier, On a Minkowski-type inequality for multiplicities -- II. In C.P. Ramanujam -- a tribute, pp. 347--61. Tata Inst. Studies in Math. 8, Springer, 1978.
\bibitem[1]{Tr1} Trung, Ngo Viet, On the symbolic powers of determinantal ideals. J. Alg. 58 (1979), 361--9.
\bibitem[2]{Tr2} Trung, Ngo Viet, A class of imperfect prime ideals having the equality of ordinary and symbolic powers. J. Math. Kyoto Univ. 21 (1981), 239--50.
\bibitem[3]{Tr3} Trung, Ngo Viet, Absolutely superficial sequences. Math. Proc. Camb. Phil. Soc. 93 (1983), 35--47.
\bibitem[1]{Ul1} B. Ulrich, Gorenstein rings as specializations of unique factorization domains. J. Alg. 86 (1984), 129--40.
\bibitem[2]{Ul2} B. Ulrich, Gorenstein rings and modules with high number of generators. Math. Z. 188 (1984), 23--32.
\bibitem[3]{Ul3} B. Ulrich, Liaison and deformation. J. Pure Appl. Alg. 39 (1986), 165--75.
\bibitem[4]{Ul4} B. Ulrich, Rings of invariants and linkage of determinantal ideals. Math. Ann. 274 (1986), 1--17.
\bibitem[1]{Va1} P. Valabrega, On the excellent property for power series rings over polynomial rings. J. Math. Kyoto Univ. 15 (1975), 387--95.
\bibitem[2]{Va2} P. Valabrega, A few theorems on completion of excellent rings. Nagoya Math. J. 61 (1976), 127--33.
\bibitem[3]{Va3} P. Valabrega, Formal fibres and openness of loci. J. Math. Kyoto Univ. 18 (1978), 199--208.
\bibitem[1]{Val1} G. Valla, Certain graded algebras are always Cohen-Macaulay. J. Alg. 42 (1976), 537--48.
\bibitem[2]{Val2} G. Valla, On the symmetric and Rees algebras of an ideal. manuscripta math. 30 (1980), 239--55.
\bibitem[1]{VV1} P. Valabrega and G. Valla, Form rings and regular sequences. Nagoya Math. J. 72 (1978), 93--101.
\bibitem[1]{Ve1} U. Vetter, Generische determinantielle Singularit\"aten: homologische Eigenschaften des Derivationenmoduls. manuscripta math. 45 (1984), 161--91.
\bibitem[1]{Vas1} W. Vasconcelos, Ideals defined by R-sequences. J. Alg. 6 (1967), 309--16.
\bibitem[2]{Vas2} W. Vasconcelos, On finitely generated flat modules. Trans. AMS 138 (1969), 505--12.
\bibitem[1]{Vog1} W. Vogel, \"Uber eine Vermutung von D.A. Buchsbaum. J. Alg. 25 (1973), 106--12.
\bibitem[2]{Vog2} W. Vogel, A non-zero-divisor characterization of Buchsbaum modules. Michigan Math. J. 28 (1981), 147--55.
\bibitem[1]{JW1} J. Watanabe, Some remarks on Cohen-Macaulay rings with many zero-divisors. J. Alg. 39 (1976), 1--14.
\bibitem[2]{JW2} J. Watanabe, A note on Gorenstein rings of embedding codimension three. Nagoya Math. J. 50 (1973), 227--32.
\bibitem[3]{JW3} J. Watanabe, m-full ideals. Nagoya Math. J. 106 (1987), 101--11.
\bibitem[4]{JW4} J. Watanabe, The Dilworth number of artinian rings and finite posets with rank function. In Commutative Algebra and Combinatorics, Advanced Studies in Pure Math. 11 (1987), pp. 303--12. North-Holland, Amsterdam.
\bibitem[1]{KWat1} K. Watanabe, Some examples of one dimensional Gorenstein domains. Nagoya Math. J. 49 (1973), 101--9.
\bibitem[2]{KWat2} K. Watanabe, Certain invariant subrings are Gorenstein, I, II. I Osaka J. Math. 11 (1974), 1--8; II Osaka J. Math. 11 (1974), 379--88.
\bibitem[3]{KWat3} K. Watanabe, Invariant subrings which are complete intersections, I. Invariant subrings of finite Abelian groups. Nagoya Math. J. 77 (1980), 89--98.
\bibitem[1]{WITO1} K. Watanabe, T. Ishikawa, S. Tachibana and K. Otsuka, On tensor products of Gorenstein rings. J. Math. Kyoto Univ. 15 (1975), 387--95.
\bibitem[1]{WR1} K. Watanabe and D. Rotillon, Invariant subrings of $C[X,Y,Z]$ which are complete intersections. manuscripta math. 39 (1982), 339--57.
\end{thebibliography}


% 索引
\printindex
\addcontentsline{toc}{chapter}{Index}

\end{document}
